% Options for packages loaded elsewhere
\PassOptionsToPackage{unicode}{hyperref}
\PassOptionsToPackage{hyphens}{url}
\PassOptionsToPackage{dvipsnames,svgnames,x11names}{xcolor}
%
\documentclass[
  a4paper,
  13pt]{extreport}
\usepackage{amsmath,amssymb}
\usepackage{setspace}
\usepackage{iftex}
\ifPDFTeX
  \usepackage[T1]{fontenc}
  \usepackage[utf8]{inputenc}
  \usepackage{textcomp} % provide euro and other symbols
\else % if luatex or xetex
  \usepackage{unicode-math} % this also loads fontspec
  \defaultfontfeatures{Scale=MatchLowercase}
  \defaultfontfeatures[\rmfamily]{Ligatures=TeX,Scale=1}
\fi
\usepackage{lmodern}
\ifPDFTeX\else
  % xetex/luatex font selection
  \setmainfont[]{TeX Gyre Termes}
\fi
% Use upquote if available, for straight quotes in verbatim environments
\IfFileExists{upquote.sty}{\usepackage{upquote}}{}
\IfFileExists{microtype.sty}{% use microtype if available
  \usepackage[]{microtype}
  \UseMicrotypeSet[protrusion]{basicmath} % disable protrusion for tt fonts
}{}
\usepackage{xcolor}
\usepackage[left=3cm,right=2cm,top=2cm,bottom=2cm,headsep=1cm,footskip=1cm]{geometry}
\usepackage{color}
\usepackage{fancyvrb}
\newcommand{\VerbBar}{|}
\newcommand{\VERB}{\Verb[commandchars=\\\{\}]}
\DefineVerbatimEnvironment{Highlighting}{Verbatim}{commandchars=\\\{\}}
% Add ',fontsize=\small' for more characters per line
\newenvironment{Shaded}{}{}
\newcommand{\AlertTok}[1]{\textcolor[rgb]{1.00,0.00,0.00}{\textbf{#1}}}
\newcommand{\AnnotationTok}[1]{\textcolor[rgb]{0.38,0.63,0.69}{\textbf{\textit{#1}}}}
\newcommand{\AttributeTok}[1]{\textcolor[rgb]{0.49,0.56,0.16}{#1}}
\newcommand{\BaseNTok}[1]{\textcolor[rgb]{0.25,0.63,0.44}{#1}}
\newcommand{\BuiltInTok}[1]{\textcolor[rgb]{0.00,0.50,0.00}{#1}}
\newcommand{\CharTok}[1]{\textcolor[rgb]{0.25,0.44,0.63}{#1}}
\newcommand{\CommentTok}[1]{\textcolor[rgb]{0.38,0.63,0.69}{\textit{#1}}}
\newcommand{\CommentVarTok}[1]{\textcolor[rgb]{0.38,0.63,0.69}{\textbf{\textit{#1}}}}
\newcommand{\ConstantTok}[1]{\textcolor[rgb]{0.53,0.00,0.00}{#1}}
\newcommand{\ControlFlowTok}[1]{\textcolor[rgb]{0.00,0.44,0.13}{\textbf{#1}}}
\newcommand{\DataTypeTok}[1]{\textcolor[rgb]{0.56,0.13,0.00}{#1}}
\newcommand{\DecValTok}[1]{\textcolor[rgb]{0.25,0.63,0.44}{#1}}
\newcommand{\DocumentationTok}[1]{\textcolor[rgb]{0.73,0.13,0.13}{\textit{#1}}}
\newcommand{\ErrorTok}[1]{\textcolor[rgb]{1.00,0.00,0.00}{\textbf{#1}}}
\newcommand{\ExtensionTok}[1]{#1}
\newcommand{\FloatTok}[1]{\textcolor[rgb]{0.25,0.63,0.44}{#1}}
\newcommand{\FunctionTok}[1]{\textcolor[rgb]{0.02,0.16,0.49}{#1}}
\newcommand{\ImportTok}[1]{\textcolor[rgb]{0.00,0.50,0.00}{\textbf{#1}}}
\newcommand{\InformationTok}[1]{\textcolor[rgb]{0.38,0.63,0.69}{\textbf{\textit{#1}}}}
\newcommand{\KeywordTok}[1]{\textcolor[rgb]{0.00,0.44,0.13}{\textbf{#1}}}
\newcommand{\NormalTok}[1]{#1}
\newcommand{\OperatorTok}[1]{\textcolor[rgb]{0.40,0.40,0.40}{#1}}
\newcommand{\OtherTok}[1]{\textcolor[rgb]{0.00,0.44,0.13}{#1}}
\newcommand{\PreprocessorTok}[1]{\textcolor[rgb]{0.74,0.48,0.00}{#1}}
\newcommand{\RegionMarkerTok}[1]{#1}
\newcommand{\SpecialCharTok}[1]{\textcolor[rgb]{0.25,0.44,0.63}{#1}}
\newcommand{\SpecialStringTok}[1]{\textcolor[rgb]{0.73,0.40,0.53}{#1}}
\newcommand{\StringTok}[1]{\textcolor[rgb]{0.25,0.44,0.63}{#1}}
\newcommand{\VariableTok}[1]{\textcolor[rgb]{0.10,0.09,0.49}{#1}}
\newcommand{\VerbatimStringTok}[1]{\textcolor[rgb]{0.25,0.44,0.63}{#1}}
\newcommand{\WarningTok}[1]{\textcolor[rgb]{0.38,0.63,0.69}{\textbf{\textit{#1}}}}
\usepackage{longtable,booktabs,array}
\usepackage{calc} % for calculating minipage widths
% Correct order of tables after \paragraph or \subparagraph
\usepackage{etoolbox}
\makeatletter
\patchcmd\longtable{\par}{\if@noskipsec\mbox{}\fi\par}{}{}
\makeatother
% Allow footnotes in longtable head/foot
\IfFileExists{footnotehyper.sty}{\usepackage{footnotehyper}}{\usepackage{footnote}}
\makesavenoteenv{longtable}
\usepackage{graphicx}
\makeatletter
\def\maxwidth{\ifdim\Gin@nat@width>\linewidth\linewidth\else\Gin@nat@width\fi}
\def\maxheight{\ifdim\Gin@nat@height>\textheight\textheight\else\Gin@nat@height\fi}
\makeatother
% Scale images if necessary, so that they will not overflow the page
% margins by default, and it is still possible to overwrite the defaults
% using explicit options in \includegraphics[width, height, ...]{}
\setkeys{Gin}{width=\maxwidth,height=\maxheight,keepaspectratio}
% Set default figure placement to htbp
\makeatletter
\def\fps@figure{htbp}
\makeatother
\setlength{\emergencystretch}{3em} % prevent overfull lines
\providecommand{\tightlist}{%
  \setlength{\itemsep}{0pt}\setlength{\parskip}{0pt}}
\setcounter{secnumdepth}{-\maxdimen} % remove section numbering
\ifLuaTeX
\usepackage[bidi=basic]{babel}
\else
\usepackage[bidi=default]{babel}
\fi
\babelprovide[main,import]{vietnamese}
\ifPDFTeX
\else
\babelfont{rm}[]{TeX Gyre Termes}
\fi
% get rid of language-specific shorthands (see #6817):
\let\LanguageShortHands\languageshorthands
\def\languageshorthands#1{}
% ctu_xelatex_header.tex — header-includes cho build LaTeX CTU (xelatex).
% Ánh xạ các ký tự Unicode mà TeX Gyre Termes thiếu glyph (chỉ số dưới/trên,
% ký tự kẻ hộp) sang lệnh LaTeX tương đương để không bị rớt glyph khi biên dịch.
\usepackage{newunicodechar}
\usepackage{pifont}
\usepackage{amssymb}
\newunicodechar{✓}{\ding{51}}
\newunicodechar{✗}{\ding{55}}
\newunicodechar{≥}{\ensuremath{\geq}}
\newunicodechar{≤}{\ensuremath{\leq}}
\newunicodechar{≈}{\ensuremath{\approx}}
\newunicodechar{↔}{\ensuremath{\leftrightarrow}}
\newunicodechar{₀}{\textsubscript{0}}
\newunicodechar{₁}{\textsubscript{1}}
\newunicodechar{₂}{\textsubscript{2}}
\newunicodechar{₃}{\textsubscript{3}}
\newunicodechar{₄}{\textsubscript{4}}
\newunicodechar{₅}{\textsubscript{5}}
\newunicodechar{₆}{\textsubscript{6}}
\newunicodechar{₇}{\textsubscript{7}}
\newunicodechar{₈}{\textsubscript{8}}
\newunicodechar{₉}{\textsubscript{9}}
\newunicodechar{²}{\textsuperscript{2}}
\newunicodechar{³}{\textsuperscript{3}}
\newunicodechar{─}{\rule[0.5ex]{0.9em}{0.4pt}}

% --- Bảng rộng: thu nhỏ cỡ chữ + khoảng cột, cho phép ngắt token dài để
%     bảng không tràn lề phải (overfull hbox). Áp dụng cho mọi longtable
%     (pandoc sinh toàn bộ bảng dạng longtable).
\usepackage{etoolbox}
\usepackage{seqsplit}      % ngắt được token rất dài (mã biến, số dài)
% Lưu ý: KHÔNG đổi \tabcolsep — pandoc tính bề rộng cột theo công thức
% (\columnwidth - 4\tabcolsep)*frac (tham chiếu \tabcolsep trực tiếp), nên
% giảm \tabcolsep lại làm cột RỘNG hơn và tràn nặng hơn.
\AtBeginEnvironment{longtable}{%
  \footnotesize
  \emergencystretch=3em
  \hbadness=10000          % không cảnh báo underfull do RaggedRight
}
% Cho phép TeX bẻ dòng linh hoạt hơn trong ô p{} hẹp
\setlength{\emergencystretch}{2em}

% URL/DOI dài: cho phép ngắt ở bất kỳ ký tự nào (tránh tràn lề ở danh mục TLTK).
% Kết hợp với pandoc +autolink_bare_uris để URL trần được bọc \url{}.
\usepackage{xurl}
\ifLuaTeX
  \usepackage{selnolig}  % disable illegal ligatures
\fi
\IfFileExists{bookmark.sty}{\usepackage{bookmark}}{\usepackage{hyperref}}
\IfFileExists{xurl.sty}{\usepackage{xurl}}{} % add URL line breaks if available
\urlstyle{same}
\hypersetup{
  pdflang={vi},
  colorlinks=true,
  linkcolor={black},
  filecolor={Maroon},
  citecolor={black},
  urlcolor={blue},
  pdfcreator={LaTeX via pandoc}}

\author{}
\date{}

\begin{document}

\setstretch{1.2}
\thispagestyle{empty}
\begin{center}
{\fontsize{14}{18}\selectfont\bfseries BỘ GIÁO DỤC VÀ ĐÀO TẠO\\TRƯỜNG ĐẠI HỌC CẦN THƠ}

\vspace{2.2cm}
{\fontsize{14}{18}\selectfont\bfseries ĐỖ THÙY HƯƠNG}

\vspace{2.4cm}
{\fontsize{20}{26}\selectfont\bfseries QUỐC TẾ HÓA VÀ HIỆU QUẢ HOẠT ĐỘNG\\KINH DOANH CỦA CÁC DOANH NGHIỆP Ở CHÂU Á}

\vspace{2.4cm}
{\fontsize{14}{18}\selectfont\bfseries LUẬN ÁN TIẾN SĨ}

{\fontsize{14}{18}\selectfont\bfseries NGÀNH QUẢN TRỊ KINH DOANH}

{\fontsize{14}{18}\selectfont\bfseries MÃ SỐ: 9340101}

\vspace{2.8cm}
{\fontsize{14}{18}\selectfont\bfseries CẦN THƠ, NĂM 2026}
\end{center}
\newpage

\thispagestyle{empty}
\begin{center}
{\fontsize{14}{18}\selectfont\bfseries BỘ GIÁO DỤC VÀ ĐÀO TẠO\\TRƯỜNG ĐẠI HỌC CẦN THƠ}

\vspace{1.6cm}
{\fontsize{14}{18}\selectfont\bfseries ĐỖ THÙY HƯƠNG}

{\fontsize{14}{18}\selectfont\bfseries MÃ SỐ NCS: \rule{3cm}{0.4pt}}

\vspace{1.6cm}
{\fontsize{20}{26}\selectfont\bfseries QUỐC TẾ HÓA VÀ HIỆU QUẢ HOẠT ĐỘNG\\KINH DOANH CỦA CÁC DOANH NGHIỆP Ở CHÂU Á}

\vspace{1.6cm}
{\fontsize{14}{18}\selectfont\bfseries LUẬN ÁN TIẾN SĨ}

{\fontsize{14}{18}\selectfont\bfseries NGÀNH QUẢN TRỊ KINH DOANH}

{\fontsize{14}{18}\selectfont\bfseries MÃ SỐ: 9340101}

\vspace{1.4cm}
{\fontsize{13}{17}\selectfont\bfseries NGƯỜI HƯỚNG DẪN KHOA HỌC}

{\fontsize{13}{17}\selectfont\bfseries PGS.TS. PHAN ANH TÚ}

\vspace{1.4cm}
{\fontsize{14}{18}\selectfont\bfseries CẦN THƠ, NĂM 2026}
\end{center}
\newpage

\thispagestyle{empty}
\begin{center}
{\fontsize{14}{18}\selectfont\bfseries CHẤP THUẬN CỦA HỘI ĐỒNG}
\end{center}

\vspace{0.5cm}
\hspace*{1cm}Luận án này với tựa đề là ``Quốc tế hóa và hiệu quả hoạt động kinh doanh của các doanh nghiệp ở châu Á'', do nghiên cứu sinh Đỗ Thùy Hương thực hiện theo sự hướng dẫn của PGS.TS. Phan Anh Tú. Luận án đã báo cáo và được Hội đồng đánh giá luận án tiến sĩ thông qua ngày \rule{1.2cm}{0.4pt} tháng \rule{1.2cm}{0.4pt} năm \rule{1.6cm}{0.4pt}. Luận án đã được chỉnh sửa theo góp ý và được Hội đồng đánh giá luận án xem lại.

\vspace{1.0cm}
\vspace{0.8cm}
\noindent\begin{minipage}[t]{0.48\textwidth}\centering Thư ký\\[2cm](ký tên)\end{minipage}\hfill\begin{minipage}[t]{0.48\textwidth}\centering Ủy viên\\[2cm](ký tên)\end{minipage}

\vspace{1.2cm}
\noindent\begin{minipage}[t]{0.48\textwidth}\centering Ủy viên\\[2cm](ký tên)\end{minipage}\hfill\begin{minipage}[t]{0.48\textwidth}\centering Phản biện 3\\[2cm](ký tên)\end{minipage}

\vspace{1.2cm}
\noindent\begin{minipage}[t]{0.48\textwidth}\centering Phản biện 2\\[2cm](ký tên)\end{minipage}\hfill\begin{minipage}[t]{0.48\textwidth}\centering Phản biện 1\\[2cm](ký tên)\end{minipage}

\vspace{1.2cm}
\noindent\begin{minipage}[t]{0.48\textwidth}\centering Người hướng dẫn\\[2cm](ký tên)\end{minipage}\hfill\begin{minipage}[t]{0.48\textwidth}\centering Chủ tịch Hội đồng\\[2cm](ký tên)\end{minipage}

\vspace{0.8cm}
\noindent{\footnotesize *: Ghi đầy đủ học hàm, học vị và họ tên.}

\newpage

\hypertarget{lux1eddi-cux1ea3m-ux1a1n}{%
\section*{LỜI CẢM ƠN}\label{lux1eddi-cux1ea3m-ux1a1n}}
\addcontentsline{toc}{section}{LỜI CẢM ƠN}

\emph{(NCS hoàn thiện: cảm ơn người hướng dẫn khoa học, Trường Kinh tế
và Trường Đại học Cần Thơ, Trường Đại học Kỹ thuật -- Công nghệ Vĩnh
Long, gia đình và đồng nghiệp.)}

\begin{center}\rule{0.5\linewidth}{0.5pt}\end{center}

\hypertarget{tuxf3m-tux1eaft}{%
\section*{TÓM TẮT}\label{tuxf3m-tux1eaft}}
\addcontentsline{toc}{section}{TÓM TẮT}

Luận án nghiên cứu mối quan hệ giữa quốc tế hóa và hiệu quả hoạt động
của doanh nghiệp (internationalization--firm performance, I--P) tại châu
Á, tập trung làm rõ các điều kiện biên giải thích sự dị biệt cao trong
mối quan hệ này. Trên cơ sở tích hợp bốn lý thuyết kinh điển, mô hình
Uppsala, lý thuyết dựa trên nguồn lực (RBV), lý thuyết thể chế và lý
thuyết thượng tầng quản trị, cùng lớp mở rộng về năng lực số, luận án
xây dựng khung giải thích đa tầng CDCM (Country--Digital--Capability
Moderation) với ba tầng điều tiết: bối cảnh thể chế (vận hành hóa bằng
khung phân loại ICRV sáu nhóm), năng lực số (phân tách năng lực công
nghệ TCI và chỉ số chấp nhận số DAI), và đặc điểm nhà quản trị cấp cao.

Luận án sử dụng thiết kế hỗn hợp tổng hợp--thực nghiệm (mixed
synthesis-empirical): (i) một phân tích tổng hợp (meta-analysis) ba tầng
giai đoạn 1977--2026 với k = 238 nghiên cứu; và (ii) phân tích thực
nghiệm dữ liệu doanh nghiệp từ Khảo sát Doanh nghiệp của Ngân hàng Thế
giới (WBES) trên 50 nền kinh tế châu Á và Thái Bình Dương (gồm Nhật Bản;
mẫu hồi quy chính N = 81.022 doanh nghiệp), bổ sung bằng các phân tích
quốc gia chuyên sâu (Việt Nam, Singapore, Trung Quốc, Ấn Độ) và trường
hợp biên là bảy nền kinh tế đảo nhỏ đang phát triển (SIDS) Thái Bình
Dương.

Kết quả chính: (1) quan hệ I--P có dạng phi tuyến chữ U ngược với điểm
uốn gộp 43,6\% (mô hình đầy đủ, 50 nền kinh tế); chữ U ngược định hình
rõ nét nhất ở nhóm chuyển đổi thu nhập trung bình--thấp (43,0\%), duỗi
thẳng gần tuyến tính ở nhóm thể chế mạnh nhất và mất cấu trúc ở các nhóm
thể chế yếu nhất; (2) cả TCI và DAI đều \textbf{nâng mặt bằng hiệu quả
một cách phổ quát} (TCI là năng lực nền tảng, DAI là nguồn lực tình
huống), nhưng vai trò \textbf{uốn đường cong} của chúng không phổ quát:
trên toàn mẫu 50 nền các tương tác với đường cong không đạt ý nghĩa
thống kê, và chỉ định hình rõ ở bối cảnh đơn quốc gia đặc thù (điển hình
là hiệu ứng ``lá chắn số'' của DAI tại Singapore), với xu hướng mạnh hơn
ở môi trường thể chế yếu nhưng chỉ ở mức biên ý nghĩa; (3) tại nhóm
SIDS, quan hệ I--P chuyển thành âm đơn điệu không có điểm uốn, hiện
tượng được luận án đặt tên là \textbf{gánh nặng quốc tế hóa bắt buộc
(Forced Internationalization Penalty, FIP)}. Luận án đóng góp khung tích
hợp đa tầng, khung phân loại thể chế ICRV, và cấu trúc lý thuyết mới FIP
cho dòng nghiên cứu kinh doanh quốc tế, đồng thời cung cấp hàm ý chính
sách và quản trị theo từng bối cảnh thể chế.

\textbf{Từ khóa:} quốc tế hóa; hiệu quả doanh nghiệp; thể chế; năng lực
số; nhà quản trị cấp cao; châu Á; WBES; phi tuyến chữ U ngược.

\begin{center}\rule{0.5\linewidth}{0.5pt}\end{center}

\hypertarget{abstract}{%
\section*{ABSTRACT}\label{abstract}}
\addcontentsline{toc}{section}{ABSTRACT}

This dissertation examines the internationalization--firm performance
(I--P) relationship across Asia, focusing on the boundary conditions
that explain its well-documented heterogeneity. Integrating four
foundational theories, the Uppsala model, the resource-based view (RBV),
institutional theory, and upper echelons theory, with a
digital-capability extension, the dissertation develops a multi-tier
moderation framework, CDCM (Country--Digital--Capability Moderation),
comprising three moderating tiers: institutional context
(operationalized through a six-group ICRV taxonomy), digital capability
(separating technological capability, TCI, from digital adoption, DAI),
and top-manager characteristics.

The study adopts a mixed synthesis-empirical design: (i) a three-level
meta-analysis spanning 1977--2026 (k = 238 studies); and (ii) firm-level
empirical analysis of World Bank Enterprise Survey (WBES) microdata
across 50 Asian and Pacific economies (including Japan; main regression
sample N = 81,022 firms), complemented by country-specific analyses
(Vietnam, Singapore, China, India) and a boundary case of seven Pacific
Small Island Developing States (SIDS).

Key findings: (1) the I--P relationship is inverted-U shaped with a
pooled turning point of 43.6\% (full model, 50 economies); the
inverted-U is sharpest in the lower-middle transition regime (43.0\%),
flattens to near-linearity in the strongest-institution regime, and
dissolves in the weakest regimes; (2) both TCI (a foundational
capability) and DAI (a situational resource) \textbf{universally shift
the performance level}, but their curve-\textbf{reshaping} role is not
universal: across the 50-economy pool the curvature interactions are not
statistically significant, and curve reshaping appears only in specific
single-economy contexts (notably DAI's ``digital shield'' in Singapore),
with a tendency toward stronger moderation in weaker institutions that
is only marginally significant; (3) in the SIDS group the relationship
becomes monotonically negative with no turning point, a phenomenon the
dissertation names the \textbf{Forced Internationalization Penalty
(FIP)}. The dissertation contributes an integrated multi-tier framework,
the ICRV institutional taxonomy, and the novel FIP construct to
international-business research, alongside context-contingent policy and
managerial implications.

\textbf{Keywords:} internationalization; firm performance; institutions;
digital capability; top management; Asia; WBES; inverted-U relationship.

\newpage

\hypertarget{lux1eddi-cam-ux111oan}{%
\section*{LỜI CAM ĐOAN}\label{lux1eddi-cam-ux111oan}}
\addcontentsline{toc}{section}{LỜI CAM ĐOAN}

Tôi xin cam đoan luận án tiến sĩ này là công trình nghiên cứu của riêng
tôi dưới sự hướng dẫn khoa học của người hướng dẫn. Các kết quả phân
tích, số liệu, bảng biểu và hình vẽ trình bày trong luận án đều trung
thực và chưa từng được công bố trong bất kỳ công trình nào khác, ngoại
trừ các bài báo khoa học do chính tôi là tác giả/đồng tác giả được liệt
kê trong \emph{Danh mục công trình đã công bố của nghiên cứu sinh} (luận
án được thực hiện theo định hướng kết hợp các công trình thành phần).
Luận án được phát triển trên cơ sở \textbf{hai chuyên đề tiến sĩ do
chính tôi thực hiện trong khuôn khổ chương trình đào tạo tiến sĩ tại
Trường Đại học Cần Thơ} (theo Quyết định giao chuyên đề của Nhà trường);
do đó, một phần nội dung của luận án kế thừa và có thể trùng lặp với hai
chuyên đề này, đây là phần công trình của chính tôi, không phải sao chép
từ nguồn của người khác. Dữ liệu sử dụng là dữ liệu thứ cấp công khai từ
Khảo sát Doanh nghiệp của Ngân hàng Thế giới (World Bank Enterprise
Surveys). Mọi trích dẫn từ các nghiên cứu khác đều được nêu rõ nguồn
theo chuẩn APA 7th và được liệt kê đầy đủ trong danh mục tài liệu tham
khảo. Tôi xin chịu trách nhiệm hoàn toàn về nội dung khoa học của luận
án này.

\emph{Cần Thơ, ngày \ldots{} tháng \ldots{} năm 2026}

\textbf{Nghiên cứu sinh}

Đỗ Thùy Hương

\newpage

\renewcommand{\contentsname}{MỤC LỤC}
\setcounter{tocdepth}{3}
\tableofcontents

\begin{center}\rule{0.5\linewidth}{0.5pt}\end{center}

\hypertarget{danh-mux1ee5c-bux1ea3ng}{%
\section*{DANH MỤC BẢNG}\label{danh-mux1ee5c-bux1ea3ng}}
\addcontentsline{toc}{section}{DANH MỤC BẢNG}

\textbf{Bảng đánh số:} - Bảng 2.1. Ánh xạ nghiên cứu thành phần, mục
trình bày, giả thuyết (Chương 2, Mục 2.5) - Bảng 2.2. Tổng hợp hệ giả
thuyết: phát biểu, dấu kỳ vọng, cơ sở lý thuyết và biến kiểm định
(Chương 2, Mục 2.5.8) - Bảng 3.1. Tổng hợp biến nghiên cứu của khung
thực nghiệm đa quốc gia (Chương 3, Mục 3.3.6) - Bảng 4.1. Phân tán năng
suất lao động và đặc điểm quốc tế hóa theo phân nhóm con ICRV (Chương 4,
Mục 4.1) - Bảng 4.2. Đổi mới sáng tạo và chấp nhận số theo phân nhóm con
ICRV (\%) (Chương 4, Mục 4.1) - Bảng 4.3. Đặc trưng doanh nghiệp Nhật
Bản 2025 so với năm nền còn lại của Nhóm I (Chương 4, Mục 4.1.5) - Bảng
4.4. Trung vị tỷ suất lợi nhuận trên doanh thu (ROS) theo nhóm ICRV
(Chương 4, Mục 4.1.8) - Bảng 4.5. Tiến triển hệ số theo chuỗi mô hình
phân tầng, Singapore 2023 (Chương 4, Mục 4.3.8) - Bảng 4.6. Tiến triển
mô hình của P7 trên khung 50 nền: điểm uốn và kiểm định Lind--Mehlum
(Chương 4, Mục 4.6.1) - Bảng 4.7. Điểm uốn theo nhóm ICRV, khung 50 nền
(Chương 4, Mục 4.6.2) - Bảng 4.8. Điểm uốn theo đợt khảo sát, Ấn Độ
(Chương 4, Mục 4.8.1) - Bảng 4.9. Hiệu ứng cường độ xuất khẩu lên năng
suất qua chuỗi mô hình, Nhật Bản 2025 (Chương 4, Mục 4.9.1) - Bảng 4.10.
Tổng hợp ước lượng các nghiên cứu thành phần thực nghiệm (Chương 4, Mục
4.11)

\textbf{Bảng chuyên đề (không đánh số, theo mục):} - Bảng định nghĩa
biến, Nghiên cứu P3 Việt Nam (Chương 3, Mục 3.4.5.1) - Bảng định nghĩa
biến, Nghiên cứu P4 Singapore (Chương 3, Mục 3.4.5.2) - Bảng định nghĩa
biến, Nghiên cứu P5 Trung Quốc (Chương 3, Mục 3.4.5.3) - Bảng định nghĩa
biến, Nghiên cứu P7 đa quốc gia (Chương 3, Mục 3.4.5.4) - Bảng định
nghĩa biến, Nghiên cứu P8 SIDS (Chương 3, Mục 3.4.5.5) - Bảng tóm tắt
phương pháp: Kế thừa và đóng góp mới (Chương 3, Mục 3.6) - Bảng tổng hợp
kết quả kiểm định hệ giả thuyết H1--H6, H1b, H1c (Chương 4, Mục 4.11)

\hypertarget{danh-mux1ee5c-huxecnh}{%
\section*{DANH MỤC HÌNH}\label{danh-mux1ee5c-huxecnh}}
\addcontentsline{toc}{section}{DANH MỤC HÌNH}

\begin{itemize}
\tightlist
\item
  Hình 2.1. Khung khái niệm tích hợp đa tầng CDCM (Chương 2, Mục 2.5)
\item
  Hình 3.1. Quy trình nghiên cứu của luận án (Chương 3, Mục 3.1.2)
\item
  Hình 4.1. Thành phần mẫu gộp theo nhóm ICRV (Chương 4, Mục 4.1)
\item
  Hình 4.2. Quy mô mẫu gộp theo năm khảo sát (Chương 4, Mục 4.1)
\item
  Hình 4.3. Quỹ đạo tăng trưởng theo nhóm ICRV (Chương 4, Mục 4.1)
\item
  Hình 4.4. Phân bố quy mô doanh nghiệp theo nhóm ICRV (Chương 4, Mục
  4.1)
\item
  Hình 4.5. Gradient năng lực công nghệ theo sáu nhóm ICRV (Chương 4,
  Mục 4.1)
\item
  Hình 4.6. Phân bố mật độ năng suất lao động theo nhóm ICRV (Chương 4,
  Mục 4.1)
\item
  Hình 4.7. Biểu đồ radar năng lực và quốc tế hóa theo nhóm ICRV (Chương
  4, Mục 4.1)
\item
  Hình 4.8. Họ đường cong quốc tế hóa và hiệu quả phụ thuộc chế độ thể
  chế ICRV (Chương 4, Mục 4.6.2)
\end{itemize}

\begin{center}\rule{0.5\linewidth}{0.5pt}\end{center}

\hypertarget{danh-mux1ee5c-tux1eeb-viux1ebft-tux1eaft}{%
\section*{DANH MỤC TỪ VIẾT
TẮT}\label{danh-mux1ee5c-tux1eeb-viux1ebft-tux1eaft}}
\addcontentsline{toc}{section}{DANH MỤC TỪ VIẾT TẮT}

\begin{longtable}[]{@{}
  >{\raggedright\arraybackslash}p{(\columnwidth - 4\tabcolsep) * \real{0.3333}}
  >{\raggedright\arraybackslash}p{(\columnwidth - 4\tabcolsep) * \real{0.3333}}
  >{\raggedright\arraybackslash}p{(\columnwidth - 4\tabcolsep) * \real{0.3333}}@{}}
\toprule\noalign{}
\begin{minipage}[b]{\linewidth}\raggedright
Viết tắt
\end{minipage} & \begin{minipage}[b]{\linewidth}\raggedright
Tiếng Anh
\end{minipage} & \begin{minipage}[b]{\linewidth}\raggedright
Tiếng Việt
\end{minipage} \\
\midrule\noalign{}
\endhead
\bottomrule\noalign{}
\endlastfoot
APA & American Psychological Association & Chuẩn trích dẫn APA \\
CDCM & Country--Digital--Capability Moderation & Khung điều tiết thể
chế--số--năng lực (khung tích hợp của luận án) \\
DAI & Digital Adoption Index & Chỉ số chấp nhận số \\
EMNE & Emerging Market Multinational Enterprise & Doanh nghiệp đa quốc
gia thị trường mới nổi \\
FDI & Foreign Direct Investment & Đầu tư trực tiếp nước ngoài \\
FGLS & Feasible Generalized Least Squares & Bình phương nhỏ nhất tổng
quát khả thi \\
FIP & Forced Internationalization Penalty & Gánh nặng quốc tế hóa bắt
buộc (cấu trúc lý thuyết mới của luận án) \\
FSTS & Foreign Sales to Total Sales & Tỷ lệ doanh thu nước ngoài trên
tổng doanh thu (cường độ xuất khẩu) \\
FATA & Foreign Assets to Total Assets & Tỷ lệ tài sản nước ngoài trên
tổng tài sản \\
HC1/HC3 & Heteroskedasticity-Consistent SE & Sai số chuẩn vững với
phương sai sai số thay đổi \\
ICRV & Institutional Context Regime Variation & Biến thiên chế độ bối
cảnh thể chế (khung phân loại sáu nhóm của luận án) \\
I--P & Internationalization--Performance & Quan hệ quốc tế hóa--hiệu
quả \\
ISIC & International Standard Industrial Classification & Phân loại
ngành chuẩn quốc tế \\
LM & Lind--Mehlum (U-test) & Kiểm định U của Lind--Mehlum \\
LP & Labor Productivity & Năng suất lao động \\
MARA & Multilevel/Multivariate Random-effects (meta-analysis) & Phân
tích tổng hợp đa tầng hiệu ứng ngẫu nhiên \\
MNC/MNE & Multinational Corporation/Enterprise & Công ty/doanh nghiệp đa
quốc gia \\
OLS & Ordinary Least Squares & Bình phương nhỏ nhất thông thường \\
PPP & Purchasing Power Parity & Ngang giá sức mua \\
PRISMA & Preferred Reporting Items for Systematic Reviews and
Meta-Analyses & Chuẩn báo cáo tổng quan hệ thống và phân tích tổng
hợp \\
R\&D & Research and Development & Nghiên cứu và phát triển \\
RBV & Resource-Based View & Lý thuyết dựa trên nguồn lực \\
ROA/ROE/ROS & Return on Assets/Equity/Sales & Tỷ suất sinh lời trên tài
sản/vốn chủ sở hữu/doanh thu \\
SIDS & Small Island Developing States & Các quốc gia đảo nhỏ đang phát
triển \\
SME & Small and Medium-sized Enterprise & Doanh nghiệp nhỏ và vừa \\
TCI & Technological Capability Index & Chỉ số năng lực công nghệ \\
TP & Turning Point & Điểm uốn (của quan hệ phi tuyến) \\
UPI & Unified Payments Interface & Hạ tầng thanh toán hợp nhất (Ấn
Độ) \\
VIF & Variance Inflation Factor & Hệ số phóng đại phương sai \\
VRIN & Valuable, Rare, Inimitable, Non-substitutable & Có giá trị, hiếm,
khó sao chép, khó thay thế \\
WBES & World Bank Enterprise Surveys & Khảo sát Doanh nghiệp của Ngân
hàng Thế giới \\
WGI & Worldwide Governance Indicators & Bộ chỉ số quản trị toàn cầu \\
\end{longtable}

\begin{center}\rule{0.5\linewidth}{0.5pt}\end{center}

\hypertarget{danh-mux1ee5c-cuxf4ng-truxecnh-cux1ee7a-nghiuxean-cux1ee9u-sinh-liuxean-quan-ux111ux1ebfn-luux1eadn-uxe1n}{%
\section*{DANH MỤC CÔNG TRÌNH CỦA NGHIÊN CỨU SINH LIÊN QUAN ĐẾN LUẬN
ÁN}\label{danh-mux1ee5c-cuxf4ng-truxecnh-cux1ee7a-nghiuxean-cux1ee9u-sinh-liuxean-quan-ux111ux1ebfn-luux1eadn-uxe1n}}
\addcontentsline{toc}{section}{DANH MỤC CÔNG TRÌNH CỦA NGHIÊN CỨU SINH
LIÊN QUAN ĐẾN LUẬN ÁN}

Luận án được phát triển gắn với chuỗi nghiên cứu thành phần (ký hiệu nội
bộ P1--P11) và một chương sách của nghiên cứu sinh. Các công trình được
phân tách rõ theo trạng thái công bố.

\textbf{A. Đã công bố (đã bình duyệt)}

\begin{enumerate}
\def\labelenumi{\arabic{enumi}.}
\tightlist
\item
  Do, T. H., \& Phan, A. T. (2025). Internationalization and firm
  performance: A meta-analysis review. In \emph{Proceedings of the Sixth
  International Conference on Sustainable Development in Economics,
  Business, and Finance (ICBEF 2024)}. Can Tho University, Vietnam.
  \emph{(Phân tích tổng hợp cơ sở, k = 113; nền tảng cho P1/P2/P6.)}
\item
  Do, T. H., \& Phan, A. T. (2025). Internationalization and firm
  performance of firms in India: The role of top management. In M.
  Bartekova (Ed.), \emph{International business research: Traditional
  and creative approaches}. IntechOpen.
  \url{https://doi.org/10.5772/intechopen.1011012} \emph{(Công trình
  tiền đề đã công bố về Ấn Độ, 380 doanh nghiệp WBES, ước lượng FGLS,
  lăng kính thượng tầng quản trị; được nghiên cứu thành phần P9 trích
  dẫn như công trình tiền đề của tác giả.)}
\item
  Do, T. H., \& Phan, A. T. (2026). Firm performance heterogeneity in
  emerging Asia: Evidence from World Bank Enterprise Surveys.
  \emph{Vietnam Economic \& Financial Review, 2}(1), 111--115.
  \emph{(P1.)}
\item
  Do, T. H., \& Phan, A. T. (2026). Unveiling the impact of Chinese
  manufacturing SMEs' internationalization on performance. \emph{Journal
  of Finance \& Accounting Research, 39}(2), 287--291. \emph{(P2.)}
\end{enumerate}

\textbf{B. Bản thảo đang trong quá trình bình duyệt (CHƯA công bố)}

\begin{enumerate}
\def\labelenumi{\arabic{enumi}.}
\setcounter{enumi}{4}
\tightlist
\item
  Nghiên cứu thành phần P3, Việt Nam (bối cảnh chuyển đổi thu nhập trung
  bình thấp). \emph{{[}Bản thảo đang bình duyệt.{]}}
\item
  Nghiên cứu thành phần P4, Singapore (bối cảnh thể chế tiên tiến--đổi
  mới). \emph{{[}Bản thảo đang bình duyệt.{]}}
\item
  Nghiên cứu thành phần P5, Trung Quốc 2012--2024 (bối cảnh trung bình
  cao). \emph{{[}Bản thảo đang bình duyệt.{]}}
\item
  Nghiên cứu thành phần P6, phân tích tổng hợp ba tầng mở rộng
  1977--2026. \emph{{[}Bản thảo đang bình duyệt.{]}}
\item
  Nghiên cứu thành phần P7, Kiểm định toàn mẫu 50 nền kinh tế châu Á và
  Thái Bình Dương (nghiên cứu tổng hợp đỉnh). \emph{{[}Bản thảo đang
  chỉnh sửa theo phản biện.{]}}
\item
  Nghiên cứu thành phần P8, Trường hợp biên SIDS, gánh nặng quốc tế hóa
  bắt buộc (FIP). \emph{{[}Bản thảo đang chỉnh sửa theo phản biện.{]}}
\item
  Nghiên cứu thành phần P9, Ấn Độ: độ bền ngưỡng I--P qua thập kỷ chuyển
  đổi thể chế 2014--2025 (ba đợt WBES, \textasciitilde29.000 doanh
  nghiệp; phát hiện sự \emph{tan biến} điểm uốn chữ U ngược dưới áp lực
  cải cách thể chế và bão hòa hạ tầng số UPI). Đích: Management
  International Review (Springer)/IJOEM. \emph{{[}Bản thảo đang bình
  duyệt.{]}}
\item
  Nghiên cứu thành phần P10, Nhật Bản: kiểm định biên trên của gradient
  thể chế trên đợt khảo sát WBES khai mở 2025. \emph{{[}Bản thảo chuẩn
  bị gửi Asian Business \& Management.{]}}
\item
  Nghiên cứu thành phần P11 (sản phẩm phái sinh theo định hướng phát
  triển), Khoảng cách số tiền-AI và năng suất doanh nghiệp qua các chế
  độ thể chế châu Á (50 nền WBES): khai thác bằng chứng DAI×ICRV của
  luận án cho chủ đề AI--phát triển. Đích: \emph{Journal of Economics
  and Development}, số đặc biệt ``Sustainable Growth and Welfare in the
  Age of AI'' (Emerald; Scopus/ESCI/ABDC; diamond OA, không phí; hạn
  31/10/2026). \emph{{[}Bản thảo chuẩn bị gửi.{]}}
\end{enumerate}

\emph{Ghi chú về liêm chính học thuật: Theo định hướng luận án kết hợp
công trình thành phần, các nghiên cứu P1--P11 là một phần nội dung
nghiên cứu của chính luận án (P11 là sản phẩm phái sinh theo định hướng
phát triển, khai thác lại bằng chứng của luận án); do đó trong các
chương, chúng được dẫn chiếu bằng nhãn nội bộ (P1--P10) thay vì trích
dẫn tác giả--năm. Các bản thảo ở mục B CHƯA được công bố/bình duyệt
chính thức, không được diễn giải như bằng chứng đã qua bình duyệt độc
lập; tên tạp chí mục tiêu được nêu để minh bạch định hướng công bố;
trạng thái bình duyệt sẽ được cập nhật tại thời điểm bảo vệ. NCS cập
nhật trạng thái tại thời điểm bảo vệ.}

\begin{center}\rule{0.5\linewidth}{0.5pt}\end{center}

\newpage

\hypertarget{chux1b0ux1a1ng-1-giux1edbi-thiux1ec7u}{%
\chapter{CHƯƠNG 1: GIỚI
THIỆU}\label{chux1b0ux1a1ng-1-giux1edbi-thiux1ec7u}}

\begin{center}\rule{0.5\linewidth}{0.5pt}\end{center}

\hypertarget{bux1ed1i-cux1ea3nh-nghiuxean-cux1ee9u}{%
\section{1.1 Bối cảnh nghiên
cứu}\label{bux1ed1i-cux1ea3nh-nghiuxean-cux1ee9u}}

Quốc tế hóa doanh nghiệp là một trong những chủ đề trung tâm của kinh
doanh quốc tế trong hơn nửa thế kỷ qua. Kể từ khi Johanson và Vahlne
(1977) đề xuất mô hình Uppsala về quá trình quốc tế hóa tuần tự, hàng
thập kỷ nghiên cứu đã tích lũy nhằm trả lời câu hỏi cốt lõi: liệu quốc
tế hóa có thực sự cải thiện hiệu quả hoạt động của doanh nghiệp hay
không, và trong điều kiện nào điều đó xảy ra? Lu và Beamish (2004) cung
cấp bằng chứng quan trọng về quan hệ hình chữ S (S-curve) giữa mức độ
quốc tế hóa và hiệu quả, trong khi Contractor et al.~(2003) và nhiều
nghiên cứu tiếp theo tiếp tục thách thức, bổ sung và điều chỉnh kết luận
đó. Dù tranh luận học thuật vẫn chưa ngã ngũ, sự đồng thuận hiện nay cho
thấy mối quan hệ quốc tế hóa--hiệu quả doanh nghiệp (I--P) là có thực
nhưng phức tạp, phụ thuộc vào nhiều điều kiện biên.

Ở châu Á, bối cảnh này càng trở nên đặc biệt đáng quan tâm vì khu vực
đang chứng kiến tốc độ hội nhập kinh tế và chuyển đổi số hết sức không
đồng đều giữa các quốc gia. Từ Singapore, một nền kinh tế nhỏ mở, dẫn
đầu thế giới về năng lực đổi mới và kết nối toàn cầu, đến Việt Nam, một
nền kinh tế chuyển đổi đang đẩy nhanh cải cách thể chế và hội nhập quốc
tế, rồi đến Trung Quốc với quy mô thị trường khổng lồ và cấu trúc thể
chế đặc thù, thậm chí đến các quốc đảo nhỏ Thái Bình Dương đang phải đối
mặt với thách thức sinh tồn kép của biến đổi khí hậu và bẫy phụ thuộc
xuất khẩu, sự đa dạng này tạo ra một phòng thí nghiệm thực địa lý tưởng
để kiểm định các điều kiện biên của quan hệ quốc tế hóa và hiệu quả
(Khanna \& Palepu, 2010; Peng, 2003).

Chuyển đổi số đang làm thay đổi căn bản cách thức doanh nghiệp quốc tế
hóa và cách thức quốc tế hóa ảnh hưởng đến hiệu quả hoạt động. Banalieva
và Dhanaraj (2019) lập luận rằng lý thuyết nội địa hóa cần được cập nhật
cho nền kinh tế số, vì dữ liệu số và tương tác qua nền tảng làm giảm sự
phụ thuộc vào hiện diện vật lý. Stallkamp và Schotter (2021) chỉ ra rằng
các nền tảng số quốc tế hóa theo logic hoàn toàn khác các công ty đa
quốc gia truyền thống. Verhoef et al.~(2021) cung cấp khung phân tầng ba
cấp từ số hóa thông tin đến số hóa quy trình và chuyển đổi số toàn diện.
Tuy nhiên, phần lớn bằng chứng thực nghiệm hiện có về vai trò của năng
lực số trong quan hệ quốc tế hóa và hiệu quả vẫn tập trung vào các nền
kinh tế phát triển, đặc biệt là Bắc Mỹ và Tây Âu, trong khi bằng chứng
cho châu Á mới nổi còn rất hạn chế (Bhandari et al., 2023).

Không thể bỏ qua vai trò của thể chế trong việc hình thành kết quả quốc
tế hóa. Lý thuyết thể chế (North, 1990; Khanna \& Palepu, 2010) khẳng
định rằng thể chế tạo ra cấu trúc khuyến khích và chi phí giao dịch khác
nhau giữa các quốc gia. Ở các thị trường mới nổi, ``khoảng trống thể
chế'', tức sự thiếu hụt các thể chế trung gian hỗ trợ giao dịch thị
trường, tạo ra chi phí bổ sung đáng kể cho doanh nghiệp quốc tế hóa
(Peng, 2003). Sự đa dạng thể chế giữa các quốc gia châu Á, từ những nền
kinh tế với quản trị nhà nước ổn định, pháp quyền mạnh như Singapore đến
các nền kinh tế đang phát triển thể chế như Campuchia hay các quốc đảo
Thái Bình Dương, tạo ra gradient thể chế đặc biệt phong phú để kiểm
định. Luận án này tiếp cận gradient đó một cách hệ thống thông qua khung
phân loại thể chế 6 nhóm được phát triển cho luận án này, gọi tắt là
ICRV (Institutional Context Regime Variation, Biến thiên chế độ bối cảnh
thể chế).

Cuối cùng, nhà quản trị cấp cao đóng vai trò không thể bỏ qua trong việc
hình thành chiến lược và kết quả quốc tế hóa. Lý thuyết thượng tầng quản
trị (Hambrick \& Mason, 1984) khẳng định rằng quyết định chiến lược phản
ánh đặc điểm nhận thức, kinh nghiệm và giá trị của nhà quản trị. Kinh
nghiệm quốc tế giúp nhà quản trị nhận diện cơ hội và giảm chi phí của
gánh nặng người nước ngoài, trong khi sự đa dạng giới tính ở cấp lãnh
đạo có thể tăng cường quản trị rủi ro và ra quyết định đa chiều (Hsu et
al., 2013; Post \& Byron, 2015). Tích hợp tầng nhà quản trị vào khung
giải thích đa tầng là cần thiết để phản ánh đầy đủ sự phức tạp của quan
hệ quốc tế hóa và hiệu quả trong bối cảnh châu Á.

\hypertarget{bux1ed1i-cux1ea3nh-ux111ux1b0ux1a1ng-ux111ux1ea1i-thuxe1ch-thux1ee9c-viux1ec7c-luxe0m-gradient-thux1ec3-chux1ebf-vuxe0-kux1ef7-nguyuxean-sux1ed1}{%
\subsection{1.1.1 Bối cảnh đương đại: thách thức việc làm, gradient thể
chế và kỷ nguyên
số}\label{bux1ed1i-cux1ea3nh-ux111ux1b0ux1a1ng-ux111ux1ea1i-thuxe1ch-thux1ee9c-viux1ec7c-luxe0m-gradient-thux1ec3-chux1ebf-vuxe0-kux1ef7-nguyuxean-sux1ed1}}

Tính cấp thiết của câu hỏi nghiên cứu được khuếch đại bởi ba dịch chuyển
vĩ mô đang định hình lại bối cảnh của các nền kinh tế đang phát triển và
mới nổi trong giai đoạn 2025--2026. Thứ nhất là \textbf{thách thức việc
làm và cơ cấu dân số}. Theo phân tích chuyên đề mới nhất của Ngân hàng
Thế giới về nhóm thị trường biên (frontier markets), các nền kinh tế
biên hiện chiếm hơn một phần năm dân số toàn cầu nhưng chỉ tạo ra khoảng
5\% sản lượng thế giới, và dự kiến có hơn 230 triệu người bước vào độ
tuổi lao động trong giai đoạn 2025--2035 (World Bank, 2026c). Khoảng
cách năng suất ngày càng giãn rộng: thu nhập bình quân đầu người ở thị
trường biên trung vị hiện thấp hơn một phần ba so với thị trường mới nổi
điển hình, rộng hơn mức chênh lệch của năm 2000 (World Bank, 2026c).
Việc chuyển hóa cơ hội dân số thành ``lợi tức dân số'' (demographic
dividend) phụ thuộc vào năng lực tạo việc làm năng suất cao của khu vực
doanh nghiệp, điều này khiến câu hỏi về cơ chế chuyển hóa quốc tế hóa
thành hiệu quả doanh nghiệp trở thành một vấn đề phát triển cấp bách chứ
không chỉ là tranh luận học thuật.

Thứ hai là \textbf{vai trò trung tâm của chất lượng thể chế trong phân
hóa kết quả tăng trưởng}. Bằng chứng vĩ mô độc lập của Ngân hàng Thế
giới củng cố trực tiếp tiền đề lý thuyết của luận án: trong số các thị
trường biên, những nền kinh tế tăng trưởng nhanh nhất một phần tư thế kỷ
qua là những nền đã cải thiện quản trị, thu hẹp chênh lệch lãi suất cho
vay--huy động của ngân hàng, kiểm soát được nợ công và tích lũy vốn trên
đầu người mạnh hơn; ngược lại, các điểm yếu thể chế và quản trị đi kèm
chênh lệch lãi suất rộng và tăng trưởng vốn đáng thất vọng (World Bank,
2026c). Các con đường thành công lại đa dạng, từ thu hút đầu tư vào năng
lượng và khoáng sản (Kazakhstan), đến sản xuất chế tạo giá trị gia tăng
và xuất khẩu (Việt Nam), đến dịch vụ và du lịch (Rwanda), cho thấy không
tồn tại một mô hình phát triển phổ quát mà kết quả phụ thuộc vào sự
tương tác giữa chiến lược doanh nghiệp và bối cảnh thể chế (World Bank,
2026c). Mô hình phân hóa theo chế độ thể chế này chính là hiện tượng mà
khung ICRV của luận án được thiết kế để vận hành hóa và kiểm định ở cấp
doanh nghiệp.

Thứ ba là \textbf{kỷ nguyên trí tuệ nhân tạo và biến động giá hàng hóa},
hai lực lượng tạo ra bất định mới cho doanh nghiệp quốc tế hóa. Sự lan
tỏa nhanh của trí tuệ nhân tạo đang định hình lại năng lực số của cả khu
vực công lẫn khu vực tư, đặt ra yêu cầu mới về hạ tầng số công (digital
public infrastructure) và quản trị dữ liệu như những điều kiện nền tảng
cho năng lực cạnh tranh quốc gia (World Bank, 2026e). Đồng thời, cú sốc
xung đột Trung Đông năm 2026 gây gián đoạn nguồn cung dầu lớn nhất trong
lịch sử ghi nhận, đẩy giá năng lượng dự báo tăng 24\% và giá hàng hóa
nói chung tăng 16\% trong năm 2026, làm chậm tăng trưởng và đẩy lạm phát
của các nền kinh tế đang phát triển lên mức cao nhất bốn năm (World
Bank, 2026d). Hai lực lượng này tác động không đồng đều theo chế độ thể
chế: chúng vừa mở ra cơ hội ``nhảy vọt số'' cho các nền kinh tế đang
nổi, vừa khuếch đại lợi thế của các nền kinh tế tài nguyên dẫn dắt, qua
đó càng làm nổi bật sự cần thiết của một khung phân tích phân biệt được
bối cảnh thay vì giả định hiệu ứng đồng nhất.

Ba dịch chuyển trên hội tụ vào một hàm ý chung: hiểu được khi nào và
trong điều kiện nào quốc tế hóa cải thiện hiệu quả doanh nghiệp không
còn là câu hỏi thuần lý thuyết mà là điều kiện tiên quyết để thiết kế
chính sách phát triển hữu hiệu cho khu vực châu Á--Thái Bình Dương đa
dạng về thể chế. Luận án này định vị trong khoảng giao đó, kết nối tranh
luận học thuật về quan hệ quốc tế hóa và hiệu quả với nhu cầu chính sách
đương đại bằng một khung đa tầng và bằng chứng cấp doanh nghiệp quy mô
lớn nhất từng có cho khu vực.

\begin{center}\rule{0.5\linewidth}{0.5pt}\end{center}

\hypertarget{hux1ed9i-tux1ee5-cux1ee7a-ba-lux1ef1c-dux1ecbch-chuyux1ec3n-vuxe0o-cuxe2u-hux1ecfi-cux1ea5p-doanh-nghiux1ec7p}{%
\subsection{1.1.2 Hội tụ của ba lực dịch chuyển vào câu hỏi cấp doanh
nghiệp}\label{hux1ed9i-tux1ee5-cux1ee7a-ba-lux1ef1c-dux1ecbch-chuyux1ec3n-vuxe0o-cuxe2u-hux1ecfi-cux1ea5p-doanh-nghiux1ec7p}}

Ba lực dịch chuyển vĩ mô nêu trên không tác động độc lập mà giao thoa
theo cách làm sắc nét thêm câu hỏi nghiên cứu ở cấp doanh nghiệp. Áp lực
việc làm ở các nền kinh tế biên chỉ được hóa giải nếu khu vực doanh
nghiệp tạo ra được việc làm năng suất cao, mà nguồn năng suất quan trọng
nhất ở các nền kinh tế nhỏ và mở chính là sự tham gia hiệu quả vào thị
trường quốc tế. Do đó, câu hỏi về điều kiện để quốc tế hóa cải thiện
hiệu quả doanh nghiệp gắn trực tiếp với năng lực hấp thụ lực lượng lao
động trẻ đang gia tăng nhanh trong khu vực.

Gradient thể chế (institutional gradient) là trục tổ chức xuyên suốt của
luận án. Khu vực châu Á và lân cận trải dài từ những nền kinh tế đổi mới
dẫn dắt với pháp quyền mạnh đến những nền kinh tế đảo nhỏ có hạ tầng thể
chế sơ khai, tạo ra một dải biến thiên thể chế hiếm có để kiểm định liệu
cùng một mức độ quốc tế hóa có tạo ra kết quả khác nhau theo chế độ thể
chế hay không. Chính dải biến thiên này, chứ không phải một quốc gia đơn
lẻ, là điều kiện cho phép phân biệt giữa giả thuyết về một quy luật phổ
quát và giả thuyết về một quan hệ có điều kiện.

Kỷ nguyên số bổ sung một chiều phức tạp mới. Khi công cụ số trở nên phổ
cập, ranh giới giữa chiều sâu năng lực công nghệ nội tại và mức độ chấp
nhận công cụ số bề mặt dễ bị xóa nhòa trong đo lường, dẫn đến nguy cơ
quy hai cơ chế khác nhau về một hiệu ứng đồng nhất. Vì hai loại năng lực
số này có thể tác động trái chiều theo chế độ thể chế, việc giữ chúng
tách biệt trở thành điều kiện tiên quyết để rút ra hàm ý chính sách
chính xác cho từng nhóm nền kinh tế trong khu vực.

\hypertarget{vux1ea5n-ux111ux1ec1-nghiuxean-cux1ee9u}{%
\section{1.2 Vấn đề nghiên
cứu}\label{vux1ea5n-ux111ux1ec1-nghiuxean-cux1ee9u}}

Mặc dù đã có hàng trăm nghiên cứu thực nghiệm, bằng chứng về mối quan hệ
giữa quốc tế hóa và hiệu quả hoạt động của doanh nghiệp vẫn chưa thống
nhất và tản mạn. Các phân tích tổng hợp (meta-analyses) về quan hệ quốc
tế hóa và hiệu quả liên tục cho thấy tác động tổng hợp là dương nhưng
nhỏ, hệ số tương quan tổng hợp r trong khoảng 0,07--0,13, đồng thời mức
độ dị biệt giữa các nghiên cứu rất cao với I² thường vượt 80\% (Bausch
\& Krist, 2007; Kirca et al., 2012; Marano et al., 2016; Wu et al.,
2022). Tình trạng dị biệt cao này không phải là tín hiệu của sai sót
phương pháp, mà là bằng chứng cho thấy quan hệ quốc tế hóa và hiệu quả
phụ thuộc sâu sắc vào điều kiện bối cảnh. Arte và Larimo (2022) lập luận
rằng dị biệt cao gợi ý sự cần thiết phải chuyển dịch từ mô hình đơn biến
sang khung đa tầng có tính đến thể chế, năng lực và đặc điểm nhà quản
trị.

Ba khoảng trống nghiên cứu quan trọng được xác định từ tổng quan tài
liệu. Khoảng trống thứ nhất là thiếu khung lý thuyết tích hợp đủ sức
giải thích dị biệt cao đó. Các nghiên cứu hiện tại thường chọn một hoặc
hai lý thuyết làm nền, dẫn đến mô hình không đủ đa tầng để nắm bắt sự
phức tạp thực tế. Mô hình Uppsala giải thích chi phí học hỏi nhưng không
đủ để xử lý vai trò của thể chế; lý thuyết dựa trên nguồn lực giải thích
sự khác biệt năng lực nhưng bỏ qua bối cảnh vĩ mô; lý thuyết thể chế
giải thích bối cảnh vĩ mô nhưng không tích hợp đặc điểm cá nhân nhà quản
trị. Không một lý thuyết nào đứng riêng có thể giải thích đầy đủ dị biệt
đó. Khoảng trống thứ hai là sự đồng nhất không hợp lý giữa hai cấu trúc
số khác nhau về bản chất: năng lực công nghệ, chiều sâu năng lực công
nghệ nhập vào tổ chức, và chấp nhận số, mức độ áp dụng công cụ số bề mặt
trong hoạt động kinh doanh. Nhiều nghiên cứu trước gộp hai cấu trúc này
thành một chỉ số chỉ số hợp thành duy nhất, điều này làm giảm độ chính
xác của kết luận và có thể che khuất các cơ chế tác động khác nhau
(Bharadwaj et al., 2013; Coltman et al., 2008; Verhoef et al., 2021).
Khoảng trống thứ ba là thiếu bằng chứng đa quốc gia quy mô lớn cho châu
Á. Phần lớn nghiên cứu hiện có tập trung vào một quốc gia hoặc một tiểu
vùng, và thậm chí khi có nghiên cứu đa quốc gia, phạm vi thường hẹp và
thiếu sự đa dạng thể chế đủ lớn để kiểm định các điều kiện biên có ý
nghĩa.

Ba khoảng trống trên cần được phát biểu sắc nét hơn để làm rõ rằng chúng
không phải là ba thiếu sót rời rạc mà là ba biểu hiện của cùng một hạn
chế nền tảng: tài liệu hiện hành thiếu một cấu trúc giải thích đủ phân
tầng để nắm bắt mức dị biệt cao của quan hệ quốc tế hóa và hiệu quả. Về
khoảng trống thứ nhất, không một lý thuyết đơn lẻ nào, dù là mô hình
Uppsala, lý thuyết dựa trên nguồn lực hay lý thuyết thể chế, cung cấp
được một lý giải đầy đủ cho sự dị biệt đó, bởi mỗi lý thuyết chỉ soi
sáng một loại dị biệt và bỏ trống các tầng còn lại. Hennart (2007) cảnh
báo rằng quan hệ giữa mức độ đa quốc gia và hiệu quả không phải là quy
luật phổ quát tự thân mà phụ thuộc vào bản chất tài sản đặc thù được
khai thác xuyên biên giới, hàm ý rằng việc kiểm định một dạng hàm duy
nhất cho mọi bối cảnh là không thỏa đáng về mặt lý thuyết.

Về khoảng trống thứ hai, sự đồng nhất không hợp lý giữa năng lực công
nghệ và mức độ chấp nhận số không chỉ là vấn đề tinh chỉnh đo lường mà
có hệ quả lý thuyết trực tiếp. Khi hai cấu trúc có chiều sâu tổ chức
khác nhau bị gộp thành một chỉ số tổng hợp duy nhất, hai cơ chế tác động
có thể trái chiều sẽ bị trung bình hóa, làm che khuất chính cơ chế mà
nhà quản trị và nhà hoạch định chính sách cần phân biệt khi quyết định
đầu tư vào chiều sâu năng lực hay vào công cụ giao diện.

Về khoảng trống thứ ba, sự thiếu đại diện của các nền kinh tế chuyển đổi
và các nền kinh tế đảo nhỏ đang phát triển không chỉ thu hẹp phạm vi mẫu
mà còn làm suy yếu khả năng kiểm định điều kiện biên của lý thuyết. Các
nỗ lực tái lặp quy mô lớn với chiến lược nhận diện chặt chẽ cho thấy
dạng hàm chữ U ngược suy yếu đáng kể, thậm chí biến mất, trong các mẫu
gộp xuyên quốc gia (Berry \& Kaul, 2016; Pisani et al., 2020), củng cố
lập luận rằng chính dị biệt bối cảnh, chứ không phải một quy luật phổ
quát, mới chi phối kết quả quốc tế hóa và hiệu quả. Đúng các địa điểm
cực trị về thể chế, nơi điều kiện biên dễ bộc lộ nhất, lại là nơi tài
liệu thiếu vắng bằng chứng nhất.

Luận án này được thiết kế để giải quyết ba khoảng trống đó bằng cách xây
dựng và kiểm định một khung giải thích đa tầng cho quan hệ quốc tế hóa
và hiệu quả tại châu Á, với thiết kế hỗn hợp tổng hợp--thực nghiệm kết
hợp phân tích tổng hợp và phân tích đa quốc gia quy mô lớn. Trọng tâm
khoa học của luận án là vai trò điều tiết của thể chế, năng lực số và
đặc điểm nhà quản trị, ba tầng điều tiết mà, theo lập luận của luận án,
cùng nhau tạo nên cấu trúc giải thích cần thiết cho dị biệt cao trong
quan hệ quốc tế hóa và hiệu quả.

\begin{center}\rule{0.5\linewidth}{0.5pt}\end{center}

\hypertarget{mux1ee5c-tiuxeau-nghiuxean-cux1ee9u}{%
\section{1.3 Mục tiêu nghiên
cứu}\label{mux1ee5c-tiuxeau-nghiuxean-cux1ee9u}}

Mục tiêu tổng quát của luận án là xây dựng và kiểm định một khung giải
thích đa tầng cho mối quan hệ giữa quốc tế hóa và hiệu quả hoạt động của
doanh nghiệp tại các quốc gia châu Á, trong đó làm rõ vai trò điều tiết
của thể chế, năng lực số và đặc điểm nhà quản trị.

Từ mục tiêu tổng quát đó, luận án triển khai bốn mục tiêu cụ thể. Mục
tiêu cụ thể thứ nhất là tổng hợp và cập nhật bằng chứng phân tích tổng
hợp về mối quan hệ quốc tế hóa và hiệu quả trong giai đoạn 1977--2026,
xác định tác động tổng hợp, mức độ dị biệt và các nhân tố điều tiết (các
biến điều tiết) chính giải thích dị biệt trong tài liệu. Mục tiêu cụ thể
thứ hai là kiểm định cơ chế cụ thể trong các bối cảnh quốc gia khác
nhau, bao gồm Singapore với tư cách nền kinh tế nhỏ mở tiên tiến, Việt
Nam với tư cách nền kinh tế chuyển đổi đang mở rộng, và Trung Quốc với
tư cách nền kinh tế mới nổi quy mô lớn, nhằm làm rõ cách thức các cơ chế
quốc tế hóa và hiệu quả vận hành khác nhau theo bối cảnh quốc gia. Mục
tiêu cụ thể thứ ba là mở rộng phân tích sang 50 nền kinh tế châu Á và
Thái Bình Dương (bao gồm 7 nền kinh tế Pacific SIDS) để kiểm định khả
năng khái quát hóa của các phát hiện trên phạm vi rộng, kế thừa và mở
rộng từ mẫu gộp 17 nền kinh tế trong nghiên cứu thành phần P1. Mục tiêu
cụ thể thứ tư là phân tích vai trò điều tiết của ba nhóm biến: thể chế
(bao gồm các sub-regime ICRV và các chỉ số rào cản thể chế), năng lực số
(phân biệt rõ technological capability index, TCI và chỉ số chấp nhận
số, digital adoption index, DAI), và đặc điểm nhà quản trị cấp cao (kinh
nghiệm và giới tính).

\begin{center}\rule{0.5\linewidth}{0.5pt}\end{center}

\hypertarget{cuxe2u-hux1ecfi-nghiuxean-cux1ee9u}{%
\section{1.4 Câu hỏi nghiên
cứu}\label{cuxe2u-hux1ecfi-nghiuxean-cux1ee9u}}

Câu hỏi nghiên cứu trung tâm của luận án là: Quốc tế hóa ảnh hưởng đến
hiệu quả hoạt động của doanh nghiệp ở châu Á như thế nào, và những điều
kiện nào, về thể chế, năng lực số và đặc điểm nhà quản trị, giải thích
sự dị biệt cao trong mối quan hệ đó?

Từ câu hỏi trung tâm đó, bốn câu hỏi nghiên cứu cụ thể được đặt ra. Câu
hỏi nghiên cứu thứ nhất (RQ1) là: Tác động tổng hợp của quốc tế hóa lên
hiệu quả hoạt động của doanh nghiệp trong tài liệu 1977--2026 là gì, mức
độ dị biệt ở mức nào, và những các biến điều tiết nào giải thích dị biệt
đó một cách có hệ thống? Câu hỏi nghiên cứu thứ hai (RQ2) là: Trong từng
bối cảnh quốc gia cụ thể bao gồm Singapore, Việt Nam và Trung Quốc, mối
quan hệ quốc tế hóa và hiệu quả vận hành theo những cơ chế đặc thù nào
và sự phi tuyến có biểu hiện ra sao? Câu hỏi nghiên cứu thứ ba (RQ3) là:
Trên phạm vi đa quốc gia châu Á bao gồm 50 nền kinh tế châu Á và Thái
Bình Dương (bao gồm 7 nền kinh tế Pacific SIDS), vai trò điều tiết của
chỉ số chấp nhận số (DAI) và năng lực công nghệ index (TCI) có khái quát
hóa được không, và hai cấu trúc này có thực sự vận hành theo các cơ chế
khác nhau không? Câu hỏi nghiên cứu thứ tư (RQ4) là: Vai trò của thể
chế, đo bằng chế độ con ICRV, cùng với sự phân biệt giữa TCI và DAI, và
đặc điểm nhà quản trị trong việc điều tiết mối quan hệ quốc tế hóa và
hiệu quả là gì, và sự tương tác ba chiều giữa ba tầng đó có ý nghĩa như
thế nào?

\begin{center}\rule{0.5\linewidth}{0.5pt}\end{center}

\hypertarget{ux111ux1ed1i-tux1b0ux1ee3ng-vuxe0-phux1ea1m-vi-nghiuxean-cux1ee9u}{%
\section{1.5 Đối tượng và phạm vi nghiên
cứu}\label{ux111ux1ed1i-tux1b0ux1ee3ng-vuxe0-phux1ea1m-vi-nghiuxean-cux1ee9u}}

\hypertarget{ux111ux1ed1i-tux1b0ux1ee3ng-nghiuxean-cux1ee9u}{%
\subsection{1.5.1 Đối tượng nghiên
cứu}\label{ux111ux1ed1i-tux1b0ux1ee3ng-nghiuxean-cux1ee9u}}

Đối tượng nghiên cứu của luận án là mối quan hệ giữa quốc tế hóa và hiệu
quả hoạt động của doanh nghiệp (internationalization--firm performance
relationship) cùng với các yếu tố điều tiết của mối quan hệ đó, bao gồm
bối cảnh thể chế, năng lực số và đặc điểm nhà quản trị cấp cao. Đơn vị
phân tích là doanh nghiệp ở cấp độ tổ chức (firm-level analysis), không
phải cấp độ ngành hay cấp độ quốc gia.

\hypertarget{phux1ea1m-vi-khuxf4ng-gian}{%
\subsection{1.5.2 Phạm vi không gian}\label{phux1ea1m-vi-khuxf4ng-gian}}

Phạm vi không gian chính của nghiên cứu bao gồm \textbf{50 nền kinh tế
châu Á và Thái Bình Dương} có dữ liệu trong Khảo sát Doanh nghiệp của
Ngân hàng Thế giới (World Bank Enterprise Surveys, WBES), được phân loại
và tổ chức theo 6 chế độ con trong khung ICRV. Sáu chế độ con bao gồm:
(i) tiên tiến đổi mới, gồm Singapore, Hong Kong SAR, Hàn Quốc, Đài Loan,
Israel và Nhật Bản; (ii) tiên tiến tài nguyên-Driven, đại diện là các
nền kinh tế GCC; (iii) trung bình cao, đại diện là Trung Quốc và
Malaysia; (iv) chuyển đổi thu nhập trung bình thấp, bao gồm Việt Nam và
các nền kinh tế ASEAN đang phát triển; (v) đang nổi, bao gồm các nền
kinh tế kém phát triển nhất châu Á như Campuchia, Nepal và Afghanistan;
và (vi) SIDS (Small Island Developing States), bao gồm 7 nền kinh tế đảo
nhỏ Thái Bình Dương (mẫu chính P8; mở rộng 9 nền gồm Comoros và
Timor-Leste khi kiểm định độ vững) được đưa vào phân tích để kiểm định
các điều kiện biên cực đoan của quan hệ quốc tế hóa và hiệu quả. Nhật
Bản được Khảo sát Doanh nghiệp của Ngân hàng Thế giới đưa vào lần đầu
năm 2025 (N = 2.168) và là thành viên đầy đủ của Nhóm I trong khung phân
tích, đồng thời là đối tượng của nghiên cứu thành phần P10.

\hypertarget{phux1ea1m-vi-thux1eddi-gian}{%
\subsection{1.5.3 Phạm vi thời gian}\label{phux1ea1m-vi-thux1eddi-gian}}

Phạm vi thời gian của nghiên cứu có hai mức. Đối với giai đoạn phân tích
tổng hợp, luận án tổng hợp tài liệu trong giai đoạn 1977--2026, bao quát
toàn bộ chương trình nghiên cứu về quan hệ quốc tế hóa và hiệu quả từ
khi các nghiên cứu thực nghiệm đầu tiên xuất hiện đến thời điểm hoàn
thành luận án. Đối với giai đoạn phân tích thực nghiệm đa quốc gia, luận
án sử dụng dữ liệu WBES trong giai đoạn 2006--2026, bao gồm 88.869 doanh
nghiệp với 103 cặp nền kinh tế-năm. Dữ liệu WBES từ ba thế hệ lược đồ
(PICS3 giai đoạn 2009--2012, Standardized giai đoạn 2013--2017, và
BREADY/BEE giai đoạn 2018--2025) đã được hòa hợp để đảm bảo tính nhất
quán trong đo lường biến qua các đợt khảo sát.

\begin{center}\rule{0.5\linewidth}{0.5pt}\end{center}

\hypertarget{giux1edbi-hux1ea1n-phux1ea1m-vi}{%
\subsection{1.5.4 Giới hạn phạm
vi}\label{giux1edbi-hux1ea1n-phux1ea1m-vi}}

Để bảo đảm tính khả thi và độ chặt chẽ, luận án xác định rõ một số giới
hạn phạm vi có chủ đích. Thứ nhất, đơn vị phân tích là doanh nghiệp ở
cấp độ tổ chức, do đó luận án không nhằm giải thích kết quả ở cấp ngành
hay cấp quốc gia, mặc dù bối cảnh ngành và quốc gia được đưa vào như
biến điều tiết và biến kiểm soát. Thứ hai, quốc tế hóa được vận hành hóa
chủ yếu qua tỷ lệ doanh thu nước ngoài trên tổng doanh thu, phù hợp với
cấu trúc dữ liệu khảo sát; các phương thức quốc tế hóa khác như đầu tư
trực tiếp ra nước ngoài hay liên doanh không nằm trong trọng tâm đo
lường, dù được thừa nhận về mặt khái niệm.

Thứ ba, dữ liệu thực nghiệm là chuỗi cắt ngang lặp lại theo từng đợt
khảo sát, không phải dữ liệu bảng theo dõi cùng một doanh nghiệp qua
thời gian, nên các suy luận về động thái theo thời gian được giới hạn ở
so sánh giữa các đợt khảo sát thay vì theo dõi quỹ đạo của từng doanh
nghiệp. Thứ tư, do ràng buộc của bộ dữ liệu, chỉ số chấp nhận số được đo
bằng các chỉ báo nền tảng ở giao diện ngoại tại; đây là ràng buộc đo
lường có chủ đích chứ không phải thiếu sót khái niệm, và được xử lý nhất
quán trong toàn bộ luận án. Cuối cùng, phạm vi không gian giới hạn ở khu
vực châu Á và lân cận; việc khái quát hóa các phát hiện sang các khu vực
khác cần được kiểm chứng thêm và được nêu như một hướng nghiên cứu tiếp
theo trong Chương 5.

\hypertarget{phux1b0ux1a1ng-phuxe1p-luux1eadn-vuxe0-phux1b0ux1a1ng-phuxe1p-nghiuxean-cux1ee9u}{%
\section{1.6 Phương pháp luận và phương pháp nghiên
cứu}\label{phux1b0ux1a1ng-phuxe1p-luux1eadn-vuxe0-phux1b0ux1a1ng-phuxe1p-nghiuxean-cux1ee9u}}

Luận án áp dụng thiết kế nghiên cứu hỗn hợp tổng hợp--thực nghiệm, sự
kết hợp giữa phân tích tổng hợp và phân tích thực nghiệm đa quốc gia,
nhằm trả lời bốn câu hỏi nghiên cứu đã đặt ra. Thiết kế này không chỉ
đơn thuần là sử dụng hai phương pháp song song, mà tạo ra một đối thoại
hai chiều: phân tích tổng hợp cung cấp tổng quan có hệ thống về trạng
thái bằng chứng và xác định các biến điều tiết quan trọng trong tài
liệu, trong khi phân tích thực nghiệm kiểm định trực tiếp các giả thuyết
phát sinh từ khung lý thuyết trong bối cảnh châu Á. Sự tương tác này
tăng cường đáng kể sức mạnh đóng góp so với việc chỉ sử dụng một phương
pháp duy nhất (Borenstein et al., 2009; Hunter \& Schmidt, 2004).

Về phương pháp cụ thể, giai đoạn phân tích tổng hợp sử dụng mô hình hiệu
ứng ngẫu nhiên với phép biến đổi \(z\) của Fisher, phân tích nhóm con
theo chế độ thể chế, và các kiểm định lệch lạc công bố bao gồm kiểm định
Egger và kiểm định tương quan hạng Begg--Mazumdar. Giai đoạn thực nghiệm
sử dụng hồi quy bình phương nhỏ nhất với sai số chuẩn vững HC1 và HC3,
mô hình bậc hai để kiểm định quan hệ phi tuyến dạng chữ U ngược, tương
tác hai chiều và ba chiều để kiểm định điều tiết, cùng với hiệu ứng cố
định theo quốc gia và năm. Kiểm định chữ U của Lind và Mehlum được áp
dụng để xác nhận dạng chữ U ngược và xác định điểm uốn trong quan hệ phi
tuyến (Haans et al., 2016; Lind \& Mehlum, 2010).

\begin{center}\rule{0.5\linewidth}{0.5pt}\end{center}

\hypertarget{giux1ea3-thuyux1ebft-khoa-hux1ecdc-tuxf3m-tux1eaft}{%
\section{1.7 Giả thuyết khoa học (tóm
tắt)}\label{giux1ea3-thuyux1ebft-khoa-hux1ecdc-tuxf3m-tux1eaft}}

Dựa trên khung lý thuyết tích hợp đa tầng bao gồm mô hình Uppsala, lý
thuyết dựa trên nguồn lực, lý thuyết thể chế, Lý thuyết thượng tầng quản
trị và lăng kính năng lực số, luận án đề xuất hệ giả thuyết khoa học với
sáu giả thuyết trọng tâm. Quan hệ nền giữa quốc tế hóa và hiệu quả hoạt
động doanh nghiệp mang dạng phi tuyến, cụ thể là dạng chữ U ngược, với
chi phí thâm nhập và thích nghi cao ở giai đoạn đầu và lợi ích quy mô,
học hỏi và đa dạng hóa thị trường ở giai đoạn sau, đến khi chi phí điều
phối vượt quá lợi ích sẽ dẫn đến sự suy giảm hiệu quả. năng lực công
nghệ index (TCI) điều tiết dương quan hệ quốc tế hóa và hiệu quả bằng
cách mở rộng biên thích nghi và tăng cường khả năng khai thác lợi thế
cạnh tranh ở thị trường quốc tế. Chỉ số chấp nhận số (DAI), được hiểu là
biến đại diện foundational chấp nhận số, có tác động điều tiết không đơn
điệu và phụ thuộc vào chế độ thể chế: ở các nền kinh tế có thể chế tốt,
DAI khuếch đại lợi ích của quốc tế hóa; ở các nền kinh tế với thể chế
yếu, DAI có thể khuếch đại chi phí điều phối khi mức độ quốc tế hóa vượt
ngưỡng. Bối cảnh thể chế, đo bằng chế độ con ICRV, điều tiết cả dạng và
độ lớn của quan hệ quốc tế hóa và hiệu quả, tạo ra gradient từ tác động
dương lớn ở tiên tiến đổi mới đến tác động âm ở đang nổi và SIDS
extreme. Đặc điểm nhà quản trị cấp cao, cụ thể là kinh nghiệm quốc tế và
đa dạng giới tính, điều tiết dương quan hệ quốc tế hóa và hiệu quả thông
qua cơ chế tăng cường nhận diện cơ hội và quản trị rủi ro. Ngoài sáu giả
thuyết trọng tâm H1--H6, luận án phát biểu thêm giả thuyết điều kiện
biên H1b: tại các nền kinh tế SIDS (ICRV Nhóm VI), nơi ba điều kiện cấu
trúc tối thiểu đồng thời bị vi phạm, thị trường nội địa quá nhỏ, chi phí
thương mại cực cao, và thiếu hỗ trợ thể chế, quan hệ quốc tế hóa và hiệu
quả chuyển từ dạng chữ U ngược sang quan hệ \textbf{âm đơn điệu không có
điểm uốn}, được gọi là gánh nặng quốc tế hóa bắt buộc (Forced
Internationalization Penalty, FIP). Chi tiết toàn bộ hệ giả thuyết
H1--H6 và H1b được trình bày đầy đủ trong Chương 2 của luận án.

\begin{center}\rule{0.5\linewidth}{0.5pt}\end{center}

\hypertarget{uxfd-nghux129a-khoa-hux1ecdc-vuxe0-thux1ef1c-tiux1ec5n}{%
\section{1.8 Ý nghĩa khoa học và thực
tiễn}\label{uxfd-nghux129a-khoa-hux1ecdc-vuxe0-thux1ef1c-tiux1ec5n}}

\hypertarget{ux111uxf3ng-guxf3p-luxfd-thuyux1ebft}{%
\subsection{1.8.1 Đóng góp lý
thuyết}\label{ux111uxf3ng-guxf3p-luxfd-thuyux1ebft}}

Về mặt lý thuyết, luận án đóng góp vào tài liệu kinh doanh quốc tế theo
ba hướng chính. Hướng thứ nhất là tích hợp bốn lý thuyết kinh điển, mô
hình Uppsala, lý thuyết dựa trên nguồn lực, lý thuyết thể chế và Lý
thuyết thượng tầng quản trị, vào một khung giải thích đa tầng thống
nhất, bổ sung bởi lăng kính năng lực số để phản ánh chuyển đổi số đương
đại. Khung này không chỉ là sự ghép nối cơ học các lý thuyết mà là một
kiến trúc lý thuyết có tính nhất quán nội tại, trong đó mỗi tầng giải
thích một loại dị biệt khác nhau trong quan hệ quốc tế hóa và hiệu quả.
Hướng thứ hai là phân tách rõ ràng năng lực công nghệ (TCI) và chỉ số
chấp nhận số (DAI) như hai cấu trúc lý thuyết độc lập, dựa trên các nền
tảng khái niệm khác nhau từ Bharadwaj et al.~(2013), Verhoef et
al.~(2021), và Coltman et al.~(2008). Sự phân tách này bổ sung vào thiếu
hụt khái niệm quan trọng trong tài liệu về số hóa và quốc tế hóa. Hướng
thứ ba là tái định vị mô hình Uppsala trong bối cảnh kỷ nguyên AI bằng
cách mở rộng cơ chế học hỏi để bao gồm ``data-augmented learning'', phù
hợp với xu hướng digitally-facilitated quốc tế hóa (Banalieva \&
Dhanaraj, 2019; Yang et al., 2025). Hướng thứ tư là đề xuất và đặt tên
lý thuyết cho \textbf{gánh nặng quốc tế hóa bắt buộc (Forced
Internationalization Penalty, FIP)}, một cấu trúc lý thuyết mới chưa
được định danh trong tài liệu kinh doanh quốc tế, như điều kiện biên cực
đoan của mô hình chữ U ngược khi cả ba tầng của khung CDCM (thể chế,
năng lực, quản trị) đồng thời ở mức tối thiểu. Bằng chứng thực nghiệm từ
bảy nền kinh tế đảo nhỏ Thái Bình Dương xác nhận quan hệ âm đơn điệu
không có điểm uốn, mở ra hướng nghiên cứu mới về điều kiện biên của lý
thuyết quốc tế hóa.

Bốn hướng đóng góp lý thuyết trên hội tụ vào một kiến trúc thống nhất là
khung điều tiết bối cảnh, năng lực và quản trị
(Country--Digital--Capability Moderation, CDCM), trong đó mỗi tầng giải
thích một loại dị biệt riêng của quan hệ quốc tế hóa và hiệu quả thay vì
cạnh tranh loại trừ lẫn nhau. Tầng vĩ mô được vận hành hóa bằng khung
biến thiên chế độ bối cảnh thể chế (Institutional Context Regime
Variation, ICRV), công cụ phân loại sáu nhóm cho phép chuyển một
gradient thể chế liên tục thành một biến điều tiết kiểm định được ở cấp
doanh nghiệp. Đóng góp của ICRV không nằm ở việc thêm một biến phân
loại, mà ở việc cung cấp một cách hệ thống để dự đoán vị trí điểm uốn
dịch chuyển theo chế độ thể chế, qua đó biến luận điểm ``bối cảnh quan
trọng'' từ một nhận định định tính thành một giả thuyết định lượng có
thể bác bỏ.

Đỉnh của kiến trúc lý thuyết là việc định danh gánh nặng quốc tế hóa bắt
buộc (Forced Internationalization Penalty, FIP) như một điều kiện biên
cực trị của mô hình chữ U ngược. FIP phát biểu rằng khi cả ba tầng của
khung CDCM đồng thời ở mức tối thiểu, quan hệ quốc tế hóa và hiệu quả
không chỉ mất phần lợi ích giai đoạn đầu mà chuyển hẳn sang dạng âm đơn
điệu không có điểm uốn. Việc định danh một điều kiện biên như vậy có giá
trị lý thuyết vượt ra ngoài một phát hiện thực nghiệm đơn lẻ, bởi nó xác
định ranh giới ứng dụng của cả một họ lý thuyết quốc tế hóa vốn ngầm giả
định những điều kiện cấu trúc tối thiểu luôn được thỏa mãn.

\hypertarget{ux111uxf3ng-guxf3p-phux1b0ux1a1ng-phuxe1p}{%
\subsection{1.8.2 Đóng góp phương
pháp}\label{ux111uxf3ng-guxf3p-phux1b0ux1a1ng-phuxe1p}}

Về mặt phương pháp, luận án đóng góp theo hai hướng. Hướng thứ nhất là
thiết kế hỗn hợp tổng hợp--thực nghiệm kết hợp phân tích tổng hợp và
phân tích đa quốc gia, tạo ra khả năng đối chiếu giữa bằng chứng tổng
hợp từ tài liệu toàn cầu và bằng chứng thực nghiệm từ bối cảnh châu Á.
Hướng thứ hai là phát triển và áp dụng khung phân loại thể chế ICRV cho
50 nền kinh tế, cung cấp một công cụ phân loại mới có ứng dụng tiềm năng
trong nhiều nghiên cứu tương lai về moderating effect của thể chế trong
tài liệu kinh doanh quốc tế.

\hypertarget{ux111uxf3ng-guxf3p-thux1ef1c-tiux1ec5n}{%
\subsection{1.8.3 Đóng góp thực
tiễn}\label{ux111uxf3ng-guxf3p-thux1ef1c-tiux1ec5n}}

Về mặt thực tiễn, luận án cung cấp căn cứ có hệ thống cho hai nhóm đối
tượng sử dụng chính. Đối với doanh nghiệp, kết quả phân tích của luận án
cho thấy tồn tại ngưỡng tối ưu trong mối quan hệ giữa mức độ xuất khẩu
và hiệu quả hoạt động, và ngưỡng này thay đổi theo bối cảnh thể chế và
năng lực số của từng doanh nghiệp. Nhà quản trị cần nhận thức được
ngưỡng này và đầu tư có chiến lược vào năng lực công nghệ thay vì chỉ
tập trung vào chấp nhận số bề mặt. Đối với nhà hoạch định chính sách,
bằng chứng từ luận án gợi ý rằng cải cách thể chế và hỗ trợ chuyển đổi
số có thể có hiệu ứng bổ trợ tích cực, được nghiên cứu về 17 nền kinh tế
châu Á mới nổi đặt tên là ``hiệu ứng lá chắn số'', trong đó năng lực số
làm giảm tác động tiêu cực của rào cản thể chế đối với doanh nghiệp quốc
tế hóa (nghiên cứu thành phần P1). Cơ chế này được luận án kiểm định
trên phạm vi rộng hơn với 50 nền kinh tế châu Á và Thái Bình Dương.

\begin{center}\rule{0.5\linewidth}{0.5pt}\end{center}

Hàm ý chính sách của luận án mang tính phân biệt theo chế độ thể chế
thay vì đưa ra một khuyến nghị đồng nhất cho toàn khu vực. Ở các nền
kinh tế thuộc nhóm thể chế mạnh, nơi quan hệ quốc tế hóa và hiệu quả duy
trì độ dốc dương trong phần lớn miền quan sát, chính sách hỗ trợ nên tập
trung vào việc gỡ bỏ các trần mở rộng và nâng trần năng lực hấp thụ công
nghệ. Ở các nền kinh tế chuyển đổi thu nhập trung bình--thấp, nơi đường
cong chữ U ngược biểu hiện rõ nét nhất, hàm ý trọng tâm là nhận diện
ngưỡng tối ưu và tránh mở rộng vượt ngưỡng khi năng lực điều phối chưa
theo kịp.

Ở cực còn lại của gradient, các nền kinh tế đảo nhỏ đang phát triển đòi
hỏi một logic chính sách khác hẳn. Khi quốc tế hóa là gánh nặng sinh tồn
chứ không phải lựa chọn chiến lược dựa trên lợi thế cạnh tranh, việc
khuyến khích tăng cường độ xuất khẩu một cách máy móc có thể phản tác
dụng; ưu tiên hợp lý hơn là khắc phục các ràng buộc cấu trúc về chi phí
thương mại và hỗ trợ thể chế trước khi kỳ vọng quốc tế hóa cải thiện
hiệu quả. Sự phân biệt này cho thấy giá trị thực tiễn cốt lõi của luận
án nằm ở việc cung cấp một bản đồ chính sách theo chế độ thể chế, thay
vì một đơn thuốc duy nhất áp dụng cho mọi nền kinh tế bất kể vị trí trên
gradient.

\hypertarget{tuxednh-mux1edbi-cux1ee7a-ux111ux1ec1-tuxe0i}{%
\section{1.9 Tính mới của đề
tài}\label{tuxednh-mux1edbi-cux1ee7a-ux111ux1ec1-tuxe0i}}

Luận án có sáu điểm mới so với các công trình nghiên cứu trước đây. Điểm
mới thứ nhất là đề xuất khung tích hợp đa tầng cho quan hệ quốc tế hóa
và hiệu quả trong bối cảnh châu Á, kết hợp bốn lý thuyết kinh điển với
lăng kính năng lực số thành một kiến trúc lý thuyết thống nhất và có khả
năng kiểm định. Không có nghiên cứu trước nào trong bối cảnh châu Á xây
dựng khung tích hợp đầy đủ ba tầng điều tiết này theo cách có hệ thống.
Điểm mới thứ hai là phân tách TCI và chỉ số chấp nhận số (DAI) như hai
cấu trúc không trùng lặp, kiểm định chúng riêng biệt trong cùng một mô
hình để làm rõ cơ chế tác động khác nhau của hai loại năng lực số đó. Sự
phân tách này khắc phục hạn chế của các nghiên cứu trước thường gộp
chung hai cấu trúc có bản chất khác nhau. Điểm mới thứ ba là phát triển
và áp dụng khung phân loại thể chế ICRV với 6 chế độ con cho 50 nền kinh
tế, bao gồm cả 7 nền kinh tế đảo nhỏ Thái Bình Dương liền kề như
boundary case cực đoan, đây là lần đầu tiên một hệ thống phân loại đủ
mịn như vậy được áp dụng trong nghiên cứu thực nghiệm về quan hệ quốc tế
hóa và hiệu quả tại châu Á. Điểm mới thứ tư là mẫu gộp dữ liệu lớn nhất
từ trước đến nay cho nghiên cứu đa quốc gia về quan hệ quốc tế hóa và
hiệu quả tại châu Á, với 88.869 doanh nghiệp từ 50 nền kinh tế châu Á và
Thái Bình Dương, mở rộng đáng kể từ mẫu gộp 17 nước (nghiên cứu thành
phần P1) là nền tảng hiện có lớn nhất trước khi luận án hoàn thành. Điểm
mới thứ năm là kiểm định temporal stability của điểm uốn trong quan hệ
quốc tế hóa và hiệu quả tại Trung Quốc, cho phép đánh giá liệu cấu trúc
phi tuyến đó có ổn định qua hơn một thập kỷ chuyển đổi số hay không. Kết
quả phân tích của luận án cho thấy điểm uốn tương đối ổn định qua hai
thời điểm quan sát, cung cấp bằng chứng về tính bền vững của cấu trúc
quốc tế hóa và hiệu quả phi tuyến ở Trung Quốc dù bối cảnh kinh tế có
nhiều thay đổi. Điểm mới thứ sáu là phát hiện và đặt tên lý thuyết cho
\textbf{gánh nặng quốc tế hóa bắt buộc (Forced Internationalization
Penalty, FIP)} tại bảy nền kinh tế đảo nhỏ Thái Bình Dương, lần đầu tiên
trong tài liệu kinh doanh quốc tế, một trường hợp quan hệ quốc tế hóa và
hiệu quả âm đơn điệu không có điểm uốn được ghi nhận và biện minh lý
thuyết bằng ba điều kiện cấu trúc vi phạm đồng thời, mở ra hướng nghiên
cứu mới về điều kiện biên của các lý thuyết quốc tế hóa phổ quát.

\begin{center}\rule{0.5\linewidth}{0.5pt}\end{center}

Sáu điểm mới trên có thể được cô đọng thành một tuyên bố tính mới trung
tâm: luận án là công trình đầu tiên trong tài liệu kinh doanh quốc tế
chuyển luận điểm ``quan hệ quốc tế hóa và hiệu quả phụ thuộc bối cảnh''
từ một nhận định định tính thành một cấu trúc kiểm định được trên toàn
bộ gradient thể chế của một khu vực, đồng thời xác định và định danh
điểm mà tại đó dạng hàm chữ U ngược không còn ứng dụng được. Khác với
các nghiên cứu trước thường dừng lại ở việc chứng minh thể chế điều tiết
độ lớn của hiệu ứng, luận án chứng minh thể chế có thể điều tiết cả dấu
của hiệu ứng, qua đó nâng luận điểm từ ``phụ thuộc bối cảnh'' lên ``có
điều kiện theo bối cảnh''. Tính mới này được nâng đỡ bởi sự phân tách
hai cấu trúc năng lực số và bởi nền dữ liệu đa quốc gia rộng nhất từng
có cho khu vực, hai yếu tố cùng bảo đảm rằng tuyên bố trên dựa trên độ
thuần khiết cấu trúc và độ bao phủ thể chế chứ không phải trên một mẫu
hẹp hoặc một đo lường gộp.

\hypertarget{cux1ea5u-truxfac-cux1ee7a-luux1eadn-uxe1n}{%
\section{1.10 Cấu trúc của luận
án}\label{cux1ea5u-truxfac-cux1ee7a-luux1eadn-uxe1n}}

Luận án được tổ chức thành năm chương theo chuẩn luận án tiến sĩ truyền
thống tại Trường Kinh tế, Đại học Cần Thơ. Chương 1, Giới thiệu, đặt vấn
đề nghiên cứu, xác định khoảng trống, trình bày mục tiêu, câu hỏi nghiên
cứu, phạm vi, phương pháp luận, giả thuyết khoa học tóm tắt, ý nghĩa,
tính mới và cấu trúc của luận án. Chương 2, Cơ sở lý thuyết và tổng quan
tài liệu, xây dựng khung lý thuyết tích hợp đa tầng, tổng quan bằng
chứng thực nghiệm đã có, và phát triển hệ giả thuyết H1--H6 dựa trên
khoảng trống và logic lý thuyết. Chương 3, Phương pháp nghiên cứu, trình
bày chi tiết thiết kế nghiên cứu, nguồn dữ liệu, cách đo lường biến, mô
hình phân tích và chiến lược kiểm định độ vững. Chương 4, Kết quả nghiên
cứu, trình bày kết quả ở ba cấp độ phân tích: phân tích tổng hợp cập
nhật 1977--2026, phân tích theo bối cảnh quốc gia cụ thể, và phân tích
đa quốc gia tổng hợp trên mẫu gộp đầy đủ 50 nền kinh tế. Chương 5, Thảo
luận và kết luận, tích hợp các phát hiện từ Chương 4 vào khung lý thuyết
ở Chương 2, thảo luận đóng góp khoa học, hàm ý quản trị và chính sách,
hạn chế của nghiên cứu và hướng nghiên cứu tiếp theo.

\textbf{Lộ trình tám nghiên cứu thành phần (P3--P10):} Luận án tích hợp
tám nghiên cứu thành phần được công bố hoặc đang trong quá trình xem xét
tại các tạp chí ISI/Scopus. \textbf{P3} (Việt Nam, JABS, đang phản
biện): kiểm định chữ U ngược tại Nhóm IV ICRV, điểm uốn = 39,3--46,2\%
FSTS, TCI là yếu tố nâng mặt bằng có ý nghĩa. \textbf{P4} (Singapore,
MIR, đang phản biện): kiểm định tại Nhóm I ICRV, điểm uốn ≈ 88,6\%
(khoảng tin cậy rất rộng và vượt ngưỡng khả thi của FSTS, nên không định
vị tin cậy, xem Mục 4.3), DAI Tầng 1 và 2 khuếch đại đường cong quốc tế
hóa và hiệu quả. \textbf{P5} (Trung Quốc, IJOEM, đang phản biện): kiểm
định temporal stability, điểm uốn ổn định qua 2012--2024 (≈ 48,78\%,
Paternoster z = 0,82, p = ,412). \textbf{P6} (meta-analysis ba tầng,
IBR, đang phản biện): k = 238 nghiên cứu, r̄ = 0,074, I² = 62,4\%,
gradient ICRV Q\_M = 17,35 (df = 4, p = ,002). \textbf{P7} (50 nền kinh
tế, mục tiêu JIBS): N = 81.022 quan sát, điểm uốn gộp 43,6\% ở mô hình
đầy đủ, chữ U ngược rõ nét nhất ở nhóm chuyển đổi thu nhập trung
bình--thấp (43,0\%), kiểm định tích hợp CDCM đầu tiên ở quy mô liên quốc
gia lớn. \textbf{P8} (SIDS, World Development, đang phản biện vòng hai):
điều kiện biên FIP, quan hệ âm đơn điệu
(\(\hat{\beta}_\text{FSTS} = -1{,}339\), \(p < {,}001\)), xác nhận H1b
tại bảy nền kinh tế đảo nhỏ Thái Bình Dương. \textbf{P9} (Ấn Độ,
MIR/IJOEM, đang bình duyệt): độ bền ngưỡng quốc tế hóa và hiệu quả qua
ba đợt WBES 2014--2025 (\textasciitilde29.000 doanh nghiệp), phát hiện
\textbf{tan biến} điểm uốn chữ U ngược dưới biến động thể chế cực đoan
(TP 61,8\% rồi 40,7\% sang tan biến; Paternoster \(z = -7{,}94\),
\(p < {,}0001\)), kèm tương tác DAI Tầng 2/UPI âm. Kết quả P3--P9 được
trình bày chi tiết trong Chương 4 và tích hợp vào bàn luận CDCM trong
Chương 5. \textbf{P10} (Nhật Bản, mục tiêu Asian Business \&
Management): khai thác đợt WBES khai mở của Nhật Bản 2025 (N = 2.168)
kiểm định dự đoán biên trên của gradient thể chế, quan hệ quốc tế hóa và
hiệu quả gần tuyến tính dương, không có điểm uốn trong miền quan sát;
TCI và DAI nâng mặt bằng; phần bù trường thọ doanh nghiệp.

\begin{center}\rule{0.5\linewidth}{0.5pt}\end{center}

\emph{Chú thích về ký hiệu và thuật ngữ: Trong toàn bộ luận án, ``quốc
tế hóa'' (internationalization) được đo chủ yếu bằng tỷ lệ doanh thu
nước ngoài trên tổng doanh thu (FSTS) hoặc export intensity; ``hiệu quả
hoạt động'' (firm performance) được đo chủ yếu bằng năng suất lao động
(labor productivity), bổ sung bằng ROS và sales growth trong các kiểm
định độ vững; TCI là năng lực công nghệ index; DAI là chỉ số chấp nhận
số (digital adoption index), được hiểu như biến đại diện foundational
chấp nhận số; ICRV là khung phân loại thể chế 6 nhóm được phát triển cho
luận án này; WBES là World Bank Enterprise Surveys.}

\newpage

\hypertarget{chux1b0ux1a1ng-2-tux1ed5ng-quan-tuxe0i-liux1ec7u}{%
\chapter{CHƯƠNG 2: TỔNG QUAN TÀI
LIỆU}\label{chux1b0ux1a1ng-2-tux1ed5ng-quan-tuxe0i-liux1ec7u}}

Chương 2 có nhiệm vụ hệ thống hóa các khái niệm cốt lõi và lý thuyết nền
tảng, lược khảo và đánh giá các công trình thực nghiệm đã công bố về mối
quan hệ giữa quốc tế hóa và hiệu quả hoạt động doanh nghiệp, xác định
những khoảng trống còn tồn tại trong tài liệu, và từ đó xây dựng khung
khái niệm tích hợp cùng hệ giả thuyết nghiên cứu H1--H6 cho luận án.
Logic của chương đi từ nền tảng khái niệm đến lý thuyết, từ bằng chứng
thực nghiệm đến khoảng trống, và từ khoảng trống đến mô hình nghiên cứu
và giả thuyết.

\begin{center}\rule{0.5\linewidth}{0.5pt}\end{center}

\hypertarget{cuxe1c-khuxe1i-niux1ec7m-cux1ed1t-luxf5i}{%
\section{2.1 Các khái niệm cốt
lõi}\label{cuxe1c-khuxe1i-niux1ec7m-cux1ed1t-luxf5i}}

\hypertarget{quux1ed1c-tux1ebf-huxf3a-doanh-nghiux1ec7p}{%
\subsection{2.1.1 Quốc tế hóa doanh
nghiệp}\label{quux1ed1c-tux1ebf-huxf3a-doanh-nghiux1ec7p}}

Quốc tế hóa doanh nghiệp (internationalization) được hiểu là mức độ và
phạm vi tham gia của một doanh nghiệp vào các hoạt động kinh tế có yếu
tố quốc tế, bao gồm xuất khẩu hàng hóa và dịch vụ, đầu tư trực tiếp ra
nước ngoài, liên doanh quốc tế, cấp phép công nghệ xuyên biên giới, và
thiết lập mạng lưới phân phối ở thị trường nước ngoài (Sullivan, 1994;
Hitt et al., 2006). Khái niệm này không đồng nhất với ``toàn cầu hóa'',
vốn mang nghĩa một xu hướng kinh tế vĩ mô, mà là một thuộc tính đo lường
được ở cấp độ doanh nghiệp, phản ánh chiến lược mở rộng quốc tế và mức
độ phụ thuộc vào thị trường nước ngoài. Nền tảng lý thuyết của lĩnh vực
này bắt nguồn từ các công trình kinh điển về doanh nghiệp đa quốc gia và
đầu tư trực tiếp nước ngoài: lý thuyết vòng đời sản phẩm (Vernon, 1966),
lý thuyết lợi thế độc quyền của doanh nghiệp (Hymer, 1976; Caves, 1971),
và lý thuyết nội bộ hóa chi phí giao dịch (Buckley \& Casson, 1976), về
sau được tổng hợp trong khung chiết trung OLI (Dunning, 1988).

Trong tài liệu, mức độ quốc tế hóa được vận hành hóa bằng nhiều thước đo
khác nhau. Sullivan (1994) đề xuất một chỉ số tổng hợp đa chiều bao gồm
tỷ trọng doanh thu từ nước ngoài trên tổng doanh thu (foreign sales to
total sales, FSTS), tỷ trọng tài sản ở nước ngoài trên tổng tài sản
(foreign assets to total assets, FATA), và phạm vi địa lý (geographic
scope). Trong những nghiên cứu sử dụng dữ liệu doanh nghiệp từ các nền
kinh tế đang phát triển và từ khảo sát doanh nghiệp của Ngân hàng Thế
giới (World Bank Enterprise Surveys, WBES), thước đo phổ biến nhất là
cường độ xuất khẩu (export intensity), tức là tỷ lệ phần trăm doanh thu
xuất khẩu trên tổng doanh thu, hay tương đương với FSTS theo nghĩa rộng
(Johanson \& Vahlne, 1977; Lu \& Beamish, 2004). Luận án này sử dụng
FSTS và bình phương FSTS (FSTS²) làm thước đo chính về mức độ quốc tế
hóa, với lý do phù hợp với cấu trúc dữ liệu WBES và nhất quán với các
nghiên cứu quốc tế hóa và hiệu quả sử dụng bộ dữ liệu tương tự (Bhandari
et al., 2023; nghiên cứu thành phần P1).

Điều quan trọng cần lưu ý là quốc tế hóa không phải là một trạng thái
nhị phân, mà là một biến liên tục biểu thị cường độ tham gia thị trường
quốc tế, có thể thay đổi theo thời gian và phản ánh quá trình cam kết
tích lũy của doanh nghiệp vào các thị trường nước ngoài (Johanson \&
Vahlne, 1977, 2009).

Việc lựa chọn cường độ xuất khẩu làm thước đo trung tâm còn được củng cố
bởi một truyền thống nghiên cứu kinh tế học thương mại quốc tế song hành
với truyền thống kinh doanh quốc tế. Melitz (2003) chứng minh trên cơ sở
lý thuyết rằng việc mở cửa thương mại tạo ra sự tái phân bổ nguồn lực
trong nội bộ ngành, theo đó chỉ những doanh nghiệp có năng suất đủ cao
mới vượt qua được chi phí cố định của xuất khẩu và tự lựa chọn tham gia
thị trường quốc tế. Helpman và cộng sự (2004) mở rộng logic này bằng
cách phân biệt giữa lựa chọn xuất khẩu và lựa chọn đầu tư trực tiếp, cho
thấy rằng những doanh nghiệp năng suất cao nhất có xu hướng phục vụ thị
trường nước ngoài thông qua hiện diện sản xuất thay vì xuất khẩu thuần
túy. Hai luận điểm này có hàm ý đo lường trực tiếp đối với luận án:
cường độ xuất khẩu không phải là một biến ngoại sinh thuần túy mà gắn
liền với năng suất nội tại của doanh nghiệp, do đó quan hệ giữa quốc tế
hóa và hiệu quả có thể vừa phản ánh cơ chế nhân quả vừa phản ánh cơ chế
tự lựa chọn. Tổng quan thực nghiệm cấp doanh nghiệp của Wagner (2007)
xác nhận rằng bằng chứng về việc doanh nghiệp năng suất cao tự lựa chọn
xuất khẩu (self-selection) vững chắc hơn nhiều so với bằng chứng về việc
xuất khẩu trực tiếp nâng cao năng suất (learning-by-exporting), một sự
bất đối xứng có hệ quả phương pháp luận quan trọng được luận án xử lý
tường minh trong chiến lược nhận diện ở Chương 3. Nhận thức về tính nội
sinh tiềm tàng này không làm suy yếu giá trị của thước đo cường độ xuất
khẩu, mà ngược lại đòi hỏi lý thuyết hóa cẩn trọng về cơ chế đứng sau
quan hệ quan sát được, một yêu cầu mà khung tích hợp đa tầng của luận án
đáp ứng bằng cách phân tách các tầng cơ chế thay vì gộp chúng vào một hệ
số tổng hợp duy nhất.

\hypertarget{hiux1ec7u-quux1ea3-houx1ea1t-ux111ux1ed9ng-doanh-nghiux1ec7p}{%
\subsection{2.1.2 Hiệu quả hoạt động doanh
nghiệp}\label{hiux1ec7u-quux1ea3-houx1ea1t-ux111ux1ed9ng-doanh-nghiux1ec7p}}

Hiệu quả hoạt động doanh nghiệp (firm performance) là một khái niệm đa
chiều không có một định nghĩa hoặc thước đo thống nhất trong tài liệu.
Combs et al.~(2005) và Richard et al.~(2009) phân biệt ba nhóm thước đo
chính: (i) thước đo dựa trên kế toán (accounting-based), bao gồm tỷ suất
sinh lời trên tổng tài sản (ROA), tỷ suất sinh lời trên vốn chủ sở hữu
(ROE), và tỷ suất sinh lời trên doanh thu (ROS); (ii) thước đo dựa trên
thị trường tài chính (market-based), bao gồm Tobin's Q và tổng lợi nhuận
cổ đông; và (iii) thước đo dựa trên năng suất và tăng trưởng
(productivity-based), bao gồm năng suất lao động (labor productivity) và
tăng trưởng doanh thu.

Trong bối cảnh nghiên cứu sử dụng dữ liệu WBES, vốn không cung cấp thông
tin về giá cổ phiếu hay giá trị thị trường, thước đo năng suất lao động
là lựa chọn phổ biến nhất và phù hợp nhất về mặt kỹ thuật. Năng suất lao
động được tính bằng logarithm tự nhiên của tỷ số doanh thu hàng năm trên
số lao động toàn thời gian (ln(doanh thu/lao động)), cho phép so sánh
chuẩn hóa giữa các doanh nghiệp thuộc các quốc gia, ngành nghề và quy mô
khác nhau (nghiên cứu thành phần P1). Luận án sử dụng thước đo này làm
biến phụ thuộc chính, với ROS, tăng trưởng doanh thu, tăng trưởng việc
làm và ln(doanh thu) được sử dụng như các thước đo độ vững để kiểm tra
tính nhất quán của kết quả (xem Chương 3, Mục 3.5).

Cần nhấn mạnh rằng việc chọn thước đo hiệu quả không chỉ là quyết định
kỹ thuật, mà còn phản ánh quan niệm lý thuyết về hiệu quả. Trong bối
cảnh quốc tế hóa ở các nền kinh tế đang phát triển, năng suất lao động
phản ánh trực tiếp khả năng chuyển hóa nguồn lực đầu vào thành kết quả
đầu ra, qua đó thể hiện năng lực cạnh tranh thực chất của doanh nghiệp,
khác với ROA vốn bị ảnh hưởng bởi cấu trúc tài sản và hệ thống kế toán
đặc thù của từng quốc gia.

\hypertarget{nux103ng-lux1ef1c-cuxf4ng-nghux1ec7-vuxe0-chux1ec9-sux1ed1-chux1ea5p-nhux1eadn-sux1ed1}{%
\subsection{2.1.3 Năng lực công nghệ và chỉ số chấp nhận
số}\label{nux103ng-lux1ef1c-cuxf4ng-nghux1ec7-vuxe0-chux1ec9-sux1ed1-chux1ea5p-nhux1eadn-sux1ed1}}

Một trong những đóng góp khái niệm trọng tâm của luận án là sự phân biệt
rõ ràng giữa hai cấu trúc: năng lực công nghệ (Technological Capability
Index, TCI) và chỉ số chấp nhận số (Digital Adoption Index, DAI). Trong
nhiều nghiên cứu trước đây, hai cấu trúc này thường bị gộp lại thành một
khái niệm chung về ``năng lực số'' (digital capability), dẫn đến sự mơ
hồ lý thuyết và sai lệch trong đo lường (Bhandari et al., 2023).

Năng lực công nghệ (TCI) đề cập đến chiều sâu của năng lực công nghệ
được nhập vào tổ chức, phản ánh khả năng hấp thụ, khai thác và phát
triển tri thức công nghệ tiên tiến. Các chỉ báo tiêu biểu bao gồm việc
sở hữu chứng nhận chất lượng quốc tế (ISO hoặc tương đương) và sử dụng
công nghệ được cấp phép từ các đối tác nước ngoài. Về bản chất lý
thuyết, TCI gần với khái niệm năng lực hấp thụ (absorptive capacity) của
Cohen và Levinthal (1990) và năng lực công nghệ của Lall (1992), đây là
chiều sâu nội tại của tổ chức, khó sao chép và tạo ra lợi thế bền vững
theo tinh thần lý thuyết dựa trên nguồn lực (Barney, 1991).

Chỉ số chấp nhận số (Digital Adoption Index, DAI), ngược lại, phản ánh
mức độ áp dụng các công cụ số tại giao diện ngoại tại của doanh nghiệp,
bao gồm việc sở hữu website, sử dụng email trong giao dịch kinh doanh,
thanh toán điện tử, và bán hàng trực tuyến. Đây là các công cụ số có
tính bề mặt (surface digital tools) theo phân loại của Verhoef et
al.~(2021), tức là nằm ở cấp độ số hóa (digitization) và số hóa quy
trình (digitalization), chưa đạt đến cấp độ chuyển đổi số toàn diện
(digital transformation). Trong phạm vi luận án, DAI được đo bằng biến
đại diện ``foundational website adoption'' từ dữ liệu WBES, đây là ràng
buộc đo lường có chủ đích do giới hạn của bộ dữ liệu, không phải thiếu
sót lý thuyết.

Hai cấu trúc TCI và DAI không đồng nhất và vận hành theo cơ chế khác
nhau trong mối quan hệ quốc tế hóa, hiệu quả: TCI hoạt động chủ yếu như
một nhân tố nâng mặt bằng hiệu quả (level shifter), trong khi DAI có xu
hướng làm thay đổi độ dốc hoặc độ cong của đường quan hệ quốc tế hóa và
hiệu quả tùy theo bối cảnh (Bhandari et al., 2023; Chen \& Meng, 2022).
Việc duy trì sự phân biệt rõ ràng này trong suốt luận án là điều kiện
tiên quyết để đảm bảo tính thuần khiết của cấu trúc (construct purity)
theo tiêu chuẩn của Coltman et al.~(2008).

\hypertarget{bux1ed1i-cux1ea3nh-thux1ec3-chux1ebf-vuxe0-khung-icrv}{%
\subsection{2.1.4 Bối cảnh thể chế và khung
ICRV}\label{bux1ed1i-cux1ea3nh-thux1ec3-chux1ebf-vuxe0-khung-icrv}}

Thể chế (institutions) theo North (1990) là ``luật chơi'' của một xã
hội, bao gồm các ràng buộc chính thức (luật pháp, quy định, hợp đồng) và
phi chính thức (chuẩn mực, giá trị, văn hóa ứng xử). Thể chế tạo ra cấu
trúc khuyến khích và chi phí giao dịch, do đó ảnh hưởng trực tiếp đến
chiến lược và kết quả kinh doanh của doanh nghiệp.

Để vận hành hóa bối cảnh thể chế đa cấp và đa quốc gia trong luận án,
khung ICRV (Institutional Context Regime Variation, Biến thiên chế độ
bối cảnh thể chế) được phát triển như một công cụ phân loại có hệ thống.
Khung ICRV chia 50 nền kinh tế trong phạm vi nghiên cứu thành sáu nhóm
(sub-regime), dựa trên tổ hợp các chỉ số quản trị toàn cầu (World
Governance Indicators, WGI) của Kaufmann et al.~(2011), mức độ phát
triển kinh tế, và đặc trưng cấu trúc thể chế: Nhóm I (tiên tiến đổi mới,
gồm Singapore, Hong Kong SAR, Hàn Quốc, Đài Loan, Israel và Nhật Bản),
Nhóm II (tiên tiến tài nguyên, bao gồm các nền kinh tế vùng Vịnh), Nhóm
III (trung bình cao, điển hình là Trung Quốc), Nhóm IV (chuyển đổi thu
nhập trung bình thấp, bao gồm Việt Nam và Ấn Độ), Nhóm V (đang nổi), và
Nhóm VI (SIDS, Small Island Developing States). Khung này mở rộng từ
phân loại thể chế của Khanna và Palepu (2010) và được điều chỉnh để phản
ánh thực tiễn đa dạng của châu Á và khu vực lân cận.

ICRV không chỉ là một biến phân loại mà còn là một gradient thể chế liên
tục, từ môi trường thể chế mạnh, ít rủi ro (Nhóm I) đến môi trường thể
chế yếu, nhiều rủi ro và khoảng trống thể chế (Nhóm V và VI). Gradient
này tạo ra điều kiện biên quan trọng để kiểm định khả năng tổng quát hóa
của mối quan hệ quốc tế hóa và hiệu quả, đặc biệt liên quan đến giả
thuyết phi tuyến.

\begin{center}\rule{0.5\linewidth}{0.5pt}\end{center}

\hypertarget{cuxe1c-luxfd-thuyux1ebft-nux1ec1n}{%
\section{2.2 Các lý thuyết
nền}\label{cuxe1c-luxfd-thuyux1ebft-nux1ec1n}}

\hypertarget{muxf4-huxecnh-quux1ed1c-tux1ebf-huxf3a-uppsala}{%
\subsection{2.2.1 Mô hình quốc tế hóa
Uppsala}\label{muxf4-huxecnh-quux1ed1c-tux1ebf-huxf3a-uppsala}}

Mô hình Uppsala (Johanson \& Vahlne, 1977) là một trong những lý thuyết
có ảnh hưởng sâu rộng nhất trong nghiên cứu về quốc tế hóa doanh nghiệp.
Được xây dựng trên cơ sở quan sát quá trình quốc tế hóa của các doanh
nghiệp Thụy Điển, mô hình đề xuất rằng doanh nghiệp quốc tế hóa theo một
quá trình tuần tự và tích lũy, trong đó tri thức thị trường nước ngoài
và cam kết nguồn lực vào thị trường đó tác động lẫn nhau theo vòng lặp
phản hồi. Doanh nghiệp bắt đầu bằng xuất khẩu gián tiếp, tiến đến đại lý
độc lập, sau đó là chi nhánh kinh doanh, và cuối cùng là sản xuất tại
nước ngoài, phản ánh mức độ cam kết và tri thức ngày càng tăng.

Một trong những khái niệm trọng tâm của mô hình Uppsala là ``gánh nặng
của người nước ngoài'' (liability of foreignness), chi phí và bất lợi mà
doanh nghiệp phải gánh chịu khi hoạt động ở môi trường văn hóa, thể chế,
và pháp lý xa lạ (Zaheer, 1995). Johanson và Vahlne (2009) cập nhật mô
hình bằng cách bổ sung khái niệm ``gánh nặng của người ngoài mạng lưới''
(liability of outsidership), nhấn mạnh rằng trong nền kinh tế toàn cầu
hiện đại, rào cản lớn nhất không phải là khoảng cách địa lý hay tâm lý,
mà là vị trí ngoài lề của các mạng lưới kinh doanh địa phương. Hai loại
gánh nặng này giải thích vì sao chi phí quốc tế hóa giai đoạn đầu thường
cao, tạo ra giai đoạn hiệu quả giảm trước khi doanh nghiệp tích lũy đủ
tri thức và quan hệ để thu được lợi ích từ mở rộng quốc tế (Lu \&
Beamish, 2004; Contractor et al., 2003).

Tuy nhiên, mô hình Uppsala cổ điển được xây dựng trong bối cảnh tiền số
và tiền nền tảng (pre-platform), nên đã bị thách thức đáng kể trong thập
kỷ gần đây. Các nghiên cứu về born-global firms, tức doanh nghiệp quốc
tế hóa ngay từ giai đoạn đầu mà không theo trình tự tuần tự, và doanh
nghiệp nền tảng, tức doanh nghiệp công nghệ tận dụng nền tảng số để vươn
ra thị trường quốc tế không cần hiện diện vật lý, cho thấy logic của
Uppsala không còn phổ quát (Stallkamp \& Schotter, 2021; Yang et al.,
2025). Trong luận án này, mô hình Uppsala không bị bác bỏ, mà được tái
định vị như một lớp nền giải thích chi phí học hỏi và bất định trong
giai đoạn đầu quốc tế hóa, còn cơ chế học hỏi được mở rộng để bao gồm
``học hỏi tăng cường bằng dữ liệu số'' (data-augmented learning) trong
kỷ nguyên chuyển đổi số, nơi dữ liệu từ thị trường nước ngoài có thể
được thu thập, xử lý và phân tích nhanh hơn nhiều so với trước
(Banalieva \& Dhanaraj, 2019).

Việc tái định vị mô hình Uppsala trong luận án dựa trên một phân biệt
khái niệm quan trọng giữa cơ chế học hỏi và trình tự thâm nhập. Vahlne
(2020), trong bản cập nhật gần đây nhất do chính đồng tác giả của mô
hình gốc thực hiện, nhấn mạnh rằng cốt lõi bền vững của truyền thống
Uppsala không nằm ở trình tự bốn bước mang tính mô tả, mà nằm ở cơ chế
tiến hóa: doanh nghiệp tích lũy tri thức trải nghiệm, điều chỉnh cam
kết, và tái cấu trúc vị trí mạng lưới theo thời gian. Cách diễn giải này
tách phần lõi lý thuyết khỏi phần biểu hiện thực nghiệm, do đó giải
phóng mô hình khỏi sự phụ thuộc vào bối cảnh tiền số. Tổng quan hệ thống
của Wach (2021) về quá trình tiến hóa của mô hình Uppsala qua bảy phiên
bản từ năm 1977 đến 2017 củng cố cách đọc này: dù các mô hình theo giai
đoạn thường bị phê phán, chuỗi điều chỉnh kế tiếp của chính các tác giả
gốc cho thấy tính bền vững và khả năng thích ứng của truyền thống
Uppsala, qua đó biện minh cho quyết định của luận án giữ mô hình như một
lớp nền giải thích thay vì bác bỏ nó. Đối với các doanh nghiệp ở những
nền kinh tế chuyển đổi và mới nổi tại châu Á, cơ chế học hỏi trải nghiệm
có hàm ý trực tiếp với hình dạng phi tuyến của quan hệ quốc tế hóa và
hiệu quả: giai đoạn đầu của đường cong, nơi hiệu quả suy giảm, phản ánh
chính chi phí học hỏi và chi phí của gánh nặng người nước ngoài chưa
được khấu hao, trong khi giai đoạn giữa phản ánh lợi ích thu được khi
tri thức trải nghiệm đạt ngưỡng tới hạn (Lu \& Beamish, 2004). Hennart
(2007) bổ sung một cảnh báo lý thuyết quan trọng rằng quan hệ giữa mức
độ đa quốc gia và hiệu quả không phải là một quy luật phổ quát tự thân,
mà phụ thuộc vào bản chất của tài sản đặc thù mà doanh nghiệp khai thác
xuyên biên giới; lập luận này củng cố cách tiếp cận của luận án, theo đó
cơ chế học hỏi của Uppsala cần được đặt trong tương tác với năng lực và
bối cảnh, thay vì được xem là động lực độc lập.

\hypertarget{luxfd-thuyux1ebft-dux1ef1a-truxean-nguux1ed3n-lux1ef1c}{%
\subsection{2.2.2 Lý thuyết dựa trên nguồn
lực}\label{luxfd-thuyux1ebft-dux1ef1a-truxean-nguux1ed3n-lux1ef1c}}

Lý thuyết dựa trên nguồn lực (Resource-Based View, RBV) do Wernerfelt
(1984) đặt nền và được Barney (1991) hệ thống hóa thành một khung lý
thuyết có ảnh hưởng bền vững. lý thuyết dựa trên nguồn lực cho rằng
doanh nghiệp không chỉ là tập hợp các hoạt động giá trị (value chain) mà
còn là bó nguồn lực (bundle of resources), và sự khác biệt về hiệu quả
lâu dài giữa các doanh nghiệp bắt nguồn từ sự khác biệt trong sở hữu và
khai thác các nguồn lực thỏa mãn tiêu chí VRIN: có giá trị (Valuable),
hiếm có (Rare), khó sao chép (Inimitable), và khó thay thế
(Non-substitutable).

Áp dụng vào bối cảnh quốc tế hóa, lý thuyết dựa trên nguồn lực giải
thích vì sao cùng một mức độ FSTS nhưng các doanh nghiệp lại đạt hiệu
quả rất khác nhau: những doanh nghiệp sở hữu năng lực công nghệ, năng
lực quản trị vượt trội, và năng lực học hỏi mạnh hơn sẽ chuyển hóa quốc
tế hóa thành hiệu quả tốt hơn do họ có thể giảm chi phí gánh nặng người
nước ngoài và khai thác lợi thế cạnh tranh trên phạm vi quốc tế hiệu quả
hơn (Hitt et al., 2006; Peng, 2001). Cùng mức độ quốc tế hóa, doanh
nghiệp có nguồn lực yếu hơn sẽ gánh chịu chi phí điều phối và rủi ro lớn
hơn, trong khi doanh nghiệp có nguồn lực VRIN mạnh hơn có thể biến những
thách thức đó thành lợi thế cạnh tranh.

Trong bối cảnh đương đại, lý thuyết dựa trên nguồn lực được mở rộng bằng
lý thuyết năng lực động (dynamic capabilities) của Teece et al.~(1997),
nhấn mạnh khả năng của doanh nghiệp trong việc tích hợp, xây dựng và tái
cấu hình các năng lực nội bộ để ứng phó với môi trường cạnh tranh thay
đổi nhanh chóng. Trong kỷ nguyên số, năng lực công nghệ (TCI), theo phân
biệt của luận án, được xem là một dạng năng lực động quan trọng, cho
phép doanh nghiệp tái cấu hình nguồn lực và mô hình kinh doanh để ứng
phó với bất định trong quốc tế hóa (Bhandari et al., 2023; Verhoef et
al., 2021). DAI, chỉ số áp dụng số nền tảng (biến đại diện Tầng 1) đo
lường sự hiện diện website, phản ánh mức độ chấp nhận kỹ thuật số cơ bản
chứ không phải năng lực động theo nghĩa Teece et al.~(1997). lý thuyết
dựa trên nguồn lực cung cấp nền lý thuyết trực tiếp cho giả thuyết H2 về
vai trò điều tiết của TCI, trong khi H3 về DAI được neo đậu trong Khung
chấp nhận số nền tảng (Banalieva \& Dhanaraj, 2019).

Lý thuyết năng lực động cần được phân tích sâu hơn để làm rõ vì sao TCI
được định vị như một dạng năng lực bậc cao trong khung của luận án.
Eisenhardt và Martin (2000) tinh chỉnh khái niệm năng lực động bằng cách
lập luận rằng các năng lực này về bản chất là những quy trình tổ chức có
thể nhận diện được, và giá trị cạnh tranh của chúng phụ thuộc vào tính
kịp thời cùng khả năng tái cấu hình nguồn lực trong môi trường biến
động. Teece (2007) làm rõ thêm vi nền tảng của năng lực động qua ba nhóm
hoạt động: cảm nhận cơ hội, nắm bắt cơ hội, và tái cấu hình tài sản; ba
nhóm hoạt động này tương ứng trực tiếp với những đòi hỏi mà doanh nghiệp
phải đáp ứng khi mở rộng ra thị trường nước ngoài bất định. Theo logic
đó, năng lực công nghệ không chỉ là một kho tài sản tĩnh thỏa mãn tiêu
chí VRIN, mà còn là cơ chế cho phép doanh nghiệp hấp thụ tri thức mới và
tái triển khai nguồn lực khi điều kiện thị trường quốc tế thay đổi
(Cohen \& Levinthal, 1990; Teece et al., 1997). Sự phân biệt giữa năng
lực bậc một, tức năng lực vận hành thường nhật, và năng lực bậc cao, tức
năng lực tái cấu hình, cung cấp cơ sở lý thuyết cho việc tách TCI khỏi
DAI: TCI mang đặc trưng của năng lực bậc cao gắn với chiều sâu tổ chức,
trong khi DAI gần với năng lực vận hành bề mặt, dễ tiếp cận và dễ sao
chép hơn, do đó khó tạo ra lợi thế bền vững theo tinh thần lý thuyết dựa
trên nguồn lực (Barney, 1991). Phân tầng này giải thích vì sao luận án
dự báo hai cấu trúc số vận hành qua hai cơ chế khác nhau trong quan hệ
quốc tế hóa và hiệu quả, thay vì giả định một hiệu ứng đồng nhất.

Một bổ sung lý thuyết cần thiết cho lý thuyết dựa trên nguồn lực trong
bối cảnh quốc tế hóa là sự phân biệt giữa tạo ra giá trị và chiếm giữ
giá trị. Teece (1986), trong lập luận kinh điển về việc thu lợi từ đổi
mới công nghệ, chỉ ra rằng việc sở hữu một năng lực có giá trị chưa đủ
để bảo đảm doanh nghiệp chiếm giữ được phần lợi ích kinh tế từ năng lực
đó; khả năng chiếm giữ giá trị còn phụ thuộc vào chế độ bảo hộ quyền sở
hữu (appropriability regime) và quyền kiểm soát các tài sản bổ trợ
chuyên biệt. Trong bối cảnh quốc tế hóa ở các nền kinh tế đang phát
triển, nơi chế độ bảo hộ thường yếu và các tài sản bổ trợ như kênh phân
phối hay tài chính thương mại còn khan hiếm, lập luận của Teece (1986)
hàm ý rằng cùng một năng lực công nghệ có thể chuyển hóa thành hiệu quả
ở mức độ rất khác nhau tùy theo môi trường thể chế. Luận điểm này tạo
cầu nối lý thuyết trực tiếp giữa lý thuyết dựa trên nguồn lực và lý
thuyết thể chế, đồng thời báo trước cơ chế thay thế giữa thể chế và năng
lực được phát biểu trong giả thuyết H5: nơi thể chế yếu, biên đóng góp
của năng lực doanh nghiệp lớn hơn vì nó bù đắp cho khoảng trống thể chế
trong việc bảo đảm chiếm giữ giá trị. Bằng chứng cấp doanh nghiệp gần
đây của Purkayastha và cộng sự (2024) củng cố hướng lập luận này khi cho
thấy rằng năng lực quản lý đầu tư nguồn lực và năng lực triển khai nguồn
lực chiến lược điều tiết đáng kể cách quốc tế hóa chuyển hóa thành hiệu
quả, nghĩa là hiệu quả của quốc tế hóa phụ thuộc không chỉ vào việc
doanh nghiệp sở hữu nguồn lực gì mà còn vào việc doanh nghiệp triển khai
và tái phân bổ nguồn lực đó hiệu quả đến đâu, đúng theo tinh thần năng
lực động của Teece và cộng sự (1997).

\hypertarget{luxfd-thuyux1ebft-thux1ec3-chux1ebf}{%
\subsection{2.2.3 Lý thuyết thể
chế}\label{luxfd-thuyux1ebft-thux1ec3-chux1ebf}}

Lý thuyết thể chế (Institutional Theory) trong kinh doanh quốc tế rút từ
nhiều nguồn lý luận: kinh tế học thể chế mới của North (1990), xã hội
học thể chế của Scott (1995), và lý thuyết về khoảng trống thể chế của
Khanna và Palepu (2010). Điểm chung của các tiếp cận này là nhấn mạnh
vai trò của môi trường thể chế, cả chính thức lẫn phi chính thức, trong
việc định hình chiến lược và kết quả kinh doanh của doanh nghiệp.

North (1990) chỉ ra rằng thể chế hiệu quả làm giảm chi phí giao dịch,
bảo vệ quyền sở hữu trí tuệ, và tạo ra môi trường khuyến khích đầu tư
dài hạn. Ngược lại, ở các nền kinh tế đang phát triển, sự vắng mặt hoặc
yếu kém của các thể chế trung gian thị trường, từ hệ thống pháp lý đến
thị trường vốn, từ tổ chức đánh giá tín nhiệm đến hiệp hội ngành nghề,
tạo ra ``khoảng trống thể chế'' (institutional voids) làm tăng chi phí
giao dịch và rủi ro cho doanh nghiệp khi hoạt động hoặc mở rộng ra nước
ngoài (Khanna \& Palepu, 2010).

Trong tài liệu về quốc tế hóa và hiệu quả, bằng chứng phân tích tổng hợp
của Marano et al.~(2016), phân tích 170 nghiên cứu với 32 quốc gia giai
đoạn 1972--2012, cho thấy thể chế quốc gia nguồn điều tiết đáng kể mối
quan hệ quốc tế hóa, hiệu quả, với doanh nghiệp từ các nền kinh tế có
thể chế yếu hơn có xu hướng thu được lợi ích thấp hơn hoặc thậm chí tổn
thất từ quốc tế hóa. Cuervo-Cazurra et al.~(2018) bổ sung rằng bất định
thể chế ở quốc gia nguồn (home country uncertainty) không chỉ tạo ra
thách thức mà còn có thể xây dựng năng lực quản lý bất định đặc thù,
giúp doanh nghiệp ở các thị trường mới nổi cạnh tranh hiệu quả hơn ở các
thị trường nước ngoài có đặc điểm tương tự. Lý thuyết thể chế cung cấp
nền tảng lý luận cho giả thuyết H5 và cho khung ICRV trong luận án.

Để vận hành hóa bối cảnh thể chế một cách chặt chẽ, luận án dựa trên một
nền tảng lý luận rộng hơn về kinh tế học thể chế và quan điểm dựa trên
thể chế trong chiến lược quốc tế. Williamson (2000) phân tầng các cấp độ
phân tích thể chế từ những ràng buộc phi chính thức ăn sâu vào văn hóa,
đến môi trường thể chế chính thức, đến quản trị giao dịch, và sau cùng
là phân bổ nguồn lực; cách phân tầng này làm rõ rằng chất lượng thể chế
tác động đến doanh nghiệp qua nhiều kênh đồng thời, chứ không chỉ qua
một chỉ số tổng hợp duy nhất. Peng và cộng sự (2008) phát triển quan
điểm dựa trên thể chế như một trụ cột thứ ba của chiến lược quốc tế, bên
cạnh quan điểm dựa trên ngành và quan điểm dựa trên nguồn lực, lập luận
rằng ở các nền kinh tế mới nổi, thể chế không chỉ là bối cảnh nền mà là
yếu tố trực tiếp định hình lựa chọn chiến lược của doanh nghiệp. Peng
(2003) bổ sung góc nhìn động về quá trình chuyển đổi thể chế, theo đó
các doanh nghiệp ở những nền kinh tế đang chuyển đổi điều chỉnh chiến
lược từ dựa trên quan hệ sang dựa trên quy tắc khi khung thể chế chính
thức trưởng thành; góc nhìn này cung cấp cơ sở lý thuyết cho giả thuyết
H6 về dị biệt thời gian của quan hệ quốc tế hóa và hiệu quả ở các nền
kinh tế chuyển đổi nhanh. Ngoài ra, Kim và cộng sự (2026) định nghĩa
năng lực thể chế như khả năng thực thi chính sách của các tổ chức công,
phân tách thành chiều tổ chức và chiều quản trị, cung cấp một khung phân
tích chính thức làm cơ sở cho việc phân loại sáu chế độ thể chế trong
khung ICRV. Tổng hợp các tiếp cận này, luận án xem khung ICRV không phải
là một biến phân loại tùy nghi, mà là sự vận hành hóa có hệ thống của
một gradient thể chế đa chiều, neo đậu trong lý thuyết thể chế đương
đại.

\hypertarget{luxfd-thuyux1ebft-buxe0n-ux111ux1ea1p-quux1ed1c-tux1ebf-huxf3a-cux1ee7a-doanh-nghiux1ec7p-nux1ec1n-kinh-tux1ebf-mux1edbi-nux1ed5i}{%
\subsection{2.2.3.1 Lý thuyết bàn đạp: quốc tế hóa của doanh nghiệp nền
kinh tế mới
nổi}\label{luxfd-thuyux1ebft-buxe0n-ux111ux1ea1p-quux1ed1c-tux1ebf-huxf3a-cux1ee7a-doanh-nghiux1ec7p-nux1ec1n-kinh-tux1ebf-mux1edbi-nux1ed5i}}

Lý thuyết bàn đạp (springboard theory) do Luo và Tung (2007) đề xuất bổ
sung một góc nhìn đặc thù cho quốc tế hóa của doanh nghiệp từ các nền
kinh tế mới nổi mà các lý thuyết có gốc ở nước phát triển chưa nắm bắt
đầy đủ. Theo lý thuyết này, doanh nghiệp mới nổi sử dụng mở rộng quốc tế
như một bàn đạp chiến lược để nhanh chóng thu nhận các nguồn lực chiến
lược còn thiếu, gồm công nghệ, thương hiệu, tri thức quản trị và tiếp
cận thị trường, qua đó bù đắp bất lợi của người đến sau và rút ngắn
khoảng cách năng lực với các doanh nghiệp dẫn đầu, tức nhảy cóc về năng
lực. Khác với logic khai thác lợi thế sẵn có trong các lý thuyết quốc tế
hóa cổ điển, lý thuyết bàn đạp nhấn mạnh động cơ tìm kiếm tài sản: doanh
nghiệp quốc tế hóa không chỉ để bán cái đã có mà còn để học hỏi và thâu
nhận cái chưa có.

Lý thuyết bàn đạp đặc biệt phù hợp với bối cảnh của luận án, vốn tập
trung vào các nền kinh tế mới nổi và chuyển đổi ở châu Á như Trung Quốc,
Việt Nam và Ấn Độ, nơi doanh nghiệp đối mặt đồng thời với bất lợi của
người đến sau và khoảng trống thể chế quê nhà. Nó bổ sung cho khung của
luận án theo hai hướng. Thứ nhất, nó soi sáng lại nhánh đi lên của quan
hệ chữ U ngược: ở giai đoạn đầu, hiệu quả chưa cao không chỉ vì chi phí
học hỏi theo mô hình Uppsala mà còn vì doanh nghiệp đang chủ động đầu tư
để xây dựng và thâu nhận năng lực, một quá trình mà phần thưởng chỉ hiện
thực hóa khi tài sản chiến lược được tích lũy và khai thác. Thứ hai, nó
kết nối trực tiếp với trục thể chế: lý thuyết bàn đạp lập luận rằng
chính khoảng trống và bất lợi thể chế quê nhà thúc đẩy doanh nghiệp mới
nổi vươn ra ngoài để bù đắp, cộng hưởng với lập luận điều tiết thể chế
ICRV và với quan sát của Cuervo-Cazurra et al.~(2018) rằng bất định thể
chế quê nhà có thể bồi đắp năng lực quản lý bất định đặc thù.

Cần phân định rõ phạm vi áp dụng để tránh suy luận quá đà. Lý thuyết bàn
đạp được phát triển chủ yếu cho đầu tư trực tiếp ra nước ngoài mang tính
thâu tóm tài sản, trong khi luận án vận hành hóa quốc tế hóa bằng cường
độ xuất khẩu (FSTS), tức chiều bán hàng ra nước ngoài chứ không phải
chiều đầu tư. Do đó luận án sử dụng lý thuyết bàn đạp như một lăng kính
động cơ bổ trợ giải thích vì sao doanh nghiệp mới nổi theo đuổi quốc tế
hóa bất chấp chi phí giai đoạn đầu, chứ không khẳng định rằng cường độ
xuất khẩu đo lường trực tiếp hành vi bàn đạp. Ở chiều xuất khẩu, một cơ
chế tương đồng vận hành qua học hỏi nhờ xuất khẩu và tích lũy năng lực
từ tiếp xúc thị trường quốc tế, một con đường thu nhận tri thức song
song với con đường thâu tóm tài sản mà lý thuyết bàn đạp mô tả.

\hypertarget{luxfd-thuyux1ebft-thux1b0ux1ee3ng-tux1ea7ng-quux1ea3n-trux1ecb}{%
\subsection{2.2.4 Lý thuyết thượng tầng quản
trị}\label{luxfd-thuyux1ebft-thux1b0ux1ee3ng-tux1ea7ng-quux1ea3n-trux1ecb}}

Lý thuyết thượng tầng quản trị (Upper Echelons Theory) do Hambrick và
Mason (1984) đề xuất và được Hambrick (2007) cập nhật, cho rằng các
quyết định chiến lược lớn và kết quả của tổ chức không phải là sản phẩm
của các quy trình duy lý hoàn toàn, mà phản ánh đặc điểm nhận thức, giá
trị và kinh nghiệm của nhà quản trị cấp cao. Nói cách khác, ``tổ chức là
sự phản chiếu của nhà quản trị đỉnh cao'' (Hambrick \& Mason, 1984, tr.
193).

Trong bối cảnh quốc tế hóa, lý thuyết thượng tầng quản trị có hai hàm ý
trực tiếp. Thứ nhất, kinh nghiệm quốc tế của nhà quản trị cấp cao, được
đo bằng số năm kinh nghiệm làm việc hoặc học tập ở nước ngoài, hay tiếp
xúc với thị trường quốc tế, giúp giảm nhận thức về gánh nặng người nước
ngoài và nâng cao khả năng đánh giá cơ hội, quản trị rủi ro ở thị trường
quốc tế (Hsu et al., 2013; Sambharya, 1996). Thứ hai, sự đa dạng giới
tính trong lãnh đạo, đặc biệt là sự hiện diện của quản trị viên nữ ở các
vị trí lãnh đạo, liên quan đến nhận thức chiến lược đa chiều hơn và thực
tiễn quản trị rủi ro thận trọng hơn, qua đó có tác động tích cực đến
hiệu quả doanh nghiệp trong quá trình quốc tế hóa (Post \& Byron, 2015).
Những hàm ý này cung cấp cơ sở lý thuyết cho giả thuyết H4 về vai trò
điều tiết của đặc điểm nhà quản trị cấp cao.

Cơ sở lý luận của lý thuyết thượng tầng quản trị trong bối cảnh quốc tế
hóa được củng cố thêm bởi luận điểm trung tâm rằng quyết định chiến lược
lớn được lọc qua trường nhận thức hạn chế và hệ giá trị của nhà quản trị
cấp cao, đặc biệt khi nhiệm vụ mang tính bất định cao và quá tải thông
tin (Hambrick, 2007). Mở rộng quốc tế là một trong những quyết định
chiến lược điển hình cho điều kiện này, bởi nó đòi hỏi nhà quản trị đánh
giá cơ hội và rủi ro ở những thị trường mà thông tin không đầy đủ và
khung tham chiếu quen thuộc không còn phù hợp. Kinh nghiệm quốc tế của
đội ngũ quản trị cấp cao, theo Sambharya (1996), gắn với mức độ và phạm
vi của chiến lược đa quốc gia, vì kinh nghiệm này mở rộng trường nhận
thức và giảm thiên lệch nhận thức gắn với gánh nặng người nước ngoài.
Nielsen (2010) tinh chỉnh thêm bằng cách chỉ ra rằng tác động của tính
quốc tế của đội ngũ quản trị cấp cao lên hiệu quả có thể vận hành qua
những cơ chế trung gian liên quan đến lựa chọn phương thức thâm nhập,
hàm ý rằng đặc điểm nhà quản trị không tác động trực tiếp lên hiệu quả
mà điều tiết chất lượng của các quyết định quốc tế hóa. Trong khung của
luận án, hai hàm ý này hội tụ vào giả thuyết H4: kinh nghiệm và đa dạng
giới tính của nhà quản trị cấp cao điều tiết khả năng doanh nghiệp
chuyển hóa quốc tế hóa thành hiệu quả, thông qua việc nâng cao chất
lượng nhận thức chiến lược và sự thận trọng trong quản trị rủi ro quốc
tế (Post \& Byron, 2015).

Một bước tiến lý thuyết gần đây làm sâu sắc thêm nền tảng vi mô của giả
thuyết H4. Grego và cộng sự (2025), qua tổng quan hệ thống hơn hai thập
kỷ nghiên cứu cấp cá nhân trong kinh doanh quốc tế, chỉ ra rằng các lý
thuyết truyền thống, gồm cả lý thuyết nội hóa và mô hình Uppsala, đều
giả định nhà quản trị là tác nhân duy lý có giới hạn, ra quyết định quốc
tế hóa chủ yếu bằng tính toán lợi ích chi phí lạnh lùng, do đó bỏ qua
vai trò của nhận thức cảm tính. Các tác giả đề xuất một mô hình quyết
định quốc tế hóa được kiến tạo bởi cảm xúc, theo đó nhận thức cảm tính
và nhận thức phi cảm tính cùng định hình tri thức, kinh nghiệm, thái độ
và đánh giá rủi ro của nhà quản trị. Luận điểm này có hai hàm ý cho luận
án. Thứ nhất, nó củng cố cách đọc rằng tác động của đặc điểm nhà quản
trị lên quan hệ quốc tế hóa và hiệu quả là không đồng nhất và không quy
giản về một cơ chế duy lý đơn giản, nhất quán với phát hiện thực nghiệm
của luận án rằng hệ số của biến quản trị nữ phụ thuộc thước đo và bối
cảnh. Thứ hai, nó cảnh báo rằng các thước đo đặc điểm nhà quản trị quan
sát được ở dữ liệu doanh nghiệp, như kinh nghiệm và giới tính, chỉ là
chỉ báo bề mặt của một quá trình nhận thức phức tạp hơn, nên các hệ số
ước lượng cần được diễn giải thận trọng thay vì gán nhân quả trực tiếp,
đúng theo lập trường mà luận án giữ trong Chương 4 và Chương 5. Bằng
chứng của Pan, Verbeke và Yuan (2021) về lãnh đạo doanh nghiệp tại Trung
Quốc minh họa cùng nguyên tắc ở một góc khác: tác động của lãnh đạo
chuyển hóa lên tinh thần khởi nghiệp doanh nghiệp không có ý nghĩa thống
kê khi bỏ qua các điều kiện biên, nhưng trở nên có ý nghĩa khi đưa vào
các biến điều tiết về đội ngũ quản trị, môi trường và thiết kế tổ chức,
một mô thức song hành với luận điểm trung tâm của luận án rằng hiệu ứng
chỉ bộc lộ đúng khi được điều kiện hóa theo bối cảnh.

\hypertarget{lux1edbp-mux1edf-rux1ed9ng-lux103ng-kuxednh-nux103ng-lux1ef1c-sux1ed1}{%
\subsection{2.2.5 Lớp mở rộng: Lăng kính năng lực
số}\label{lux1edbp-mux1edf-rux1ed9ng-lux103ng-kuxednh-nux103ng-lux1ef1c-sux1ed1}}

Bốn lý thuyết nền trên được xây dựng chủ yếu trong bối cảnh tiền số
(pre-digital). Sự phát triển mạnh mẽ của kinh tế số và trí tuệ nhân tạo
từ cuối thập niên 2010 đến nay đặt ra yêu cầu cập nhật và mở rộng các lý
thuyết kinh doanh quốc tế để phản ánh logic mới của quốc tế hóa trong kỷ
nguyên số.

Banalieva và Dhanaraj (2019) là một trong những nghiên cứu đầu tiên đề
xuất cập nhật lý thuyết nội hóa (internalization theory) cho nền kinh tế
số, lập luận rằng dữ liệu số và tương tác dựa trên nền tảng làm giảm sự
phụ thuộc vào hiện diện vật lý và thay đổi cơ bản logic về quyết định
nội hóa. Stallkamp và Schotter (2021) mở rộng luận điểm này bằng cách
chỉ ra rằng các doanh nghiệp nền tảng (platform firms) quốc tế hóa theo
logic đa diện (multi-sided) khác hẳn với các MNCs truyền thống, với hiệu
ứng mạng lưới tạo ra lợi thế quy mô không bị ràng buộc bởi biên giới địa
lý. Verhoef et al.~(2021) cung cấp khung khái niệm phân tầng ba cấp độ:
số hóa thông tin (digitization), số hóa quy trình (digitalization), và
chuyển đổi số toàn diện (digital transformation), nhấn mạnh rằng các cấp
độ này khác nhau về chiều sâu tổ chức và tác động kinh doanh. Yang et
al.~(2025) bổ sung bằng tổng quan hệ thống về quốc tế hóa của các doanh
nghiệp số, phác thảo các lý thuyết cần điều chỉnh và hướng nghiên cứu
tiếp theo. Đặc biệt liên quan đến luận án, Drori và cộng sự (2024), cũng
từ lăng kính lý thuyết nội hóa, kiểm định trực tiếp trên mẫu 2.370 doanh
nghiệp Mỹ trong giai đoạn 18 năm và phát hiện rằng số hóa làm giảm khiếm
khuyết thị trường nhưng mức độ giảm thay đổi theo cường độ số hóa của
ngành, đồng thời các hiệu ứng tương tác này biến đổi theo thời gian khi
đổi mới số lan tỏa rộng hơn. Phát hiện này cung cấp bằng chứng thực
nghiệm gần đây và trực tiếp cho hai luận điểm cốt lõi của luận án: vai
trò của năng lực số đối với quốc tế hóa là phụ thuộc bối cảnh (ở đây là
cường độ số hóa của ngành, song hành với lập luận về phụ thuộc thể chế
của luận án) và phụ thuộc thời gian, chứ không phải một hiệu ứng đồng
nhất và bất biến.

Sự phân biệt giữa hai cấu trúc số có nền tảng đo lường vững chắc trong
tài liệu. Bharadwaj et al.~(2013) lập luận rằng năng lực công nghệ thông
tin (IT capability) và chiến lược kinh doanh số (digital business
strategy) có mạng quan hệ lý thuyết khác nhau, nghĩa là chúng liên hệ
với các tiền tố và hệ quả khác nhau, do đó không nên gộp thành một cấu
trúc. Coltman et al.~(2008) cung cấp bốn tiêu chí để phân biệt cấu trúc
formative với reflective, làm cơ sở phương pháp luận cho việc xây dựng
TCI và DAI như hai chỉ số hợp thành formative độc lập thay vì các chỉ
báo của một cấu trúc tiềm ẩn chung. Trên nền đó, luận án định vị TCI ở
truyền thống năng lực công nghệ hóa công nghiệp và năng lực hấp thụ, vốn
xem năng lực công nghệ là chiều sâu năng lực nội tại được tích lũy qua
đầu tư và học hỏi (Cohen \& Levinthal, 1990; Lall, 1992), trong khi DAI
phản ánh mức độ áp dụng công cụ số ở giao diện ngoại tại. Sự tách biệt
khái niệm này không chỉ là tinh tế đo lường mà có hệ quả lý thuyết trực
tiếp: nếu hai cấu trúc vận hành qua hai cơ chế khác nhau, thì việc gộp
chúng sẽ trung bình hóa hai hiệu ứng trái chiều và che khuất chính cơ
chế mà nhà nghiên cứu và nhà hoạch định chính sách cần phân biệt.

Sự trỗi dậy của trí tuệ nhân tạo trong những năm gần đây làm sâu sắc
thêm, chứ không làm thay đổi, logic phân tách hai cấu trúc số của luận
án. Babina và cộng sự (2024) cung cấp bằng chứng quy mô lớn rằng việc áp
dụng trí tuệ nhân tạo gắn với tăng trưởng doanh nghiệp và đổi mới sản
phẩm, song chủ yếu thông qua mở rộng quy mô và phạm vi sản phẩm ở những
doanh nghiệp vốn đã có chiều sâu năng lực, chứ không phân bổ đồng đều
cho mọi doanh nghiệp áp dụng công cụ số. Phát hiện này nhất quán với
``đường cong J năng suất'' của Brynjolfsson và cộng sự (2021), theo đó
lợi ích năng suất của các công nghệ đa dụng chỉ hiện thực hóa sau một độ
trễ đầu tư vào tài sản vô hình bổ trợ, hàm ý rằng việc chấp nhận công cụ
số bề mặt (DAI) và việc tích lũy năng lực tổ chức sâu (TCI) tách biệt cả
về thời gian lẫn cơ chế tạo giá trị. Đồng thời, một khối bằng chứng vĩ
mô cho thấy năng lực số phân bố cực kỳ không đồng đều giữa các nền kinh
tế: khoảng cách kết nối internet vẫn còn rất lớn ở các nền kinh tế thu
nhập thấp (Pirlea \& Corso, 2026), và bất bình đẳng trong tiếp cận cũng
như sử dụng trí tuệ nhân tạo còn sâu sắc hơn (Mahler, Wang \& Weber,
2026). Katz và Callorda (2018) định lượng đóng góp kinh tế của băng
thông rộng và số hóa, cho thấy hiệu ứng của hạ tầng số lên tăng trưởng
phụ thuộc vào ngưỡng phát triển và chất lượng quy định. Những bằng chứng
vĩ mô này không phải là bối cảnh ngoài lề mà có hàm ý lý thuyết trực
tiếp đối với khung của luận án: vì hạ tầng và kỹ năng số là điều kiện
cần để công cụ số phát huy tác dụng phối hợp xuyên biên giới, hiệu ứng
điều tiết của DAI được dự báo là phụ thuộc bối cảnh thể chế chứ không
phổ quát, một dự đoán có điều kiện được phát biểu tường minh trong giả
thuyết H3 và H5.

Trong luận án, lớp mở rộng về năng lực số đóng vai trò mở rộng lý thuyết
dựa trên nguồn lực bằng cách xác định TCI và DAI như hai loại tài sản số
có chiều sâu và cơ chế vận hành khác nhau, đồng thời mở rộng Uppsala
bằng cơ chế học hỏi dựa trên dữ liệu số. Sự tích hợp này tạo thành lớp
thứ năm của khung lý thuyết, bổ sung cho bốn lý thuyết kinh điển để tạo
thành một khung giải thích hoàn chỉnh và phù hợp với thực tiễn kinh
doanh đương đại. Cách tích hợp năm lớp này, thay vì chọn một hoặc hai lý
thuyết làm nền như đa số nghiên cứu trước, chính là phản hồi trực tiếp
cho khoảng trống thứ nhất đã nêu ở Chương 1: không một lý thuyết đơn lẻ
nào đủ sức giải thích mức dị biệt cao của quan hệ quốc tế hóa và hiệu
quả, và chỉ một khung đa tầng tích hợp mới nắm bắt được đồng thời chi
phí học hỏi (Uppsala), khác biệt năng lực (RBV), bối cảnh thể chế
(Institutional Theory), đặc điểm nhà quản trị (Upper Echelons) và logic
số (Digital Capability Lens).

\hypertarget{tux1ed5ng-hux1ee3p-khung-luxfd-thuyux1ebft-nguyuxean-tux1eafc-tuxedch-hux1ee3p-ux111a-tux1ea7ng}{%
\subsection{2.2.6 Tổng hợp khung lý thuyết: nguyên tắc tích hợp đa
tầng}\label{tux1ed5ng-hux1ee3p-khung-luxfd-thuyux1ebft-nguyuxean-tux1eafc-tuxedch-hux1ee3p-ux111a-tux1ea7ng}}

Việc đặt cạnh nhau năm lý thuyết nền không tự động tạo thành một khung
tích hợp; điều phân biệt một khung tích hợp với một danh sách lý thuyết
rời rạc là sự minh định về cách các lý thuyết liên kết với nhau theo cấp
độ phân tích và về điều kiện biên áp dụng của từng lý thuyết. Tiểu mục
này tổng hợp năm lớp lý thuyết thành một logic tích hợp duy nhất, làm
nền cho khung CDCM được trình bày ở Mục 2.5.

Nguyên tắc tổ chức thứ nhất là phân tầng theo cấp độ phân tích. Mỗi lý
thuyết soi sáng một tầng nguyên nhân riêng biệt nhưng vận hành đồng
thời: mô hình Uppsala định vị cơ chế ở cấp quá trình, giải thích chi phí
học hỏi và động thái tích lũy tri thức theo thời gian; lý thuyết dựa
trên nguồn lực định vị cơ chế ở cấp doanh nghiệp, giải thích vì sao cùng
một mức quốc tế hóa lại tạo ra hiệu quả khác nhau giữa các doanh nghiệp;
lý thuyết thể chế định vị cơ chế ở cấp bối cảnh quốc gia, giải thích
biến thiên theo môi trường nền; lý thuyết thượng tầng quản trị định vị
cơ chế ở cấp cá nhân người ra quyết định; và lăng kính năng lực số bổ
sung một chiều ngang cắt qua cả bốn tầng, phản ánh cách công nghệ số tái
định hình chi phí học hỏi, bản chất nguồn lực, hiệu lực thể chế và năng
lực nhận thức của nhà quản trị. Quan điểm dựa trên thể chế trong chiến
lược quốc tế của Peng và cộng sự (2008), vốn đề xuất thể chế như trụ cột
thứ ba bên cạnh quan điểm dựa trên ngành và quan điểm dựa trên nguồn
lực, cung cấp chính cái khớp nối lý thuyết cho phép các tầng này tương
tác thay vì tồn tại song song; trong khung của luận án, thể chế không
chỉ là biến bối cảnh mà là yếu tố định hình hiệu lực của mọi tầng cơ chế
còn lại.

Nguyên tắc tổ chức thứ hai là bổ sung thay vì cạnh tranh. Các tranh luận
trong tài liệu thường khung hóa các lý thuyết này như những giải thích
đối kháng, ví dụ tranh luận giữa giải thích dựa trên năng lực và giải
thích dựa trên thể chế cho hiệu quả của doanh nghiệp ở các nền kinh tế
mới nổi. Luận án lập luận rằng sự đối kháng này phần lớn là biểu kiến:
vì các cơ chế vận hành ở các cấp độ khác nhau, chúng không loại trừ nhau
mà tương tác. Kano và cộng sự (2020), trong tổng quan đa ngành về chuỗi
giá trị toàn cầu, minh họa rõ rằng kết quả của doanh nghiệp khi tham gia
hoạt động xuyên biên giới là sản phẩm đồng thời của năng lực nội tại,
cấu trúc quản trị và bối cảnh thể chế, chứ không quy giản về bất kỳ yếu
tố đơn lẻ nào. Logic tương tác này là cơ sở để luận án mô hình hóa các
biến điều tiết một cách đồng thời thay vì kiểm định rời rạc từng biến.

Tuy nhiên, bổ sung thay vì cạnh tranh không có nghĩa là gộp các lý
thuyết một cách hình thức; việc kết hợp các lý thuyết có tiền đề khác
nhau đòi hỏi minh định chính các tiền đề đó để tránh một khung tích hợp
gượng ép (Santangelo \& Verbeke, 2022). Năm lý thuyết nền của luận án
dựa trên những tiền đề khác nhau về đơn vị phân tích và bản chất hành
vi: mô hình Uppsala giả định doanh nghiệp học hỏi tiệm tiến và e ngại
bất định; lý thuyết dựa trên nguồn lực giả định tính dị biệt và khó di
chuyển của nguồn lực; lý thuyết thể chế giả định hành vi bị ràng buộc
bởi cấu trúc khuyến khích bên ngoài; lý thuyết thượng tầng quản trị giả
định tính duy lý có giới hạn và vai trò của nhận thức cá nhân; và lăng
kính năng lực số giả định công nghệ làm thay đổi chi phí giao dịch và
bản chất tài sản. Luận án hòa giải các tiền đề này không phải bằng cách
xem chúng loại trừ nhau mà bằng cách định vị chúng như các lát cắt bổ
sung của cùng một quá trình: thể chế định khung không gian lựa chọn,
năng lực và nhà quản trị quyết định cách doanh nghiệp di chuyển trong
không gian đó, và logic số dịch chuyển các đường biên chi phí. Lý thuyết
bàn đạp (Luo \& Tung, 2007) được tích hợp đúng theo nguyên tắc này: nó
không thay thế mà bổ sung một tiền đề động cơ đặc thù cho doanh nghiệp
nền kinh tế mới nổi, rằng quốc tế hóa đồng thời là phương tiện thâu nhận
năng lực còn thiếu, và tiền đề này tương thích với chứ không mâu thuẫn
với logic tích lũy năng lực của lý thuyết dựa trên nguồn lực, miễn là
giữ rõ phân biệt giữa chiều đầu tư thâu tóm tài sản và chiều cường độ
xuất khẩu mà luận án vận hành hóa.

Nguyên tắc tổ chức thứ ba là điều kiện biên. Một khung tích hợp chỉ chặt
chẽ khi nó xác định được không chỉ nơi lý thuyết áp dụng mà cả nơi lý
thuyết thất bại. Hennart (2007) đã cảnh báo rằng quan hệ giữa mức độ đa
quốc gia và hiệu quả không phải là quy luật phổ quát mà phụ thuộc vào
bản chất tài sản đặc thù; điều kiện biên FIP tại các nền kinh tế đảo
nhỏ, phát biểu trong giả thuyết H1b, chính là sự cụ thể hóa cảnh báo này
thành một mệnh đề kiểm định được, xác định bối cảnh nơi dạng hàm chữ U
ngược sụp đổ thành quan hệ âm đơn điệu. Việc khung tích hợp bao gồm cả
điều kiện biên cực trị, thay vì chỉ khẳng định tính phổ quát của một
dạng hàm, là điều phân biệt một lý thuyết có khả năng phản nghiệm với
một mô tả thiếu ràng buộc. Tổng hợp lại, ba nguyên tắc phân tầng theo
cấp độ, bổ sung thay vì cạnh tranh, và điều kiện biên tường minh cấu
thành logic tích hợp của khung CDCM, và mỗi nguyên tắc được phản ánh
trực tiếp trong cấu trúc của hệ giả thuyết H1 đến H6 cùng điều kiện biên
H1b.

\begin{center}\rule{0.5\linewidth}{0.5pt}\end{center}

\hypertarget{tux1ed5ng-quan-thux1ef1c-nghiux1ec7m}{%
\section{2.3 Tổng quan thực
nghiệm}\label{tux1ed5ng-quan-thux1ef1c-nghiux1ec7m}}

\hypertarget{giai-ux111oux1ea1n-phuxe1t-triux1ec3n-cux1ee7a-tuxe0i-liux1ec7u-vux1ec1-quan-hux1ec7-quux1ed1c-tux1ebf-huxf3a-vuxe0-hiux1ec7u-quux1ea3}{%
\subsection{2.3.1 Giai đoạn phát triển của tài liệu về quan hệ quốc tế
hóa và hiệu
quả}\label{giai-ux111oux1ea1n-phuxe1t-triux1ec3n-cux1ee7a-tuxe0i-liux1ec7u-vux1ec1-quan-hux1ec7-quux1ed1c-tux1ebf-huxf3a-vuxe0-hiux1ec7u-quux1ea3}}

Nghiên cứu về mối quan hệ giữa quốc tế hóa và hiệu quả hoạt động doanh
nghiệp (I--P relationship) đã trải qua hơn bốn thập kỷ phát triển, với
những thay đổi đáng kể về cả câu hỏi nghiên cứu, phương pháp luận, và
kết luận lý thuyết.

Giai đoạn đầu (từ thập niên 1980 đến giữa thập niên 1990) đặc trưng bởi
giả định về quan hệ tuyến tính dương giữa quốc tế hóa và hiệu quả. Grant
(1987) và Kim et al.~(1989) cho rằng quốc tế hóa giúp doanh nghiệp khai
thác lợi thế theo quy mô (economies of scale), lợi thế phạm vi
(economies of scope), và phân tán rủi ro kinh tế vĩ mô, do đó dẫn đến
hiệu quả cao hơn. Logic này trực quan và phù hợp với quan sát ban đầu về
các MNCs thành công từ Mỹ, Châu Âu và Nhật Bản.

Giai đoạn thứ hai (từ cuối thập niên 1990 đến đầu thập niên 2000) chứng
kiến sự xuất hiện của mô hình phi tuyến, thách thức giả định tuyến tính
giản đơn. Hitt et al.~(1997) là nghiên cứu đầu tiên kiểm định và tìm
thấy bằng chứng cho quan hệ chữ U ngược (chữ U ngược) giữa đa dạng hóa
sản phẩm quốc tế và hiệu quả, lập luận rằng lợi ích từ quốc tế hóa tăng
đến một điểm tối ưu rồi suy giảm do chi phí điều phối và phức tạp tổ
chức vượt qua lợi ích. Gomes và Ramaswamy (1999) cung cấp bằng chứng
tương tự trong bối cảnh doanh nghiệp đa ngành. Contractor et al.~(2003)
và Lu và Beamish (2004) tinh chỉnh mô hình phi tuyến với đề xuất ``đường
cong chữ S ba giai đoạn'' (three-stage S-curve), theo đó hiệu quả giảm ở
giai đoạn thâm nhập ban đầu (chi phí thiết lập), tăng ở giai đoạn khai
thác (học hỏi tích lũy và quy mô), rồi lại giảm ở giai đoạn siêu mở rộng
(over-extension) khi chi phí phức tạp vượt quá lợi ích. Ruigrok và
Wagner (2003) bổ sung bằng chứng cho cơ chế học hỏi tổ chức làm nền cho
giai đoạn giữa của đường cong, trong khi Contractor (2012) tổng hợp lý
thuyết đa giai đoạn này cùng các phê phán phương pháp luận về cách đo
lường và kiểm định quan hệ phi tuyến.

Giai đoạn thứ ba (từ cuối thập niên 2000 đến hiện tại) chuyển trọng tâm
sang nghiên cứu các yếu tố điều tiết (các biến điều tiết), tức điều kiện
bối cảnh làm thay đổi hình dạng, cường độ, hoặc dấu của mối quan hệ quốc
tế hóa và hiệu quả. Các biến điều tiết được nghiên cứu bao gồm đặc điểm
quốc gia nguồn (Marano et al., 2016), đặc điểm ngành (Kirca et al.,
2012), năng lực lãnh đạo cấp cao (Hsu et al., 2013), và năng lực kỹ
thuật số (Bhandari et al., 2023; Wu et al., 2022). Sự chuyển dịch này
phản ánh nhận thức ngày càng sâu sắc rằng không có quan hệ quốc tế hóa
và hiệu quả phổ quát, kết quả quốc tế hóa phụ thuộc vào điều kiện cụ thể
của doanh nghiệp, ngành và quốc gia (Arte \& Larimo, 2022).

Một vấn đề xuyên suốt bốn thập kỷ phát triển của tài liệu về quan hệ
quốc tế hóa và hiệu quả là sự thiếu thống nhất về hình dạng hàm số của
mối quan hệ, mà bản thân nó phần nào bắt nguồn từ khác biệt phương pháp
luận hơn là từ khác biệt thực chất. Glaum và Oesterle (2007), trong bài
tổng kết bốn mươi năm nghiên cứu, kết luận rằng tài liệu để lại nhiều
câu hỏi hơn câu trả lời, một phần do sự đa dạng trong cách đo lường mức
độ quốc tế hóa và cách định dạng hàm hồi quy. Thomas và Eden (2004) chỉ
ra trực tiếp rằng hình dạng quan hệ phụ thuộc vào việc nghiên cứu giả
định bậc hai hay bậc ba, và vào khoảng giá trị của biến quốc tế hóa được
quan sát trong mẫu. Berry và Kaul (2016), qua một nỗ lực tái lặp quy mô
lớn, đặt câu hỏi về tính phổ quát của đường cong chữ S, cho thấy kết quả
nhạy cảm với đặc tả mô hình và cấu thành mẫu. Những phê phán này có hệ
quả phương pháp luận quan trọng đối với luận án: kiểm định quan hệ phi
tuyến đòi hỏi quy trình chặt chẽ, bao gồm kiểm định sự tồn tại của điểm
cực trị nằm trong khoảng dữ liệu (Lind \& Mehlum, 2010) và lý thuyết hóa
cẩn trọng về cơ chế tạo ra độ cong thay vì chỉ đưa số hạng bình phương
vào mô hình một cách máy móc (Haans et al., 2016). Cách tiếp cận của
luận án, theo đó điểm uốn của đường cong được lý thuyết hóa như hệ quả
của tương tác giữa chi phí học hỏi, năng lực và bối cảnh thể chế, là một
phản hồi trực tiếp cho các phê phán phương pháp luận này, đồng thời
tránh hiện tượng ``quá nhiều điều tốt'' được khái quát hóa bởi Pierce và
Aguinis (2013).

Sự đa dạng về dạng hàm được đề xuất trong tài liệu phi tuyến đáng được
lược khảo chi tiết hơn, vì chính sự đa dạng này là một phần của khoảng
trống mà luận án giải quyết. Bên cạnh dạng chữ U ngược bậc hai và dạng
chữ S ba giai đoạn, một số nghiên cứu đề xuất các dạng hàm phức tạp hơn.
Almodóvar và Rugman (2014) đề xuất ``đường cong chữ M'' cho các doanh
nghiệp khởi nghiệp quốc tế mới, hàm ý hai điểm uốn phản ánh hai chu kỳ
học hỏi liên tiếp. Tallman và Li (1996) phân tách tác động của đa dạng
quốc tế khỏi đa dạng sản phẩm và cho thấy hai chiều này có quan hệ khác
nhau với hiệu quả, một lời nhắc rằng cách vận hành hóa quốc tế hóa quyết
định dạng hàm quan sát được. Riahi-Belkaoui (1998) cung cấp bằng chứng
ban đầu về dạng phi tuyến trên mẫu doanh nghiệp đa quốc gia, trong khi
Miller và cộng sự (2016) lập luận rằng việc gộp ba thành phần cường độ,
đa dạng và khoảng cách của quốc tế hóa vào một thước đo duy nhất che
khuất các quan hệ trái chiều mà mỗi thành phần có với hiệu quả. Tsionas
và Tzeremes (2021) đo lường trực tiếp quan hệ giữa mức độ quốc tế hóa và
năng suất doanh nghiệp trên mẫu các doanh nghiệp đa quốc gia lớn, cung
cấp một tiền lệ phương pháp luận cho việc luận án sử dụng năng suất lao
động làm biến phụ thuộc thay vì các thước đo dựa trên kế toán. Đáng chú
ý, Schmuck và cộng sự (2022) đảo ngược chiều nhân quả thường được giả
định và kiểm định khả năng hiệu quả cao dẫn đến mức độ đa quốc gia cao,
một lời nhắc rằng vấn đề nội sinh trong quan hệ quốc tế hóa và hiệu quả
không thể bỏ qua. Tổng hợp lại, sự phân tán về dạng hàm và về chiều nhân
quả trong tài liệu củng cố luận điểm trung tâm của luận án rằng không
tồn tại một dạng hàm phổ quát; hình dạng quan sát được là sản phẩm của
bối cảnh, cách đo lường và chiến lược nhận diện, và do đó phải được lý
thuyết hóa thay vì giả định.

\hypertarget{bux1eb1ng-chux1ee9ng-tux1eeb-phuxe2n-tuxedch-tux1ed5ng-hux1ee3p}{%
\subsection{2.3.2 Bằng chứng từ phân tích tổng
hợp}\label{bux1eb1ng-chux1ee9ng-tux1eeb-phuxe2n-tuxedch-tux1ed5ng-hux1ee3p}}

Một trong những đóng góp quan trọng nhất trong giai đoạn gần đây là sự
phát triển của các các phân tích tổng hợp về quan hệ quốc tế hóa và hiệu
quả, cho phép tổng hợp và định lượng hóa kết quả từ hàng trăm nghiên cứu
thực nghiệm riêng lẻ. Tổng hợp từ các các phân tích tổng hợp chính cho
thấy một bức tranh nhất quán nhưng phức tạp.

Bausch và Krist (2007) phân tích 65 nghiên cứu và tìm thấy tác động tổng
hợp dương nhưng nhỏ (r ≈ 0,10), đồng thời xác nhận rằng mức độ dị biệt
(heterogeneity) giữa các nghiên cứu rất cao với I² \textgreater{} 80\%,
hàm ý rằng các biến điều tiết là nguồn giải thích chủ yếu cho sự khác
biệt. Kirca et al.~(2012) mở rộng mẫu gộp lên 154 mẫu và tìm thấy tác
động tổng hợp tương tự, đồng thời xác nhận vai trò của bối cảnh quốc gia
và ngành. Marano et al.~(2016), nghiên cứu phân tích tổng hợp toàn diện
nhất trong tài liệu với 170 nghiên cứu và 32 quốc gia, tìm thấy r ≈
0,07--0,13 và chứng minh rằng thể chế quốc gia nguồn là biến điều tiết
quan trọng nhất trong tất cả các biến bối cảnh được kiểm định. Wu et
al.~(2022) cập nhật với mẫu gộp 218 quy mô hiệu ứng tập trung vào các
thị trường mới nổi multinationals (EMNEs) và tìm thấy mô hình biến điều
tiết tương tự. Tại hội nghị ICBEF 2024, kết quả sơ bộ từ mẫu gộp nghiên
cứu về châu Á với k = 113 nghiên cứu cho thấy r = 0,07, nhất quán với
các các phân tích tổng hợp toàn cầu nhưng cụ thể hơn về bối cảnh châu Á.
Luận án này mở rộng mẫu gộp đó thông qua Nghiên cứu P6, một phân tích
tổng hợp ba tầng (three-level MARA) với k = 238 nghiên cứu và K = 288
quy mô hiệu ứng từ châu Á và Thái Bình Dương, ước lượng r̄ = 0,074 {[}KTC
95\%: 0,060; 0,088{]}, I² = 62,4\%, Q\_M(ICRV) = 17,35 (df = 4, p =
,002). Kết quả này xác nhận tác động tổng hợp dương nhỏ nhưng bền vững,
đồng thời chứng minh rằng gradient ICRV, điểm uốn giảm dần khi chuyển từ
thể chế yếu sang thể chế mạnh, là biến điều tiết thực chất của mối quan
hệ quốc tế hóa và hiệu quả trong bối cảnh châu Á (nghiên cứu thành phần
P6). Tuy vậy, một số nghiên cứu gần đây sử dụng chiến lược nhận diện
chặt chẽ hơn đặt câu hỏi về tính phổ quát của dạng hàm chữ U ngược:
Pisani, Garcia-Bernardo và Heemskerk (2020) cho thấy quan hệ này suy yếu
đáng kể, thậm chí biến mất, trong các mẫu gộp xuyên quốc gia khi áp dụng
nhận diện nghiêm ngặt hơn, củng cố lập luận rằng chính dị biệt bối cảnh,
chứ không phải một quy luật hàm phổ quát, mới chi phối kết quả quốc tế
hóa và hiệu quả. Ở cấp doanh nghiệp, Arbelo, Arbelo-Pérez và Pérez-Gómez
(2024) chứng minh quan hệ quốc tế hóa và hiệu quả phụ thuộc vào tài sản
đặc thù của doanh nghiệp theo quan điểm dựa trên nguồn lực, nhất quán
với vai trò điều tiết của năng lực công nghệ được nhấn mạnh trong luận
án này.

Nhìn chung, bằng chứng phân tích tổng hợp cho thấy ba kết luận quan
trọng: tác động tổng hợp của quốc tế hóa lên hiệu quả là dương nhưng
khiêm tốn; mức dị biệt cực kỳ cao; và các biến điều tiết bối cảnh, đặc
biệt là thể chế quốc gia, năng lực doanh nghiệp, và đặc điểm ngành, là
những nguồn giải thích chính cho sự khác biệt giữa các nghiên cứu. Kết
luận về vai trò trung tâm của thể chế quốc gia tiếp tục được củng cố bởi
bằng chứng rất gần đây: Kalubandi và cộng sự (2025), qua một phân tích
tổng hợp quy mô lớn trên 292 mẫu độc lập với 598.525 quan sát doanh
nghiệp và năm, dùng cùng họ phương pháp hồi quy phân tích tổng hợp như
Nghiên cứu thành phần P6, phát hiện rằng thể chế quốc gia nguồn điều
tiết dương quan hệ giữa tinh thần kinh doanh quốc tế và hiệu quả, song
hành trực tiếp với giả thuyết H5 của luận án về vai trò điều tiết của
chất lượng thể chế. Đây là luận điểm cơ bản biện hộ cho sự cần thiết của
một khung giải thích đa tầng.

Giá trị của bằng chứng phân tích tổng hợp không chỉ nằm ở ước lượng tác
động tổng hợp, mà còn ở khả năng phân tách các nguồn tạo nên mức dị biệt
cao giữa các nghiên cứu, qua đó định hướng việc xây dựng lý thuyết.
Combs và cộng sự (2011) lập luận rằng tổng hợp tích lũy bằng phương pháp
phân tích tổng hợp nên được ưu tiên hơn so với đếm phiếu, vì phân tích
tổng hợp cho phép đánh giá cả độ lớn lẫn tính ổn định của hiệu ứng qua
nhiều bối cảnh. Aguinis và cộng sự (2011) cảnh báo rằng nhiều lựa chọn
phương pháp và phán đoán trong quá trình tổng hợp, từ cách mã hóa biến
điều tiết đến cách xử lý phụ thuộc giữa các quy mô hiệu ứng, có thể ảnh
hưởng đáng kể đến kết luận, do đó đòi hỏi tính minh bạch và phân tích độ
vững. Những nguyên tắc này có hàm ý trực tiếp đối với thiết kế của
Nghiên cứu thành phần P6: việc sử dụng mô hình ba tầng cho phép tách
phương sai trong nội bộ nghiên cứu khỏi phương sai giữa các nghiên cứu,
qua đó xử lý đúng cấu trúc phụ thuộc mà Aguinis và cộng sự (2011) lưu ý.
Karna và cộng sự (2016), qua một phân tích tổng hợp về quan hệ giữa năng
lực và hiệu quả tài chính, cung cấp bằng chứng rằng môi trường điều tiết
đáng kể cách năng lực chuyển hóa thành hiệu quả; phát hiện này song song
với luận điểm trung tâm của luận án rằng bối cảnh thể chế điều tiết cơ
chế chuyển hóa quốc tế hóa thành hiệu quả. Tổng hợp lại, bằng chứng phân
tích tổng hợp không chỉ định lượng mức dị biệt mà còn xác nhận rằng các
biến điều tiết bối cảnh là đối tượng lý thuyết hóa chính, từ đó biện hộ
cho việc xây dựng một khung tích hợp đa tầng thay vì kiểm định rời rạc
từng biến điều tiết đơn lẻ.

Một chiều quan trọng khác mà bằng chứng phân tích tổng hợp mở ra là vấn
đề nhận diện nhân quả trong một tài liệu chủ yếu dựa trên dữ liệu quan
sát. Shaver (2020) lập luận rằng trong nghiên cứu chiến lược và tổ chức,
nhận diện nhân quả đáng tin cậy thường đến từ một khối nghiên cứu tích
lũy hơn là từ một thiết kế đơn lẻ hoàn hảo, vì mỗi nghiên cứu xử lý một
mối đe dọa nội sinh khác nhau và sự hội tụ của bằng chứng qua nhiều bối
cảnh nâng cao độ tin cậy của suy luận. Quan điểm này biện hộ trực tiếp
cho kiến trúc nghiên cứu thành phần của luận án, trong đó nhiều nghiên
cứu đơn quốc gia với chiến lược nhận diện khác nhau hội tụ về cùng một
mô thức phi tuyến và cùng một logic điều tiết thể chế. Eden và Nielsen
(2020) bổ sung một cảnh báo phương pháp luận rằng nghiên cứu kinh doanh
quốc tế đối mặt với độ phức tạp đặc thù do tính lồng ghép đa cấp của dữ
liệu, từ doanh nghiệp đến ngành đến quốc gia, đòi hỏi các kỹ thuật ước
lượng tôn trọng cấu trúc phân cấp này; yêu cầu đó được luận án đáp ứng
bằng mô hình ba tầng trong Nghiên cứu thành phần P6 và bằng hiệu ứng cố
định nền kinh tế và năm trong các mô hình hồi quy cấp doanh nghiệp. Abdi
và Aulakh (2018) làm sâu sắc thêm khái niệm quốc tế hóa bằng cách phân
tách ba chiều mức độ, thời lượng và quy mô hoạt động ở thị trường nước
ngoài, cho thấy rằng kết luận về quan hệ quốc tế hóa và hiệu quả nhạy
cảm với chiều nào của quốc tế hóa được đo lường. Cùng hướng này, phân
tích tổng hợp gần đây của Fariborzi, Verbeke và Steel (2025), sử dụng mô
hình hóa phương trình cấu trúc trong phân tích tổng hợp trên 378 nghiên
cứu, phân biệt rõ tiền tố cấp cá nhân với tiền tố cấp doanh nghiệp và ba
chiều tốc độ, phạm vi và cường độ của quốc tế hóa, đồng thời cho thấy hệ
quả hiệu quả phụ thuộc vào chiều được xét và vào cấp độ phân tích; kết
quả này vừa biện minh cho lựa chọn đo lường nhất quán của luận án, vừa
kết nối tài liệu phân tích tổng hợp với chương trình nghị sự nền tảng vi
mô về phân tách hiệu ứng cấp cá nhân và cấp tổ chức. Nguyen và Kim
(2020), trong một lược khảo và mở rộng tài liệu về quan hệ đa quốc gia
và hiệu quả, nhấn mạnh rằng những kết luận trái ngược trong tài liệu
phần lớn bắt nguồn từ khác biệt trong thước đo, mẫu và đặc tả mô hình
hơn là từ sự không tồn tại của một quan hệ thực chất. Những đóng góp
phương pháp luận này không chỉ định khung cho việc đọc bằng chứng phân
tích tổng hợp một cách thận trọng mà còn định hình trực tiếp các lựa
chọn thiết kế của luận án, được trình bày chi tiết ở Chương 3.

\hypertarget{bux1ed1i-cux1ea3nh-chuxe2u-uxe1-vuxe0-cuxe1c-nux1ec1n-kinh-tux1ebf-mux1edbi-nux1ed5i}{%
\subsection{2.3.3 Bối cảnh châu Á và các nền kinh tế mới
nổi}\label{bux1ed1i-cux1ea3nh-chuxe2u-uxe1-vuxe0-cuxe1c-nux1ec1n-kinh-tux1ebf-mux1edbi-nux1ed5i}}

Châu Á là bối cảnh đặc biệt quan trọng và đặc thù trong tài liệu về quan
hệ quốc tế hóa và hiệu quả vì nhiều lý do. Thứ nhất, sự đa dạng thể chế
trong khu vực là ngoại lệ không có ở bất kỳ khu vực nào khác trên thế
giới: từ Singapore với chất lượng quản trị đứng hàng đầu thế giới đến
Afghanistan với chất lượng quản trị cực kỳ yếu, từ Nhật Bản với nền kinh
tế phát triển hoàn thiện đến các SIDS với quy mô siêu nhỏ và hạ tầng thể
chế còn sơ khai (Khanna \& Palepu, 2010; Peng, 2003). Thứ hai, tốc độ
chuyển đổi số và hội nhập quốc tế rất không đồng đều giữa các quốc gia
trong khu vực, tạo ra ``thí nghiệm tự nhiên'' về ảnh hưởng của bối cảnh
thể chế và năng lực số đến kết quả quốc tế hóa.

Một góc nhìn lý thuyết gần đây bổ sung chiều giải thích cho mô thức
cường độ xuất khẩu thấp đặc trưng của khu vực. Verbeke, Oh và Jain
(2025), kế thừa truyền thống nghiên cứu doanh nghiệp đa quốc gia theo
vùng, lập luận rằng phần lớn các doanh nghiệp lớn trên thế giới vẫn định
hướng nội vùng chứ không thực sự toàn cầu, với hơn một nửa doanh số tập
trung ở khu vực bản địa. Quan điểm này giúp giải thích vì sao ngay cả ở
các nền kinh tế thể chế mạnh nhất trong mẫu của luận án, như Singapore
và Nhật Bản, cường độ xuất khẩu trung bình cấp doanh nghiệp vẫn thấp và
đường cong quốc tế hóa và hiệu quả có xu hướng duỗi thẳng về gần tuyến
tính với điểm uốn nằm ngoài miền quan sát: phần lớn doanh nghiệp tập
trung ở nhánh tham gia thấp của phân phối cường độ xuất khẩu thay vì ở
vùng thâm dụng xuất khẩu cao. Các tác giả đồng thời đặt nghiên cứu doanh
nghiệp đa quốc gia trước bốn thách thức lớn đương đại là chuyển đổi số,
bất ổn địa chính trị, biến đổi khí hậu và khả năng chống chịu trước gián
đoạn quy mô lớn, bốn lực mà luận án cũng định vị là bối cảnh cấp bách
cho khu vực châu Á và Thái Bình Dương trong Mục 1.1.1.

Nghiên cứu thành phần P1 là nghiên cứu thực nghiệm đa quốc gia quy mô
lớn nhất hiện có về quan hệ quốc tế hóa và hiệu quả trong bối cảnh châu
Á mới nổi, phân tích 40.633 doanh nghiệp từ 17 nền kinh tế châu Á sử
dụng dữ liệu WBES. Nghiên cứu tìm thấy hai phát hiện chính. Thứ nhất,
tồn tại phân tầng năng suất rõ rệt giữa nhóm Đại Trung Hoa, Nam Á và
ASEAN, phản ánh sự dị biệt thể chế và năng lực giữa các nhóm. Thứ hai,
việc áp dụng công nghệ vừa làm tăng năng suất lao động trực tiếp vừa làm
suy giảm đáng kể tác động bất lợi của các rào cản thể chế, một phát hiện
được gọi là hiệu ứng lá chắn số. Nghiên cứu thành phần P2 kiểm định cụ
thể với dữ liệu doanh nghiệp nhỏ và vừa Trung Quốc và tìm thấy dạng chữ
U ngược bậc hai với ngưỡng cường độ xuất khẩu tối ưu khoảng 47,8\%, tức
điểm uốn khá thấp so với các nền kinh tế tiên tiến. Luận án kế thừa và
mở rộng các kết quả này trên mẫu đầy đủ 50 nền kinh tế, bao gồm 7 nền
kinh tế đảo nhỏ Thái Bình Dương trong mẫu chính và 9 nền khi mở rộng
kiểm định độ vững.

Phần lớn các nghiên cứu về quốc tế hóa và hiệu quả trong bối cảnh châu Á
tập trung vào Nhật Bản, Hàn Quốc, Đài Loan, và Trung Quốc, trong khi các
nền kinh tế chuyển đổi như Việt Nam, Campuchia, Bangladesh, và đặc biệt
là các SIDS, về cơ bản thiếu vắng trong tài liệu. Ngay cả trong nhóm
tiên tiến Đông Á, sự đồng nhất cũng không nên được giả định: Hemmert và
Jackson (2016), qua so sánh doanh nghiệp Nhật Bản và Hàn Quốc, cho thấy
không tồn tại một mô hình quốc tế hóa Đông Á duy nhất mà có những khác
biệt hệ thống giữa hai nền, một phát hiện củng cố cách luận án phân biệt
nội bộ nhóm thể chế thay vì gộp các nền tiên tiến thành một khối. Ở
chiều ASEAN, bằng chứng gần đây của Murhadi và cộng sự (2025) trên doanh
nghiệp Indonesia và Malaysia cho thấy quản trị doanh nghiệp và mức độ
quốc tế hóa cùng tác động đến hiệu quả tài chính, bổ sung một chiều quản
trị cho bức tranh khu vực mà luận án bao phủ ở cấp doanh nghiệp. Sự
thiếu đại diện này tạo ra khoảng trống nghiêm trọng trong khả năng khái
quát hóa lý thuyết, vì các nền kinh tế biên thể chế yếu và đảo nhỏ có
thể là địa điểm quan trọng để kiểm định điều kiện biên của các lý thuyết
về quốc tế hóa và hiệu quả.

Tài liệu về quan hệ quốc tế hóa và hiệu quả trong bối cảnh các nền kinh
tế mới nổi của châu Á đã tích lũy đủ để cho thấy rằng vai trò của bối
cảnh thể chế và năng lực doanh nghiệp là đặc biệt nổi bật ở khu vực này.
Kafouros và cộng sự (2023) chứng minh rằng chất lượng thể chế và tính
năng động của ngành cùng giải thích biến thiên hiệu quả doanh nghiệp ở
các nền kinh tế mới nổi, hàm ý rằng tác động của quốc tế hóa không thể
tách rời khỏi điều kiện thể chế và ngành mà doanh nghiệp hoạt động. Xu
(2024) tinh chỉnh quan điểm về chênh lệch thể chế bằng cách phân tách
giữa mức độ phơi nhiễm trên luật định và mức độ thực thi trên thực tế,
cho thấy rằng khoảng cách giữa quy định chính thức và thực thi thực tế
là một nguồn dị biệt quan trọng mà các thước đo thể chế tổng hợp dễ bỏ
sót. Chao và Kumar (2010) cung cấp bằng chứng rằng khoảng cách thể chế
giữa quốc gia nguồn và quốc gia đích điều tiết quan hệ giữa đa dạng quốc
tế và hiệu quả, củng cố luận điểm rằng cùng một mức độ quốc tế hóa tạo
ra kết quả khác nhau tùy theo bối cảnh thể chế. Da Cunha và cộng sự
(2023) bổ sung bằng chứng rằng thể chế chính thức ở quốc gia nguồn điều
tiết quan hệ giữa phương thức quốc tế hóa và hiệu quả, trong khi
Cuervo-Cazurra và cộng sự (2018) lập luận rằng bất định thể chế ở quốc
gia nguồn có thể bồi đắp năng lực quản lý bất định đặc thù. Tổng hợp các
bằng chứng này định vị châu Á như một bối cảnh lý tưởng để kiểm định
khung điều tiết đa tầng, đồng thời làm nổi bật rằng sự thiếu đại diện
của các nền kinh tế chuyển đổi và các nền kinh tế đảo nhỏ là một khoảng
trống thực chất chứ không chỉ là vấn đề về phạm vi mẫu.

Bằng chứng cấp doanh nghiệp ở châu Á cũng cho thấy rằng cấu trúc sở hữu
và quy mô doanh nghiệp là những nguồn dị biệt quan trọng trong quan hệ
quốc tế hóa và hiệu quả, củng cố sự cần thiết của một khung lấy biến
điều tiết làm trung tâm. Đối với doanh nghiệp nhỏ và vừa, Lu và Beamish
(2001) cùng Pangarkar (2008) cho thấy quốc tế hóa có thể nâng cao hiệu
quả nhưng theo một quan hệ phụ thuộc vào nguồn lực và mức độ cam kết thị
trường, nhất quán với logic phi tuyến của luận án. Đối với doanh nghiệp
gia đình và doanh nghiệp thuộc tập đoàn kinh doanh, vốn là hình thức tổ
chức phổ biến ở Nam Á và Đông Á, Mondal và cộng sự (2022) cùng Gaur và
Kumar (2009) cho thấy quan hệ phụ thuộc vào việc liên kết tập đoàn cung
cấp nguồn lực nội bộ thay thế cho thị trường bên ngoài còn kém phát
triển, một biểu hiện trực tiếp của logic khoảng trống thể chế ở cấp
doanh nghiệp. Pattnaik và Elango (2009) cùng Singla và George (2013)
cung cấp bằng chứng từ doanh nghiệp Ấn Độ rằng nguồn lực doanh nghiệp
điều tiết quan hệ quốc tế hóa và hiệu quả và rằng bối cảnh thể chế quốc
gia định hình quan hệ này một cách hệ thống. Ở chiều đầu tư ra nước
ngoài, Buckley và cộng sự (2007) phân tích các yếu tố quyết định đầu tư
trực tiếp ra nước ngoài của doanh nghiệp Trung Quốc, làm nổi bật vai trò
của yếu tố thể chế quốc gia nguồn trong việc định hình chiến lược quốc
tế hóa của các doanh nghiệp đa quốc gia thị trường mới nổi, trong khi
Likitwongkajon và Vithessonthi (2020) cho thấy bằng chứng từ Nhật Bản
rằng tác động của đầu tư nước ngoài lên giá trị và hiệu quả doanh nghiệp
không nhất thiết dương và phụ thuộc vào điều kiện cụ thể. Sự phong phú
nhưng phân mảnh của bằng chứng đơn quốc gia này, mỗi nghiên cứu xử lý
một loại doanh nghiệp và một bối cảnh thể chế, chính là điều biện hộ cho
cách tiếp cận tích hợp đa quốc gia của luận án: chỉ khi đặt các bối cảnh
này trên cùng một gradient thể chế và cùng một khung đo lường mới có thể
phân tách được phần dị biệt do bối cảnh thể chế tạo ra khỏi phần do đặc
trưng doanh nghiệp.

\hypertarget{vai-truxf2-cux1ee7a-nux103ng-lux1ef1c-sux1ed1-trong-mux1ed1i-quan-hux1ec7-quux1ed1c-tux1ebf-huxf3a-vuxe0-hiux1ec7u-quux1ea3}{%
\subsection{2.3.4 Vai trò của năng lực số trong mối quan hệ quốc tế hóa
và hiệu
quả}\label{vai-truxf2-cux1ee7a-nux103ng-lux1ef1c-sux1ed1-trong-mux1ed1i-quan-hux1ec7-quux1ed1c-tux1ebf-huxf3a-vuxe0-hiux1ec7u-quux1ea3}}

Sự tích hợp giữa nghiên cứu về năng lực số và quan hệ quốc tế hóa và
hiệu quả là một xu hướng tương đối mới, phát triển mạnh từ 2019 đến nay.
Bhandari et al.~(2023) phân tích 571 doanh nghiệp sản xuất của Mỹ và tìm
thấy rằng năng lực kỹ thuật số (được đo như một composite duy nhất) điều
tiết tích cực mối quan hệ quốc tế hóa và hiệu quả, với doanh nghiệp có
năng lực số cao hơn thu được lợi ích lớn hơn từ quốc tế hóa. Banalieva
và Dhanaraj (2019) phát triển lý thuyết nội bộ hóa cho nền kinh tế số,
lập luận rằng công nghệ số làm giảm chi phí phối hợp và chi phí giao
dịch xuyên biên giới, qua đó hạ thấp rào cản chi phí khi doanh nghiệp
thâm nhập thị trường nước ngoài. Chen và Meng (2022) nghiên cứu doanh
nghiệp Trung Quốc từ WBES và phát hiện rằng ràng buộc thể chế làm giảm
cường độ xuất khẩu, nhưng doanh nghiệp có năng lực số cao hơn có thể
vượt qua ràng buộc này hiệu quả hơn.

Tuy nhiên, phần lớn các nghiên cứu hiện tại không phân biệt TCI và DAI,
xử lý năng lực số như một khái niệm đơn nhất. Hậu quả lý thuyết của sự
đồng nhất không hợp lý này là không thể phân tách được hai cơ chế vận
hành khác nhau: chiều sâu năng lực tổ chức (TCI) và khả năng phối hợp kỹ
thuật số bề mặt (DAI). Sự không phân biệt này làm mờ đi cơ chế tác động
và dẫn đến kết luận không thể áp dụng trực tiếp vào thực tiễn chính sách
và quản trị doanh nghiệp.

Nền tảng lý luận cho việc phân tách TCI và DAI cần được trình bày sâu
hơn để làm rõ rằng đây không phải là một tinh chỉnh đo lường tùy nghi,
mà là một yêu cầu lý thuyết có cơ sở. Bharadwaj và cộng sự (2013) lập
luận rằng chiến lược kinh doanh số là một cấu trúc ở cấp độ chiến lược,
khác biệt với năng lực công nghệ thông tin ở cấp độ chức năng, và việc
gộp hai cấu trúc này che khuất chính cơ chế tạo giá trị mà nhà nghiên
cứu cần phân tách. Nambisan và cộng sự (2019) làm rõ thêm rằng chuyển
đổi số tác động đến đổi mới và khởi nghiệp qua nhiều cơ chế phân tầng,
từ số hóa công cụ đến tái cấu trúc mô hình kinh doanh, hàm ý rằng các
chỉ báo số ở những tầng khác nhau không thể được xem là tương đương về
mặt lý thuyết. Phân tầng này nhất quán với khung ba cấp độ của Verhoef
và cộng sự (2021), theo đó số hóa thông tin và số hóa quy trình nằm ở
tầng nông hơn so với chuyển đổi số toàn diện. Trong bối cảnh quốc tế
hóa, Li và cộng sự (2022) cung cấp bằng chứng từ doanh nghiệp sản xuất
Trung Quốc rằng số hóa điều tiết quan hệ quốc tế hóa và hiệu quả, song
nghiên cứu này, giống như đa số tài liệu hiện hành, vẫn xử lý năng lực
số như một cấu trúc tổng hợp. Knight và Cavusgil (2004) bổ sung một góc
nhìn lịch sử quan trọng rằng năng lực tổ chức và đổi mới cho phép một số
doanh nghiệp quốc tế hóa sớm và nhanh, cho thấy rằng chiều sâu năng lực
nội tại và công cụ giao diện đóng vai trò khác nhau trong quá trình quốc
tế hóa. Tổng hợp các luận điểm này, luận án lập luận rằng nếu TCI vận
hành như chiều sâu năng lực hấp thụ và tái cấu hình, còn DAI vận hành
như công cụ phối hợp bề mặt, thì việc gộp hai cấu trúc sẽ trung bình hóa
hai hiệu ứng có thể trái chiều, dẫn đến ước lượng chệch và kết luận
chính sách không thể áp dụng (Coltman et al., 2008; Bharadwaj et al.,
2013).

Cần đặt cơ chế của năng lực số trong một bức tranh rộng hơn về các ràng
buộc thực chất mà doanh nghiệp ở các nền kinh tế đang phát triển đối mặt
khi quốc tế hóa, vì chính những ràng buộc này tạo ra dư địa để công cụ
số phát huy hoặc không phát huy tác dụng. Tài liệu kinh tế học thương
mại cho thấy ràng buộc tín dụng là một trong những rào cản nghiêm trọng
nhất đối với xuất khẩu: Manova (2013) chứng minh trên cơ sở lý thuyết và
thực nghiệm rằng các doanh nghiệp ở những quốc gia có hệ thống tài chính
kém phát triển bị hạn chế khả năng tham gia thương mại quốc tế do không
huy động được vốn lưu động cho chi phí cố định của xuất khẩu, trong khi
Demir và Javorcik (2018) cho thấy tín dụng thương mại có thể thay thế
một phần cho tín dụng ngân hàng và nới lỏng ràng buộc này. Trong khung
của luận án, các phát hiện này có hàm ý cơ chế rõ ràng: ở những bối cảnh
mà ràng buộc tài chính và khoảng trống thể chế bóp nghẹt khả năng chuyển
hóa quốc tế hóa thành hiệu quả, công cụ số ở giao diện ngoại tại có thể
giảm chi phí giao dịch và mở rộng khả năng tiếp cận thông tin thị
trường, qua đó bù đắp một phần cho hạ tầng thể chế và tài chính còn
thiếu, đúng theo logic lá chắn số được phát biểu trong giả thuyết H5.
Avenyo và cộng sự (2021) bổ sung bằng chứng từ doanh nghiệp châu Phi
rằng năng lực sản xuất nội tại quyết định kết quả xuất khẩu, song song
với phát hiện của Ayyagari và cộng sự (2011) rằng đổi mới ở các nền kinh
tế mới nổi gắn chặt với năng lực hấp thụ và môi trường thể chế. Tổng hợp
các bằng chứng này, luận án định vị TCI và DAI không phải như những yếu
tố tách rời khỏi cấu trúc kinh tế nền, mà như hai cơ chế tương tác với
ràng buộc tài chính và thể chế đặc thù của từng chế độ ICRV, qua đó cung
cấp cơ sở lý thuyết cho dự đoán rằng hiệu ứng điều tiết của năng lực số
mạnh nhất ở biên dưới của gradient thể chế.

\hypertarget{bux1eb1ng-chux1ee9ng-vux129-muxf4-vux1ec1-chux1ea5t-lux1b0ux1ee3ng-thux1ec3-chux1ebf-vuxe0-phuxe2n-huxf3a-kux1ebft-quux1ea3-tux103ng-trux1b0ux1edfng}{%
\subsection{2.3.5 Bằng chứng vĩ mô về chất lượng thể chế và phân hóa kết
quả tăng
trưởng}\label{bux1eb1ng-chux1ee9ng-vux129-muxf4-vux1ec1-chux1ea5t-lux1b0ux1ee3ng-thux1ec3-chux1ebf-vuxe0-phuxe2n-huxf3a-kux1ebft-quux1ea3-tux103ng-trux1b0ux1edfng}}

Bên cạnh dòng nghiên cứu cấp doanh nghiệp, một khối bằng chứng vĩ mô
đang tích lũy cung cấp nền tảng bối cảnh quan trọng cho luận điểm điều
tiết thể chế của luận án. Phân tích hệ thống đầu tiên của Ngân hàng Thế
giới về nhóm thị trường biên định nghĩa nhóm này là tập con của các nền
kinh tế đang phát triển và mới nổi có mức độ tiếp cận thị trường tài
chính quốc tế ``có ý nghĩa nhưng hạn chế'', nằm giữa các thị trường mới
nổi và các nền kinh tế đang phát triển khác (World Bank, 2026c). Điều
đáng chú ý về mặt phương pháp luận là chính báo cáo này thừa nhận rằng
``thị trường biên với tư cách một tiểu nhóm có xu hướng bị bỏ qua trong
các phân tích kinh tế'' (World Bank, 2026c), một sự thừa nhận trực tiếp
về khoảng trống mà luận án này lấp ở cấp doanh nghiệp đối với khu vực
châu Á--Thái Bình Dương.

Phát hiện cốt lõi của khối bằng chứng vĩ mô này là chất lượng thể chế
phân hóa kết quả tăng trưởng một cách có hệ thống. Các thị trường biên
tăng trưởng nhanh nhất trong một phần tư thế kỷ qua chia sẻ một số đặc
điểm chung: cải thiện quản trị, thu hẹp chênh lệch lãi suất cho vay--huy
động của ngân hàng, kiểm soát nợ công và gánh nặng trả nợ, cùng với tích
lũy vốn trên đầu người mạnh hơn (World Bank, 2026c). Ngược lại, các điểm
yếu thể chế và quản trị, biểu hiện qua thị trường tài chính nội địa nông
và chênh lệch lãi suất rộng, đi kèm với tăng trưởng vốn trên đầu người
đáng thất vọng và làm khuếch đại tác động bất lợi của các cú sốc (World
Bank, 2026c). Cơ chế vĩ mô này nhất quán với lập luận
institutional-khoảng trống ở cấp doanh nghiệp (Khanna \& Palepu, 2010):
thể chế yếu làm tăng chi phí giao dịch và bất định, từ đó bào mòn khả
năng chuyển hóa đầu tư và quốc tế hóa thành hiệu quả thực.

Quan trọng không kém, bằng chứng vĩ mô cho thấy không tồn tại một con
đường phát triển phổ quát: các thị trường biên thành công đã áp dụng
những chiến lược rất khác nhau, từ thu hút đầu tư vào năng lượng và
khoáng sản (Kazakhstan), đến sản xuất chế tạo giá trị gia tăng và xuất
khẩu (Việt Nam), đến tăng trưởng dựa trên dịch vụ và du lịch (Rwanda);
và bốn nền kinh tế (Bulgaria, Costa Rica, Panama, Romania) đã đạt thu
nhập cao và ``tốt nghiệp'' khỏi nhóm biên trong giai đoạn 2012--2025
(World Bank, 2026c). Sự đa dạng con đường này củng cố trực tiếp tiền đề
của khung ICRV: cơ chế chuyển hóa quốc tế hóa thành hiệu quả phụ thuộc
vào chế độ thể chế và mô hình tăng trưởng đặc thù, chứ không tuân theo
một dạng hàm phổ quát. Tuy nhiên, bằng chứng vĩ mô này dừng ở cấp quốc
gia và để ngỏ câu hỏi về cơ chế cấp doanh nghiệp, chính là khoảng trống
mà luận án giải quyết bằng dữ liệu WBES cấp cơ sở trên 50 nền kinh tế.

\begin{center}\rule{0.5\linewidth}{0.5pt}\end{center}

\hypertarget{khoux1ea3ng-trux1ed1ng-nghiuxean-cux1ee9u}{%
\section{2.4 Khoảng trống nghiên
cứu}\label{khoux1ea3ng-trux1ed1ng-nghiuxean-cux1ee9u}}

Trên cơ sở tổng quan lý thuyết và thực nghiệm ở các phần trước, ba
khoảng trống nghiên cứu chính được xác định như sau.

\textbf{Khoảng trống thứ nhất: Thiếu khung tích hợp đa tầng để giải
thích dị biệt trong quan hệ quốc tế hóa và hiệu quả.} Các phân tích tổng
hợp trước đây thường báo cáo mức dị biệt rất cao (\(I^2\) thường trên
80\%) (Bausch \& Krist, 2007; Marano et al., 2016), và ngay phân tích
tổng hợp cập nhật của luận án (nghiên cứu thành phần P6, Mục 4.2) cũng
ghi nhận mức dị biệt đáng kể (\(I^2 = 62{,}4\%\)); ở cả hai trường hợp,
phần lớn biến thiên trong cỡ ảnh hưởng không được giải thích bởi sai số
ngẫu nhiên mà bởi sự khác biệt có hệ thống giữa các bối cảnh. Tuy nhiên,
hầu hết các nghiên cứu thực nghiệm hiện tại chỉ kiểm định một hoặc hai
biến điều tiết từ một lý thuyết đơn lẻ, không có khả năng giải thích
đồng thời dị biệt do bối cảnh thể chế, năng lực doanh nghiệp và đặc điểm
lãnh đạo tạo ra. Không một lý thuyết đơn lẻ nào, dù là mô hình Uppsala,
lý thuyết dựa trên nguồn lực, lý thuyết thể chế hay lý thuyết thượng
tầng quản trị, đủ để giải thích đầy đủ sự phức tạp của quan hệ này trong
bối cảnh đa quốc gia đa dạng như châu Á. Mỗi lý thuyết soi sáng một tầng
nguyên nhân khác nhau: mô hình Uppsala giải thích chi phí học hỏi giai
đoạn đầu, lý thuyết dựa trên nguồn lực giải thích khác biệt năng lực, lý
thuyết thể chế giải thích biến thiên theo môi trường nền, và lý thuyết
thượng tầng quản trị giải thích vai trò của người ra quyết định. Vì các
tầng nguyên nhân này vận hành đồng thời, một khung tích hợp đa tầng,
trong đó các lý thuyết bổ sung cho nhau thay vì cạnh tranh nhau, là điều
kiện cần để giải thích mức dị biệt cao mà các phân tích tổng hợp ghi
nhận.

\textbf{Khoảng trống thứ hai: Sự đồng nhất không hợp lý giữa TCI và DAI
trong phần lớn tài liệu trước.} Phần lớn các nghiên cứu về năng lực số
trong bối cảnh quốc tế hóa và hiệu quả hoặc là không đo lường năng lực
số, hoặc là gộp tất cả các chỉ báo số vào một chỉ số hợp thành duy nhất,
không phân biệt giữa chiều sâu năng lực công nghệ tổ chức (TCI) và mức
độ chấp nhận công cụ số bề mặt (DAI). Hậu quả thực tiễn của sự đồng nhất
này là: (i) không thể trả lời câu hỏi liệu nên đầu tư vào năng lực công
nghệ nội tại hay vào việc áp dụng công cụ số giao diện; (ii) không thể
phân biệt được hai cơ chế tác động có hàm ý chính sách khác nhau. Phân
biệt rõ TCI và DAI là điều kiện cần để rút ra hàm ý quản trị và chính
sách có tính cụ thể và ứng dụng cao.

\textbf{Khoảng trống thứ ba: Thiếu bằng chứng đa quốc gia quy mô lớn,
toàn diện và có cấu trúc cho bối cảnh châu Á.} Mặc dù châu Á là khu vực
kinh tế năng động và đa dạng nhất thế giới, tài liệu về quan hệ quốc tế
hóa và hiệu quả vẫn thiếu bằng chứng đa quốc gia bao quát đủ rộng để
kiểm định khả năng tổng quát hóa của các lý thuyết. Mẫu 17 nền kinh tế
của nghiên cứu thành phần P1 là cơ sở lớn nhất hiện có trong tài liệu về
châu Á, nhưng vẫn chưa bao gồm các nền kinh tế biên thể chế yếu và các
nền kinh tế đảo nhỏ đang phát triển, những địa điểm cực kỳ quan trọng để
kiểm định điều kiện biên của lý thuyết. Thiếu vắng các bằng chứng này
khiến các lý thuyết về quan hệ quốc tế hóa và hiệu quả, vốn được xây
dựng chủ yếu trên dữ liệu từ Nhật Bản, Hàn Quốc và các nền kinh tế tiên
tiến, có nguy cơ bị áp dụng sai trong bối cảnh các nền kinh tế đang phát
triển và chuyển đổi ở châu Á. Điều kiện biên là nơi một lý thuyết bộc lộ
giới hạn của nó; do đó việc đưa các nền kinh tế ở hai cực của gradient
thể chế vào khung phân tích không chỉ mở rộng phạm vi mà còn là phép thử
nghiêm ngặt nhất đối với tính phổ quát của dạng hàm chữ U ngược.

Ba khoảng trống này hội tụ vào một hướng giải quyết chung: cần phát
triển và kiểm định một khung giải thích tích hợp đa tầng trên dữ liệu đa
quốc gia đủ rộng và đủ đa dạng về thể chế để bao phủ toàn bộ gradient
thể chế của châu Á. Luận án đề xuất khung CDCM (điều tiết theo quốc gia,
năng lực số và năng lực doanh nghiệp), là khung lý thuyết tích hợp đa
tầng kết hợp bối cảnh thể chế (tầng vĩ mô, vận hành hóa bằng ICRV), năng
lực doanh nghiệp với TCI và DAI tách biệt (tầng tổ chức), đặc điểm nhà
quản trị cấp cao (tầng cá nhân), và lớp mở rộng lăng kính năng lực số.
Trong khung này, cường độ xuất khẩu (FSTS) là biến độc lập chính và năng
suất lao động là biến phụ thuộc.

\textbf{Từ khoảng trống đến nghiên cứu P7: kiểm định CDCM trên 50 nền
kinh tế châu Á.} Khoảng trống thứ nhất và thứ ba hội tụ vào một câu hỏi
kiểm định trực tiếp: liệu khung CDCM có thể giải thích đồng thời dị biệt
của quan hệ quốc tế hóa và hiệu quả khi kiểm soát đầy đủ cả ba tầng
trong một mô hình tích hợp ở quy mô đủ rộng để bao phủ toàn bộ gradient
ICRV hay không. Tài liệu hiện tại chưa có bằng chứng kiểm định đồng thời
năm biến điều tiết chính, gồm TCI, DAI, kinh nghiệm quản trị, giới tính
nhà quản trị và ICRV, trong một mô hình hồi quy đơn lẻ trên dữ liệu đa
quốc gia quy mô lớn. Nghiên cứu thành phần P7 lấp đầy khoảng trống này
bằng dữ liệu WBES từ 50 nền kinh tế (mẫu hồi quy bậc hai \(N = 81.022\);
mẫu mô hình đầy đủ \(N = 79.080\) sau khi loại bỏ theo danh sách các
biến điều tiết khuyết, 103 cặp nền kinh tế và năm), bao phủ toàn bộ
gradient ICRV từ nhóm tiên tiến đến nhóm đảo nhỏ, kiểm định đồng thời H1
đến H6 với hiệu ứng cố định nền kinh tế và năm. Điểm uốn hiệu chỉnh năng
lực của P7 là 43,6\% FSTS, nằm trong vùng chuyển đổi của gradient thể
chế, nhất quán với P3 Việt Nam (khoảng 40\%) và thấp hơn P5 Trung Quốc
(khoảng 49\%); cấu trúc ba vùng theo nhóm ICRV xác nhận thể chế điều
tiết cả vị trí lẫn sự tồn tại của điểm uốn theo gradient dự đoán lý
thuyết, đóng góp bằng chứng kiểm định tích hợp đầu tiên cho khung CDCM ở
cấp độ liên quốc gia quy mô lớn.

\textbf{Từ khoảng trống đến nghiên cứu P8: kiểm định điều kiện biên FIP
tại các nền kinh tế đảo nhỏ.} Khoảng trống thứ ba không chỉ đòi hỏi mở
rộng phạm vi địa lý mà còn đòi hỏi kiểm định điều kiện biên lý thuyết
cực trị: liệu có tồn tại bối cảnh nơi chữ U ngược hoàn toàn sụp đổ thành
quan hệ âm đơn điệu hay không. Không có nghiên cứu nào trong tài liệu
hiện tại kiểm định hình phạt quốc tế hóa cưỡng bức (FIP), bối cảnh mà
quốc tế hóa là gánh nặng sinh tồn chứ không phải lựa chọn chiến lược dựa
trên lợi thế cạnh tranh. Nghiên cứu thành phần P8 kiểm định H1b bằng dữ
liệu WBES từ các nền kinh tế đảo nhỏ Thái Bình Dương (Fiji, Papua New
Guinea, Quần đảo Solomon, Tonga, Vanuatu, Samoa, Kiribati), nơi ba điều
kiện cấu trúc cơ bản của chữ U ngược đồng thời bị vi phạm: thị trường
nội địa cực nhỏ, chi phí vận tải cao bất thường, và hỗ trợ thể chế cho
xuất khẩu gần như không tồn tại. Kết quả P8 xác định điều kiện biên cực
trị của khung CDCM, tức điểm nơi lý thuyết đường cong chữ U ngược không
còn ứng dụng được.

\begin{center}\rule{0.5\linewidth}{0.5pt}\end{center}

\hypertarget{muxf4-huxecnh-nghiuxean-cux1ee9u-vuxe0-hux1ec7-giux1ea3-thuyux1ebft}{%
\section{2.5 Mô hình nghiên cứu và hệ giả
thuyết}\label{muxf4-huxecnh-nghiuxean-cux1ee9u-vuxe0-hux1ec7-giux1ea3-thuyux1ebft}}

\hypertarget{muxf4-huxecnh-khuxe1i-niux1ec7m}{%
\subsection{2.5.1 Mô hình khái
niệm}\label{muxf4-huxecnh-khuxe1i-niux1ec7m}}

Mô hình khái niệm của luận án vận hành theo logic đa tầng, trong đó mối
quan hệ trung tâm giữa quốc tế hóa (FSTS) và hiệu quả hoạt động doanh
nghiệp (labor productivity) được điều tiết đồng thời bởi ba tầng can
thiệp.

Tầng vĩ mô, bối cảnh thể chế, được vận hành hóa bằng khung ICRV với sáu
chế độ con theo gradient từ tiên tiến đổi mới đến SIDS. Ở tầng này, ICRV
điều tiết mối quan hệ quốc tế hóa và hiệu quả thông qua hai cơ chế: (i)
mức độ khoảng trống thể chế ảnh hưởng đến chi phí giao dịch xuyên biên
giới; và (ii) chất lượng thể chế ảnh hưởng đến khả năng doanh nghiệp
khai thác lợi ích từ quốc tế hóa. Ở các nền kinh tế biên thể chế yếu và
đảo nhỏ, khoảng trống thể chế cực kỳ lớn có thể đảo ngược quan hệ quốc
tế hóa và hiệu quả, biến lợi ích tiềm năng thành tổn thất thực tế
(nghiên cứu thành phần P8).

Tầng tổ chức, năng lực doanh nghiệp, bao gồm TCI và DAI như hai biến
điều tiết độc lập với cơ chế tác động khác nhau. TCI đóng vai trò như
một ``nâng đỡ mặt bằng'' (level shifter): doanh nghiệp có TCI cao sẽ đạt
hiệu quả cao hơn ở mọi mức độ FSTS, và đặc biệt có khả năng khai thác
lợi ích quốc tế hóa tốt hơn nhờ năng lực hấp thụ tri thức và công nghệ
từ thị trường nước ngoài. DAI, ngược lại, đóng vai trò là ``bộ đệm phối
hợp số'' (digital coordination buffer): ở mức FSTS cao, DAI giúp giảm
chi phí điều phối xuyên biên giới, làm chuyển dịch điểm uốn sang phải và
kéo dài giai đoạn lợi ích dương của quan hệ quốc tế hóa và hiệu quả.

Tầng cá nhân, đặc điểm nhà quản trị cấp cao, bao gồm kinh nghiệm quản
trị và giới tính nhà quản trị như các biến điều tiết vi mô. Kinh nghiệm
quốc tế của nhà quản trị giảm chi phí nhận thức của gánh nặng người nước
ngoài; sự hiện diện của quản trị viên nữ tăng cường đa dạng nhận thức và
thận trọng trong quản trị rủi ro quốc tế.

Ngoài các cơ chế điều tiết, mô hình còn bao gồm một yếu tố động thái
quan trọng: hình dạng phi tuyến của quan hệ quốc tế hóa và hiệu quả
(điểm uốn của chữ U ngược) không cố định theo thời gian mà có thể thay
đổi khi bối cảnh thể chế và năng lực doanh nghiệp thay đổi theo chiều
dọc, đặc biệt rõ rệt ở các nền kinh tế chuyển đổi nhanh như Trung Quốc.

\begin{figure}
\centering
\includegraphics{figures/fig_2_1_cdcm_conceptual_model.png}
\caption{Hình 2.1. Khung khái niệm tích hợp đa tầng CDCM cho quan hệ
Quốc tế hóa--Hiệu quả tại châu Á}
\end{figure}

\textbf{Hình 2.1.} \emph{Khung khái niệm tích hợp đa tầng CDCM.} Mũi tên
liền thể hiện tác động trực tiếp (H1, quan hệ phi tuyến chữ U ngược);
mũi tên đứt thể hiện tác động điều tiết của ba tầng (thể chế ICRV, H5;
năng lực số TCI/DAI, H2, H3; đặc điểm nhà quản trị, H4); mũi tên chấm
thể hiện điều kiện biên FIP (H1b) tại nhóm SIDS. Biến kiểm soát không
thể hiện để giữ rõ ràng. \emph{Nguồn: tác giả xây dựng dựa trên mô hình
Uppsala, lý thuyết dựa trên nguồn lực, lý thuyết thể chế, lý thuyết
thượng tầng quản trị và lăng kính năng lực số.}

\hypertarget{giux1ea3-thuyux1ebft-h1-phi-tuyux1ebfn-chux1eef-u-ngux1b0ux1ee3c}{%
\subsection{2.5.2 Giả thuyết H1: Phi tuyến chữ U
ngược}\label{giux1ea3-thuyux1ebft-h1-phi-tuyux1ebfn-chux1eef-u-ngux1b0ux1ee3c}}

Trên cơ sở lý thuyết từ Uppsala về chi phí thâm nhập giai đoạn đầu, lý
thuyết dựa trên nguồn lực về lợi thế cạnh tranh tích lũy ở giai đoạn
trung, và lý thuyết chi phí giao dịch về sự phức tạp điều phối ở giai
đoạn mở rộng quá mức, giả thuyết H1 được phát biểu như sau. Quan hệ giữa
mức độ quốc tế hóa và hiệu quả hoạt động của các doanh nghiệp ở các quốc
gia châu Á có dạng phi tuyến chữ U ngược, với một điểm uốn tối ưu trước
khi chi phí điều phối và bất định vượt qua lợi ích từ mở rộng thị trường
quốc tế. Cơ sở lý thuyết trực tiếp bao gồm mô hình đường cong chữ S của
Lu và Beamish (2004), lý thuyết ba giai đoạn của Contractor et
al.~(2003), và bằng chứng chữ U ngược ban đầu của Hitt et al.~(1997).

\hypertarget{b-ux111iux1ec1u-kiux1ec7n-biuxean-h1b-guxe1nh-nux1eb7ng-quux1ed1c-tux1ebf-huxf3a-bux1eaft-buux1ed9c-fip-tux1ea1i-sids}{%
\subsection{2.5.2b Điều kiện biên H1b: Gánh nặng quốc tế hóa bắt buộc
(FIP) tại
SIDS}\label{b-ux111iux1ec1u-kiux1ec7n-biuxean-h1b-guxe1nh-nux1eb7ng-quux1ed1c-tux1ebf-huxf3a-bux1eaft-buux1ed9c-fip-tux1ea1i-sids}}

H1 dự báo quan hệ chữ U ngược, tức là tồn tại một mức FSTS tối ưu mà
dưới ngưỡng đó quốc tế hóa tạo ra lợi ích, và trên ngưỡng đó chi phí
điều phối vượt qua lợi ích học hỏi. Tuy nhiên, lý thuyết dự đoán này
ngầm giả định rằng ba điều kiện cấu trúc cơ bản được đáp ứng: (i) tồn
tại thị trường nội địa đủ lớn để doanh nghiệp có nền tảng hiệu suất và
thanh khoản trước khi xuất khẩu; (ii) chi phí thương mại ở mức chấp nhận
được, cho phép cân nhắc kinh tế hợp lý; và (iii) tồn tại hỗ trợ thể chế
tối thiểu cho hoạt động trao đổi quốc tế.

Tại các nền kinh tế SIDS (Nhóm VI ICRV), cả ba điều kiện này đồng thời
bị vi phạm: thị trường nội địa cực nhỏ buộc các doanh nghiệp phải xuất
khẩu không phải vì có lợi thế cạnh tranh mà vì thiếu cầu nội địa đủ lớn
để duy trì hoạt động (xuất khẩu bắt buộc); chi phí hậu cần và vận chuyển
cực cao do địa lý đảo xa và phân tán làm suy giảm biên lợi nhuận ở mọi
mức FSTS; và hỗ trợ thể chế (logistics, tài chính thương mại, năng lực
xúc tiến xuất khẩu) thiếu hụt nghiêm trọng. Khi ba điều kiện này đồng
thời không được đáp ứng, quan hệ quốc tế hóa và hiệu quả không chỉ mất
đi phần lợi ích giai đoạn đầu của chữ U ngược mà có thể chuyển thành
\textbf{quan hệ âm đơn điệu không có điểm uốn}, tức là điều kiện biên lý
thuyết nơi mọi mức độ quốc tế hóa đều gây tổn hại đến hiệu quả doanh
nghiệp.

Giả thuyết H1b được phát biểu như sau. Tại các nền kinh tế SIDS (ICRV
Nhóm VI), mối quan hệ giữa mức độ quốc tế hóa (FSTS) và hiệu quả hoạt
động doanh nghiệp (năng suất lao động) là \textbf{âm đơn điệu}, không có
điểm uốn trong phạm vi dữ liệu, phản ánh gánh nặng quốc tế hóa bắt buộc
(Forced Internationalization Penalty, FIP) nơi doanh nghiệp xuất khẩu để
sinh tồn chứ không phải để khai thác lợi thế cạnh tranh. Cơ sở lý thuyết
bao gồm lý thuyết thể chế (North, 1990), lý thuyết chi phí giao dịch
(Williamson, 1985), Khanna và Palepu (2010) về khoảng trống thể chế, và
Lall (2004) về các ràng buộc cấu trúc của các nền kinh tế đảo nhỏ.

Một yêu cầu lý thuyết quan trọng đối với H1b là phân biệt FIP như một
điều kiện biên thực chất với khả năng quan hệ âm đơn điệu chỉ là sản
phẩm của cỡ mẫu nhỏ hoặc miền FSTS bị giới hạn. Sự phân biệt này có thể
bác bỏ được về mặt thực nghiệm: nếu FIP là một cấu trúc thực, kiểm định
Lind và Mehlum (2010) sẽ không phát hiện một cực đại nội vùng ngay cả
khi độ phủ dữ liệu đủ rộng, và dấu của hiệu ứng biên phải âm ở cả hai
biên của miền quan sát; ngược lại, nếu quan hệ âm chỉ phản ánh thiếu sức
mạnh thống kê hay cắt cụt phạm vi, một cực đại nội vùng sẽ lộ ra khi mở
rộng độ phủ. Đặc thù phân phối FSTS tại SIDS củng cố cách đọc cấu trúc:
khối lượng doanh nghiệp tập trung ở vùng xuất khẩu cao do cưỡng bức xuất
khẩu chứ không phải ở vùng FSTS thấp, nên một kết quả âm đơn điệu không
thể quy cho việc thiếu quan sát ở đuôi trên của miền. Do đó H1b được
phát biểu như một mệnh đề có thể kiểm sai: việc phát hiện một điểm uốn
nội vùng có ý nghĩa tại SIDS khi dữ liệu đủ phủ sẽ phủ định FIP, đặt
khái niệm này dưới một tiêu chí nhận dạng minh bạch chứ không phải một
diễn giải hậu nghiệm.

\hypertarget{c-ux111iux1ec1u-kiux1ec7n-biuxean-h1c-tan-biux1ebfn-ngux1b0ux1ee1ng-dux1b0ux1edbi-biux1ebfn-ux111ux1ed9ng-thux1ec3-chux1ebf-tuxedch-lux169y}{%
\subsection{2.5.2c Điều kiện biên H1c: Tan biến ngưỡng dưới biến động
thể chế tích
lũy}\label{c-ux111iux1ec1u-kiux1ec7n-biuxean-h1c-tan-biux1ebfn-ngux1b0ux1ee1ng-dux1b0ux1edbi-biux1ebfn-ux111ux1ed9ng-thux1ec3-chux1ebf-tuxedch-lux169y}}

Trong khi H6 dự báo ngưỡng (điểm uốn) \textbf{dịch chuyển} một cách có
trật tự theo quá trình trưởng thành năng lực và thể chế, một khả năng
đối lập về mặt lý thuyết là ngưỡng \textbf{tan biến} dưới một chuỗi cú
sốc thể chế dồn dập, khi đó cấu trúc chữ U ngược không còn được nhận
diện. Giả thuyết H1c được phát biểu như sau. Ở nền kinh tế trải qua biến
động thể chế tích lũy cực đoan qua các giai đoạn (điển hình là Ấn Độ
trong nghiên cứu thành phần P9), cấu trúc chữ U ngược của quan hệ quốc
tế hóa và hiệu quả \textbf{tan biến theo thời gian} (điểm uốn dịch
chuyển rồi biến mất trong miền dữ liệu), phản ánh sự phá vỡ cấu trúc chứ
không phải dịch chuyển tiệm tiến. H1c là điều kiện biên thứ hai của H1
(bên cạnh H1b), phân biệt rõ với H6 ở chỗ H6 mô tả dịch chuyển có trật
tự còn H1c mô tả tan rã cấu trúc. Cơ sở lý thuyết bao gồm North (1990)
về bất định thể chế, Banalieva và Dhanaraj (2019), và bằng chứng đa đợt
của P9.

\hypertarget{giux1ea3-thuyux1ebft-h2-ux111iux1ec1u-tiux1ebft-bux1edfi-nux103ng-lux1ef1c-cuxf4ng-nghux1ec7-tci}{%
\subsection{2.5.3 Giả thuyết H2: Điều tiết bởi năng lực công nghệ
(TCI)}\label{giux1ea3-thuyux1ebft-h2-ux111iux1ec1u-tiux1ebft-bux1edfi-nux103ng-lux1ef1c-cuxf4ng-nghux1ec7-tci}}

Dựa trên lý thuyết dựa trên nguồn lực (Barney, 1991) và lý thuyết năng
lực động (Teece et al., 1997), năng lực công nghệ là một nguồn lực thỏa
mãn tiêu chí VRIN giúp doanh nghiệp khai thác lợi ích từ quốc tế hóa
hiệu quả hơn. Giả thuyết H2 được phát biểu như sau. Năng lực công nghệ
(TCI) làm gia tăng tác động tích cực của quốc tế hóa lên hiệu quả hoạt
động của doanh nghiệp, bằng cách nâng mặt bằng hiệu quả tổng thể và làm
giảm chi phí học hỏi thị trường quốc tế, do đó làm dịch chuyển đường
quan hệ quốc tế hóa và hiệu quả lên cao hơn và có thể điều chỉnh vị trí
của điểm uốn. Cơ sở lý thuyết bao gồm Barney (1991), Teece et
al.~(1997), và Bhandari et al.~(2023).

Cơ chế lý thuyết đứng sau H2 là năng lực hấp thụ: doanh nghiệp có năng
lực công nghệ cao hơn xử lý tín hiệu thị trường quốc tế hiệu quả hơn,
chuyển hóa tri thức bên ngoài thành cải thiện năng suất nhanh hơn, và do
đó khấu hao chi phí gánh nặng người nước ngoài với tốc độ cao hơn (Cohen
\& Levinthal, 1990). Vì cơ chế này vận hành chủ yếu qua việc nâng chiều
sâu năng lực nội tại chứ không qua việc thay đổi bản chất đánh đổi của
quốc tế hóa, dự đoán mạnh nhất là một hiệu ứng nâng mặt bằng phổ quát,
trong khi hiệu ứng uốn lại đường cong, nếu có, sẽ phụ thuộc vào sự bổ
trợ thể chế. Dự đoán này phù hợp với bằng chứng định hướng từ các nghiên
cứu thành phần đơn quốc gia, nơi năng lực công nghệ liên tục xuất hiện
như yếu tố nâng mặt bằng dương và có ý nghĩa, sẽ được kiểm định hệ thống
ở Chương 4. Về mặt nhận dạng thực nghiệm, H2 hàm ý một mẫu hệ số phân
biệt được: hệ số chính của TCI dương và có ý nghĩa (nâng mặt bằng),
trong khi số hạng tương tác độ cong TCI × FSTS² không có ý nghĩa hoặc
yếu (không uốn lại đường cong). Chính sự kết hợp giữa hiệu ứng mức mạnh
và hiệu ứng độ cong yếu này là tiêu chí thực nghiệm cho phép tách H2
(nâng mặt bằng) khỏi H3 (uốn đường cong) ngay trong cùng một mô hình
tương tác.

\hypertarget{giux1ea3-thuyux1ebft-h3-ux111iux1ec1u-tiux1ebft-bux1edfi-chux1ec9-sux1ed1-chux1ea5p-nhux1eadn-sux1ed1-dai}{%
\subsection{2.5.4 Giả thuyết H3: Điều tiết bởi chỉ số chấp nhận số
(DAI)}\label{giux1ea3-thuyux1ebft-h3-ux111iux1ec1u-tiux1ebft-bux1edfi-chux1ec9-sux1ed1-chux1ea5p-nhux1eadn-sux1ed1-dai}}

Phân biệt rõ với H2, DAI không có chiều sâu năng lực tổ chức như TCI
nhưng có vai trò giảm thiểu chi phí phối hợp xuyên biên giới ở mức xuất
khẩu cao. Giả thuyết H3 được phát biểu như sau. Chỉ số chấp nhận số
(DAI) điều tiết mối quan hệ quốc tế hóa và hiệu quả, với vai trò bổ trợ
rõ rệt hơn ở mức cường độ xuất khẩu cao khi DAI giúp giảm chi phí phối
hợp giao dịch xuyên biên giới, theo đó DAI làm thay đổi độ dốc của đường
cong chứ không nhất thiết nâng mặt bằng hiệu quả như TCI. Cơ sở lý
thuyết bao gồm Verhoef et al.~(2021), Stallkamp và Schotter (2021), và
Banalieva và Dhanaraj (2019). Sự phân biệt cơ chế tác động giữa H2 và H3
là đóng góp lý thuyết quan trọng, cho phép phân tách TCI và DAI trong mô
hình hồi quy thực nghiệm.

Cơ chế đứng sau H3 là sự nén chi phí phối hợp ở nhánh đi xuống của đường
cong. Khi cường độ xuất khẩu tăng cao, chi phí điều phối nhiều thị
trường, múi giờ và chuỗi cung ứng phân tán là một nguồn chính kéo năng
suất đi xuống; công cụ số làm giảm chính loại chi phí này thông qua kênh
giao tiếp, theo dõi đơn hàng và phối hợp số hóa. Vì lý do đó, hiệu ứng
dự kiến của DAI tập trung ở việc làm thoải nhánh đi xuống và đẩy điểm
uốn về phía phải, tức một hiệu ứng uốn đường cong, khác về bản chất với
hiệu ứng nâng mặt bằng của TCI. Tương ứng, H3 dự đoán một mẫu hệ số đối
ngược với H2: số hạng tương tác DAI × FSTS² (hoặc DAI × FSTS) có ý
nghĩa, biểu thị thay đổi độ cong hoặc độ dốc, trong khi hiệu ứng mức
chính của DAI không nhất thiết mạnh. Cặp tiêu chí nhận dạng đối ngược
này, hiệu ứng mức không kèm uốn cong cho H2 so với uốn cong không nhất
thiết kèm hiệu ứng mức cho H3, khiến hai giả thuyết có thể phân biệt và
bác bỏ độc lập trong khung tương tác, thay vì chỉ là hai cách diễn đạt
cho cùng một hiệu ứng năng lực số. Tuy nhiên, vì cơ chế này phụ thuộc
vào việc hạ tầng số có định hướng xuyên biên giới hay không, sức mạnh
của hiệu ứng uốn được dự báo là phụ thuộc bối cảnh thể chế chứ không phổ
quát, một dự đoán có điều kiện sẽ được kiểm chứng trong khung CDCM.

\hypertarget{giux1ea3-thuyux1ebft-h4-ux111iux1ec1u-tiux1ebft-bux1edfi-ux111ux1eb7c-ux111iux1ec3m-nhuxe0-quux1ea3n-trux1ecb-cux1ea5p-cao}{%
\subsection{2.5.5 Giả thuyết H4: Điều tiết bởi đặc điểm nhà quản trị cấp
cao}\label{giux1ea3-thuyux1ebft-h4-ux111iux1ec1u-tiux1ebft-bux1edfi-ux111ux1eb7c-ux111iux1ec3m-nhuxe0-quux1ea3n-trux1ecb-cux1ea5p-cao}}

Dựa trên lý thuyết thượng tầng quản trị (Hambrick \& Mason, 1984;
Hambrick, 2007), quyết định chiến lược quốc tế hóa và khả năng chuyển
hóa nó thành hiệu quả phụ thuộc vào đặc điểm nhận thức và kinh nghiệm
của nhà quản trị. Giả thuyết H4 được phát biểu như sau. Kinh nghiệm quản
trị và giới tính nữ của nhà quản trị cấp cao làm mạnh hơn tác động tích
cực của quốc tế hóa lên hiệu quả hoạt động, thông qua cơ chế nâng cao
nhận thức chiến lược quốc tế và quản trị rủi ro thận trọng hơn trong quá
trình mở rộng quốc tế. Cơ sở lý thuyết bao gồm Hambrick và Mason (1984),
Hsu et al.~(2013), và Post và Byron (2015); bằng chứng thực nghiệm ở cấp
doanh nghiệp cho thấy kinh nghiệm và giới tính của nhà quản trị có thể
điều tiết quan hệ quốc tế hóa và hiệu quả (Phan, Do, \& Phan, 2020b;
Phan et al., 2021; Do \& Phan, 2025).

Cơ chế đứng sau H4 là vai trò của trường nhận thức quản trị: kinh nghiệm
và sự đa dạng trong đội ngũ lãnh đạo cấp cao định hình cách doanh nghiệp
diễn giải rủi ro và cơ hội quốc tế, qua đó điều tiết chất lượng của các
quyết định quốc tế hóa thay vì tác động trực tiếp lên mức hiệu quả
(Nielsen, 2010). Cần lưu ý rằng dự đoán về dấu của hiệu ứng giới tính
phụ thuộc vào thước đo hiệu quả và đặc tả mô hình: khi đo bằng lợi nhuận
và mô hình hóa như điều tiết độ dốc, hiệu ứng dự kiến là dương; khi đo
bằng năng suất và ước lượng như hiệu ứng mức ở mẫu gộp đa quốc gia, hệ
số có thể âm do chọn lọc cấu thành. Sự phụ thuộc thước đo và bối cảnh
này được xử lý tường minh trong phần tổng hợp bằng chứng ở Chương 4 và
Chương 5, nên H4 được khung lại như một giả thuyết về điều tiết chất
lượng quyết định chứ không phải về một hiệu ứng mức phổ quát.

\hypertarget{giux1ea3-thuyux1ebft-h5-ux111iux1ec1u-tiux1ebft-bux1edfi-bux1ed1i-cux1ea3nh-thux1ec3-chux1ebf-icrv}{%
\subsection{2.5.6 Giả thuyết H5: Điều tiết bởi bối cảnh thể chế
(ICRV)}\label{giux1ea3-thuyux1ebft-h5-ux111iux1ec1u-tiux1ebft-bux1edfi-bux1ed1i-cux1ea3nh-thux1ec3-chux1ebf-icrv}}

Trên cơ sở lý thuyết thể chế (North, 1990; Khanna \& Palepu, 2010), chất
lượng thể chế tạo ra chi phí giao dịch và rủi ro khác nhau, do đó điều
tiết đáng kể mối quan hệ quốc tế hóa và hiệu quả. Giả thuyết H5 được
phát biểu như sau. Chất lượng thể chế theo khung ICRV điều tiết tác động
của các rào cản thể chế lên hiệu quả hoạt động trong quá trình quốc tế
hóa, với chỉ số chấp nhận số (DAI) đóng vai trò lá chắn số giúp bù đắp
một phần khoảng trống thể chế; cường độ điều tiết này giảm dần theo
gradient ICRV từ các nền kinh tế biên thể chế yếu (Nhóm V và VI, điều
tiết mạnh nhất) đến các nền kinh tế tiên tiến (Nhóm I, điều tiết yếu
nhất). Cơ sở lý thuyết bao gồm North (1990), Khanna và Palepu (2010), và
Marano et al.~(2016).

Cơ chế đứng sau H5 là sự thay thế giữa thể chế và năng lực: nơi thể chế
nền yếu, mỗi đơn vị năng lực doanh nghiệp, đặc biệt là năng lực số, tạo
ra khác biệt cạnh tranh lớn hơn vì nó bù đắp cho hạ tầng thể chế còn
thiếu; nơi thể chế đã mạnh, biên đóng góp của cùng năng lực đó nhỏ hơn
vì môi trường đã cung cấp phần lớn sự bảo đảm cần thiết cho giao dịch
xuyên biên giới. Logic thay thế này hàm ý một dự đoán gradient có hướng:
vai trò điều tiết của thể chế, và của lá chắn số gắn với nó, mạnh nhất ở
biên dưới của gradient ICRV và yếu dần khi tiến lên biên trên. Đây chính
là dự đoán định hướng mà luận án kiểm định bằng cả phân tích tổng hợp
theo nhóm thể chế và ước lượng điểm uốn trong từng chế độ ICRV ở Chương
4.

Một phản biện hiển nhiên đối với logic thay thế này là tính bổ trợ hạ
tầng: chính những nền kinh tế có thể chế yếu nhất cũng thường có hạ tầng
số kém phát triển nhất, nên có thể lập luận rằng lá chắn số phát huy tác
dụng yếu nhất chứ không phải mạnh nhất ở biên dưới của gradient. Luận án
giải quyết căng thẳng giữa thay thế và bổ trợ này bằng cách lưu ý rằng
thành phần DAI được vận hành hóa ở tầng nền tảng nhất (sự hiện diện web
và giao dịch số cơ bản), tức loại công cụ ít đòi hỏi hạ tầng dày nhất và
vẫn khả thi trên hạ tầng mỏng; ở tầng này, lợi ích chủ yếu đến từ việc
hạ chi phí tìm kiếm và thông tin trong trao đổi xuyên biên giới, vốn
chính là những chi phí mà khoảng trống thể chế khuếch đại. Vì vậy cơ chế
thay thế được dự báo trội hơn cơ chế bổ trợ trong dải năng lực số nền
tảng mà luận án đo lường; song chính sự phụ thuộc tầng đo lường này hàm
ý rằng hiệu ứng lá chắn có thể đảo chiều ở các tầng số đòi hỏi hạ tầng
cao hơn, một ranh giới phạm vi áp dụng được thừa nhận tường minh và sẽ
được nối lại trong phần hạn chế đo lường DAI ở Chương 5.

\hypertarget{giux1ea3-thuyux1ebft-h6-dux1ecb-biux1ec7t-thux1eddi-gian-trong-quan-hux1ec7-quux1ed1c-tux1ebf-huxf3a-vuxe0-hiux1ec7u-quux1ea3}{%
\subsection{2.5.7 Giả thuyết H6: Dị biệt thời gian trong quan hệ quốc tế
hóa và hiệu
quả}\label{giux1ea3-thuyux1ebft-h6-dux1ecb-biux1ec7t-thux1eddi-gian-trong-quan-hux1ec7-quux1ed1c-tux1ebf-huxf3a-vuxe0-hiux1ec7u-quux1ea3}}

Dựa trên quan sát về tốc độ chuyển đổi số và thay đổi thể chế rất nhanh
ở một số nền kinh tế châu Á, đặc biệt Trung Quốc, quan hệ quốc tế hóa và
hiệu quả không nhất thiết ổn định theo thời gian. Giả thuyết H6 được
phát biểu như sau. Hình dạng phi tuyến của quan hệ, cụ thể là vị trí của
điểm uốn và độ cong của đường quan hệ, thay đổi theo thời gian trong các
nền kinh tế chuyển đổi nhanh, phản ánh sự phát triển của năng lực doanh
nghiệp và cải thiện điều kiện thể chế qua các giai đoạn. Cơ sở lý thuyết
bao gồm Wu et al.~(2022), Xiao et al.~(2013), và Banalieva và Dhanaraj
(2019).

Việc kiểm định H6 đòi hỏi phân biệt hai khả năng đối lập về mặt lý
thuyết: dịch chuyển ngưỡng do trưởng thành năng lực và thể chế (một quá
trình tiệm tiến, có trật tự) so với tan rã ngưỡng do biến động thể chế
tích lũy (một quá trình phá vỡ cấu trúc). Bằng chứng từ hai bối cảnh
tương phản trong luận án phục vụ chính sự phân biệt này. Tại Trung Quốc,
nghiên cứu thành phần P5 cho thấy điểm uốn ổn định tương đối giữa 2012
và 2024 (kiểm định Paternoster \(z = +0{,}82\), \(p = {,}412\), không
bác bỏ tính bằng nhau), nhất quán với một quá trình trưởng thành có trật
tự. Ngược lại, bằng chứng Ấn Độ (P9) cho thấy ngưỡng có thể tan rã dưới
một chuỗi cú sốc thể chế dồn dập. Đặt cạnh nhau, hai trường hợp này tinh
chỉnh H6 thành một mệnh đề có điều kiện: dị biệt thời gian của quan hệ
là hàm của tốc độ và biên độ thay đổi thể chế tích lũy, chứ không phải
một xu hướng đơn điệu theo thời gian.

\hypertarget{tux1ed5ng-hux1ee3p-hux1ec7-giux1ea3-thuyux1ebft}{%
\subsection{2.5.8 Tổng hợp hệ giả
thuyết}\label{tux1ed5ng-hux1ee3p-hux1ec7-giux1ea3-thuyux1ebft}}

Bảy giả thuyết nêu trên được tổng hợp trong Bảng 2.2, làm rõ với mỗi giả
thuyết bốn yếu tố: phát biểu tóm tắt, dấu kỳ vọng của hệ số, cơ sở lý
thuyết nền và biến đo lường cùng kiểm định tương ứng. Cách trình bày này
bảo đảm mỗi giả thuyết đều được suy ra từ lý thuyết (không phải rút ra
từ quan sát thực nghiệm) và đều gắn với một thao tác kiểm định cụ thể
trong Chương 3 và Chương 4. Ba giả thuyết điều tiết H2, H3, H5 cùng vận
hành trên một đường quan hệ chính (H1), tương ứng ba tầng của khung
CDCM; H1b là điều kiện biên của H1 tại nhóm thể chế yếu nhất; H4 và H6
bổ sung tầng cá nhân và chiều thời gian.

\textbf{Bảng 2.2.} \emph{Tổng hợp hệ giả thuyết nghiên cứu: phát biểu,
dấu kỳ vọng, cơ sở lý thuyết và biến kiểm định.}

\begin{longtable}[]{@{}
  >{\raggedright\arraybackslash}p{(\columnwidth - 8\tabcolsep) * \real{0.2000}}
  >{\raggedright\arraybackslash}p{(\columnwidth - 8\tabcolsep) * \real{0.2000}}
  >{\raggedright\arraybackslash}p{(\columnwidth - 8\tabcolsep) * \real{0.2000}}
  >{\raggedright\arraybackslash}p{(\columnwidth - 8\tabcolsep) * \real{0.2000}}
  >{\raggedright\arraybackslash}p{(\columnwidth - 8\tabcolsep) * \real{0.2000}}@{}}
\toprule\noalign{}
\begin{minipage}[b]{\linewidth}\raggedright
Giả thuyết
\end{minipage} & \begin{minipage}[b]{\linewidth}\raggedright
Phát biểu tóm tắt
\end{minipage} & \begin{minipage}[b]{\linewidth}\raggedright
Dấu kỳ vọng
\end{minipage} & \begin{minipage}[b]{\linewidth}\raggedright
Cơ sở lý thuyết
\end{minipage} & \begin{minipage}[b]{\linewidth}\raggedright
Biến / kiểm định
\end{minipage} \\
\midrule\noalign{}
\endhead
\bottomrule\noalign{}
\endlastfoot
\textbf{H1} & Quan hệ quốc tế hóa và hiệu quả có dạng phi tuyến chữ U
ngược & \(\beta_1>0\); \(\beta_2<0\) & Mô hình Uppsala; Lu và Beamish
(2004) & FSTS, FSTS²; kiểm định Lind--Mehlum \\
\textbf{H1b} & Tại nhóm SIDS (ICRV VI), quan hệ âm đơn điệu, không điểm
uốn (gánh nặng quốc tế hóa bắt buộc, FIP) & \(\beta(\text{FSTS})<0\) &
Lý thuyết thể chế (North, 1990); chi phí giao dịch (Williamson, 1985);
Lall (2004) & FSTS trên mẫu SIDS; kiểm tra không tồn tại điểm uốn \\
\textbf{H2} & Năng lực công nghệ (TCI) điều tiết dương, nâng mặt bằng
hiệu quả và có thể dịch chuyển điểm uốn & \(+\) & Lý thuyết dựa trên
nguồn lực (Barney, 1991); năng lực động (Teece et al., 1997) & TCI; TCI
× FSTS; TCI × FSTS² \\
\textbf{H3} & Chấp nhận số (DAI) điều tiết, uốn đường cong rõ hơn ở
cường độ xuất khẩu cao (giảm chi phí phối hợp) & \(+\) (đổi độ dốc) &
Verhoef et al.~(2021); Stallkamp và Schotter (2021); Banalieva và
Dhanaraj (2019) & DAI; DAI × FSTS; DAI × FSTS² \\
\textbf{H4} & Kinh nghiệm và giới tính nữ của nhà quản trị cấp cao điều
tiết dương quan hệ quốc tế hóa và hiệu quả & \(+\) & Lý thuyết thượng
tầng quản trị (Hambrick và Mason, 1984); Post và Byron (2015) & Kinh
nghiệm, giới tính × FSTS \\
\textbf{H5} & Chất lượng thể chế (gradient ICRV) điều tiết; tác động
mạnh nhất ở nhóm thể chế yếu, giảm dần đến nhóm tiên tiến & \(+/-\) theo
gradient & Lý thuyết thể chế (North, 1990; Khanna và Palepu, 2010);
Marano et al.~(2016) & ICRV × FSTS; ICRV × FSTS² \\
\textbf{H6} & Dạng hàm (vị trí điểm uốn, độ cong) dịch chuyển theo thời
gian ở nền kinh tế chuyển đổi nhanh & Thay đổi theo thời gian & Wu et
al.~(2022); Xiao et al.~(2013); Banalieva và Dhanaraj (2019) & So sánh
đa đợt khảo sát; kiểm định Paternoster \\
\end{longtable}

\emph{Ghi chú: FSTS = tỷ lệ doanh thu nước ngoài trên tổng doanh thu
(cường độ xuất khẩu); TCI = chỉ số năng lực công nghệ; DAI = chỉ số chấp
nhận số; ICRV = khung phân loại chế độ bối cảnh thể chế sáu nhóm; FIP =
gánh nặng quốc tế hóa bắt buộc. Dấu kỳ vọng tham chiếu hệ số của số hạng
tương tác tương ứng. Nguồn: tác giả tổng hợp.}

\begin{center}\rule{0.5\linewidth}{0.5pt}\end{center}

\hypertarget{uxe1nh-xux1ea1-giux1ea3-thuyux1ebft-vuxe0-cuxe1c-nghiuxean-cux1ee9u-thuxe0nh-phux1ea7n}{%
\subsection{2.5.9 Ánh xạ giả thuyết và các nghiên cứu thành
phần}\label{uxe1nh-xux1ea1-giux1ea3-thuyux1ebft-vuxe0-cuxe1c-nghiuxean-cux1ee9u-thuxe0nh-phux1ea7n}}

Luận án được phát triển gắn với một chuỗi nghiên cứu thành phần (ký hiệu
P1--P10 và một chương sách), mỗi nghiên cứu kiểm định một phần của hệ
giả thuyết H1--H6 và điều kiện biên H1b trong một bối cảnh thể chế cụ
thể. Bảng 2.1 ánh xạ các nghiên cứu thành phần với mục trình bày kết quả
và các giả thuyết tương ứng.

\textbf{Bảng 2.1.} \emph{Ánh xạ nghiên cứu thành phần, mục trình bày,
giả thuyết.}

\begin{longtable}[]{@{}
  >{\raggedright\arraybackslash}p{(\columnwidth - 6\tabcolsep) * \real{0.2500}}
  >{\raggedright\arraybackslash}p{(\columnwidth - 6\tabcolsep) * \real{0.2500}}
  >{\raggedright\arraybackslash}p{(\columnwidth - 6\tabcolsep) * \real{0.2500}}
  >{\raggedright\arraybackslash}p{(\columnwidth - 6\tabcolsep) * \real{0.2500}}@{}}
\toprule\noalign{}
\begin{minipage}[b]{\linewidth}\raggedright
Nghiên cứu thành phần
\end{minipage} & \begin{minipage}[b]{\linewidth}\raggedright
Bối cảnh (nhóm ICRV)
\end{minipage} & \begin{minipage}[b]{\linewidth}\raggedright
Trình bày tại
\end{minipage} & \begin{minipage}[b]{\linewidth}\raggedright
Giả thuyết kiểm định
\end{minipage} \\
\midrule\noalign{}
\endhead
\bottomrule\noalign{}
\endlastfoot
P1: Cơ sở đa quốc gia (17 nền kinh tế) & Châu Á (đa nhóm) & Mục 2.3.3;
Mục 4.1.7 & Nền tảng cho H1, H2, H5 \\
P2: Doanh nghiệp nhỏ và vừa Trung Quốc & Trung bình cao (Nhóm III) & Mục
2.3.3; Mục 4.4.8 & Nền tảng cho H1 \\
P3: Việt Nam & Chuyển đổi thu nhập trung bình thấp (Nhóm IV) & Mục 4.5 &
H1, H2, H3 \\
P4: Singapore & Tiên tiến đổi mới (Nhóm I) & Mục 4.3 & H1, H3 \\
P5: Trung Quốc 2012--2024 & Trung bình cao (Nhóm III) & Mục 4.4 & H1,
H6 \\
P6: Phân tích tổng hợp ba tầng & Toàn bộ tài liệu châu Á & Mục 4.2 & H1,
H5 \\
P7: Toàn mẫu đa quốc gia (50 nền kinh tế) & Toàn bộ gradient ICRV & Mục
4.6 & H1, H2, H3, H4, H5 (và tương tác tổng hợp ba chiều) \\
P8: Đảo nhỏ Thái Bình Dương (mẫu chính 7 nước, N=959; kiểm định độ vững
mở rộng 9 nước Nhóm VI) & SIDS (Nhóm VI) & Mục 4.7 & H1b (FIP) \\
P9: Ấn Độ 2014--2025 & Chuyển đổi thu nhập trung bình thấp (Nhóm IV) &
Mục 4.8 & H1c (độ bền ngưỡng dưới biến động thể chế) \\
P10: Nhật Bản 2025 & Tiên tiến đổi mới (Nhóm I) & Mục 4.1.5; Mục 4.9 &
H1 (dự đoán biên trên: gần tuyến tính) \\
Chương sách: Ấn Độ, quản trị cấp cao & Chuyển đổi thu nhập trung bình
thấp (Nhóm IV) & Mục 4.8.4 & H4 (điều tiết bởi đặc trưng quản trị) \\
\end{longtable}

\emph{Ghi chú về ký hiệu: Hệ giả thuyết của luận án gồm sáu giả thuyết
trọng tâm H1--H6 và một giả thuyết điều kiện biên H1b. Trong các mô hình
thực nghiệm chi tiết ở Chương 3, một số nghiên cứu thành phần sử dụng
các giả thuyết con (sub-hypotheses) lồng dưới hệ giả thuyết chính của
luận án, cụ thể: H2b (kiểm định tính ổn định xuyên sóng của năng lực
điều tiết, thuộc P5, lồng dưới H2/H6); H4a--H4b (kiểm định điều tiết độ
cong theo năng lực, thuộc P5, lồng dưới H2); tương tác tổng hợp ba chiều
của P7 (lồng dưới H3 và H5); E1a--E1b (kiểm định định hướng gradient
ICRV, lồng dưới H5); và H1c (kiểm định tan biến ngưỡng chữ U ngược theo
thời gian tại Ấn Độ, thuộc P9, lồng dưới H1 và H6). Các ký hiệu con này
phục vụ độ chính xác kỹ thuật của từng mô hình và không làm thay đổi hệ
giả thuyết trọng tâm H1--H6/H1b của luận án.}

\hypertarget{tuxf3m-tux1eaft-chux1b0ux1a1ng}{%
\section{2.6 Tóm tắt chương}\label{tuxf3m-tux1eaft-chux1b0ux1a1ng}}

Chương 2 đã trình bày hệ thống nền tảng lý thuyết và thực nghiệm cho
luận án theo một trục logic nhất quán: từ khái niệm đến lý thuyết, từ
bằng chứng thực nghiệm đến khoảng trống, và từ khoảng trống đến khung
nghiên cứu và giả thuyết.

Phần đầu chương xác lập bốn khái niệm cốt lõi: quốc tế hóa doanh nghiệp
(đo bằng FSTS), hiệu quả hoạt động (đo chủ yếu bằng năng suất lao động),
năng lực công nghệ (TCI) và chỉ số chấp nhận số (DAI) như hai cấu trúc
độc lập, và bối cảnh thể chế được vận hành hóa bằng khung ICRV sáu nhóm.
Phần lý thuyết nền trình bày và phân tích bốn lý thuyết kinh điển, gồm
mô hình Uppsala, lý thuyết dựa trên nguồn lực, lý thuyết thể chế và lý
thuyết thượng tầng quản trị, cùng với lớp mở rộng lăng kính năng lực số,
xác định vai trò của từng lý thuyết trong khung giải thích tổng thể của
luận án. Điểm nhấn lý thuyết là lập luận rằng năm lớp này bổ sung cho
nhau theo cấp độ phân tích, từ chi phí học hỏi, khác biệt năng lực, bối
cảnh thể chế, đặc điểm người ra quyết định, đến logic số, và do đó cần
được tích hợp thay vì lựa chọn loại trừ.

Tổng quan thực nghiệm cho thấy tài liệu về quan hệ quốc tế hóa và hiệu
quả đã phát triển qua ba giai đoạn, từ tuyến tính dương đến phi tuyến
đến tập trung vào các biến điều tiết; các phân tích tổng hợp nhất quán
cho thấy tác động gộp là dương nhưng nhỏ với mức dị biệt cao (\(I^2\)
trên 80\% ở các tổng hợp trước, \(62{,}4\%\) ở phân tích tổng hợp cập
nhật P6), đòi hỏi cách tiếp cận lấy biến điều tiết làm trung tâm; và bối
cảnh châu Á, đặc biệt là các nền kinh tế biên thể chế yếu và các nền đảo
nhỏ, đang thiếu đại diện nghiêm trọng trong tài liệu. Phần này cũng ghi
nhận làn sóng nhân rộng quy mô lớn gần đây vốn đặt câu hỏi về tính phổ
quát của dạng hàm, qua đó củng cố luận điểm rằng chính dị biệt bối cảnh,
chứ không phải một quy luật hàm phổ quát, mới chi phối kết quả.

Ba khoảng trống được xác định: thiếu khung tích hợp đa tầng, sự đồng
nhất không hợp lý giữa TCI và DAI, và thiếu bằng chứng đa quốc gia quy
mô lớn có cấu trúc cho châu Á. Để lấp đầy ba khoảng trống này, luận án
đề xuất khung CDCM và phát triển sáu giả thuyết trọng tâm (H1 đến H6)
cùng điều kiện biên H1b, bao phủ phi tuyến cơ bản, vai trò của TCI và
DAI, vai trò của đặc điểm nhà quản trị, vai trò của thể chế ICRV, và dị
biệt thời gian. Mỗi giả thuyết được neo vào một hoặc nhiều nghiên cứu
thành phần và một bối cảnh thể chế cụ thể, như Bảng 2.1 trình bày, để
bảo đảm rằng khung lý thuyết được kiểm định trên toàn bộ gradient thể
chế của khu vực chứ không chỉ ở một bối cảnh đơn lẻ. Chương 3 sẽ trình
bày thiết kế phương pháp và chiến lược kiểm định hệ giả thuyết này.

\newpage

\hypertarget{chux1b0ux1a1ng-3-phux1b0ux1a1ng-phuxe1p-nghiuxean-cux1ee9u}{%
\chapter{CHƯƠNG 3: PHƯƠNG PHÁP NGHIÊN
CỨU}\label{chux1b0ux1a1ng-3-phux1b0ux1a1ng-phuxe1p-nghiuxean-cux1ee9u}}

\begin{center}\rule{0.5\linewidth}{0.5pt}\end{center}

\hypertarget{thiux1ebft-kux1ebf-nghiuxean-cux1ee9u-tux1ed5ng-thux1ec3}{%
\section{3.1 Thiết kế nghiên cứu tổng
thể}\label{thiux1ebft-kux1ebf-nghiuxean-cux1ee9u-tux1ed5ng-thux1ec3}}

\hypertarget{loux1ea1i-thiux1ebft-kux1ebf}{%
\subsection{3.1.1 Loại thiết kế}\label{loux1ea1i-thiux1ebft-kux1ebf}}

Luận án sử dụng \textbf{thiết kế hỗn hợp tổng hợp--thực nghiệm} gồm hai
giai đoạn bổ sung cho nhau. Giai đoạn thứ nhất là phân tích tổng hợp tài
liệu để trả lời RQ1, còn giai đoạn thứ hai là phân tích thực nghiệm
doanh nghiệp đa quốc gia để trả lời RQ2--RQ4. Hai giai đoạn đối thoại
hai chiều: phân tích tổng hợp cung cấp cơ sở và gợi ý biến điều tiết cần
kiểm định, còn thực nghiệm cung cấp bằng chứng bổ sung cho phân tích
tổng hợp về bối cảnh châu Á và Pacific.

\textbf{Cơ sở kế thừa}: Thiết kế hỗn hợp kết hợp phân tích tổng hợp và
dữ liệu sơ cấp đã được áp dụng trong quản trị chiến lược và IB
(Borenstein et al., 2009; Hunter \& Schmidt, 2004; Combs et al., 2011).

\textbf{Đóng góp mới}: Áp dụng thiết kế hỗn hợp cho một luận án duy nhất
về quốc tế hóa--hiệu quả doanh nghiệp trong bối cảnh châu Á và Thái Bình
Dương với \textbf{50 nền kinh tế} là thiết kế hiếm gặp; đa số các luận
án IB chỉ làm một trong hai giai đoạn. Quy mô 50 nền kinh tế/88.869
doanh nghiệp là phạm vi địa lý lớn nhất từng có cho khu vực, mở rộng từ
mẫu gộp 17 nước trong nghiên cứu thành phần P1.

\hypertarget{quy-truxecnh-nghiuxean-cux1ee9u}{%
\subsection{3.1.2 Quy trình nghiên
cứu}\label{quy-truxecnh-nghiuxean-cux1ee9u}}

Quy trình đi theo sáu bước: (i) xác định khoảng trống nghiên cứu từ tổng
quan và phân tích tổng hợp trước đây; (ii) xây dựng khung lý thuyết tích
hợp đa tầng và hệ giả thuyết H1--H6; (iii) tiến hành phân tích tổng hợp
cập nhật 1977--2026; (iv) thu thập và chuẩn hóa dữ liệu WBES cho
\textbf{50 nền kinh tế châu Á và Thái Bình Dương} sau khi hòa hợp 3 thế
hệ lược đồ (PICS3, Standardized, BREADY/BEE); (v) ước lượng mô hình thực
nghiệm theo bối cảnh quốc gia và đa quốc gia; (vi) kiểm định độ vững và
tích hợp kết quả để đóng góp lý thuyết. Toàn bộ quy trình và quan hệ bổ
sung hai chiều giữa hai giai đoạn được trình bày ở Hình 3.1.

\includegraphics{figures/fig_3_1_research_process.png}

\textbf{Hình 3.1.} \emph{Quy trình nghiên cứu của luận án.} Mũi tên liền
thể hiện trình tự thực hiện; mũi tên hai chiều thể hiện quan hệ bổ sung
giữa hai giai đoạn (phân tích tổng hợp gợi ý biến điều tiết cần kiểm
định; phân tích thực nghiệm cung cấp bằng chứng châu Á và Thái Bình
Dương cho phân tích tổng hợp). \emph{Nguồn: tác giả xây dựng.}

\hypertarget{giai-ux111oux1ea1n-1-phuxe2n-tuxedch-tux1ed5ng-hux1ee3p-19772026}{%
\section{3.2 Giai đoạn 1: Phân tích tổng hợp
1977--2026}\label{giai-ux111oux1ea1n-1-phuxe2n-tuxedch-tux1ed5ng-hux1ee3p-19772026}}

\hypertarget{cux1a1-sux1edf-kux1ebf-thux1eeba}{%
\subsection{3.2.1 Cơ sở kế
thừa}\label{cux1a1-sux1edf-kux1ebf-thux1eeba}}

phân tích tổng hợp trong IB đã được thiết lập vững chắc qua nhiều thập
kỷ (Hunter \& Schmidt, 2004; Borenstein et al., 2009). Phương pháp chuẩn
bao gồm: xây dựng tiêu chí lựa chọn theo PRISMA (Page et al., 2021),
chuyển dạng hiệu ứng bằng Fisher's \(z\), tính hiệu ứng gộp size theo
hiệu ứng ngẫu nhiên, kiểm định dị biệt bằng \(Q\)-test và \(I^2\), và
đánh giá lệch lạc công bố bằng Egger và Begg-Mazumdar (Egger et al.,
1997; Begg \& Mazumdar, 1994). Các các phân tích tổng hợp trước đây về
quốc tế hóa và hiệu quả (Bausch \& Krist, 2007; Kirca et al., 2012;
Marano et al., 2016; Wu et al., 2022; Arte \& Larimo, 2022) cung cấp quy
trình mẫu mà luận án kế thừa và cập nhật.

\hypertarget{ux111uxf3ng-guxf3p-mux1edbiux111iux1ec1u-chux1ec9nh}{%
\subsection{3.2.2 Đóng góp mới/điều
chỉnh}\label{ux111uxf3ng-guxf3p-mux1edbiux111iux1ec1u-chux1ec9nh}}

\begin{itemize}
\tightlist
\item
  \textbf{Mở rộng coverage time}: từ 1977--2022 (bài ICBEF 2025) sang
  \textbf{1977--2026}, bổ sung các nghiên cứu 2023--2026 liên quan đến
  chuyển đổi số và doanh nghiệp đa quốc gia thị trường mới nổi.
\item
  \textbf{Mở rộng biến điều tiết coverage}: thêm cụm biến điều tiết về
  năng lực số và 6 chế độ con ICRV, chưa được xem xét đầy đủ trong các
  các phân tích tổng hợp trước.
\item
  \textbf{Tập tổng hợp}: K=288 cỡ ảnh hưởng / k=238 nghiên cứu, được xây
  dựng theo chiến lược tìm kiếm hệ thống nhưng có giới hạn (neo theo
  trích dẫn: quét trích dẫn xuôi và ngược từ 5 phân tích tổng hợp mỏ
  neo, bổ sung bằng cơ sở dữ liệu), báo cáo theo biến thể PRISMA 2020
  ``studies identified via other methods''; mở rộng so với bản nền ICBEF
  2025 (K=200 / k=113).
\item
  \textbf{Subgroup theo châu Á và Pacific}: tách riêng nhóm nghiên cứu
  về doanh nghiệp châu Á và Pacific để cung cấp cơ sở trực tiếp cho giai
  đoạn thực nghiệm, đối chiếu với mẫu gộp 17 nước trong nghiên cứu thành
  phần P1.
\end{itemize}

\hypertarget{quy-truxecnh-prisma}{%
\subsection{3.2.3 Quy trình PRISMA}\label{quy-truxecnh-prisma}}

Bốn bước Nhận diện, Sàng lọc, Đủ điều kiện và Đưa vào theo hướng dẫn
PRISMA 2020 (Page et al., 2021). tiêu chí lựa chọn: nghiên cứu thực
nghiệm ở cấp doanh nghiệp, có đo lường quốc tế hóa và hiệu quả doanh
nghiệp, có đủ thống kê để tính cỡ ảnh hưởng (Pearson \(r\), \(\beta\),
\(t\)-stat).

\hypertarget{giai-ux111oux1ea1n-2-phuxe2n-tuxedch-thux1ef1c-nghiux1ec7m-ux111a-quux1ed1c-gia}{%
\section{3.3 Giai đoạn 2: Phân tích thực nghiệm đa quốc
gia}\label{giai-ux111oux1ea1n-2-phuxe2n-tuxedch-thux1ef1c-nghiux1ec7m-ux111a-quux1ed1c-gia}}

\hypertarget{nguux1ed3n-dux1eef-liux1ec7u}{%
\subsection{3.3.1 Nguồn dữ liệu}\label{nguux1ed3n-dux1eef-liux1ec7u}}

Luận án sử dụng World Bank Enterprise Surveys cho \textbf{50 nền kinh tế
châu Á và Thái Bình Dương} trong giai đoạn \textbf{2006--2026} (88.869
doanh nghiệp trong khung phân tích, 103 cặp nền kinh tế và năm; bao gồm
Nhật Bản 2025; pool phân loại ICRV 96.415 doanh nghiệp/52 nhãn nền)
(World Bank, n.d., 2019, 2023, 2026). WBES là bộ dữ liệu doanh nghiệp
chuẩn hóa, phù hợp cho nghiên cứu so sánh đa quốc gia về cấp doanh
nghiệp kết quả (Aterido et al., 2011). Quy trình hợp nhất và hài hòa hóa
chi tiết các bộ khảo sát quốc gia--năm thành bộ dữ liệu phân tích thống
nhất (sáu bước S1--S6, xử lý tính so sánh tiền tệ, chiến lược lát cắt
ngang gộp với hiệu ứng cố định hai chiều, cùng dòng dữ liệu PRISMA và mã
tái lập) được trình bày đầy đủ ở \textbf{Phụ lục A}.

\textbf{Cơ sở kế thừa}: WBES đã được sử dụng rộng rãi trong nghiên cứu
IB về thị trường mới nổi (Ayyagari et al., 2011; Cuervo-Cazurra et al.,
2018; Chen \& Meng, 2022), đồng thời đã được sử dụng trong bài nghiên
cứu thực trạng châu Á mới nổi với 17 nền kinh tế và
\textasciitilde40.633 quan sát doanh nghiệp (nghiên cứu thành phần P1).

\textbf{Đóng góp mới}: mở rộng quy mô từ 17 sang \textbf{50 nền kinh tế
châu Á và Thái Bình Dương}, với 6 chế độ con ICRV (Advanced
innovation-driven, tiên tiến tài nguyên, trung bình cao, chuyển đổi thu
nhập trung bình thấp, đang nổi, SIDS). Quy mô này là lớn nhất cho khu
vực trong tài liệu quốc tế hóa--performance, gồm cả trường hợp biên đảo
nhỏ Thái Bình Dương (7 nước, n=959 (mẫu phân tích sau lọc missing)) cho
phép kiểm định forced quốc tế hóa penalty (nghiên cứu thành phần P8).

\begin{quote}
\textbf{Ghi chú về trọng số khảo sát (survey sampling weights):} WBES
cung cấp sampling weights để đại diện cho tổng thể doanh nghiệp trong
từng quốc gia. Tuy nhiên, luận án \emph{không áp dụng sampling weights}
trong ước lượng mô hình hồi quy vì ba lý do: (1) Mục tiêu của luận án là
kiểm định cơ chế lý thuyết (mechanism testing) chứ không phải ước lượng
mô tả đại diện tổng thể (population-representative descriptive
estimation), trong bối cảnh này, không áp dụng weights là thực hành
chuẩn trong tài liệu IB sử dụng WBES (Aterido et al., 2011;
Cuervo-Cazurra et al., 2018); (2) hiệu ứng cố định ở cấp cặp nền kinh tế
và năm đã kiểm soát đặc điểm cấu trúc của từng sóng khảo sát; (3) Việc
áp dụng weights có thể làm sai lệch ước lượng khi mẫu được gộp chung từ
nhiều quốc gia có quy mô tổng thể doanh nghiệp khác nhau nhiều bậc (ví
dụ: Trung Quốc \textasciitilde100 triệu so với Tonga \textasciitilde500
doanh nghiệp). Phân tích độ nhạy có áp dụng survey weights cho từng quốc
gia riêng lẻ (P3/P4/P5) không thay đổi hướng hoặc mức độ ý nghĩa của các
hệ số chính (xem Mục 3.5).
\end{quote}

\hypertarget{huxe0i-huxf2a-huxf3a-ba-thux1ebf-hux1ec7-lux1b0ux1ee3c-ux111ux1ed3-bux1ea3ng-hux1ecfi-wbes}{%
\subsection{3.3.1.1 Hài hòa hóa ba thế hệ lược đồ bảng hỏi
WBES}\label{huxe0i-huxf2a-huxf3a-ba-thux1ebf-hux1ec7-lux1b0ux1ee3c-ux111ux1ed3-bux1ea3ng-hux1ecfi-wbes}}

Bộ công cụ WBES đã trải qua ba thế hệ lược đồ với hệ thống mã biến và
phạm vi module khác nhau: thế hệ PICS3 đời đầu, thế hệ Standardized toàn
cầu áp dụng từ năm 2006, và thế hệ BREADY/BEE gần đây. Sự thay đổi này
tạo ra thách thức then chốt cho hợp nhất đa quốc gia, vì cùng một khái
niệm có thể được hỏi bằng trường khác nhau hoặc thiếu hẳn ở một số đợt.
Luận án giải quyết bằng một lược đồ biến chung ánh xạ mọi trường dị biệt
về các module cốt lõi bất biến của bộ công cụ toàn cầu, với tiêu chí vận
hành để một đợt được nhận vào khung so sánh là sự hiện diện của các biến
lõi đã hài hòa, cụ thể là doanh thu (d2), lao động (l1) và xuất khẩu
trực tiếp cùng gián tiếp (d3b và d3c). Đợt khảo sát đáp ứng năm khảo sát
từ 2006 nhưng chạy trên bộ công cụ khu vực cũ không chứa các biến lõi
này bị loại do bất tương thích bộ công cụ chứ không phải do thiếu dữ
liệu.

Quy trình hợp nhất tuân theo nguyên tắc minh bạch kiểu PRISMA cho dữ
liệu thứ cấp, gồm sáu bước tuần tự: lọc loại khảo sát, khử trùng lặp một
khảo sát cho mỗi cặp nền kinh tế và năm, hài hòa hóa biến, mã hóa mã
không phản hồi của WBES thành giá trị khuyết, phân tầng chế độ thể chế,
đến tạo mẫu phân tích bằng xóa theo danh sách. Mỗi bước ghi rõ tiêu chí
và số quan sát còn lại nhằm cho phép tái dựng và kiểm chứng. Một hệ quả
của xóa theo danh sách là cỡ mẫu giảm dần khi thêm biến kiểm soát, rõ
nhất ở biến sở hữu nước ngoài vốn khuyết nhiều ở các nền nhỏ và sóng
sớm; luận án báo cáo minh bạch sự suy giảm này và kiểm tra tính ổn định
của điểm uốn qua các đặc tả để bảo đảm kết quả không bị chi phối bởi cơ
chế khuyết. Toàn bộ thao tác chi tiết của quy trình sáu bước, cùng dòng
dữ liệu kiểu PRISMA và mã tái lập, được trình bày đầy đủ ở Phụ lục A.

\hypertarget{ux111o-lux1b0ux1eddng-biux1ebfn-phux1ee5-thuux1ed9c-hiux1ec7u-quux1ea3-houx1ea1t-ux111ux1ed9ng}{%
\subsection{3.3.2 Đo lường biến phụ thuộc: Hiệu quả hoạt
động}\label{ux111o-lux1b0ux1eddng-biux1ebfn-phux1ee5-thuux1ed9c-hiux1ec7u-quux1ea3-houx1ea1t-ux111ux1ed9ng}}

\textbf{Biến chính}: năng suất lao động đo bằng
\(\ln(\text{annual sales} / \text{permanent full-time employees})\).

\textbf{Cơ sở kế thừa}: năng suất lao động được sử dụng rộng rãi trong
nghiên cứu WBES và trong tài liệu về firm productivity (Bloom et al.,
2012; Hsieh \& Klenow, 2009). Chỉ tiêu này đã được vận hành hóa thành
công trong bài thực trạng châu Á mới nổi (nghiên cứu thành phần P1).

\textbf{Lập luận chọn}: WBES không phải bộ financial statement đầy đủ,
nên năng suất lao động phù hợp hơn ROA/ROE/ROS trong bối cảnh đa quốc
gia. Ngoài ra, năng suất lao động ít bị ảnh hưởng bởi chuẩn mực kế toán
khác biệt giữa các nước hơn so với các chỉ tiêu kế toán (Combs et al.,
2005; Richard et al., 2009).

\textbf{Nhất quán đo lường xuyên các nghiên cứu thành phần và điều kiện
so sánh tiền tệ}: cùng một cấu trúc \(\ln(\text{LP})\) từ trường WBES
d2/l1 được dùng thống nhất ở P3, P4, P5 và P7, nhưng cách xử lý đơn vị
tiền tệ khác nhau \textbf{theo đúng bản chất thiết kế} của từng nghiên
cứu: (i) các nghiên cứu \textbf{đơn quốc gia} (P3 Việt Nam, P4
Singapore, P5 Trung Quốc, P9 Ấn Độ) đo doanh thu bằng \textbf{một nội tệ
duy nhất} nên mức \(\ln(\text{LP})\) so sánh được trực tiếp trong nội bộ
mẫu; khác biệt mặt bằng giá giữa các đợt khảo sát được hấp thụ bằng biến
giả đợt/hiệu ứng cố định năm (P5: \(D_{2024}\); P3 bổ sung hiệu chỉnh
PPP cho khoảng cách 14 năm 2009--2023), kèm winsorize 1\%/99\% (P4);
(ii) nghiên cứu \textbf{đa quốc gia} (P7) gộp doanh thu bằng
\textbf{nhiều nội tệ}, do đó mức \(\ln(\text{LP})\) thô không so sánh
được giữa các nước, tính so sánh được bảo đảm bằng chuẩn hóa z within
country--year và hiệu ứng cố định hai chiều, như trình bày chi tiết ở
Phụ lục A (Mục A.5); mọi so sánh \textbf{mức} xuyên quốc gia trong luận
án dùng tỷ suất không thứ nguyên (ROS) hoặc giá trị đã hiệu chỉnh PPP.
Như vậy tính so sánh xuyên quốc gia của thước đo là \textbf{có điều
kiện}, đạt được qua thiết kế ước lượng, không phải tự động ở mức thô. Về
phía biến quốc tế hóa, khung ước lượng đa quốc gia của luận án (P7, P8,
P9) và thống kê mô tả ở Chuyên đề 1 đo FSTS bằng \textbf{tổng xuất khẩu
trực tiếp và gián tiếp} (trường WBES d3b + d3c), định nghĩa tham gia
xuất khẩu chuẩn của bộ công cụ WBES; hai nghiên cứu đơn quốc gia (P3
Việt Nam, P5 Trung Quốc) dùng \textbf{riêng xuất khẩu trực tiếp} (d3c),
vì xuất khẩu gián tiếp ủy thác chức năng quốc tế hóa cho trung gian
thương mại với cơ chế học hỏi và yêu cầu nguồn lực khác về bản chất
(Hessels \& Terjesen, 2010; Peng \& Ilinitch, 1998), trong khi tỷ trọng
doanh thu xuất khẩu là thước đo cường độ quốc tế hóa chuẩn của tài liệu
về quan hệ quốc tế hóa và hiệu quả (Lu \& Beamish, 2001; Sullivan,
1994); lựa chọn thao tác hóa được ghi rõ trong phần dữ liệu của từng
nghiên cứu thành phần.

\textbf{Biến độ vững}: ROS, sales growth, employment growth,
\(\ln\)(annual sales).

\textbf{Đóng góp mới}: lập luận rõ ràng về hiệu quả doanh nghiệp là khái
niệm đa chiều và việc lựa chọn thước đo phải phù hợp với bối cảnh dữ
liệu (xem Chương 2).

\hypertarget{vux1eadn-huxe0nh-huxf3a-thux1b0ux1edbc-ux111o-nux103ng-suux1ea5t-vuxe0-quy-truxecnh-xux1eed-luxfd-phuxe2n-phux1ed1i}{%
\subsection{3.3.2.1 Vận hành hóa thước đo năng suất và quy trình xử lý
phân
phối}\label{vux1eadn-huxe0nh-huxf3a-thux1b0ux1edbc-ux111o-nux103ng-suux1ea5t-vuxe0-quy-truxecnh-xux1eed-luxfd-phuxe2n-phux1ed1i}}

Năng suất lao động được dựng từ hai trường khảo sát cốt lõi của bộ công
cụ WBES: tổng doanh thu năm tài chính gần nhất (trường d2) và số lao
động thường xuyên toàn thời gian cuối kỳ (trường l1). Biến phụ thuộc
nhận dạng logarit tự nhiên của thương số hai trường này, nghĩa là
\(\ln(\text{LP}) = \ln(d2 / l1)\), một phép biến đổi vừa nén đuôi phải
của phân phối doanh thu vừa cho phép diễn giải hệ số theo phần trăm. Do
doanh thu thô được báo cáo bằng nội tệ của từng nền kinh tế, mức năng
suất thô không so sánh trực tiếp giữa các nước, một mối đe dọa hiệu lực
được xử lý qua thiết kế ước lượng chứ không qua quy đổi tỷ giá danh
nghĩa.

Trước khi ước lượng, năng suất lao động trải qua hai bước xử lý phân
phối nhằm bảo đảm tính bền vững của ước lượng trước quan sát cực trị.
Bước thứ nhất là cắt đuôi (winsorize) tại ngưỡng phân vị 1 và 99 áp dụng
trong từng cặp nền kinh tế và năm, qua đó giới hạn ảnh hưởng của giá trị
ngoại lai mà không loại bỏ quan sát khỏi mẫu. Bước thứ hai, chỉ áp dụng
cho khung ước lượng đa quốc gia, là chuẩn hóa z trong từng cặp nền kinh
tế và năm, đưa mọi quan sát về cùng thang vị thế tương đối trong phân
phối nội bộ nước--năm. Phép chuẩn hóa này khử triệt để khác biệt đơn vị
tiền tệ và mặt bằng giá, đồng thời bảo toàn cấu trúc biến thiên bên
trong mỗi nước, nơi diễn ra việc nhận dạng quan hệ quốc tế hóa--hiệu
quả. Đối với các nghiên cứu đơn quốc gia, doanh thu được biểu thị bằng
một nội tệ duy nhất nên mức \(\ln(\text{LP})\) so sánh được trực tiếp
trong nội bộ mẫu; khác biệt mặt bằng giá giữa các đợt khảo sát được hấp
thụ bằng biến giả đợt hoặc hiệu ứng cố định năm, kèm cắt đuôi 1 và 99.

\hypertarget{ux111o-lux1b0ux1eddng-biux1ebfn-ux111ux1ed9c-lux1eadp-quux1ed1c-tux1ebf-huxf3a}{%
\subsection{3.3.3 Đo lường biến độc lập: Quốc tế
hóa}\label{ux111o-lux1b0ux1eddng-biux1ebfn-ux111ux1ed9c-lux1eadp-quux1ed1c-tux1ebf-huxf3a}}

\textbf{Biến chính}: cường độ xuất khẩu (FSTS, tỷ lệ doanh thu xuất khẩu
trên tổng doanh thu). Mô hình đưa vào đồng thời \(FSTS\) và \(FSTS^2\)
để kiểm định quan hệ phi tuyến dạng chữ U ngược.

\textbf{Cơ sở kế thừa}: FSTS là thước đo quốc tế hóa phổ biến nhất trong
tài liệu về quan hệ quốc tế hóa và hiệu quả (Sullivan, 1994; Hitt et
al., 2006). Việc bổ sung \(FSTS^2\) để kiểm định phi tuyến đã được áp
dụng trong Lu và Beamish (2004), Contractor et al.~(2003), và Hitt et
al.~(1997). Dạng bậc hai được nghiên cứu thành phần P2 xác nhận có ý
nghĩa thống kê ở Trung Quốc với điểm uốn khoảng 47,8\% FSTS, làm cơ sở
cho việc kiểm định mở rộng trong luận án.

\hypertarget{vux1eadn-huxe0nh-huxf3a-cux1b0ux1eddng-ux111ux1ed9-quux1ed1c-tux1ebf-huxf3a-vuxe0-quy-ux1b0ux1edbc-trung-tuxe2m-huxf3a}{%
\subsection{3.3.3.1 Vận hành hóa cường độ quốc tế hóa và quy ước trung
tâm
hóa}\label{vux1eadn-huxe0nh-huxf3a-cux1b0ux1eddng-ux111ux1ed9-quux1ed1c-tux1ebf-huxf3a-vuxe0-quy-ux1b0ux1edbc-trung-tuxe2m-huxf3a}}

Cường độ quốc tế hóa được đo bằng tỷ trọng doanh thu xuất khẩu trên tổng
doanh thu (Foreign Sales To Total Sales, FSTS), giới hạn miền giá trị từ
0 đến 1. Doanh nghiệp không xuất khẩu nhận giá trị FSTS bằng 0, một giá
trị hợp lệ phản ánh trạng thái phi xuất khẩu chứ không phải dữ liệu
khuyết. Khung đa quốc gia cộng gộp xuất khẩu trực tiếp và gián tiếp
(tổng hai trường d3b và d3c), đúng định nghĩa tham gia xuất khẩu chuẩn
của bộ công cụ WBES; hai nghiên cứu đơn quốc gia (Việt Nam, Trung Quốc)
dùng riêng xuất khẩu trực tiếp (trường d3c), vì xuất khẩu gián tiếp ủy
thác chức năng quốc tế hóa cho trung gian thương mại với cơ chế học hỏi
và yêu cầu nguồn lực khác về bản chất.

Để kiểm định phi tuyến, mô hình đưa đồng thời số hạng bậc nhất và số
hạng bậc hai của cường độ xuất khẩu. Trước khi tạo số hạng bậc hai,
cường độ xuất khẩu được trung tâm hóa quanh trung bình trong từng đợt
khảo sát hoặc trong từng mẫu gộp, một quy ước làm giảm tương quan giả
tạo giữa số hạng tuyến tính và số hạng bình phương, qua đó cải thiện độ
ổn định số học của ước lượng và làm rõ diễn giải hệ số tại điểm trung
bình của mẫu. Điểm uốn của quan hệ chữ U ngược được khôi phục về thang
gốc bằng cách cộng trung bình mẫu vào tỷ số
\(-\hat\beta_1 / (2\hat\beta_2)\), cho phép báo cáo ngưỡng tối ưu theo
tỷ trọng xuất khẩu thực tế. Dạng bậc hai được dùng nhất quán trong mọi
nghiên cứu thành phần, vì nó đủ để nhận dạng một điểm uốn nội miền và
phù hợp với mật độ phân phối cường độ xuất khẩu của các mẫu doanh nghiệp
WBES; luận án không mở rộng lên dạng bậc ba để tránh nhận dạng điểm uốn
không ổn định ở vùng cường độ xuất khẩu cao vốn thưa quan sát.

\hypertarget{vux1eadn-huxe0nh-huxf3a-chi-tiux1ebft-cux1b0ux1eddng-ux111ux1ed9-xuux1ea5t-khux1ea9u-xuyuxean-cuxe1c-nghiuxean-cux1ee9u-thuxe0nh-phux1ea7n}{%
\subsection{3.3.3.2 Vận hành hóa chi tiết cường độ xuất khẩu xuyên các
nghiên cứu thành
phần}\label{vux1eadn-huxe0nh-huxf3a-chi-tiux1ebft-cux1b0ux1eddng-ux111ux1ed9-xuux1ea5t-khux1ea9u-xuyuxean-cuxe1c-nghiuxean-cux1ee9u-thuxe0nh-phux1ea7n}}

Sự thống nhất khái niệm của cường độ xuất khẩu xuyên các nghiên cứu
thành phần không kéo theo một định nghĩa thao tác duy nhất, mà là một họ
định nghĩa được hiệu chỉnh theo bản chất bộ công cụ của từng đợt và theo
câu hỏi nghiên cứu. Khung đa quốc gia cộng gộp xuất khẩu trực tiếp và
gián tiếp rồi chia cho một trăm để đưa về tỷ trọng, cụ thể
\(FSTS = (d3b + d3c)/100\) với miền giá trị từ 0 đến 1, đúng định nghĩa
tham gia xuất khẩu chuẩn của bộ công cụ toàn cầu; trung bình mẫu của tỷ
trọng này trong mẫu gộp năm mươi nền là 9,0\%. Ngược lại, hai nghiên cứu
đơn quốc gia dùng riêng xuất khẩu trực tiếp dựng từ một trường khảo sát:
ở Trung Quốc cường độ xuất khẩu là tỷ trọng doanh thu xuất khẩu trực
tiếp, còn ở đảo nhỏ Thái Bình Dương cường độ xuất khẩu được dựng từ
trường \texttt{d3c} chia cho một trăm. Lý do lựa chọn riêng xuất khẩu
trực tiếp đã nêu ở Mục 3.3.3.1: xuất khẩu gián tiếp ủy thác chức năng
quốc tế hóa cho trung gian thương mại, mang cơ chế học hỏi và yêu cầu
nguồn lực khác về bản chất.

Doanh nghiệp không xuất khẩu nhận giá trị cường độ bằng 0 và được giữ
trong mẫu, vì giá trị này phản ánh trạng thái phi xuất khẩu hợp lệ chứ
không phải dữ liệu khuyết; điểm uốn của quan hệ do đó được nhận dạng
trên toàn mẫu chứ không chỉ trên tập doanh nghiệp xuất khẩu. Quy ước này
đặc biệt hệ trọng ở đảo nhỏ Thái Bình Dương, nơi chỉ 132 trong số 959
quan sát phân tích, tương đương 13,8\%, là doanh nghiệp có xuất khẩu với
cường độ dương: loại bỏ nhóm phi xuất khẩu sẽ làm sụp đổ công suất thống
kê và làm sai lệch ước lượng ở biên tham gia.

Quy ước trung tâm hóa được áp dụng nhất quán nhưng theo hai cấp khác
nhau. Ở mẫu gộp đa quốc gia, cường độ xuất khẩu được trung tâm hóa quanh
trung bình toàn mẫu trước khi lập bình phương. Ở các nghiên cứu đa đợt,
cường độ xuất khẩu được trung tâm hóa trong từng đợt nhằm chặn khác biệt
thành phần mẫu giữa các đợt làm nhiễu vị trí điểm uốn. Trung tâm hóa làm
giảm tương quan giả tạo giữa số hạng tuyến tính và số hạng bình phương
mà không làm dịch chuyển vị trí điểm uốn, vì công thức điểm uốn
\(-\hat\beta_1/(2\hat\beta_2)\) bất biến với phép trung tâm hóa; điểm
uốn trên thang trung tâm hóa được quy đổi về thang gốc bằng cách cộng
trung bình mẫu, tức
\(TP = -\hat\beta_1/(2\hat\beta_2) + \overline{FSTS}\).

\hypertarget{ux111o-lux1b0ux1eddng-biux1ebfn-ux111iux1ec1u-tiux1ebft}{%
\subsection{3.3.4 Đo lường biến điều
tiết}\label{ux111o-lux1b0ux1eddng-biux1ebfn-ux111iux1ec1u-tiux1ebft}}

\hypertarget{digital-constructs-tci-vuxe0-dai}{%
\subsubsection{Digital constructs (TCI và
DAI)}\label{digital-constructs-tci-vuxe0-dai}}

Luận án phân biệt rõ hai cấu trúc:

\begin{itemize}
\tightlist
\item
  \textbf{Technological Capability Index (TCI)}: chiều sâu năng lực công
  nghệ, tổng hợp từ cấp phép công nghệ nước ngoài, chi tiêu nghiên cứu
  và phát triển, đổi mới sản phẩm, và chứng nhận chất lượng.
\item
  \textbf{Digital Adoption Index (DAI)}: mức độ chấp nhận số bề mặt,
  tổng hợp từ website, thư điện tử, bán hàng trực tuyến (DAI\_thin) hoặc
  bổ sung thanh toán điện tử và nhà cung cấp điện tử (DAI\_rich).
\end{itemize}

\textbf{Nguyên tắc không chồng lấn}: TCI và DAI không chia sẻ thành phần
chỉ báo để bảo đảm độ thuần cấu trúc. Đặc biệt, cấp phép công nghệ nước
ngoài thuộc TCI chứ không được đồng thời tính vào DAI.

\textbf{Biến thể đo lường (rút gọn/đầy đủ).} Do bộ công cụ WBES thay đổi
qua ba thế hệ bảng hỏi (PICS3, Standardized, BREADY) khiến không phải
mọi mục đều khả dụng ở mọi sóng, TCI được vận hành hóa thống nhất theo
hai biến thể của cùng một chỉ số hợp thành cấu thành:

\begin{itemize}
\tightlist
\item
  \textbf{TCI\_thin} = chuẩn hóa z của trung bình hai mục cốt lõi: chứng
  chỉ chất lượng quốc tế (b8) và cấp phép công nghệ nước ngoài (e6).
  Dùng trong các nghiên cứu Việt Nam (P3), Singapore (P4), đa quốc gia
  (P7) và đảo nhỏ Thái Bình Dương (P8).
\item
  \textbf{TCI\_full} = chuẩn hóa z của trung bình bốn mục: b8, e6, đổi
  mới sản phẩm (h1), và chi R\&D (h8). Dùng trong nghiên cứu Trung Quốc
  (P5), nơi cả hai sóng 2012 và 2024 đều có đủ module đổi mới.
\end{itemize}

Bộ mục chính xác của mỗi nghiên cứu thành phần được nêu lại trong bảng
định nghĩa biến của nghiên cứu đó (Mục 3.4.5) và Phụ lục A; khác biệt
giữa biến thể rút gọn và đầy đủ phản ánh ràng buộc khả dụng dữ liệu
khách quan, không phải lựa chọn tùy ý. Vì các mục là chỉ báo nhị phân
của cùng một cấu trúc cấu thành (năng lực công nghệ), hai biến thể tương
quan cao và cho kết luận điều tiết nhất quán về dấu; phần độ vững (Mục
3.5.1) kiểm tra độ nhạy của kết quả với lựa chọn biến thể. Bộ mục TCI
của nghiên cứu đa quốc gia (P7) đã được đối chiếu trực tiếp với mã xây
dựng dữ liệu thực tế (quy trình xử lý dữ liệu thô WBES 50 nền kinh tế,
gồm Nhật Bản 2025): biến \texttt{tci\_z} được dựng đúng từ hai mục b8 và
e6, thống nhất với khung rút gọn nêu trên.

\textbf{Cơ sở kế thừa}: Bhandari et al.~(2023) phân biệt số hóa, quốc tế
hóa và năng lực theo logic điều phối nguồn lực; Verhoef et al.~(2021)
phân biệt ba giai đoạn (số hóa dữ liệu, số hóa quy trình, chuyển đổi
số). Khuôn mẫu ``hiệu ứng lá chắn số'' được nghiên cứu thành phần P1
phát hiện trên 17 nước châu Á mới nổi.

\textbf{Đóng góp mới}: Đưa ra \textbf{giao thức hài hòa hóa đo lường}
trong bối cảnh WBES, đảm bảo độ thuần cấu trúc ở 50 nền kinh tế châu Á
và Thái Bình Dương, với quy trình hòa hợp 3 thế hệ lược đồ. Phân tách
TCI và DAI theo (i) Bharadwaj et al.~(2013), năng lực công nghệ thông
tin so với chiến lược kinh doanh số có mạng quan hệ lý thuyết khác nhau;
(ii) phân tầng 3 cấp của Verhoef et al.~(2021), số hóa dữ liệu, số hóa
quy trình, chuyển đổi số; (iii) truyền thống Lall (1992) và Cohen \&
Levinthal (1990) coi năng lực công nghệ là chiều sâu năng lực nội tại
(absorptive capacity), khác về bản chất với chấp nhận số ở giao diện
ngoại tại; (iv) logic điều phối nguồn lực của Bhandari et al.~(2023) cho
quan hệ quốc tế hóa và hiệu quả; và (v) 4 tiêu chí chỉ số hợp thành cấu
thành của Coltman et al.~(2008), đảm bảo TCI và DAI là hai chỉ số hợp
thành độc lập.

\hypertarget{biux1ebfn-thux1ec3-chux1ebf}{%
\subsubsection{Biến thể chế}\label{biux1ebfn-thux1ec3-chux1ebf}}

\textbf{Biến chính}: business obstacle index (tổng hợp từ các câu hỏi
WBES về obstacles), country-level governance proxies, chỉ số pháp quyền
WGI (Kaufmann et al., 2011), 6 chế độ con ICRV (classification từ Khanna
\& Palepu, 2010 mở rộng).

\textbf{Cơ sở kế thừa}: Khanna và Palepu (2010) với khoảng trống thể
chế; North (1990) với institutional analysis; Marano et al.~(2016) với
phân tích tổng hợp evidence về institutional các biến điều tiết.

\hypertarget{ux111ux1eb7c-ux111iux1ec3m-nhuxe0-quux1ea3n-trux1ecb-cux1ea5p-cao}{%
\subsubsection{Đặc điểm nhà quản trị cấp
cao}\label{ux111ux1eb7c-ux111iux1ec3m-nhuxe0-quux1ea3n-trux1ecb-cux1ea5p-cao}}

\textbf{Biến chính}: top manager experience (năm kinh nghiệm), top
manager gender (female / male).

\textbf{Cơ sở kế thừa}: Hambrick và Mason (1984), Hambrick (2007), Hsu
et al.~(2013), Post và Byron (2015).

\textbf{Điều kiện sử dụng}: chỉ trong subset có đủ dữ liệu; nếu thiếu,
đưa vào độ vững chứ không bắt buộc trong mô hình chính.

\hypertarget{phuxe2n-tux1ea7ng-chux1ec9-sux1ed1-chux1ea5p-nhux1eadn-sux1ed1-vuxe0-quy-tux1eafc-xuxe2y-dux1ef1ng-chux1ec9-sux1ed1-hux1ee3p-thuxe0nh-cux1ea5u-thuxe0nh}{%
\subsection{3.3.4.1 Phân tầng chỉ số chấp nhận số và quy tắc xây dựng
chỉ số hợp thành cấu
thành}\label{phuxe2n-tux1ea7ng-chux1ec9-sux1ed1-chux1ea5p-nhux1eadn-sux1ed1-vuxe0-quy-tux1eafc-xuxe2y-dux1ef1ng-chux1ec9-sux1ed1-hux1ee3p-thuxe0nh-cux1ea5u-thuxe0nh}}

Chỉ số chấp nhận số (Digital Adoption Index, DAI) được vận hành hóa theo
nguyên tắc phân tầng nhằm phản ánh chiều sâu tăng dần của số hóa doanh
nghiệp. Tầng 1 ghi nhận hiện diện số bề mặt qua chỉ báo website (trường
c22b), một mục nhị phân khả dụng ở hầu hết các đợt khảo sát. Tầng 2 bổ
sung các chức năng số kích hoạt giao dịch, cụ thể là thanh toán điện tử
và tỷ trọng doanh thu qua kênh điện tử (trường k33), những mục chỉ xuất
hiện ở các đợt khảo sát gần đây. Việc phân tầng tạo ra ba biến thể đo
lường: biến thể chỉ Tầng 1 (đơn biến nhị phân) dùng ở Việt Nam và đảo
nhỏ Thái Bình Dương, biến thể Tầng 1 cộng Tầng 2 dùng ở Singapore và
khung đa quốc gia khi trường k33 có mặt, và biến thể thu gọn về Tầng 1
khi trường Tầng 2 thiếu. Khác biệt giữa các biến thể phản ánh ràng buộc
khả dụng dữ liệu khách quan qua ba thế hệ bảng hỏi chứ không phải lựa
chọn tùy ý.

Cả chỉ số năng lực công nghệ (Technological Capability Index, TCI) và
DAI được xử lý như chỉ số hợp thành cấu thành (formative composite) chứ
không phải chỉ số phản ánh, vì mỗi chỉ báo nhị phân định nghĩa một khía
cạnh độc lập của khái niệm chứ không phải là biểu hiện tương quan của
một biến tiềm ẩn chung. Quy tắc xây dựng gồm ba bước: chuyển mỗi chỉ báo
thành dạng nhị phân 0--1, lấy trung bình cộng các chỉ báo trong cùng một
tầng, rồi chuẩn hóa z giá trị trung bình trong từng đợt khảo sát. Nguyên
tắc không chồng lấn được áp đặt nghiêm ngặt: cấp phép công nghệ nước
ngoài thuộc TCI và không được đồng thời tính vào DAI, bảo đảm hai chỉ số
là hai cấu trúc độc lập với mạng quan hệ lý thuyết khác nhau. Vì các chỉ
báo là biểu hiện nhị phân của cùng một cấu trúc cấu thành, các biến thể
của một chỉ số tương quan cao và cho kết luận điều tiết nhất quán về
dấu; phần độ vững kiểm tra trực tiếp độ nhạy của suy luận với lựa chọn
biến thể.

\hypertarget{uxe1nh-xux1ea1-mux1ee5c-khux1ea3o-suxe1t-cux1ee7a-chux1ec9-sux1ed1-nux103ng-lux1ef1c-cuxf4ng-nghux1ec7-vuxe0-chux1ec9-sux1ed1-chux1ea5p-nhux1eadn-sux1ed1}{%
\subsection{3.3.4.2 Ánh xạ mục khảo sát của chỉ số năng lực công nghệ và
chỉ số chấp nhận
số}\label{uxe1nh-xux1ea1-mux1ee5c-khux1ea3o-suxe1t-cux1ee7a-chux1ec9-sux1ed1-nux103ng-lux1ef1c-cuxf4ng-nghux1ec7-vuxe0-chux1ec9-sux1ed1-chux1ea5p-nhux1eadn-sux1ed1}}

Để bảo đảm tính minh bạch và tái lập, mỗi biến thể của hai chỉ số hợp
thành được truy nguyên về đúng bộ mục khảo sát WBES dùng để dựng nó.
Biến thể mỏng của chỉ số năng lực công nghệ được dựng từ đúng hai mục
nhị phân: chứng chỉ chất lượng được công nhận quốc tế (trường
\texttt{b8}) và cấp phép công nghệ nước ngoài (trường \texttt{e6}); giá
trị chỉ số là chuẩn hóa z của trung bình hai chỉ báo nhị phân này, xử lý
như một chỉ số hợp thành cấu thành. Bộ mục hai trường này được dùng nhất
quán ở khung đa quốc gia, ở Việt Nam, Singapore và đảo nhỏ Thái Bình
Dương. Biến thể đầy đủ chỉ dùng ở Trung Quốc, nơi cả hai đợt đều có đủ
module đổi mới, mở rộng lên bốn mục để nắm bắt chiều sâu năng lực nội
tại phong phú hơn so với biến thể hai mục dùng ở các nghiên cứu khác.
Khác biệt giữa biến thể mỏng và đầy đủ phản ánh ràng buộc khả dụng dữ
liệu khách quan qua ba thế hệ bảng hỏi chứ không phải lựa chọn tùy ý; vì
các mục là chỉ báo nhị phân của cùng một cấu trúc cấu thành, hai biến
thể cho kết luận điều tiết nhất quán về dấu, và phần độ vững (Mục 3.5.1)
kiểm tra trực tiếp độ nhạy với lựa chọn biến thể.

Chỉ số chấp nhận số được vận hành hóa theo nguyên tắc phân tầng đã trình
bày ở Mục 3.3.4.1, với bộ mục cụ thể cũng được truy nguyên về trường
khảo sát. Biến thể Tầng 1 dùng đúng một chỉ báo nhị phân là hiện diện
website vận hành (trường \texttt{c22b}); đây là biến thể duy nhất khả
dụng ở đảo nhỏ Thái Bình Dương, nơi bộ công cụ không thu thập các mục
Tầng 2 thanh toán điện tử và Tầng 3 đám mây/hệ hoạch định nguồn lực
doanh nghiệp. Khung đa quốc gia dùng biến thể Tầng 1 cộng Tầng 2 khi
trường thanh toán điện tử có mặt và thu gọn về Tầng 1 khi thiếu, đúng
như nguyên tắc nhúng tầng theo khả dụng bộ công cụ. Nguyên tắc không
chồng lấn được áp đặt nghiêm ngặt: cấp phép công nghệ nước ngoài thuộc
chỉ số năng lực công nghệ và không bao giờ được tính đồng thời vào chỉ
số chấp nhận số, bảo đảm hai chỉ số là hai cấu trúc độc lập với mạng
quan hệ lý thuyết khác nhau. Bộ mục năng lực công nghệ của khung đa quốc
gia đã được đối chiếu trực tiếp với mã xây dựng dữ liệu thực tế trên kho
năm mươi nền, gồm Nhật Bản 2025, xác nhận biến \texttt{tci\_z} được dựng
đúng từ hai mục \texttt{b8} và \texttt{e6}.

\hypertarget{biux1ebfn-kiux1ec3m-souxe1t}{%
\subsection{3.3.5 Biến kiểm soát}\label{biux1ebfn-kiux1ec3m-souxe1t}}

Khung thực nghiệm đưa vào một tập biến kiểm soát chuẩn trong nghiên cứu
dữ liệu doanh nghiệp và kinh doanh quốc tế, nhằm tách hiệu ứng của quốc
tế hóa và năng lực khỏi các yếu tố cấu trúc đồng biến với cả cường độ
xuất khẩu lẫn năng suất. Quy mô doanh nghiệp được đo bằng logarit tự
nhiên của số lao động toàn thời gian (mục l1 của WBES), một biến kiểm
soát thiết yếu vì doanh nghiệp lớn vừa có xác suất xuất khẩu cao hơn vừa
đạt năng suất cao hơn nhờ lợi thế quy mô; lấy logarit để xử lý phân phối
lệch phải đặc trưng của biến quy mô. Tuổi doanh nghiệp được tính bằng số
năm từ năm thành lập đến năm khảo sát, kiểm soát hiệu ứng học hỏi tích
lũy và hiệu ứng sống sót, theo đó doanh nghiệp tồn tại lâu hơn có thể đã
vượt qua các rào cản gia nhập và tích lũy năng lực.

Sở hữu nước ngoài được đưa vào dưới dạng biến giả nhận giá trị một khi
tỷ lệ sở hữu của nhà đầu tư nước ngoài đạt từ 10\% trở lên, ngưỡng chuẩn
của đầu tư trực tiếp nước ngoài, vì doanh nghiệp có vốn nước ngoài
thường tiếp cận mạng lưới quốc tế, công nghệ và thị trường theo cách
khác biệt về chất so với doanh nghiệp trong nước. Ngành được kiểm soát
bằng hệ biến giả ngành nhằm hấp thụ chênh lệch năng suất và cường độ
thương mại có hệ thống giữa các lĩnh vực. Quan trọng nhất, mọi đặc tả
gộp đều mang hiệu ứng cố định nền kinh tế và hiệu ứng cố định năm: hiệu
ứng cố định nền kinh tế hấp thụ mọi khác biệt bất biến theo thời gian
giữa các nền kinh tế, gồm mặt bằng giá, chế độ tỷ giá, chất lượng thể
chế nền và quy mô thị trường nội địa; hiệu ứng cố định năm hấp thụ các
cú sốc vĩ mô chung như chu kỳ kinh tế toàn cầu và biến động thương mại
theo năm. Cấu trúc hiệu ứng cố định hai chiều này là then chốt vì nó
loại bỏ chính các nhiễu cấp nền kinh tế và cấp thời gian dễ tạo tương
quan giả giữa quốc tế hóa và hiệu quả, qua đó cô lập biến thiên trong
nội bộ nền kinh tế và năm để nhận dạng quan hệ trọng tâm.

\textbf{Cơ sở kế thừa}: tập biến kiểm soát và cấu trúc hiệu ứng cố định
tuân theo chuẩn trong nghiên cứu dữ liệu doanh nghiệp WBES và kinh doanh
quốc tế (Aterido et al., 2011; Ayyagari et al., 2011).

\hypertarget{vux1eadn-huxe0nh-huxf3a-chux1ebf-ux111ux1ed9-thux1ec3-chux1ebf-icrv-vuxe0-ux111ux1eb7c-ux111iux1ec3m-nhuxe0-quux1ea3n-trux1ecb}{%
\subsection{3.3.5.1 Vận hành hóa chế độ thể chế ICRV và đặc điểm nhà
quản
trị}\label{vux1eadn-huxe0nh-huxf3a-chux1ebf-ux111ux1ed9-thux1ec3-chux1ebf-icrv-vuxe0-ux111ux1eb7c-ux111iux1ec3m-nhuxe0-quux1ea3n-trux1ecb}}

Chế độ thể chế được vận hành hóa thành một biến phân loại số nguyên nhận
giá trị từ một đến sáu theo phân loại ICRV, đi từ chế độ tiên tiến đổi
mới đến chế độ đảo nhỏ Thái Bình Dương, trật tự tăng dần của số nhóm
tương ứng với chất lượng thể chế giảm dần. Biến này được dùng theo hai
vai trò trong khung đa quốc gia: làm biến điều tiết liên tục trong mô
hình tương tác, nơi hiệu ứng chính của nó bị hấp thụ bởi hiệu ứng cố
định nền kinh tế, và làm biến phân nhóm cho phân tích điểm uốn theo từng
chế độ. Lựa chọn xử lý ICRV như biến số nguyên liên tục thay vì hệ biến
giả là có chủ đích, nhằm giữ cỡ mẫu đủ lớn trong các mô hình điều tiết
và ba chiều, nơi việc tách thành sáu biến giả tương tác sẽ làm phân mảnh
mẫu và làm bất ổn ước lượng.

Phân tích điểm uốn theo từng chế độ làm lộ rõ cấu trúc ba vùng đặc trưng
của khu vực. Ở chế độ chuyển đổi thu nhập trung bình--thấp, vốn là chế
độ neo ước lượng gộp, quan hệ chữ U ngược định hình sắc nét với điểm uốn
quanh 43\%. Ở chế độ tiên tiến đổi mới, gồm cả Nhật Bản, điểm uốn nằm
gần hoặc vượt biên dữ liệu quan sát nên trong vùng dữ liệu quan hệ gần
như tuyến tính dương, với điểm uốn xấp xỉ 80\%. Ở các chế độ yếu nhất,
chữ U ngược tan rã; riêng ở chế độ đảo nhỏ Thái Bình Dương số hạng bậc
hai đổi dấu, không còn cực đại nội miền, nhất quán với gánh nặng quốc tế
hóa bắt buộc ghi nhận ở nghiên cứu thành phần đảo nhỏ.

Đặc điểm nhà quản trị cấp cao được vận hành hóa từ hai mục khảo sát: số
năm kinh nghiệm của nhà quản trị cấp cao trong ngành (trường
\texttt{b7}) và một biến nhị phân chỉ nhà quản trị cấp cao là nữ (trường
\texttt{b7a} hoặc \texttt{b6a} tùy đợt). Hai biến này được đưa vào mô
hình đa quốc gia để kiểm định hiệu ứng trực tiếp và hiệu ứng điều tiết
của đặc điểm quản trị lên hình dạng đường quan hệ; ở các nghiên cứu
thiếu dữ liệu quản trị, hai biến chỉ được dùng trong phân tích độ vững
chứ không bắt buộc trong mô hình chính, đúng điều kiện sử dụng đã nêu.

\hypertarget{tux1ed5ng-hux1ee3p-biux1ebfn-nghiuxean-cux1ee9u}{%
\subsection{3.3.6 Tổng hợp biến nghiên
cứu}\label{tux1ed5ng-hux1ee3p-biux1ebfn-nghiuxean-cux1ee9u}}

Toàn bộ biến của khung thực nghiệm được tổng hợp trong Bảng 3.1, nêu rõ
với mỗi biến: vai trò trong mô hình, ký hiệu, trường khảo sát WBES
nguồn, cách vận hành hóa và giả thuyết liên quan. Bảng này bảo đảm tính
minh bạch đo lường và khả năng tái lập (Aguinis et al., 2019); chi tiết
vận hành hóa theo từng nghiên cứu thành phần được trình bày trong các
bảng định nghĩa biến ở Mục 3.4.5.

\textbf{Bảng 3.1.} \emph{Tổng hợp biến nghiên cứu của khung thực nghiệm
đa quốc gia.}

\begin{longtable}[]{@{}
  >{\raggedright\arraybackslash}p{(\columnwidth - 8\tabcolsep) * \real{0.2000}}
  >{\raggedright\arraybackslash}p{(\columnwidth - 8\tabcolsep) * \real{0.2000}}
  >{\raggedright\arraybackslash}p{(\columnwidth - 8\tabcolsep) * \real{0.2000}}
  >{\raggedright\arraybackslash}p{(\columnwidth - 8\tabcolsep) * \real{0.2000}}
  >{\raggedright\arraybackslash}p{(\columnwidth - 8\tabcolsep) * \real{0.2000}}@{}}
\toprule\noalign{}
\begin{minipage}[b]{\linewidth}\raggedright
Vai trò
\end{minipage} & \begin{minipage}[b]{\linewidth}\raggedright
Biến (ký hiệu)
\end{minipage} & \begin{minipage}[b]{\linewidth}\raggedright
Trường WBES
\end{minipage} & \begin{minipage}[b]{\linewidth}\raggedright
Vận hành hóa
\end{minipage} & \begin{minipage}[b]{\linewidth}\raggedright
Giả thuyết
\end{minipage} \\
\midrule\noalign{}
\endhead
\bottomrule\noalign{}
\endlastfoot
Biến phụ thuộc & Năng suất lao động, ln(LP) & \texttt{d2}, \texttt{l1} &
Logarit tự nhiên của doanh thu trên số lao động thường xuyên & --- \\
Biến độc lập & Cường độ xuất khẩu (FSTS) & \texttt{d3c} (và
\texttt{d3b}) & Tỷ lệ doanh thu nước ngoài trên tổng doanh thu (0--1),
trung tâm hóa & H1 \\
Biến độc lập & FSTS bình phương (FSTS²) & (suy ra) & Bình phương của
FSTS đã trung tâm hóa, kiểm định phi tuyến & H1, H1b \\
Điều tiết -- tầng doanh nghiệp & Năng lực công nghệ (TCI) & \texttt{b8},
\texttt{e6} & Chỉ số hợp thành chuẩn hóa z (chứng chỉ chất lượng + công
nghệ ngoại nhập) & H2 \\
Điều tiết -- tầng doanh nghiệp & Chấp nhận số (DAI) & \texttt{c22b},
\texttt{k33}/\texttt{k38} & Tầng 1: website (nhị phân); Tầng 2: hợp
thành z (website + thanh toán điện tử) & H3 \\
Điều tiết -- tầng cá nhân & Kinh nghiệm nhà quản trị & \texttt{b7} & Số
năm kinh nghiệm của nhà quản trị cấp cao trong ngành & H4 \\
Điều tiết -- tầng cá nhân & Giới tính nhà quản trị &
\texttt{b7a}/\texttt{b6a} & Biến nhị phân: nhà quản trị cấp cao là nữ &
H4 \\
Điều tiết -- tầng quốc gia & Chế độ thể chế (ICRV) & (phân loại) & Biến
số nguyên 1--6 theo khung ICRV (thể chế giảm dần) & H5 \\
Điều tiết -- thời gian & Đợt khảo sát & (năm khảo sát) & Biến giả
đợt/năm; kiểm định ổn định hệ số Paternoster & H6 \\
Kiểm soát & Quy mô doanh nghiệp & \texttt{l1} & Logarit số lao động
thường xuyên & --- \\
Kiểm soát & Tuổi doanh nghiệp & \texttt{b5} & Năm khảo sát trừ năm thành
lập & --- \\
Kiểm soát & Sở hữu nước ngoài & \texttt{b2b} & Tỷ lệ vốn sở hữu nước
ngoài (hoặc biến nhị phân ≥10\%) & --- \\
Kiểm soát & Hiệu ứng cố định & --- & Ngành (FE), nền kinh tế (FE), năm
(FE) & --- \\
\end{longtable}

\emph{Ghi chú: trường WBES tham chiếu mã biến của bộ công cụ khảo sát;
một số mã thay đổi theo thế hệ lược đồ (PICS3 / Standardized / BREADY)
và được hài hòa hóa theo Phụ lục A. Vận hành hóa chi tiết theo từng
nghiên cứu thành phần xem Mục 3.4.5. Nguồn: tác giả tổng hợp.}

\hypertarget{muxf4-huxecnh-phuxe2n-tuxedch}{%
\section{3.4 Mô hình phân tích}\label{muxf4-huxecnh-phuxe2n-tuxedch}}

\hypertarget{muxf4-huxecnh-phuxe2n-tuxedch-tux1ed5ng-hux1ee3p}{%
\subsection{3.4.1 Mô hình phân tích tổng
hợp}\label{muxf4-huxecnh-phuxe2n-tuxedch-tux1ed5ng-hux1ee3p}}

Mô hình hiệu ứng ngẫu nhiên dựa trên phép biến đổi \(z\) của Fisher:

\[z_i = \frac{1}{2} \ln\left(\frac{1+r_i}{1-r_i}\right)\]

với \(r_i\) là cỡ ảnh hưởng của nghiên cứu \(i\). Cỡ ảnh hưởng gộp được
tính theo trọng số nghịch đảo phương sai (Borenstein et al., 2009).

Dị biệt được đo bằng kiểm định \(Q\) và chỉ số \(I^2\) (Higgins et al.,
2003). Phân tích theo nhóm con và hồi quy phân tích tổng hợp được dùng
để kiểm định các biến điều tiết.

\textbf{Cơ sở kế thừa}: chuẩn phân tích tổng hợp trong quản trị chiến
lược (Combs et al., 2011; Aguinis et al., 2011).

\hypertarget{muxf4-huxecnh-thux1ef1c-nghiux1ec7m-tuyux1ebfn-tuxednh-cux1a1-bux1ea3n}{%
\subsection{3.4.2 Mô hình thực nghiệm: Tuyến tính cơ
bản}\label{muxf4-huxecnh-thux1ef1c-nghiux1ec7m-tuyux1ebfn-tuxednh-cux1a1-bux1ea3n}}

\[\ln(LP)_i = \beta_0 + \beta_1 FSTS_i + \boldsymbol{\gamma} \mathbf{X}_i + \varepsilon_i\]

với \(\mathbf{X}_i\) là vector control variables.

\hypertarget{muxf4-huxecnh-phi-tuyux1ebfn-h1}{%
\subsection{3.4.3 Mô hình phi tuyến
(H1)}\label{muxf4-huxecnh-phi-tuyux1ebfn-h1}}

\[\ln(LP)_i = \beta_0 + \beta_1 FSTS_i + \beta_2 FSTS_i^2 + \boldsymbol{\gamma} \mathbf{X}_i + \varepsilon_i\]

Nếu \(\beta_1 > 0\) và \(\beta_2 < 0\), chữ U ngược được xác nhận. Điểm
uốn tính theo \(-\beta_1 / (2\beta_2)\).

Đối với Trung Quốc, mô hình giữ nguyên dạng bậc hai trên và được ước
lượng riêng theo từng đợt khảo sát (2012 và 2024) cũng như trên mẫu gộp,
kèm các kiểm định ổn định hệ số xuyên đợt để kiểm tra giả thuyết dị biệt
thời gian H6, thay vì mở rộng lên dạng bậc ba; dạng bậc hai đủ để nhận
dạng điểm uốn và phù hợp với mật độ phân phối cường độ xuất khẩu của mẫu
doanh nghiệp Trung Quốc.

\textbf{Cơ sở kế thừa}: Lu và Beamish (2004), Hitt et al.~(1997),
Contractor et al.~(2003); nghiên cứu thành phần P2 cho dạng bậc hai ở
Trung Quốc với điểm uốn khoảng 47,8\% FSTS, và nghiên cứu thành phần P5
xác nhận dạng bậc hai ổn định theo thời gian với điểm uốn quanh 48\%
FSTS.

\textbf{Kiểm định chữ U của Lind và Mehlum} (Lind \& Mehlum, 2010; Haans
et al., 2016) để xác nhận chặt chẽ chữ U ngược: kiểm định độ dốc dương ở
bên trái và độ dốc âm ở bên phải của miền giá trị FSTS.

\hypertarget{thux1ee7-tux1ee5c-kiux1ec3m-ux111ux1ecbnh-chux1eef-u-ngux1b0ux1ee3c-vuxe0-ux1b0ux1edbc-lux1b0ux1ee3ng-khoux1ea3ng-tin-cux1eady-ux111iux1ec3m-uux1ed1n}{%
\subsection{3.4.3.1 Thủ tục kiểm định chữ U ngược và ước lượng khoảng
tin cậy điểm
uốn}\label{thux1ee7-tux1ee5c-kiux1ec3m-ux111ux1ecbnh-chux1eef-u-ngux1b0ux1ee3c-vuxe0-ux1b0ux1edbc-lux1b0ux1ee3ng-khoux1ea3ng-tin-cux1eady-ux111iux1ec3m-uux1ed1n}}

Sự xuất hiện đồng thời của hệ số bậc nhất dương và hệ số bậc hai âm là
điều kiện cần nhưng chưa đủ để khẳng định quan hệ chữ U ngược, vì cực
trị toán học có thể nằm ngoài miền dữ liệu quan sát, khiến quan hệ thực
tế chỉ đơn điệu trong phạm vi giá trị có ý nghĩa. Luận án vì vậy áp dụng
kiểm định U của Lind và Mehlum (2010) như chuẩn xác nhận chính thức. Thủ
tục đòi hỏi ba điều kiện đồng thời: hệ số bậc nhất dương và hệ số bậc
hai âm; độ dốc của hàm tại biên trái của miền dương và độ dốc tại biên
phải âm; và điểm uốn nằm chặt bên trong khoảng giá trị quan sát của
cường độ xuất khẩu. Kiểm định gộp ba ràng buộc trên thành một kiểm định
giao điểm theo nguyên lý Sasabuchi, cho ra một giá trị p duy nhất; chỉ
khi giá trị p này đủ nhỏ thì hình chữ U ngược mới được xác nhận về mặt
thống kê chứ không chỉ về mặt mô tả.

Điểm uốn được tính theo công thức
\(\text{TP} = -\hat\beta_1 / (2\hat\beta_2)\) trên thang trung tâm hóa
rồi quy đổi về thang gốc bằng cách cộng trung bình mẫu. Vì điểm uốn là
tỷ số phi tuyến của hai hệ số ước lượng, sai số chuẩn của nó không thể
đọc trực tiếp từ bảng hồi quy. Luận án báo cáo khoảng tin cậy 95 phần
trăm theo phương pháp delta (delta method), khai triển tuyến tính bậc
nhất quanh ước lượng điểm để xấp xỉ phương sai của tỷ số (Haans et al.,
2016). Ở các nghiên cứu có mật độ thưa quanh điểm uốn, luận án bổ sung
khoảng tin cậy theo phương pháp lấy mẫu lại bootstrap nhằm nới lỏng giả
định chuẩn của phương pháp delta và phản ánh tính bất đối xứng của phân
phối điểm uốn trong mẫu hữu hạn. Khi kiểm định không bác bỏ tính đơn
điệu trong phạm vi dữ liệu, kết quả được diễn giải là thông tin tích cực
theo khung bão hòa chứ không phải là thất bại của giả thuyết phi tuyến,
vì nó cho biết doanh nghiệp trong mẫu chưa vượt qua ngưỡng nơi lợi suất
quốc tế hóa bắt đầu suy giảm.

Cách vận hành thủ tục này thể hiện rõ ở các nghiên cứu thành phần. Ở
Trung Quốc, kiểm định của Lind và Mehlum được kiểm chứng cho từng đợt
qua điều kiện độ dốc dương có ý nghĩa tại biên trái khi cường độ xuất
khẩu bằng 0 và độ dốc âm có ý nghĩa tại biên phải khi cường độ bằng 1;
điểm uốn được khôi phục bằng phương pháp delta cùng khoảng tin cậy 95
phần trăm theo gradient của phương pháp delta, cho ra điểm uốn 49,4\%
với khoảng tin cậy từ 43,1\% đến 55,7\% năm 2012 và 47,2\% với khoảng từ
40,8\% đến 53,6\% năm 2024, hai khoảng chồng lấn rộng; trên mẫu gộp điểm
uốn là 48,8\% với khoảng tin cậy từ 44,2\% đến 53,4\%. Ở khung đa quốc
gia, ba điều kiện của Lind và Mehlum, gồm dấu hệ số, điểm uốn nằm chặt
bên trong miền quan sát, và dấu độ dốc tại hai biên, được gộp thành một
kiểm định giao điểm duy nhất. Trường hợp đảo nhỏ Thái Bình Dương minh
họa giá trị thông tin của một kiểm định không bác bỏ: cả hai số hạng mất
ý nghĩa khi đưa vào đồng thời do cộng tuyến cao trong mẫu cường độ thấp
với trung bình cường độ chỉ 0,048, và hệ số bậc hai dương hàm ý dạng chữ
U xuôi thay vì chữ U ngược, một sự đảo dạng bản thân nó nhất quán với
quan hệ đơn điệu âm không có điểm quay nội miền.

\hypertarget{muxf4-huxecnh-ux111iux1ec1u-tiux1ebft-h2h6}{%
\subsection{3.4.4 Mô hình điều tiết
(H2--H6)}\label{muxf4-huxecnh-ux111iux1ec1u-tiux1ebft-h2h6}}

\hypertarget{tux1b0ux1a1ng-tuxe1c-hai-chiux1ec1u-h2-h3}{%
\subsubsection{Tương tác hai chiều (H2,
H3)}\label{tux1b0ux1a1ng-tuxe1c-hai-chiux1ec1u-h2-h3}}

\[\ln(LP)_i = \beta_0 + \beta_1 FSTS_i + \beta_2 FSTS_i^2 + \beta_3 M_i + \beta_4 (FSTS_i \times M_i) + \beta_5 (FSTS_i^2 \times M_i) + \boldsymbol{\gamma} \mathbf{X}_i + \varepsilon_i\]

với \(M_i\) là biến điều tiết (TCI hoặc DAI).

\hypertarget{tux1b0ux1a1ng-tuxe1c-ba-chiux1ec1u-synthesis}{%
\subsubsection{Tương tác ba chiều
(synthesis)}\label{tux1b0ux1a1ng-tuxe1c-ba-chiux1ec1u-synthesis}}

\[\ln(LP)_i = \beta_0 + \beta_1 FSTS_i + \beta_2 FSTS_i^2 + \beta_3 D_i + \beta_4 I_i + \beta_5 M_i + \sum \text{interactions} + \boldsymbol{\gamma} \mathbf{X}_i + \varepsilon_i\]

với \(D\) = digital, \(I\) = institutional, \(M\) = manager.

\textbf{Cơ sở kế thừa}: chuẩn điều tiết trong quản trị chiến lược (Aiken
\& West, 1991; Dawson, 2014).

\textbf{Đóng góp mới}: three-way điều tiết digital × institutional ×
manager trong bối cảnh châu Á và Pacific với 50 nước chưa được kiểm định
trước đây.

\hypertarget{muxf4-huxecnh-dux1ecb-biux1ec7t-theo-thux1eddi-gian-h6}{%
\subsection{3.4.5 Mô hình dị biệt theo thời gian
(H6)}\label{muxf4-huxecnh-dux1ecb-biux1ec7t-theo-thux1eddi-gian-h6}}

\[\ln(LP)_i = \beta_0 + \beta_1 FSTS_i + \beta_2 FSTS_i^2 + \beta_3 Y_{2024,i} + \beta_4 (FSTS_i \times Y_{2024,i}) + \beta_5 (FSTS_i^2 \times Y_{2024,i}) + \boldsymbol{\gamma} \mathbf{X}_i + \varepsilon_i\]

Nếu \(\beta_4\) hoặc \(\beta_5\) có ý nghĩa, hình dạng đường quan hệ đã
dịch chuyển giữa hai mốc thời gian.

\textbf{Cơ sở kế thừa}: Wu et al.~(2022) với tiến hóa hai mươi năm; Xiao
et al.~(2013) với các điều kiện đặc thù của Trung Quốc; nghiên cứu thành
phần P2 làm cơ sở dạng bậc hai năm 2012.

\textbf{Đóng góp mới}: áp dụng cho Trung Quốc các năm 2012 và 2024 để
kiểm định tính ổn định của điểm uốn dạng bậc hai (khoảng 47,8\% FSTS)
sau hơn một thập kỷ chuyển đổi số.

\hypertarget{kiux1ec3m-ux111ux1ecbnh-ux1ed5n-ux111ux1ecbnh-hux1ec7-sux1ed1-xuyuxean-ux111ux1ee3t-khux1ea3o-suxe1t-theo-paternoster}{%
\subsection{3.4.5.0 Kiểm định ổn định hệ số xuyên đợt khảo sát theo
Paternoster}\label{kiux1ec3m-ux111ux1ecbnh-ux1ed5n-ux111ux1ecbnh-hux1ec7-sux1ed1-xuyuxean-ux111ux1ee3t-khux1ea3o-suxe1t-theo-paternoster}}

Để đánh giá xem hình dạng quan hệ quốc tế hóa--hiệu quả có dịch chuyển
giữa các đợt khảo sát hay không, luận án không chỉ dựa vào biến tương
tác giữa cường độ xuất khẩu và biến giả đợt mà còn áp dụng kiểm định z
của Paternoster và cộng sự (1998) cho tính bằng nhau của hệ số hồi quy
giữa hai mẫu độc lập. Kiểm định này phù hợp với cấu trúc dữ liệu lát cắt
ngang lặp lại, nơi mỗi đợt khảo sát rút một mẫu doanh nghiệp độc lập chứ
không theo dõi cùng doanh nghiệp qua thời gian, khiến so sánh hệ số giữa
các đợt là so sánh giữa hai ước lượng không tương quan. Thống kê kiểm
định nhận dạng chênh lệch hai hệ số chia cho căn bậc hai của tổng phương
sai hai ước lượng, cụ thể
\(z = (\hat\beta_A - \hat\beta_B) / \sqrt{\text{SE}(\hat\beta_A)^2 + \text{SE}(\hat\beta_B)^2}\),
với giá trị p hai phía từ phân phối chuẩn tiêu chuẩn.

Kiểm định được áp dụng riêng cho từng cặp hệ số trọng tâm, gồm số hạng
bậc nhất, số hạng bậc hai và các hệ số trực tiếp của năng lực công nghệ
và chấp nhận số. Việc tách kiểm định theo từng tham số cho phép phân
biệt hai loại dịch chuyển có ý nghĩa lý thuyết khác nhau: dịch chuyển ở
tham số độ cong báo hiệu hình dạng quan hệ thay đổi cấu trúc, còn dịch
chuyển ở hệ số trực tiếp báo hiệu vai trò của một chiều năng lực biến
đổi theo vòng đời số hóa mà không làm đổi dạng đường cong. Khung đa đợt
của luận án vì vậy có thể kết luận, chẳng hạn, rằng ngưỡng tối ưu bền
vững cấu trúc trong khi đóng góp trực tiếp của chấp nhận số lại biến
thiên rõ rệt giữa các đợt. Ở nghiên cứu có sẵn lõi panel nhỏ gồm các
doanh nghiệp xuất hiện ở cả hai đợt, kiểm định Paternoster được bổ trợ
bằng sai số chuẩn phân cụm theo định danh doanh nghiệp nhằm xử lý tương
quan nội nhóm.

Thủ tục này cho ra hai kết cục đối lập về mặt lý thuyết tùy bối cảnh thể
chế. Ở Trung Quốc, kiểm định không bác bỏ bình đẳng hệ số giữa hai đợt
cho cả số hạng bậc nhất và bậc hai, với thống kê dịch chuyển bậc nhất
bằng +0,82 (giá trị p bằng 0,412) và dịch chuyển bậc hai bằng −0,61 (giá
trị p bằng 0,545), hỗ trợ cách diễn giải rằng ngưỡng chữ U ngược là một
quy luật cấu trúc bền vững với điểm uốn ổn định quanh 48\% bất kể thập
niên. Tính bền vững này còn vững trước giới hạn mẫu: trên mẫu con chế
biến chế tạo, điểm uốn là 48,1\% năm 2012 và 46,4\% năm 2024 với các
thống kê Paternoster vẫn không có ý nghĩa. Ngược lại, ở Ấn Độ, kiểm định
Paternoster xuyên đợt bác bỏ bình đẳng hệ số ở mức rất nhỏ hơn 0,0001
cho cả số hạng tuyến tính và bậc hai, một sự sụp đổ trùng với thập niên
thay đổi thể chế ngoại lệ gồm bãi bỏ tiền mặt, thuế hàng hóa và dịch vụ,
luật phá sản, hạ tầng thanh toán hợp nhất, chương trình khuyến khích sản
xuất và đại dịch. Sự đối lập giữa bền vững ở Trung Quốc và sụp đổ ở Ấn
Độ chính là giá trị nhận dạng của kiểm định xuyên đợt: nó phân biệt được
dịch chuyển cấu trúc thực do cú sốc thể chế khỏi biến thiên ngẫu nhiên
của mẫu.

\hypertarget{muxf4-huxecnh-cux1ee5-thux1ec3-nghiuxean-cux1ee9u-3-viux1ec7t-nam-wbes-200920152023}{%
\subsection{3.4.5.1 Mô hình cụ thể: Nghiên cứu 3 (Việt Nam, WBES
2009/2015/2023)}\label{muxf4-huxecnh-cux1ee5-thux1ec3-nghiuxean-cux1ee9u-3-viux1ec7t-nam-wbes-200920152023}}

Nghiên cứu 3 sử dụng chuỗi mô hình lồng nhau M0--M8 ước lượng riêng theo
từng sóng khảo sát (2009, 2015, 2023) và trên mẫu gộp. Ký hiệu tiếng
Việt:

\begin{itemize}
\tightlist
\item
  \textbf{lnNSLD\_it} = lnLP\_it: log năng suất lao động (ln doanh thu
  PPP / lao động thường xuyên)
\item
  \textbf{CDDXK\_c\_it} = FSTS\_c\_it: cường độ xuất khẩu điều chỉnh
  trung bình trong sóng (mean-centred)
\item
  \textbf{CDDXK\_c²\_it} = FSTS\_c²\_it: bình phương cường độ xuất khẩu
  điều chỉnh
\item
  \textbf{NLCN\_z\_it} = TCI\_z\_it: năng lực công nghệ (chuẩn hoá z
  trong sóng)
\item
  \textbf{CSS\_z\_it} = DAI\_z\_it: chỉ số số hoá cơ sở, Tầng 1 (chuẩn
  hoá z trong sóng)
\item
  \textbf{lnLD, TuoiDN, SoHuuNuocNgoai}: kiểm soát quy mô, tuổi, sở hữu
\item
  \textbf{δ\_s}: hiệu ứng cố định ngành (1-digit ISIC); \textbf{λ\_t}:
  hiệu ứng cố định sóng (chỉ pooled)
\end{itemize}

\textbf{M0, Mô hình cơ sở (biến kiểm soát):}
\[\ln NSLD_{it} = \alpha + \gamma_1 \ln LD_{it} + \gamma_2 TuoiDN_{it} + \gamma_3 SoHuuNN_{it} + \delta_s + [\lambda_t] + \varepsilon_{it}\]

\textbf{M1, Quốc tế hoá tuyến tính:}
\[\ln NSLD_{it} = \alpha + \beta_1 CDDXK\_c_{it} + \boldsymbol{\gamma}\mathbf{X}_{it} + \delta_s + [\lambda_t] + \varepsilon_{it}\]

\textbf{M2, Quốc tế hoá phi tuyến (kiểm định H1, chữ U ngược):}
\[\ln NSLD_{it} = \alpha + \beta_1 CDDXK\_c_{it} + \beta_2 CDDXK\_c^2_{it} + \boldsymbol{\gamma}\mathbf{X}_{it} + \delta_s + [\lambda_t] + \varepsilon_{it}\]

H1: \(\beta_1 > 0\) và \(\beta_2 < 0\); điểm quay
\(TP^* = -\beta_1 / (2\beta_2)\), xác nhận bởi kiểm định Lind--Mehlum
(2010).

\textbf{M3, Điều tiết NLCN (kiểm định H2):} \[\begin{aligned}
\ln NSLD_{it} = {} & \alpha + \beta_1 CDDXK\_c + \beta_2 CDDXK\_c^2 + \beta_3 NLCN\_z \\
& + \beta_4(CDDXK\_c \times NLCN\_z) + \beta_5(CDDXK\_c^2 \times NLCN\_z) \\
& + \boldsymbol{\gamma}\mathbf{X} + \delta_s + [\lambda_t] + \varepsilon
\end{aligned}\]

\textbf{M4, Điều tiết CSS/DAI (H3 thăm dò):} \[\begin{aligned}
\ln NSLD_{it} = {} & \alpha + \beta_1 CDDXK\_c + \beta_2 CDDXK\_c^2 + \beta_3 CSS\_z \\
& + \beta_4(CDDXK\_c \times CSS\_z) + \beta_5(CDDXK\_c^2 \times CSS\_z) \\
& + \boldsymbol{\gamma}\mathbf{X} + \delta_s + [\lambda_t] + \varepsilon
\end{aligned}\]

\textbf{M5, NLCN trực tiếp (không tương tác):}
\[\ln NSLD_{it} = \alpha + \beta_1 CDDXK\_c + \beta_2 CDDXK\_c^2 + \beta_3 NLCN\_z + \boldsymbol{\gamma}\mathbf{X} + \delta_s + [\lambda_t] + \varepsilon\]

\textbf{M6, CSS/DAI trực tiếp (không tương tác):}
\[\ln NSLD_{it} = \alpha + \beta_1 CDDXK\_c + \beta_2 CDDXK\_c^2 + \beta_3 CSS\_z + \boldsymbol{\gamma}\mathbf{X} + \delta_s + [\lambda_t] + \varepsilon\]

\textbf{M7, Cả hai trực tiếp, không tương tác:}
\[\ln NSLD_{it} = \alpha + \beta_1 CDDXK\_c + \beta_2 CDDXK\_c^2 + \beta_3 NLCN\_z + \beta_4 CSS\_z + \boldsymbol{\gamma}\mathbf{X} + \delta_s + [\lambda_t] + \varepsilon\]

\textbf{M8, Mô hình đầy đủ (trực tiếp + điều tiết CSS):}
\[\begin{aligned}
\ln NSLD_{it} = {} & \alpha + \beta_1 CDDXK\_c + \beta_2 CDDXK\_c^2 + \beta_3 NLCN\_z + \beta_4 CSS\_z \\
& + \beta_5(CDDXK\_c \times CSS\_z) + \beta_6(CDDXK\_c^2 \times CSS\_z) \\
& + \boldsymbol{\gamma}\mathbf{X} + \delta_s + [\lambda_t] + \varepsilon
\end{aligned}\]

\emph{Tất cả mô hình đều ước lượng bằng OLS với sai số chuẩn vững HC1.
Hiệu ứng cố định sóng \(\lambda_t\) chỉ áp dụng trong mô hình gộp. Mô
hình được ước lượng riêng cho 3 sóng và gộp (N = 989 / 956 / 1.013 /
2.958).}

\textbf{Bảng định nghĩa biến, Nghiên cứu 3 (Việt Nam):}

\begin{longtable}[]{@{}
  >{\raggedright\arraybackslash}p{(\columnwidth - 6\tabcolsep) * \real{0.2500}}
  >{\raggedright\arraybackslash}p{(\columnwidth - 6\tabcolsep) * \real{0.2500}}
  >{\raggedright\arraybackslash}p{(\columnwidth - 6\tabcolsep) * \real{0.2500}}
  >{\raggedright\arraybackslash}p{(\columnwidth - 6\tabcolsep) * \real{0.2500}}@{}}
\toprule\noalign{}
\begin{minipage}[b]{\linewidth}\raggedright
Ký hiệu
\end{minipage} & \begin{minipage}[b]{\linewidth}\raggedright
Biến WBES
\end{minipage} & \begin{minipage}[b]{\linewidth}\raggedright
Cách tính
\end{minipage} & \begin{minipage}[b]{\linewidth}\raggedright
Vai trò
\end{minipage} \\
\midrule\noalign{}
\endhead
\bottomrule\noalign{}
\endlastfoot
lnNSLD & d2, l1 & ln(d2 / l1): ln(doanh thu PPP / lao động thường xuyên)
& Biến phụ thuộc \\
CDDXK & d3c & d3c / 100: cường độ xuất khẩu trực tiếp (0--1) & Biến độc
lập \\
CDDXK\_c & d3c & CDDXK − trung bình sóng: chuẩn hoá điều chỉnh & Biến
độc lập (centred) \\
CDDXK\_c² & d3c & CDDXK\_c bình phương: hạng phi tuyến & Kiểm định
inverted-U (H1) \\
NLCN\_z & b8, e6 & chuẩn hóa z trong sóng của TB(b8₀₁, e6₀₁): chứng chỉ
chất lượng + công nghệ ngoại & Năng lực công nghệ (H2) \\
CSS\_z & c22b & chuẩn hóa z trong sóng của c22b₀₁: hiện diện website &
Số hoá Tầng 1 (H3 thăm dò) \\
lnLD & l1 & ln(l1): log số lao động thường xuyên & Kiểm soát: quy mô
doanh nghiệp \\
TuoiDN & b5 & năm khảo sát − b5: số năm hoạt động & Kiểm soát: tuổi
doanh nghiệp \\
SoHuuNN & b2b & 1 nếu b2b \textgreater{} 0: có vốn nước ngoài & Kiểm
soát: hình thức sở hữu \\
δ\_s & a4b / a4a & ngành 1-chữ số ISIC (a4b cho 2009/2015; a4a cho 2023)
& Hiệu ứng cố định ngành \\
λ\_t & đợt & chỉ báo sóng khảo sát (chỉ pooled) & Hiệu ứng cố định thời
kỳ \\
\end{longtable}

\hypertarget{muxf4-huxecnh-cux1ee5-thux1ec3-nghiuxean-cux1ee9u-4-singapore-wbes-2023}{%
\subsection{3.4.5.2 Mô hình cụ thể: Nghiên cứu 4 (Singapore, WBES
2023)}\label{muxf4-huxecnh-cux1ee5-thux1ec3-nghiuxean-cux1ee9u-4-singapore-wbes-2023}}

Nghiên cứu 4 sử dụng dữ liệu mặt cắt ngang WBES Singapore 2023 (N =
623). Thiết kế đơn sóng, không có thành phần λ\_t. Ký hiệu tiếng Việt:

\begin{itemize}
\tightlist
\item
  \textbf{lnNSLD\_i} = lnLP\_i: log năng suất lao động (ln doanh thu PPP
  / lao động thường xuyên)
\item
  \textbf{CDDXK\_i} = FSTS\_i: cường độ xuất khẩu trực tiếp (d3c / 100)
\item
  \textbf{CDDXK\_c\_i} = FSTS\_c\_i: CDDXK trung bình mẫu; CDDXK\_c² là
  bình phương
\item
  \textbf{NLCN\_z\_i} = TCI\_z\_i: năng lực công nghệ, chuẩn hóa z của
  trung bình(b8₀₁, e6₀₁)
\item
  \textbf{CSS\_z\_i} = DAI\_z\_i: chỉ số số hoá \textbf{Tầng 1 và 2},
  chuẩn hóa z của trung bình(c22b₀₁, k33₀₁, k38₀₁); khác P3 ở chỗ bao
  gồm thanh toán điện tử hai chiều (Tầng 2)
\item
  \textbf{lnLD\_i}, \textbf{TuoiDN\_i}, \textbf{SoHuuNN\_i}: biến kiểm
  soát tương tự P3
\item
  \textbf{δ\_s}: ngành hiệu ứng cố định (ISIC 1-chữ số)
\item
  \textbf{IMR\_i} = nghịch đảo tỷ lệ Mills từ mô hình probit tham gia
  xuất khẩu (kiểm tra độ nhạy Heckman)
\end{itemize}

\textbf{Chuỗi mô hình lồng nhau M0--M5:}

\begin{quote}
M0 (Baseline): lnNSLD\_i = α + γ₁ lnLD\_i + γ₂ TuoiDN\_i + γ₃ SoHuuNN\_i
+ δ\_s + ε\_i
\end{quote}

\begin{quote}
M1 (Tuyến tính FSTS): lnNSLD\_i = α + β₁ CDDXK\_c\_i + γ·X\_i + δ\_s +
ε\_i
\end{quote}

\begin{quote}
M2 (Bậc hai FSTS, kiểm định H1): lnNSLD\_i = α + β₁ CDDXK\_c\_i + β₂
CDDXK\_c²\_i + γ·X\_i + δ\_s + ε\_i H1: β₁ \textgreater{} 0, β₂
\textless{} 0; TP* = −β₁/(2β₂) ≈ 88,6\% FSTS {[}CI bootstrap {[}53\%,
253\%{]}{]} Lưu ý: Lind--Mehlum p = 0,303, không bác bỏ tuyến tính trong
phạm vi dữ liệu quan sát; đây là kết quả thông tin tích cực theo khung
bão hòa (saturation framework)
\end{quote}

\begin{quote}
M3 (+ NLCN, H1-TCI trực tiếp): lnNSLD\_i = α + β₁ CDDXK\_c + β₂
CDDXK\_c² + β₃ NLCN\_z + γ·X + δ\_s + ε
\end{quote}

\begin{quote}
M4 (+ CSS trực tiếp): lnNSLD\_i = α + β₁ CDDXK\_c + β₂ CDDXK\_c² + β₃
NLCN\_z + β₄ CSS\_z + γ·X + δ\_s + ε
\end{quote}

\begin{quote}
M5 (Mô hình đầy đủ, kiểm định H3): lnNSLD\_i = α + β₁ CDDXK\_c + β₂
CDDXK\_c² + β₃ NLCN\_z + β₄ CSS\_z + β₅(CDDXK\_c × CSS\_z) +
β₆(CDDXK\_c² × CSS\_z) + γ·X + δ\_s + ε H3: β₆ \textgreater{} 0, DAI
khuếch đại lợi nhuận quốc tế hóa và hiệu quả ở cường độ xuất khẩu cao
(coordination platform mechanism; Stallkamp \& Schotter, 2021) Kết quả:
β₆ = +3,119 (p = 0,005) trong mẫu đầy đủ; β₆ = +2,821 (p = 0,003,
F-test) trong mẫu chỉ xuất khẩu (N = 84, lưu ý: công suất thống kê ≈
16\%)
\end{quote}

\textbf{Bảng định nghĩa biến, Nghiên cứu 4 (Singapore):}

\begin{longtable}[]{@{}
  >{\raggedright\arraybackslash}p{(\columnwidth - 6\tabcolsep) * \real{0.2500}}
  >{\raggedright\arraybackslash}p{(\columnwidth - 6\tabcolsep) * \real{0.2500}}
  >{\raggedright\arraybackslash}p{(\columnwidth - 6\tabcolsep) * \real{0.2500}}
  >{\raggedright\arraybackslash}p{(\columnwidth - 6\tabcolsep) * \real{0.2500}}@{}}
\toprule\noalign{}
\begin{minipage}[b]{\linewidth}\raggedright
Ký hiệu VN
\end{minipage} & \begin{minipage}[b]{\linewidth}\raggedright
Mã WBES
\end{minipage} & \begin{minipage}[b]{\linewidth}\raggedright
Cách tính
\end{minipage} & \begin{minipage}[b]{\linewidth}\raggedright
Vai trò
\end{minipage} \\
\midrule\noalign{}
\endhead
\bottomrule\noalign{}
\endlastfoot
lnNSLD & d2, l1 & ln(d2 / l1) & Biến phụ thuộc \\
CDDXK\_c & d3c & FSTS − mean(FSTS): centred & Biến độc lập \\
CDDXK\_c² & d3c & CDDXK\_c bình phương & Phi tuyến (H1) \\
NLCN\_z & b8, e6 & chuẩn hóa z của TB(b8₀₁, e6₀₁) & Năng lực công nghệ
(H2) \\
CSS\_z & c22b, k33, k38 & chuẩn hóa z của TB(c22b₀₁, k33₀₁, k38₀₁):
\textbf{Tầng 1 và 2} & Số hoá Tầng 1 và 2 (H3) \\
lnLD & l1 & ln(l1) & Quy mô doanh nghiệp \\
TuoiDN & b5 & năm khảo sát − b5 & Tuổi doanh nghiệp \\
SoHuuNN & b2b & 1 nếu b2b \textgreater{} 0 & Hình thức sở hữu \\
IMR & probit selection & Nghịch đảo tỷ lệ Mills (kiểm tra Heckman) &
Kiểm soát chọn lựa tham gia xuất khẩu \\
δ\_s & a4b & ngành FE (ISIC 1-chữ số) & Hiệu ứng cố định ngành \\
\end{longtable}

\begin{center}\rule{0.5\linewidth}{0.5pt}\end{center}

\hypertarget{muxf4-huxecnh-cux1ee5-thux1ec3-nghiuxean-cux1ee9u-5-trung-quux1ed1c-wbes-2012-vuxe0-2024}{%
\subsection{3.4.5.3 Mô hình cụ thể: Nghiên cứu 5 (Trung Quốc, WBES 2012
và
2024)}\label{muxf4-huxecnh-cux1ee5-thux1ec3-nghiuxean-cux1ee9u-5-trung-quux1ed1c-wbes-2012-vuxe0-2024}}

Nghiên cứu 5 sử dụng hai sóng WBES Trung Quốc (2012: N = 2.610; 2024: N
= 1.934; pooled N = 4.544). Thiết kế đa sóng không panel (217 doanh
nghiệp xuất hiện cả hai sóng tạo thành ``panel core'' được sử dụng cho
cluster-robust SE). Ký hiệu tiếng Việt:

\begin{itemize}
\tightlist
\item
  \textbf{lnNSLD\_it} = lnLP\_it: log năng suất lao động (biến phụ
  thuộc)
\item
  \textbf{CDDXK\_c\_it} = FSTS\_c\_it: cường độ xuất khẩu centred theo
  trung bình sóng
\item
  \textbf{NLCN\_z\_it} = TCI\_full\_z\_it: năng lực công nghệ toàn diện,
  chuẩn hóa z của TB(b8₀₁, e6₀₁, h1₀₁, h8₀₁); mở rộng hơn P3
\item
  \textbf{CSS\_z\_it} = DAI\_core\_it: chỉ số số hoá cơ bản, chỉ c22b₀₁
  (Tầng 1 mỏng, một chỉ báo nhị phân)
\item
  \textbf{lnLD\_it}, \textbf{TuoiDN\_it}, \textbf{SoHuuNN\_it}: biến
  kiểm soát
\item
  \textbf{δ\_s}: ngành hiệu ứng cố định; \textbf{λ\_t}: đợt hiệu ứng cố
  định (pooled)
\end{itemize}

\textbf{Chuỗi mô hình lồng nhau M0--M6:}

\begin{quote}
M0 (Baseline): lnNSLD\_it = α + γ₁ lnLD\_it + γ₂ TuoiDN\_it + γ₃
SoHuuNN\_it + δ\_s + λ\_t + ε\_it
\end{quote}

\begin{quote}
M1 (Tuyến tính FSTS): lnNSLD\_it = α + β₁ CDDXK\_c\_it + γ·X\_it + δ\_s
+ λ\_t + ε\_it
\end{quote}

\begin{quote}
M2 (Bậc hai FSTS, kiểm định H1): lnNSLD\_it = α + β₁ CDDXK\_c\_it + β₂
CDDXK\_c²\_it + γ·X\_it + δ\_s + λ\_t + ε\_it H1: β₁ \textgreater{} 0,
β₂ \textless{} 0; TP* = −β₁/(2β₂), điểm uốn 49,4\% (2012), 47,2\%
(2024), và 48,8\% (pooled) Kiểm định ổn định cấu trúc: Paternoster et
al.~(1998) z-test: z(FSTS) = +0,82 (p = 0,412); z(FSTS²) = −0,61 (p =
0,545), không bác bỏ bình đẳng hệ số giữa hai sóng
\end{quote}

\begin{quote}
M3 (+ NLCN trực tiếp, H2 level-shift): lnNSLD\_it = α + β₁ CDDXK\_c + β₂
CDDXK\_c² + β₃ NLCN\_z + γ·X + δ\_s + λ\_t + ε Kết quả: β₃ = +0,28
(2012, p \textless{} 0,001), +0,43 (2024, p \textless{} 0,001), TCI là
``bộ tăng mức'' (level-shifter), không điều tiết độ cong
\end{quote}

\begin{quote}
M4 (+ CSS trực tiếp): lnNSLD\_it = α + β₁ CDDXK\_c + β₂ CDDXK\_c² + β₃
NLCN\_z + β₄ CSS\_z + γ·X + δ\_s + λ\_t + ε
\end{quote}

\begin{quote}
M5 (Kiểm định ổn định xuyên sóng, H2b): Ước lượng M2 riêng biệt cho 2012
và 2024, sau đó áp dụng: z = (β̂₁,2012 − β̂₁,2024) / √(SE²₁,2012 +
SE²₁,2024) (Paternoster et al., 1998), tương tự cho β₂ (FSTS²)
\end{quote}

\begin{quote}
M6 (Điều tiết ba chiều, kiểm định H3/H4a/H4b): lnNSLD\_it = α + β₁
CDDXK\_c + β₂ CDDXK\_c² + β₃ NLCN\_z + β₄ CSS\_z + β₅(CDDXK\_c ×
NLCN\_z) + β₆(CDDXK\_c² × NLCN\_z) + β₇(CDDXK\_c × wave) + β₈(CDDXK\_c²
× wave) + γ·X + δ\_s + λ\_t + ε F-tests: F1 (dịch chuyển độ cong xuyên
sóng), F2 (điều tiết năng lực), F3 (dịch chuyển điều kiện) Kết quả: F2 =
3,26 (p = 0,039, không qua Bonferroni α* = 0,017), do đó H4b (không có
điều tiết độ cong theo năng lực) được chấp nhận
\end{quote}

\textbf{Bảng định nghĩa biến, Nghiên cứu 5 (Trung Quốc):}

\begin{longtable}[]{@{}
  >{\raggedright\arraybackslash}p{(\columnwidth - 6\tabcolsep) * \real{0.2500}}
  >{\raggedright\arraybackslash}p{(\columnwidth - 6\tabcolsep) * \real{0.2500}}
  >{\raggedright\arraybackslash}p{(\columnwidth - 6\tabcolsep) * \real{0.2500}}
  >{\raggedright\arraybackslash}p{(\columnwidth - 6\tabcolsep) * \real{0.2500}}@{}}
\toprule\noalign{}
\begin{minipage}[b]{\linewidth}\raggedright
Ký hiệu VN
\end{minipage} & \begin{minipage}[b]{\linewidth}\raggedright
Mã WBES
\end{minipage} & \begin{minipage}[b]{\linewidth}\raggedright
Cách tính
\end{minipage} & \begin{minipage}[b]{\linewidth}\raggedright
Vai trò
\end{minipage} \\
\midrule\noalign{}
\endhead
\bottomrule\noalign{}
\endlastfoot
lnNSLD & d2, l1 & ln(d2 / l1) & Biến phụ thuộc \\
CDDXK\_c & d3c & FSTS − mean\_wave(FSTS) & Biến độc lập \\
CDDXK\_c² & d3c & CDDXK\_c bình phương & Phi tuyến (H1) \\
NLCN\_z & b8, e6, h1, h8 & chuẩn hóa z của TB(b8₀₁, e6₀₁, h1₀₁, h8₀₁):
TCI toàn diện & Năng lực công nghệ (H2) \\
CSS\_z & c22b & c22b₀₁: \textbf{Tầng 1 đơn biến} (binary) & Số hoá Tầng
1 mỏng (kiểm soát) \\
lnLD & l1 & ln(l1) & Quy mô doanh nghiệp \\
TuoiDN & b5 & năm khảo sát − b5 & Tuổi doanh nghiệp \\
SoHuuNN & b2b & 1 nếu b2b \textgreater{} 0 & Hình thức sở hữu \\
δ\_s & ISIC ngành & ngành FE & Hiệu ứng cố định ngành \\
λ\_t & đợt (2012/2024) & đợt FE (chỉ pooled) & Hiệu ứng cố định thời
kỳ \\
\end{longtable}

\emph{Ghi chú: 217 doanh nghiệp xuất hiện cả hai sóng (``panel core'')
tạo thêm nhận dạng biến thiên trong mẫu trong khi clustered SE theo
idstd xử lý tương quan nội nhóm. DAI Trung Quốc chỉ là Tầng 1 đơn biến,
không đủ để kiểm định CDCM; hàm ý chính sách và lý thuyết khác với P4
Singapore (Tầng 1 và 2).}

\begin{center}\rule{0.5\linewidth}{0.5pt}\end{center}

\hypertarget{muxf4-huxecnh-cux1ee5-thux1ec3-nghiuxean-cux1ee9u-7-ux111a-quux1ed1c-gia-chuxe2u-uxe1-wbes-20062026}{%
\subsection{3.4.5.4 Mô hình cụ thể: Nghiên cứu 7 (Đa quốc gia châu Á,
WBES
2006--2026)}\label{muxf4-huxecnh-cux1ee5-thux1ec3-nghiuxean-cux1ee9u-7-ux111a-quux1ed1c-gia-chuxe2u-uxe1-wbes-20062026}}

Nghiên cứu 7 là kiểm định quy mô lớn nhất của luận án, sử dụng dữ liệu
WBES từ \textbf{50 nền kinh tế châu Á và Thái Bình Dương} (103 cặp nền
kinh tế × năm, 2006--2026, gồm Nhật Bản 2025). Mẫu phân tích đạt N =
81.022 (M2) đến N = 79.080 (M5 với kiểm soát đầy đủ). Thiết kế hồi quy
hiệu ứng cố định hai chiều (nền kinh tế và năm) với sai số chuẩn cụm
theo nền kinh tế. Ký hiệu tiếng Việt:

\begin{itemize}
\tightlist
\item
  \textbf{lnNSLD\_it}: ln(năng suất lao động) = ln(doanh thu / lao động
  thường trực), biến phụ thuộc
\item
  \textbf{CDDXK\_c\_it}: cường độ xuất khẩu mean-centred, FSTS\_c =
  (d3c/100) − mean\_pool(FSTS)
\item
  \textbf{CDDXK\_c²\_it}: CDDXK\_c bình phương, kiểm định phi tuyến
\item
  \textbf{NLCN\_z\_it}: năng lực công nghệ z-standardised, TCI\_z =
  chuẩn hóa z(b8, e6)
\item
  \textbf{CSS\_z\_it}: chỉ số số hoá z-standardised, DAI\_z = chuẩn hóa
  z(website + e-pay, Tầng 1 và 2: c22b/e1, k33)
\item
  \textbf{KNQLy\_it}: kinh nghiệm nhà quản lý trong ngành (năm, b7)
\item
  \textbf{NQL\_nu\_it}: top manager là nữ (binary, b7a/b6a)
\item
  \textbf{ICRV\_j}: phân loại chế độ thể chế ICRV (integer 1--6,
  Advanced sang SIDS)
\item
  \textbf{lnLD\_it}: ln(lao động thường trực), kiểm soát quy mô doanh
  nghiệp
\item
  \textbf{TuoiDN\_it}: tuổi doanh nghiệp (năm\_khảo\_sát − b5)
\item
  \textbf{SoHuuNN\_it}: tỷ lệ sở hữu nước ngoài (b6a/100)
\item
  \textbf{δ\_ct}: country × year hiệu ứng cố định (M5)
\end{itemize}

\textbf{Chuỗi mô hình lồng nhau M0--M11:}

\textbf{M0} (tuyến tính baseline):
\[\ln NSLD_{it} = \alpha + \beta_1 CDDXK\_c_{it} + \varepsilon_{it}\]

\textbf{M2} (phi tuyến, kiểm định H1):
\[\ln NSLD_{it} = \alpha + \beta_1 CDDXK\_c_{it} + \beta_2 CDDXK\_c^2_{it} + \varepsilon_{it}\]
\[TP^* = -\beta_1 / (2\beta_2); \quad \text{Lind–Mehlum } p < {,}001;\ TP = 51{,}5\%\ (N = 81.022)\]

\textbf{M3} (M2 + kiểm soát cơ bản):
\[\ln NSLD_{it} = \alpha + \beta_1 CDDXK\_c + \beta_2 CDDXK\_c^2 + \gamma \cdot X_{it} + \varepsilon_{it}\]

\textbf{M5} (M3 + country-year FE, kiểm định vững mạnh nhất):
\[\ln NSLD_{it} = \alpha + \beta_1 CDDXK\_c + \beta_2 CDDXK\_c^2 + \gamma \cdot X_{it} + \delta_{ct} + \varepsilon_{it}\]
\[TP = 43{,}6\%\ (N = 79.080;\ \text{Lind–Mehlum } p < {,}001)\]

\textbf{M7} (M3 + điều tiết NLCN, kiểm định H2):
\[\ln NSLD_{it} = \alpha + \beta_1 CDDXK\_c + \beta_2 CDDXK\_c^2 + \beta_3 NLCN\_z\]
\[+ \beta_4 (CDDXK\_c \times NLCN\_z) + \beta_5 (CDDXK\_c^2 \times NLCN\_z) + \gamma \cdot X + \varepsilon\]

\textbf{M8} (M7 + điều tiết CSS, kiểm định H3):
\[\ln NSLD_{it} = \alpha + \beta_1 CDDXK\_c + \beta_2 CDDXK\_c^2 + \beta_3 NLCN\_z + \beta_6 CSS\_z\]
\[+ \beta_7 (CDDXK\_c \times CSS\_z) + \beta_8 (CDDXK\_c^2 \times CSS\_z) + \gamma \cdot X + \varepsilon\]
\[\begin{aligned}
\text{H3: } & \hat\beta_7 = -0{,}271\ (p = {,}149);\ \hat\beta_8 = +0{,}252\ (p = {,}367), \\
& \text{cùng chiều nén nhưng không ý nghĩa toàn mẫu; hiệu ứng mức DAI} = +0{,}201\ (p < {,}001)
\end{aligned}\]

\textbf{M9} (M8 + đặc điểm nhà quản lý, kiểm định H4):
\[+ \beta_9 KNQLy + \beta_{10} NQL\_nu + \beta_{11}(CDDXK\_c \times KNQLy) + \gamma \cdot X + \varepsilon\]
\[\hat\beta_{10} = -0{,}088^{*} \ (p = {,}026) \text{ (nữ quản lý cấp cao: hệ số mức âm nhẹ; tương tác với FSTS không ý nghĩa)}\]

\textbf{M10} (M3 + điều tiết ICRV, kiểm định H5):
\[\ln NSLD_{it} = \alpha + \beta_1 CDDXK\_c + \beta_2 CDDXK\_c^2 + \beta_{11} ICRV_j\]
\[+ \beta_{12}(CDDXK\_c \times ICRV) + \beta_{13}(CDDXK\_c^2 \times ICRV) + \gamma \cdot X + \varepsilon\]
\[\begin{aligned}
\text{H5 (ba vùng): } & \text{chữ U ngược định hình rõ ở Nhóm IV } (TP = 43{,}0\%,\ p_{LM} < {,}001); \\
& \text{gần tuyến tính ở Nhóm I; tan rã ở Nhóm V–VI}
\end{aligned}\]

\textbf{M11} (full three-way, kiểm định tương tác tổng hợp ba chiều,
tổng hợp H3 và H5; P7):
\[\ln NSLD_{it} = \alpha + \beta_1 CDDXK\_c + \beta_2 CDDXK\_c^2 + \beta_3 NLCN\_z + \beta_6 CSS\_z + \beta_{11} ICRV\]
\[+ \beta_{7}(CDDXK\_c \times CSS\_z) + \beta_{8}(CDDXK\_c^2 \times CSS\_z)\]
\[+ \beta_{12}(CDDXK\_c \times ICRV) + \beta_{13}(CDDXK\_c^2 \times ICRV)\]
\[+ \beta_{14}(CDDXK\_c \times CSS\_z \times ICRV) + \beta_9 KNQLy + \beta_{10} NQL\_nu + \gamma \cdot X + \delta_{ct} + \varepsilon\]
\[\begin{aligned}
& \text{Tương tác chế độ-yếu: FSTS} \times \text{Weak} = -0{,}523\ (p = {,}087); \\
& \text{ba chiều không ý nghĩa; } TP\ (M10) = 45{,}9\%,\ \text{LM } p < {,}001
\end{aligned}\]

\textbf{Bảng định nghĩa biến, Nghiên cứu 7 (đa quốc gia)}

\begin{longtable}[]{@{}
  >{\raggedright\arraybackslash}p{(\columnwidth - 6\tabcolsep) * \real{0.2250}}
  >{\raggedright\arraybackslash}p{(\columnwidth - 6\tabcolsep) * \real{0.2750}}
  >{\raggedright\arraybackslash}p{(\columnwidth - 6\tabcolsep) * \real{0.2750}}
  >{\raggedright\arraybackslash}p{(\columnwidth - 6\tabcolsep) * \real{0.2250}}@{}}
\toprule\noalign{}
\begin{minipage}[b]{\linewidth}\raggedright
Ký hiệu
\end{minipage} & \begin{minipage}[b]{\linewidth}\raggedright
WBES item
\end{minipage} & \begin{minipage}[b]{\linewidth}\raggedright
Cách tính
\end{minipage} & \begin{minipage}[b]{\linewidth}\raggedright
Vai trò
\end{minipage} \\
\midrule\noalign{}
\endhead
\bottomrule\noalign{}
\endlastfoot
lnNSLD & d2/n3, l1 & ln(doanh thu/lao động) & Biến phụ thuộc \\
CDDXK\_c & d3c & (d3c/100) − mean\_pool(FSTS): centred & Biến độc lập
(I) \\
CDDXK\_c² & d3c & CDDXK\_c bình phương & I, phi tuyến \\
NLCN\_z & b8, e6 & chuẩn hóa z(cert chất lượng + CN nước ngoài) & Năng
lực công nghệ \\
CSS\_z & c22b/e1, k33 & chuẩn hóa z(website + thanh toán điện tử): Tầng
1 và 2 & Số hoá (DAI) \\
KNQLy & b7 & năm kinh nghiệm nhà quản lý trong ngành & Quản trị \\
NQL\_nu & b7a/b6a & 1 nếu top manager là nữ & Quản trị \\
ICRV & phân loại & integer 1--6 (tiên tiến đổi mới đến SIDS) & Thể
chế \\
lnLD & l1 & ln(lao động thường trực) & Control (quy mô) \\
TuoiDN & b5 & năm\_khảo\_sát − b5 & Control (tuổi DN) \\
SoHuuNN & b6a & tỷ lệ sở hữu nước ngoài (0--1) & Control (sở hữu) \\
δ\_ct & n/a & country × year FE & hiệu ứng cố định \\
\end{longtable}

\emph{Ghi chú: Mẫu phân tích giảm nhẹ từ N = 81.022 (M2) xuống 79.080
(M5, thêm kiểm soát); các mô hình điều tiết M7, M8 và M10 giữ N tương
đương 79.080--79.220. DAI P7 là Tầng 1 và 2 khi k33 có mặt, Tầng 1 only
khi thiếu, khác với P3 Vietnam (chỉ Tầng 1) và đồng nhất với P4
Singapore (Tầng 1 và 2). ICRV là biến điều tiết liên tục (integer),
không phải dummies, nhằm giữ N đủ lớn trong M10--M11.}

\begin{center}\rule{0.5\linewidth}{0.5pt}\end{center}

\hypertarget{muxf4-huxecnh-cux1ee5-thux1ec3-nghiuxean-cux1ee9u-8-ux111ux1ea3o-nhux1ecf-thuxe1i-buxecnh-dux1b0ux1a1ng-n-959-nhuxf3m-vi}{%
\subsubsection{3.4.5.5 Mô hình cụ thể: Nghiên cứu 8: đảo nhỏ Thái Bình
Dương (N = 959, Nhóm
VI)}\label{muxf4-huxecnh-cux1ee5-thux1ec3-nghiuxean-cux1ee9u-8-ux111ux1ea3o-nhux1ecf-thuxe1i-buxecnh-dux1b0ux1a1ng-n-959-nhuxf3m-vi}}

Nghiên cứu 8 phân tích mẫu gồm \textbf{N = 959 doanh nghiệp} tại
\textbf{7 nền kinh tế Pacific SIDS} (Fiji, Kiribati, Papua New Guinea,
Samoa, Solomon Islands, Tonga, Vanuatu) từ các sóng WBES 2009--2025. Đây
là nhóm thể chế ICRV Nhóm VI, cấp thấp nhất trong phân loại 6 chế độ,
đặc trưng bởi thị trường nội địa cực nhỏ, chi phí thương mại cao và hỗ
trợ thể chế yếu. Trong tổng mẫu, chỉ \textbf{132 doanh nghiệp có hoạt
động xuất khẩu} (13,8\%), phản ánh cấu trúc thị trường đảo nhỏ.

Phát hiện trọng tâm: \textbf{gánh nặng quốc tế hóa bắt buộc (Forced
Internationalization Penalty, FIP)}, mối quan hệ quốc tế hóa--hiệu quả
là \textbf{đơn điệu âm} tại Nhóm VI, trái ngược với hình chữ U ngược
được quan sát ở P3--P7. Lind--Mehlum U-test không bác bỏ đơn điệu (p
\textgreater,10), xác nhận không có điểm quay.

\textbf{Ký hiệu biến (tiếng Việt ↔ mã WBES):}

\begin{longtable}[]{@{}
  >{\raggedright\arraybackslash}p{(\columnwidth - 6\tabcolsep) * \real{0.2000}}
  >{\raggedright\arraybackslash}p{(\columnwidth - 6\tabcolsep) * \real{0.3556}}
  >{\raggedright\arraybackslash}p{(\columnwidth - 6\tabcolsep) * \real{0.2000}}
  >{\raggedright\arraybackslash}p{(\columnwidth - 6\tabcolsep) * \real{0.2444}}@{}}
\toprule\noalign{}
\begin{minipage}[b]{\linewidth}\raggedright
Ký hiệu
\end{minipage} & \begin{minipage}[b]{\linewidth}\raggedright
Tên tiếng Việt
\end{minipage} & \begin{minipage}[b]{\linewidth}\raggedright
Mã WBES
\end{minipage} & \begin{minipage}[b]{\linewidth}\raggedright
Cách tính
\end{minipage} \\
\midrule\noalign{}
\endhead
\bottomrule\noalign{}
\endlastfoot
lnNSLD & Log năng suất lao động & d2, l1 & ln(doanh thu PPP / lao động
thường trực) \\
CDDXK\_c & Cường độ xuất khẩu (điều chỉnh) & d3c & (d3c/100) −
mean\_wave(FSTS): centred theo sóng \\
CDDXK\_c² & CDDXK\_c bình phương & d3c & CDDXK\_c × CDDXK\_c \\
NLCN\_z & Năng lực công nghệ & b8, e6 & chuẩn hóa z(trung bình cert chất
lượng + CN nước ngoài) \\
CSS\_z & Chỉ số số hoá (Tầng 1) & c22b & chuẩn hóa z(website binary):
\textbf{chỉ Tầng 1, không dynamic} \\
lnLD & Log quy mô lao động & l1 & ln(lao động thường trực) \\
TuoiDN & Tuổi doanh nghiệp & b5 & năm\_khảo\_sát − b5 \\
SoHuuNN & Tỷ lệ sở hữu nước ngoài & b6a & tỷ lệ 0--1 \\
δ\_c & Country hiệu ứng cố định & a3b & biến giả theo quốc gia \\
λ\_t & đợt hiệu ứng cố định & đợt & biến giả theo sóng điều tra \\
\end{longtable}

\emph{Ghi chú: CSS\_z (DAI) là biến đại diện tĩnh Tầng 1 (website binary
c22b). WBES không thu thập đủ thang đo cho số hoá động tại các đảo nhỏ,
KHÔNG gọi là ``năng lực số hoá động''.}

\textbf{Chuỗi mô hình M0--M3 (P8):}

M0, cơ sở kiểm soát:

\[lnNSLD_{it} = \alpha + \gamma_1 lnLD_{it} + \gamma_2 TuoiDN_{it} + \gamma_3 SoHuuNN_{it} + \delta_c + \lambda_t + \varepsilon_{it}\]

M1, Kiểm định FIP (H1b): thêm cường độ xuất khẩu tuyến tính

\[lnNSLD_{it} = \alpha + \beta_1 CDDXK\_c_{it} + \gamma \cdot X_{it} + \delta_c + \lambda_t + \varepsilon_{it}\]

\[\beta_1 = -1{,}339,\ p < {,}001\ (\text{country+year FE, N} = 959) \Rightarrow \textbf{FIP xác nhận}\]

M2, Kiểm định phi tuyến: thêm CDDXK\_c²

\[lnNSLD_{it} = \alpha + \beta_1 CDDXK\_c_{it} + \beta_2 CDDXK\_c^2_{it} + \gamma \cdot X_{it} + \delta_c + \lambda_t + \varepsilon_{it}\]

\[\begin{aligned}
& \beta_1\ \text{NS},\ \beta_2\ \text{NS};\ \text{Lind–Mehlum U-test NS} \\
& \Rightarrow \text{không có điểm quay, đơn điệu âm}
\end{aligned}\]

M3, Kiểm định năng lực điều tiết (H2 null cho SIDS):

\[\begin{aligned}
lnNSLD_{it} = {} & \alpha + \beta_1 CDDXK\_c_{it} + \beta_2 CDDXK\_c^2_{it} + \beta_3 NLCN\_z_{it} \\
& + \beta_4 CSS\_z_{it} + \gamma \cdot X_{it} + \delta_c + \lambda_t + \varepsilon_{it}
\end{aligned}\]

\[\begin{aligned}
& \beta_3\ (NLCN\_z):\ p = {,}003\ (\text{dương, có ý nghĩa});\ \beta_4\ (CSS\_z):\ p = {,}285\ \text{NS} \\
& \Rightarrow \text{TCI nâng mặt bằng năng suất; không có biến điều tiết uốn đường cong (H2 null)}
\end{aligned}\]

\textbf{Phân tích độ vững (kiểm tra độ vững):}

\begin{longtable}[]{@{}llll@{}}
\toprule\noalign{}
Đặc tả & β(CDDXK\_c) & SE & p \\
\midrule\noalign{}
\endhead
\bottomrule\noalign{}
\endlastfoot
M1 country+year FE (N = 959) & −1,339 & 0,386 & \textless,001 \\
Year FE only (không country FE) & −3,351 & 0,808 & \textless,001 \\
Bivariate (không controls) & −0,864 & 0,441 & ,050 \\
chỉ doanh nghiệp xuất khẩu (N = 26) & −1,176 & 0,763 & ,130 (NS) \\
\end{longtable}

\emph{Hệ số âm nhất quán về dấu qua các specification; suy luận FIP dựa
chủ yếu vào các mô hình fixed-effects (M1, year-FE), trong khi chỉ doanh
nghiệp xuất khẩu (N=26) thiếu công suất thống kê.}

\textbf{Bảng định nghĩa biến, Nghiên cứu 8 (SIDS)}

\begin{longtable}[]{@{}
  >{\raggedright\arraybackslash}p{(\columnwidth - 6\tabcolsep) * \real{0.2250}}
  >{\raggedright\arraybackslash}p{(\columnwidth - 6\tabcolsep) * \real{0.2750}}
  >{\raggedright\arraybackslash}p{(\columnwidth - 6\tabcolsep) * \real{0.2750}}
  >{\raggedright\arraybackslash}p{(\columnwidth - 6\tabcolsep) * \real{0.2250}}@{}}
\toprule\noalign{}
\begin{minipage}[b]{\linewidth}\raggedright
Ký hiệu
\end{minipage} & \begin{minipage}[b]{\linewidth}\raggedright
WBES item
\end{minipage} & \begin{minipage}[b]{\linewidth}\raggedright
Cách tính
\end{minipage} & \begin{minipage}[b]{\linewidth}\raggedright
Vai trò
\end{minipage} \\
\midrule\noalign{}
\endhead
\bottomrule\noalign{}
\endlastfoot
lnNSLD & d2, l1 & ln(doanh thu PPP / lao động) & Biến phụ thuộc \\
CDDXK\_c & d3c & (d3c/100) − mean\_wave: centred & Biến độc lập (I) \\
CDDXK\_c² & d3c & CDDXK\_c bình phương & I, kiểm định phi tuyến \\
NLCN\_z & b8, e6 & chuẩn hóa z(cert + CN nước ngoài) & Năng lực công
nghệ \\
CSS\_z & c22b & chuẩn hóa z(website binary): Tầng 1 only & Số hoá, biến
đại diện tĩnh \\
lnLD & l1 & ln(lao động thường trực) & Control (quy mô) \\
TuoiDN & b5 & năm\_KS − năm\_thành\_lập & Control (tuổi DN) \\
SoHuuNN & b6a & tỷ lệ sở hữu nước ngoài & Control (sở hữu) \\
δ\_c & a3b & country FE & hiệu ứng cố định \\
λ\_t & đợt & đợt FE & hiệu ứng cố định \\
\end{longtable}

\emph{Ghi chú: FIP là điều kiện biên (boundary condition) của mô hình
chữ U ngược, áp dụng đặc thù tại ICRV Nhóm VI nơi 3 điều kiện cấu trúc
bị vi phạm đồng thời: (1) không có thị trường nội địa đủ lớn, (2) chi
phí thương mại không kiểm soát được, (3) hỗ trợ thể chế không đủ chức
năng. DAI (CSS\_z) Tầng 1 only, WBES không thu thập thang đo đầy đủ cho
số hoá động tại SIDS.}

\begin{center}\rule{0.5\linewidth}{0.5pt}\end{center}

\hypertarget{muxf4-huxecnh-cux1ee5-thux1ec3-nghiuxean-cux1ee9u-9-ux1ea5n-ux111ux1ed9-ux111ux1ed9-bux1ec1n-ngux1b0ux1ee1ng-qua-chuyux1ec3n-ux111ux1ed5i-thux1ec3-chux1ebf-p9}{%
\subsubsection{3.4.5.6 Mô hình cụ thể, Nghiên cứu 9: Ấn Độ, độ bền
ngưỡng qua chuyển đổi thể chế
(P9)}\label{muxf4-huxecnh-cux1ee5-thux1ec3-nghiuxean-cux1ee9u-9-ux1ea5n-ux111ux1ed9-ux111ux1ed9-bux1ec1n-ngux1b0ux1ee1ng-qua-chuyux1ec3n-ux111ux1ed5i-thux1ec3-chux1ebf-p9}}

Nghiên cứu 9 sử dụng ba đợt Khảo sát Doanh nghiệp Ngân hàng Thế giới cho
Ấn Độ (2014 PICS3, \(N_{\text{phân tích}} = 8.941\); 2022 BEE,
\(N_{\text{phân tích}} = 9.300\); 2025 BREADY,
\(N_{\text{phân tích}} = 10.476\); mẫu gộp \(N = 28.717\)), mẫu gộp đơn
quốc gia lớn nhất của luận án. Ấn Độ thuộc Nhóm IV ICRV (Chuyển đổi thu
nhập trung bình--thấp); cửa sổ 2014--2025 bao trùm chuỗi cú sốc thể chế
(bãi bỏ tiền mặt 2016, GST 2017, IBC 2016, UPI 2016, PLI và Atmanirbhar
Bharat 2020, COVID-19), tạo điều kiện stress-test cho giả thuyết bền
vững ngưỡng.

Mô hình ước lượng OLS với sai số chuẩn vững HC1, gồm số hạng FSTS bậc
nhất và bậc hai (centred theo đợt) cùng các tương tác năng lực (TCI,
DAI). Ba kiểm định bổ sung: (i) Lind--Mehlum U-test xác nhận U ngược
theo từng đợt; (ii) z-test Paternoster (1998) so sánh hệ số giữa 2014 và
2025 (cửa sổ 11 năm chính), báo cáo cả sai số chuẩn HC1 và cụm theo
bang; (iii) đặc tả gộp ba đợt với biến giả/xu thế đợt khảo sát làm kiểm
định độ vững. Việc UPI ra mắt tháng 4/2016 cung cấp thí nghiệm gần tự
nhiên cho bão hòa DAI; riêng đợt 2025 (BREADY) bổ sung chỉ báo Tầng 2
(\texttt{k33}, tỷ trọng doanh thu qua thanh toán điện tử) cho phép kiểm
định điều tiết DAI Tầng 2, được xử lý như phân tích thăm dò vì chỉ có ở
một đợt.

\hypertarget{sai-sux1ed1-chuux1ea9n}{%
\subsection{3.4.6 Sai số chuẩn}\label{sai-sux1ed1-chuux1ea9n}}

Tất cả các mô hình đều sử dụng HC1/HC3 sai số chuẩn vững (Long \& Ervin,
2000; White, 1980) để xử lý phương sai thay đổi. Đối với multi-country,
bổ sung clustered SE theo country.

\hypertarget{lux1ef1a-chux1ecdn-cux1ea5u-truxfac-sai-sux1ed1-chuux1ea9n-theo-thiux1ebft-kux1ebf-dux1eef-liux1ec7u}{%
\subsection{3.4.6.1 Lựa chọn cấu trúc sai số chuẩn theo thiết kế dữ
liệu}\label{lux1ef1a-chux1ecdn-cux1ea5u-truxfac-sai-sux1ed1-chuux1ea9n-theo-thiux1ebft-kux1ebf-dux1eef-liux1ec7u}}

Cấu trúc sai số chuẩn được lựa chọn theo bản chất thiết kế của từng
nghiên cứu chứ không áp dụng một quy ước đồng nhất. Các nghiên cứu đơn
quốc gia trên lát cắt ngang dùng sai số chuẩn vững với phương sai thay
đổi loại HC1, một hiệu chỉnh chuẩn xử lý phương sai sai số không đồng
nhất giữa các doanh nghiệp mà không đòi hỏi giả định về dạng cụ thể của
phương sai. Khung đa quốc gia bổ sung phân cụm sai số chuẩn theo nền
kinh tế, vì thiết kế chọn mẫu phân tầng của WBES tạo ra tương quan nội
cụm giữa các doanh nghiệp trong cùng một nền kinh tế, khiến giả định độc
lập của HC1 không còn thỏa đáng. Phân cụm theo nền kinh tế cũng phản ánh
đúng cấp độ mà các cú sốc thể chế và chính sách tác động đồng thời lên
nhiều doanh nghiệp.

Cụ thể, mọi mô hình trong khung đa quốc gia đều mang hiệu ứng cố định
hai chiều theo nền kinh tế và theo năm, cùng sai số chuẩn vững phân cụm
theo nền kinh tế theo hiệu chỉnh CRV1. Cấp phân cụm theo nền kinh tế
được chọn vì thiết kế chọn mẫu phân tầng của bộ công cụ tạo tương quan
nội cụm giữa các doanh nghiệp trong cùng một nền, và vì các cú sốc thể
chế cùng chính sách tác động đồng thời lên nhiều doanh nghiệp ở cùng cấp
này.

Một giới hạn cần ghi nhận của phân cụm theo nền kinh tế là số cụm hữu
hạn có thể làm sai số chuẩn tiệm cận kém chính xác khi quy mô cụm chênh
lệch lớn, đặc biệt giữa các nền kinh tế đông doanh nghiệp và các đảo nhỏ
chỉ có vài trăm quan sát. Trong các tình huống này, luận án diễn giải
suy luận một cách thận trọng và neo kết luận chính vào tính nhất quán về
dấu qua nhiều đặc tả thay vì vào ý nghĩa của một hệ số đơn lẻ. Việc luận
án không áp dụng trọng số chọn mẫu trong ước lượng chính tuân theo
khuyến nghị của Solon, Haider và Wooldridge (2015): khi mô hình đã kiểm
soát các chiều phân tầng qua biến kiểm soát và qua hiệu ứng cố định, ước
lượng bình phương nhỏ nhất không trọng số cho hệ số nhất quán và thường
hiệu quả hơn, trong khi trọng số được dành cho phân tích độ vững nhằm
kiểm tra tính ổn định.

\hypertarget{kiux1ec3m-ux111ux1ecbnh-ux111ux1ed9-vux1eefng}{%
\section{3.5 Kiểm định độ
vững}\label{kiux1ec3m-ux111ux1ecbnh-ux111ux1ed9-vux1eefng}}

\hypertarget{thay-ux111ux1ed5i-thux1b0ux1edbc-ux111o}{%
\subsection{3.5.1 Thay đổi thước
đo}\label{thay-ux111ux1ed5i-thux1b0ux1edbc-ux111o}}

Lớp độ vững thứ nhất hoán đổi thước đo các biến trọng tâm. Biến phụ
thuộc năng suất lao động dạng logarit được thay bằng tỷ suất lợi nhuận
trên doanh thu, tăng trưởng doanh thu và tăng trưởng việc làm để kiểm
tra liệu hình dạng quan hệ có phụ thuộc vào một định nghĩa hiệu quả duy
nhất hay không. Biến độc lập cường độ xuất khẩu được thay bằng trạng
thái xuất khẩu nhị phân nhằm phân biệt hiệu ứng ở biên tham gia khỏi
hiệu ứng ở biên cường độ. Biến điều tiết số được hoán đổi giữa biến thể
chỉ gồm Tầng 1 và biến thể gồm cả Tầng 2 trên mẫu con từ năm 2024 trở
đi, nơi các chỉ báo số phong phú hơn đã sẵn có.

\hypertarget{mux1eabu-phux1ee5}{%
\subsection{3.5.2 Mẫu phụ}\label{mux1eabu-phux1ee5}}

Lớp thứ hai giới hạn mẫu để loại trừ các giải thích thay thế dựa trên
thành phần mẫu. Các mẫu con gồm: loại doanh nghiệp siêu nhỏ dưới năm lao
động, chỉ giữ doanh nghiệp vừa và nhỏ, chỉ giữ ngành chế biến chế tạo,
chỉ giữ doanh nghiệp có xuất khẩu với cường độ dương, và tách theo từng
chế độ trong sáu chế độ thể chế ICRV. Việc điểm uốn và dấu của độ cong
được bảo toàn qua các mẫu con này là bằng chứng rằng quan hệ trọng tâm
không phải tạo tác của một nhóm doanh nghiệp đặc thù.

\hypertarget{giuxe1-trux1ecb-ngoux1ea1i-lai}{%
\subsection{3.5.3 Giá trị ngoại
lai}\label{giuxe1-trux1ecb-ngoux1ea1i-lai}}

Lớp thứ ba kiểm soát ảnh hưởng của giá trị ngoại lai bằng cách cắt đuôi
năng suất lao động và số lao động tại các ngưỡng phân vị 1 và 99, rồi
lặp lại tại ngưỡng 5 và 95, áp dụng riêng trong từng cặp nền kinh tế và
năm để tránh nhầm chênh lệch mặt bằng giữa các nền kinh tế với giá trị
ngoại lai thực.

\hypertarget{ux111a-cux1ed9ng-tuyux1ebfn-vuxe0-phux1b0ux1a1ng-sai-thay-ux111ux1ed5i}{%
\subsection{3.5.4 Đa cộng tuyến và phương sai thay
đổi}\label{ux111a-cux1ed9ng-tuyux1ebfn-vuxe0-phux1b0ux1a1ng-sai-thay-ux111ux1ed5i}}

Hai mối lo của hồi quy có số hạng phi tuyến và tương tác là đa cộng
tuyến và phương sai sai số thay đổi. Để xử lý đa cộng tuyến giữa số hạng
bậc nhất và số hạng bậc hai của cường độ xuất khẩu, cũng như giữa các
thành phần của số hạng tương tác, biến cường độ xuất khẩu được căn giữa
quanh trung bình mẫu trước khi lập bình phương và lập tương tác; căn
giữa làm giảm mạnh tương quan giả tạo giữa một biến và bình phương của
chính nó mà không làm thay đổi hệ số bậc hai hay vị trí điểm uốn. Hệ số
phóng đại phương sai được tính cho mọi mô hình với ngưỡng cảnh báo dưới
năm (Hair et al., 2019); trong các đặc tả đã căn giữa, hệ số phóng đại
phương sai nằm an toàn dưới ngưỡng này, xác nhận rằng các ước lượng hệ
số không bị bất ổn do cộng tuyến. Đối với phương sai sai số thay đổi,
kiểm định Breusch và Pagan (1979) được dùng để phát hiện, và quan trọng
hơn, mọi suy luận đều dựa trên sai số chuẩn vững với phương sai thay đổi
và được gom cụm theo nền kinh tế, do đó kết luận thống kê có hiệu lực
ngay cả khi phương sai sai số không đồng nhất, một cấu trúc sai số phù
hợp với thiết kế dữ liệu gộp nhiều nền kinh tế đã trình bày ở Mục 3.4.6.

\hypertarget{phuxe2n-tuxedch-ux111ux1ed9-nhux1ea1y-cho-phuxe2n-tuxedch-tux1ed5ng-hux1ee3p}{%
\subsection{3.5.5 Phân tích độ nhạy cho phân tích tổng
hợp}\label{phuxe2n-tuxedch-ux111ux1ed9-nhux1ea1y-cho-phuxe2n-tuxedch-tux1ed5ng-hux1ee3p}}

Khung phân tích tổng hợp (P6) áp dụng một bộ kiểm định độ nhạy riêng phù
hợp với cấu trúc dữ liệu ba tầng. Kiểm định bỏ một nghiên cứu lần lượt
loại từng nghiên cứu khỏi mẫu và tái ước lượng cỡ ảnh hưởng gộp, nhằm
bảo đảm kết luận không bị chi phối bởi một nghiên cứu đơn lẻ có trọng số
lớn hay giá trị cực đoan. Lệch lạc công bố được đánh giá bằng nhiều công
cụ bổ trợ: kiểm tra trực quan biểu đồ phễu, kiểm định bất đối xứng của
Egger, và thủ tục cắt và điền của Duval và Tweedie (2000) để ước lượng
và hiệu chỉnh số nghiên cứu khuyết có thể có. Ngoài ra, cấu trúc ba tầng
cho phép so sánh mô hình ba tầng với mô hình hai tầng truyền thống nhằm
kiểm tra xem việc tách riêng phương sai trong nội bộ nghiên cứu khỏi
phương sai giữa các nghiên cứu có thực sự cải thiện độ phù hợp hay
không, đồng thời cho phép phân rã \(I^2\) thành thành phần trong và giữa
nghiên cứu để định vị nguồn dị biệt. Kết hợp các kiểm định này bảo đảm
rằng cỡ ảnh hưởng gộp và kết luận điều tiết theo ICRV là vững trước cả
lựa chọn mẫu nghiên cứu lẫn giả định cấu trúc phương sai.

\hypertarget{giao-thux1ee9c-ux111ux1ed9-vux1eefng-ux111a-lux1edbp-vuxe0-ux111a-ux111ux1ee3t-cho-khung-thux1ef1c-nghiux1ec7m}{%
\subsection{3.5.5.1 Giao thức độ vững đa lớp và đa đợt cho khung thực
nghiệm}\label{giao-thux1ee9c-ux111ux1ed9-vux1eefng-ux111a-lux1edbp-vuxe0-ux111a-ux111ux1ee3t-cho-khung-thux1ef1c-nghiux1ec7m}}

Khung thực nghiệm áp dụng một giao thức độ vững có cấu trúc nhằm bảo đảm
suy luận trọng tâm không phụ thuộc vào một quyết định đo lường, một mẫu
con, hay một đặc tả đơn lẻ. Lớp thứ nhất hoán đổi thước đo: biến phụ
thuộc được thay bằng tỷ suất lợi nhuận trên doanh thu, tăng trưởng doanh
thu và tăng trưởng việc làm; biến độc lập được thay bằng trạng thái xuất
khẩu nhị phân; biến điều tiết số được hoán đổi giữa biến thể Tầng 1 và
biến thể Tầng 1 cộng Tầng 2 trên mẫu con có đủ chỉ báo. Lớp thứ hai giới
hạn mẫu con: loại doanh nghiệp siêu nhỏ, chỉ giữ doanh nghiệp vừa và
nhỏ, chỉ giữ ngành chế tạo, chỉ giữ doanh nghiệp có xuất khẩu, và tách
theo từng chế độ thể chế. Lớp thứ ba kiểm soát phân phối qua cắt đuôi
tại hai ngưỡng phân vị khác nhau trong từng cặp nền kinh tế và năm.

Lớp thứ tư kiểm tra các giả định kinh tế lượng phụ trợ. Hệ số phóng đại
phương sai (variance inflation factor, VIF) được tính cho mọi mô hình
với ngưỡng cảnh báo dưới năm để phát hiện đa cộng tuyến đe dọa độ ổn
định hệ số; kiểm định Breusch--Pagan đánh giá phương sai thay đổi. Đối
với cấu phần phân tích tổng hợp, giao thức bổ sung kiểm định loại từng
nghiên cứu (leave-one-out) tái tính hiệu ứng gộp nhằm phát hiện ảnh
hưởng quá mức của một nghiên cứu đơn lẻ, cùng thủ tục cắt và điền
(trim-and-fill) đánh giá thiên lệch công bố. Luận án không áp dụng hiệu
chỉnh kiểm định bội chính thức cho toàn bộ bảng độ vững, vì các bảng
kiểm tra các quan ngại nhận diện khác nhau thay vì lặp lại cùng một giả
thuyết; thay vào đó, suy luận thực chất được neo vào mẫu hình nhất quán
qua các bảng và tính nhất quán về dấu của ước lượng trọng tâm chứ không
vào ý nghĩa của một bảng độ vững đơn lẻ.

\hypertarget{chiux1ebfn-lux1b0ux1ee3c-nhux1eadn-dux1ea1ng-vuxe0-ranh-giux1edbi-suy-luux1eadn-nhuxe2n-quux1ea3}{%
\subsection{3.5.6 Chiến lược nhận dạng và ranh giới suy luận nhân
quả}\label{chiux1ebfn-lux1b0ux1ee3c-nhux1eadn-dux1ea1ng-vuxe0-ranh-giux1edbi-suy-luux1eadn-nhuxe2n-quux1ea3}}

Vì dữ liệu chủ đạo của luận án là lát cắt ngang lặp lại (repeated
cross-sections) chứ không phải bảng cân bằng theo dõi cùng doanh nghiệp
qua thời gian, luận án thừa nhận minh bạch các ranh giới suy luận và áp
dụng một chiến lược nhận dạng nhiều lớp để củng cố tính giá trị nội tại
trong giới hạn cho phép. Thứ nhất, mọi mô hình đa quốc gia đều bao gồm
hiệu ứng cố định hai chiều (nền kinh tế và năm), qua đó hấp thụ toàn bộ
thành phần bất biến theo nước (tiền tệ, mặt bằng giá, chế độ kế toán,
đặc điểm thể chế ổn định) và các cú sốc chung theo năm; hệ số quan hệ
quốc tế hóa và hiệu quả do đó được nhận dạng từ biến thiên nội bộ nền
kinh tế--năm, không bị nhiễu bởi khác biệt cấu trúc cố định giữa các
nước (Wooldridge, 2010). Thứ hai, ở các nghiên cứu đơn quốc gia, luận án
bổ sung biến công cụ (instrumental variable) cho năng lực công nghệ, với
các kiểm định độ mạnh công cụ (first-stage F) và over-identification phù
hợp, nhằm tách hiệu ứng nhân quả khỏi nội sinh do lựa chọn và đồng thời
(P3, P5). Cần phân biệt nội sinh của năng lực, vốn được xử lý bằng biến
công cụ nói trên, với một mối đe dọa nhận dạng riêng là tự chọn lọc tham
gia xuất khẩu: theo logic Melitz (2003), doanh nghiệp năng suất cao tự
chọn vào xuất khẩu, nên hệ số FSTS có thể trộn lẫn hiệu ứng nhân quả của
quốc tế hóa với sự chọn lọc có sẵn. Luận án khoanh vùng đe dọa này theo
ba cách: (i) hiệu ứng cố định hai chiều cùng chuẩn hóa within
country--year hấp thụ chênh lệch mức năng suất cố định giữa doanh nghiệp
xuất khẩu và không xuất khẩu ở tầng nền kinh tế--năm; (ii) đối tượng suy
luận trọng tâm là hình dạng phi tuyến (độ cong và vị trí điểm uốn) chứ
không chỉ là mức hệ số, mà một câu chuyện chọn lọc mức đơn thuần khó tạo
ra một cực đại nội vùng có hệ thống; và (iii) nghiên cứu Singapore (P4)
bổ sung hiệu chỉnh chọn lọc kiểu Heckman với tỷ số Mills nghịch đảo như
một kiểm tra cục bộ. Tuy vậy, do hiệu chỉnh chọn lọc tham gia không được
thực hiện đồng bộ ở mọi nghiên cứu thành phần, luận án thừa nhận đây là
một ranh giới suy luận còn lại của thiết kế lát cắt ngang. Thứ ba, độ ổn
định của ước lượng trước thay đổi do biến bỏ sót được đánh giá bằng giới
hạn δ-stability theo Oster (2019), một chuẩn mực ngày càng được áp dụng
trong nghiên cứu thực nghiệm khi không có thiết kế thực nghiệm thuần
túy. Cuối cùng, luận án tuân thủ nguyên tắc báo cáo trung thực: các kết
quả từ mẫu nhỏ hoặc thiếu công suất (ví dụ mẫu chỉ gồm doanh nghiệp xuất
khẩu tại SIDS) được trình bày như ``ranh giới suy luận'' (inferential
bounds) chứ không phải bằng chứng nhân quả khẳng định, và mọi suy luận
về điều tiết thể chế được neo đồng thời vào hai nguồn độc lập là kiểm
định điều tiết của phân tích tổng hợp và gradient điểm uốn của mô hình
đa quốc gia (xem Mục 4.2.3, Mục 4.6.2). Cách tiếp cận đa lớp này tuân
theo các khuyến nghị về suy luận từ dữ liệu quan sát quy mô lớn (Angrist
\& Pischke, 2009). Một giới hạn đo lường cần lưu ý là \textbf{tính bất
biến đo lường (measurement invariance)} của các cấu trúc TCI và DAI qua
ba thế hệ bảng hỏi WBES: do cấu trúc câu hỏi thay đổi giữa PICS3,
Standardized và BREADY, một phần khác biệt theo thời gian có thể phản
ánh hiệu ứng lược đồ thay vì thay đổi thực chất (ví dụ mức R\&D đo được
ở một số nhóm). Luận án giảm thiểu rủi ro này bằng hiệu ứng cố định
năm/đợt, chuẩn hóa within country--year, và kiểm tra độ nhạy với biến
thể thước đo (Mục 3.5.1); tuy vậy luận án chưa thực hiện kiểm định bất
biến đo lường chính thức, nên mọi so sánh xuyên-lược đồ cần được đọc một
cách thận trọng.

\hypertarget{chiux1ebfn-lux1b0ux1ee3c-biux1ebfn-cuxf4ng-cux1ee5-vuxe0-kiux1ec3m-ux111ux1ecbnh-ux111ux1ed9-mux1ea1nh-cuxf4ng-cux1ee5}{%
\subsection{3.5.6.1 Chiến lược biến công cụ và kiểm định độ mạnh công
cụ}\label{chiux1ebfn-lux1b0ux1ee3c-biux1ebfn-cuxf4ng-cux1ee5-vuxe0-kiux1ec3m-ux111ux1ecbnh-ux111ux1ed9-mux1ea1nh-cuxf4ng-cux1ee5}}

Để tách hiệu ứng nhân quả của năng lực khỏi nội sinh do lựa chọn và đồng
thời, các nghiên cứu đơn quốc gia bổ sung ước lượng bình phương nhỏ nhất
hai bước (two-stage least squares, 2SLS). Công cụ được xây dựng từ tỷ lệ
áp dụng đồng cấp trong các ô ngành, vùng và đợt khảo sát theo nguyên tắc
loại bỏ chính doanh nghiệp khỏi trung bình ô (leave-one-out), nhằm bảo
đảm giá trị công cụ của một doanh nghiệp không chứa thông tin của chính
nó. Logic nhận dạng dựa trên hai điều kiện: tính liên quan, vì tỷ lệ áp
dụng đồng cấp dự báo mạnh quyết định áp dụng của doanh nghiệp riêng lẻ;
và loại trừ xấp xỉ, vì tỷ lệ áp dụng trung bình của các doanh nghiệp
khác trong cùng ô khó tác động đến năng suất của một doanh nghiệp qua
kênh nào khác ngoài chính quyết định áp dụng, sau khi đã kiểm soát ô
ngành và đợt.

Độ mạnh công cụ được đánh giá bằng thống kê F của hồi quy giai đoạn một,
với kết quả nằm trong khoảng 22 đến 35, vượt xa ngưỡng kinh nghiệm về
công cụ yếu. Một phát hiện phương pháp luận đáng chú ý là ước lượng 2SLS
cho hai chiều năng lực phân kỳ rõ rệt: hệ số công cụ hóa của năng lực
công nghệ dương mạnh (ước lượng 2SLS bằng 1,64, p \textless{} 0,001,
dưới một công cụ mạnh với F giai đoạn một bằng 22,1), trong khi hệ số
công cụ hóa của chấp nhận số ở mức Tầng 1 không khác biệt thống kê với
không, một kết quả độ vững củng cố cách diễn giải tính lỗi thời của biến
đại diện Tầng 1 trong môi trường số đã lan tỏa gần phổ quát.

\hypertarget{giux1edbi-hux1ea1n-ux1ed5n-ux111ux1ecbnh-hux1ec7-sux1ed1-oster-trux1b0ux1edbc-biux1ebfn-bux1ecf-suxf3t-khuxf4ng-quan-suxe1t}{%
\subsection{3.5.6.2 Giới hạn ổn định hệ số Oster trước biến bỏ sót không
quan
sát}\label{giux1edbi-hux1ea1n-ux1ed5n-ux111ux1ecbnh-hux1ec7-sux1ed1-oster-trux1b0ux1edbc-biux1ebfn-bux1ecf-suxf3t-khuxf4ng-quan-suxe1t}}

Vì dữ liệu chủ đạo là lát cắt ngang lặp lại chứ không phải thiết kế thực
nghiệm thuần túy, luận án đánh giá độ vững của hệ số trọng tâm trước
thiên lệch biến bỏ sót không quan sát bằng giới hạn ổn định delta của
Oster (2019). Phương pháp dựa trên trực giác rằng mức độ một hệ số dịch
chuyển khi thêm biến kiểm soát quan sát được mang thông tin về mức độ nó
có thể tiếp tục dịch chuyển nếu kiểm soát thêm các biến không quan sát
có cùng cấu trúc liên hệ. Tham số delta lượng hóa tỷ lệ tương đối giữa
lựa chọn trên biến không quan sát và lựa chọn trên biến quan sát; giá
trị delta bằng một biểu thị giả định bảo thủ rằng các biến không quan
sát quan trọng ngang với các biến quan sát.

Việc tính giới hạn đòi hỏi đặt cận trên cho mức độ giải thích tối đa của
một mô hình giả định đầy đủ biến. Luận án dùng quy ước cận trên này bằng
1,3 lần hệ số xác định của mô hình có kiểm soát, một mức được khuyến
nghị rộng rãi trong thực hành thực nghiệm. Kết luận phương pháp luận là
không hệ số trọng tâm nào đổi dấu hoặc sụp đổ về không dưới các độ lớn
khả tín của lựa chọn không quan sát, qua đó cung cấp một bằng chứng độc
lập rằng các mối liên hệ trọng tâm không phải là tạo phẩm của biến bỏ
sót có thể thấy trước.

\hypertarget{tuxednh-nhux1ea5t-quuxe1n-vux1ec1-dux1ea5u-qua-nhiux1ec1u-ux111ux1eb7c-tux1ea3-nhux1b0-nguyuxean-tux1eafc-suy-luux1eadn-nux1ec1n-tux1ea3ng}{%
\subsection{3.5.6.3 Tính nhất quán về dấu qua nhiều đặc tả như nguyên
tắc suy luận nền
tảng}\label{tuxednh-nhux1ea5t-quuxe1n-vux1ec1-dux1ea5u-qua-nhiux1ec1u-ux111ux1eb7c-tux1ea3-nhux1b0-nguyuxean-tux1eafc-suy-luux1eadn-nux1ec1n-tux1ea3ng}}

Triết lý suy luận xuyên suốt khung độ vững là neo kết luận thực chất vào
tính nhất quán về dấu của ước lượng trọng tâm qua nhiều đặc tả, chứ
không vào ý nghĩa của một hệ số đơn lẻ, vì các bảng độ vững kiểm tra các
quan ngại nhận diện khác nhau thay vì lặp lại cùng một giả thuyết.
Nguyên tắc này được minh họa rõ nhất ở nghiên cứu đảo nhỏ Thái Bình
Dương, nơi hệ số cường độ xuất khẩu giữ dấu âm nhất quán qua mọi đặc tả:
âm và có ý nghĩa rất cao trong mô hình hiệu ứng cố định nền kinh tế và
năm với hệ số −1,339, vẫn âm và có ý nghĩa khi chỉ giữ hiệu ứng cố định
năm hoặc khi bỏ toàn bộ biến kiểm soát, và chỉ mất ý nghĩa ở mẫu con chỉ
gồm doanh nghiệp xuất khẩu vốn quá nhỏ. Suy luận về gánh nặng quốc tế
hóa bắt buộc do đó được neo chủ yếu vào các mô hình hiệu ứng cố định chứ
không vào mẫu con thiếu công suất, đúng nguyên tắc báo cáo trung thực đã
nêu ở Mục 3.5.6.

Nguyên tắc tương tự áp dụng cho các kết quả điều tiết không có ý nghĩa
trong khung đa quốc gia. Năng lực công nghệ vào mô hình với hiệu ứng mức
dương và rất có ý nghĩa, cho thấy doanh nghiệp năng lực cao hơn năng
suất hơn ở mọi mức quốc tế hóa, nhưng các số hạng tương tác với cường độ
xuất khẩu đều không có ý nghĩa và điểm uốn gần như không đổi, nghĩa là
năng lực nâng nền mà không bẻ cong đường. Kết quả không có ý nghĩa của
các số hạng tương tác vì vậy không bị diễn giải là thất bại đo lường mà
là bằng chứng tích cực cho cách phân vai năng lực công nghệ là bộ tăng
mức chứ không phải bộ điều tiết độ cong, một kết luận được củng cố thêm
bởi sự nhất quán của hiệu ứng mức dương của năng lực công nghệ ở các
nghiên cứu đơn quốc gia. Việc đối chiếu chéo giữa hiệu ứng mức ổn định
và hiệu ứng điều tiết không có ý nghĩa, lặp lại qua nhiều nghiên cứu
thành phần và nhiều chế độ thể chế, là chính cấu trúc bằng chứng cho
phép luận án phân biệt một quy luật cấu trúc bền vững khỏi một hiệu ứng
ngẫu nhiên của mẫu.

\hypertarget{tuxf3m-tux1eaft-phux1b0ux1a1ng-phuxe1p-kux1ebf-thux1eeba-vuxe0-ux111uxf3ng-guxf3p-mux1edbi}{%
\section{3.6 Tóm tắt phương pháp: Kế thừa và đóng góp
mới}\label{tuxf3m-tux1eaft-phux1b0ux1a1ng-phuxe1p-kux1ebf-thux1eeba-vuxe0-ux111uxf3ng-guxf3p-mux1edbi}}

\begin{longtable}[]{@{}
  >{\raggedright\arraybackslash}p{(\columnwidth - 4\tabcolsep) * \real{0.3333}}
  >{\raggedright\arraybackslash}p{(\columnwidth - 4\tabcolsep) * \real{0.3333}}
  >{\raggedright\arraybackslash}p{(\columnwidth - 4\tabcolsep) * \real{0.3333}}@{}}
\toprule\noalign{}
\begin{minipage}[b]{\linewidth}\raggedright
Thành phần
\end{minipage} & \begin{minipage}[b]{\linewidth}\raggedright
Cơ sở kế thừa
\end{minipage} & \begin{minipage}[b]{\linewidth}\raggedright
Đóng góp mới
\end{minipage} \\
\midrule\noalign{}
\endhead
\bottomrule\noalign{}
\endlastfoot
Thiết kế hỗn hợp & Borenstein et al.~(2009); Hunter \& Schmidt (2004) &
Áp dụng cho \textbf{50 nền kinh tế châu Á + Thái Bình Dương} trong một
luận án \\
phân tích tổng hợp 1977--2026 & Bausch \& Krist (2007); Kirca et
al.~(2012); Marano et al.~(2016) & Mở rộng phạm vi, bổ sung digital và 6
chế độ con ICRV moderators \\
Mẫu thực nghiệm gộp & nghiên cứu thành phần P1, 17 nền kinh tế châu Á
mới nổi & Mở rộng từ 17 lên \textbf{50 nền kinh tế}, 103 cặp nền kinh tế
và năm (gồm Nhật Bản 2025) \\
Biến phụ thuộc năng suất lao động & Bloom et al.~(2012); Hsieh \& Klenow
(2009); nghiên cứu thành phần P1 & Lập luận rõ về hiệu quả doanh nghiệp
đa chiều trong bối cảnh WBES \\
FSTS + FSTS² & Lu \& Beamish (2004); Hitt et al.~(1997); nghiên cứu
thành phần P2 và P5 & Đưa vào cùng khung điều tiết ba chiều \\
TCI vs DAI & Bhandari et al.~(2023); Verhoef et al.~(2021) & Không chồng
lấn, độ thuần cấu trúc protocol cho WBES; lá chắn số mở rộng (nghiên cứu
thành phần P1) \\
Lind--Mehlum U-test & Lind \& Mehlum (2010); Haans et al.~(2016) & Áp
dụng cho từng country sample và multi-country \\
Three-way điều tiết & Aiken \& West (1991); Dawson (2014) & digital ×
institutional × manager mới cho châu Á và Pacific \\
dị biệt theo thời gian & Wu et al.~(2022); Xiao et al.~(2013); nghiên
cứu thành phần P2 & China 2012--2024 stability test (turning point
\textasciitilde47,8\% FSTS) \\
sai số chuẩn vững & Long \& Ervin (2000); White (1980) & Chuẩn \\
lệch lạc công bố & Egger et al.~(1997); Begg \& Mazumdar (1994) &
Chuẩn \\
\end{longtable}

\hypertarget{ux111ux1ea1o-ux111ux1ee9c-nghiuxean-cux1ee9u-vuxe0-khux1ea3-nux103ng-tuxe1i-lux1eadp}{%
\section{3.7 Đạo đức nghiên cứu và khả năng tái
lập}\label{ux111ux1ea1o-ux111ux1ee9c-nghiuxean-cux1ee9u-vuxe0-khux1ea3-nux103ng-tuxe1i-lux1eadp}}

Luận án sử dụng dữ liệu thứ cấp từ WBES, bộ dữ liệu công khai được World
Bank cung cấp miễn phí cho mục đích nghiên cứu (World Bank Enterprise
Surveys, 2025). Tất cả các trích dẫn theo APA 7th để ghi nhận đầy đủ
công sức của các tác giả trước, tránh đạo văn. Code và dữ liệu phân tích
được lưu trữ trên các nhánh chuyên biệt của kho mã để bảo đảm khả năng
tái lập kết quả.

Về bảo mật và quyền riêng tư, dữ liệu WBES đã được World Bank ẩn danh ở
nguồn: bộ dữ liệu vi mô không chứa tên doanh nghiệp hay thông tin nhận
dạng cá nhân của đáp viên, và mọi định danh chỉ là mã số nội bộ phục vụ
ghép nối. Luận án không thu thập dữ liệu sơ cấp từ con người nên không
phát sinh nghĩa vụ xin chấp thuận đạo đức từ hội đồng đối với đối tượng
nghiên cứu là người; tuy vậy luận án vẫn tuân thủ điều khoản sử dụng của
World Bank, không thực hiện bất kỳ thao tác nào nhằm tái nhận dạng doanh
nghiệp, và chỉ báo cáo kết quả ở cấp tổng hợp hoặc cấp hệ số. Việc xử lý
mã không phản hồi của bộ công cụ, gồm mã không biết, mã từ chối và mã
không áp dụng, thành giá trị khuyết được thực hiện minh bạch theo quy
ước nêu ở Phụ lục A nhằm tránh diễn giải sai các mã kỹ thuật thành giá
trị thực.

Về khả năng tái lập, luận án áp dụng nguyên tắc minh bạch kiểu PRISMA
cho dữ liệu thứ cấp: toàn bộ quy trình hợp nhất sáu bước, từ lọc loại
khảo sát đến tạo mẫu phân tích bằng xóa theo danh sách, được ghi rõ tiêu
chí và số quan sát còn lại ở từng bước để cho phép tái dựng và kiểm
chứng, cùng dòng dữ liệu kiểu PRISMA và mã tái lập trình bày đầy đủ ở
Phụ lục A. Các con số trọng tâm của khung thực nghiệm được khóa và có
thể tái lập đúng, gồm mô hình bậc hai gốc với cỡ mẫu hồi quy 81.022 và
điểm uốn 51,5\%, cùng mô hình đầy đủ thêm năng lực công nghệ và chấp
nhận số với cỡ mẫu 79.080 và điểm uốn 43,6\%. Tập lệnh xây dựng dữ liệu
gộp tái lập đúng quy trình sáu bước S1 đến S6 trên kho mã hiện hành, bảo
đảm rằng mỗi cập nhật bổ sung nền kinh tế hay đợt khảo sát mới, chẳng
hạn Nhật Bản 2025, đều đi qua cùng một chuỗi quyết định lọc và hài hòa
hóa có thể kiểm chứng.

\newpage

\hypertarget{chux1b0ux1a1ng-4-kux1ebft-quux1ea3-nghiuxean-cux1ee9u}{%
\chapter{CHƯƠNG 4: KẾT QUẢ NGHIÊN
CỨU}\label{chux1b0ux1a1ng-4-kux1ebft-quux1ea3-nghiuxean-cux1ee9u}}

\begin{center}\rule{0.5\linewidth}{0.5pt}\end{center}

\hypertarget{giux1edbi-thiux1ec7u-chux1b0ux1a1ng}{%
\section{Giới thiệu chương}\label{giux1edbi-thiux1ec7u-chux1b0ux1a1ng}}

Chương 4 trình bày toàn bộ kết quả nghiên cứu của luận án theo logic đi
từ bức tranh thực trạng mô tả (Mục 4.1) đến bằng chứng tổng hợp ở cấp
phân tích tổng hợp (Mục 4.2), sau đó đến bằng chứng quốc gia cụ thể và
cuối cùng là kiểm định toàn mẫu đa quốc gia. Cấu trúc trình bày gồm mười
phần chính, bắt đầu bằng thực trạng quốc tế hóa và hiệu quả hoạt động
kinh doanh tại châu Á theo các phân nhóm con thể chế ICRV, tiếp theo là:
(ii) kết quả phân tích tổng hợp ba tầng (P6); (iii) bối cảnh thể chế
tiên tiến--đổi mới tại Singapore (P4); (iv) bối cảnh thể chế trung bình
cao tại Trung Quốc (P5); (v) bối cảnh chuyển đổi bậc thấp--trung tại
Việt Nam (P3); (vi) kiểm định toàn mẫu châu Á đa bối cảnh (P7); (vii)
trường hợp biên là các nền kinh tế đảo nhỏ đang phát triển ở Thái Bình
Dương (P8); và (viii) tổng hợp kết quả theo hệ giả thuyết.

\begin{center}\rule{0.5\linewidth}{0.5pt}\end{center}

\hypertarget{thux1ef1c-trux1ea1ng-quux1ed1c-tux1ebf-huxf3a-vuxe0-hiux1ec7u-quux1ea3-houx1ea1t-ux111ux1ed9ng-kinh-doanh-tux1ea1i-chuxe2u-uxe1}{%
\section{4.1 Thực trạng quốc tế hóa và hiệu quả hoạt động kinh doanh tại
châu
Á}\label{thux1ef1c-trux1ea1ng-quux1ed1c-tux1ebf-huxf3a-vuxe0-hiux1ec7u-quux1ea3-houx1ea1t-ux111ux1ed9ng-kinh-doanh-tux1ea1i-chuxe2u-uxe1}}

\hypertarget{bux1ee9c-tranh-tux1ed5ng-thux1ec3-vux1ec1-hiux1ec7u-quux1ea3-doanh-nghiux1ec7p-chuxe2u-uxe1}{%
\subsection{4.1.1 Bức tranh tổng thể về hiệu quả doanh nghiệp châu
Á}\label{bux1ee9c-tranh-tux1ed5ng-thux1ec3-vux1ec1-hiux1ec7u-quux1ea3-doanh-nghiux1ec7p-chuxe2u-uxe1}}

Phân tích mô tả trên khung 50 nền kinh tế châu Á và Thái Bình Dương (gồm
Nhật Bản, thành viên thứ sáu của Nhóm I được WBES khảo sát lần đầu năm
2025; xem Mục 4.1.5) cho thấy bức tranh hiệu quả hoạt động kinh doanh có
sự phân tán lớn và không hội tụ. Các bảng phân loại 4.1 dưới đây dựa
trên bộ dữ liệu ước lượng đa quốc gia (pool phân loại 96.415 doanh
nghiệp / 52 nhãn nền; khung phân tích 88.869 / 50 nền, gồm Nhật Bản; mẫu
hồi quy chính P7 M2 N = 81.022), nhất quán với thống kê mô tả chi tiết
của Chuyên đề 1 trên đầy đủ 50 nền. Vì doanh thu WBES được báo cáo bằng
nội tệ, độ phân tán năng suất được tính \textbf{trong từng cặp nền kinh
tế × năm} (bất biến với đơn vị tiền tệ): độ lệch chuẩn của ln(doanh
thu/lao động) trong đợt tăng từ 1,04 (Nhóm I, Advanced đổi mới) lên
1,31--1,39 (Nhóm III--V), và tỷ số P90/P10 leo từ khoảng 12 lần (Nhóm I)
lên khoảng 30 lần (Nhóm V) theo chiều từ nhóm thể chế cao đến nhóm thể
chế thấp.

\textbf{Bảng 4.1.} \emph{Phân tán năng suất lao động và đặc điểm quốc tế
hóa theo phân nhóm con ICRV (mẫu gộp WBES 2006--2026; số doanh nghiệp
phân loại ICRV ≈ 96.415; khung phân tích 88.869/50 nền; mẫu hồi quy
chính P7 M2 N = 81.022).}

\begin{longtable}[]{@{}
  >{\raggedright\arraybackslash}p{(\columnwidth - 10\tabcolsep) * \real{0.1667}}
  >{\raggedright\arraybackslash}p{(\columnwidth - 10\tabcolsep) * \real{0.1667}}
  >{\raggedright\arraybackslash}p{(\columnwidth - 10\tabcolsep) * \real{0.1667}}
  >{\raggedright\arraybackslash}p{(\columnwidth - 10\tabcolsep) * \real{0.1667}}
  >{\raggedright\arraybackslash}p{(\columnwidth - 10\tabcolsep) * \real{0.1667}}
  >{\raggedright\arraybackslash}p{(\columnwidth - 10\tabcolsep) * \real{0.1667}}@{}}
\toprule\noalign{}
\begin{minipage}[b]{\linewidth}\raggedright
Phân nhóm con ICRV
\end{minipage} & \begin{minipage}[b]{\linewidth}\raggedright
n
\end{minipage} & \begin{minipage}[b]{\linewidth}\raggedright
sd log LP (trong đợt)¹
\end{minipage} & \begin{minipage}[b]{\linewidth}\raggedright
FSTS (\%)
\end{minipage} & \begin{minipage}[b]{\linewidth}\raggedright
Xuất khẩu (\%)
\end{minipage} & \begin{minipage}[b]{\linewidth}\raggedright
FDI≥10\% (\%)
\end{minipage} \\
\midrule\noalign{}
\endhead
\bottomrule\noalign{}
\endlastfoot
Nhóm I: Advanced đổi mới (SG, HK, KOR, TWN, ISR, JPN) & 6.390 &
\textbf{1,00} & 10,3 & 24,5 & 7,7 \\
Nhóm II: Advanced tài nguyên (GCC) & 2.269 & 1,14² & \textasciitilde3 &
\textasciitilde8 & n/a \\
Nhóm III: trung bình cao (CHN, MYS, THA, KAZ\ldots) & 17.905 & 1,37 &
10,3 & 21,7 & 8,4 \\
Nhóm IV: chuyển đổi thu nhập trung bình thấp (VNM, IND, IDN, PHL, BGD,
PAK, MNG) & 50.926 & 1,31 & 8,6 & 15,5 & 4,7 \\
Nhóm V: đang nổi (LAO, KHM, NPL, LKA, JOR\ldots) & 18.569 &
\textbf{1,39} & 10,1 & 16,6 & 5,9 \\
Nhóm VI: SIDS (FJI, PNG, SLB\ldots) & 2.038 & 1,32 & 6,3 & 16,3 &
23,5 \\
\textbf{Tổng (phân loại ICRV)} & \textbf{96.415} & n/a & n/a & n/a &
n/a \\
\end{longtable}

\emph{Nguồn: Phân tích của tác giả từ World Bank Enterprise Surveys
(WBES), 2006--2026.}

\emph{Ghi chú: Cột n là số doanh nghiệp được phân loại vào sáu nhóm ICRV
trong mẫu gộp WBES (tổng 96.415; ngoài ra còn \textasciitilde10.350
doanh nghiệp chưa gán nhóm ICRV do thiếu chỉ số thể chế). Mẫu hồi quy
chính của P7 là N = 81.022 (mô hình M2, khung 50 nền gồm Nhật Bản) và
giảm còn 79.080 ở M5 khi bổ sung biến kiểm soát (xem Mục 4.6.1), do loại
các quan sát thiếu biến phụ thuộc/FSTS. ¹ Độ lệch chuẩn ln(năng suất lao
động) tính trong từng cặp nền kinh tế × năm rồi lấy trung vị giữa các
đợt của nhóm, bất biến với đơn vị tiền tệ (thừa số quy đổi triệt tiêu
trên thang log); năng suất được winsorize 1/99 trong cụm country-year
(xem Mục 3.5.3). ² Trung vị Nhóm II bị kéo lên bởi các nền nhỏ; hai nền
lớn nhất khối có phân tán hẹp nhất toàn pool (Saudi Arabia 0,49; Qatar
0,33).}

\emph{Quy ước nhãn nhóm ICRV (căn theo nhãn dữ liệu P7): Nhóm I =
Advanced innovation-driven; Nhóm II = tiên tiến tài nguyên; Nhóm III =
trung bình cao; \textbf{Nhóm IV = chuyển đổi thu nhập trung bình thấp}
(\texttt{Lower\_mid\_transition}: Bangladesh, Ấn Độ, Indonesia, Mông Cổ,
Pakistan, Philippines, Việt Nam); \textbf{Nhóm V = đang nổi}
(\texttt{Emerging}: 17 nền gồm Campuchia, Lào, Nepal, Sri Lanka,
Jordan\ldots); Nhóm VI = SIDS (\texttt{SIDS\_small}). Hệ nhãn này được
dùng thống nhất trong toàn bộ luận án và hai chuyên đề.}

Phát hiện quan trọng nhất từ thực trạng mô tả là sự phân biệt nội bộ
trong nhóm Advanced. Hai phân nhóm cùng mức thu nhập cao nhưng đối lập
về cấu trúc: nhóm đổi mới sáng tạo dẫn dắt có R\&D 21,0\%, ISO 32,5\% và
24,5\% doanh nghiệp xuất khẩu, trong khi nhóm tài nguyên dẫn dắt chỉ có
R\&D 7,0\% và 11,7\% doanh nghiệp xuất khẩu (Chuyên đề 1). Về phân phối,
hai nền lớn nhất Vùng Vịnh có phân tán năng suất trong đợt hẹp nhất toàn
mẫu gộp (Saudi Arabia 0,49; Qatar 0,33, so với Singapore 1,08, chênh hơn
2 lần), biểu hiện của kinh tế nhà nước tô (rentier economy) nơi phân
phối lại từ xuất khẩu dầu mỏ thu hẹp khoảng cách năng suất giữa doanh
nghiệp, nhưng không nhất thiết phản ánh hiệu quả doanh nghiệp đáng kể.
Điều này gợi ý rằng việc gộp hai loại Advanced vào một nhóm duy nhất
(như các nghiên cứu trước thường làm) sẽ che khuất cơ chế quan trọng
này.

Sự gia tăng của độ phân tán năng suất trong đợt theo gradient ICRV (từ
\textasciitilde12 lần lên \textasciitilde30 lần ở tỷ số P90/P10) không
chỉ là một quan sát mô tả mà còn nhất quán với khung lý thuyết phân bổ
sai nguồn lực (resource misallocation) của Hsieh và Klenow (2009): trong
môi trường thể chế yếu, các rào cản thị trường, tiếp cận tài chính hạn
chế, thực thi hợp đồng kém, méo mó chính sách, ngăn nguồn lực dịch
chuyển từ doanh nghiệp kém hiệu quả sang doanh nghiệp hiệu quả hơn,
khiến phương sai năng suất giữa các doanh nghiệp nới rộng. Theo logic
này, độ phân tán năng suất cao ở Nhóm V--VI \textbf{có thể} là dấu hiệu
của tổn thất phân bổ (allocative loss) thay vì phản ánh sự đa dạng lành
mạnh. Cần lưu ý rằng độ lệch chuẩn của log năng suất lao động chỉ là
\textbf{biến đại diện gián tiếp} cho cơ chế phân tán TFPR (năng suất
doanh thu) mà Hsieh và Klenow (2009) dùng để định danh phân bổ sai nguồn
lực, vì nó còn chịu ảnh hưởng của dị biệt công nghệ, mức thâm dụng vốn
và điều kiện cầu; do đó đây là một cách diễn giải \emph{nhất quán với}
khung phân bổ sai nguồn lực chứ không phải bằng chứng phân rã trực tiếp.
Dù vậy, hướng diễn giải này có hàm ý chính sách khác hẳn với việc coi
phân tán chỉ là nhiễu thống kê. Cần nhấn mạnh rằng các so sánh giữa các
nhóm ICRV trong Mục 4.1 mang bản chất \textbf{mô tả}; xác nhận
\textbf{suy diễn} rằng khác biệt giữa các chế độ thể chế có ý nghĩa
thống kê được cung cấp ở hai cấp độ độc lập trong các phần sau của
chương: kiểm định điều tiết ICRV của phân tích tổng hợp P6
(\(Q_M = 17{,}35\), \(df = 4\), \(p = {,}002\); Mục 4.2.3) và gradient
điểm uốn theo ICRV của mô hình hồi quy đa quốc gia P7 (Mục 4.6.2). Hai
nguồn bằng chứng này bảo đảm bức tranh thực trạng mô tả được neo vào
bằng chứng suy diễn nhất quán, thay vì dừng lại ở so sánh trung bình--độ
lệch chuẩn.

\begin{figure}
\centering
\includegraphics{figures/fig_4_1_pool_composition.png}
\caption{Hình 4.1. Thành phần mẫu gộp theo nhóm ICRV}
\end{figure}

\begin{figure}
\centering
\includegraphics{figures/fig_4_0_pool_by_year.png}
\caption{Hình 4.2. Quy mô mẫu gộp theo năm khảo sát}
\end{figure}

\hypertarget{thux1ef1c-trux1ea1ng-quux1ed1c-tux1ebf-huxf3a-phuxe2n-cux1ef1c-tham-gia-xuux1ea5t-khux1ea9u}{%
\subsection{4.1.2 Thực trạng quốc tế hóa: phân cực tham gia xuất
khẩu}\label{thux1ef1c-trux1ea1ng-quux1ed1c-tux1ebf-huxf3a-phuxe2n-cux1ef1c-tham-gia-xuux1ea5t-khux1ea9u}}

Trung vị FSTS bằng 0\% trong tất cả các phân nhóm, chỉ 15--23\% doanh
nghiệp tham gia xuất khẩu tùy theo bối cảnh. Phân phối FSTS có dạng phân
cực mạnh: đại đa số không xuất khẩu, trong khi một thiểu số nhỏ có tỷ lệ
xuất khẩu rất cao. Phân cực này đặt ra vấn đề phương pháp quan trọng cho
nghiên cứu quốc tế hóa và hiệu quả: kết quả hồi quy trên toàn mẫu bị kéo
bởi hành vi của thiểu số xuất khẩu tích cực.

Cường độ quốc tế hóa (FSTS trung bình, tính trên toàn mẫu gồm cả doanh
nghiệp không xuất khẩu) dao động từ 6,3\% (SIDS) đến 10,3\% (trung bình
cao), với mức trung bình khu vực khoảng 8,8\%. Điều đáng chú ý là trung
bình cao và Advanced có cùng FSTS trung bình (\textasciitilde10\%) nhưng
cơ chế xuất khẩu hoàn toàn khác: Advanced chủ yếu dịch vụ và hàng hóa
hàm lượng tri thức cao; trung bình cao chủ yếu sản xuất theo chuỗi giá
trị toàn cầu với biên lợi nhuận khác nhau.

\begin{figure}
\centering
\includegraphics{figures/fig_4_4_growth_pathway.png}
\caption{Hình 4.3. Quỹ đạo tăng trưởng theo nhóm ICRV}
\end{figure}

\hypertarget{thux1ef1c-trux1ea1ng-nux103ng-lux1ef1c-sux1ed1-vuxe0-cuxf4ng-nghux1ec7}{%
\subsection{4.1.3 Thực trạng năng lực số và công
nghệ}\label{thux1ef1c-trux1ea1ng-nux103ng-lux1ef1c-sux1ed1-vuxe0-cuxf4ng-nghux1ec7}}

Phân tích năm thành phần năng lực (đổi mới sản phẩm, đổi mới quy trình,
R\&D, ISO, website) theo phân nhóm con ICRV cho thấy sự đa dạng đáng kể.

\textbf{Bảng 4.2.} \emph{Đổi mới sáng tạo và chấp nhận số theo phân nhóm
con ICRV (\%).}

\begin{longtable}[]{@{}
  >{\raggedright\arraybackslash}p{(\columnwidth - 10\tabcolsep) * \real{0.1667}}
  >{\raggedright\arraybackslash}p{(\columnwidth - 10\tabcolsep) * \real{0.1667}}
  >{\raggedright\arraybackslash}p{(\columnwidth - 10\tabcolsep) * \real{0.1667}}
  >{\raggedright\arraybackslash}p{(\columnwidth - 10\tabcolsep) * \real{0.1667}}
  >{\raggedright\arraybackslash}p{(\columnwidth - 10\tabcolsep) * \real{0.1667}}
  >{\raggedright\arraybackslash}p{(\columnwidth - 10\tabcolsep) * \real{0.1667}}@{}}
\toprule\noalign{}
\begin{minipage}[b]{\linewidth}\raggedright
Nhóm ICRV (nhãn dữ liệu)
\end{minipage} & \begin{minipage}[b]{\linewidth}\raggedright
Sản phẩm mới
\end{minipage} & \begin{minipage}[b]{\linewidth}\raggedright
Quy trình mới
\end{minipage} & \begin{minipage}[b]{\linewidth}\raggedright
R\&D
\end{minipage} & \begin{minipage}[b]{\linewidth}\raggedright
ISO
\end{minipage} & \begin{minipage}[b]{\linewidth}\raggedright
Website (DAI Tầng 1)
\end{minipage} \\
\midrule\noalign{}
\endhead
\bottomrule\noalign{}
\endlastfoot
I, Adv. đổi mới & 26,9 & 18,8 & \textbf{21,0} & 32,5 & \textbf{69,2} \\
II, Adv. tài nguyên (GCC) & 17,0 & 8,6 & 7,0 & 15,2 & 50,2 \\
III, Trung bình cao & 26,7 & 18,9 & 21,5 & 25,4 & 53,9 \\
IV, Chuyển đổi TB-thấp & 17,1 & 14,7 & 15,5 & 25,3 & 48,3 \\
V, Đang nổi & 24,4 & 16,3 & 17,5 & 12,7 & 40,9 \\
VI, SIDS & \textbf{34,2} & \textbf{24,7} & 10,2 & 12,2 & 41,5 \\
\end{longtable}

\emph{Nguồn: Phân tích của tác giả từ WBES, 2006--2026.}

\emph{Ghi chú: Bảng 4.2 được tính trực tiếp từ dữ liệu thô WBES trên
khung 50 nền (một cross-section chuẩn cho mỗi cặp nền kinh tế × năm, đợt
khảo sát ≥ 2006), với các chỉ báo chuẩn của bộ công cụ WBES toàn cầu:
đổi mới sản phẩm (h1), đổi mới quy trình (h5), chi R\&D (h8), chứng nhận
chất lượng quốc tế (b8) và website (c22b). Toàn bộ sáu nhóm ICRV đều có
độ phủ dữ liệu đầy đủ, kể cả Nhóm II (đủ 6/6 nền GCC). Gradient năng lực
thể chế thể hiện rõ ở các biến nền tảng: R\&D 21,0\% (Nhóm I) xuống
10,2\% (Nhóm VI); ISO 32,5\% xuống 12,2\%; website 69,2\% xuống 41,5\%.
Bộ số này nhất quán với phân tích thực trạng ở Chuyên đề 1.}

Phát hiện đáng chú ý là khuôn mẫu ``thích nghi và nhảy vọt'' (adaptive
leapfrogging) ở SIDS: với dữ liệu raw đủ 8/8 nền Nhóm VI, SIDS có
\textbf{tỷ lệ đổi mới sản phẩm cao nhất mọi nhóm (34,2\%)} và đổi mới
quy trình cao nhất (24,7\%), trong khi website (41,5\%) ngang nhóm Đang
nổi (40,9\%) dù GNI thấp hơn đáng kể, phản ánh đổi mới bằng nguồn lực
hạn chế (frugal innovation) như công cụ sinh tồn chứ không phải tối ưu
hóa hiệu quả. Kết quả này hỗ trợ lập luận lý thuyết về sự phân biệt giữa
DAI (chấp nhận số bề mặt) và TCI (năng lực công nghệ chiều sâu): SIDS
đổi mới/áp dụng số ở bề mặt cao nhưng năng lực công nghệ chiều sâu (R\&D
10,2\%, ISO 12,2\%) thấp nhất, và đây là nhóm duy nhất có quan hệ quốc
tế hóa và hiệu quả âm đơn điệu (FIP).

Một biểu hiện khác của bão hòa Tầng 1 nằm ở tương quan giữa website và
năng suất: phần bù tương quan của DAI Tầng 1 thấp nhất ở nhóm Advanced
đổi mới (+0,090, so với +0,180--0,201 ở các nhóm chuyển đổi và tài
nguyên, Chuyên đề 1, Bảng 2.3.8.1) dù tỷ lệ có website của nhóm này cao
(69,2\%). Khi một chỉ báo nhị phân tiệm cận bão hòa, phương sai của nó
co lại và nhóm ``không website'' còn lại mang tính chọn lọc đặc thù, nên
hệ số đo được phản ánh giới hạn của thước đo Tầng 1 nhiều hơn là giới
hạn của số hóa, một vấn đề đo lường cần phân biệt với quan hệ nhân quả
thực.

\begin{figure}
\centering
\includegraphics{figures/fig_4_6_firm_size.png}
\caption{Hình 4.4. Phân bố quy mô doanh nghiệp theo nhóm ICRV}
\end{figure}

\hypertarget{bux1ed1n-ruxe0o-cux1ea3n-cux1ea5u-truxfac-ux1ea3nh-hux1b0ux1edfng-ux111ux1ebfn-quux1ed1c-tux1ebf-huxf3a-vuxe0-hiux1ec7u-quux1ea3}{%
\subsection{4.1.4 Bốn rào cản cấu trúc ảnh hưởng đến quốc tế hóa và hiệu
quả}\label{bux1ed1n-ruxe0o-cux1ea3n-cux1ea5u-truxfac-ux1ea3nh-hux1b0ux1edfng-ux111ux1ebfn-quux1ed1c-tux1ebf-huxf3a-vuxe0-hiux1ec7u-quux1ea3}}

Phân tích thực trạng xác định bốn rào cản cấu trúc phổ biến nhất ảnh
hưởng đến khả năng quốc tế hóa và hiệu quả doanh nghiệp tại châu Á: (i)
tiếp cận tài chính hạn chế, đặc biệt nghiêm trọng ở đang nổi và chuyển
đổi thu nhập trung bình thấp (Nhóm V--IV); (ii) thiếu hụt lao động có kỹ
năng; (iii) cạnh tranh từ doanh nghiệp phi chính thức; và (iv) nguồn
cung điện không đáng tin cậy. Bốn rào cản này hình thành ``khoảng trống
thể chế'' (Khanna \& Palepu, 2010) tạo ra chi phí giao dịch bổ sung cho
doanh nghiệp khi mở rộng hoạt động ra ngoài biên giới. Sự hiện diện đồng
thời của nhiều khoảng trống tại đang nổi và SIDS giải thích tại sao quan
hệ quốc tế hóa và hiệu quả ở các nhóm này yếu hoặc âm.

\hypertarget{nhux1eadt-bux1ea3n-2025-lux1ea7n-ux111ux1ea7u-tiuxean-ux111ux1b0ux1ee3c-khux1ea3o-suxe1t-vuxe0-vux1ecb-truxed-trong-khung-icrv}{%
\subsection{4.1.5 Nhật Bản 2025: Lần đầu tiên được khảo sát và vị trí
trong khung
ICRV}\label{nhux1eadt-bux1ea3n-2025-lux1ea7n-ux111ux1ea7u-tiuxean-ux111ux1b0ux1ee3c-khux1ea3o-suxe1t-vuxe0-vux1ecb-truxed-trong-khung-icrv}}

Một sự kiện dữ liệu đáng chú ý của giai đoạn nghiên cứu là việc Nhật Bản
lần đầu tiên được đưa vào Khảo sát Doanh nghiệp của Ngân hàng Thế giới
năm 2025, với mẫu \(N = 2.168\) doanh nghiệp. Nhật Bản là \textbf{thành
viên thứ sáu của Nhóm I} (Advanced đổi mới sáng tạo dẫn dắt) và là nền
kinh tế tiên tiến lớn cuối cùng gia nhập bộ dữ liệu WBES, cung cấp một
điểm neo quý giá cho biên trên của gradient thể chế. Nhật Bản được tích
hợp đầy đủ vào khung mô tả 50 nền của chuyên đề thực trạng (Chuyên đề 1)
và mọi thống kê Nhóm I trong luận án. Về phía ước lượng kinh tế lượng,
\textbf{điểm uốn toàn mẫu P7 đã được tái ước lượng trên khung 50 nền bao
gồm Nhật Bản} (2.168 doanh nghiệp Nhật trong khung phân tích 88.869; mẫu
hồi quy chính M2 N = 81.022, M5 N = 79.080; xem Mục 4.6), nên các kết
quả P7 trong luận án đã phản ánh Nhật Bản. Hai nghiên cứu thành phần còn
lại là mẫu gộp con riêng không liên quan Nhật Bản: P8 (FIP) trên 7 nền
đảo nhỏ Thái Bình Dương và P9 (ngưỡng) trên Ấn Độ. Ngoài ra, nghiên cứu
thành phần P10 kiểm định chuyên sâu quan hệ quốc tế hóa và hiệu quả tại
Nhật Bản (xem Mục 4.9). Toàn bộ số liệu dưới đây được tính trực tiếp từ
tệp vi mô gốc.

\textbf{Bảng 4.3.} \emph{Đặc trưng doanh nghiệp Nhật Bản 2025 so với năm
nền còn lại của Nhóm I (Advanced đổi mới).}

\begin{longtable}[]{@{}lrr@{}}
\toprule\noalign{}
Chỉ tiêu & Nhật Bản 2025 & 5 nền Nhóm I còn lại \\
\midrule\noalign{}
\endhead
\bottomrule\noalign{}
\endlastfoot
n doanh nghiệp & 2.168 & 4.222 \\
sd ln(LP) trong đợt & \textbf{0,89} & 1,04 \\
P90/P10 năng suất & \textbf{8,3} & 12,2 \\
FSTS trung bình (\%) & 4,1 & 13,4 \\
Doanh nghiệp xuất khẩu (\%) & 16,5 & 28,4 \\
FDI ≥ 10\% (\%) & \textbf{1,9} & 10,6 \\
Website (\%) & \textbf{83,8} & 61,7 \\
R\&D (\%) & 21,9 & 20,5 \\
ISO (\%) & 27,7 & 35,0 \\
Đổi mới sản phẩm (\%) & 32,2 & 24,3 \\
Nữ quản lý cấp cao (\%) & \textbf{7,3} & 27,5 \\
Tuổi doanh nghiệp trung bình (năm) & \textbf{50,4} & 23,0 \\
\end{longtable}

\emph{Nguồn: tính trực tiếp từ \texttt{Japan-2025-full-data.dta} (WBES)
và pool mô tả của năm nền còn lại Nhóm I. FSTS = d3b + d3c.}

Bốn đặc trưng nổi bật định vị Nhật Bản trong khung ICRV. Thứ nhất,
\textbf{phân tán năng suất hẹp nhất}: độ lệch chuẩn ln năng suất trong
đợt của Nhật (0,89) thấp hơn mức trung vị của năm nền còn lại Nhóm I
(1,04), và tỷ số P90/P10 chỉ 8,3 lần so với 12,2 lần, phản ánh một nền
kinh tế tiên tiến trưởng thành với sàn năng suất cao và đồng đều. Đặc
điểm này củng cố trực tiếp logic ``bão hòa Tầng 1'' của khung CDCM: ở
nền kinh tế đồng đều và số hóa cao, dư địa để hiện diện số cơ bản tạo
khác biệt năng suất là rất hẹp. Thứ hai, \textbf{số hóa trưởng thành
nhất}: tỷ lệ có website 83,8\% là cao nhất trong toàn bộ mẫu gộp, vượt
xa mức 61,7\% của năm nền còn lại Nhóm I. Thứ ba, \textbf{trường thọ
doanh nghiệp}: tuổi trung bình 50,4 năm (so với 23,0 năm của năm nền còn
lại Nhóm I) phản ánh cấu trúc doanh nghiệp lâu đời đặc thù của Nhật Bản.
Thứ tư, \textbf{nghịch lý đổi mới--quốc tế hóa và khoảng cách giới}: dù
cường độ R\&D và đổi mới cao, cường độ quốc tế hóa cấp cơ sở lại thấp
(FSTS 4,1\%) cùng tỷ lệ sở hữu nước ngoài rất thấp (1,9\%), cho thấy
hoạt động quốc tế hóa tập trung ở các tập đoàn đa quốc gia lớn chứ không
phân bố đều ở mẫu cơ sở; đồng thời tỷ lệ nữ quản lý cấp cao chỉ 7,3\%,
thấp hơn nhiều so với mức dẫn đầu khu vực Đông Á--Thái Bình Dương, nhất
quán với tài liệu về trần kính trong quản trị doanh nghiệp Nhật Bản. Các
đặc trưng này được khai thác đầy đủ trong nghiên cứu thành phần P10,
kiểm định chuyên sâu quan hệ quốc tế hóa và hiệu quả tại Nhật Bản: kết
quả xác nhận quan hệ gần tuyến tính dương (hiệu ứng tuyến tính +0,671, p
\textless{} 0,001; số hạng bậc hai không ý nghĩa; điểm uốn đại số
90--108\% nằm ngoài miền dữ liệu), TCI (+0,100, p \textless{} 0,001) và
DAI (+0,110, p = 0,028) nâng mặt bằng không uốn đường cong, cùng phần bù
trường thọ doanh nghiệp (+0,073, p = 0,007), đúng như dự đoán biên trên
của khung ICRV (trình bày đầy đủ ở Mục 4.9).

\begin{center}\rule{0.5\linewidth}{0.5pt}\end{center}

\hypertarget{so-suxe1nh-xuyuxean-nhuxf3m-icrv-gradient-tux1b0ux1a1ng-quan-vuxe0-ux111ux1ed9-phuxe2n-tuxe1n-hiux1ec7u-quux1ea3}{%
\subsection{4.1.6 So sánh xuyên nhóm ICRV: gradient tương quan và độ
phân tán hiệu
quả}\label{so-suxe1nh-xuyuxean-nhuxf3m-icrv-gradient-tux1b0ux1a1ng-quan-vuxe0-ux111ux1ed9-phuxe2n-tuxe1n-hiux1ec7u-quux1ea3}}

Khi đặt cạnh nhau toàn bộ các chỉ báo mô tả theo sáu nhóm ICRV, một bức
tranh có cấu trúc nổi lên trong đó cường độ và đôi khi cả dấu của các
tương quan thô biến thiên theo chiều suy yếu của thể chế. Hệ số tương
quan Pearson giữa năng suất lao động chuẩn hóa và cường độ quốc tế hóa
giảm đơn điệu theo gradient thể chế: từ \(+0{,}139\) ở Nhóm I, qua
\(+0{,}131\) ở Nhóm II, \(+0{,}094\) ở Nhóm III, \(+0{,}047\) ở Nhóm IV,
\(+0{,}018\) ở Nhóm V, rồi mất ý nghĩa thống kê và đổi dấu ở SIDS
(\(-0{,}009\), không có ý nghĩa). Mô thức này cho thấy quan hệ thô giữa
quốc tế hóa và năng suất gần như biến mất ở các môi trường thể chế yếu,
gợi ý một quan hệ phi tuyến và bị điều tiết thay vì một hiệu ứng tuyến
tính phổ quát; SIDS là ngoại lệ duy nhất về dấu, nhất quán với giả
thuyết chi phí buộc phải quốc tế hóa, theo đó doanh nghiệp đảo nhỏ xuất
khẩu vì bắt buộc chứ không phải vì được chọn lọc theo năng suất.

\begin{figure}
\centering
\includegraphics{figures/fig_icrv_capability_gradient.png}
\caption{Hình 4.5. Gradient năng lực công nghệ theo sáu nhóm ICRV}
\end{figure}

Ngược với gradient của quốc tế hóa, năng lực công nghệ là yếu tố tương
quan mạnh và đồng đều nhất ở mọi nhóm, dao động từ \(+0{,}117\) ở SIDS
đến \(+0{,}240\) ở Nhóm III, củng cố cách đọc rằng năng lực công nghệ
nền tảng là một yếu tố dịch chuyển mức năng suất có tính phổ quát chứ
không phụ thuộc chế độ thể chế. Tương quan của sở hữu nước ngoài dương ở
mọi nhóm nhưng mạnh hơn ở hai cực của phổ thể chế (Nhóm I \(+0{,}104\)
và SIDS \(+0{,}122\)), phản ánh hai cơ chế đầu tư trực tiếp khác nhau đã
nêu ở Mục 4.1.1: trung tâm khu vực của doanh nghiệp đa quốc gia tại nhóm
tiên tiến so với đầu tư ngoại sinh hướng du lịch và viễn thông tại các
đảo nhỏ.

\begin{figure}
\centering
\includegraphics{figures/fig_4_2_productivity_density.png}
\caption{Hình 4.6. Phân bố mật độ năng suất lao động theo nhóm ICRV}
\end{figure}

Độ phân tán hiệu quả cũng thể hiện gradient thể chế rõ rệt khi đo bằng
các tỷ số phân vị bất biến với đơn vị tiền tệ. Tỷ số liên tứ phân vị
\(P75/P25\) của năng suất tăng từ \(3{,}4\) lần ở Nhóm I lên \(5{,}8\)
lần ở Nhóm V, một mô thức song hành với tỷ số \(P90/P10\) đã trình bày ở
Bảng 4.1. Bổ sung cho thước đo phân tán, mức sinh lời cũng mỏng dần theo
chiều thể chế yếu: tỷ suất sinh lời trên doanh thu trung vị giảm từ
\(0{,}50\) ở Nhóm I xuống \(0{,}45\) ở Nhóm III rồi \(0{,}35\) ở Nhóm
IV, nhất quán với lập luận của lý thuyết thể chế rằng chi phí giao dịch
cao trong môi trường thể chế yếu bào mòn biên lợi nhuận của doanh
nghiệp.

\begin{figure}
\centering
\includegraphics{figures/fig_4_7_spider_chart.png}
\caption{Hình 4.7. Biểu đồ radar năng lực và quốc tế hóa theo nhóm ICRV}
\end{figure}

Một điểm tương phản đáng chú ý về số hóa nằm ở chính nhóm tiên tiến.
Phần bù tương quan của chấp nhận số bề mặt (website, thuộc Tầng 1) lại
thấp nhất ở Nhóm I (\(+0{,}090\)), trong khi cao nhất ở Nhóm II
(\(+0{,}201\)) và Nhóm IV (\(+0{,}180\)), dù tỷ lệ có website ở Nhóm I
đã đạt \(69{,}2\%\). Đây là biểu hiện trực tiếp của hiện tượng bão hòa
thước đo Tầng 1: khi một chỉ báo nhị phân tiệm cận trần, phương sai của
nó co lại và lợi suất biên của hiện diện số cơ bản giảm dần, một vấn đề
đo lường cần phân biệt với giới hạn của số hóa thực chất. Cuối cùng, về
cấu trúc lãnh đạo, khu vực Đông Á và Thái Bình Dương dẫn đầu toàn cầu về
tỷ lệ nữ trong nhóm quản trị cấp cao (\(33{,}4\%\)) lẫn tỷ lệ nữ sở hữu
đa số (\(25{,}7\%\)), bỏ xa khu vực Trung Đông và Bắc Phi (tương ứng
\(3{,}3\%\) và \(1{,}7\%\)), một khác biệt thể chế kép gắn với bối cảnh
nhà nước tô ở khối tài nguyên dẫn dắt và làm phong phú thêm bức tranh đa
dạng nội bộ của các nhóm tiên tiến.

\hypertarget{bux1eb1ng-chux1ee9ng-nux1ec1n-tux1ea3ng-ux111a-quux1ed1c-gia-phuxe2n-tux1ea7ng-nux103ng-suux1ea5t-vuxf9ng-vuxe0-hiux1ec7u-ux1ee9ng-luxe1-chux1eafn-sux1ed1-p1}{%
\subsection{4.1.7 Bằng chứng nền tảng đa quốc gia: phân tầng năng suất
vùng và hiệu ứng lá chắn số
(P1)}\label{bux1eb1ng-chux1ee9ng-nux1ec1n-tux1ea3ng-ux111a-quux1ed1c-gia-phuxe2n-tux1ea7ng-nux103ng-suux1ea5t-vuxf9ng-vuxe0-hiux1ec7u-ux1ee9ng-luxe1-chux1eafn-sux1ed1-p1}}

Trước khi đi vào các kiểm định đường cong theo từng bối cảnh, nghiên cứu
thành phần P1 cung cấp một nền tảng thực nghiệm cấp doanh nghiệp về dị
biệt năng suất và vai trò của năng lực trong khu vực, trên mẫu phân tích
\textbf{40.633 doanh nghiệp} thuộc \textbf{17 nền kinh tế châu Á} khảo
sát trong giai đoạn 2009--2024. Mẫu được nhóm thành ba cụm vùng: Đại
Trung Hoa (5.251 doanh nghiệp), ASEAN (12.068 doanh nghiệp) và Nam Á
(23.314 doanh nghiệp). Biến phụ thuộc là năng suất lao động dạng
logarit, winsorize tại bách phân vị 1 và 99 để kiểm soát giá trị ngoại
lai.

Kết quả hồi quy bình phương nhỏ nhất phân tầng theo từng khối biến cho
thấy một cấu trúc rõ ràng. Mô hình chỉ với biến vùng (M1) xác nhận phân
tầng năng suất sâu sắc: so với ASEAN, doanh nghiệp Đại Trung Hoa có năng
suất cao hơn đáng kể (\(\hat{\beta} = 0{,}990\), \(p < {,}001\), tương
ứng khoảng 169,1\% cao hơn) và doanh nghiệp Nam Á cũng vượt ASEAN
(\(\hat{\beta} = 0{,}506\), \(p < {,}001\), phần bù khoảng 65,9\%);
riêng biến vùng giải thích 5,4\% phương sai năng suất
(\(R^2 = 0{,}054\)). Khi bổ sung khối năng lực doanh nghiệp (M2), chấp
nhận công nghệ gắn mạnh và dương với năng suất
(\(\hat{\beta} = 0{,}409\), \(p < {,}001\), tương ứng khoảng 50,5\% cao
hơn ở doanh nghiệp có áp dụng), và chứng nhận chất lượng quốc tế có hiệu
ứng tương tự (\(\hat{\beta} = 0{,}322\), \(p < {,}001\), khoảng 38,0\%);
độ phù hợp mô hình tăng gần gấp đôi, từ \(R^2 = 0{,}054\) lên
\(0{,}105\).

Phát hiện then chốt nằm ở mô hình tương tác (M3). Chỉ số rào cản kinh
doanh tác động âm và có ý nghĩa lên năng suất
(\(\hat{\beta} = -0{,}076\), \(p < {,}001\) ở hiệu ứng chính), nhưng
tương tác giữa rào cản và chấp nhận công nghệ dương và có ý nghĩa cao
(\(\hat{\beta} = +0{,}110\), \(p < {,}001\)). Đây là bằng chứng định
lượng trực tiếp cho \textbf{hiệu ứng lá chắn số}: chấp nhận công nghệ
làm giảm khoảng 83\% tổn thất năng suất do rào cản thể chế gây ra, theo
đó ở doanh nghiệp không áp dụng công nghệ, mỗi đơn vị tăng của rào cản
làm giảm năng suất khoảng 12,5\%, trong khi ở doanh nghiệp có áp dụng,
tổn thất biên gần như bị triệt tiêu. Khoảng cách năng suất của doanh
nghiệp nhỏ và vừa cũng phụ thuộc vùng: bất lợi quy mô rõ rệt ở ASEAN
(tương tác \(\hat{\beta} = -0{,}390\), \(p < {,}001\)) nhưng gần như
không tồn tại ở Đại Trung Hoa (\(\hat{\beta} = 0{,}043\), không ý
nghĩa). Bằng chứng nền tảng này định khung cho toàn bộ chương: dị biệt
năng suất giữa các nền kinh tế châu Á không phải ngẫu nhiên mà gắn với
năng lực doanh nghiệp và chất lượng thể chế, và năng lực số vận hành như
cơ chế bù đắp cho thiếu hụt thể chế, đúng theo logic của khung CDCM.

\hypertarget{gradient-biuxean-lux1ee3i-nhuux1eadn-vuxe0-hux1ed3-sux1a1-ux111ux1eb7c-trux1b0ng-doanh-nghiux1ec7p-theo-nhuxf3m-icrv}{%
\subsection{4.1.8 Gradient biên lợi nhuận và hồ sơ đặc trưng doanh
nghiệp theo nhóm
ICRV}\label{gradient-biuxean-lux1ee3i-nhuux1eadn-vuxe0-hux1ed3-sux1a1-ux111ux1eb7c-trux1b0ng-doanh-nghiux1ec7p-theo-nhuxf3m-icrv}}

Phân tích phân tán năng suất ở các mục trước dựa trên thước đo bất biến
tiền tệ trong từng cặp nền kinh tế và năm, nên không cho phép so sánh
\emph{mức} hiệu quả giữa các nhóm. Để bổ sung một so sánh mức hợp lệ
giữa các nền kinh tế, luận án sử dụng tỷ suất lợi nhuận trên doanh thu
(ROS), một tỷ số vô đơn vị không chịu ảnh hưởng của khác biệt đơn vị
tiền tệ. Bảng 4.4 trình bày trung vị ROS theo sáu nhóm ICRV.

\textbf{Bảng 4.4.} \emph{Trung vị tỷ suất lợi nhuận trên doanh thu (ROS)
theo nhóm ICRV (mẫu con có dữ liệu chi phí).}

\begin{longtable}[]{@{}lrr@{}}
\toprule\noalign{}
Nhóm ICRV & n & Trung vị ROS \\
\midrule\noalign{}
\endhead
\bottomrule\noalign{}
\endlastfoot
I, Tiên tiến đổi mới & 3.337 & 0,503 \\
II, Tiên tiến tài nguyên & 118¹ & 0,470 \\
III, Trung bình cao & 13.401 & 0,455 \\
IV, Chuyển đổi thu nhập trung bình thấp & 37.792 & 0,351 \\
V, Đang nổi & 12.523 & 0,408 \\
VI, Đảo nhỏ & 520¹ & 0,474 \\
\end{longtable}

\emph{Ghi chú: trung vị được ưu tiên vì trung bình ở Nhóm V và VI bị méo
bởi giá trị ngoại lai. ¹ Mẫu nhỏ do ít doanh nghiệp thuộc khối tài
nguyên và đảo nhỏ có dữ liệu chi phí đầy đủ. Nguồn: ước lượng của tác
giả từ dữ liệu thô WBES (tái lập: \texttt{scripts/relock\_ros.py}).}

Mô thức ROS bổ sung một chiều thông tin cho bức tranh phân tán. Biên lợi
nhuận cao nhất ở chế độ tiên tiến đổi mới (khoảng 0,50) và thấp nhất ở
chế độ chuyển đổi thu nhập trung bình thấp (khoảng 0,35); đây là
gradient mức hợp lệ giữa các nền kinh tế, nhất quán với lý thuyết thể
chế theo đó thể chế mạnh kéo theo chi phí giao dịch thấp và biên lợi
nhuận cao hơn. Cần lưu ý rằng gradient này không đơn điệu hoàn toàn: các
nhóm thể chế yếu nhất (V và VI) có trung vị ROS cao hơn Nhóm IV, một
phần do mẫu nhỏ và do thành phần ngành dịch vụ biên cao chiếm tỷ trọng
lớn ở các nền đảo nhỏ. Điểm cốt lõi là sự khác biệt mức giữa các chế độ
thể chế hiện ra ở thước đo lợi nhuận hợp lệ, trong khi ở thước đo năng
suất thô nó bị nhiễu tiền tệ che khuất, nên luận án đọc ROS như chỉ báo
mức và đọc độ phân tán năng suất như chỉ báo phân bổ.

Hồ sơ đặc trưng doanh nghiệp theo nhóm cũng phơi bày những tương phản
đáng chú ý không trùng với thứ bậc thu nhập. Cường độ xuất khẩu trung
bình cao nhất ở Nhóm I (13,3\%) và thấp nhất ở Nhóm II tài nguyên
(3,5\%), phản ánh định hướng nội địa và tài nguyên của khối sau. Tỷ lệ
doanh nghiệp có hiện diện số bề mặt giảm dần khá đều theo gradient thể
chế, từ 61,7\% ở Nhóm I xuống 41,2\% ở Nhóm V, phù hợp với khoảng cách
hạ tầng số. Đáng chú ý nhất, tỷ lệ nữ trong quản trị cấp cao không theo
gradient thu nhập mà cao ở cả Nhóm I (27,5\%), Nhóm III (29,1\%) và Nhóm
VI đảo nhỏ (31,8\%) trong khi rất thấp ở Nhóm II tài nguyên (4,0\%), một
khác biệt thể chế kép gắn với bối cảnh văn hóa và nhà nước tô của khối
tài nguyên dẫn dắt. Hồ sơ này củng cố lập luận xuyên suốt rằng các nhóm
ICRV khác nhau về chất chứ không chỉ về mức phát triển, và do đó đòi hỏi
cách đọc theo từng chế độ thay vì một thang tuyến tính duy nhất.

\begin{center}\rule{0.5\linewidth}{0.5pt}\end{center}

\hypertarget{kux1ebft-quux1ea3-tux1ed5ng-hux1ee3p-ux111ux1ecbnh-lux1b0ux1ee3ng-phuxe2n-tuxedch-tux1ed5ng-hux1ee3p-ba-tux1ea7ng-p6}{%
\section{4.2 Kết quả tổng hợp định lượng: phân tích tổng hợp ba tầng
(P6)}\label{kux1ebft-quux1ea3-tux1ed5ng-hux1ee3p-ux111ux1ecbnh-lux1b0ux1ee3ng-phuxe2n-tuxedch-tux1ed5ng-hux1ee3p-ba-tux1ea7ng-p6}}

\hypertarget{chux1ec9-sux1ed1-hiux1ec7u-ux1ee9ng-tux1ed5ng-hux1ee3p}{%
\subsection{4.2.1 Chỉ số hiệu ứng tổng
hợp}\label{chux1ec9-sux1ed1-hiux1ec7u-ux1ee9ng-tux1ed5ng-hux1ee3p}}

Phân tích phân tích tổng hợp tổng hợp ba tầng (three-level multilevel
meta-analysis, MARA) được tiến hành trên tập dữ liệu gồm \(k = 238\)
nghiên cứu với tổng cộng K = 288 cỡ ảnh hưởng đa dạng theo vùng địa lý,
giai đoạn và phương pháp. Kết quả cho thấy hiệu ứng tổng hợp (pooled
correlation) là \(r = 0{,}074\), với khoảng tin cậy 95\% dương và có ý
nghĩa thống kê. Giá trị \(r\) này khẳng định bức tranh tổng quát từ các
các phân tích tổng hợp trước đây của Bausch và Krist (2007), Kirca et
al.~(2012), Marano et al.~(2016) và Arte và Larimo (2022): mối quan hệ
tổng hợp giữa quốc tế hóa và hiệu quả hoạt động là dương nhưng có biên
độ khiêm tốn.

Đặc biệt đáng lưu ý, kết quả P6 hoàn toàn nhất quán với cơ sở phân tích
tổng hợp đã công bố tại hội nghị ICBEF 2024, với \(k = 113\) nghiên cứu
và \(r = 0{,}07\), sự hội tụ này tạo ra bằng chứng chéo vững chắc khi mở
rộng mẫu gộp lên gần gấp đôi số nghiên cứu.

\hypertarget{phuxe2n-tuxedch-dux1ecb-biux1ec7t-ba-tux1ea7ng}{%
\subsection{4.2.2 Phân tích dị biệt ba
tầng}\label{phuxe2n-tuxedch-dux1ecb-biux1ec7t-ba-tux1ea7ng}}

Kết quả quan trọng nhất của phân tích MARA không phải là giá trị điểm
của \(r\) mà là cấu trúc của dị biệt (\(I^2\)). Tổng \(I^2 = 62{,}4\%\),
cho thấy mức dị biệt đáng kể trong tập hợp 238 nghiên cứu không phải do
sai số lấy mẫu. Quan trọng hơn, khi phân tách theo cấu trúc ba tầng:

\begin{itemize}
\tightlist
\item
  \textbf{Tầng 2} (between-effect within-study, tức dị biệt giữa các
  effect-size trong cùng một nghiên cứu): chiếm \(54{,}1\%\) tổng phương
  sai, đây là nguồn dị biệt chủ đạo.
\item
  \textbf{Tầng 3} (between-study, tức dị biệt giữa các nghiên cứu): chỉ
  chiếm \(8{,}4\%\).
\end{itemize}

Phát hiện này có hàm ý lý thuyết sâu sắc: dị biệt trong tài liệu về quan
hệ quốc tế hóa và hiệu quả xuất phát chủ yếu từ sự biến thiên
\textbf{bên trong các nghiên cứu} (ví dụ: cùng một mẫu doanh nghiệp
nhưng dùng các thước đo khác nhau cho cùng cấu trúc cho ra kết quả khác
nhau), chứ không phải giữa các nghiên cứu với nhau. Điều này gợi ý rằng
sự không nhất quán trong tài liệu phần lớn là vấn đề \textbf{đo lường và
bối cảnh}, không phải là bằng chứng về sự thiếu nhất quán thực sự trong
quan hệ nhân quả giữa quốc tế hóa và hiệu quả.

\hypertarget{ux111iux1ec1u-tiux1ebft-bux1edfi-chux1ebf-ux111ux1ed9-thux1ec3-chux1ebf-icrv-ux111iux1ec1u-tiux1ebft}{%
\subsection{4.2.3 Điều tiết bởi chế độ thể chế: ICRV điều
tiết}\label{ux111iux1ec1u-tiux1ebft-bux1edfi-chux1ebf-ux111ux1ed9-thux1ec3-chux1ebf-icrv-ux111iux1ec1u-tiux1ebft}}

Kiểm định điều tiết bằng nhóm chế độ thể chế ICRV cho kết quả có ý nghĩa
thống kê: \(Q_M = 17{,}35\), \(df = 4\), \(p = {,}002\). Điều này xác
nhận rằng chế độ thể chế, được phân loại thành các nhóm từ tiên tiến đổi
mới (Nhóm I) đến nhóm biên thể chế yếu nhất, là biến điều tiết có ý
nghĩa thực sự cho mối quan hệ tổng hợp giữa quốc tế hóa và hiệu quả
trong tài liệu. Kiểm định \(Q_M\) vượt ngưỡng tới hạn (\(p = {,}002\)),
hàm ý rằng sự khác biệt giữa các nhóm ICRV không thể giải thích bằng sai
số lấy mẫu.

Tuy điều tiết có ý nghĩa thống kê, mô thức cỡ ảnh hưởng theo nhóm
\textbf{không} tạo thành một gradient đơn điệu. Cỡ ảnh hưởng trung bình
theo nhóm là: Nhóm I tiên tiến đổi mới \(\bar{r} = 0{,}079\)
(\(k = 139\)), Nhóm III \(\bar{r} = 0{,}069\) (\(k = 91\)), Nhóm II
\(\bar{r} = 0{,}065\) (\(k = 25\)), và nhóm hỗn hợp
\(\bar{r} = 0{,}053\) (\(k = 30\)); nhóm biên thể chế yếu nhất cho giá
trị bất thường cao (\(\bar{r} = 0{,}349\), \(k = 3\)) nhưng bị chi phối
bởi một giá trị ngoại lai duy nhất (Pouresmaeili và cộng sự, 2018) nên
không đáng tin cậy. Vì vậy, kết luận đúng từ P6 là chế độ thể chế điều
tiết cỡ ảnh hưởng một cách có ý nghĩa nhưng không theo trật tự đơn điệu
rõ ràng ở cấp tổng hợp; nhóm thể chế mạnh nhất có cỡ ảnh hưởng tin cậy
cao nhất, song các nhóm còn lại không xếp hạng theo gradient kỳ vọng, và
nhóm đảo nhỏ không có đủ cỡ ảnh hưởng để phân tích riêng. Cách đọc này
nhất quán với lý thuyết thể chế ở chiều tổng quát (North, 1990; Marano
et al., 2016) nhưng tinh chỉnh kỳ vọng quan trọng: bằng chứng về một
\emph{trật tự} thể chế đơn điệu không hiện ra ở cỡ ảnh hưởng cấp tổng
hợp mà ở \textbf{vị trí điểm uốn} trong phân tích hồi quy cấp doanh
nghiệp của P7, nơi cấu trúc ba vùng định hình rõ (Mục 4.6.2). Sự phân
công này, P6 xác nhận \emph{sự tồn tại} của điều tiết thể chế còn P7
định vị \emph{hình thức} của nó, chính là một lý do để ưu tiên bằng
chứng cấp doanh nghiệp trên khung phân loại thể chế nhất quán.
\emph{(Lưu ý phân loại: P6, phân tích tổng hợp cấp nghiên cứu, dùng hệ
ICRV năm chế độ rút gọn {[}I tiên tiến đổi mới, II, III, FR biên thể chế
yếu, MX hỗn hợp{]}, khác với hệ sáu nhóm cấp doanh nghiệp dùng ở Mục 4.1
và Mục 4.6; nhóm biên thể chế yếu nhất của P6 tương ứng Nhóm V và VI
trong hệ sáu nhóm.)}

\hypertarget{cuxe1c-biux1ebfn-ux111iux1ec1u-tiux1ebft-nux103ng-lux1ef1c-sux1ed1-pha-sux1ed1-huxf3a-vuxe0-thux1b0ux1edbc-ux111o}{%
\subsection{4.2.3.1 Các biến điều tiết năng lực số, pha số hóa và thước
đo}\label{cuxe1c-biux1ebfn-ux111iux1ec1u-tiux1ebft-nux103ng-lux1ef1c-sux1ed1-pha-sux1ed1-huxf3a-vuxe0-thux1b0ux1edbc-ux111o}}

Ngoài chế độ thể chế, P6 kiểm định sáu biến điều tiết khác, ba biến mới
của luận án và bốn biến chuẩn của tài liệu, qua đó định vị rõ điều tiết
nào là thực chất. Trái với kỳ vọng, chấp nhận số ngữ cảnh ở cấp tổng hợp
\textbf{không} điều tiết đáng kể cỡ ảnh hưởng (\(Q_M = 1{,}23\),
\(df = 2\), \(p = {,}541\); hệ số hồi quy phân tích tổng hợp
\(\beta = +0{,}003\), \(p = {,}744\)), nghĩa là phần thưởng số không
hiện ra như một điều tiết phổ quát ở dữ liệu thứ cấp, nhất quán với phát
hiện cấp doanh nghiệp rằng vai trò uốn đường cong của DAI phụ thuộc bối
cảnh chứ không phổ quát. Tương tự, pha phát triển số hóa không điều tiết
đáng kể (\(Q_M = 0{,}56\), \(df = 2\), \(p = {,}755\)), với cỡ ảnh hưởng
gần như đồng đều qua các pha (trước, trong và sau làn sóng số hóa).

Hai biến điều tiết về thước đo cho thấy lựa chọn vận hành hóa có ảnh
hưởng đáng kể đến cỡ ảnh hưởng quan sát, củng cố cách đọc rằng phần lớn
dị biệt nằm ở tầng nội bộ nghiên cứu. Theo cách đo mức độ quốc tế hóa,
cỡ ảnh hưởng dao động từ 0,059 khi đo bằng cường độ xuất khẩu trên doanh
thu (FSTS, \(K = 137\)) đến 0,097 khi đo bằng phạm vi địa lý
(\(K = 50\)) và 0,111 khi đo bằng chỉ số tổng hợp (\(K = 31\)). Theo
cách đo hiệu quả, cỡ ảnh hưởng theo thước đo kế toán là 0,078
(\(K = 246\)) nhưng giảm xuống gần không và đổi dấu nhẹ khi đo bằng năng
suất lao động (\(-0{,}019\), \(K = 12\)), một mẫu rất nhỏ. Mô thức này
có hàm ý phương pháp trực tiếp: việc luận án dùng nhất quán cường độ
xuất khẩu làm thước đo quốc tế hóa và năng suất lao động làm thước đo
hiệu quả ở cấp doanh nghiệp loại bỏ chính nguồn dị biệt do thước đo mà
phân tích tổng hợp cấp tổng hợp không thể kiểm soát, qua đó cho ước
lượng điều tiết thể chế sạch hơn.

\hypertarget{kiux1ec3m-ux111ux1ecbnh-lux1ec7ch-lux1ea1c-cuxf4ng-bux1ed1}{%
\subsection{4.2.4 Kiểm định lệch lạc công
bố}\label{kiux1ec3m-ux111ux1ecbnh-lux1ec7ch-lux1ea1c-cuxf4ng-bux1ed1}}

Phân tích lệch lạc công bố gồm ba kiểm định bổ sung nhau. \textbf{Kiểm
định Egger} cho kết quả \(b = 0{,}475\) (\(p = {,}057\)), sát ngưỡng
\(\alpha = 0{,}05\) nhưng không đạt ý nghĩa thống kê theo tiêu chuẩn
thông thường. \textbf{Kiểm định Begg--Mazumdar} (Begg \& Mazumdar, 1994)
cho kết quả \(\tau = -0{,}134\) (\(p = {,}0007\)), có ý nghĩa thống kê
rõ ràng, cho thấy các nghiên cứu với sai số chuẩn lớn hơn (mẫu nhỏ hơn)
có xu hướng báo cáo cỡ ảnh hưởng thấp hơn, phù hợp với lệch lạc công bố
ưu tiên kết quả dương tính. \textbf{Phân tích cắt và điền} điều chỉnh
ước lượng từ \(\bar{r} = 0{,}074\) xuống \(\bar{r} = 0{,}035\) (với
\(k = 58\) nghiên cứu được điền khuyết), mức điều chỉnh đáng kể nhưng
không đổi chiều hướng tích cực của quan hệ tổng hợp. \textbf{Số an toàn
\(N\) = 45.848} vượt xa ngưỡng tới hạn \(5k + 10 = 1.200\), khẳng định
lệch lạc công bố không thể bác bỏ hoàn toàn kết quả tổng hợp. Tổng hợp
lại, bằng chứng lệch lạc công bố được xác nhận qua kiểm định
Begg--Mazumdar có ý nghĩa (\(p = {,}0007\)): cỡ ảnh hưởng gộp
\(\bar{r} = 0{,}074\) có thể bị thổi phồng so với hiệu ứng tổng thể
thực, với ước lượng bảo thủ hơn là \(\bar{r} = 0{,}035\).

Phát hiện lệch lạc công bố này có một hàm ý lý thuyết quan trọng vượt ra
ngoài một cảnh báo phương pháp đơn thuần. Trong phân tách dị biệt ba
tầng, phần lớn dị biệt nằm ở \textbf{tầng nội bộ nghiên cứu}
(\(I^2_{\text{within}} = 54{,}1\%\)) hơn là tầng giữa các nghiên cứu
(\(I^2_{\text{between}} = 8{,}4\%\)), nghĩa là biến thiên hệ số chủ yếu
đến từ cách mỗi nghiên cứu vận hành hóa quốc tế hóa, lựa chọn thước đo
hiệu quả và đặc tả biến kiểm soát, chứ không phải từ khác biệt bối cảnh
quốc gia. Kết hợp với bằng chứng lệch lạc công bố, cấu trúc dị biệt này
gợi ý một cách diễn giải lại ``câu đố dị biệt'' trong tài liệu: phần lớn
phương sai chưa giải thích được có thể phản ánh lựa chọn ở phía công bố
và đặc tả mô hình nhiều hơn là các điều kiện bối cảnh thể chế hay số
hóa. Phát hiện này không phủ định vai trò điều tiết của thể chế, vốn
được xác nhận độc lập qua kiểm định \(Q_M\) giữa các nhóm ICRV, nhưng nó
định vị lại kỳ vọng: phân tích tổng hợp cấp tổng hợp khó nắm bắt được
điều tiết thể chế vì các nghiên cứu sơ cấp hiếm khi báo cáo đủ chi tiết
theo chế độ thể chế, và vì vậy bằng chứng cấp doanh nghiệp trực tiếp
trên một khung phân loại thể chế nhất quán, như luận án thực hiện, là
con đường hữu hiệu hơn để bóc tách dị biệt so với việc tích lũy thêm
phân tích tổng hợp trên dữ liệu thứ cấp không đồng nhất.

\hypertarget{phuxe2n-ruxe3-phux1b0ux1a1ng-sai-ba-tux1ea7ng-vuxe0-huxe0m-uxfd-vux1ec1-nguux1ed3n-dux1ecb-biux1ec7t}{%
\subsection{4.2.5 Phân rã phương sai ba tầng và hàm ý về nguồn dị
biệt}\label{phuxe2n-ruxe3-phux1b0ux1a1ng-sai-ba-tux1ea7ng-vuxe0-huxe0m-uxfd-vux1ec1-nguux1ed3n-dux1ecb-biux1ec7t}}

Cấu trúc ba tầng của mô hình MARA cho phép quy phương sai hệ thống về
đúng nguồn gốc của nó thay vì gộp chung như mô hình hiệu ứng ngẫu nhiên
một tầng. Hai thành phần phương sai ước lượng bằng REML là phương sai
trong nghiên cứu \(\hat{\sigma}^2_{(2)} = 0{,}00878\) và phương sai giữa
các nghiên cứu \(\hat{\sigma}^2_{(3)} = 0{,}00136\). Phương sai trong
nghiên cứu lớn gấp khoảng sáu lần phương sai giữa các nghiên cứu, một tỷ
lệ trực tiếp chuyển thành phần dị biệt \(I^2_{(2)} = 54{,}1\%\) áp đảo
phần \(I^2_{(3)} = 8{,}4\%\) đã trình bày ở Mục 4.2.2. Kiểm định Cochran
tổng cho \(Q = 1.909{,}42\) (\(df = 287\), \(p < {,}001\)), bác bỏ giả
thuyết đồng nhất một cách rõ ràng và do đó biện minh cho lựa chọn mô
hình hiệu ứng ngẫu nhiên ba tầng.

Hàm ý phương pháp luận của phân rã này khác biệt căn bản so với một báo
cáo \(I^2\) tổng đơn thuần. Vì biến thiên hệ số chủ yếu phát sinh từ
những lựa chọn vận hành hóa bên trong cùng một bài báo, cụ thể là cách
đo lường cường độ quốc tế hóa (tỷ lệ doanh thu xuất khẩu trên tổng doanh
thu, chỉ số entropy, số thị trường nước ngoài), cách chọn thước đo hiệu
quả (dựa trên kế toán, dựa trên thị trường, dựa trên năng suất) và đặc
tả biến kiểm soát, nên việc bổ sung thêm các nghiên cứu sơ cấp mới khó
có khả năng thu hẹp dị biệt nếu các nghiên cứu đó tiếp tục báo cáo nhiều
hệ số không đồng nhất cho cùng một mẫu. Mô hình hai tầng quy ước (ước
lượng DerSimonian--Laird) cho cùng giá trị điểm \(r = 0{,}074\) với
khoảng tin cậy \([0{,}061;\ 0{,}087]\), xác nhận rằng việc bỏ qua cấu
trúc lồng ghép không làm chệch ước lượng gộp mà chỉ làm mất thông tin về
nguồn dị biệt. Đây là lý do trung tâm để luận án ưu tiên bằng chứng cấp
doanh nghiệp trên một khung phân loại thể chế nhất quán hơn là tiếp tục
tích lũy phân tích tổng hợp trên dữ liệu thứ cấp không đồng nhất.

\hypertarget{gradient-hiux1ec7u-ux1ee9ng-theo-nhuxf3m-thux1ec3-chux1ebf-vuxe0-giux1edbi-hux1ea1n-cux1ee7a-bux1eb1ng-chux1ee9ng-ux111ux1ecbnh-hux1b0ux1edbng}{%
\subsection{4.2.6 Gradient hiệu ứng theo nhóm thể chế và giới hạn của
bằng chứng định
hướng}\label{gradient-hiux1ec7u-ux1ee9ng-theo-nhuxf3m-thux1ec3-chux1ebf-vuxe0-giux1edbi-hux1ea1n-cux1ee7a-bux1eb1ng-chux1ee9ng-ux111ux1ecbnh-hux1b0ux1edbng}}

Bảng hệ số gộp theo từng nhóm ICRV cho thấy một bức tranh tinh tế hơn so
với kết luận ``có điều tiết'' đơn thuần. Nhóm I (Advanced đổi mới sáng
tạo) có hệ số gộp lớn nhất trong các nhóm đủ dữ liệu,
\(\bar{r} = 0{,}079\) (khoảng tin cậy \([0{,}058;\ 0{,}099]\),
\(k = 139\) hiệu ứng); nhóm II (Trung bình cao) có \(\bar{r} = 0{,}065\)
(\(k = 25\)); nhóm III (Đang nổi) có \(\bar{r} = 0{,}069\) (\(k = 91\));
và nhóm đa bối cảnh có \(\bar{r} = 0{,}053\) (\(k = 30\)). Bốn nhóm này
đều cho hệ số dương có ý nghĩa và xếp theo chiều phù hợp với kỳ vọng lý
thuyết, song khác biệt giữa chúng không đạt ý nghĩa thống kê trong các
so sánh cặp đôi: Nhóm I so với Nhóm III cho \(b = -0{,}011\)
(\(p = {,}502\)), và Nhóm I so với nhóm đa bối cảnh cho \(b = -0{,}026\)
(\(p = {,}269\)).

Do đó cần phân biệt rạch ròi giữa kết quả của kiểm định tổng thể và kết
quả của các so sánh định hướng. Kiểm định \(Q_M\) giữa các nhóm có ý
nghĩa (\(p = {,}002\)) khẳng định rằng hệ số quốc tế hóa và hiệu quả
biến thiên có hệ thống giữa các chế độ thể chế, nhưng ý nghĩa này được
dẫn dắt gần như toàn bộ bởi một nhóm cận biên: nhóm Frontier chỉ gồm ba
nghiên cứu lại cho hệ số gộp dị thường \(\bar{r} = 0{,}349\)
(\(k = 3\)), và tương phản duy nhất có ý nghĩa là Nhóm I so với Frontier
(\(b = +0{,}285\), \(p < {,}001\)). Giá trị Frontier này bị chi phối bởi
một nghiên cứu ngoại lai duy nhất (Pouresmaeili và cộng sự, 2018, với
\(r = 0{,}69\) trên 226 doanh nghiệp) và vì vậy không thể coi là một ước
lượng đáng tin cậy cho chế độ thể chế yếu nhất. Hệ quả là gradient định
hướng không được xác nhận ở cấp phân tích tổng hợp với cỡ mẫu hiện tại;
sự khác biệt giữa Nhóm I và Nhóm III chỉ là \(0{,}010\) và không có ý
nghĩa thống kê. Phát hiện này định vị lại kỳ vọng một cách quan trọng:
bằng chứng về điều tiết thể chế tồn tại, nhưng việc bóc tách gradient cụ
thể đòi hỏi bằng chứng cấp doanh nghiệp trên một khung phân loại nhất
quán, đúng như con đường mà các nghiên cứu thành phần P3--P8 và mô hình
đa quốc gia P7 của luận án theo đuổi.

Hai biến điều tiết bối cảnh khác được kiểm định trong cùng khung MARA
nhưng không đạt ý nghĩa thống kê. Chấp nhận số cấp quốc gia
(country-level Digital Adoption Index --- cDAI) cho \(Q_M = 1{,}23\)
(\(df = 2\), \(p = {,}541\)), với hệ số phân tích tổng hợp-hồi quy liên
tục \(\beta_{\text{cDAI}} = +0{,}003\) (\(p = {,}744\)) và thứ tự ba
phân nhóm không đơn điệu (Thấp \(0{,}075\), Trung bình \(0{,}063\), Cao
\(0{,}091\)). Chu kỳ nghịch lý số (Digital Paradox Lifecycle --- DPL)
cho \(Q_M = 0{,}56\) (\(df = 2\), \(p = {,}755\)), với thứ tự ba giai
đoạn ngược với dự đoán lý thuyết. Việc chỉ có một trong ba biến điều
tiết bối cảnh đạt ý nghĩa, và biến đó lại bị chi phối bởi một ô dữ liệu
cận biên, củng cố cách diễn giải rằng dị biệt quốc tế hóa và hiệu quả ở
cấp tổng hợp khó được giải thích bằng các biến bối cảnh quốc gia khi các
nghiên cứu sơ cấp hiếm khi báo cáo đủ chi tiết theo chế độ thể chế.

\hypertarget{tuxednh-vux1eefng-cux1ee7a-hiux1ec7u-ux1ee9ng-gux1ed9p-vuxe0-phux1ea1m-vi-cux1ee7a-bux1eb1ng-chux1ee9ng-lux1ec7ch-lux1ea1c-cuxf4ng-bux1ed1}{%
\subsection{4.2.7 Tính vững của hiệu ứng gộp và phạm vi của bằng chứng
lệch lạc công
bố}\label{tuxednh-vux1eefng-cux1ee7a-hiux1ec7u-ux1ee9ng-gux1ed9p-vuxe0-phux1ea1m-vi-cux1ee7a-bux1eb1ng-chux1ee9ng-lux1ec7ch-lux1ea1c-cuxf4ng-bux1ed1}}

Hiệu ứng gộp \(r = 0{,}074\) giữ ổn định qua toàn bộ các kiểm định độ
vững có sẵn. Khi chỉ giữ các hệ số tương quan được báo cáo trực tiếp, hệ
số gộp là \(0{,}077\) (\(k = 241\)); khi loại các hiệu ứng từ mẫu rất
nhỏ (\(n < 30\)), hệ số là \(0{,}073\) (\(k = 286\)); khi chỉ giữ thước
đo hiệu quả dựa trên kế toán, hệ số là \(0{,}075\) (\(k = 247\)). Kiểm
định bảo thủ nhất là giới hạn ở các nghiên cứu chỉ dùng tỷ lệ doanh thu
xuất khẩu trên tổng doanh thu làm thước đo quốc tế hóa, cho hệ số suy
giảm còn \(\bar{r} = 0{,}060\) (\(k = 138\)), vẫn dương và có ý nghĩa,
gợi ý rằng tính không đồng nhất của thước đo quốc tế hóa làm phóng đại
một phần hiệu ứng gộp khi gộp chung các thước đo rộng hơn. Phân tích
loại bỏ từng nghiên cứu (leave-one-out) cho dải hệ số
\([0{,}071;\ 0{,}076]\) và không một nghiên cứu nào trong 288 hiệu ứng
làm đổi chiều kết quả, xác nhận rằng hiệu ứng dương khiêm tốn không phải
là sản phẩm của bất kỳ quan sát ảnh hưởng đơn lẻ nào.

Về lệch lạc công bố, cần nhấn mạnh sự không thống nhất giữa hai kiểm
định bất đối xứng phễu đã nêu ở Mục 4.2.4: kiểm định Egger không vượt
ngưỡng quy ước (\(b = 0{,}475\), \(p = {,}057\)) trong khi kiểm định thứ
hạng Begg--Mazumdar lại có ý nghĩa (\(\tau = -0{,}134\),
\(p = {,}0007\)), nên bằng chứng bất đối xứng dựa trên Egger đơn lẻ
không được khẳng định. Quy trình cắt-và-điền (trim-and-fill) gán thêm 58
nghiên cứu ở phía trái của phễu và đưa hệ số điều chỉnh xuống
\(\bar{r} = 0{,}035\) (khoảng tin cậy \([0{,}019;\ 0{,}051]\)); đây là
cận dưới bảo thủ, hàm ý hiệu ứng dân số thực có thể gần
\(r \approx 0{,}035\) hơn là các con số trong dải \(0{,}07\)--\(0{,}10\)
thường được trích dẫn. Cùng lúc, chỉ số an toàn fail-safe \(N = 45.848\)
vượt xa ngưỡng tới hạn \(5k + 10 = 1.200\), hàm ý rằng lệch lạc công bố
tuy hiện diện nhưng không thể đảo ngược hoàn toàn kết luận về dấu của
quan hệ. Tổng hợp lại, các kiểm định độ vững và lệch lạc công bố hội tụ
về một thông điệp duy nhất: quan hệ quốc tế hóa và hiệu quả ở cấp tổng
hợp là dương, ổn định về dấu, nhưng nhỏ hơn nhiều so với ấn tượng từ tài
liệu đã công bố.

Một hàm ý phương pháp luận trực tiếp nối kết phân tích tổng hợp P6 với
chiến lược thiết kế cấp doanh nghiệp của luận án. Vì hai biến điều tiết
thước đo cho thấy cỡ ảnh hưởng quan sát phụ thuộc mạnh vào lựa chọn vận
hành hóa, dao động từ \(0{,}059\) khi đo quốc tế hóa bằng cường độ xuất
khẩu trên doanh thu (\(K = 137\)) tới \(0{,}097\) khi đo bằng phạm vi
địa lý (\(K = 50\)) và \(0{,}111\) khi đo bằng chỉ số tổng hợp
(\(K = 31\)), và vì cỡ ảnh hưởng theo thước đo kế toán (\(0{,}078\),
\(K = 246\)) đổi dấu nhẹ khi chuyển sang năng suất lao động
(\(-0{,}019\), \(K = 12\)), một phần lớn dị biệt tầng nội bộ nghiên cứu
(\(I^2_{\text{within}} = 54{,}1\%\)) phản ánh chính sự không đồng nhất
của thước đo. Hệ quả thiết kế là rõ ràng: việc luận án dùng nhất quán
cường độ xuất khẩu (FSTS) làm thước đo quốc tế hóa và năng suất lao động
làm thước đo hiệu quả ở cấp doanh nghiệp loại bỏ chính nguồn dị biệt do
thước đo mà phân tích tổng hợp cấp tổng hợp không thể kiểm soát, cho ước
lượng điều tiết thể chế sạch hơn và biến bằng chứng cấp doanh nghiệp
trên khung phân loại thể chế nhất quán thành con đường hữu hiệu hơn để
bóc tách dị biệt so với việc tích lũy thêm phân tích tổng hợp trên dữ
liệu thứ cấp không đồng nhất.

\begin{center}\rule{0.5\linewidth}{0.5pt}\end{center}

\hypertarget{kux1ebft-quux1ea3-bux1ed1i-cux1ea3nh-thux1ec3-chux1ebf-tiuxean-tiux1ebfn-ux111ux1ed5i-mux1edbi-singapore-p4}{%
\section{4.3 Kết quả bối cảnh thể chế tiên tiến: Đổi mới: Singapore
(P4)}\label{kux1ebft-quux1ea3-bux1ed1i-cux1ea3nh-thux1ec3-chux1ebf-tiuxean-tiux1ebfn-ux111ux1ed5i-mux1edbi-singapore-p4}}

\hypertarget{xuxe1c-nhux1eadn-quan-hux1ec7-phi-tuyux1ebfn-chux1eef-u-ngux1b0ux1ee3c}{%
\subsection{4.3.1 Xác nhận quan hệ phi tuyến chữ U
ngược}\label{xuxe1c-nhux1eadn-quan-hux1ec7-phi-tuyux1ebfn-chux1eef-u-ngux1b0ux1ee3c}}

Phân tích trên mẫu doanh nghiệp Singapore từ dữ liệu WBES với
\(N = 623\) cơ sở (trong đó 84 cơ sở xuất khẩu dùng cho kiểm định mẫu
phụ) cho thấy quan hệ giữa FSTS và hiệu quả gần tuyến tính dương; thành
phần bậc hai chỉ gợi ý đường cong rất thoải với điểm uốn ngoài miền dữ
liệu. Mô hình bậc hai với biến FSTS và FSTS\(^2\) cho hệ số bậc hai âm,
phù hợp với dạng chữ U ngược như kỳ vọng từ H1 trong khung lý thuyết
CDCM.

Điểm điểm uốn ước lượng có giá trị trung tâm \textbf{TP ≈ 88,6\% FSTS},
nhưng khoảng tin cậy bootstrap rất rộng và \textbf{vượt ngưỡng khả thi
của FSTS} (CI: {[}53\%, 253\%{]}; trong khi FSTS theo định nghĩa không
thể vượt 100\%). Vì cận trên của khoảng tin cậy nằm ngoài miền giá trị
khả thi, điểm uốn \textbf{không được định vị một cách tin cậy trong miền
dữ liệu quan sát}: trên thực tế, trong khoảng FSTS quan sát được
(0--100\%), quan hệ FSTS--hiệu quả của doanh nghiệp Singapore là
\textbf{tăng gần như đơn điệu}, và phần suy giảm bên phải của đường chữ
U ngược nằm ngoài (hoặc ở rìa) miền dữ liệu nên không được kiểm chứng
thực nghiệm. Do đó, kết quả Singapore nên được diễn giải thận trọng là
\emph{``chưa quan sát thấy ngưỡng bất lợi của quốc tế hóa trong miền dữ
liệu''} thay vì khẳng định một điểm uốn cụ thể ở 88,6\%; điều này nhất
quán với kỳ vọng lý thuyết rằng ở bối cảnh thể chế mạnh (Nhóm I), dư địa
quốc tế hóa có lợi là rất rộng (xem Hình 4.8).

\hypertarget{ux111iux1ec1u-tiux1ebft-bux1edfi-chux1ec9-sux1ed1-chux1ea5p-nhux1eadn-sux1ed1-dai}{%
\subsection{4.3.2 Điều tiết bởi chỉ số chấp nhận số
(DAI)}\label{ux111iux1ec1u-tiux1ebft-bux1edfi-chux1ec9-sux1ed1-chux1ea5p-nhux1eadn-sux1ed1-dai}}

Mô hình M5 bổ sung tương tác \(FSTS \times DAI\) (bao gồm cả Tầng 1 và
Tầng 2 của chỉ số DAI) cho thấy hiệu ứng khuếch đại (amplification) tại
mức FSTS cao: hệ số tương tác \(FSTS^2 \times DAI = +3{,}119\) và có ý
nghĩa thống kê (\(p = {,}005\)). Cụ thể, ở mức độ quốc tế hóa cao, doanh
nghiệp Singapore có năng lực chấp nhận số (digital adoption) cao hơn duy
trì hiệu quả tốt hơn so với doanh nghiệp có mức DAI thấp hơn, cho thấy
DAI đóng vai trò nguồn lực bổ trợ (situational resource) làm chậm đà suy
giảm phía bên phải của đường phi tuyến.

\hypertarget{giux1edbi-hux1ea1n-suy-luux1eadn-vuxe0-cux1ea3nh-buxe1o-thux1ed1ng-kuxea}{%
\subsection{4.3.3 Giới hạn suy luận và cảnh báo thống
kê}\label{giux1edbi-hux1ea1n-suy-luux1eadn-vuxe0-cux1ea3nh-buxe1o-thux1ed1ng-kuxea}}

Cần nhấn mạnh rằng kết quả Singapore là \textbf{ranh giới suy luận}
(inferential bounds) do hạn chế về quy mô mẫu: \(N = 623\) quan sát toàn
mẫu và chỉ 84 cơ sở xuất khẩu cho các kiểm định mẫu phụ, dẫn đến power
thống kê ước tính chỉ đạt khoảng 16\%. Luận án diễn giải kết quả
Singapore như \emph{bằng chứng phù hợp nhưng chưa kết luận} (consistent
but underpowered evidence), không bác bỏ giả thuyết nhưng cũng không
tuyên bố xác nhận dứt khoát.

\hypertarget{phuxe2n-phux1ed1i-mux1eabu-ux111ux1ecbnh-vux1ecb-ux111iux1ec3m-uux1ed1n-bux1eb1ng-bootstrap-vuxe0-ux111a-cux1ed9ng-tuyux1ebfn}{%
\subsection{4.3.4 Phân phối mẫu, định vị điểm uốn bằng bootstrap và đa
cộng
tuyến}\label{phuxe2n-phux1ed1i-mux1eabu-ux111ux1ecbnh-vux1ecb-ux111iux1ec3m-uux1ed1n-bux1eb1ng-bootstrap-vuxe0-ux111a-cux1ed9ng-tuyux1ebfn}}

Bức tranh phân phối của mẫu Singapore củng cố cách diễn giải thận trọng
nói trên. Cường độ xuất khẩu trung bình chỉ đạt
\(\overline{\text{FSTS}} = 0{,}045\) với độ lệch chuẩn \(0{,}144\),
trong đó khoảng \(82\%\) doanh nghiệp báo cáo không xuất khẩu và chỉ
\(3\%\) vượt ngưỡng \(50\%\) doanh thu nước ngoài. Chính sự dồn nén ở
đầu thấp của miền FSTS, kết hợp với mức bão hòa số rất cao (khoảng
\(67\%\) doanh nghiệp có hiện diện website), làm cho nhánh suy giảm bên
phải của đường cong khó được nhận diện thống kê trong cỡ mẫu 623 cơ sở.
Kiểm định Lind--Mehlum do đó không bác bỏ tính đơn điệu ở mức ý nghĩa
thông thường (\(p = {,}303\)), một kết quả không có ý nghĩa nhưng giàu
thông tin (informative null), nhất quán với giả thuyết bão hòa số thay
vì một thất bại của khung chữ U ngược. Về kỹ thuật, điểm uốn trung tâm
\(88{,}6\%\) được tính từ công thức \(-\hat{\beta}_1/(2\hat{\beta}_2)\)
trên FSTS đã trừ trung bình (cho giá trị \(\approx 84{,}1\%\)) rồi cộng
lại trung bình mẫu \(0{,}045\); bootstrap 5.000 lần tái lập tái hiện
dạng chữ U ngược ở \(96{,}3\%\) số lần lặp nhưng khoảng tin cậy phân vị
\(95\%\) của điểm uốn trải rộng \([53\%; 253\%]\) (trung vị \(80\%\),
khoảng tứ phân vị \([68\%; 102\%]\)), khẳng định vị trí điểm uốn chỉ
được định vị lỏng lẻo. Đa cộng tuyến giữa số hạng bậc nhất và bậc hai
không phải là vấn đề thực chất: hệ số phóng đại phương sai (variance
inflation factor --- VIF) duy trì dưới 3 ở mọi đặc tả cuối.

\hypertarget{tuxe1ch-bux1ea1ch-hai-cux1ea5u-truxfac-tci-vuxe0-dai-hiux1ec7u-ux1ee9ng-trux1ef1c-tiux1ebfp-vuxe0-tiux1ebfn-triux1ec3n-ux111ux1ed9-phuxf9-hux1ee3p-muxf4-huxecnh}{%
\subsection{4.3.5 Tách bạch hai cấu trúc TCI và DAI: hiệu ứng trực tiếp
và tiến triển độ phù hợp mô
hình}\label{tuxe1ch-bux1ea1ch-hai-cux1ea5u-truxfac-tci-vuxe0-dai-hiux1ec7u-ux1ee9ng-trux1ef1c-tiux1ebfp-vuxe0-tiux1ebfn-triux1ec3n-ux111ux1ed9-phuxf9-hux1ee3p-muxf4-huxecnh}}

Bằng chứng Singapore cho phép phân tách rạch ròi hai cấu trúc năng lực
thường bị gộp chung trong tài liệu số hóa. Năng lực công nghệ thể hiện
một hiệu ứng trực tiếp dương và bền vững trên năng suất lao động: ở mô
hình hiệu ứng trực tiếp (M5), \(\hat{\beta}(\text{TCI}) = +0{,}168\)
(\(SE = 0{,}040\); \(p < {,}001\)), đi kèm cải thiện độ phù hợp mô hình
khi \(R^2\) tăng từ \(0{,}178\) ở M2 lên \(0{,}199\) ở M5. Trong đặc tả
điều tiết TCI bổ sung (M3), hệ số trực tiếp vẫn dương và có ý nghĩa
(\(\hat{\beta} = +0{,}188\), \(p < {,}001\)) nhưng các số hạng tương tác
với \(\text{FSTS}\) và \(\text{FSTS}^2\) không có ý nghĩa thống kê khi
kiểm định chung, củng cố cách đọc rằng năng lực công nghệ nâng mặt bằng
năng suất thay vì uốn hình dạng đường cong.

Ngược lại, năng lực chấp nhận số nền tảng thể hiện hành vi mang tính bối
cảnh. Ở mô hình trực tiếp (M6), \(\hat{\beta}(\text{DAI}) = +0{,}104\)
(\(SE = 0{,}038\); \(p = {,}007\)), nhưng khi đưa đồng thời TCI vào (M7)
hệ số DAI suy giảm còn \(+0{,}077\) và chỉ còn ở biên ý nghĩa
(\(p = {,}048\)), cho thấy một phần phương sai năng suất mà DAI nắm bắt
đã trùng lặp với TCI. Sự suy giảm này quan trọng về nội dung vì nó bác
bỏ cách hiểu DAI như một phần thưởng năng suất phổ quát cho mọi doanh
nghiệp.

\hypertarget{cux1a1-chux1ebf-luxe1-chux1eafn-sux1ed1-tux1b0ux1a1ng-tuxe1c-tux1ea7ng-1tux1ea7ng-2-vuxe0-hiux1ec7u-ux1ee9ng-biuxean-theo-cux1b0ux1eddng-ux111ux1ed9-xuux1ea5t-khux1ea9u}{%
\subsection{4.3.6 Cơ chế lá chắn số: tương tác Tầng 1--Tầng 2 và hiệu
ứng biên theo cường độ xuất
khẩu}\label{cux1a1-chux1ebf-luxe1-chux1eafn-sux1ed1-tux1b0ux1a1ng-tuxe1c-tux1ea7ng-1tux1ea7ng-2-vuxe0-hiux1ec7u-ux1ee9ng-biuxean-theo-cux1b0ux1eddng-ux111ux1ed9-xuux1ea5t-khux1ea9u}}

Bằng chứng mạnh nhất về vai trò của DAI xuất hiện ở các đặc tả điều
tiết. Trong mô hình đầy đủ M8, số hạng trực tiếp DAI nhỏ và không có ý
nghĩa (\(\hat{\beta} = +0{,}019\); \(p = {,}705\)), tương tác bậc nhất
\(\text{FSTS} \times \text{DAI}\) âm và ở biên
(\(\hat{\beta} = -1{,}177\); \(p = {,}083\)), trong khi tương tác bậc
hai \(\text{FSTS}^2 \times \text{DAI}\) dương và có ý nghĩa thống kê
(\(\hat{\beta} = +3{,}119\); \(SE = 1{,}124\); \(p = {,}005\)). M8 cũng
đạt sức giải thích cao nhất trong các đặc tả chính, với
\(R^2 = 0{,}211\) và \(R^2\) hiệu chỉnh \(= 0{,}196\). Đây chính là cơ
chế lá chắn số: tương tác bậc hai dương làm tù đi nhánh suy giảm bên
phải, tức là DAI bảo vệ doanh nghiệp khỏi đà giảm hiệu quả ở mức quốc tế
hóa cao.

Bảng hiệu ứng biên của DAI theo từng mức FSTS làm rõ tính bối cảnh này.
Tại doanh nghiệp thuần nội địa (\(\text{FSTS} = 0\)), hiệu ứng biên của
một độ lệch chuẩn DAI chỉ là \(+0{,}080\) (\(p = {,}045\)); trên dải
xuất khẩu thấp và trung bình, hiệu ứng không phân biệt được với 0 (ví dụ
tại \(\text{FSTS} = 20\%\): \(-0{,}087\), \(p = {,}444\)); nhưng tại
nhóm xuất khẩu cường độ cao, hiệu ứng trở nên dương mạnh và có ý nghĩa:
\(+0{,}660\) tại \(\text{FSTS} = 70\%\) (\(p = {,}025\)) và \(+1{,}818\)
tại \(\text{FSTS} = 100\%\) (\(p = {,}002\)). Mô thức này xác nhận DAI
là một nguồn lực mở rộng quy mô có điều kiện, phát huy giá trị chính khi
doanh nghiệp đối mặt với nhu cầu phối hợp xuyên biên giới dày đặc hơn.
Dấu dương của số hạng điều tiết bậc hai được giữ nguyên qua sáu đặc tả
độ vững. Một kiểm định đảo chỗ chỉ báo (item-swap) củng cố ranh giới cấu
trúc giữa TCI và DAI: khi chuyển chỉ báo website (c22b) từ DAI sang TCI,
ý nghĩa chung của khối điều tiết DAI sụp đổ, với thống kê F chung giảm
từ \(4{,}56\) (\(p = {,}011\)) xuống \(1{,}88\) (\(p = {,}154\)). Ở mẫu
phụ chỉ gồm doanh nghiệp xuất khẩu (R5, \(N = 84\)), tương tác bậc hai
vẫn dương (\(\hat{\beta} = +2{,}821\)) với kiểm định F chung có ý nghĩa
(\(F = 6{,}32\); \(p = {,}003\)).

\hypertarget{ux111ux1ecbnh-lux1b0ux1ee3ng-cuxf4ng-suux1ea5t-thux1ed1ng-kuxea}{%
\subsection{4.3.7 Định lượng công suất thống
kê}\label{ux111ux1ecbnh-lux1b0ux1ee3ng-cuxf4ng-suux1ea5t-thux1ed1ng-kuxea}}

Cảnh báo về công suất thống kê cần được định lượng cụ thể. Hệ số Cohen's
\(f^2\) của hiệu ứng trực tiếp TCI (so sánh M2 với M5) xấp xỉ
\(0{,}036\), thuộc dải nhỏ-đến-trung bình; còn \(f^2\) của khối điều
tiết DAI (so sánh M7 với M8) chỉ khoảng \(0{,}018\), dưới ngưỡng hiệu
ứng nhỏ quy ước \(0{,}02\) của Cohen (1988). Với \(f^2 = 0{,}018\),
\(\alpha = {,}05\) và một bậc tự do tử số, để đạt công suất \(80\%\) cần
khoảng \(N > 430\): toàn mẫu phân tích (\(N = 617\)) đủ điều kiện, nhưng
mẫu phụ chỉ gồm doanh nghiệp xuất khẩu (\(N = 84\)) chỉ đạt công suất
khoảng \(16\%\). Do đó, các kết quả không có ý nghĩa về điều tiết DAI
trong mẫu phụ xuất khẩu không nên được hiểu là bằng chứng phản bác cơ
chế mở rộng quy mô có điều kiện, mà là hệ quả của cỡ mẫu chưa đủ để phát
hiện hiệu ứng nhỏ này.

\hypertarget{bux1ed9-hux1ec7-sux1ed1-ux111ux1ea7y-ux111ux1ee7-theo-chuux1ed7i-muxf4-huxecnh-phuxe2n-tux1ea7ng-vuxe0-tiux1ebfn-triux1ec3n-ux111ux1ed9-phuxf9-hux1ee3p}{%
\subsection{4.3.8 Bộ hệ số đầy đủ theo chuỗi mô hình phân tầng và tiến
triển độ phù
hợp}\label{bux1ed9-hux1ec7-sux1ed1-ux111ux1ea7y-ux111ux1ee7-theo-chuux1ed7i-muxf4-huxecnh-phuxe2n-tux1ea7ng-vuxe0-tiux1ebfn-triux1ec3n-ux111ux1ed9-phuxf9-hux1ee3p}}

Để minh bạch hóa cách hai cấu trúc năng lực tách dần khỏi đường cong
cường độ qua chuỗi mô hình lồng nhau, Bảng 4.5 hợp nhất bộ hệ số trọng
tâm của bảy đặc tả chính trên mẫu Singapore. Hạng cường độ xuất khẩu
tuyến tính dương và có ý nghĩa cao ở mọi đặc tả, dao động từ
\(+1{,}952\) (M7) đến \(+2{,}768\) (M4), trong khi hạng bậc hai chỉ đạt
ý nghĩa thống kê trong hai đặc tả có khối tương tác DAI (M4:
\(-2{,}959\), \(p < {,}01\); M8: \(-2{,}543\), \(p < {,}05\)) và ở biên
trong đặc tả nền M2 (\(-1{,}705\), \(p = {,}068\)).

\textbf{Bảng 4.5.} \emph{Tiến triển hệ số trọng tâm theo chuỗi mô hình
phân tầng, Singapore 2023 (biến phụ thuộc: ln năng suất lao động; sai số
chuẩn vững HC1).}

\begin{longtable}[]{@{}
  >{\raggedright\arraybackslash}p{(\columnwidth - 12\tabcolsep) * \real{0.1429}}
  >{\raggedleft\arraybackslash}p{(\columnwidth - 12\tabcolsep) * \real{0.1429}}
  >{\raggedleft\arraybackslash}p{(\columnwidth - 12\tabcolsep) * \real{0.1429}}
  >{\raggedleft\arraybackslash}p{(\columnwidth - 12\tabcolsep) * \real{0.1429}}
  >{\raggedleft\arraybackslash}p{(\columnwidth - 12\tabcolsep) * \real{0.1429}}
  >{\raggedleft\arraybackslash}p{(\columnwidth - 12\tabcolsep) * \real{0.1429}}
  >{\raggedleft\arraybackslash}p{(\columnwidth - 12\tabcolsep) * \real{0.1429}}@{}}
\toprule\noalign{}
\begin{minipage}[b]{\linewidth}\raggedright
Hệ số
\end{minipage} & \begin{minipage}[b]{\linewidth}\raggedleft
M2 (U ngược)
\end{minipage} & \begin{minipage}[b]{\linewidth}\raggedleft
M5 (+TCI)
\end{minipage} & \begin{minipage}[b]{\linewidth}\raggedleft
M6 (+DAI)
\end{minipage} & \begin{minipage}[b]{\linewidth}\raggedleft
M7 (T+D)
\end{minipage} & \begin{minipage}[b]{\linewidth}\raggedleft
M4 (DAI×)
\end{minipage} & \begin{minipage}[b]{\linewidth}\raggedleft
M8 (đầy đủ)
\end{minipage} \\
\midrule\noalign{}
\endhead
\bottomrule\noalign{}
\endlastfoot
FSTS (định tâm) & +2,652*** & +2,165** & +2,322*** & +1,952** &
+2,768*** & +2,409** \\
FSTS² & −1,705† & −1,168 & −1,389 & −0,965 & −2,959** & −2,543* \\
TCI (z) & --- & +0,168*** & --- & +0,153*** & --- & +0,153*** \\
DAI (z) & --- & --- & +0,104** & +0,077* & +0,046 & +0,019 \\
FSTS × DAI & --- & --- & --- & --- & −1,144 & −1,167† \\
FSTS² × DAI & --- & --- & --- & --- & +3,098** & +3,119** \\
\(R^2\) & 0,178 & 0,199 & 0,184 & 0,202 & 0,193 & 0,211 \\
\(R^2\) hiệu chỉnh & 0,168 & 0,189 & 0,174 & 0,190 & 0,180 & 0,196 \\
N & 623 & 623 & 617 & 617 & 617 & 617 \\
\end{longtable}

\emph{Nguồn: P4 Singapore, \texttt{p4/p4\_singapore\_en\_clean.md} (Bảng
2 bản thảo). † \(p < {,}10\); * \(p < {,}05\); ** \(p < {,}01\); ***
\(p < {,}001\). Sai số chuẩn vững HC1; FSTS định tâm trước khi bình
phương; mọi mô hình có quy mô, tuổi, sở hữu nước ngoài và hiệu ứng cố
định ngành. Các đặc tả chứa DAI dùng \(N = 617\) do sáu cơ sở thiếu giá
trị biến liên quan DAI.}

Tiến triển độ phù hợp mô hình củng cố trật tự đóng góp của hai cấu trúc.
Việc bổ sung TCI nâng \(R^2\) từ \(0{,}178\) (M2) lên \(0{,}199\) (M5),
một mức nhảy lớn hơn so với khi bổ sung riêng DAI (lên \(0{,}184\) ở
M6); mô hình đầy đủ M8 đạt sức giải thích cao nhất trong các đặc tả
chính (\(R^2 = 0{,}211\); \(R^2\) hiệu chỉnh \(= 0{,}196\)). Quỹ đạo suy
giảm của hệ số DAI từ M6 qua M7 đến M8
(\(+0{,}104 \to +0{,}077 \to +0{,}019\)) cho thấy phần lớn phương sai
năng suất mà DAI nắm bắt ở dạng hiệu ứng mức thực ra trùng lặp với TCI
hoặc được chuyển vào khối tương tác, bác bỏ cách hiểu DAI như một phần
thưởng năng suất phổ quát. Đặc trưng mẫu nền tảng cho cách đọc này:
trong \(N = 623\) cơ sở, năng suất lao động dạng log trung bình
\(11{,}42\) (độ lệch chuẩn \(1{,}18\)),
\(\overline{\mathrm{FSTS}} = 0{,}045\), và tương quan TCI--DAI chỉ ở mức
\(r = 0{,}31\), đủ thấp để biện minh cho việc giữ hai cấu trúc tách biệt
thay vì gộp thành một thước đo duy nhất.

\hypertarget{phuxe2n-xux1eed-ba-cux1a1-chux1ebf-cux1ea1nh-tranh-vuxe0-ux111ux1ed9-vux1eefng-trux1b0ux1edbc-hiux1ec7u-chux1ec9nh-chux1ecdn-lux1ef1a}{%
\subsection{4.3.9 Phân xử ba cơ chế cạnh tranh và độ vững trước hiệu
chỉnh chọn
lựa}\label{phuxe2n-xux1eed-ba-cux1a1-chux1ebf-cux1ea1nh-tranh-vuxe0-ux111ux1ed9-vux1eefng-trux1b0ux1edbc-hiux1ec7u-chux1ec9nh-chux1ecdn-lux1ef1a}}

Dấu dương của số hạng điều tiết bậc hai cho phép phân xử giữa ba cơ chế
cạnh tranh: thay thế (substitution), khuếch đại (amplification) và chọn
lựa (selection). Chữ ký thực nghiệm phân biệt khuếch đại với thay thế
nằm ở dấu của số hạng tương tác \emph{bậc hai}: hệ số dương
\(\mathrm{FSTS}^2 \times \mathrm{DAI} = +3{,}119\) (\(p = {,}005\)) hàm
ý đóng góp năng suất của DAI tăng siêu tuyến tính theo quy mô cường độ
xuất khẩu, điều mà một cơ chế thay thế dịch chuyển mức (chỉ dự đoán dịch
chuyển tuyến tính \(\mathrm{FSTS} \times \mathrm{DAI}\)) không thể tạo
ra. Tính bền vững của tương tác bậc hai dương qua các hạn chế mẫu (loại
siêu nhỏ R3, chỉ doanh nghiệp nhỏ và vừa R4, chỉ doanh nghiệp xuất khẩu
R5) lại bác bỏ một câu chuyện chọn lựa thuần túy, vì dưới chọn lựa,
tương quan đồng thời năng suất--DAI--FSTS lẽ ra phải suy giảm đáng kể
khi loại các doanh nghiệp dẫn dắt chọn lựa, song dấu dương được bảo toàn
xuyên suốt. Cụ thể, dấu dương và độ lớn của
\(\mathrm{FSTS}^2 \times \mathrm{DAI}\) giữ ổn định ở \(+3{,}521\) khi
loại doanh nghiệp siêu nhỏ (\(N = 464\)), \(+3{,}505\) khi giới hạn
doanh nghiệp nhỏ và vừa (\(N = 595\)) và \(+2{,}821\) ở mẫu chỉ doanh
nghiệp xuất khẩu với kiểm định F liên hợp có ý nghĩa (\(F = 6{,}32\),
\(p = {,}003\); \(N = 84\)).

Cơ chế khuếch đại còn được kiểm tra trước hiệu chỉnh chọn lựa biên gia
nhập. Vì thiết kế một đợt khảo sát Singapore 2023 không cung cấp công cụ
thỏa mãn ràng buộc loại trừ, một chẩn đoán dạng rút gọn ước lượng probit
giai đoạn một về tham gia xuất khẩu rồi đưa tỉ số Mills nghịch đảo vào
đặc tả M8: hệ số IMR (\(\hat{\beta} = 0{,}264\), \(SE = 0{,}138\),
\(t = 1{,}92\), \(p = {,}055\)) chỉ vừa vượt ngưỡng, cho thấy chọn lựa
nhẹ ở biên gia nhập, nhưng tương tác trọng tâm
\(\mathrm{FSTS}^2 \times \mathrm{DAI}\) và hệ số TCI gần như không đổi
(\(|\Delta| < 0{,}02\) trong mỗi trường hợp). Phù hợp với cảnh báo trong
tài liệu kinh tế lượng rằng hiệu chỉnh Heckman thiếu công cụ loại trừ
khả tín có thể chệch hơn cả ước lượng OLS cần sửa, đặc tả OLS được giữ
làm đặc tả chính. Tổng hợp lại, bằng chứng Singapore ủng hộ khuếch đại
là cơ chế trọng tâm, giữ chọn lựa như một quan ngại dư có giới hạn, và
đặt Singapore như một trường hợp biên ở trần thể chế nơi bão hòa số nội
địa xóa hiệu ứng phần thưởng phổ quát còn tầng giao dịch Tầng 2 sẵn có
về mặt thể chế làm cơ chế mở rộng quy mô có điều kiện trở nên nhận diện
được về thực nghiệm.

\begin{center}\rule{0.5\linewidth}{0.5pt}\end{center}

\hypertarget{kux1ebft-quux1ea3-bux1ed1i-cux1ea3nh-thux1ec3-chux1ebf-trung-buxecnh-cao-trung-quux1ed1c-20122024-p5}{%
\section{4.4 Kết quả bối cảnh thể chế trung bình cao: Trung Quốc,
2012--2024
(P5)}\label{kux1ebft-quux1ea3-bux1ed1i-cux1ea3nh-thux1ec3-chux1ebf-trung-buxecnh-cao-trung-quux1ed1c-20122024-p5}}

\hypertarget{xuxe1c-nhux1eadn-quan-hux1ec7-phi-tuyux1ebfn-vuxe0-ux111iux1ec3m-uux1ed1n}{%
\subsection{4.4.1 Xác nhận quan hệ phi tuyến và điểm
uốn}\label{xuxe1c-nhux1eadn-quan-hux1ec7-phi-tuyux1ebfn-vuxe0-ux111iux1ec3m-uux1ed1n}}

Phân tích hai giai đoạn trên dữ liệu WBES Trung Quốc (Nhóm III, trung
bình cao) giai đoạn 2012--2024 xác nhận dạng quan hệ phi tuyến chữ U
ngược giữa FSTS và hiệu quả hoạt động doanh nghiệp. Điểm điểm uốn được
ước lượng tại:

\begin{itemize}
\tightlist
\item
  \textbf{Năm 2012}: TP = 49,37\% FSTS
\item
  \textbf{Năm 2024}: TP = 47,19\% FSTS
\item
  \textbf{gộp (toàn mẫu)}: TP = 48,78\% FSTS
\end{itemize}

Các con số này cho thấy doanh nghiệp Trung Quốc đạt hiệu quả tối ưu khi
cường độ quốc tế hóa (FSTS) ở mức gần 49\%, một mức độ tập trung quốc tế
vừa phải trong bối cảnh nền kinh tế nội địa khổng lồ của Trung Quốc cũng
cung cấp cơ hội tăng trưởng song song.

\hypertarget{kiux1ec3m-ux111ux1ecbnh-tuxednh-ux1ed5n-ux111ux1ecbnh-cux1ee7a-ux111iux1ec3m-uux1ed1n-theo-thux1eddi-gian}{%
\subsection{4.4.2 Kiểm định tính ổn định của điểm uốn theo thời
gian}\label{kiux1ec3m-ux111ux1ecbnh-tuxednh-ux1ed5n-ux111ux1ecbnh-cux1ee7a-ux111iux1ec3m-uux1ed1n-theo-thux1eddi-gian}}

Luận án đặt câu hỏi thực nghiệm quan trọng: liệu hình dạng phi tuyến và
điểm uốn có thay đổi qua thập kỷ 2012--2024, giai đoạn Trung Quốc chứng
kiến nhiều thay đổi chính sách thương mại, căng thẳng địa chính trị và
chuyển đổi số mạnh mẽ? Kiểm định Paternoster cho kết quả
\(z = +0{,}82\), \(p = {,}412\) (vế FSTS) và \(z = -0{,}61\),
\(p = {,}545\) (vế FSTS²), không có ý nghĩa thống kê. Điều này xác nhận
rằng \textbf{điểm uốn ổn định qua thời gian}: sự dịch chuyển từ TP =
49,37\% (2012) xuống TP = 47,19\% (2024) là không đáng kể về mặt thống
kê.

Kiểm định tương tác thời gian (Temporal H6): \(F = 1{,}83\),
\(p = {,}176\), không có ý nghĩa thống kê, nghĩa là hình dạng tổng thể
của quan hệ phi tuyến không biến đổi đáng kể qua giai đoạn phân tích.
Phát hiện này cho thấy trong bối cảnh trung bình cao như Trung Quốc, cơ
chế căn bản của quan hệ quốc tế hóa và hiệu quả duy trì tính ổn định
đáng kể, có thể phản ánh sự bù trừ giữa chi phí thích ứng tăng lên do
chính sách thương mại khó đoán và năng lực tổ chức ngày càng trưởng
thành của các doanh nghiệp xuất khẩu Trung Quốc.

\hypertarget{vai-truxf2-cux1ee7a-tci-vuxe0-dai}{%
\subsection{4.4.3 Vai trò của TCI và
DAI}\label{vai-truxf2-cux1ee7a-tci-vuxe0-dai}}

Tương tác \(FSTS \times TCI\) không đạt ý nghĩa thống kê (null
moderation, \(p > {,}10\)), cho thấy chỉ số năng lực công nghệ (TCI)
không phát huy vai trò điều tiết đáng kể ở bối cảnh Trung Quốc theo mô
hình P5. Tương tự, DAI cũng không cho thấy điều tiết có ý nghĩa trong
bối cảnh này. kết quả không có ý nghĩa về điều tiết tại Trung Quốc có
thể phản ánh thực tế rằng năng lực công nghệ đã được phân phối tương đối
đồng đều hơn trong mẫu doanh nghiệp xuất khẩu Trung Quốc, làm giảm
phương sai đủ để phát hiện hiệu ứng điều tiết.

\hypertarget{bux1ed9-hux1ec7-sux1ed1-muxf4-huxecnh-ngux1b0ux1ee1ng-chuxednh-m2-vuxe0-ux111ux1ed9-lux1edbn-kinh-tux1ebf}{%
\subsection{4.4.4 Bộ hệ số mô hình ngưỡng chính (M2) và độ lớn kinh
tế}\label{bux1ed9-hux1ec7-sux1ed1-muxf4-huxecnh-ngux1b0ux1ee1ng-chuxednh-m2-vuxe0-ux111ux1ed9-lux1edbn-kinh-tux1ebf}}

Mô hình ngưỡng chính M2 (đặc tả
\(\ln(\mathrm{LP}) \sim \mathrm{FSTS} + \mathrm{FSTS}^2 + \ln(\mathrm{Emp}) + \text{firmage} + \text{foreigndummy}\))
được ước lượng riêng cho từng đợt khảo sát và mẫu gộp, củng cố cho hình
dạng chữ U ngược bằng một bộ hệ số đầy đủ. Hạng cường độ xuất khẩu tuyến
tính dương trong mọi mẫu: \(\hat{\beta}_1 = +2{,}07\) năm 2012
(\(p < {,}001\)), \(+1{,}50\) năm 2024 (\(p = {,}010\)) và \(+1{,}78\)
trong mẫu gộp (\(p < {,}001\)). Hạng bậc hai âm tương ứng:
\(\hat{\beta}_2 = -2{,}09\) năm 2012 (\(p < {,}001\)), \(-1{,}59\) năm
2024 (\(p = {,}026\)) và \(-1{,}83\) gộp (\(p < {,}001\)). Kiểm định U
Lind--Mehlum xác nhận chữ U ngược một cách chính thức với \(p < {,}001\)
trong cả mẫu 2012 và mẫu gộp, và \(p = {,}037\) trong mẫu 2024.

Khoảng tin cậy 95\% theo phương pháp delta của điểm uốn là
\([43{,}2\%; 55{,}6\%]\) năm 2012, \([34{,}5\%; 59{,}9\%]\) năm 2024 và
\([42{,}7\%; 54{,}9\%]\) trong mẫu gộp. Sự chồng lấp chặt chẽ của ba
khoảng tin cậy này là cơ sở định lượng trực tiếp cho kết luận về tính ổn
định theo thời gian đã trình bày ở Mục 4.4.2. Về độ lớn kinh tế, tại
điểm uốn đặc thù đợt khảo sát, \(\ln(\mathrm{LP})\) dự đoán cao hơn
khoảng \(0{,}51\) điểm log năm 2012 và \(0{,}35\) điểm log năm 2024 so
với mức cơ sở \(\mathrm{FSTS} = 0\), tức tương ứng khoảng \(67\%\) và
\(42\%\) chênh lệch năng suất ở mức kiểm soát trung bình hình học. Ngược
lại, tại \(\mathrm{FSTS} = 1\) (một nhà xuất khẩu 100\% giả định), năng
suất dự đoán gần như ngang bằng mức cơ sở trong cả hai đợt, minh hoạ
trực quan logic ``tối ưu có giới hạn'' nằm ở trung tâm của giả thuyết
H1.

\hypertarget{quy-muxf4-mux1eabu-ux111ux1eb7c-trux1b0ng-nhuxf3m-doanh-nghiux1ec7p-vuxe0-phuxe2n-tuxe1ch-tux103ng-trux1b0ux1edfng-nux103ng-suux1ea5t}{%
\subsection{4.4.5 Quy mô mẫu, đặc trưng nhóm doanh nghiệp và phân tách
tăng trưởng năng
suất}\label{quy-muxf4-mux1eabu-ux111ux1eb7c-trux1b0ng-nhuxf3m-doanh-nghiux1ec7p-vuxe0-phuxe2n-tuxe1ch-tux103ng-trux1b0ux1edfng-nux103ng-suux1ea5t}}

Mẫu phân tích cơ sở gồm \(N = 2.610\) quan sát năm 2012 và \(N = 1.934\)
năm 2024, với mẫu gộp \(N = 4.544\) trong đặc tả ngưỡng chính. Cường độ
xuất khẩu trung bình thấp và lệch mạnh: \(6{,}9\%\) năm 2012 và
\(5{,}2\%\) năm 2024, với chỉ \(15{,}4\%\) doanh nghiệp năm 2012 và
\(12{,}6\%\) năm 2024 báo cáo cường độ xuất khẩu dương. Năng suất lao
động trung bình dạng log tăng từ \(12{,}52\) (2012) lên \(13{,}01\)
(2024), một dịch chuyển mức \(0{,}49\) điểm log tương đương năng suất
cao hơn khoảng \(64\%\) về mức (\(\exp(0{,}49) \approx 1{,}64\)). Một
phân tách Oaxaca--Blinder không chính thức quy khoảng \(70\%\) dịch
chuyển này cho tăng trưởng năng suất trong cùng ngành và khoảng \(30\%\)
cho thay đổi thành phần ngành.

Năng lực công nghệ vận hành bền vững như một yếu tố dịch chuyển mức trực
tiếp chứ không phải yếu tố điều tiết độ cong. Hệ số TCI chuẩn hoá z
trong cùng đợt là \(\hat{\beta}_z = +0{,}260\) năm 2012, \(+0{,}426\)
năm 2024 và \(+0{,}380\) trong mẫu gộp (đều \(p < {,}001\)), với kiểm
định Paternoster về thay đổi chéo đợt khảo sát cho \(z = -2{,}55\),
\(p = {,}011\), cho thấy liên hệ giữa năng lực công nghệ và năng suất
được củng cố có ý nghĩa giữa hai đợt.

\hypertarget{cux1ea5u-truxfac-kiux1ec3m-ux111ux1ecbnh-ux111iux1ec1u-tiux1ebft-ba-chiux1ec1u-vuxe0-cux1a1-chux1ebf-vux1ed1n-lux1b0u-ux111ux1ed9ng}{%
\subsection{4.4.6 Cấu trúc kiểm định điều tiết ba chiều và cơ chế vốn
lưu
động}\label{cux1ea5u-truxfac-kiux1ec3m-ux111ux1ecbnh-ux111iux1ec1u-tiux1ebft-ba-chiux1ec1u-vuxe0-cux1a1-chux1ebf-vux1ed1n-lux1b0u-ux111ux1ed9ng}}

Kết quả điều tiết không có ý nghĩa trình bày trong Mục 4.4.3 được nâng
đỡ bởi một bộ ba kiểm định F liên hợp trên đặc tả điều tiết ba chiều gộp
(\(N = 3.559\)). Kiểm định dịch chuyển chéo đợt khảo sát cho
\(F(2, 3.558) = 2{,}24\), \(p = {,}107\) (không bác bỏ tính bằng nhau);
kiểm định điều tiết độ cong theo năng lực cho \(F(2, 3.558) = 3{,}26\),
\(p = {,}039\) (chỉ ở mức biên); và kiểm định điều tiết động phụ thuộc
năng lực cho \(F(2, 3.558) = 0{,}27\), \(p = {,}760\) (không hỗ trợ).
Đáng chú ý, dù kiểm định F liên hợp về điều tiết độ cong đạt ý nghĩa
biên, không một hệ số tương tác riêng lẻ nào phân biệt được với 0:
\(\text{TCI} \times \mathrm{FSTS}\) cho \(\beta = -0{,}414\),
\(p = {,}443\) và \(\text{TCI} \times \mathrm{FSTS}^2\) cho
\(\beta = +0{,}051\), \(p = {,}934\), củng cố cách diễn giải rằng năng
lực công nghệ là yếu tố dịch chuyển mức chứ không phải yếu tố điều tiết
ngưỡng.

Việc khảo sát cơ chế vốn lưu động (working capital) cũng không nhận dạng
được một yếu tố điều tiết bền vững. Tương tác thấu chi
\(\times \mathrm{FSTS}^2\) đảo dấu giữa hai đợt (\(\beta = +0{,}27\),
\(p = {,}050\) năm 2012; \(\beta = -0{,}38\), \(p = {,}081\) năm 2024);
tương tác tín dụng thương mại \(\times \mathrm{FSTS}^2\) âm có ý nghĩa
năm 2012 (\(\beta = -0{,}017\), \(p = {,}009\)) nhưng không có ý nghĩa
năm 2024 (\(p = {,}660\)). Như vậy, cơ chế bẫy vốn lưu động vẫn là cơ sở
lý thuyết nhất quán nhất cho đoạn suy giảm sau ngưỡng nhưng chưa được
nhận dạng trực tiếp với công cụ WBES hiện có, đòi hỏi dữ liệu vi mô tài
chính như thời gian chu kỳ chuyển đổi tiền mặt hoặc số ngày khoản phải
thu chưa thanh toán cho nghiên cứu tương lai.

\hypertarget{kiux1ec3m-ux111ux1ecbnh-ux111ux1ed9-vux1eefng-mux1eabu-phux1ee5-chux1ec9-nguxe0nh-sux1ea3n-xuux1ea5t}{%
\subsection{4.4.7 Kiểm định độ vững: mẫu phụ chỉ ngành sản
xuất}\label{kiux1ec3m-ux111ux1ecbnh-ux111ux1ed9-vux1eefng-mux1eabu-phux1ee5-chux1ec9-nguxe0nh-sux1ea3n-xuux1ea5t}}

Để loại trừ quan ngại trộn lẫn ngành, mô hình M2 được tái ước lượng trên
mẫu phụ chỉ ngành sản xuất (\(N = 1.656\) năm 2012, \(N = 1.062\) năm
2024, gộp \(N = 2.718\)). Dưới ràng buộc này, điểm uốn dịch về khoảng
\(42\%\) năm 2012 (khoảng tin cậy \([37{,}8; 46{,}8]\)) và \(30\%\) năm
2024 (khoảng tin cậy \([15{,}1; 44{,}2]\)), khoảng cuối rộng hơn nhiều
do sức mạnh thống kê giảm, song kiểm định Paternoster vẫn không bác bỏ
tính bằng nhau chéo đợt khảo sát (\(z = +1{,}51\), \(p = {,}130\) cho
FSTS; \(z = -0{,}96\), \(p = {,}337\) cho \(\mathrm{FSTS}^2\)). Khi bổ
sung tầng ngành như hiệu ứng cố định vào M2 gộp, hệ số FSTS giảm từ
\(+1{,}78\) xuống \(+1{,}56\) (giảm \(12{,}8\%\)) và hệ số
\(\mathrm{FSTS}^2\) từ \(-1{,}83\) xuống \(-1{,}55\) (giảm
\(15{,}2\%\)), nhưng chữ U ngược được bảo toàn và điểm uốn vẫn nằm trong
dải vận hành an toàn 30--60\%.

\hypertarget{bux1eb1ng-chux1ee9ng-bux1ed5-trux1ee3-cux1ea5p-doanh-nghiux1ec7p-nhux1ecf-vuxe0-vux1eeba-chux1eef-u-ngux1b0ux1ee3c-vuxe0-cux1a1-chux1ebf-bux1eaby-vux1ed1n-lux1b0u-ux111ux1ed9ng-p2}{%
\subsection{4.4.8 Bằng chứng bổ trợ cấp doanh nghiệp nhỏ và vừa: chữ U
ngược và cơ chế bẫy vốn lưu động
(P2)}\label{bux1eb1ng-chux1ee9ng-bux1ed5-trux1ee3-cux1ea5p-doanh-nghiux1ec7p-nhux1ecf-vuxe0-vux1eeba-chux1eef-u-ngux1b0ux1ee3c-vuxe0-cux1a1-chux1ebf-bux1eaby-vux1ed1n-lux1b0u-ux111ux1ed9ng-p2}}

Nghiên cứu thành phần P2 cung cấp một kiểm định độc lập về quan hệ chữ U
ngược trong bối cảnh Trung Quốc, trên một lát cắt khác của tổng thể:
\textbf{4.290 quan sát doanh nghiệp và năm} thuộc các doanh nghiệp sản
xuất nhỏ và vừa (mã ngành ISIC 15--37). Cường độ xuất khẩu trung bình
của mẫu là 23,7\% (độ lệch chuẩn 28,4\%), trong đó 18\% doanh nghiệp có
cường độ xuất khẩu vượt 50\%, tạo nên một phân bố trải rộng từ thuần nội
địa đến định hướng xuất khẩu cao, phù hợp để định vị điểm uốn.

Ước lượng đặc tả bậc hai với hiệu ứng cố định ngành và năm xác nhận chữ
U ngược: hệ số nhánh đi lên dương và có ý nghĩa cao
(\(\hat{\beta}_1 = +1{,}704\), \(p < {,}001\)), hệ số nhánh đi xuống âm
và có ý nghĩa, định vị điểm uốn ở mức \textbf{47,8\% cường độ xuất
khẩu}. Giá trị này nằm sát ngưỡng của nghiên cứu Trung Quốc theo chuỗi
thời gian (P5: 48,78\%) và trong dải vận hành an toàn 30--60\% được xác
nhận xuyên suốt luận án. Đóng góp đặc thù của P2 là cơ chế giải thích
nhánh đi xuống: \textbf{bẫy vốn lưu động}. Khi cường độ xuất khẩu vượt
47,8\%, chi phí biên của việc tài trợ chu kỳ chuyển đổi tiền mặt kéo
dài, do doanh nghiệp xuất khẩu phải ứng vốn cho sản xuất, đóng gói và
vận chuyển từ rất lâu trước khi thu được khoản phải thu nước ngoài, bắt
đầu lấn át lợi ích từ quy mô và học hỏi qua xuất khẩu, kéo năng suất đi
xuống. Cơ chế này bổ sung cho cách đọc dựa trên chi phí điều phối và
năng lực, làm rõ một kênh tài chính cụ thể đứng sau độ cong ở các doanh
nghiệp nhỏ và vừa thâm dụng vốn lưu động. Nhất quán với phần còn lại của
luận án, năng lực số và đặc điểm giới tính của nhà quản trị trong P2 vận
hành như yếu tố nâng mặt bằng năng suất chứ không uốn lại hình dạng
đường cong.

\begin{center}\rule{0.5\linewidth}{0.5pt}\end{center}

\hypertarget{kux1ebft-quux1ea3-bux1ed1i-cux1ea3nh-chuyux1ec3n-ux111ux1ed5i-bux1eadc-thux1ea5ptrung-viux1ec7t-nam-p3}{%
\section{4.5 Kết quả bối cảnh chuyển đổi bậc thấp--trung: Việt Nam
(P3)}\label{kux1ebft-quux1ea3-bux1ed1i-cux1ea3nh-chuyux1ec3n-ux111ux1ed5i-bux1eadc-thux1ea5ptrung-viux1ec7t-nam-p3}}

\hypertarget{xuxe1c-nhux1eadn-quan-hux1ec7-phi-tuyux1ebfn-vuxe0-hux1ec7-sux1ed1-hux1ed3i-quy}{%
\subsection{4.5.1 Xác nhận quan hệ phi tuyến và hệ số hồi
quy}\label{xuxe1c-nhux1eadn-quan-hux1ec7-phi-tuyux1ebfn-vuxe0-hux1ec7-sux1ed1-hux1ed3i-quy}}

Phân tích hồi quy bảng (panel regression) trên ba làn sóng khảo sát WBES
Việt Nam (2009, 2015, 2023) xác nhận dạng phi tuyến chữ U ngược. Mô hình
M2, mô hình phi tuyến chính (mẫu gộp, \(N = 2.958\)), cho kết quả:

\[\hat{\beta}(\text{FSTS}_c) = +0{,}984, \quad \hat{\beta}(\text{FSTS}_c^2) = -1{,}909\]

Cả hai hệ số đều có ý nghĩa thống kê, với hệ số bậc nhất dương và bậc
hai âm xác nhận đường cong chữ U ngược. Kiểm định hình thức bằng phương
pháp Lind--Mehlum (U-test) cho kết quả \(p < {,}001\), bằng chứng thống
kê mạnh nhất trong tất cả các mẫu quốc gia của luận án.

\hypertarget{ux111iux1ec3m-uux1ed1n-theo-luxe0n-suxf3ng}{%
\subsection{4.5.2 Điểm uốn theo làn
sóng}\label{ux111iux1ec3m-uux1ed1n-theo-luxe0n-suxf3ng}}

Điểm uốn được ước lượng riêng cho từng làn sóng và cho gộp:

\begin{longtable}[]{@{}ll@{}}
\toprule\noalign{}
Làn sóng & điểm uốn \\
\midrule\noalign{}
\endhead
\bottomrule\noalign{}
\endlastfoot
2009 & 46,2\% FSTS \\
2015 & 39,3\% FSTS \\
2023 & 41,6\% FSTS \\
gộp & 39,7\% FSTS \\
\end{longtable}

Điểm điểm uốn trong khoảng \textbf{39--46\% FSTS} đặt Việt Nam ở mức
``thấp hơn đáng kể'' so với Singapore (TP ≈ 88,6\%) và tương đương hoặc
thấp hơn Trung Quốc (TP ≈ 48,78\%). Kiểm định Paternoster so sánh điểm
uốn giữa 2009 và 2015 cho kết quả \(z = 3{,}353\), \(p < {,}001\), có ý
nghĩa thống kê, cho thấy điểm uốn đã dịch chuyển đáng kể từ 46,2\% xuống
39,3\% giữa hai làn sóng.

\hypertarget{ux111iux1ec1u-tiux1ebft-bux1edfi-tci-tuxedch-cux1ef1c-vuxe0-cuxf3-uxfd-nghux129a}{%
\subsection{4.5.3 Điều tiết bởi TCI: Tích cực và có ý
nghĩa}\label{ux111iux1ec1u-tiux1ebft-bux1edfi-tci-tuxedch-cux1ef1c-vuxe0-cuxf3-uxfd-nghux129a}}

Mô hình M3 với tương tác \(FSTS \times TCI\) cho kết quả \textbf{dương
và có ý nghĩa thống kê}: TCI làm nâng cao mặt bằng hiệu quả và khuếch
đại tác động tích cực của quốc tế hóa. Doanh nghiệp Việt Nam có năng lực
công nghệ tốt hơn (đo bằng TCI) đạt hiệu quả cao hơn từ quá trình quốc
tế hóa, nhất quán với H2 trong khung lý thuyết CDCM.

\hypertarget{ux111iux1ec1u-tiux1ebft-bux1edfi-dai-khuxf4ng-cuxf3-uxfd-nghux129a}{%
\subsection{4.5.4 Điều tiết bởi DAI: Không có ý
nghĩa}\label{ux111iux1ec1u-tiux1ebft-bux1edfi-dai-khuxf4ng-cuxf3-uxfd-nghux129a}}

Mô hình M4 với tương tác \(FSTS \times DAI\) (chỉ bao gồm Tầng 1 DAI do
giới hạn đo lường WBES trong các làn sóng này) cho kết quả yếu dương
nhưng không có ý nghĩa thống kê. Kết quả này có thể phản ánh hạn chế về
đo lường: Tầng 1 DAI chỉ nắm bắt mức độ chấp nhận số bề mặt (website,
email, online sales) chứ không đo lường được chiều sâu năng lực số. Với
dữ liệu WBES Việt Nam chỉ cung cấp Tầng 1, không đủ phân biệt để phát
hiện hiệu ứng điều tiết tinh tế hơn.

\hypertarget{bux1ed9-hux1ec7-sux1ed1-ux111ux1ea7y-ux111ux1ee7-theo-luxe0n-suxf3ng-vuxe0-quux1ef9-ux111ux1ea1o-ux111iux1ec1u-tiux1ebft}{%
\subsection{4.5.5 Bộ hệ số đầy đủ theo làn sóng và quỹ đạo điều
tiết}\label{bux1ed9-hux1ec7-sux1ed1-ux111ux1ea7y-ux111ux1ee7-theo-luxe0n-suxf3ng-vuxe0-quux1ef9-ux111ux1ea1o-ux111iux1ec1u-tiux1ebft}}

Bộ hệ số chữ U ngược M2 nhất quán qua cả ba làn sóng. Hạng tuyến tính
\(\mathrm{FSTS}^{c}\) dương trong mọi đợt (\(\hat{\beta}_1 = +1{,}045\)
năm 2009, \(p = {,}015\); \(+1{,}159\) năm 2015, \(p = {,}029\);
\(+0{,}962\) năm 2023, \(p = {,}039\); \(+0{,}984\) gộp,
\(p < {,}001\)), trong khi hạng bậc hai âm (\(\hat{\beta}_2 = -1{,}823\)
năm 2009, \(p = {,}005\); \(-2{,}115\) năm 2015, \(p = {,}004\);
\(-1{,}686\) năm 2023, \(p = {,}008\); \(-1{,}909\) gộp,
\(p < {,}001\)). Độ cong sắc nét nhất quan sát được trong làn sóng 2015
(\(\hat{\beta}_2 = -2{,}115\)), nhất quán với một giai đoạn chuyển đổi
trong đó kênh số bị nén hoàn toàn. Khoảng tin cậy 95\% theo phương pháp
delta cho từng điểm uốn lần lượt là \([37{,}4\%; 55{,}1\%]\) năm 2009,
\([30{,}3\%; 48{,}4\%]\) năm 2015, \([31{,}7\%; 51{,}5\%]\) năm 2023 và
\([34{,}0\%; 45{,}5\%]\) trong mẫu gộp.

Năng lực công nghệ \(\tilde{\mathrm{TCI}}\) có liên hệ trực tiếp thuận
chiều và bền vững qua cả ba đợt (\(\hat{\beta} = +0{,}215\) năm 2009,
\(+0{,}128\) năm 2015, \(+0{,}123\) năm 2023, đều có ý nghĩa) và trong
mẫu gộp (\(\hat{\beta} = +0{,}179\), \(p < {,}001\)), với vai trò điều
tiết phân biệt được trong ba trong bốn bảng (kiểm định liên hợp M3:
\(p = {,}040\) năm 2009, \(p = {,}027\) năm 2023, \(p = {,}003\) gộp,
chỉ \(p = {,}713\) năm 2015). Trong mẫu gộp, tương tác tuyến tính
\(\mathrm{FSTS}^{c} \times \tilde{\mathrm{TCI}} = -0{,}587\)
(\(p = {,}003\)) và bậc hai
\((\mathrm{FSTS}^{c})^2 \times \tilde{\mathrm{TCI}} = +0{,}640\)
(\(p = {,}031\)), cho thấy chữ U ngược phẳng lại đối với doanh nghiệp có
năng lực cao chứ không dịch chuyển về mức.

\hypertarget{tuxednh-phux1ee5-thuux1ed9c-giai-ux111oux1ea1n-cux1ee7a-dai-vuxe0-kiux1ec3m-ux111ux1ecbnh-paternoster-chuxe9o-luxe0n-suxf3ng}{%
\subsection{4.5.6 Tính phụ thuộc giai đoạn của DAI và kiểm định
Paternoster chéo làn
sóng}\label{tuxednh-phux1ee5-thuux1ed9c-giai-ux111oux1ea1n-cux1ee7a-dai-vuxe0-kiux1ec3m-ux111ux1ecbnh-paternoster-chuxe9o-luxe0n-suxf3ng}}

Trái với năng lực công nghệ, áp dụng công nghệ số nền tảng
\(\tilde{\mathrm{DAI}}\) theo một quỹ đạo không đơn điệu: mạnh nhất năm
2009 (\(\hat{\beta} = +0{,}175\), \(p < {,}001\)), mất ý nghĩa năm 2015
(\(\hat{\beta} = -0{,}044\), \(p = {,}377\)), tái xuất hiện năm 2023
(\(\hat{\beta} = +0{,}095\), \(p = {,}038\)) và dương trong mẫu gộp
(\(\hat{\beta} = +0{,}078\), \(p = {,}004\)). Kiểm định Paternoster xác
nhận các dịch chuyển này phân biệt được về mặt thống kê: mức giảm từ
2009 đến 2015 cho \(z = 3{,}35\) (\(p = {,}001\)) và sự phục hồi từ 2015
đến 2023 cho \(z = -2{,}05\) (\(p = {,}040\)), trong khi các tham số độ
cong \(\mathrm{FSTS}^{c}\) và \((\mathrm{FSTS}^{c})^2\) không phân biệt
được giữa các đợt (mọi \(p > {,}33\)). Điều tiết DAI tập trung hoàn toàn
trong làn sóng 2023, nơi
\(\mathrm{FSTS}^{c} \times \tilde{\mathrm{DAI}} = -0{,}912\)
(\(p = {,}043\)) và kiểm định liên hợp M8 ở mức cận ngưỡng
(\(p = {,}062\)). Theo phân loại Haans, Pieters và He (2016), mẫu hình
2023 phù hợp với điều tiết Loại I (làm phẳng độ dốc, giữ nguyên hình
dạng chữ U ngược) chứ không phải Loại II.

\hypertarget{phuxe2n-ruxe3-tham-giacux1b0ux1eddng-ux111ux1ed9-vuxe0-bux1eb1ng-chux1ee9ng-nux1ed9i-sinh}{%
\subsection{4.5.7 Phân rã tham gia--cường độ và bằng chứng nội
sinh}\label{phuxe2n-ruxe3-tham-giacux1b0ux1eddng-ux111ux1ed9-vuxe0-bux1eb1ng-chux1ee9ng-nux1ed9i-sinh}}

Một phát hiện thực nghiệm cốt lõi của P3 là chữ U ngược toàn mẫu được
nhận diện chủ yếu thông qua biên tham gia chứ không qua độ cong trong
nhóm doanh nghiệp xuất khẩu. Tỉ lệ doanh nghiệp báo cáo xuất khẩu trực
tiếp dương giảm từ \(28{,}4\%\) năm 2009 xuống \(20{,}7\%\) năm 2015 và
\(18{,}8\%\) năm 2023, và chỉ khoảng \(1{,}0\%\) doanh nghiệp gộp (29
trong 2.958) nằm trong dải \(\pm 5\) điểm phần trăm quanh điểm uốn. Khi
tái ước lượng trên nhóm phụ chỉ doanh nghiệp xuất khẩu
(\(\mathrm{FSTS} > 0\); gộp \(N = 669\)), hạng tuyến tính trở nên âm
(\(\mathrm{FSTS}^{c}\) \(\hat{\beta} = -0{,}861\), \(p < {,}001\)) còn
hạng bậc hai mất ý nghĩa (\((\mathrm{FSTS}^{c})^2\)
\(\hat{\beta} = -0{,}200\), \(p = {,}660\); kiểm định liên hợp
\(p = {,}462\)). Do đó độ cong toàn mẫu phản ánh chủ yếu bước nhảy năng
suất khi doanh nghiệp chuyển từ không xuất khẩu sang xuất khẩu (H1a),
chứ không phải độ cong cường độ trong nhóm doanh nghiệp xuất khẩu (H1b).

Bằng chứng nội sinh và chọn lựa làm rõ thêm sự phân biệt giữa hai kênh
năng lực. Hiệu chỉnh Heckman hai bước cho tỉ số Mills nghịch đảo
(inverse Mills ratio) không có ý nghĩa qua mọi đợt (mọi
\(|\lambda| < 0{,}84\), \(p > {,}25\)), cho thấy không phát hiện thiên
lệch chọn lựa trên nhóm phụ doanh nghiệp xuất khẩu. Ước lượng biến công
cụ hai giai đoạn (2SLS) với công cụ là tỉ lệ áp dụng đồng cấp ngành ×
vùng × đợt (thống kê F giai đoạn một 22--35, vượt xa ngưỡng
Staiger--Stock) cho hệ số DAI được công cụ hoá bằng không
(\(\hat{\beta} = 0{,}018\), \(p = {,}942\)) nhưng hệ số TCI được công cụ
hoá dương mạnh (\(\hat{\beta} = 1{,}639\), \(p < {,}001\)). Hệ số DAI
2SLS bằng không xác nhận rằng hiện diện website không gây ra cải thiện
năng suất một cách khả tín ở Việt Nam 2023, củng cố cách diễn giải
``tính lỗi thời của biến đại diện Tầng 1'': đến năm 2023 tỉ lệ doanh
nghiệp có website đạt \(49{,}8\%\) so với \(42{,}5\%\) năm 2009, khiến
chỉ báo này gần như phổ quát và mất sức mạnh phân biệt. Đối sánh điểm xu
hướng (propensity score matching) xác nhận thêm: tác động bình quân lên
nhóm được xử lý cho sở hữu website là \(0{,}298\)--\(0{,}321\)
(\(p < {,}001\)) và cho công nghệ nước ngoài hoặc chứng nhận là
\(0{,}609\)--\(0{,}637\) (\(p < {,}001\)).

\hypertarget{cux1ea5u-truxfac-lux1ea1m-phuxe1t-sux1ed1-khuxf4ng-cux1ee7a-fsts-vuxe0-hux1ec7-quux1ea3-nhux1eadn-dux1ea1ng-hai-biuxean}{%
\subsection{4.5.8 Cấu trúc lạm phát số không của FSTS và hệ quả nhận
dạng hai
biên}\label{cux1ea5u-truxfac-lux1ea1m-phuxe1t-sux1ed1-khuxf4ng-cux1ee7a-fsts-vuxe0-hux1ec7-quux1ea3-nhux1eadn-dux1ea1ng-hai-biuxean}}

Bản chất nhận dạng của chữ U ngược Việt Nam chỉ trở nên minh bạch khi
đặt vào cấu trúc phân phối của biến cường độ xuất khẩu. Biến FSTS bị
chặn dưới ở 0 và lạm phát số không (zero-inflated) nghiêm trọng: tỉ
trọng doanh nghiệp không xuất khẩu là \(71{,}6\%\) năm 2009,
\(79{,}3\%\) năm 2015 và \(81{,}2\%\) năm 2023, mức gộp \(77{,}4\%\). Vì
khối lượng quan sát dồn ở \(\mathrm{FSTS} = 0\) trong khi điểm uốn nằm
quanh 40\%, quan hệ quốc tế hóa và hiệu quả phân tách về mặt khái niệm
thành hai biên độc lập: biên tham gia, tức bước nhảy từ
\(\mathrm{FSTS} = 0\) sang \(\mathrm{FSTS} > 0\) nơi doanh nghiệp vượt
ngưỡng gia nhập và hấp thụ chi phí gia nhập cố định; và biên cường độ,
tức biến thiên FSTS trong nội bộ nhóm doanh nghiệp xuất khẩu. Lý thuyết
dự đoán phần thưởng năng suất có thể phát sinh ở cả hai biên, nhưng
trong các nền kinh tế chuyển đổi với hạ tầng xuất khẩu mỏng, biên tham
gia là biên nâng năng suất chủ đạo, còn biên cường độ thể hiện lợi suất
giảm dần hoặc bằng không một khi đã vượt ngưỡng tham gia. Cơ chế lý
thuyết đứng sau nhánh đi xuống là chi phí giao dịch thể chế (Williamson,
1985; Hennart, 2007): trong nền kinh tế chuyển đổi, khe hở thực thi hợp
đồng, thiếu hụt logistics và bất cân xứng thông tin làm tăng chi phí
biên quản lý mỗi thị trường nước ngoài bổ sung, khiến nhánh suy giảm sau
ngưỡng bám ở mức cường độ xuất khẩu thấp hơn so với các bối cảnh nơi hạ
tầng số và chất lượng thực thi hợp đồng đã hấp thụ phần lớn ma sát này.
Đây chính là cơ sở lý thuyết giải thích vì sao điểm uốn Việt Nam
(39--46\%) nằm thấp hơn đáng kể so với điểm uốn Singapore
(\textasciitilde88,6\%) ở biên thể chế mạnh nhất.

\hypertarget{bux1ed9-hux1ec7-sux1ed1-biuxean-ux111ux1ed9-kinh-tux1ebf-tux1eebng-mux1ee9c-fsts-vuxe0-ux111ux1ecbnh-lux1b0ux1ee3ng-bux1b0ux1edbc-nhux1ea3y-tham-gia}{%
\subsection{4.5.9 Bộ hệ số biên độ kinh tế từng mức FSTS và định lượng
bước nhảy tham
gia}\label{bux1ed9-hux1ec7-sux1ed1-biuxean-ux111ux1ed9-kinh-tux1ebf-tux1eebng-mux1ee9c-fsts-vuxe0-ux111ux1ecbnh-lux1b0ux1ee3ng-bux1b0ux1edbc-nhux1ea3y-tham-gia}}

Bộ hiệu ứng biên theo từng mức cường độ xuất khẩu trong đặc tả gộp M2
(\(N = 2.958\)) cho phép định lượng độ lớn kinh tế của đường cong thay
vì chỉ đọc dấu hệ số. Đạo hàm biên
\(\mathrm{d}y/\mathrm{d}\mathrm{FSTS} = \beta_1 + 2\beta_2(\mathrm{FSTS} - \overline{\mathrm{FSTS}})\)
với \(\beta_1 = 0{,}984\), \(\beta_2 = -1{,}909\) và
\(\overline{\mathrm{FSTS}} \approx 13{,}9\%\) cho độ dốc \(+1{,}13\) tại
mức 10\% FSTS, \(+0{,}75\) tại 20\%, \(+0{,}37\) tại 30\%, triệt tiêu
tại điểm uốn \(39{,}7\%\), rồi đảo âm \(-0{,}39\) tại 50\% và
\(-0{,}78\) tại 60\%. Tại hai biên dữ liệu, độ dốc là \(+1{,}52\) ở
\(\mathrm{FSTS} = 0\) và \(-2{,}30\) ở \(\mathrm{FSTS} = 100\%\), hai
dấu trái ngược mà kiểm định Lind--Mehlum xác nhận là không tương thích
với quan hệ đơn điệu (\(p < {,}001\)). Bốn điều kiện nhận diện chữ U
ngược của Haans và cộng sự (2016) đều thỏa trong mẫu gộp: C1 độ dốc
dương ở biên dưới (\(\beta_1 = +0{,}984\), \(p < {,}001\)); C2 độ cong
âm (\(\beta_2 = -1{,}909\), \(p < {,}001\)); C3 điểm uốn \(39{,}7\%\)
nằm trong miền dữ liệu \([0\%, 100\%]\); và C4 dấu độ dốc đảo chiều giữa
hai biên. Độ lớn của bước nhảy biên tham gia được thước đo đối sánh điểm
xu hướng xác nhận độc lập: tác động bình quân lên nhóm được xử lý của
trạng thái có công nghệ nước ngoài hoặc chứng nhận đạt
\(0{,}609\)--\(0{,}637\) điểm log (\(p < {,}001\)), tức khoảng 84--89\%
chênh lệch năng suất, lớn hơn nhiều so với tác động của riêng sở hữu
website (\(0{,}298\)--\(0{,}321\) điểm log).

\hypertarget{pin-chux1eb7t-ux111ux1ed9-vux1eefng-tuxe1ch-nguxe0nh-ux111ux1eb7c-tux1ea3-bux1eadc-ba-vuxe0-batterie-kiux1ec3m-ux111ux1ecbnh-nux1ed9i-sinh}{%
\subsection{4.5.10 Pin chặt độ vững: tách ngành, đặc tả bậc ba và
batterie kiểm định nội
sinh}\label{pin-chux1eb7t-ux111ux1ed9-vux1eefng-tuxe1ch-nguxe0nh-ux111ux1eb7c-tux1ea3-bux1eadc-ba-vuxe0-batterie-kiux1ec3m-ux111ux1ecbnh-nux1ed9i-sinh}}

Diễn giải chữ U ngược và sự phân biệt hai cấu trúc năng lực được nâng đỡ
bởi một loạt kiểm định độ vững có hệ thống. Thứ nhất, tách theo ngành
cho thấy chữ U ngược tồn tại ở cả hai mảng nhưng sắc nét hơn ở chế biến
chế tạo (\(\hat{\beta}_1 = +0{,}971\), \(p = {,}001\);
\(\hat{\beta}_2 = -1{,}883\), \(p < {,}001\); \(N = 1.854\)) so với phi
chế tạo (\(\hat{\beta}_1 = +1{,}615\), \(p = {,}064\);
\(\hat{\beta}_2 = -2{,}479\), \(p = {,}046\); \(N = 1.104\)); quan trọng
hơn, các kênh năng lực và số hóa vận hành chủ yếu trong chế biến chế
tạo, nơi \(\tilde{\mathrm{TCI}}\) trực tiếp \(= +0{,}223\)
(\(p < {,}001\)) và \(\tilde{\mathrm{DAI}}\) trực tiếp \(= +0{,}087\)
(\(p = {,}009\)), trong khi ở phi chế tạo \(\tilde{\mathrm{TCI}}\) chỉ ở
biên (\(+0{,}090\), \(p = {,}096\)) và \(\tilde{\mathrm{DAI}}\) không có
ý nghĩa (\(+0{,}068\), \(p = {,}133\)). Thứ hai, kiểm định dạng hàm bằng
đặc tả bậc ba bác bỏ đường cong chữ S: số hạng bậc ba không có ý nghĩa
(\(\hat{\beta} = -1{,}763\), \(p = {,}287\)), và cả AIC lẫn BIC đều ưu
tiên đặc tả bậc hai (AIC giảm \(0{,}99\), BIC giảm \(7{,}0\) khi quay về
đặc tả bậc hai). Thứ ba, loại doanh nghiệp siêu nhỏ (giữ
\(\ell_1 \geq 10\); \(N = 2.473\)) bảo toàn chữ U ngược
(\(\hat{\beta}_1 = +0{,}794\), \(p = {,}008\);
\(\hat{\beta}_2 = -1{,}647\), \(p < {,}001\)) và hiệu ứng TCI
(\(+0{,}188\), \(p < {,}001\)). Thứ tư, kiểm định giới hạn ổn định hệ số
của Oster (2019) với giả định
\(R^2_{\max} = 1{,}3 \times R^2_{\text{controlled}}\) cho thấy không hệ
số trọng tâm nào đổi dấu hay sụp về 0 dưới các độ lớn hợp lý của chọn
lựa không quan sát được, loại trừ khả năng thiên lệch biến bỏ sót giải
thích phát hiện chính. Bộ kiểm định này cộng hưởng với hệ số TCI 2SLS
dương mạnh (\(1{,}639\), lớn hơn ước lượng OLS, nhất quán với hiệu chỉnh
thiên lệch suy giảm do sai số đo) để xác lập rằng năng lực công nghệ là
cơ chế năng suất bền vững nhất trong bối cảnh thể chế Việt Nam, trong
khi hiện diện website mất sức phân biệt nhân quả.

\hypertarget{bux1ed1i-cux1ea3nh-vux129-muxf4-ux111ux1ed3ng-huxe0nh-vuxe0-ux111ux1ecbnh-vux1ecb-thiux1ebft-kux1ebf-ux111ux1a1n-quux1ed1c-gia-theo-thux1eddi-gian}{%
\subsection{4.5.11 Bối cảnh vĩ mô đồng hành và định vị thiết kế đơn quốc
gia theo thời
gian}\label{bux1ed1i-cux1ea3nh-vux129-muxf4-ux111ux1ed3ng-huxe0nh-vuxe0-ux111ux1ecbnh-vux1ecb-thiux1ebft-kux1ebf-ux111ux1a1n-quux1ed1c-gia-theo-thux1eddi-gian}}

Tính bền vững của dải điểm uốn 39--46\% qua 14 năm trở nên thuyết phục
hơn khi đối chiếu với bức tranh vĩ mô đồng hành ba làn sóng. Tại thời
điểm cuối cửa sổ quan sát, tăng trưởng GDP Việt Nam đạt \(7{,}09\%\) năm
2024, dòng vốn đầu tư trực tiếp nước ngoài đạt \(20{,}17\) tỷ đô-la Mỹ
(tương đương \(4{,}23\%\) GDP), và cấu trúc xuất khẩu dịch chuyển mạnh
sang máy móc và điện tử ở mức \(45{,}8\%\) tổng giá trị thương mại, soán
ngôi cụm dệt may từng thống trị (World Bank, 2025b). Sự nâng cấp cấu
trúc sản xuất này diễn ra song hành với ba đợt khảo sát, củng cố cách
đọc rằng dải điểm uốn ổn định không phải do nền kinh tế tĩnh mà do các
ràng buộc cấp doanh nghiệp (năng lực điều phối, thực thi hợp đồng, tiếp
cận tài trợ thương mại) thay đổi chậm hơn môi trường vận hành rộng. Về
định vị thiết kế, tài liệu phân tích tổng hợp gần đây cho thấy quan hệ
chữ U ngược giữa quốc tế hóa và hiệu quả suy yếu đáng kể, đôi khi biến
mất, dưới nhận dạng nghiêm ngặt hơn trong các mẫu gộp xuyên quốc gia
(Pisani và cộng sự, 2020), gợi ý rằng việc gộp các bối cảnh thể chế
không đồng nhất thổi phồng tính quy luật biểu kiến của dạng hàm. Thiết
kế dọc theo thời gian trong nội bộ một quốc gia, giữ cố định bối cảnh
thể chế trong khi thay đổi thời gian, do đó cung cấp kiểm định khả tín
hơn về việc liệu chữ U ngược là tạo tác của gộp xuyên quốc gia hay một
quan hệ thực bền vững khi môi trường thể chế tiến hóa, và bảng ba đợt
WBES Việt Nam (2009, 2015, 2023) đáp ứng đúng yêu cầu thiết kế này. Cách
đọc này định vị đóng góp của P3 vượt khỏi một báo cáo điểm uốn đơn lẻ:
nó chuyển trọng tâm tranh luận từ ``quốc tế hóa có giúp ích không?''
sang ``dưới điều kiện thể chế nào thì ngưỡng dịch chuyển?''.

\begin{center}\rule{0.5\linewidth}{0.5pt}\end{center}

\hypertarget{kux1ebft-quux1ea3-kiux1ec3m-ux111ux1ecbnh-touxe0n-mux1eabu-ux111a-bux1ed1i-cux1ea3nh-chuxe2u-uxe1-p7}{%
\section{4.6 Kết quả kiểm định toàn mẫu đa bối cảnh châu Á
(P7)}\label{kux1ebft-quux1ea3-kiux1ec3m-ux111ux1ecbnh-touxe0n-mux1eabu-ux111a-bux1ed1i-cux1ea3nh-chuxe2u-uxe1-p7}}

\hypertarget{mux1eabu-vuxe0-muxf4-huxecnh-chuxednh}{%
\subsection{4.6.1 Mẫu và mô hình
chính}\label{mux1eabu-vuxe0-muxf4-huxecnh-chuxednh}}

Phân tích P7 là kiểm định quy mô lớn nhất của luận án, kết hợp tất cả
các bối cảnh ICRV trong khu vực châu Á và Thái Bình Dương trên khung
\textbf{50 nền kinh tế} (gồm Nhật Bản, đợt khảo sát đầu tiên năm 2025),
103 cặp nền kinh tế và năm, với \textbf{\(N = 81.022\) quan sát} ở đặc
tả bậc hai cơ sở (M2) và \(N = 79.080\) khi bổ sung đầy đủ biến kiểm
soát (M5). Biến phụ thuộc là năng suất lao động chuẩn hóa z trong từng
cặp nền kinh tế và năm (trung hòa tiền tệ, winsorize 1/99 trong đợt);
mọi mô hình có hiệu ứng cố định nền kinh tế và năm, sai số chuẩn cụm
theo nền kinh tế. Mô hình M5, cấu hình đầy đủ với các biến kiểm soát sở
hữu nước ngoài, tuổi doanh nghiệp, TCI và DAI, ước lượng điểm uốn gộp:

\[\text{TP(M5)} = 43{,}6\% \text{ FSTS} \quad (p_{\text{LM}} < {,}001)\]

Điểm uốn dao động nhất quán trong dải \textbf{43,6--51,5\%} FSTS qua các
đặc tả M2 đến M10, với kiểm định Lind và Mehlum có ý nghĩa thống kê cao
ở mọi cấu hình. Tính nhất quán này bác bỏ giả thuyết rằng chữ U ngược
chỉ là sản phẩm phụ của mẫu nhỏ hay thiếu biến kiểm soát.

Bảng 4.6 trình bày tiến triển các mô hình lồng nhau, từ đặc tả bậc hai
cơ sở đến các mô hình điều tiết, cho thấy điểm uốn ổn định và quan hệ
chữ U ngược được xác nhận qua mọi cấu hình.

\textbf{Bảng 4.6.} \emph{Tiến triển mô hình của P7 trên khung 50 nền
kinh tế: điểm uốn và kiểm định Lind và Mehlum.}

\begin{longtable}[]{@{}
  >{\raggedright\arraybackslash}p{(\columnwidth - 12\tabcolsep) * \real{0.1429}}
  >{\raggedleft\arraybackslash}p{(\columnwidth - 12\tabcolsep) * \real{0.1429}}
  >{\raggedright\arraybackslash}p{(\columnwidth - 12\tabcolsep) * \real{0.1429}}
  >{\raggedright\arraybackslash}p{(\columnwidth - 12\tabcolsep) * \real{0.1429}}
  >{\raggedleft\arraybackslash}p{(\columnwidth - 12\tabcolsep) * \real{0.1429}}
  >{\raggedleft\arraybackslash}p{(\columnwidth - 12\tabcolsep) * \real{0.1429}}
  >{\raggedright\arraybackslash}p{(\columnwidth - 12\tabcolsep) * \real{0.1429}}@{}}
\toprule\noalign{}
\begin{minipage}[b]{\linewidth}\raggedright
Mô hình
\end{minipage} & \begin{minipage}[b]{\linewidth}\raggedleft
N
\end{minipage} & \begin{minipage}[b]{\linewidth}\raggedright
β₁ (p)
\end{minipage} & \begin{minipage}[b]{\linewidth}\raggedright
β₂ (p)
\end{minipage} & \begin{minipage}[b]{\linewidth}\raggedleft
Điểm uốn
\end{minipage} & \begin{minipage}[b]{\linewidth}\raggedleft
\(p_{\text{LM}}\)
\end{minipage} & \begin{minipage}[b]{\linewidth}\raggedright
Đặc tả bổ sung
\end{minipage} \\
\midrule\noalign{}
\endhead
\bottomrule\noalign{}
\endlastfoot
M2 & 81.022 & +1,189 (\textless,001) & −1,398 (\textless,001) & 51,5\% &
\textless,001 & FSTS + FSTS², FE hai chiều \\
\textbf{M5} & \textbf{79.080} & \textbf{+0,500 (\textless,001)} &
\textbf{−0,721 (\textless,001)} & \textbf{43,6\%} &
\textbf{\textless,001} & \textbf{+ kiểm soát (mô hình chính)} \\
M7 & 79.220 & +0,627 (\textless,001) & −0,836 (\textless,001) & 46,4\% &
\textless,001 & + TCI và tương tác TCI \\
M8 & 79.080 & +0,704 (\textless,001) & −0,924 (\textless,001) & 47,1\% &
\textless,001 & + DAI và tương tác DAI \\
M10 & 79.080 & +0,615 (\textless,001) & −0,833 (\textless,001) & 45,9\%
& \textless,001 & + tương tác chế độ thể chế yếu \\
\end{longtable}

\emph{Nguồn: ước lượng của tác giả trên khung 50 nền (tái lập:
\texttt{scripts/p7\_run\_50econ.py}). Biến kiểm soát: sở hữu nước ngoài
≥10\%, ln(tuổi doanh nghiệp), TCI chuẩn hóa, DAI. Điểm uốn xác nhận bằng
kiểm định Lind và Mehlum (2010), một phía. Sai số chuẩn cụm theo nền
kinh tế.}

\textbf{Lưu ý xác thực.} Bộ kết quả này được ước lượng bằng công cụ hồi
quy hiệu ứng cố định tương đương reghdfe (pyfixest) với spec được tài
liệu hóa đầy đủ và tái lập được; bộ do-file Stata tương ứng
(\texttt{scripts/stata/}) cho phép xác thực lại từng hệ số trên máy có
Stata trước khi nộp tạp chí.

\hypertarget{ux111iux1ec3m-uux1ed1n-theo-chux1ebf-ux111ux1ed9-thux1ec3-chux1ebf-tuxe1i-ux111ux1ecbnh-huxecnh-h5}{%
\subsection{4.6.2 Điểm uốn theo chế độ thể chế: tái định hình
H5}\label{ux111iux1ec3m-uux1ed1n-theo-chux1ebf-ux111ux1ed9-thux1ec3-chux1ebf-tuxe1i-ux111ux1ecbnh-huxecnh-h5}}

Ước lượng dạng M5 trong từng nhóm ICRV cho thấy một bức tranh điều tiết
thể chế \textbf{tinh tế hơn} giả thuyết ban đầu:

\textbf{Bảng 4.7.} \emph{Điểm uốn theo nhóm ICRV (khung 50 nền, dạng M5
trong từng nhóm).}

\begin{longtable}[]{@{}
  >{\raggedright\arraybackslash}p{(\columnwidth - 12\tabcolsep) * \real{0.1429}}
  >{\raggedleft\arraybackslash}p{(\columnwidth - 12\tabcolsep) * \real{0.1429}}
  >{\raggedright\arraybackslash}p{(\columnwidth - 12\tabcolsep) * \real{0.1429}}
  >{\raggedright\arraybackslash}p{(\columnwidth - 12\tabcolsep) * \real{0.1429}}
  >{\raggedleft\arraybackslash}p{(\columnwidth - 12\tabcolsep) * \real{0.1429}}
  >{\raggedleft\arraybackslash}p{(\columnwidth - 12\tabcolsep) * \real{0.1429}}
  >{\raggedright\arraybackslash}p{(\columnwidth - 12\tabcolsep) * \real{0.1429}}@{}}
\toprule\noalign{}
\begin{minipage}[b]{\linewidth}\raggedright
Nhóm ICRV
\end{minipage} & \begin{minipage}[b]{\linewidth}\raggedleft
N
\end{minipage} & \begin{minipage}[b]{\linewidth}\raggedright
β₁ (p)
\end{minipage} & \begin{minipage}[b]{\linewidth}\raggedright
β₂ (p)
\end{minipage} & \begin{minipage}[b]{\linewidth}\raggedleft
Điểm uốn
\end{minipage} & \begin{minipage}[b]{\linewidth}\raggedleft
\(p_{\text{LM}}\)
\end{minipage} & \begin{minipage}[b]{\linewidth}\raggedright
Kết luận
\end{minipage} \\
\midrule\noalign{}
\endhead
\bottomrule\noalign{}
\endlastfoot
I, Adv. đổi mới (gồm Nhật Bản) & 5.581 & +0,698 (,17) & −0,504 (,31) &
{[}79,1\%{]} & ,30 & gần tuyến tính dương; điểm uốn ngoài miền, không
định vị \\
II, Adv. tài nguyên & 2.075 & +0,854 (,045) & −0,730 (,082) & 62,0\% &
,081 & chữ U ngược ở biên ý nghĩa \\
III, Trung bình cao & 12.055 & +0,168 (,42) & −0,189 (,48) &
{[}55,0\%{]} & ,26 & không xác định ở cấp nhóm \\
\textbf{IV, Chuyển đổi TB-thấp} & \textbf{42.094} & \textbf{+0,709
(\textless,001)} & \textbf{−1,012 (\textless,001)} & \textbf{43,0\%} &
\textbf{\textless,001} & \textbf{chữ U ngược rõ nét nhất} \\
V, Đang nổi & 15.457 & +0,218 (,44) & −0,455 (,21) & {[}34,8\%{]} & ,18
& quan hệ tan rã \\
VI, SIDS & 1.818 & −0,728 (,010) & +0,870 (,013) & n/a & n/a & không có
chữ U ngược \\
\end{longtable}

\emph{Ghi chú: giá trị trong ngoặc vuông là điểm uốn đại số khi đường
cong không được xác nhận thống kê, chỉ mang tính tham khảo. Nguồn: như
Bảng 4.6.}

Kết quả này \textbf{tái định hình H5} theo hướng giàu thông tin hơn: thể
chế không chỉ dịch chuyển vị trí điểm uốn mà quyết định \textbf{chính sự
tồn tại của dạng hàm chữ U ngược}. Ba vùng được nhận diện rõ. Ở
\textbf{vùng giữa của gradient thể chế} (Nhóm IV), chữ U ngược định hình
rõ nét và tin cậy nhất, với điểm uốn 43,0\% trùng khớp với các ước lượng
đơn quốc gia độc lập trong cùng nhóm (P3 Việt Nam 39--46\%; P9 Ấn Độ
2022 ở mức 40,7\%). Ở \textbf{biên trên của gradient} (Nhóm I, nay gồm
cả Nhật Bản), đường cong duỗi thẳng thành quan hệ gần tuyến tính dương
với điểm uốn nằm ngoài miền dữ liệu và không định vị được, trùng khớp
với phát hiện đơn quốc gia tại Singapore (P4: gần tuyến tính, điểm uốn
ước lượng \textasciitilde88,6\% không định vị chắc chắn). Ở \textbf{biên
dưới của gradient} (Nhóm V và VI), quan hệ mất cấu trúc: các hệ số không
còn ý nghĩa thống kê ở Nhóm V và Nhóm VI không có dạng chữ U ngược, nhất
quán với việc tương quan FSTS và năng suất mất ý nghĩa ở SIDS trong phân
tích mô tả (Chuyên đề 1, Bảng 2.3.8.1) và với phát hiện FIP của P8.

Tương tác chế độ thể chế yếu trong mô hình gộp củng cố cách đọc này: hệ
số FSTS×(chế độ yếu) = −0,523 (\(p = {,}087\)) cho thấy độ dốc nhánh đi
lên của đường cong yếu hơn rõ rệt ở các chế độ thể chế yếu, nghĩa là
cùng một mức tăng cường độ xuất khẩu mang lại phần thưởng năng suất nhỏ
hơn khi đệm thể chế mỏng. Kiểm định điều tiết tổng thể về ICRV của phân
tích tổng hợp (\(Q_M = 17{,}35\), \(p = {,}002\), P6) cung cấp nguồn xác
nhận độc lập thứ hai rằng chế độ thể chế là biến điều tiết có ý nghĩa
của quan hệ quốc tế hóa và hiệu quả.

\begin{figure}
\centering
\includegraphics{figures/fig_4_x_ip_curves_icrv_gradient.png}
\caption{Hình 4.8. Họ đường cong quốc tế hóa và hiệu quả phụ thuộc chế
độ thể chế ICRV}
\end{figure}

\textbf{Hình 4.8.} \emph{Họ đường cong quốc tế hóa và hiệu quả phụ thuộc
chế độ thể chế ICRV.} Chữ U ngược định hình rõ nét nhất ở vùng giữa
gradient thể chế (Nhóm IV, điểm uốn 43,0\%); ở biên trên (Nhóm I) đường
cong duỗi thẳng gần tuyến tính dương; ở biên dưới (Nhóm V và VI) quan hệ
mất cấu trúc, riêng SIDS chuyển sang quan hệ âm theo mức (FIP, H1b, theo
spec mức của P8). \emph{Nguồn: minh họa của tác giả từ kết quả ước lượng
50 nền và P4, P8.}

\hypertarget{ux111iux1ec1u-tiux1ebft-bux1edfi-tci-vuxe0-dai-trong-touxe0n-mux1eabu}{%
\subsection{4.6.3 Điều tiết bởi TCI và DAI trong toàn
mẫu}\label{ux111iux1ec1u-tiux1ebft-bux1edfi-tci-vuxe0-dai-trong-touxe0n-mux1eabu}}

\begin{itemize}
\tightlist
\item
  \textbf{TCI}: hiệu ứng chính dương và có ý nghĩa cao,
  \(\hat{\beta}(\text{TCI}) = +0{,}141\) (\(p < {,}001\)): một độ lệch
  chuẩn TCI nâng năng suất chuẩn hóa thêm 0,141 độ lệch chuẩn. Các tương
  tác FSTS×TCI và FSTS²×TCI không có ý nghĩa thống kê, xác nhận TCI là
  \textbf{năng lực nâng mặt bằng} chứ không uốn hình dạng đường cong,
  nhất quán với kết luận từ các nghiên cứu đơn quốc gia.
\item
  \textbf{DAI}: hiệu ứng chính dương và có ý nghĩa cao,
  \(\hat{\beta}(\text{DAI}) = +0{,}201\) (\(p < {,}001\)): doanh nghiệp
  có hiện diện số cơ bản đạt năng suất chuẩn hóa cao hơn 0,2 độ lệch
  chuẩn. Các tương tác với đường cong (FSTS×DAI = −0,271,
  \(p = {,}149\); FSTS²×DAI = +0,252, \(p = {,}367\)) mang \textbf{cùng
  chiều nén đường cong} như dự đoán của khung CDCM nhưng không đạt ý
  nghĩa thống kê trên toàn mẫu 50 nền; bằng chứng điều tiết DAI rõ nét
  nhất vẫn là ở cấp bối cảnh đơn quốc gia (P4 Singapore: FSTS²×DAI =
  +3,119, \(p = {,}005\)). H3 do đó được hỗ trợ ở hiệu ứng mức và ở bối
  cảnh đặc thù, chưa được xác nhận như hiệu ứng uốn đường cong phổ quát.
  Khoảng cách độ lớn giữa hai ước lượng (3,119 ở Singapore so với 0,252
  ở mẫu gộp) không phải mâu thuẫn mà phản ánh bản chất trung bình hóa
  của mẫu gộp: phần lớn 50 nền chỉ đo được DAI ở Tầng 1 (hiện diện web),
  nơi công cụ số chưa đủ chiều sâu để uốn đường cong, trong khi hiệu ứng
  uốn mạnh ở Singapore gắn với năng lực số Tầng 2 hiếm có trong mẫu; gộp
  hai loại bối cảnh này lại làm pha loãng hiệu ứng uốn về một trung bình
  không có ý nghĩa, đúng như logic phụ thuộc bối cảnh của khung CDCM dự
  báo chứ không phủ định cơ chế lá chắn số.
\end{itemize}

\hypertarget{diux1ec5n-tiux1ebfn-ux111iux1ec3m-uux1ed1n-vuxe0-bux1ed1n-ux111iux1ec1u-kiux1ec7n-nhux1eadn-diux1ec7n-chux1eef-u-ngux1b0ux1ee3c}{%
\subsection{4.6.4 Diễn tiến điểm uốn và bốn điều kiện nhận diện chữ U
ngược}\label{diux1ec5n-tiux1ebfn-ux111iux1ec3m-uux1ed1n-vuxe0-bux1ed1n-ux111iux1ec1u-kiux1ec7n-nhux1eadn-diux1ec7n-chux1eef-u-ngux1b0ux1ee3c}}

Khi đọc theo chuỗi mô hình lồng nhau, điểm uốn dịch chuyển có quy luật
chứ không ngẫu nhiên. Đặc tả tuyến tính cơ sở cho hệ số FSTS dương và có
ý nghĩa cao (\(\hat{\beta} = +0{,}219\), \(p < {,}001\)), xác nhận rằng
ở mức trung bình toàn mẫu quan hệ quốc tế hóa và năng suất là dương. Khi
bổ sung số hạng bậc hai (đặc tả M2), đường cong định hình chữ U ngược
với điểm uốn ở mức 51,5\% FSTS. Việc đưa vào lần lượt các biến kiểm soát
nhân khẩu doanh nghiệp rồi khối năng lực (TCI và DAI) kéo điểm uốn xuống
43,8\% rồi 43,6\%, chính là điểm uốn của mô hình chính (M5). Sự dịch
chuyển xuống khoảng bảy điểm phần trăm này có ý nghĩa thực chất: một
phần độ cong biểu kiến ở đặc tả thô phản ánh chênh lệch năng lực giữa
các doanh nghiệp xuất khẩu nhiều và ít; sau khi kiểm soát năng lực,
ngưỡng tối ưu thực của cường độ xuất khẩu thấp hơn và nằm sát các ước
lượng đơn quốc gia trong vùng chuyển đổi (Việt Nam 39--46\%, Ấn Độ 2022
ở mức 40,7\%).

Chữ U ngược được xác nhận chính thức theo bốn điều kiện của Haans và
cộng sự (2016) chứ không chỉ dựa vào dấu của hệ số bậc hai. Thứ nhất, hệ
số nhánh đi lên dương (\(\beta_1 > 0\)). Thứ hai, hệ số nhánh đi xuống
âm (\(\beta_2 < 0\)). Thứ ba, điểm uốn nằm hẳn trong miền dữ liệu quan
sát của FSTS. Thứ tư, độ dốc ở hai biên thỏa mãn dự đoán hướng, dương ở
biên dưới và âm ở biên trên. Kiểm định Lind và Mehlum (2010), vốn kiểm
tra đồng thời các điều kiện này, cho giá trị \(p_{\text{LM}} < {,}001\)
ở mô hình chính, nghĩa là chữ U ngược được nhận diện chắc chắn chứ không
phải tạo tác của một quan hệ đơn điệu bị ép vào dạng bậc hai.

\hypertarget{tuxednh-ux1ed5n-ux111ux1ecbnh-cux1ee7a-ux111iux1ec3m-uux1ed1n-trux1b0ux1edbc-ux111ux1eb7c-trux1b0ng-quux1ea3n-luxfd-vuxe0-tux1b0ux1a1ng-tuxe1c-thux1ec3-chux1ebf}{%
\subsection{4.6.5 Tính ổn định của điểm uốn trước đặc trưng quản lý và
tương tác thể
chế}\label{tuxednh-ux1ed5n-ux111ux1ecbnh-cux1ee7a-ux111iux1ec3m-uux1ed1n-trux1b0ux1edbc-ux111ux1eb7c-trux1b0ng-quux1ea3n-luxfd-vuxe0-tux1b0ux1a1ng-tuxe1c-thux1ec3-chux1ebf}}

Điểm uốn gộp tỏ ra trơ trước việc bổ sung các đặc trưng quản lý. Trong
mô hình có biến quản lý (\(N = 75.029\)), kinh nghiệm ngành của người
quản lý không có ý nghĩa thống kê (\(\hat{\beta} = +0{,}002\),
\(p > {,}10\)), còn hệ số người quản lý cấp cao là nữ mang dấu âm và có
ý nghĩa cao (\(\hat{\beta} = -0{,}104\), \(p < {,}001\)); điểm uốn gần
như không đổi ở mức 44,0\%. Dấu âm của hệ số quản lý nữ nhất quán với
phần còn lại của luận án (nghiên cứu đơn quốc gia Nhật Bản P10 và phân
tích gộp đều ước lượng hệ số âm) và được đọc thận trọng như một quy luật
mô tả phản ánh chọn lọc cấu thành, theo đó các cơ sở do nữ lãnh đạo tập
trung nhiều hơn vào các phân khúc nhỏ, biên lợi nhuận thấp ở những nền
kinh tế mà lãnh đạo nữ cấp cao còn hiếm, chứ không phải bằng chứng nhân
quả về hiệu quả của lãnh đạo nữ. Vì các tương tác giữa FSTS và đặc trưng
quản lý không có ý nghĩa, đặc trưng quản lý dịch chuyển mặt bằng năng
suất chứ không giúp doanh nghiệp đi qua đánh đổi quốc tế hóa hiệu quả
hơn.

Tương tự, chế độ thể chế ICRV không biểu hiện như một dịch chuyển tuyến
tính đơn của đường cong gộp. Vì biến ICRV bất biến theo thời gian trong
mỗi nền kinh tế, hiệu ứng chính của nó bị hấp thụ bởi hiệu ứng cố định
nền kinh tế; chỉ các số hạng tương tác được nhận diện. Trong mô hình
tương tác thể chế gộp, cả FSTS×ICRV (\(\hat{\beta} = -0{,}086\)) lẫn
FSTS²×ICRV (\(\hat{\beta} = -0{,}013\)) đều không có ý nghĩa, khiến điểm
uốn gộp không định vị được. Đây không phải bằng chứng phủ định H5 mà là
chỉ dấu phương pháp luận: cấu trúc thể chế của quan hệ chỉ hiện ra khi
ước lượng điểm uốn trong từng nhóm ICRV (Bảng 4.7), nơi nó biểu hiện
thành ba vùng rời rạc chứ không phải một gradient trơn. Nói cách khác,
điều tiết thể chế là điều tiết về sự tồn tại của dạng hàm, không phải
điều tiết về vị trí của một điểm uốn phổ quát.

\hypertarget{ux111ux1eb7c-trux1b0ng-mux1eabu-ux111a-cux1ed9ng-tuyux1ebfn-vuxe0-uxfd-nghux129a-tux1ed5ng-hux1ee3p-ba-vuxf9ng}{%
\subsection{4.6.6 Đặc trưng mẫu, đa cộng tuyến và ý nghĩa tổng hợp ba
vùng}\label{ux111ux1eb7c-trux1b0ng-mux1eabu-ux111a-cux1ed9ng-tuyux1ebfn-vuxe0-uxfd-nghux129a-tux1ed5ng-hux1ee3p-ba-vuxf9ng}}

Cường độ xuất khẩu trung bình toàn mẫu là 9,0\% FSTS, phản ánh một quần
thể doanh nghiệp thiên về thị trường nội địa với một đuôi xuất khẩu
mạnh, đặc điểm điển hình của dữ liệu doanh nghiệp WBES ở các nền kinh tế
đang phát triển và mới nổi. Để giảm đa cộng tuyến giữa số hạng tuyến
tính và bậc hai, FSTS được căn giữa quanh trung bình mẫu trước khi lập
số hạng bình phương và các số hạng tương tác; nhờ vậy hệ số nhánh đi lên
và nhánh đi xuống được ước lượng ổn định qua toàn bộ chuỗi mô hình. Mọi
mô hình mang hiệu ứng cố định nền kinh tế và năm cùng sai số chuẩn cụm
theo nền kinh tế (CRV1), kiểm soát đồng thời chênh lệch mặt bằng giá và
năng suất giữa các nền kinh tế và xu hướng theo thời gian.

Tổng hợp lại, ba mức bằng chứng, gồm diễn tiến điểm uốn ổn định trong
dải 43,6--51,5\%, tính trơ của điểm uốn trước năng lực và đặc trưng quản
lý, và cấu trúc ba vùng theo chế độ thể chế, hội tụ về một cách đọc
thống nhất: năng lực công nghệ và số hóa cùng đặc trưng quản lý nâng
\emph{mặt bằng} năng suất nhưng không uốn lại hình dạng đường cong gộp;
cái uốn lại đường cong là chế độ thể chế, và nó tác động không như một
gradient trơn mà như ba vùng rời rạc. Ở giữa gradient thể chế (vùng
chuyển đổi thu nhập trung bình thấp), chữ U ngược định hình sắc nét
quanh điểm uốn 43\%; ở biên trên (thể chế mạnh nhất, gồm Nhật Bản),
nhánh đi xuống bị đẩy ra ngoài miền quan sát và quan hệ gần như tuyến
tính dương; ở biên dưới (đang nổi và SIDS Thái Bình Dương), chữ U ngược
tan rã và đảo dấu ở nhóm SIDS thành hình phạt quốc tế hóa cưỡng bức.
Cách đọc ba vùng này dung hòa các kết quả đơn quốc gia tưởng như mâu
thuẫn của luận án trong một khung ước lượng hài hòa trên 50 nền kinh tế.

\hypertarget{hux1ed9i-tux1ee5-chuxe9o-nghiuxean-cux1ee9u-cux1ee7a-ux111iux1ec3m-uux1ed1n-vuxf9ng-giux1eefa-vuxe0-logic-chux1ecdn-lux1ecdc-melitz-ux1edf-biuxean-truxean}{%
\subsection{4.6.7 Hội tụ chéo nghiên cứu của điểm uốn vùng giữa và logic
chọn lọc Melitz ở biên
trên}\label{hux1ed9i-tux1ee5-chuxe9o-nghiuxean-cux1ee9u-cux1ee7a-ux111iux1ec3m-uux1ed1n-vuxf9ng-giux1eefa-vuxe0-logic-chux1ecdn-lux1ecdc-melitz-ux1edf-biuxean-truxean}}

Sức thuyết phục của cấu trúc ba vùng nằm ở chỗ ước lượng gộp 50 nền và
các ước lượng đơn quốc gia độc lập hội tụ chặt chẽ ở từng vùng dù được
ước lượng trên các mẫu, công cụ khảo sát và đặc tả khác nhau. Ở vùng
giữa, điểm uốn cấp nhóm của Nhóm IV (43,0\%) trùng khớp đáng kể với điểm
uốn đơn quốc gia của Việt Nam, vốn được kiểm định Lind--Mehlum xác nhận
trong cả ba đợt (2009 \(p = {,}006\), 2015 \(p = {,}009\), 2023
\(p = {,}013\)) và mẫu gộp (\(p < {,}001\)) với dải điểm uốn
\(39{,}3\)--\(46{,}2\%\), và với điểm uốn Ấn Độ 2022 ở mức \(40{,}7\%\).
Sự trùng khớp ba ước lượng độc lập quanh mốc 40--43\% ở cùng chế độ thể
chế ICRV IV cung cấp bằng chứng hội tụ mạnh rằng điểm uốn vùng giữa là
một quy luật cấu trúc chứ không phải tạo tác của một mẫu đơn lẻ.

Ở biên trên của gradient, hình thái gần tuyến tính dương của Nhóm I được
làm rõ về cơ chế nhờ bằng chứng đơn quốc gia. Tại Singapore, nhánh suy
giảm bên phải không được nhận diện vì khoảng \(82\%\) doanh nghiệp không
xuất khẩu và chỉ \(3\%\) vượt ngưỡng \(50\%\) FSTS, khiến điểm uốn
\(\approx 88{,}6\%\) rơi vào vùng dữ liệu thưa với khoảng tin cậy
bootstrap rộng \([53\%; 253\%]\). Cơ chế đứng sau biên trên là chọn lọc
theo năng suất kiểu Melitz (2003): ở chế độ thể chế mạnh nhất, câu hỏi
ràng buộc không phải doanh nghiệp xuất khẩu \emph{bao nhiêu} mà là
\emph{có} xuất khẩu hay không, và các doanh nghiệp vượt ngưỡng tham gia
là những doanh nghiệp được trang bị năng lực, đúng như logic tự chọn lọc
mà mô hình doanh nghiệp dị biệt dự đoán cho biên thâm dụng tri thức
(Helpman và cộng sự, 2004). Ba vùng do đó không chỉ mô tả ba hình thái
đường cong mà ánh xạ ba cơ chế kinh tế khác nhau: chi phí điều phối bám
ngưỡng ở vùng giữa, chọn lọc theo năng suất ở biên trên, và hình phạt
quốc tế hóa cưỡng bức ở biên dưới.

\begin{center}\rule{0.5\linewidth}{0.5pt}\end{center}

\hypertarget{kux1ebft-quux1ea3-trux1b0ux1eddng-hux1ee3p-biuxean-sids-cuxe1c-nux1ec1n-kinh-tux1ebf-ux111ux1ea3o-nhux1ecf-ux111ang-phuxe1t-triux1ec3n-p8}{%
\section{4.7 Kết quả trường hợp biên SIDS: Các nền kinh tế đảo nhỏ đang
phát triển
(P8)}\label{kux1ebft-quux1ea3-trux1b0ux1eddng-hux1ee3p-biuxean-sids-cuxe1c-nux1ec1n-kinh-tux1ebf-ux111ux1ea3o-nhux1ecf-ux111ang-phuxe1t-triux1ec3n-p8}}

\hypertarget{ux111ux1eb7c-ux111iux1ec3m-mux1eabu-cuxe1c-nux1ec1n-kinh-tux1ebf-ux111ux1ea3o-nhux1ecf}{%
\subsection{4.7.1 Đặc điểm mẫu các nền kinh tế đảo
nhỏ}\label{ux111ux1eb7c-ux111iux1ec3m-mux1eabu-cuxe1c-nux1ec1n-kinh-tux1ebf-ux111ux1ea3o-nhux1ecf}}

Phân tích P8 tập trung vào bảy nền kinh tế đảo nhỏ đang phát triển Thái
Bình Dương thuộc Nhóm VI, gồm Fiji, Kiribati, Papua New Guinea, Samoa,
Quần đảo Solomon, Tonga và Vanuatu, với dữ liệu WBES các đợt khảo sát
giai đoạn 2009--2025. Mẫu phân tích, gồm các quan sát có đủ giá trị về
năng suất lao động, cường độ xuất khẩu và tuổi doanh nghiệp, có tổng
\(N = 959\) doanh nghiệp, trong đó chỉ 132 doanh nghiệp (13,8\% mẫu) có
xuất khẩu với cường độ dương. Cường độ xuất khẩu trung bình rất thấp, ở
mức 0,048 (tức 4,8\%), phản ánh đúng đặc trưng của nhóm: phần lớn doanh
nghiệp định hướng nội địa, còn một thiểu số xuất khẩu thì làm vậy do bắt
buộc cấu trúc chứ không phải lựa chọn chiến lược.

Ba đặc điểm cấu trúc của các nền này định khung cho toàn bộ phân tích.
Thứ nhất, thị trường nội địa siêu nhỏ, với quy mô dân số từ khoảng
94.000 (Micronesia, ngoài mẫu) và Tonga 104.000, Kiribati 121.000 đến
tối đa khoảng 10 triệu (Papua New Guinea, nền lớn nhất trong mẫu), khiến
nhiều doanh nghiệp không thể đạt quy mô hiệu quả tối thiểu nếu chỉ phục
vụ thị trường trong nước. Thứ hai, sự cô lập địa lý cực độ khuếch đại
chi phí thương mại thêm khoảng 30 đến 210\% so với mặt bằng. Thứ ba, các
khoảng trống thể chế toàn diện hạn chế khả năng doanh nghiệp chiếm giữ
lợi ích từ hoạt động quốc tế. Trong bối cảnh này, cường độ xuất khẩu
không còn là chỉ báo của cam kết tự nguyện với thị trường nước ngoài mà
trở thành đại lượng đại diện cho \emph{mức độ bất cập của thị trường nội
địa}: doanh nghiệp có cường độ xuất khẩu cao hơn không hẳn là doanh
nghiệp có năng lực quốc tế vượt trội, mà có thể chỉ là doanh nghiệp chịu
ràng buộc thị trường nội địa gay gắt hơn. Chính sự đảo nghĩa của thước
đo này là nền tảng cho dự đoán hình phạt quốc tế hóa cưỡng bức.

\hypertarget{hiux1ec7u-ux1ee9ng-uxe2m-ux111ux1a1n-ux111iux1ec7u-huxecnh-phux1ea1t-quux1ed1c-tux1ebf-huxf3a-cux1b0ux1ee1ng-bux1ee9c-fip}{%
\subsection{4.7.2 Hiệu ứng âm đơn điệu: Hình phạt quốc tế hóa cưỡng bức
(FIP)}\label{hiux1ec7u-ux1ee9ng-uxe2m-ux111ux1a1n-ux111iux1ec7u-huxecnh-phux1ea1t-quux1ed1c-tux1ebf-huxf3a-cux1b0ux1ee1ng-bux1ee9c-fip}}

Đây là phát hiện nổi bật nhất và khác biệt nhất của luận án so với toàn
bộ kết quả từ các bối cảnh ICRV khác. Mô hình M1 với country hiệu ứng cố
định và year hiệu ứng cố định cho:

\[\hat{\beta}(\text{FSTS}_c) = -1{,}339, \quad SE = 0{,}386, \quad p < {,}001^{***}\]

Hệ số hồi quy âm và có ý nghĩa thống kê, nghĩa là tại các nền kinh tế
SIDS, quốc tế hóa càng cao thì hiệu quả hoạt động doanh nghiệp càng
\textbf{thấp hơn}, quan hệ âm đơn điệu. Trong mẫu chỉ gồm doanh nghiệp
xuất khẩu (exporters-only, N = 26, rất nhỏ), hệ số vẫn âm nhưng không
đạt ý nghĩa thống kê (thiếu công suất):

\[\hat{\beta}(\text{FSTS}_c) = -1{,}176, \quad p = {,}130\ (\text{NS})\]

Mô hình M2 bổ sung FSTS\(^2\) không cho kết quả có ý nghĩa thống kê, xác
nhận quan hệ là \textbf{đơn điệu} (monotone decline), không có điểm uốn
trong phạm vi dữ liệu.

\hypertarget{giux1ea3i-thuxedch-cux1a1-chux1ebf-huxecnh-phux1ea1t-quux1ed1c-tux1ebf-huxf3a-cux1b0ux1ee1ng-bux1ee9c}{%
\subsection{4.7.3 Giải thích cơ chế hình phạt quốc tế hóa cưỡng
bức}\label{giux1ea3i-thuxedch-cux1a1-chux1ebf-huxecnh-phux1ea1t-quux1ed1c-tux1ebf-huxf3a-cux1b0ux1ee1ng-bux1ee9c}}

Luận án đặt tên cho hiện tượng này là \textbf{Forced
Internationalization Penalty (FIP)}, gánh nặng quốc tế hóa bắt buộc, một
cấu trúc lý thuyết mới chưa có trong tài liệu IB trước đây. Cơ chế giải
thích gồm ba yếu tố cấu trúc đặc thù của SIDS: (1) thị trường nội địa
quá nhỏ, buộc doanh nghiệp SIDS xuất khẩu không phải vì có lợi thế cạnh
tranh mà vì nhu cầu nội địa không đủ để duy trì hoạt động; (2) chi phí
hậu cần và tiếp cận thị trường cực cao do vị trí địa lý xa xôi và phân
tán; và (3) thiếu hỗ trợ thể chế và năng lực xuất khẩu.

Kết quả FIP là đóng góp lý thuyết gốc của luận án, lần đầu tiên được ghi
nhận và đặt tên trong tài liệu quốc tế hóa--hiệu quả doanh nghiệp cho
khu vực Thái Bình Dương.

\hypertarget{ux111ux1ecbnh-vux1ecb-fip-trong-bux1eb1ng-chux1ee9ng-vux129-muxf4-vux1ec1-tux1ed5n-thux1b0ux1a1ng-cux1ee7a-nux1ec1n-kinh-tux1ebf-nhux1ecf}{%
\subsection{4.7.4 Định vị FIP trong bằng chứng vĩ mô về tổn thương của
nền kinh tế
nhỏ}\label{ux111ux1ecbnh-vux1ecb-fip-trong-bux1eb1ng-chux1ee9ng-vux129-muxf4-vux1ec1-tux1ed5n-thux1b0ux1a1ng-cux1ee7a-nux1ec1n-kinh-tux1ebf-nhux1ecf}}

Cơ chế FIP ở cấp doanh nghiệp ăn khớp với đặc trưng tổn thương cấu trúc
của các nền kinh tế nhỏ và thị trường biên được ghi nhận ở cấp vĩ mô.
Bằng chứng của Ngân hàng Thế giới cho thấy bản chất phơi nhiễm bên ngoài
của các thị trường biên là bất lợi so với thị trường mới nổi: dòng vốn
danh mục biến động hơn, rổ xuất khẩu tập trung hơn, và nợ chủ yếu được
định danh bằng ngoại tệ với mức dự trữ thấp, tất cả làm tăng rủi ro bảng
cân đối (World Bank, 2026c). Đối với các nền kinh tế quy mô siêu nhỏ như
SIDS Thái Bình Dương, những bất lợi này được khuếch đại bởi chi phí
logistics cực cao và sự tập trung xuất khẩu cao độ, khiến việc gia tăng
cường độ xuất khẩu không đi kèm phần thưởng năng suất mà thậm chí trở
thành gánh nặng. Như vậy, FIP không phải một ngoại lệ thống kê mà là
biểu hiện cấp doanh nghiệp của một quy luật cấu trúc rộng hơn: ở những
bối cảnh thị trường nội địa quá nhỏ và đệm thể chế quá mỏng, quốc tế hóa
mang tính bắt buộc chứ không phải lựa chọn chiến lược tối ưu hóa năng
suất. Phát hiện này bổ sung một viên gạch vi mô còn thiếu vào bức tranh
vĩ mô của Ngân hàng Thế giới về các nền kinh tế nhỏ và dễ tổn thương.

\hypertarget{cux1ea5u-truxfac-mux1eabu-bux1ea3y-nux1ec1n-kinh-tux1ebf-vuxe0-phux1b0ux1a1ng-sai-giux1eefa-quux1ed1c-gia}{%
\subsection{4.7.5 Cấu trúc mẫu bảy nền kinh tế và phương sai giữa quốc
gia}\label{cux1ea5u-truxfac-mux1eabu-bux1ea3y-nux1ec1n-kinh-tux1ebf-vuxe0-phux1b0ux1a1ng-sai-giux1eefa-quux1ed1c-gia}}

Cấu trúc bảy nền kinh tế cho thấy mức độ phơi nhiễm xuất khẩu rất mỏng
và không đồng đều: tỷ lệ doanh nghiệp xuất khẩu dao động từ \(2{,}6\%\)
(Kiribati, FSTS trung bình \(0{,}011\)) đến \(27{,}9\%\) (Fiji, FSTS
trung bình \(0{,}109\)), với mức gộp toàn mẫu là \(13{,}8\%\) và FSTS
trung bình \(0{,}048\). Chỉ ba nền kinh tế (Fiji, Papua New Guinea,
Vanuatu) có hai đợt khảo sát WBES, phần còn lại chỉ có một đợt, nên nhận
dạng chủ yếu dựa trên biến thiên chéo trong từng ô quốc gia--năm. Hiệu
ứng cố định quốc gia hấp thụ một phần phương sai rất lớn: \(R^2\) tăng
từ \(0{,}109\) (chỉ hiệu ứng cố định năm) lên \(0{,}657\) khi bổ sung
hiệu ứng cố định quốc gia, nghĩa là khoảng \(83\%\) phương sai năng suất
được giải thích là phương sai giữa các quốc gia, một đặc trưng nhất quán
với tính dị biệt cấu trúc cực đoan giữa các nền kinh tế đảo nhỏ.

\hypertarget{kiux1ec3m-ux111ux1ecbnh-bux1ed1n-ux111iux1ec1u-kiux1ec7n-ux111ux1b0ux1eddng-cong-vuxe0-hux1ec7-sux1ed1-bux1eadc-hai-ux111ux1a1n-ux111iux1ec7u-uxe2m}{%
\subsection{4.7.6 Kiểm định bốn điều kiện đường cong và hệ số bậc hai
đơn điệu
âm}\label{kiux1ec3m-ux111ux1ecbnh-bux1ed1n-ux111iux1ec1u-kiux1ec7n-ux111ux1b0ux1eddng-cong-vuxe0-hux1ec7-sux1ed1-bux1eadc-hai-ux111ux1a1n-ux111iux1ec7u-uxe2m}}

Quan hệ đơn điệu âm được xác nhận qua việc cả bốn điều kiện cần của một
đường cong chữ U ngược (Haans và cộng sự, 2016) đều thất bại, đúng như
chữ ký kỳ vọng của hiện tượng gánh nặng quốc tế hóa bắt buộc. Ở mô hình
bậc hai M2, cả hai số hạng đều mất ý nghĩa khi đưa vào đồng thời:
\(\hat{\beta}(\text{FSTS}_c) = -2{,}177\) (\(SE = 1{,}603\);
\(p = {,}176\)) và \(\hat{\beta}(\text{FSTS}_c^2) = +1{,}125\)
(\(SE = 1{,}943\); \(p = {,}563\)). Điều đáng chú ý là dấu dương của hệ
số bậc hai (nếu có ý nghĩa) sẽ hàm ý một đường cong chữ U xuôi chứ không
phải chữ U ngược, tức một sự đảo ngược hoàn toàn của dạng hàm cổ điển,
bản thân điều này nhất quán với FIP. Cụ thể, điều kiện C1 (độ dốc dương
ở FSTS thấp) không thỏa vì hệ số bậc nhất âm và không có ý nghĩa; C2 (độ
cong âm \(\beta_2 < 0\)) không thỏa vì dấu dương; C3 (điểm uốn nằm trong
miền dữ liệu) cho điểm uốn đại số \(\text{TP}^* = 0{,}968\) nhưng không
số hạng nào có ý nghĩa thống kê; và C4 (kiểm định Lind--Mehlum) không
bác bỏ giả thuyết quan hệ đơn điệu âm. Sự suy giảm ý nghĩa khi đưa hai
số hạng vào cùng lúc phản ánh đa cộng tuyến cao trong mẫu FSTS thấp
(FSTS trung bình \(0{,}048\)).

\hypertarget{ux111ux1ed9-bux1ec1n-cux1ee7a-hux1ec7-sux1ed1-fip-qua-cuxe1c-ux111ux1eb7c-tux1ea3-vuxe0-cux1a1-chux1ebf-nux103ng-lux1ef1c-khuxf4ng-ux111iux1ec1u-tiux1ebft}{%
\subsection{4.7.7 Độ bền của hệ số FIP qua các đặc tả và cơ chế năng lực
không điều
tiết}\label{ux111ux1ed9-bux1ec1n-cux1ee7a-hux1ec7-sux1ed1-fip-qua-cuxe1c-ux111ux1eb7c-tux1ea3-vuxe0-cux1a1-chux1ebf-nux103ng-lux1ef1c-khuxf4ng-ux111iux1ec1u-tiux1ebft}}

Hệ số FIP duy trì dấu âm và quy mô đáng kể qua một chuỗi đặc tả độ vững,
củng cố tính cấu trúc của phát hiện. Khi loại bỏ hiệu ứng cố định quốc
gia (chỉ giữ hiệu ứng cố định năm), hệ số âm còn lớn hơn rõ rệt:
\(\hat{\beta}(\text{FSTS}) = -3{,}351\) (\(SE = 0{,}808\);
\(t = -4{,}15\); \(p < {,}001\)), cho thấy bất lợi cấu trúc cấp quốc gia
cộng dồn vào FIP cấp doanh nghiệp. Ở đặc tả song biến không kiểm soát,
quan hệ vẫn âm và ở biên ý nghĩa (\(\hat{\beta} = -0{,}864\);
\(p = {,}050\)). Hai kiểm định loại trừ bổ sung xác nhận FIP không phải
là sản phẩm phụ của mẫu hoàn chỉnh hay của định nghĩa bảy nền kinh tế:
bỏ biến kiểm soát sở hữu nước ngoài (vốn chỉ có ở \(22\%\) doanh nghiệp)
để khôi phục mẫu đầy đủ cho \(\hat{\beta}(\text{FSTS}_c) = -0{,}835\)
(\(p = {,}010\); \(N = 950\)), và mở rộng sang nhóm chín nền kinh tế đảo
nhỏ (thêm Comoros và Timor-Leste) vẫn cho
\(\hat{\beta}(\text{FSTS}_c) = -0{,}510\) (\(p = {,}008\);
\(N = 1.469\)).

Cơ chế năng lực không điều tiết được kiểm chứng ở M3. Năng lực công nghệ
có hiệu ứng trực tiếp dương và có ý nghĩa,
\(\hat{\beta}(\text{TCI}_z) = +0{,}299\) (\(SE = 0{,}099\);
\(p = {,}003\)), trong khi chấp nhận số không có ý nghĩa,
\(\hat{\beta}(\text{DAI}_z) = +0{,}094\) (\(SE = 0{,}088\);
\(p = {,}285\)). Cả hai năng lực đều không làm đổi độ dốc âm của FSTS,
nên giả thuyết không điều tiết được giữ vững. Mô thức năng lực nâng mặt
bằng năng suất nhưng không uốn được độ dốc này chính là điểm phân biệt
then chốt giữa bối cảnh SIDS và các bối cảnh châu Á lục địa: khi ba điều
kiện tiên quyết cấu trúc đồng thời vắng mặt, năng lực cấp doanh nghiệp
không thể chuyển hướng quan hệ quốc tế hóa và hiệu quả sang vùng dương.
Hệ số FSTS trong M3 vẫn âm và lớn hơn về độ lớn
(\(\hat{\beta} = -2{,}979\); \(p = {,}056\)), với việc mất ý nghĩa quy
cho cỡ mẫu giảm mạnh (còn \(N = 205\) do thiếu biến năng lực).

\begin{center}\rule{0.5\linewidth}{0.5pt}\end{center}

\hypertarget{kux1ebft-quux1ea3-bux1ed1i-cux1ea3nh-biux1ebfn-ux111ux1ed9ng-thux1ec3-chux1ebf-cux1ef1c-ux111oan-ux1ea5n-ux111ux1ed9-20142025-p9}{%
\section{4.8 Kết quả bối cảnh biến động thể chế cực đoan: Ấn Độ,
2014--2025
(P9)}\label{kux1ebft-quux1ea3-bux1ed1i-cux1ea3nh-biux1ebfn-ux111ux1ed9ng-thux1ec3-chux1ebf-cux1ef1c-ux111oan-ux1ea5n-ux111ux1ed9-20142025-p9}}

Phân tích trên ba đợt WBES Ấn Độ (2014 PICS3,
\(N_{\text{phân tích}} = 8.941\); 2022 BEE,
\(N_{\text{phân tích}} = 9.300\); 2025 BREADY,
\(N_{\text{phân tích}} = 10.476\)) cho mẫu gộp \(N = 28.717\), lớn nhất
trong các nghiên cứu đơn quốc gia của luận án. Khoảng thời gian
2014--2025 bao trùm một loạt cú sốc thể chế chưa từng có trong lịch sử
Ấn Độ hiện đại: bãi bỏ tiền mặt tháng 11/2016, ban hành Luật Phá sản và
Mất khả năng thanh toán (IBC) tháng 5/2016, áp dụng Thuế Hàng hóa và
Dịch vụ (GST) tháng 7/2017, ra mắt Giao diện Thanh toán Hợp nhất (UPI)
tháng 4/2016 đạt mức 12 tỷ giao dịch/tháng năm 2024, triển khai sáng
kiến Sản xuất tại Ấn Độ và chương trình Liên kết Khuyến khích Sản xuất
(PLI) năm 2020, cùng chính sách Tự chủ Ấn Độ (Atmanirbhar Bharat) định
hướng lại chiến lược quốc gia theo hướng nội địa hóa giai đoạn 2020.

Kết quả thực nghiệm cho thấy một dạng phá vỡ ngưỡng (threshold collapse)
không quan sát được ở bất kỳ bối cảnh quốc gia nào khác trong portfolio.
Đợt 2014 ghi nhận quan hệ U ngược cổ điển với hệ số bậc nhất dương
(\(\beta_1 = +1{,}865\), \(p < {,}001\)) và bậc hai âm
(\(\beta_2 = -1{,}508\), \(p < {,}001\)), điểm uốn ước lượng tại
\(TP = 61{,}8\%\) FSTS (khoảng tin cậy delta-method 95\%:
\([55{,}0\%; 68{,}6\%]\)), Lind--Mehlum xác nhận đường cong U ngược ở
mức \(p < {,}001\). Đợt 2022, sau ba cú sốc thể chế đầu tiên (bãi bỏ
tiền mặt, GST, IBC), đường cong U ngược vẫn duy trì nhưng điểm uốn đã
dịch chuyển đáng kể vào trong: \(\beta_1 = +1{,}542\),
\(\beta_2 = -1{,}893\), \(TP = 40{,}7\%\) (khoảng tin cậy 95\%:
\([38{,}0\%; 43{,}5\%]\)). Đợt 2025, hệ số bậc hai mất ý nghĩa thống kê
(\(\beta_2 = -0{,}160\), \(p = {,}420\)), quan hệ trở thành đơn điệu âm
trên toàn dải cường độ xuất khẩu thực tế.

Kiểm định Paternoster so sánh hệ số giữa 2014 và 2025 bác bỏ giả thuyết
cân bằng với cường độ cực mạnh: \(z(\text{FSTS}) = -7{,}94\),
\(p < {,}0001\) và \(z(\text{FSTS}^2) = +4{,}17\), \(p < {,}0001\) dưới
sai số chuẩn HC1; dưới sai số chuẩn cụm theo bang (cluster-robust,
\(G = 24\) bang), kết quả vẫn có ý nghĩa thống kê:
\(z(\text{FSTS}) = -3{,}50\), \(p = {,}001\) và
\(z(\text{FSTS}^2) = +2{,}19\), \(p = {,}029\). Đặc tả gộp ba đợt với
biến giả đợt khảo sát xác nhận sự dịch chuyển có hệ thống:
\(\beta(\text{FSTS} \times \text{wave}\_2025) = -1{,}63\)
(\(p < {,}0001\)), cho thấy độ dốc đường cong quốc tế hóa và hiệu quả đã
thay đổi căn bản giữa 2014 và 2025.

\textbf{Bảng 4.8.} \emph{Điểm uốn theo đợt khảo sát, Ấn Độ.}

\begin{longtable}[]{@{}
  >{\raggedright\arraybackslash}p{(\columnwidth - 8\tabcolsep) * \real{0.2000}}
  >{\raggedright\arraybackslash}p{(\columnwidth - 8\tabcolsep) * \real{0.2000}}
  >{\raggedright\arraybackslash}p{(\columnwidth - 8\tabcolsep) * \real{0.2000}}
  >{\raggedright\arraybackslash}p{(\columnwidth - 8\tabcolsep) * \real{0.2000}}
  >{\raggedright\arraybackslash}p{(\columnwidth - 8\tabcolsep) * \real{0.2000}}@{}}
\toprule\noalign{}
\begin{minipage}[b]{\linewidth}\raggedright
Đợt khảo sát
\end{minipage} & \begin{minipage}[b]{\linewidth}\raggedright
\(N\) phân tích
\end{minipage} & \begin{minipage}[b]{\linewidth}\raggedright
Điểm uốn (TP)
\end{minipage} & \begin{minipage}[b]{\linewidth}\raggedright
Lind--Mehlum
\end{minipage} & \begin{minipage}[b]{\linewidth}\raggedright
Nhận xét
\end{minipage} \\
\midrule\noalign{}
\endhead
\bottomrule\noalign{}
\endlastfoot
2014 (PICS3) & 8.941 & 61,8\% FSTS & \(p < {,}001\) & U ngược cổ điển \\
2022 (BEE) & 9.300 & 40,7\% FSTS & \(p < {,}001\) & U ngược, TP dịch
trong \\
2025 (BREADY) & 10.476 & Không xác định & \(p = {,}99\) & Đường cong sụp
đổ \\
\end{longtable}

\emph{Nguồn: phân tích Ấn Độ của tác giả (nghiên cứu thành phần P9), mở
rộng công trình đã công bố Do và Phan (2025, IntechOpen, phân tích 380
doanh nghiệp WBES Ấn Độ, ước lượng FGLS); bản thảo mở rộng kiểm định độ
bền ngưỡng I--P qua ba đợt 2014--2025 đang chuẩn bị cho MIR/IJOEM.}

Phát hiện ngưỡng phá vỡ ở Ấn Độ trái ngược trực tiếp với kết quả Trung
Quốc thuộc chế độ thu nhập trung bình cao (Nhóm III ICRV), trong khi Ấn
Độ thuộc Nhóm IV (chuyển đổi thu nhập trung bình--thấp), nơi điểm uốn
của Trung Quốc được chứng minh là bền vững cấu trúc qua 12 năm tăng
trưởng tích lũy (P5 China: \(TP = 49{,}4\%\) năm 2012 so với
\(TP = 47{,}2\%\) năm 2024, Paternoster \(z = +{0,}82\), \(p = {,}412\)
chấp nhận giả thuyết cân bằng). Sự đối lập này cho thấy bền vững ngưỡng
không phải đặc tính phổ quát của chế độ thể chế ICRV III mà phụ thuộc
vào mức độ và tốc độ thay đổi thể chế tích lũy trong cửa sổ quan sát.
Cần lưu ý thêm rằng điểm uốn Ấn Độ năm 2014 (\(61{,}8\%\)) cao bất
thường so với dải đặc trưng của Nhóm IV (điểm uốn vùng giữa khoảng
\(43\%\) trong mẫu gộp P7), gần với vùng của các chế độ thể chế mạnh
hơn; chính sự \textbf{hội tụ} của điểm uốn xuống \(40{,}7\%\) năm 2022
mới là dữ kiện nhất quán với chuẩn Nhóm IV, chứ không phải mức 2014.
Cách đọc này biến một mâu thuẫn bề ngoài (một nền Nhóm IV khởi đầu với
điểm uốn kiểu Nhóm I) thành bằng chứng ủng hộ: Ấn Độ tiền cú sốc vận
hành như một chế độ thể chế cao hơn vị trí thu nhập của nó, và chỉ chuỗi
cú sốc thể chế sau 2016 mới kéo điểm uốn về dải Nhóm IV trước khi làm
tan rã hẳn cấu trúc vào 2025.

Vai trò của TCI tại Ấn Độ cho thấy hiện tượng đảo dấu (sign flip) chưa
từng quan sát ở các bối cảnh quốc gia khác. Năm 2014, hiệu ứng trực tiếp
TCI dương đáng kể (\(\beta = +0{,}12\), \(p < {,}001\)), năm 2022 tiếp
tục dương và mạnh hơn (\(\beta = +0{,}19\), \(p < {,}001\)), nhưng năm
2025 đảo sang âm (\(\beta = -0{,}08\), \(p < {,}001\)) đi kèm tương tác
FSTS\(\times\)TCI dương (\(\beta = +0{,}20\), \(p = {,}03\)). Diễn giải
nhất quán nhất với khung Mizik và Jacobson (2003) về đánh đổi chú trọng
chiến lược (strategic emphasis trade-off): khi chính sách PLI và
Atmanirbhar Bharat sau 2020 chuyển hướng năng lực công nghệ sang sản
xuất nội địa, đầu tư TCI của doanh nghiệp định hướng nội địa cạnh tranh
nguồn lực với hoạt động xuất khẩu, làm triệt tiêu lợi ích năng suất TCI
vốn quan sát được khi tích lũy năng lực không thiên kiến giữa nội địa và
xuất khẩu.

Vai trò của DAI tại Ấn Độ làm rõ một sự tinh chỉnh lý thuyết quan trọng
cho khung CDCM. Tại đợt 2025, ngoài chỉ báo Tầng 1 (trang web) chung cho
cả ba đợt, WBES BREADY cung cấp thêm chỉ báo Tầng 2 (tỷ trọng doanh thu
nhận qua thanh toán điện tử, biến \texttt{k33}) làm bối cảnh thí nghiệm
gần tự nhiên cho hạ tầng số công cộng UPI. Tương tác DAI Tầng 2 với
cường độ xuất khẩu lại âm và có ý nghĩa thống kê:
\(\text{FSTS} \times \text{DAI}_{\text{Tầng 2}} = -4{,}02\),
\(p = {,}004\), ngược với dự đoán phụ phẩm cộng hưởng năng lực số công
cộng. Diễn giải cộng hưởng (complementarity reading): hạ tầng số công
cộng UPI vận hành trên đường ray rupee nội địa và yêu cầu KYC Aadhaar
cho đối tác cư trú tại Ấn Độ; doanh nghiệp xuất khẩu cần hạ tầng tài
chính xuyên biên giới khác hẳn (SWIFT, ngân hàng đại lý, tài trợ thương
mại, bảo hiểm tín dụng xuất khẩu) mà UPI không thể thay thế. Phát hiện
này cho thấy hạ tầng công cộng số chỉ thay thế năng lực tư nhân khi hai
loại hạ tầng có cùng định hướng mục đích, không phải thay thế phổ quát
cho mọi hoạt động sản xuất doanh nghiệp.

\hypertarget{muxf4-tux1ea3-mux1eabu-vuxe0-phux1b0ux1a1ng-truxecnh-hux1ed3i-quy-tux1eebng-ux111ux1ee3t}{%
\subsection{4.8.1 Mô tả mẫu và phương trình hồi quy từng
đợt}\label{muxf4-tux1ea3-mux1eabu-vuxe0-phux1b0ux1a1ng-truxecnh-hux1ed3i-quy-tux1eebng-ux111ux1ee3t}}

Mô tả mẫu cho thấy một sự thu hẹp xuất khẩu mang tính hệ thống đồng hành
với sự sụp đổ ngưỡng. Cường độ xuất khẩu trung bình dịch chuyển từ
\(7{,}2\%\) (2014) lên \(8{,}0\%\) (2022) rồi rơi xuống \(2{,}7\%\)
(2025), nhất quán với bước ngoặt hướng nội sau 2020 dưới chính sách
Atmanirbhar Bharat; mặt bằng năng suất tăng đều, với \(\ln\text{LP}\)
trung bình lần lượt là \(13{,}93\); \(14{,}37\) và \(15{,}01\). Tỷ lệ áp
dụng website (DAI Tầng 1) đạt \(52{,}7\%\) (2014), \(60{,}6\%\) (2022)
rồi giảm còn \(41{,}8\%\) (2025), trong khi tỷ trọng doanh thu qua thanh
toán điện tử (DAI Tầng 2, chỉ đo năm 2025) trung bình \(65{,}8\%\), phản
ánh mức bão hòa thanh toán hợp nhất (Unified Payments Interface ---
UPI).

Các phương trình hồi quy bậc hai từng đợt làm rõ động học sụp đổ ngưỡng.
Năm 2014, đường cong là chữ U ngược cổ điển:
\(\ln\text{LP} = 13{,}51 + 1{,}86 \cdot \text{FSTS} - 1{,}51 \cdot \text{FSTS}^2 + \mathbf{X}\boldsymbol{\beta}\),
với điểm uốn \(61{,}8\%\). Năm 2022, đường cong vẫn duy trì nhưng điểm
uốn dịch vào trong \(21\) điểm phần trăm:
\(\ln\text{LP} = 13{,}47 + 1{,}54 \cdot \text{FSTS} - 1{,}89 \cdot \text{FSTS}^2 + \mathbf{X}\boldsymbol{\beta}\),
với điểm uốn \(40{,}7\%\) và đây là đường chữ U ngược mạnh nhất trong ba
đợt (Lind--Mehlum liên hợp \(p < 1\mathrm{e}{-}12\)). Năm 2025, độ cong
mất ý nghĩa:
\(\ln\text{LP} = 14{,}24 - 0{,}36 \cdot \text{FSTS} - 0{,}16 \cdot \text{FSTS}^2 + \mathbf{X}\boldsymbol{\beta}\),
quan hệ trở nên đơn điệu âm trên toàn dải, với độ dốc tại
\(\text{FSTS}_{\min}\) là \(-0{,}36\) (\(p = {,}02\)) và tại
\(\text{FSTS}_{\max}\) là \(-0{,}68\) (\(p = {,}005\)).

\hypertarget{ux111ux1eb7c-tux1ea3-gux1ed9p-ba-ux111ux1ee3t-vuxe0-bux1eb1ng-chux1ee9ng-dux1ecbch-chuyux1ec3n-lux169y-tiux1ebfn}{%
\subsection{4.8.2 Đặc tả gộp ba đợt và bằng chứng dịch chuyển lũy
tiến}\label{ux111ux1eb7c-tux1ea3-gux1ed9p-ba-ux111ux1ee3t-vuxe0-bux1eb1ng-chux1ee9ng-dux1ecbch-chuyux1ec3n-lux169y-tiux1ebfn}}

Mô hình gộp ba đợt với biến giả đợt khảo sát (\(N = 28.717\); \(R^2\)
hiệu chỉnh \(= 0{,}197\)) phân giải sự dịch chuyển thành các thành phần
có thể kiểm định riêng. Hệ số \(\text{Wave}\_2022 = +0{,}46\)
(\(p < {,}0001\)) và \(\text{Wave}\_2025 = +1{,}12\) (\(p < {,}0001\))
xác nhận đà nâng mặt bằng năng suất. Quan trọng hơn, tương tác
\(\text{FSTS} \times \text{wave}\_2022 = +0{,}45\) (\(p = {,}15\)) không
cho thấy dịch chuyển bền vững so với mốc 2014, trong khi
\(\text{FSTS} \times \text{wave}\_2025 = -1{,}63\) (\(p < {,}0001\)) xác
nhận sự sụp đổ độ dốc năm 2025; ở vế bậc hai,
\(\text{FSTS}^2 \times \text{wave}\_2022 = -1{,}09\) (\(p = {,}001\))
phản ánh điểm uốn 2022 dịch vào trong qua việc độ cong dốc thêm, còn
\(\text{FSTS}^2 \times \text{wave}\_2025 = +0{,}86\) (\(p = {,}007\))
phản ánh độ cong làm phẳng năm 2025. Đặc tả xu hướng tuyến tính cho
\(\text{FSTS} \times \text{trend} = -0{,}10\) (\(p < {,}001\)), nhất
quán với một sự dịch chuyển cấu trúc lũy tiến chứ không phải đứt gãy đột
ngột do thay đổi công cụ khảo sát.

\hypertarget{nux103m-kiux1ec3m-ux111ux1ecbnh-ux111ux1ed9-bux1ec1n-vuxe0-sux1ef1-cuxf4-ux111ux1eb7c-cux1ee7a-chux1eef-u-ngux1b0ux1ee3c-2014-ux1edf-doanh-nghiux1ec7p-lux1edbn}{%
\subsection{4.8.3 Năm kiểm định độ bền và sự cô đặc của chữ U ngược 2014
ở doanh nghiệp
lớn}\label{nux103m-kiux1ec3m-ux111ux1ecbnh-ux111ux1ed9-bux1ec1n-vuxe0-sux1ef1-cuxf4-ux111ux1eb7c-cux1ee7a-chux1eef-u-ngux1b0ux1ee3c-2014-ux1edf-doanh-nghiux1ec7p-lux1edbn}}

Năm đặc tả độ vững củng cố diễn giải cấu trúc. (R1) Mẫu chỉ doanh nghiệp
chế biến chế tạo làm mô thức mạnh thêm đôi chút: điểm uốn \(57{,}5\%\)
(2014), \(41{,}3\%\) (2022), độ cong 2025 không có ý nghĩa
(\(\beta_2 = -0{,}46\); \(p = {,}017\)); kiểm định Paternoster chéo đợt
trên mẫu này cho \(z(\text{FSTS}) = -6{,}89\) (\(p < {,}0001\)) và
\(z(\text{FSTS}^2) = +3{,}64\) (\(p = {,}0003\)), với độ lớn hiệu ứng
nằm trong khoảng \(12\%\) của ước lượng toàn mẫu, bác bỏ khả năng kết
quả do thành phần doanh nghiệp dịch vụ. (R2) Cắt bỏ doanh nghiệp có FSTS
\(> 0{,}95\) cho điểm uốn \(48{,}7\%\) (2014) và \(39{,}7\%\) (2022), độ
cong 2025 vẫn không có ý nghĩa (\(\beta_2 = -0{,}30\); \(p = {,}37\)).

Phân tích theo quy mô (R3) phát hiện một sắc thái tinh tế. Ở mẫu doanh
nghiệp nhỏ (dưới 100 lao động), chữ U ngược năm 2014 không được hỗ trợ
(\(\beta_2 = -0{,}07\); \(p = {,}84\)), chỉ định hình ở 2022
(\(\beta_2 = -1{,}64\); \(p < {,}0001\); điểm uốn \(41{,}2\%\)); ngược
lại ở mẫu doanh nghiệp lớn (từ 100 lao động), chữ U ngược hiện diện ở
2014 (điểm uốn \(50{,}2\%\)) và 2022 (điểm uốn \(40{,}8\%\)) với độ cong
rất lớn (\(\beta_2\) giữa \(-2{,}4\) và \(-3{,}0\)). Điều này cho thấy
chữ U ngược năm 2014 cô đặc ở doanh nghiệp lớn và chỉ lan tới phân khúc
nhỏ vào 2022, trùng khớp với sự mở rộng tham gia xuất khẩu của doanh
nghiệp nhỏ và vừa dưới chính sách thúc đẩy xuất khẩu giai đoạn
2014--2020. (R4) Thay biến phụ thuộc bằng năng suất lao động chuẩn hóa
theo mức cho điểm uốn \(54{,}4\%\) (2014, Lind--Mehlum ở biên
\(p = {,}048\)), \(41{,}5\%\) (2022) và độ cong sụp đổ ở 2025
(\(\beta_2 = +0{,}03\); \(p = {,}90\)).

Một lưu ý đo lường quan trọng cần đi kèm: đợt 2025 trùng với sự co hẹp
mạnh tham gia xuất khẩu, khi số doanh nghiệp có FSTS trên \(50\%\) giảm
từ \(536\) xuống còn \(141\), nên việc mất độ cong bậc hai năm 2025 một
phần phản ánh sự suy giảm biến thiên nhận dạng ở đuôi xuất khẩu cao. Tuy
vậy, chính dấu của độ dốc ở vùng xuất khẩu thấp cũng đảo chiều, từ nhánh
đi lên dương năm 2014 sang đơn điệu âm năm 2025, một thay đổi định tính
mà riêng việc đuôi mỏng đi không thể tạo ra. Đối với điều tiết DAI Tầng
2, bên cạnh tương tác bậc nhất âm đã nêu, số hạng bậc hai
\(\text{FSTS}^2 \times \text{DAI}_{\text{epay}} = +3{,}39\)
(\(p = {,}10\)); và DAI Tầng 1 (website) có hiệu ứng trực tiếp dương ở
cả ba đợt (\(\beta \approx +0{,}38\) đến \(+0{,}39\)), với tương tác
\(\text{FSTS} \times \text{DAI}\) trở thành dương có ý nghĩa ở 2022
(\(+1{,}43\); \(p = {,}02\)) rồi mất ý nghĩa ở 2025, nhất quán với vai
trò biên của website giảm dần khi UPI hấp thụ phần tương tác số hướng
tới khách hàng.

\hypertarget{bux1eb1ng-chux1ee9ng-tux1eadp-trung-vuxe0o-quux1ea3n-trux1ecb-cux1ea5p-cao-tux1ea1i-ux1ea5n-ux111ux1ed9-vai-truxf2-ux111iux1ec1u-tiux1ebft-cux1ee7a-kinh-nghiux1ec7m-vuxe0-giux1edbi-tuxednh-chux1b0ux1a1ng-suxe1ch}{%
\subsection{4.8.4 Bằng chứng tập trung vào quản trị cấp cao tại Ấn Độ:
vai trò điều tiết của kinh nghiệm và giới tính (chương
sách)}\label{bux1eb1ng-chux1ee9ng-tux1eadp-trung-vuxe0o-quux1ea3n-trux1ecb-cux1ea5p-cao-tux1ea1i-ux1ea5n-ux111ux1ed9-vai-truxf2-ux111iux1ec1u-tiux1ebft-cux1ee7a-kinh-nghiux1ec7m-vuxe0-giux1edbi-tuxednh-chux1b0ux1a1ng-suxe1ch}}

Bổ sung cho phân tích đường cong theo thời gian của P9, một nghiên cứu
thành phần khác cùng bối cảnh Ấn Độ, được công bố dưới dạng chương sách
(Do và Phan, 2025), tập trung riêng vào vai trò điều tiết của đặc điểm
nhà quản trị cấp cao. Nghiên cứu phân tích \textbf{380 doanh nghiệp}
thuộc cả khu vực chế biến chế tạo và dịch vụ (760 quan sát doanh nghiệp
và năm) từ dữ liệu WBES, sử dụng hồi quy bình phương nhỏ nhất tổng quát
khả thi (FGLS) để xử lý phương sai sai số thay đổi. Biến phụ thuộc ở đây
là tỷ suất lợi nhuận trên doanh thu (ROS, trung bình 30,34\%), khác với
thước đo năng suất lao động dùng ở các phần còn lại của luận án, và đặc
tả là tuyến tính thay vì bậc hai.

Kết quả xác nhận ba mệnh đề. Thứ nhất, mức độ quốc tế hóa tác động âm
lên hiệu quả ở đặc tả tuyến tính (\(\hat{\beta} = -36{,}66\),
\(p < {,}001\)), nhất quán với logic chi phí gia nhập giai đoạn đầu của
mô hình Uppsala, tức nhánh đi xuống bên trái của quan hệ phi tuyến khi
quan sát ở dải cường độ quốc tế hóa thấp đặc trưng của mẫu này. Thứ hai,
kinh nghiệm của nhà quản trị điều tiết dương quan hệ này
(\(\hat{\beta} = +1{,}55\), \(p < {,}001\)), làm dịu tác động bất lợi
của quốc tế hóa lên lợi nhuận. Thứ ba, lãnh đạo nữ cũng điều tiết dương
quan hệ (\(\hat{\beta} = +77{,}03\), \(p < {,}001\)), củng cố giả thuyết
H4 ở khía cạnh điều tiết.

Phát hiện về giới tính ở đây cần được đối chiếu thận trọng với hệ số âm
của biến quản trị nữ trong các ước lượng dựa trên năng suất (P7 toàn mẫu
và P10 Nhật Bản). Sự khác dấu này không phải mâu thuẫn nội tại mà phản
ánh ba khác biệt thiết kế: biến phụ thuộc là lợi nhuận trên doanh thu
chứ không phải năng suất lao động; tham số ước lượng là hệ số điều tiết
độ dốc của quan hệ quốc tế hóa và hiệu quả chứ không phải hiệu ứng mức
chính; và mẫu là 380 doanh nghiệp Ấn Độ chứ không phải mẫu gộp đa quốc
gia. Cách đọc hài hòa, nhất quán với phần tổng hợp giả thuyết của luận
án, là đặc điểm nhà quản trị cấp cao có thể nâng chất lượng quyết định
quốc tế hóa (biểu hiện thành điều tiết dương trên thước đo lợi nhuận
trong một bối cảnh đơn quốc gia), trong khi hệ số mức âm trên thước đo
năng suất ở mẫu gộp phản ánh chọn lọc cấu thành chứ không phải hiệu quả
lãnh đạo. Đặt cạnh nhau, hai mảnh bằng chứng cho thấy vai trò của đặc
trưng quản trị là phụ thuộc thước đo và phụ thuộc bối cảnh, đúng như kết
luận chung của luận án rằng H4 không được hỗ trợ như một hiệu ứng mức
phổ quát mà cần phân tích nhận biết chọn lọc trong nội bộ từng bối cảnh.

\hypertarget{cux1ea5u-truxfac-mux1eabu-ba-ux111ux1ee3t-vuxe0-khoux1ea3ng-tin-cux1eady-ux111iux1ec3m-uux1ed1n-theo-phux1b0ux1a1ng-phuxe1p-delta}{%
\subsection{4.8.5 Cấu trúc mẫu ba đợt và khoảng tin cậy điểm uốn theo
phương pháp
delta}\label{cux1ea5u-truxfac-mux1eabu-ba-ux111ux1ee3t-vuxe0-khoux1ea3ng-tin-cux1eady-ux111iux1ec3m-uux1ed1n-theo-phux1b0ux1a1ng-phuxe1p-delta}}

Bộ ba khoảng tin cậy điểm uốn theo phương pháp delta cung cấp bằng chứng
định lượng trực tiếp cho luận điểm dịch chuyển ngưỡng thay vì chỉ dựa
vào giá trị điểm. Mẫu phân tích thu được sau khi loại bỏ theo danh sách
trên bộ biến trọng tâm và kiểm tra tính hợp lệ của tuổi doanh nghiệp là
\(8.941\) (2014 PICS3, từ \(9.281\) thô), \(9.300\) (2022 BEE, từ
\(9.376\) thô) và \(10.476\) (2025 BREADY, từ \(10.479\) thô), cho mẫu
gộp phân tích \(28.717\) trên tổng thô \(29.136\). Điểm uốn 2014 ở mức
\(61{,}8\%\) với khoảng tin cậy delta-method \([55{,}0\%; 68{,}6\%]\),
một mức đỉnh năng suất rất cao đặc trưng cho bối cảnh thị trường mới nổi
nơi doanh nghiệp xuất khẩu chiết xuất lợi suất học hỏi qua xuất khẩu
đáng kể trước khi chi phí mở rộng quá mức bám chặt. Điểm uốn 2022 dịch
vào trong còn \(40{,}7\%\) với khoảng tin cậy \([38{,}0\%; 43{,}5\%]\),
và đây là chữ U ngược mạnh nhất ba đợt với kiểm định Lind--Mehlum liên
hợp \(p < 1\mathrm{e}{-}12\). Điểm cốt lõi là hai khoảng tin cậy 2014 và
2022 \textbf{không chồng lấn} (\([55{,}0; 68{,}6]\) so với
\([38{,}0; 43{,}5]\)), bằng chứng định lượng trực tiếp rằng dịch chuyển
ngưỡng \(21\) điểm phần trăm là thực chứ không phải nhiễu ước lượng,
tương phản sắc nét với ba khoảng tin cậy chồng lấn chặt của Trung Quốc
đã trình bày ở Mục 4.4.4 (nơi điểm uốn ổn định cấu trúc).

\hypertarget{phux1ea3n-buxe1c-giux1ea3-thuyux1ebft-tux1ea1o-tuxe1c-cuxf4ng-cux1ee5-khux1ea3o-suxe1t-qua-ux111ux1ee3t-cux1ea7u-nux1ed1i-2022}{%
\subsection{4.8.6 Phản bác giả thuyết tạo tác công cụ khảo sát qua đợt
cầu nối
2022}\label{phux1ea3n-buxe1c-giux1ea3-thuyux1ebft-tux1ea1o-tuxe1c-cuxf4ng-cux1ee5-khux1ea3o-suxe1t-qua-ux111ux1ee3t-cux1ea7u-nux1ed1i-2022}}

Một quan ngại tự nhiên là liệu sự dịch chuyển biểu kiến có phải tạo tác
của việc thay đổi công cụ khảo sát hay không, vì đợt 2014 dùng công cụ
PICS3 còn đợt 2025 dùng công cụ BREADY mới, hai bộ không đồng nhất về
diễn đạt câu hỏi và thiết kế chọn mẫu. Đợt cầu nối 2022 (BEE trung gian)
cung cấp kiểm định quyết định: nếu sự sụp đổ 2025 thuần túy là tạo tác
công cụ, một bước gián đoạn sẽ phải xuất hiện giữa 2022 và 2025 (chuyển
tiếp BEE sang BREADY), trong khi 2014 và 2022 không phân biệt được. Dữ
liệu cho thấy mô thức ngược lại: điểm uốn dịch đơn điệu vào trong từ
\(61{,}8\%\) (2014) xuống \(40{,}7\%\) (2022) qua cửa sổ bãi bỏ tiền
mặt, GST và IBC trên cùng công cụ phái sinh từ PICS, rồi quan hệ mới tan
rã vào 2025; đặc tả xu hướng tuyến tính cho
\(\mathrm{FSTS} \times \mathrm{trend} = -0{,}10\) (\(p < {,}001\)) khẳng
định một tiến trình đơn điệu thay vì đứt gãy đột ngột. Mô hình gộp ba
đợt (\(N = 28.717\); \(R^2\) hiệu chỉnh \(= 0{,}197\)) định lượng các
thành phần dịch chuyển:
\(\mathrm{FSTS}^2 \times \mathrm{wave}\_2022 = -1{,}09\)
(\(p = {,}001\)) xác nhận điểm uốn 2022 dịch vào trong qua độ cong dốc
thêm, còn \(\mathrm{FSTS}^2 \times \mathrm{wave}\_2025 = +0{,}86\)
(\(p = {,}007\)) xác nhận độ cong làm phẳng năm 2025. Tiến trình dịch
chuyển cấu trúc khớp dữ liệu; tạo tác công cụ thì không.

\hypertarget{phuxe2n-ruxe3-theo-quy-muxf4-vuxe0-nux103m-kiux1ec3m-ux111ux1ecbnh-ux111ux1ed9-vux1eefng-ux111ux1ecbnh-vux1ecb-nguux1ed3n-dux1ecbch-chuyux1ec3n}{%
\subsection{4.8.7 Phân rã theo quy mô và năm kiểm định độ vững định vị
nguồn dịch
chuyển}\label{phuxe2n-ruxe3-theo-quy-muxf4-vuxe0-nux103m-kiux1ec3m-ux111ux1ecbnh-ux111ux1ed9-vux1eefng-ux111ux1ecbnh-vux1ecb-nguux1ed3n-dux1ecbch-chuyux1ec3n}}

Năm kiểm định độ vững định vị chính xác phân khúc doanh nghiệp nơi chữ U
ngược hình thành rồi tan rã. Mẫu chỉ chế biến chế tạo (ISIC 15--37 cho
2014; ISIC 10--33 cho 2022 và 2025; \(N = 6.974\), \(5.366\), \(5.706\))
làm mô thức mạnh thêm với điểm uốn \(57{,}5\%\) (2014) và \(41{,}3\%\)
(2022), kiểm định Paternoster chéo đợt cho
\(z(\mathrm{FSTS}) = -6{,}89\) (\(p < {,}0001\)) và
\(z(\mathrm{FSTS}^2) = +3{,}64\) (\(p = {,}0003\)), với độ lớn hiệu ứng
nằm trong khoảng \(12\%\) của ước lượng toàn mẫu, bác bỏ khả năng kết
quả do thành phần doanh nghiệp dịch vụ. Phân rã theo quy mô phát hiện
sắc thái then chốt: ở mẫu doanh nghiệp nhỏ (dưới 100 lao động;
\(N = 6.844\), \(6.497\), \(7.383\)), chữ U ngược 2014 không được hỗ trợ
(\(\beta_2 = -0{,}07\), \(p = {,}84\)), chỉ định hình ở 2022
(\(\beta_2 = -1{,}64\), \(p < {,}0001\); điểm uốn \(41{,}2\%\)); ngược
lại ở mẫu doanh nghiệp lớn (từ 100 lao động; \(N = 2.097\), \(2.803\),
\(3.093\)), chữ U ngược hiện diện ở 2014 (điểm uốn \(50{,}2\%\)) và 2022
(điểm uốn \(40{,}8\%\)) với độ cong rất lớn (\(\beta_2\) giữa \(-2{,}4\)
và \(-3{,}0\)). Điều này cho thấy chữ U ngược 2014 cô đặc ở doanh nghiệp
lớn và chỉ lan tới phân khúc nhỏ vào 2022, trùng khớp với sự mở rộng
tham gia xuất khẩu của doanh nghiệp nhỏ và vừa dưới chính sách thúc đẩy
xuất khẩu giai đoạn 2014--2020 rồi đảo chiều sau 2020 dưới Atmanirbhar
Bharat. Bốn trong năm đặc tả (R1, R2, R3b, R4) xác nhận mô thức định
tính chữ U ngược hiện diện 2014--2022 rồi biến mất 2025; đặc tả thứ năm
(R3a) xác nhận hiện diện 2022 và biến mất 2025.

\hypertarget{ux111ux1ed9-vux1eefng-suy-luux1eadn-trux1b0ux1edbc-sai-sux1ed1-chuux1ea9n-cux1ee5m-vuxe0-quan-ngux1ea1i-nux1ed9i-sinh-upi}{%
\subsection{4.8.8 Độ vững suy luận trước sai số chuẩn cụm và quan ngại
nội sinh
UPI}\label{ux111ux1ed9-vux1eefng-suy-luux1eadn-trux1b0ux1edbc-sai-sux1ed1-chuux1ea9n-cux1ee5m-vuxe0-quan-ngux1ea1i-nux1ed9i-sinh-upi}}

Tính bền vững của phát hiện dịch chuyển ngưỡng được kiểm tra trước hai
cấp độ suy luận. Dưới sai số chuẩn vững HC1, kiểm định Paternoster bác
bỏ cân bằng hệ số 2014--2025 với cường độ cực mạnh
(\(z(\mathrm{FSTS}) = -7{,}94\); \(z(\mathrm{FSTS}^2) = +4{,}17\); cả
hai \(p < {,}0001\)); dưới sai số chuẩn cụm theo bang bảo thủ hơn
(khoảng 24 cụm, theo Cameron và cộng sự), kết quả vẫn có ý nghĩa
(\(z(\mathrm{FSTS}) = -3{,}50\), \(p = {,}0005\);
\(z(\mathrm{FSTS}^2) = +2{,}19\), \(p = {,}029\)) bất chấp cỡ mẫu hiệu
dụng nhỏ hơn nhiều mà 24 cụm bang hàm ý. Với số cụm ở mức biên này, biện
minh tiệm cận cho suy luận cụm-vững là cận biên, nên đây được đọc là
kiểm định bảo thủ chứ không phải đặc tả chính. Đối với tương tác UPI
Tầng 2 âm
(\(\mathrm{FSTS} \times \mathrm{DAI}_{\text{epay}} = -4{,}02\),
\(p = {,}004\)), một quan ngại nội sinh được thừa nhận rõ: ở các doanh
nghiệp 2025 có tỉ trọng thanh toán UPI cao, phân phối cường độ xuất khẩu
bị cắt cụt về phía thấp một cách cơ học, vì doanh nghiệp có khách hàng
nước ngoài không thanh toán qua UPI nên tỉ trọng UPI cao và tỉ trọng
xuất khẩu cao nghịch biến ngay trong quá trình sinh dữ liệu. Do đó việc
bác bỏ H4b (Tầng 2 lấn át Tầng 1) được đọc là gợi ý chứ không dứt khoát,
với diễn giải cộng hưởng là cách giải thích nhất quán lý thuyết nhất: hạ
tầng số công cộng chỉ thay thế năng lực tư nhân khi cùng định hướng mục
đích, và UPI định hướng nội địa không thể thay thế hạ tầng tài chính
xuyên biên giới (SWIFT, ngân hàng đại lý, tài trợ thương mại) mà doanh
nghiệp xuất khẩu cần.

\begin{center}\rule{0.5\linewidth}{0.5pt}\end{center}

\hypertarget{kux1ebft-quux1ea3-bux1ed1i-cux1ea3nh-tiuxean-tiux1ebfn-trux1b0ux1edfng-thuxe0nh-nhux1eadt-bux1ea3n-p10}{%
\section{4.9 Kết quả bối cảnh tiên tiến trưởng thành: Nhật Bản
(P10)}\label{kux1ebft-quux1ea3-bux1ed1i-cux1ea3nh-tiuxean-tiux1ebfn-trux1b0ux1edfng-thuxe0nh-nhux1eadt-bux1ea3n-p10}}

Nghiên cứu thành phần P10 khai thác đợt khảo sát WBES khai mở của Nhật
Bản năm 2025 (\(N = 2.168\) cơ sở; 339 biến hài hòa; chọn mẫu ngẫu nhiên
phân tầng theo ngành, quy mô và vùng) để kiểm định trực tiếp \textbf{dự
đoán biên trên} của khung phụ thuộc thể chế: tại chế độ thể chế mạnh
nhất, nhánh đi xuống của chữ U ngược bị đẩy ra ngoài miền cường độ xuất
khẩu quan sát được, khiến quan hệ quốc tế hóa và hiệu quả trong dữ liệu
trở nên dương và xấp xỉ tuyến tính. Bối cảnh mang sức nặng lý thuyết đặc
biệt: chính các tập đoàn đa quốc gia Nhật Bản là nơi chữ U ngược được
thiết lập ban đầu, song Nhật Bản chưa từng được quan sát trên một bộ
công cụ hài hòa cho phép định vị nó so với các nền kinh tế cùng khu vực.
Năng suất lao động (ln doanh thu trên một lao động thường xuyên,
winsorize tại phân vị 1/99) được hồi quy theo FSTS và bình phương của nó
với hiệu ứng cố định ngành và sai số chuẩn vững HC1, theo chuỗi phân
tầng M1--M5.

\hypertarget{quan-hux1ec7-cux1b0ux1eddng-ux111ux1ed9-duux1ed7i-thux1eb3ng-vux1ec1-gux1ea7n-tuyux1ebfn-tuxednh}{%
\subsection{4.9.1 Quan hệ cường độ duỗi thẳng về gần tuyến
tính}\label{quan-hux1ec7-cux1b0ux1eddng-ux111ux1ed9-duux1ed7i-thux1eb3ng-vux1ec1-gux1ea7n-tuyux1ebfn-tuxednh}}

Hiệu ứng tuyến tính của cường độ xuất khẩu lên năng suất \textbf{dương
và có ý nghĩa cao} ở mô hình nền (\(M1: \beta_1 = +0{,}671\);
\(p < 0{,}001\)). Khi bổ sung số hạng bậc hai (M2), số hạng tuyến tính
vẫn dương (\(+1{,}042\); \(p = 0{,}002\)) trong khi bậc hai âm nhưng
\textbf{xa ý nghĩa} (\(-0{,}606\); \(p = 0{,}323\)); điểm uốn đại số hàm
ý 90,1\% nằm sát rìa trên cùng của miền FSTS và không phân biệt được với
trạng thái ``không có điểm uốn''. Khi lần lượt thêm kiểm soát nhân khẩu
(M3), năng lực và chấp nhận số (M4), rồi tương tác uốn đường cong (M5),
\textbf{số hạng bậc hai không bao giờ tiệm cận ý nghĩa} và điểm uốn hàm
ý trôi tới các giá trị phi lý (108\% ở M3, vượt 700\% khi kiểm soát năng
lực), dấu hiệu điển hình của một hệ số không phân biệt được với 0. Theo
tiêu chí Lind--Mehlum (2010), \textbf{không một đặc tả nào xác nhận chữ
U ngược}: trong miền quan sát (FSTS trung bình 4,1\%, trung bình có điều
kiện 24,6\%), quốc tế hóa nâng hiệu quả một cách đơn điệu.

\textbf{Bảng 4.9.} \emph{Hiệu ứng cường độ xuất khẩu (FSTS) lên năng
suất qua chuỗi mô hình phân tầng, Nhật Bản 2025 (\(N = 2.168\); HC1).}

\begin{longtable}[]{@{}
  >{\raggedright\arraybackslash}p{(\columnwidth - 10\tabcolsep) * \real{0.1667}}
  >{\raggedleft\arraybackslash}p{(\columnwidth - 10\tabcolsep) * \real{0.1667}}
  >{\raggedleft\arraybackslash}p{(\columnwidth - 10\tabcolsep) * \real{0.1667}}
  >{\raggedleft\arraybackslash}p{(\columnwidth - 10\tabcolsep) * \real{0.1667}}
  >{\raggedleft\arraybackslash}p{(\columnwidth - 10\tabcolsep) * \real{0.1667}}
  >{\raggedleft\arraybackslash}p{(\columnwidth - 10\tabcolsep) * \real{0.1667}}@{}}
\toprule\noalign{}
\begin{minipage}[b]{\linewidth}\raggedright
Hệ số
\end{minipage} & \begin{minipage}[b]{\linewidth}\raggedleft
M1
\end{minipage} & \begin{minipage}[b]{\linewidth}\raggedleft
M2
\end{minipage} & \begin{minipage}[b]{\linewidth}\raggedleft
M3
\end{minipage} & \begin{minipage}[b]{\linewidth}\raggedleft
M4
\end{minipage} & \begin{minipage}[b]{\linewidth}\raggedleft
M5
\end{minipage} \\
\midrule\noalign{}
\endhead
\bottomrule\noalign{}
\endlastfoot
FSTS (định tâm) & +0,671*** & +1,042** & +0,872* & +0,476 & +0,486 \\
FSTS² & --- & −0,606 (ns) & (ns) & (ns) & (ns) \\
Điểm uốn đại số & --- & {[}90,1\%{]} & {[}108,3\%{]} & ngoài miền &
ngoài miền \\
Bổ sung kiểm soát & nền & +bậc hai & +nhân khẩu & +TCI/DAI & +tương
tác \\
\end{longtable}

\emph{Ghi chú: sai số chuẩn vững HC1. * \(p<0{,}05\); ** \(p<0{,}01\);
*** \(p<0{,}001\); (ns) không có ý nghĩa thống kê. Không đặc tả nào có
FSTS² âm có ý nghĩa nên không xác nhận chữ U ngược (Lind \& Mehlum,
2010). Độ vững với d3c riêng: \(\beta_1 = +1{,}255\)**;
\(\beta_2 = -0{,}674\) (\(p = 0{,}445\)), cùng kết luận. Nguồn: P10,
\texttt{p10\_japan/replication/}.}

Kết luận dương đơn điệu \textbf{vững qua sáu đặc tả độ vững}: thêm kiểm
soát quy mô (\(+0{,}656\); \(p = 0{,}046\)), giới hạn cơ sở nội địa
(\(+0{,}917\); \(p = 0{,}012\) --- loại trừ khả năng số ít doanh nghiệp
vốn nước ngoài dẫn dắt), và winsorize năng suất tại phân vị 5/95
(\(+0{,}837\); \(p = 0{,}006\)); ở mọi đặc tả số hạng bậc hai vô nghĩa.
Riêng đặc tả chỉ trên doanh nghiệp xuất khẩu (\(N = 344\)), hiệu ứng
cường độ dương nhưng không còn ý nghĩa (\(+0{,}200\); \(p = 0{,}410\)),
nhất quán với việc trọng tâm nằm ở biên rộng chứ không phải cường độ.
\textbf{H1 được ủng hộ} --- Nhật Bản nằm trên một nhánh đi lên dài mà
nhánh đi xuống vắng mặt về mặt thực nghiệm, trùng khớp kết quả Singapore
(Mục 4.3) và kết quả cấp nhóm Advanced của ước lượng 50 nền (Mục 4.6.2),
hoàn tất bức tranh ba vùng của gradient thể chế: chữ U ngược ở các chế
độ chuyển đổi, gần tuyến tính ở biên trên, tan rã ở biên dưới.

\hypertarget{nux103ng-lux1ef1c-sux1ed1-huxf3a-vuxe0-ux111ux1eb7c-ux111iux1ec3m-quux1ea3n-trux1ecb-nuxe2ng-mux1eb7t-bux1eb1ng-khuxf4ng-uux1ed1n-ux111ux1b0ux1eddng-cong}{%
\subsection{4.9.2 Năng lực, số hóa và đặc điểm quản trị: nâng mặt bằng,
không uốn đường
cong}\label{nux103ng-lux1ef1c-sux1ed1-huxf3a-vuxe0-ux111ux1eb7c-ux111iux1ec3m-quux1ea3n-trux1ecb-nuxe2ng-mux1eb7t-bux1eb1ng-khuxf4ng-uux1ed1n-ux111ux1b0ux1eddng-cong}}

Năng lực công nghệ có \textbf{hiệu ứng mặt bằng chính xác}
(\(+0{,}100\); \(p < 0{,}001\)) và hấp thụ phần lớn hiệu ứng FSTS thô
khi được đưa vào M4, nhất quán với chọn lọc theo năng lực vào hoạt động
xuất khẩu; tương tác uốn đường cong không có ý nghĩa (M5: \(-0{,}126\);
\(p = 0{,}328\)) --- năng lực nâng sàn, không uốn đường cong (\textbf{H2
được ủng hộ}). Chấp nhận số giữ phần bù dương khiêm tốn (\(+0{,}110\);
\(p = 0{,}028\)) dù tỷ lệ website đã bão hòa ở 83,8\%, không có tương
tác uốn đường cong; phần bù bằng khoảng một nửa mức của ước lượng gộp 50
nền, đúng mức nén mà trạng thái \textbf{bão hòa Tầng 1} dự đoán
(\textbf{H3 được ủng hộ}). Hai quy luật đặc thù Nhật Bản kỷ luật hóa
cách đọc hiệu ứng thượng tầng quản trị: phần bù \textbf{trường thọ doanh
nghiệp} dương (\(+0{,}073\); \(p = 0{,}007\)), mỗi đơn vị log tuổi thêm
khoảng 7\% năng suất, cộng hưởng với truyền thống \emph{shinise} ở nền
kinh tế có cơ sở lâu đời nhất mẫu gộp (tuổi trung bình 50,4 năm)
(\textbf{H5 được ủng hộ}); và hệ số \textbf{âm} của nữ quản lý cấp cao
(\(-0{,}170\); \(p = 0{,}010\)) trong nền kinh tế có biên lãnh đạo nữ
mỏng nhất khu vực (7,3\%), được đọc thận trọng như chọn lọc số ít cơ sở
do nữ điều hành vào các ngách biên thấp chứ \textbf{không diễn giải nhân
quả}.

\hypertarget{dux1ecbch-chuyux1ec3n-trux1ecdng-tuxe2m-sang-biuxean-rux1ed9ng-ai-xuux1ea5t-khux1ea9u}{%
\subsection{4.9.3 Dịch chuyển trọng tâm sang biên rộng: ai xuất
khẩu}\label{dux1ecbch-chuyux1ec3n-trux1ecdng-tuxe2m-sang-biuxean-rux1ed9ng-ai-xuux1ea5t-khux1ea9u}}

Vì đường cường độ phẳng trong miền quan sát, biến thiên có ý nghĩa nằm ở
\textbf{ai xuất khẩu}. Mô hình xác suất tuyến tính của tham gia xuất
khẩu cho thấy tham gia được chi phối bởi \textbf{sở hữu nước ngoài}
(\(+0{,}408\); \(p < 0{,}001\) --- một cơ sở có ít nhất 10\% vốn nước
ngoài có xác suất xuất khẩu cao hơn khoảng bốn mươi điểm phần trăm),
\textbf{năng lực công nghệ} (\(+0{,}083\) trên mỗi độ lệch chuẩn;
\(p < 0{,}001\)) và \textbf{chấp nhận số} (\(+0{,}046\);
\(p = 0{,}005\)); mô hình giải thích \(R^2 = 0{,}228\) biến thiên cắt
ngang của tham gia. Doanh nghiệp xuất khẩu hưởng \textbf{phần bù năng
suất 25,8\%} so với doanh nghiệp không xuất khẩu (\(p < 0{,}001\), kiểm
soát sở hữu, tuổi, giới tính và ngành). Đây chính là cơ chế \textbf{tự
chọn lọc Melitz (2003)} ở một nền kinh tế biên giới: tại biên trên của
gradient thể chế, câu hỏi ràng buộc không phải doanh nghiệp xuất khẩu
\emph{bao nhiêu} mà là \emph{có} xuất khẩu hay không --- lợi ích năng
suất của quốc tế hóa vận hành chủ yếu qua ngưỡng tham gia, và các doanh
nghiệp vượt ngưỡng là những doanh nghiệp được trang bị năng lực, công
nghệ, số và gắn với sở hữu nước ngoài (\textbf{H4 được ủng hộ}). Bằng
chứng Nhật Bản do đó vừa hoàn tất phổ ba vùng của khung CDCM ở biên
trên, vừa làm rõ rằng cơ chế quốc tế hóa và hiệu quả chuyển từ biên
cường độ (ở các chế độ chuyển đổi) sang biên rộng (ở chế độ tiên tiến
trưởng thành).

\hypertarget{tam-giuxe1c-huxf3a-kux1ebft-quux1ea3-cux1ea5p-doanh-nghiux1ec7p-vux1edbi-bux1eb1ng-chux1ee9ng-vux129-muxf4-ux111ux1ed9c-lux1eadp}{%
\section{4.10 Tam giác hóa kết quả cấp doanh nghiệp với bằng chứng vĩ mô
độc
lập}\label{tam-giuxe1c-huxf3a-kux1ebft-quux1ea3-cux1ea5p-doanh-nghiux1ec7p-vux1edbi-bux1eb1ng-chux1ee9ng-vux129-muxf4-ux111ux1ed9c-lux1eadp}}

Một thử thách then chốt đối với tính giá trị của các phát hiện cấp doanh
nghiệp là liệu chúng có nhất quán với bằng chứng vĩ mô độc lập hay
không. Mục này thực hiện tam giác hóa (triangulation) bốn kết quả trụ
cột của luận án với các báo cáo chuyên đề mới nhất của Ngân hàng Thế
giới (2026), qua đó củng cố tính giá trị hội tụ (convergent validity)
của khung ICRV.

\textbf{Thứ nhất, gradient thể chế của điểm uốn (H5) hội tụ với quy luật
vĩ mô ``thể chế tốt hơn đi liền tăng trưởng nhanh hơn''.} Luận án phát
hiện chế độ thể chế quyết định chính sự tồn tại và vị trí của chữ U
ngược: đường cong định hình rõ nét nhất ở vùng giữa gradient thể chế
(Nhóm IV, điểm uốn 43,0\%), duỗi thẳng gần tuyến tính ở biên trên và mất
cấu trúc ở biên dưới (Mục 4.6.2). Ở cấp vĩ mô, Ngân hàng Thế giới ghi
nhận rằng các thị trường biên tăng trưởng nhanh nhất là những nền đã cải
thiện quản trị, thu hẹp chênh lệch lãi suất ngân hàng và kiểm soát nợ
công, trong khi các điểm yếu thể chế đi kèm tăng trưởng vốn đáng thất
vọng (World Bank, 2026c). Hai tầng bằng chứng cùng chỉ về một cơ chế:
chất lượng thể chế quyết định mức độ và điều kiện mà tích lũy vốn và
quốc tế hóa chuyển hóa thành kết quả thực. Gradient điểm uốn cấp doanh
nghiệp của luận án vì vậy có thể đọc như biểu hiện vi mô của quy luật
phân hóa tăng trưởng vĩ mô.

\textbf{Thứ hai, hai cơ chế FDI tương phản giữa nhóm tài nguyên và nhóm
đổi mới ăn khớp với cú sốc hàng hóa 2026.} Luận án cho thấy nhóm
Advanced tài nguyên dẫn dắt (Nhóm II, các nền GCC) có phân tán năng suất
hẹp đặc trưng của kinh tế nhà nước tô, nơi phân phối lại tô lợi tài
nguyên san bằng khác biệt năng suất giữa doanh nghiệp (Mục 4.1, Chuyên
đề 1). Cú sốc xung đột Trung Đông năm 2026, với giá năng lượng dự báo
tăng 24\% và giá hàng hóa nói chung tăng 16\%, đẩy giá kim loại cơ bản
và quý lên mức cao kỷ lục (World Bank, 2026d), củng cố vị thế tô lợi của
nhóm này trong ngắn hạn nhưng đồng thời khắc sâu rủi ro phụ thuộc tài
nguyên mà mô hình tăng trưởng dựa trên đổi mới (Nhóm I) không gặp phải.
Bối cảnh hàng hóa 2026 vì vậy làm rõ hơn vì sao việc gộp hai loại
Advanced thành một nhóm duy nhất, như tài liệu trước thường làm, che
khuất hai cơ chế chuyển hóa quốc tế hóa hoàn toàn khác nhau.

\textbf{Thứ ba, Nghịch lý bão hòa số và vai trò DAI ăn khớp với chương
trình nghị sự về thể chế công trong kỷ nguyên AI.} Luận án cho thấy phần
thưởng của hiện diện số Tầng 1 co lại ở các nền kinh tế tiên tiến nơi nó
đã tiệm cận bão hòa, trong khi vẫn còn dư địa ở các nền đang chuyển đổi
(Mục 4.3.2, Mục 4.6.3). Phân tích của Ngân hàng Thế giới về thể chế công
trong kỷ nguyên trí tuệ nhân tạo nhấn mạnh rằng hạ tầng số công và quản
trị dữ liệu là điều kiện nền tảng quyết định năng lực số được chuyển hóa
thành giá trị, và rằng giai đoạn trưởng thành số khác nhau đòi hỏi công
cụ chính sách khác nhau (World Bank, 2026e). Sự tương đồng này củng cố
luận điểm cốt lõi của khung CDCM: phần thưởng của áp dụng số là hàm của
mức độ phát triển thể chế và tầng năng lực số được đo, chứ không phải
một ``phần bù số'' đồng nhất.

\textbf{Thứ tư, FIP và thách thức việc làm của các nền kinh tế biên định
vị đóng góp chính sách của luận án.} Phát hiện FIP tại SIDS (Mục 4.7) và
việc xác nhận rằng quốc tế hóa chỉ tạo phần thưởng năng suất trong điều
kiện thể chế và năng lực đủ mạnh, đặt trong bối cảnh hơn 230 triệu người
sẽ bước vào tuổi lao động ở các thị trường biên giai đoạn 2025--2035
(World Bank, 2026c), cho thấy chính sách thúc đẩy xuất khẩu đơn thuần là
chưa đủ và có thể phản tác dụng ở những bối cảnh dễ tổn thương. Hàm ý
chính sách của luận án, ưu tiên nâng cấp năng lực công nghệ trước khi mở
rộng cường độ quốc tế hóa, vì vậy nhất quán với khuyến nghị vĩ mô của
Ngân hàng Thế giới về việc đầu tư vào hạ tầng nền tảng và năng lực hấp
thụ như điều kiện để các nền kinh tế biên hiện thực hóa tiềm năng (World
Bank, 2026c).

\textbf{Thứ năm, cú sốc thương mại Việt Nam đầu năm 2026 minh họa trực
tiếp tính hệ quả của nhánh đi xuống đối với nhóm chuyển đổi.} Dữ liệu
giám sát vĩ mô của Ngân hàng Thế giới (2026) cho Việt Nam, nền kinh tế
Nhóm IV lớn nhất trong mẫu, ghi nhận chỉ số nhà quản trị mua hàng ngành
chế tạo rơi xuống 45,3 điểm (vùng thu hẹp) vào tháng 3 năm 2026 khi đơn
hàng xuất khẩu mới sụt giảm sau khi một cuộc điều tra của cơ quan đại
diện thương mại Hoa Kỳ được khởi động; nhập khẩu hàng hóa đã tăng 27,8\%
so với cùng kỳ năm trước tính đến cùng tháng, một phần do đứt gãy chuỗi
cung ứng năng lượng từ xung đột tại eo biển Hormuz. Đây chính là điều
kiện mà vế phải của đường cong chữ U ngược trở nên hệ quả: các doanh
nghiệp bị đẩy vượt ngưỡng cường độ xuất khẩu tối ưu bởi cam kết xuất
khẩu trước đó phải gánh áp lực chi phí điều phối và chi phí đầu vào
khuếch đại. Vì Việt Nam nằm trong chế độ chuyển đổi nơi chữ U ngược sắc
nét nhất (điểm uốn Nhóm IV khoảng 43\%), các doanh nghiệp xuất khẩu Việt
Nam hiện ở trên ngưỡng này thuộc nhóm chịu phơi nhiễm cao nhất trước các
cú sốc vĩ mô. Bằng chứng năng lực hàm ý rằng xây dựng năng lực công nghệ
và số nâng sàn năng suất cho các doanh nghiệp này, dù trên ước lượng gộp
năng lực không tự nó làm thoải nhánh đi xuống; ràng buộc quyết định là
độ rộng của tham gia xuất khẩu có năng lực, không phải một biện pháp
dịch chuyển đường cong.

Năm điểm hội tụ trên cho thấy các phát hiện cấp doanh nghiệp của luận án
không đứng tách biệt mà ăn khớp một cách hệ thống với bằng chứng vĩ mô
độc lập gần nhất. Sự tam giác hóa này tăng cường tính giá trị của khung
ICRV và đặt đóng góp của luận án vào đúng dòng chảy nghiên cứu phát
triển đương đại.

\begin{center}\rule{0.5\linewidth}{0.5pt}\end{center}

\hypertarget{tux1ed5ng-hux1ee3p-kux1ebft-quux1ea3-theo-hux1ec7-giux1ea3-thuyux1ebft}{%
\section{4.11 Tổng hợp kết quả theo hệ giả
thuyết}\label{tux1ed5ng-hux1ee3p-kux1ebft-quux1ea3-theo-hux1ec7-giux1ea3-thuyux1ebft}}

Trước khi tổng hợp theo hệ giả thuyết, Bảng 4.10 hợp nhất các ước lượng
trụ cột của chín nghiên cứu thành phần thực nghiệm, cho phép đối chiếu
trực tiếp dạng hàm, điểm uốn, vai trò của hai cấu trúc số (TCI, DAI) và
quy mô mẫu xuyên các bối cảnh ICRV.

\textbf{Bảng 4.10.} \emph{Tổng hợp ước lượng các nghiên cứu thành phần
thực nghiệm.}

\begin{longtable}[]{@{}
  >{\raggedright\arraybackslash}p{(\columnwidth - 12\tabcolsep) * \real{0.1429}}
  >{\raggedright\arraybackslash}p{(\columnwidth - 12\tabcolsep) * \real{0.1429}}
  >{\raggedleft\arraybackslash}p{(\columnwidth - 12\tabcolsep) * \real{0.1429}}
  >{\raggedright\arraybackslash}p{(\columnwidth - 12\tabcolsep) * \real{0.1429}}
  >{\raggedright\arraybackslash}p{(\columnwidth - 12\tabcolsep) * \real{0.1429}}
  >{\raggedright\arraybackslash}p{(\columnwidth - 12\tabcolsep) * \real{0.1429}}
  >{\raggedright\arraybackslash}p{(\columnwidth - 12\tabcolsep) * \real{0.1429}}@{}}
\toprule\noalign{}
\begin{minipage}[b]{\linewidth}\raggedright
NC
\end{minipage} & \begin{minipage}[b]{\linewidth}\raggedright
Bối cảnh (nhóm ICRV)
\end{minipage} & \begin{minipage}[b]{\linewidth}\raggedleft
N
\end{minipage} & \begin{minipage}[b]{\linewidth}\raggedright
Dạng hàm
\end{minipage} & \begin{minipage}[b]{\linewidth}\raggedright
Điểm uốn (FSTS)
\end{minipage} & \begin{minipage}[b]{\linewidth}\raggedright
Vai trò TCI
\end{minipage} & \begin{minipage}[b]{\linewidth}\raggedright
Vai trò DAI
\end{minipage} \\
\midrule\noalign{}
\endhead
\bottomrule\noalign{}
\endlastfoot
P2 & Trung Quốc, doanh nghiệp nhỏ và vừa (III) & 4.290 & U ngược &
47,8\% & nâng mặt bằng & nâng mặt bằng \\
P3 & Việt Nam (IV) & 2.958 & U ngược & 39--46\% & dương, có ý nghĩa
(β=0,179, biến công cụ) & không có ý nghĩa (Tầng 1) \\
P4 & Singapore (I) & 623 & gần tuyến tính dương & \textasciitilde88,6\%
(không định vị chắc) & yếu tố vệ sinh & FSTS²×DAI=+3,119** \\
P5 & Trung Quốc (III) & 4.544 & U ngược & 47--49\% (ổn định 2012--2024)
& dương (+0,26 lên +0,43) & không áp dụng \\
P6 & Phân tích tổng hợp (châu Á) & k=238 & dương nhỏ & không áp dụng &
không áp dụng & r̄=0,074; I²=62,4\% \\
P7 & Đa quốc gia (50 nền, I--VI) & 81.022 & U ngược & 43,6\% (M5); rõ
nhất ở Nhóm IV (43,0\%) & nâng mặt bằng (+0,141 sd) & nâng mặt bằng
(+0,201 sd) \\
P8 & Đảo nhỏ Thái Bình Dương (VI) & 959 & âm đơn điệu (FIP) & không có &
nâng mặt bằng, không điều tiết & β(FSTS)=−1,339*** \\
P9 & Ấn Độ (IV) & 28.717 & U ngược chuyển sang sụp đổ ngưỡng & 61,8\%
rồi 40,7\% rồi tan biến & không áp dụng & tương tác Tầng 2 âm \\
P10 & Nhật Bản (I) & 2.168 & gần tuyến tính dương &
\textasciitilde90,1\% (không định vị) & nâng mặt bằng & nâng mặt bằng \\
\end{longtable}

\emph{Nguồn: tổng hợp từ các nghiên cứu thành phần P2--P10. ** p
\textless,01; *** p \textless,001. Điểm uốn xác nhận bằng kiểm định
Lind--Mehlum (2010). Điểm uốn P7 = 43,6\% ở mô hình đầy đủ M5 trên khung
50 nền (dải 43,6--51,5\% giữa các đặc tả).}

Bảng 4.10 làm nổi bật ba quy luật xuyên suốt. Thứ nhất, \textbf{dạng hàm
chữ U ngược là quy luật chứ không phải ngoại lệ}, được xác nhận ở năm
trong chín nghiên cứu (P2, P3, P5, P7 và một phần P9), với các trường
hợp biên có ý nghĩa lý thuyết riêng: quan hệ gần tuyến tính dương ở biên
thể chế mạnh (P4 Singapore và P10 Nhật Bản, điểm uốn nằm ngoài miền quan
sát), quan hệ âm đơn điệu tại các nền đảo nhỏ (P8, FIP), và sự tan biến
ngưỡng dưới biến động thể chế cực đoan tại Ấn Độ (P9). Thứ hai,
\textbf{điểm uốn dịch theo chiều thể chế}: từ khoảng 88--90\% ở
Singapore và Nhật Bản (Nhóm I) xuống 47--49\% ở Trung Quốc (III) và
39--46\% ở Việt Nam (IV), nhất quán với gradient toàn mẫu của P7. Thứ
ba, \textbf{TCI và DAI vận hành khác nhau}: TCI nhất quán nâng mặt bằng
năng suất, trong khi vai trò DAI phụ thuộc bối cảnh (nén đường cong ở
thể chế yếu, khuếch đại ở Singapore), củng cố sự phân biệt cốt lõi của
khung CDCM.

\begin{longtable}[]{@{}
  >{\raggedright\arraybackslash}p{(\columnwidth - 4\tabcolsep) * \real{0.3333}}
  >{\raggedright\arraybackslash}p{(\columnwidth - 4\tabcolsep) * \real{0.3333}}
  >{\raggedright\arraybackslash}p{(\columnwidth - 4\tabcolsep) * \real{0.3333}}@{}}
\toprule\noalign{}
\begin{minipage}[b]{\linewidth}\raggedright
Giả thuyết
\end{minipage} & \begin{minipage}[b]{\linewidth}\raggedright
Nội dung
\end{minipage} & \begin{minipage}[b]{\linewidth}\raggedright
Kết quả
\end{minipage} \\
\midrule\noalign{}
\endhead
\bottomrule\noalign{}
\endlastfoot
H1 (Phi tuyến chữ U ngược) & quốc tế hóa và hiệu quả có dạng chữ U ngược
& Xác nhận tại Trung Quốc, Việt Nam và toàn mẫu P7; xác nhận một phần
tại Singapore và Nhật Bản (gần tuyến tính, điểm uốn ngoài miền); không
áp dụng tại SIDS \\
H2 (TCI điều tiết dương) & TCI khuếch đại quốc tế hóa và hiệu quả & Xác
nhận một phần: TCI nâng mặt bằng năng suất tại tất cả bối cảnh
(+0,141***, 50 nền); uốn đường cong xác nhận tại Việt Nam nhưng không ý
nghĩa tại toàn mẫu P7 \\
H3 (DAI điều tiết) & DAI thay đổi độ dốc quốc tế hóa và hiệu quả & Hỗ
trợ một phần: hiệu ứng mức phổ quát (+0,201*\textbf{, 50 nền); tương tác
đường cong cùng chiều nén nhưng không ý nghĩa toàn mẫu; bằng chứng điều
tiết rõ nhất ở bối cảnh đơn quốc gia (P4: FSTS²×DAI = +3,119}); không có
ý nghĩa tại Việt Nam (Tầng 1) tương thích với hạn chế đo lường \\
H4 (Nhà quản trị điều tiết) & Kinh nghiệm/giới tính nhà quản trị & Không
được hỗ trợ trên toàn mẫu 50 nền (nữ quản lý: mức −0,104,
p\textless,001; tương tác FSTS không ý nghĩa); hiệu ứng phụ thuộc bối
cảnh \\
H5 (ICRV gradient) & Thể chế điều tiết gradient quốc tế hóa và hiệu quả
& Xác nhận một phần: \(Q_M = 17{,}35^{**}\) (\emph{df} = 4, \emph{p}
=,002) từ P6 MARA xác nhận dị biệt giữa chế độ; gradient định hướng
(E1a/E1b) chưa xác nhận do \emph{k} = 3 tại nhóm Frontier (bin FR, hệ
5-regime của P6); gradient điểm uốn theo ICRV được xác nhận từ P7 \\
H6 (Temporal heterogeneity) & Hình dạng quốc tế hóa và hiệu quả thay đổi
theo thời gian & Không xác nhận tại Trung Quốc (\(F = 1{,}83\),
\(p = {,}176\)); TP dịch chuyển tại Việt Nam giữa 2009 và 2015 \\
H1b (FIP, SIDS boundary condition) & Quan hệ âm đơn điệu tại SIDS & Xác
nhận: \(\beta = -1{,}339^{***}\) (M1, country+year FE) đến
\(-3{,}351^{***}\) (year-FE); chỉ doanh nghiệp xuất khẩu NS do mẫu nhỏ
(N=26) \\
H1c (Tan biến ngưỡng tại Ấn Độ, P9) & Ngưỡng U ngược \emph{tan biến}
dưới biến động thể chế cực đoan 2014--2025 & Xác nhận tại Ấn Độ: TP
61,8\% (2014) rồi 40,7\% (2022) rồi tan biến (2025, \(\beta_2\) NS);
Paternoster \(z = -7{,}94\), \(p < {,}0001\); tương tác DAI Tầng 2 (UPI)
âm (\(-4{,}02\), \(p = {,}004\)) \\
\end{longtable}

\textbf{Đối chiếu với ba phản bác chính.} Để khẳng định độ tin cậy của
các phát hiện trụ cột, cần xử lý tường minh ba phản bác hiển nhiên nhất.
Thứ nhất là nhân quả ngược và nội sinh: doanh nghiệp năng suất cao có
thể tự chọn lọc vào xuất khẩu, làm thổi phồng hệ số nhánh đi lên. Tuy
nhiên, định vị điểm uốn phụ thuộc vào dấu và độ lớn của hệ số bậc hai,
vốn ít nhạy với chiều của chọn lọc hơn hệ số bậc nhất; thêm vào đó, chữ
U ngược tái lập qua nhiều đợt khảo sát và nhiều bối cảnh có môi trường
chọn lọc khác nhau, và nghiên cứu Việt Nam (P3) dùng biến công cụ cho hệ
số TCI, củng cố cách đọc rằng phát hiện không phải tạo tác của nội sinh.
Thứ hai là thiên lệch chọn lọc do dùng mẫu chỉ doanh nghiệp xuất khẩu:
phản bác này được xử lý bằng việc ước lượng trên toàn mẫu gồm cả doanh
nghiệp có cường độ xuất khẩu bằng không, bằng phân rã hai biên tham gia
và cường độ, và bằng độ vững khi loại doanh nghiệp siêu nhỏ. Thứ ba là
tạo tác đo lường, nhất là nhiễu tiền tệ: việc trung tâm hóa cường độ
xuất khẩu, cắt đuôi phân vị trong từng cặp nền kinh tế và năm, và áp
dụng hiệu ứng cố định hai chiều đã hấp thụ chênh lệch mặt bằng giá; quan
trọng hơn, cấu trúc ba vùng được xác nhận độc lập bởi hai phương pháp có
điểm yếu không trùng nhau, là phân tích tổng hợp P6 và ước lượng gộp
toàn mẫu P7. Sự hội tụ này khiến ba phản bác trên khó có thể đồng thời
giải thích toàn bộ mô thức quan sát được.

\begin{center}\rule{0.5\linewidth}{0.5pt}\end{center}

\hypertarget{tuxe0i-liux1ec7u-tham-khux1ea3o-truxedch-dux1eabn-trong-chux1b0ux1a1ng-4}{%
\section{Tài liệu tham khảo (trích dẫn trong Chương
4)}\label{tuxe0i-liux1ec7u-tham-khux1ea3o-truxedch-dux1eabn-trong-chux1b0ux1a1ng-4}}

Arte, P., \& Larimo, J. (2022). Moderating influence of product
diversification on the international diversification-performance
relationship: A meta-analysis. \emph{Journal of Business Research}, 139,
281--295.

Bausch, A., \& Krist, M. (2007). The effect of context-related
moderators on the internationalization-performance relationship:
Evidence from meta-analysis. \emph{Management International Review,
47}(3), 319--347.

Begg, C. B., \& Mazumdar, M. (1994). Operating characteristics of a rank
correlation test for publication bias. \emph{Biometrics, 50}(4),
1088--1101.

Cameron, A. C., Gelbach, J. B., \& Miller, D. L. (2008). Bootstrap-based
improvements for inference with clustered errors. \emph{Review of
Economics and Statistics, 90}(3), 414--427.

Cohen, J. (1988). \emph{Statistical power analysis for the behavioral
sciences} (2nd ed.). Lawrence Erlbaum.

Cohen, W. M., \& Levinthal, D. A. (1990). Absorptive capacity: A new
perspective on learning and innovation. \emph{Administrative Science
Quarterly, 35}(1), 128--152.

Do, T. H., \& Phan, A. T. (2024, December). \emph{Internationalization
and firm performance: A meta-analysis review} {[}Paper presentation{]}.
The Sixth International Conference on Sustainable Development in
Economics, Business, and Finance (ICBEF), Can Tho University, Vietnam.

Do, T. H., \& Phan, A. T. (2025). Internationalization and firm
performance of firms in India: The role of top management. In M.
Bartekova (Ed.), \emph{International business research: Traditional and
creative approaches}. IntechOpen.
\url{https://doi.org/10.5772/intechopen.1011012}

Haans, R. F. J., Pieters, C., \& He, Z.-L. (2016). Thinking about U:
Theorizing and testing U- and inverted U-shaped relationships in
strategy research. \emph{Strategic Management Journal, 37}(7),
1177--1195.

Helpman, E., Melitz, M. J., \& Yeaple, S. R. (2004). Export versus FDI
with heterogeneous firms. \emph{American Economic Review, 94}(1),
300--316.

Hennart, J.-F. (2007). The theoretical rationale for a
multinationality--performance relationship. \emph{Management
International Review, 47}(3), 423--452.

Hsieh, C.-T., \& Klenow, P. J. (2009). Misallocation and manufacturing
TFP in China and India. \emph{Quarterly Journal of Economics, 124}(4),
1403--1448.

Johanson, J., \& Vahlne, J.-E. (1977). The internationalization process
of the firm --- A model of knowledge development and increasing foreign
market commitments. \emph{Journal of International Business Studies,
8}(1), 23--32.

Kaufmann, D., Kraay, A., \& Mastruzzi, M. (2011). The Worldwide
Governance Indicators: Methodology and analytical issues. \emph{Hague
Journal on the Rule of Law, 3}(2), 220--246.

Khanna, T., \& Palepu, K. G. (2010). \emph{Winning in thị trường mới
nổi: A road map for strategy and execution}. Harvard Business Press.

Kirca, A. H., Roth, K., Hult, G. T. M., \& Cavusgil, S. T. (2012). The
role of context in the multinationality-performance relationship: A
meta-analytic review. \emph{Global Strategy Journal, 2}(2), 108--121.

Lall, S. (1992). Technological capabilities and industrialization.
\emph{World Development, 20}(2), 165--186.

Lind, J. T., \& Mehlum, H. (2010). With or without U? The appropriate
test for a U-shaped relationship. \emph{Oxford Bulletin of Economics and
Statistics, 72}(1), 109--118.

Marano, V., Arregle, J.-L., Hitt, M. A., Spadafora, E., \& Van Essen, M.
(2016). Home country institutions and the
internationalization-performance relationship: A meta-analytic review.
\emph{Journal of Management, 42}(5), 1075--1110.

Melitz, M. J. (2003). The impact of trade on intra-industry
reallocations and aggregate industry productivity. \emph{Econometrica,
71}(6), 1695--1725.

Mizik, N., \& Jacobson, R. (2003). Trading off between value creation
and value appropriation: The financial implications of shifts in
strategic emphasis. \emph{Journal of Marketing, 67}(1), 63--76.

North, D. C. (1990). \emph{Institutions, institutional change and
economic performance}. Cambridge University Press.

Oster, E. (2019). Unobservable selection and coefficient stability:
Theory and evidence. \emph{Journal of Business \& Economic Statistics,
37}(2), 187--204.

Paternoster, R., Brame, R., Mazerolle, P., \& Piquero, A. (1998). Using
the correct statistical test for the equality of regression
coefficients. \emph{Criminology, 36}(4), 859--866.

Pisani, N., Garcia-Bernardo, J., \& Heemskerk, E. (2020). Does it pay to
be a multinational? A large-sample, cross-national replication assessing
the multinationality--performance relationship. \emph{Strategic
Management Journal, 41}(1), 152--172.

Verhoef, P. C., Broekhuizen, T., Bart, Y., Bhattacharya, A., Dong, J.
Q., Fabian, N., \& Haenlein, M. (2021). Digital transformation: A
multidisciplinary reflection and research agenda. \emph{Journal of
Business Research}, 122, 889--901.

Wagner, J. (2007). Exports and productivity: A survey of the evidence
from firm-level data. \emph{The World Economy, 30}(1), 60--82.

Williamson, O. E. (1985). \emph{The economic institutions of
capitalism}. Free Press.

World Bank. (2025b, September). \emph{Viet Nam economic update:
Nurturing high-tech talents} (Taking Stock, September 2025). World Bank
Group.

World Bank Enterprise Surveys. (2025). \emph{Enterprise surveys data}.
World Bank Group. \url{https://www.enterprisesurveys.org}

Wu, J., Xu, L., \& Yang, Z. (2022). Multinationality and performance of
emerging economy multinational enterprises: A twenty-year study.
\emph{International Business Review, 31}(5), 102003.

\begin{center}\rule{0.5\linewidth}{0.5pt}\end{center}

\newpage

\hypertarget{chux1b0ux1a1ng-5-kux1ebft-luux1eadn-vuxe0-ux111ux1ec1-xuux1ea5t}{%
\chapter{CHƯƠNG 5: KẾT LUẬN VÀ ĐỀ
XUẤT}\label{chux1b0ux1a1ng-5-kux1ebft-luux1eadn-vuxe0-ux111ux1ec1-xuux1ea5t}}

\begin{center}\rule{0.5\linewidth}{0.5pt}\end{center}

\hypertarget{giux1edbi-thiux1ec7u-chux1b0ux1a1ng-1}{%
\section{Giới thiệu
chương}\label{giux1edbi-thiux1ec7u-chux1b0ux1a1ng-1}}

Chương 5 tổng hợp và diễn giải kết quả từ Chương 4 trong ánh sáng của
khung lý thuyết CDCM và bốn lý thuyết nền (Uppsala, RBV, Institutional
Theory, Upper Echelons Theory). Cấu trúc chương gồm sáu phần: (i) bàn
luận kết quả theo khung lý thuyết; (ii) so sánh với nghiên cứu trước;
(iii) đóng góp lý thuyết; (iv) đề xuất chính sách và quản trị; (v) hạn
chế và hướng nghiên cứu tương lai; và (vi) kết luận tổng thể của luận
án.

\begin{center}\rule{0.5\linewidth}{0.5pt}\end{center}

\hypertarget{buxe0n-luux1eadn-kux1ebft-quux1ea3-theo-khung-luxfd-thuyux1ebft-cdcm}{%
\section{5.1 Bàn luận kết quả theo khung lý thuyết
CDCM}\label{buxe0n-luux1eadn-kux1ebft-quux1ea3-theo-khung-luxfd-thuyux1ebft-cdcm}}

\hypertarget{icrv-gradient-xuxe1c-nhux1eadn-h5-vuxe0-tux1ea7ng-thux1ec3-chux1ebf-cux1ee7a-cdcm}{%
\subsection{5.1.1 ICRV Gradient: Xác nhận H5 và tầng thể chế của
CDCM}\label{icrv-gradient-xuxe1c-nhux1eadn-h5-vuxe0-tux1ea7ng-thux1ec3-chux1ebf-cux1ee7a-cdcm}}

Phát hiện nhất quán nhất và có ý nghĩa lý thuyết rộng nhất của luận án
là \textbf{gradient ICRV (Biến thiên chế độ bối cảnh thể chế,
Institutional Context Regime Variation)}: cùng một mức độ quốc tế hóa
(FSTS) mang lại hiệu quả rất khác nhau tùy thuộc vào chế độ thể chế mà
doanh nghiệp hoạt động. Bằng chứng từ hai nguồn độc lập, phân tích tổng
hợp (P6, \(Q_M = 17{,}35\), \(df = 4\), \(p = {,}002\)) và ước lượng
toàn mẫu trên khung 50 nền kinh tế (P7, \(N = 81.022\) quan sát), đều
chỉ đến cùng một kết luận: chế độ thể chế quyết định \textbf{chính sự
tồn tại và vị trí của dạng hàm chữ U ngược}. Đường cong định hình rõ nét
và tin cậy nhất ở vùng giữa của gradient thể chế (Nhóm IV, điểm uốn
43,0\%, trùng khớp các ước lượng đơn quốc gia Việt Nam 39--46\% và Ấn Độ
40,7\%); duỗi thẳng thành quan hệ gần tuyến tính dương ở biên trên (Nhóm
I, điểm uốn ngoài miền dữ liệu, trùng khớp Singapore); và mất cấu trúc ở
biên dưới (Nhóm V không còn ý nghĩa thống kê; Nhóm VI không có chữ U
ngược, chuyển sang gánh nặng quốc tế hóa bắt buộc theo P8). Hay nói cách
khác, thể chế không chỉ dịch chuyển điểm tối ưu của quốc tế hóa mà quyết
định liệu một điểm tối ưu như vậy có tồn tại hay không.

Kết quả này xác nhận H5 và tầng thể chế của khung lý thuyết CDCM theo
một cách thuyết phục nhất: không phải chỉ một biến điều tiết đơn lẻ mà
là \textbf{cả một chế độ thể chế}, được vận hành hóa thành sáu chế độ
con ICRV, xác định điều kiện biên cho quan hệ quốc tế hóa và hiệu quả.
lý thuyết thể chế của North (1990) và Khanna và Palepu (2010) dự đoán
rằng thể chế định hình chi phí giao dịch và cấu trúc khuyến khích; kết
quả ICRV gradient là biểu hiện thực nghiệm của cơ chế này ở quy mô 50
nền kinh tế châu Á và Thái Bình Dương (P7).

Về mặt lý thuyết, gradient ICRV có thể được hiểu qua hai cơ chế bổ sung
cho nhau. Cơ chế thứ nhất là \textbf{giảm chi phí giao dịch}: thể chế
chất lượng cao (hệ thống pháp lý, bảo vệ quyền sở hữu trí tuệ, thực thi
hợp đồng) làm giảm chi phí xuyên biên giới, cho phép doanh nghiệp duy
trì đà tăng trưởng hiệu quả đến mức FSTS cao hơn. Cơ chế thứ hai là
\textbf{khuếch đại năng lực}: môi trường thể chế tốt cũng có nghĩa là hệ
sinh thái dịch vụ hỗ trợ (logistics, tài chính thương mại, tư vấn pháp
lý quốc tế) phát triển, làm giảm gánh nặng điều phối nội bộ của doanh
nghiệp khi mở rộng quốc tế. Hai cơ chế này cùng giải thích vì sao đường
cong duỗi thẳng ở biên trên của gradient: khi chi phí giao dịch xuyên
biên giới đủ thấp và hệ sinh thái hỗ trợ đủ dày, nhánh đi xuống của chữ
U ngược bị đẩy lùi ra ngoài miền quan sát; ngược lại ở biên dưới, sự
thiếu vắng cả hai điều kiện khiến phần thưởng của quốc tế hóa trở nên
nhỏ và bất định đến mức dạng hàm không còn định hình được.

\hypertarget{tci-luxe0-nux103ng-lux1ef1c-nux1ec1n-tux1ea3ng-xuxe1c-nhux1eadn-h2}{%
\subsection{5.1.2 TCI là Năng Lực Nền Tảng: Xác nhận
H2}\label{tci-luxe0-nux103ng-lux1ef1c-nux1ec1n-tux1ea3ng-xuxe1c-nhux1eadn-h2}}

chỉ số năng lực công nghệ (TCI) nhất quán cho thấy \textbf{hiệu ứng nâng
mặt bằng năng suất dương} trong tất cả các mẫu (Việt Nam P3, Trung Quốc
P5, toàn mẫu P7: hệ số +0,141 độ lệch chuẩn năng suất chuẩn hóa cho mỗi
độ lệch chuẩn TCI, p \textless,001). Tuy nhiên, TCI chỉ \textbf{điều
tiết hình dạng đường cong I--P} tại bối cảnh quốc gia cụ thể (P3 Việt
Nam, môi trường chuyển đổi với phân mảnh thể chế tạo điều kiện cho TCI
reshaping hiệu ứng), không phổ quát trên toàn mẫu 50 nền kinh tế (P7:
tương tác FSTS×TCI và FSTS²×TCI đều không ý nghĩa). H2 được \textbf{xác
nhận một phần}: TCI là năng lực nền tảng phổ quát, nhưng vai trò uốn
đường cong quốc tế hóa và hiệu quả phụ thuộc bối cảnh thể chế. Điều này
xác nhận vai trò của TCI là \textbf{năng lực nền tảng} (foundational
capability) trong khung CDCM: doanh nghiệp có năng lực công nghệ sâu
hơn, được đo bằng việc tiếp cận công nghệ nước ngoài, chi tiêu R\&D, đổi
mới sản phẩm và chứng nhận chất lượng, tạo ra giá trị cao hơn từ cùng
một mức độ quốc tế hóa.

Cơ chế lý thuyết đằng sau kết quả này bắt nguồn từ lý thuyết dựa trên
nguồn lực (Barney, 1991): TCI là nguồn lực VRIN (có giá trị, hiếm, khó
sao chép, khó thay thế), do đó tạo ra lợi thế cạnh tranh bền vững ngay
cả khi doanh nghiệp mở rộng ra thị trường nước ngoài với mức độ bất định
cao. Năng lực công nghệ cũng tạo ra năng lực hấp thụ (Cohen \&
Levinthal, 1990), khả năng hấp thu và tích hợp tri thức từ thị trường
nước ngoài, giúp doanh nghiệp học hỏi hiệu quả hơn trong quá trình quốc
tế hóa, nhất quán với learning loop của Uppsala model (Johanson \&
Vahlne, 1977).

\hypertarget{dai-luxe0-nguux1ed3n-lux1ef1c-tuxecnh-huux1ed1ng-hux1ed7-trux1ee3-mux1ed9t-phux1ea7n-h3}{%
\subsection{5.1.3 DAI là Nguồn Lực Tình Huống: Hỗ trợ một phần
H3}\label{dai-luxe0-nguux1ed3n-lux1ef1c-tuxecnh-huux1ed1ng-hux1ed7-trux1ee3-mux1ed9t-phux1ea7n-h3}}

Kết quả về chỉ số chấp nhận số (DAI) phức tạp hơn và có nhiều sắc thái
hơn so với TCI. Trên toàn mẫu P7 (khung 50 nền kinh tế), DAI có hiệu ứng
nâng mặt bằng năng suất dương và có ý nghĩa cao
(\(\hat{\beta} = +0{,}201\), \(p < {,}001\)): doanh nghiệp có hiện diện
số cơ bản đạt năng suất chuẩn hóa cao hơn một phần năm độ lệch chuẩn.
Các tương tác điều tiết đường cong mang đúng chiều nén như khung CDCM dự
đoán (\(\hat{\beta}(\text{FSTS} \times \text{DAI}) = -0{,}271\);
\(\hat{\beta}(\text{FSTS}^2 \times \text{DAI}) = +0{,}252\)) nhưng không
đạt ý nghĩa thống kê ở cấp toàn mẫu (\(p = {,}149\) và \(p = {,}367\)),
nghĩa là hiệu ứng uốn đường cong của DAI không phổ quát mà tập trung ở
các bối cảnh đặc thù. Tại Singapore (P4, DAI Tầng 1 và 2), DAI phát huy
rõ nhất dưới dạng hiệu ứng khuếch đại
(\(\hat{\beta}(\text{FSTS}^2 \times \text{DAI}) = +3{,}119\),
\(p = {,}005\)). Tại Việt Nam (P3, Tầng 1 only, IV-estimated), DAI không
có ý nghĩa phản ánh hạn chế đo lường và thiên lệch chọn lựa hơn là bác
bỏ H3.

Kết quả tổng hợp hỗ trợ lập luận trong khung CDCM rằng DAI là
\textbf{nguồn lực tình huống} (situational resource) chứ không phải năng
lực nền tảng: cơ chế cụ thể của nó (compression vs amplification) phụ
thuộc vào điều kiện môi trường (thể chế, cơ sở hạ tầng số), nhưng bản
thân hiệu ứng DAI là \textbf{phổ quát} trên toàn mẫu đa bối cảnh. Đây là
điểm phân biệt quan trọng giữa TCI và DAI trong đóng góp lý thuyết của
luận án.

Cần lưu ý rằng kết quả không có ý nghĩa về DAI tại Việt Nam (P3) và các
bối cảnh thấp hơn không nhất thiết bác bỏ H3, mà còn phản ánh hạn chế đo
lường: Tầng 1 DAI (website, email, online sales) chưa nắm bắt được chiều
sâu của năng lực số. Khi DAI được đo bằng Tầng 1 và 2 (bổ sung
e-payment, e-supplier) như tại Singapore, hiệu ứng điều tiết xuất hiện
rõ ràng hơn. Quan trọng hơn, tất cả đo lường DAI trong luận án đều là
\textbf{chỉ số chấp nhận số tĩnh} (static digital adoption index) dựa
trên dữ liệu WBES tại một thời điểm, không phải ``dynamic năng lực số''
theo nghĩa của năng lực động theory (Teece et al., 1997) vì dữ liệu WBES
không có đủ thang đo để đo lường tính động của năng lực số.

Một giải thích bổ sung từ Lý thuyết thể chế: ở các bối cảnh ICRV thấp
với nhiều khoảng trống thể chế (Khanna \& Palepu, 2010), doanh nghiệp
xuất khẩu có thể phụ thuộc vào DAI nhiều hơn như một cơ chế bù đắp cho
thiếu hụt hạ tầng thể chế (hiệu ứng lá chắn số đã ghi nhận ở nghiên cứu
thành phần P1 trên 17 nền). Trên khung 50 nền, tương tác chế độ-yếu với
độ dốc đường cong ở mức biên ý nghĩa (FSTS×chế độ yếu = −0,523, p
=,087), nhất quán về chiều với lập luận này nhưng cần kiểm định nhân quả
chặt hơn trước khi khẳng định.

\hypertarget{ux111ux1eb7c-ux111iux1ec3m-nhuxe0-quux1ea3n-trux1ecb-luxe0-biux1ebfn-ux111iux1ec1u-tiux1ebft-cuxe1-nhuxe2n-hux1ed7-trux1ee3-cuxf3-ux111iux1ec1u-kiux1ec7n-h4}{%
\subsection{5.1.4 Đặc Điểm Nhà Quản Trị là Biến điều tiết Cá Nhân: Hỗ
trợ có điều kiện
H4}\label{ux111ux1eb7c-ux111iux1ec3m-nhuxe0-quux1ea3n-trux1ecb-luxe0-biux1ebfn-ux111iux1ec1u-tiux1ebft-cuxe1-nhuxe2n-hux1ed7-trux1ee3-cuxf3-ux111iux1ec1u-kiux1ec7n-h4}}

Kết quả về H4 trên khung 50 nền cho thấy đặc điểm nhà quản trị cấp cao
\textbf{không tạo hiệu ứng điều tiết phổ quát}. Hệ số mức của nữ quản lý
cấp cao là âm nhẹ (\(\hat{\beta} = -0{,}104\), \(p < {,}001\)) sau khi
kiểm soát hiệu ứng cố định nền kinh tế và năm, và tương tác với cường độ
quốc tế hóa không có ý nghĩa thống kê (FSTS×nữ quản lý: \(p = {,}847\)).
Kết quả này nhất quán với lập luận của Lý thuyết thượng tầng quản trị
(Hambrick \& Mason, 1984) ở dạng có điều kiện: tác động của đặc điểm
lãnh đạo phụ thuộc mạnh vào bối cảnh thể chế và cấu trúc ngành mà doanh
nghiệp hoạt động, thay vì là một hiệu ứng đồng nhất xuyên 50 nền kinh
tế; hệ số mức âm nhẹ nhiều khả năng phản ánh sự phân bố không ngẫu nhiên
của nữ quản lý vào các ngành thâm dụng lao động biên thấp, một dạng chọn
lọc ngành cần thiết kế nhận dạng riêng để bóc tách.

Cơ chế lý thuyết có thể qua hai kênh: thứ nhất, \textbf{đa dạng nhận
thức}, nhà quản trị nữ thường mang lại góc nhìn khác biệt trong ra quyết
định quốc tế hóa, giúp nhận diện thị trường ngách và giảm thiểu rủi ro
tập trung; thứ hai, \textbf{chuẩn mực quản trị}, sự hiện diện của lãnh
đạo nữ tương quan với thực hành quản trị doanh nghiệp tốt hơn và minh
bạch hơn, đặc biệt quan trọng trong môi trường thể chế yếu.

Tuy nhiên, cần lưu ý rằng tác động này được phát hiện ở cấp gộp toàn mẫu
(P7) và cần kiểm định riêng biệt theo từng bối cảnh ICRV để hiểu đầy đủ
điều kiện biên của biến điều tiết cấp quản trị. kết quả không có ý nghĩa
về kinh nghiệm nhà quản trị (mgr\_experience) gợi ý rằng trong bối cảnh
châu Á đa dạng, số năm kinh nghiệm đơn thuần không nắm bắt được
\textbf{chất lượng} kinh nghiệm quốc tế (breadth vs.~depth), một hạn chế
đo lường trong WBES cần được khắc phục trong nghiên cứu tương lai.

\hypertarget{fip-luxe0-ux111iux1ec1u-kiux1ec7n-biuxean-cux1ef1c-ux111oan-h1b-cho-sids}{%
\subsection{5.1.5 FIP là Điều Kiện Biên Cực Đoan: H1b cho
SIDS}\label{fip-luxe0-ux111iux1ec1u-kiux1ec7n-biuxean-cux1ef1c-ux111oan-h1b-cho-sids}}

Gánh nặng quốc tế hóa bắt buộc (Forced Internationalization Penalty,
FIP) tại các nền kinh tế SIDS (\(\beta = -1{,}339\), \(p < {,}001\); M1
với hiệu ứng cố định quốc gia và năm; trong mẫu phụ chỉ gồm doanh nghiệp
xuất khẩu, hệ số vẫn âm nhưng không ý nghĩa do N = 26) đại diện cho điều
kiện biên mà CDCM không hoàn toàn dự báo trước trong phiên bản ban đầu
của khung lý thuyết. FIP là biểu hiện của trường hợp khi \textbf{cả ba
tầng của CDCM} (thể chế, năng lực, và quản trị) đều ở mức thấp: thể chế
yếu, TCI thấp, DAI không phát huy tác dụng, thì quốc tế hóa không chỉ
dừng ở mức không mang lại lợi ích mà còn gây tổn hại thực sự đến hiệu
quả doanh nghiệp.

Quan hệ âm đơn điệu này khác về bản chất với quan hệ phi tuyến chữ U
ngược: không có điểm uốn trong phạm vi dữ liệu, nghĩa là không có mức
FSTS nào mà doanh nghiệp SIDS có thể tận dụng để cải thiện hiệu quả
thông qua xuất khẩu trong điều kiện hiện tại. Đây là ranh giới biên lý
thuyết quan trọng cho models phi tuyến phổ quát trong tài liệu về quan
hệ quốc tế hóa và hiệu quả.

\begin{center}\rule{0.5\linewidth}{0.5pt}\end{center}

\hypertarget{ux1ea5n-ux111ux1ed9-ux111iux1ec1u-kiux1ec7n-biuxean-phux1ea1m-vi-uxe1p-dux1ee5ng-cux1ee7a-khung-cdcm-p9}{%
\subsection{5.1.6 Ấn Độ, Điều kiện biên phạm vi áp dụng của khung CDCM
(P9)}\label{ux1ea5n-ux111ux1ed9-ux111iux1ec1u-kiux1ec7n-biuxean-phux1ea1m-vi-uxe1p-dux1ee5ng-cux1ee7a-khung-cdcm-p9}}

\textbf{Ấn Độ và điều kiện biên phạm vi áp dụng của khung CDCM.} Phát
hiện ngưỡng phá vỡ tại Ấn Độ (P9 (Ấn Độ)) qua giai đoạn 2014--2025 bổ
sung một điều kiện biên quan trọng cho khung CDCM. Trong khi quan hệ U
ngược được khung dự đoán giả định môi trường thể chế ổn định, sự sụp đổ
độ cong tại đợt 2025 dưới áp lực tích lũy của bãi bỏ tiền mặt, GST, IBC,
UPI, PLI, COVID-19 và Atmanirbhar Bharat cho thấy bản thân khung CDCM
cũng cần điều kiện ổn định thể chế làm tiền đề. Khung này không khẳng
định U ngược là quan hệ phổ quát mà là quan hệ có điều kiện. Sự bác bỏ
giả thuyết phổ quát của Tầng 2 thay thế Tầng 1 (qua tương tác âm
\(\beta = -4{,}02\), \(p = {,}004\) tại Ấn Độ) cũng tinh chỉnh đề xuất
phụ trợ của khung: hạ tầng số công cộng chỉ thay thế năng lực số tư nhân
khi hai loại hạ tầng có cùng định hướng mục đích. UPI hỗ trợ giao dịch
nội địa hiệu quả nhưng không thay thế hạ tầng tài chính xuyên biên giới
(SWIFT, ngân hàng đại lý, tài trợ thương mại) mà doanh nghiệp xuất khẩu
yêu cầu.

\begin{center}\rule{0.5\linewidth}{0.5pt}\end{center}

\hypertarget{tuxedch-hux1ee3p-ba-vuxf9ng-cux1ee7a-gradient-thux1ec3-chux1ebf-tux1eeb-chux1eef-u-ngux1b0ux1ee3c-ux111ux1ebfn-tan-ruxe3-vuxe0-ux111ux1ea3o-dux1ea5u}{%
\subsection{5.1.7 Tích hợp ba vùng của gradient thể chế: từ chữ U ngược
đến tan rã và đảo
dấu}\label{tuxedch-hux1ee3p-ba-vuxf9ng-cux1ee7a-gradient-thux1ec3-chux1ebf-tux1eeb-chux1eef-u-ngux1b0ux1ee3c-ux111ux1ebfn-tan-ruxe3-vuxe0-ux111ux1ea3o-dux1ea5u}}

Khi đặt cạnh nhau bảy nghiên cứu thành phần thực nghiệm, một cấu trúc ba
vùng nổi lên như trục diễn giải trung tâm của khung lý thuyết CDCM,
không phải dưới dạng ba kết quả rời rạc mà như ba đoạn liên tục của cùng
một họ đường cong phụ thuộc bối cảnh. Ở vùng giữa của gradient thể chế
(chế độ chuyển đổi bậc thấp--trung, Nhóm IV), chữ U ngược định hình rõ
nét và tin cậy nhất: điểm uốn toàn mẫu của Nhóm IV ở mức 43,0\% (P7)
trùng khớp đáng kể với các ước lượng đơn quốc gia độc lập là Việt Nam
39--46\% (P3) và Ấn Độ 40,7\% năm 2022 (P9). Sự hội tụ của ba ước lượng
từ ba nguồn dữ liệu khác nhau quanh một dải hẹp là bằng chứng nội tại
mạnh về tính vững của dạng hàm trong chế độ chuyển đổi.

Ở biên trên của gradient (chế độ tiên tiến đổi mới, Nhóm I), đường cong
duỗi thẳng thành quan hệ gần tuyến tính dương, với nhánh đi xuống bị đẩy
ra ngoài miền cường độ xuất khẩu quan sát được. Ba nguồn bằng chứng độc
lập hội tụ về kết luận này: Singapore với điểm uốn ước lượng quanh
88,6\% nhưng không định vị tin cậy (P4); ước lượng cấp nhóm Advanced
trên khung 50 nền với điểm uốn ngoài miền dữ liệu (P7); và Nhật Bản với
hiệu ứng tuyến tính dương cao (+0,671, p \textless{} 0,001) trong khi số
hạng bậc hai không bao giờ tiệm cận ý nghĩa qua mọi đặc tả (P10). Bằng
chứng Nhật Bản mang sức nặng đặc biệt vì chính các tập đoàn đa quốc gia
Nhật Bản là nơi giả thuyết phi tuyến được thiết lập ban đầu; việc nền
kinh tế này, khi được quan sát trên bộ công cụ hài hòa, lại nằm trọn
trên một nhánh đi lên dài, xác nhận rằng dạng phi tuyến cổ điển là kết
quả có điều kiện chứ không phải quy luật phổ quát.

Ở biên dưới của gradient (chế độ đang nổi và đảo nhỏ, Nhóm V và VI),
quan hệ mất cấu trúc theo hai mức độ khác nhau. Tại Nhóm V, các hệ số
không còn ý nghĩa thống kê và dạng hàm tan rã (P7). Tại Nhóm VI là các
nền kinh tế đảo nhỏ, quan hệ không chỉ tan rã mà đảo hẳn sang âm đơn
điệu, gánh nặng quốc tế hóa bắt buộc với hệ số −1,339 (p \textless{}
0,001, P8). Như vậy, ba vùng không phải là ba hiện tượng tách biệt mà là
biểu hiện liên tục của cùng một cơ chế thể chế: khi đệm thể chế dày, dư
địa quốc tế hóa có lợi mở rộng đến mức nhánh suy giảm biến mất; khi đệm
thể chế mỏng dần, điểm uốn dịch về phía trái và độ cong sắc nét lên; khi
đệm thể chế cạn kiệt, bản thân điểm tối ưu biến mất và phần thưởng của
quốc tế hóa chuyển thành tổn thất.

Một sắc thái quan trọng phân biệt hai cơ chế tạo ra chữ U ngược dọc theo
gradient. Bằng chứng Việt Nam (P3) cho thấy độ cong toàn mẫu trong chế
độ chuyển đổi được nhận diện chủ yếu qua biên tham gia, tức bước nhảy
năng suất khi doanh nghiệp chuyển từ không xuất khẩu sang xuất khẩu
(H1a), chứ không qua độ cong cường độ trong nhóm doanh nghiệp đã xuất
khẩu (H1b); khi tái ước lượng chỉ trên doanh nghiệp xuất khẩu, số hạng
bậc hai mất ý nghĩa. Bằng chứng Nhật Bản (P10) cho thấy ở biên trên, cơ
chế cũng vận hành qua biên tham gia nhưng theo logic tự chọn lọc Melitz
(2003): câu hỏi ràng buộc không phải doanh nghiệp xuất khẩu bao nhiêu mà
là có xuất khẩu hay không, với phần bù năng suất 25,8\% cho doanh nghiệp
xuất khẩu. Sự dịch chuyển trọng tâm từ biên cường độ ở chế độ chuyển đổi
sang biên tham gia ở chế độ tiên tiến là một tinh chỉnh lý thuyết mà
khung CDCM phiên bản ban đầu chưa nêu bật, và nó cung cấp một giải thích
thống nhất cho việc vì sao đường cong duỗi thẳng ở biên trên.

\hypertarget{so-suxe1nh-vux1edbi-nghiuxean-cux1ee9u-trux1b0ux1edbc-ux111uxe2y}{%
\section{5.2 So sánh với nghiên cứu trước
đây}\label{so-suxe1nh-vux1edbi-nghiuxean-cux1ee9u-trux1b0ux1edbc-ux111uxe2y}}

\hypertarget{nhux1ea5t-quuxe1n-vux1edbi-cux1a1-sux1edf-icbef-2024}{%
\subsection{5.2.1 Nhất quán với cơ sở ICBEF
2024}\label{nhux1ea5t-quuxe1n-vux1edbi-cux1a1-sux1edf-icbef-2024}}

Kết quả hiệu ứng gộp \(r = 0{,}074\) (P6) nhất quán hoàn toàn với
\(r = 0{,}07\) từ phân tích tổng hợp cơ sở đã công bố tại hội nghị ICBEF
2024 (nghiên cứu thành phần P1--P2 của luận án). Sự hội tụ này là đáng
kể vì P6 mở rộng mẫu gộp lên \(k = 238\) so với \(k = 113\), đồng thời
áp dụng cấu trúc ba tầng phức tạp hơn. Tính vững của hiệu ứng gộp khi mở
rộng mẫu gộp nghiên cứu là bằng chứng về tính ổn định của cơ sở ước
lượng.

Tuy nhiên, đóng góp của P6 so với cơ sở không phải ở gộp \(r\) mà ở cấu
trúc dị biệt: phân tách \(I^2 = 62{,}4\%\) thành 54,1\% tầng 2 và 8,4\%
tầng 3 là thông tin mới mà phân tích phân tích tổng hợp đơn tầng không
cung cấp. Điều này thay đổi hiểu biết về nguồn gốc của dị biệt trong tài
liệu, vấn đề chủ yếu là \textbf{đo lường bên trong nghiên cứu}, không
phải sự bất đồng căn bản giữa các nghiên cứu.

\hypertarget{vux1b0ux1ee3t-qua-lu-vuxe0-beamish-2004-vuxe0-contractor-et-al.-2003}{%
\subsection{5.2.2 Vượt qua Lu và Beamish (2004) và Contractor et
al.~(2003)}\label{vux1b0ux1ee3t-qua-lu-vuxe0-beamish-2004-vuxe0-contractor-et-al.-2003}}

Lu và Beamish (2004) và Contractor et al.~(2003) là hai công trình nền
tảng về quan hệ phi tuyến quốc tế hóa và hiệu quả. Lu và Beamish (2004)
tìm điểm uốn khoảng 60\% FSTS trong mẫu doanh nghiệp Nhật Bản, trong khi
Contractor et al.~(2003) đề xuất đường cong chữ S ba giai đoạn. Luận án
này vượt qua hai công trình đó theo ba hướng:

\textbf{Thứ nhất}, luận án chứng minh rằng điểm uốn \textbf{không cố
định} mà phụ thuộc vào bối cảnh ICRV: từ TP ≈ 88,6\% (Singapore, tiên
tiến đổi mới; lưu ý điểm uốn này không định vị tin cậy do CI vượt ngưỡng
khả thi, xem Mục 4.3) đến TP ≈ 39,7\% (Việt Nam pooled, chuyển đổi thu
nhập trung bình thấp), với Trung Quốc ở mức trung gian (TP ≈ 48,78\%).
Không có một điểm uốn phổ quát cho tất cả bối cảnh.

\textbf{Thứ hai}, luận án xác định trường hợp biên FIP tại SIDS, không
phải là đường cong chữ S hay chữ U ngược mà là quan hệ âm đơn điệu, điều
mà cả ba-stage theory lẫn đường cong chữ S không dự báo được.

\textbf{Thứ ba}, luận án mở rộng khung phân tích từ dữ liệu doanh nghiệp
Nhật Bản/đa quốc gia sang 50 nền kinh tế châu Á và Thái Bình Dương, bối
cảnh underrepresented trong tài liệu Lu và Beamish.

\hypertarget{mux1edf-rux1ed9ng-marano-et-al.-2016-icrv-thay-thux1ebf-phuxe2n-loux1ea1i-thux1ec3-chux1ebf-nhux1ecb-phuxe2n}{%
\subsection{5.2.3 Mở rộng Marano et al.~(2016): ICRV thay thế phân loại
thể chế nhị
phân}\label{mux1edf-rux1ed9ng-marano-et-al.-2016-icrv-thay-thux1ebf-phuxe2n-loux1ea1i-thux1ec3-chux1ebf-nhux1ecb-phuxe2n}}

Marano et al.~(2016) là phân tích tổng hợp quan trọng nhất trong tài
liệu gần đây, cung cấp bằng chứng về vai trò điều tiết của thể chế quốc
gia nguồn trong quan hệ quốc tế hóa và hiệu quả. Tuy nhiên, nghiên cứu
này sử dụng phân loại nhị phân (phát triển so với mới nổi) hoặc điểm
chất lượng thể chế liên tục.

Luận án này mở rộng Marano et al.~(2016) theo hướng \textbf{phân loại đa
chiều}: khung ICRV sáu chế độ nắm bắt được sự đa dạng trong nội bộ nhóm
mới nổi, phân biệt nhóm trung bình cao của Trung Quốc với nhóm chuyển
đổi thu nhập trung bình thấp của Việt Nam, với nhóm đang nổi của
Campuchia và nhóm đảo nhỏ; đồng thời phân biệt trong nhóm tiên tiến giữa
chế độ tiên tiến đổi mới của Singapore và chế độ tiên tiến tài nguyên
của các nền vùng Vịnh. Phân loại nhị phân bỏ sót sự biến thiên quan
trọng này, dẫn đến mất thông tin về gradient thể chế. Quan trọng hơn
việc mất thông tin, phân loại nhị phân tạo ra một sai lệch triệt tiêu
dấu khi gộp: dưới một nhãn ``mới nổi'' duy nhất, hệ số điều tiết của
Marano et al.~(2016) lấy trung bình đồng thời các nền có độ cong chữ U
ngược âm rõ (Trung Quốc, Việt Nam), các nền gần tuyến tính, và các nền
đảo nhỏ có quan hệ âm đơn điệu (FIP) với độ cong bậc hai không tạo điểm
uốn nội vùng; kết quả là một hệ số điều tiết gộp che lấp chính sự phân
kỳ mà nó cần bộc lộ. Đóng góp của ICRV do đó không nằm ở chỗ ``có nhiều
nhóm hơn'' mà ở chỗ tránh được hiện tượng triệt tiêu dấu trong nội bộ
nhóm mới nổi, qua đó phục hồi thông tin gradient mà thước đo nhị phân
làm mất.

Ngoài ra, luận án sử dụng phương pháp MARA ba tầng để phân tách cấu trúc
dị biệt, một cải tiến phương pháp so với phân tích tổng hợp hiệu ứng
ngẫu nhiên đơn tầng của Marano et al.~(2016).

\hypertarget{luxfd-giux1ea3i-cuxe1c-kux1ebft-quux1ea3-nhuxe2n-rux1ed9ng-quy-muxf4-lux1edbn-cuxf2n-phuxe2n-tuxe1n}{%
\subsection{5.2.4 Lý giải các kết quả nhân rộng quy mô lớn còn phân
tán}\label{luxfd-giux1ea3i-cuxe1c-kux1ebft-quux1ea3-nhuxe2n-rux1ed9ng-quy-muxf4-lux1edbn-cuxf2n-phuxe2n-tuxe1n}}

Một làn sóng nghiên cứu nhân rộng quy mô lớn gần đây đặt lại câu hỏi về
tính bền vững của quan hệ quốc tế hóa và hiệu quả. Pisani và cộng sự
(2020), khi nhân rộng quan hệ đa quốc gia và hiệu quả trên mẫu xuyên
quốc gia cỡ lớn, tìm thấy một quan hệ yếu và không nhất quán giữa các
nước; Berry và Kaul (2016) kiểm chứng lại đường cong chữ S và kết luận
rằng dạng hàm này không nhân rộng vững chắc; Kirca và cộng sự (2012),
qua phân tích tổng hợp, cho thấy bối cảnh điều tiết mạnh quan hệ này.
Điểm chung của ba nghiên cứu là một kết quả gộp bị suy giảm hoặc bất
nhất khi nhiều bối cảnh được gộp chung mà không phân tầng thể chế.

Bằng chứng ba vùng của luận án cung cấp một lời giải thích cấu trúc cho
mô thức phân tán này. Khi gộp các nền kinh tế thuộc những chế độ thể chế
khác nhau vào một ước lượng duy nhất, người ta đang lấy trung bình của
ba dạng quan hệ khác nhau về bản chất: chữ U ngược sắc nét ở vùng chuyển
đổi (Nhóm IV, điểm uốn quanh 43\%), quan hệ gần tuyến tính dương ở biên
thể chế mạnh (Nhóm I, điểm uốn nằm ngoài miền quan sát), và quan hệ tan
rã hoặc đảo dấu ở biên thể chế yếu (Nhóm V và VI). Phép lấy trung bình
này tất yếu tạo ra một tín hiệu gộp bị nén và bất ổn, đúng như các
nghiên cứu nhân rộng ghi nhận. Như vậy, kết quả của luận án không mâu
thuẫn với tài liệu nhân rộng mà giải thích nó: tính phân tán của các kết
quả gộp không phải bằng chứng rằng quan hệ quốc tế hóa và hiệu quả là
giả tạo, mà là hệ quả đo lường của việc bỏ qua phân tầng thể chế. Việc
một nghiên cứu đơn quốc gia có tìm thấy chữ U ngược hay không phụ thuộc
vào vùng thể chế mà nền kinh tế đó chiếm giữ, và đây chính là biến số mà
các phép nhân rộng gộp đã làm mờ đi.

Một phản biện sắc bén có thể nêu là bản thân cấu trúc ba vùng có nguy cơ
là tạo tác của việc nhà nghiên cứu tự chọn điểm cắt ICRV, tức một dạng
dò tìm đặc tả hậu nghiệm thay vì một quy luật. Hai đặc điểm của thiết kế
bác bỏ cách đọc này. Thứ nhất, các điểm cắt ICRV được xác định trước
trên các chỉ báo thể chế (chỉ số quản trị toàn cầu, mức phát triển, đặc
trưng cấu trúc), độc lập với biến kết quả quốc tế hóa và hiệu quả, nên
chúng không thể được nắn để khớp với hình dạng đường cong. Thứ hai, các
ước lượng điểm uốn đơn quốc gia độc lập, Việt Nam (khoảng 39--46\% qua
ba đợt) và Trung Quốc (khoảng 47--49\% qua hai đợt), hội tụ về cùng một
dải vùng giữa mà không hề bị ràng buộc phải làm vậy, một sự trùng khớp
ngoài mẫu khó nảy sinh từ việc dò tìm đặc tả. Hai đặc điểm này, kết hợp
với tam giác hóa giữa phân tích tổng hợp và ước lượng gộp ở Mục 5.2.5,
đặt cấu trúc ba vùng trên nền tảng xác nhận chéo chứ không phải trên tự
do bậc đặc tả.

\hypertarget{tam-giuxe1c-huxf3a-giux1eefa-phuxe2n-tuxedch-tux1ed5ng-hux1ee3p-p6-vuxe0-ux1b0ux1edbc-lux1b0ux1ee3ng-gux1ed9p-touxe0n-mux1eabu-p7}{%
\subsection{5.2.5 Tam giác hóa giữa phân tích tổng hợp (P6) và ước lượng
gộp toàn mẫu
(P7)}\label{tam-giuxe1c-huxf3a-giux1eefa-phuxe2n-tuxedch-tux1ed5ng-hux1ee3p-p6-vuxe0-ux1b0ux1edbc-lux1b0ux1ee3ng-gux1ed9p-touxe0n-mux1eabu-p7}}

Một nguồn củng cố độc lập cho luận điểm điều tiết thể chế đến từ sự hội
tụ giữa hai phương pháp khác biệt về bản chất. Phân tích tổng hợp ba
tầng (P6), vận hành trên các cỡ ảnh hưởng đã công bố từ nhiều nghiên cứu
độc lập, ước lượng cỡ ảnh hưởng trung bình \(\bar{r} = 0{,}074\) với dị
biệt đáng kể (\(I^2 = 62{,}4\%\)) và xác nhận chế độ thể chế ICRV là
biến điều tiết có ý nghĩa (\(Q_M = 17{,}35\), \(df = 4\),
\(p = {,}002\)). Ước lượng gộp toàn mẫu trên dữ liệu doanh nghiệp gốc
của 50 nền kinh tế (P7) đi đến cùng một kết luận bằng con đường hoàn
toàn khác: cấu trúc thể chế của quan hệ hiện ra thành ba vùng rời rạc
khi ước lượng điểm uốn trong từng nhóm ICRV.

Sự hội tụ này có giá trị suy luận vượt quá tổng của hai phần. Phân tích
tổng hợp và ước lượng gộp cấp doanh nghiệp chịu các nguồn thiên lệch
khác nhau: phân tích tổng hợp dễ bị thiên lệch công bố và phụ thuộc vào
chất lượng báo cáo của các nghiên cứu gốc, trong khi ước lượng gộp cấp
doanh nghiệp dễ bị thiên lệch chọn lọc tham gia xuất khẩu và chênh lệch
đo lường giữa các đợt khảo sát. Việc hai phương pháp với các điểm yếu
không trùng nhau cùng định vị chế độ thể chế là biến điều tiết then chốt
làm cho kết luận này khó có thể là tạo tác của một quy trình phân tích
đơn lẻ. Đây là dạng tam giác hóa phương pháp mà các điểm yếu bù trừ lẫn
nhau, nâng độ tin cậy của luận điểm trung tâm rằng thể chế quyết định
không chỉ vị trí mà cả sự tồn tại của dạng hàm chữ U ngược.

\hypertarget{ux111ux1ed1i-chiux1ebfu-vux1edbi-truyux1ec1n-thux1ed1ng-ux111ux1b0ux1eddng-cong-chux1eef-s-vuxe0-cuxe1c-kiux1ec3m-ux111ux1ecbnh-nhuxe2n-rux1ed9ng-chux1eef-s}{%
\subsection{5.2.6 Đối chiếu với truyền thống đường cong chữ S và các
kiểm định nhân rộng chữ
S}\label{ux111ux1ed1i-chiux1ebfu-vux1edbi-truyux1ec1n-thux1ed1ng-ux111ux1b0ux1eddng-cong-chux1eef-s-vuxe0-cuxe1c-kiux1ec3m-ux111ux1ecbnh-nhuxe2n-rux1ed9ng-chux1eef-s}}

Một dòng nghiên cứu song hành với truyền thống chữ U ngược là truyền
thống đường cong chữ S ba giai đoạn, vốn lập luận rằng quan hệ quốc tế
hóa và hiệu quả đi qua một giai đoạn đầu suy giảm (chi phí học hỏi), một
giai đoạn giữa đi lên (khai thác lợi thế đa quốc gia), rồi một giai đoạn
cuối lại suy giảm (chi phí điều phối vượt ngưỡng). Berry và Kaul (2016),
khi kiểm định nhân rộng dạng chữ S trên mẫu xuyên quốc gia cỡ lớn, kết
luận rằng dạng hàm bậc ba này không nhân rộng vững chắc và phần lớn bằng
chứng ủng hộ nó nhạy cảm với cách đặc tả và mẫu. Phát hiện ba vùng của
luận án cung cấp một cách đọc bổ sung cho kết quả nhân rộng âm này: nếu
hình dạng quan hệ thực sự thay đổi theo chế độ thể chế, từ chữ U ngược
sắc nét ở vùng chuyển đổi đến gần tuyến tính dương ở biên trên đến âm
đơn điệu ở biên dưới, thì việc ép một dạng hàm bậc ba cố định lên một
mẫu trộn lẫn nhiều chế độ thể chế tất yếu cho ra một ước lượng bất ổn.
Nói cách khác, sự thất bại nhân rộng của dạng chữ S không nhất thiết bác
bỏ tính phi tuyến mà có thể phản ánh việc bỏ qua tính không đồng nhất
theo thể chế của bản thân dạng hàm.

Bằng chứng Nhật Bản (P10) mang sức nặng đặc biệt trong đối chiếu này.
Chính các tập đoàn đa quốc gia Nhật Bản là nơi cả giả thuyết chữ U ngược
của Lu và Beamish (2004) lẫn giả thuyết chữ S được thiết lập ban đầu.
Khi nền kinh tế này được quan sát trên bộ công cụ hài hòa của luận án,
quan hệ lại nằm trọn trên một nhánh đi lên dài với phần bù năng suất cho
doanh nghiệp xuất khẩu, không có điểm uốn trong miền quan sát. Sự khác
biệt giữa kết quả này và điểm uốn khoảng 60\% mà Lu và Beamish (2004)
tìm thấy không phản ánh mâu thuẫn dữ liệu mà phản ánh khác biệt khung
lấy mẫu: mẫu WBES Nhật Bản thiên về doanh nghiệp cơ sở với cường độ quốc
tế hóa thấp, trong khi mẫu Lu và Beamish tập trung vào các tập đoàn đa
quốc gia thâm dụng xuất khẩu cao, tức quan sát hai đoạn khác nhau của
cùng một họ đường cong. Đối chiếu này củng cố luận điểm rằng dạng phi
tuyến cổ điển là kết quả có điều kiện theo cả chế độ thể chế lẫn dải
cường độ xuất khẩu được quan sát, chứ không phải một quy luật phổ quát,
nhất quán với cảnh báo phương pháp của Haans và cộng sự (2016) về việc
suy diễn dạng hàm phi tuyến từ những đoạn dữ liệu giới hạn.

\hypertarget{ux111ux1eb7t-vuxe0o-duxf2ng-nghiuxean-cux1ee9u-vux1ec1-ux111iux1ec1u-tiux1ebft-bux1edfi-muxf4i-trux1b0ux1eddng-cux1ee7a-quan-hux1ec7-nux103ng-lux1ef1c-vuxe0-hiux1ec7u-quux1ea3}{%
\subsection{5.2.7 Đặt vào dòng nghiên cứu về điều tiết bởi môi trường
của quan hệ năng lực và hiệu
quả}\label{ux111ux1eb7t-vuxe0o-duxf2ng-nghiuxean-cux1ee9u-vux1ec1-ux111iux1ec1u-tiux1ebft-bux1edfi-muxf4i-trux1b0ux1eddng-cux1ee7a-quan-hux1ec7-nux103ng-lux1ef1c-vuxe0-hiux1ec7u-quux1ea3}}

Phát hiện rằng năng lực công nghệ nâng mặt bằng năng suất phổ quát còn
vai trò uốn đường cong của chấp nhận số lại phụ thuộc bối cảnh thể chế
có thể được đặt vào dòng nghiên cứu rộng hơn về vai trò điều tiết của
môi trường trong quan hệ giữa năng lực và hiệu quả. Karna, Richter và
Riesenkampff (2016), qua phân tích tổng hợp về quan hệ năng lực và hiệu
quả tài chính, cho thấy động lực môi trường điều tiết một cách hệ thống
mức độ mà năng lực chuyển hóa thành hiệu quả. Kết quả của luận án bổ
sung một sắc thái cho kết luận đó: không phải mọi loại năng lực đều bị
điều tiết đồng đều bởi môi trường thể chế. Năng lực công nghệ chiều sâu
(TCI) hoạt động như một yếu tố nâng sàn gần như bất biến với chế độ thể
chế, trong khi chấp nhận số (DAI) là nguồn lực tình huống có hiệu ứng
uốn đường cong bị điều kiện hóa bởi tầng thể chế. Sự phân tách này tinh
chỉnh kết luận tổng hợp của Karna và cộng sự (2016) theo hướng phân biệt
loại năng lực: điều tiết môi trường tác động mạnh lên các nguồn lực bổ
trợ phụ thuộc hệ sinh thái nhưng yếu lên các năng lực nền tảng nội sinh
hóa trong doanh nghiệp.

Hướng đối chiếu này cũng kết nối với lập luận của Kafouros và cộng sự
(2023) rằng chất lượng thể chế và động lực ngành cùng định hình hiệu quả
doanh nghiệp ở các nền kinh tế mới nổi, trong đó tính động công nghệ tạo
tương tác dương còn tính động thị trường tạo tương tác âm. Bằng chứng
gradient ICRV của luận án mở rộng quan sát này từ một mẫu khu vực sang
khung 50 nền kinh tế châu Á và Thái Bình Dương, và bổ sung chiều quan
trọng rằng chính dạng hàm của quan hệ quốc tế hóa và hiệu quả, chứ không
chỉ độ lớn hệ số, thay đổi theo chất lượng thể chế. Đồng thời, lập luận
về vai trò bù đắp của chấp nhận số ở các chế độ thể chế yếu nhất quán
với tài liệu mới nổi về năng lực số như đòn bẩy quốc tế hóa cho doanh
nghiệp nhỏ và vừa trong môi trường thể chế kém hoàn thiện (Banalieva \&
Dhanaraj, 2019; Chen \& Meng, 2022), song luận án đặt một điều kiện biên
rõ hơn: hiệu ứng bù đắp này phụ thuộc chiều sâu của thước đo năng lực số
và không tự nó đảo được độ dốc âm ở những chế độ thể chế cực yếu như
nhóm đảo nhỏ.

\begin{center}\rule{0.5\linewidth}{0.5pt}\end{center}

\hypertarget{ux111uxf3ng-guxf3p-luxfd-thuyux1ebft-1}{%
\section{5.3 Đóng góp lý
thuyết}\label{ux111uxf3ng-guxf3p-luxfd-thuyux1ebft-1}}

\hypertarget{cdcm-luxe0m-ruxf5-dai-luxe0-nguux1ed3n-lux1ef1c-tuxecnh-huux1ed1ng}{%
\subsection{5.3.1 CDCM: Làm rõ DAI là nguồn lực tình
huống}\label{cdcm-luxe0m-ruxf5-dai-luxe0-nguux1ed3n-lux1ef1c-tuxecnh-huux1ed1ng}}

Đóng góp lý thuyết đầu tiên là làm rõ \textbf{bản chất của DAI} trong
khung CDCM. Phát hiện rằng DAI có hiệu ứng nâng mặt bằng phổ quát
(β=+0,201, p\textless,001, mẫu gộp N=79.080 từ 50 nền kinh tế) nhưng vai
trò điều tiết đường cong quốc tế hóa và hiệu quả là \textbf{bối cảnh-phụ
thuộc} (các tương tác đường cong cùng chiều nhưng không ý nghĩa toàn
mẫu; bằng chứng rõ nhất ở bối cảnh đơn quốc gia P4) gợi ý rằng DAI không
phải là năng lực nền tảng (như TCI) mà là \textbf{nguồn lực tình huống},
tài nguyên có vai trò thay đổi theo bối cảnh thể chế. Quan trọng: tương
tác chế độ-yếu (FSTS×Weak = −0,523, p =,087, biên ý nghĩa) cho thấy hiệu
ứng điều tiết của DAI \textbf{mạnh hơn} tại các nền kinh tế thể chế yếu
hơn (đang nổi, SIDS), ``lá chắn số'' bù đắp cho thiếu hụt thể chế.

Điều này có hàm ý lý thuyết quan trọng cho lý thuyết dựa trên nguồn lực
và tài liệu năng lực số (Bhandari et al., 2023; Verhoef et al., 2021):
không phải mọi nguồn lực số đều tạo ra lợi thế cạnh tranh theo cùng một
cơ chế trong mọi bối cảnh. Giá trị của DAI là \textbf{tài sản bổ trợ}
(Teece, 1986) với bối cảnh thể chế, nhưng theo hướng \textbf{nghịch
chiều}: DAI hoạt động như lá chắn bù đắp (compensatory asset) trong bối
cảnh thể chế yếu, nơi hạ tầng thể chế chưa đủ hỗ trợ cho doanh nghiệp
xuất khẩu. Đây là bằng chứng thực nghiệm bổ sung cho lý thuyết về vai
trò thay thế giữa thể chế và năng lực số.

\hypertarget{phuxe2n-biux1ec7t-nuxe2ng-mux1eb7t-bux1eb1ng-phux1ed5-quuxe1t-vuxe0-uux1ed1n-ux111ux1b0ux1eddng-cong-phi-phux1ed5-quuxe1t}{%
\subsection{5.3.1.1 Phân biệt nâng mặt bằng phổ quát và uốn đường cong
phi phổ
quát}\label{phuxe2n-biux1ec7t-nuxe2ng-mux1eb7t-bux1eb1ng-phux1ed5-quuxe1t-vuxe0-uux1ed1n-ux111ux1b0ux1eddng-cong-phi-phux1ed5-quuxe1t}}

Một đóng góp lý thuyết tinh tế nhưng có hệ quả sâu của luận án là phân
biệt rạch ròi giữa hai loại hiệu ứng mà các cấu trúc năng lực có thể tạo
ra trên quan hệ quốc tế hóa và hiệu quả: hiệu ứng nâng mặt bằng (dịch
chuyển toàn bộ đường cong lên trên) và hiệu ứng uốn đường cong (thay đổi
hình dạng, độ dốc hoặc vị trí điểm uốn). Sự phân biệt này, vốn được hình
thức hóa trong tài liệu kiểm định quan hệ phi tuyến của Haans, Pieters
và He (2016), trở thành lăng kính tổ chức cho toàn bộ bằng chứng về TCI
và DAI.

Đối với năng lực công nghệ, bằng chứng hội tụ mạnh mẽ rằng đây là yếu tố
nâng mặt bằng phổ quát: hiệu ứng trực tiếp dương và có ý nghĩa cao xuất
hiện nhất quán ở mọi bối cảnh (Việt Nam +0,179 dạng biến công cụ; Trung
Quốc từ +0,26 đến +0,43; toàn mẫu 50 nền +0,141; Nhật Bản +0,100; SIDS
+0,299), trong khi các tương tác uốn đường cong với cường độ xuất khẩu
không đạt ý nghĩa thống kê trên toàn mẫu và chỉ định hình ở một bối cảnh
đơn quốc gia. Ngay cả tại SIDS, nơi quan hệ tổng thể là âm đơn điệu,
năng lực công nghệ vẫn nâng mặt bằng năng suất nhưng không thể uốn độ
dốc âm sang vùng dương, một bằng chứng đặc biệt thuyết phục rằng vai trò
của năng lực công nghệ là nâng sàn chứ không phải định hình đường cong.

Đối với chấp nhận số, mô thức cũng tương tự nhưng với một khác biệt then
chốt về cường độ ngữ cảnh. Hiệu ứng nâng mặt bằng của DAI là phổ quát và
lớn hơn cả TCI trên toàn mẫu 50 nền (+0,201, p \textless{} 0,001), nhưng
vai trò uốn đường cong của nó không phổ quát: các tương tác đường cong
trên toàn mẫu mang đúng chiều nén mà khung CDCM dự đoán nhưng không đạt
ý nghĩa thống kê, và chỉ phát huy rõ rệt ở bối cảnh đặc thù (Singapore,
tương tác bậc hai +3,119, p = 0,005). Phát hiện này tinh chỉnh đóng góp
lý thuyết theo hướng quan trọng: sự khác biệt giữa TCI và DAI không nằm
ở việc cái nào nâng mặt bằng (cả hai đều nâng), mà ở chỗ vai trò uốn
đường cong của DAI bị điều kiện hóa bởi tầng phát triển thể chế và mức
độ chiều sâu của thước đo, trong khi TCI thuần túy là yếu tố nền tảng
nâng sàn. Đây là cơ sở thực nghiệm trực tiếp cho lập luận rằng năng lực
công nghệ là năng lực nền tảng còn chấp nhận số là nguồn lực tình huống
trong khung CDCM.

Hệ quả phương pháp của phân biệt này không nhỏ. Tài liệu số hóa thường
gộp hai cấu trúc thành một chỉ số tổng hợp duy nhất, làm che khuất chính
sự khác biệt về cơ chế mà luận án phát hiện. Việc tách bạch TCI và DAI
thành hai hợp phần hình thành độc lập, dựa trên các tiêu chí đo lường
của Coltman và cộng sự (2008) và phân biệt khái niệm của Bharadwaj và
cộng sự (2013), cho phép luận án nhận diện rằng hai nguồn lực số tham
gia vào quan hệ quốc tế hóa và hiệu quả qua hai con đường khác nhau, một
phát hiện không thể đạt được nếu chúng bị gộp chung.

\hypertarget{fip-khuxe1i-niux1ec7m-mux1edbi-vux1ec1-trux1b0ux1eddng-hux1ee3p-biuxean-cho-sids}{%
\subsection{5.3.2 FIP: Khái niệm mới về trường hợp biên cho
SIDS}\label{fip-khuxe1i-niux1ec7m-mux1edbi-vux1ec1-trux1b0ux1eddng-hux1ee3p-biuxean-cho-sids}}

Đóng góp lý thuyết thứ hai là đề xuất và đặt tên cho \textbf{hình phạt
quốc tế hóa cưỡng bức} (Forced Internationalization Penalty, FIP), một
cấu trúc lý thuyết mới chưa được định danh trong tài liệu kinh doanh
quốc tế, như một hiện tượng lý thuyết độc lập, khác biệt về bản chất với
nhánh đi xuống giai đoạn đầu trong đường cong chữ S. FIP không phải là
giai đoạn đầu của quá trình học hỏi với chi phí học hỏi như mô hình
Uppsala dự đoán, mà là một trạng thái cấu trúc bất lợi dài hạn: doanh
nghiệp ở các nền đảo nhỏ phải xuất khẩu để tồn tại vì thị trường nội địa
quá nhỏ, nhưng không có năng lực và môi trường để đạt lợi thế từ xuất
khẩu.

Giá trị khái niệm của FIP được làm rõ qua việc phân biệt nó với bốn cấu
trúc lân cận trong tài liệu, một phân biệt cần thiết để khẳng định FIP
không phải là cách gọi mới của một ý tưởng cũ. Thứ nhất, khác với chi
phí giai đoạn đầu trong lý thuyết ba giai đoạn (Contractor et al.,
2003), vốn là chi phí học hỏi tạm thời sẽ được bù đắp sau khi vượt điểm
uốn, FIP không có điểm uốn trong miền dữ liệu quan sát vì ràng buộc của
nó mang tính cấu trúc chứ không phải giai đoạn. Thứ hai, khác với gánh
nặng người nước ngoài (Zaheer, 1995), vốn vận hành ở phía thị trường
đích khi doanh nghiệp nước ngoài thâm nhập, FIP vận hành ở phía thị
trường nguồn, tức là phần thiệt mà chính doanh nghiệp đảo nhỏ gánh chịu
khi xuất khẩu. Thứ ba, khác với khởi nghiệp do nhu cầu thiết yếu trong
tài liệu khởi nghiệp, nơi doanh nhân trong bối cảnh phát triển có thể
đạt ngang bằng hiệu quả khi có đủ thời gian và hỗ trợ thể chế, FIP đặc
trưng cho bối cảnh nơi hỗ trợ thể chế vắng mặt một cách hệ thống, khiến
con đường quốc tế hóa do cưỡng bức bị bất lợi dai dẳng. Thứ tư, khác với
chỉ số tổn thương vĩ mô vốn đo mức phơi nhiễm của nền kinh tế trước cú
sốc bên ngoài, FIP là một cấu trúc vi mô ở cấp doanh nghiệp, định danh
chính cơ chế hiệu quả mà qua đó tổn thương cấu trúc vĩ mô chuyển hóa
thành kết quả cấp doanh nghiệp.

Đặt cạnh lý thuyết bàn đạp (Luo \& Tung, 2007), FIP còn mang một giá trị
bổ sung: nó xác định điều kiện biên nơi chính logic bàn đạp thất bại. Lý
thuyết bàn đạp giả định rằng doanh nghiệp nền kinh tế mới nổi có thể
dùng quốc tế hóa như một bàn đạp để thâu nhận tài sản chiến lược và bù
đắp bất lợi quê nhà. FIP cho thấy một chế độ thể chế, nhóm đảo nhỏ, nơi
giả định này sụp đổ: doanh nghiệp buộc phải quốc tế hóa nhưng thiếu cả
nền tảng quy mô nội địa lẫn đệm thể chế để biến quốc tế hóa thành tích
lũy năng lực, nên không có bàn đạp nào để bật lên mà chỉ còn gánh nặng.
Theo nghĩa này, luận án đóng góp một điều kiện biên cho lý thuyết bàn
đạp: cơ chế bàn đạp chỉ vận hành khi tồn tại tối thiểu một nền tảng quy
mô và một mức hỗ trợ thể chế; dưới ngưỡng đó, cùng một hành vi quốc tế
hóa chuyển từ cơ chế xây dựng năng lực sang cơ chế bào mòn hiệu quả. Đây
là một đóng góp lý thuyết kép: vừa định danh FIP như một hiện tượng độc
lập, vừa định vị nó như ranh giới áp dụng của một lý thuyết quốc tế hóa
nền mới nổi đang có ảnh hưởng, qua đó nối liền mạch lập luận từ lăng
kính bàn đạp được giới thiệu ở Mục 2.2.3.1 đến trường hợp biên thực
nghiệm của luận án.

Khái niệm FIP do đó mở ra hướng nghiên cứu mới về \textbf{điều kiện biên
của lý thuyết quốc tế hóa}: trong điều kiện nào thì quốc tế hóa nhiều
hơn luôn đồng nghĩa với hiệu quả kém hơn? Luận án đề xuất rằng FIP xuất
hiện khi ba điều kiện đồng thời thỏa mãn: thứ nhất, thể chế ICRV ở mức
thấp nhất; thứ hai, quy mô nền kinh tế quá nhỏ để tạo cơ sở nội địa; và
thứ ba, thiếu hỗ trợ năng lực xuất khẩu từ chính sách. Kết hợp ba điều
kiện này tạo ra một cái bẫy quốc tế hóa cưỡng bức mà không có điểm uốn
để thoát khỏi. Về mặt lý thuyết, phát hiện này dịch chuyển trọng tâm của
tài liệu từ mệnh đề ``thể chế điều tiết quan hệ quốc tế hóa và hiệu
quả'' (Marano et al., 2016) sang một mệnh đề mạnh hơn: ``thể chế quyết
định cả dấu của quan hệ quốc tế hóa và hiệu quả''.

\hypertarget{khung-icrv-suxe1u-chux1ebf-ux111ux1ed9-phuxe2n-loux1ea1i-thux1ec3-chux1ebf-ux111ux1eb7c-thuxf9-cho-khu-vux1ef1c-chuxe2u-uxe1-vuxe0-thuxe1i-buxecnh-dux1b0ux1a1ng}{%
\subsection{5.3.3 Khung ICRV sáu chế độ: Phân loại thể chế đặc thù cho
khu vực châu Á và Thái Bình
Dương}\label{khung-icrv-suxe1u-chux1ebf-ux111ux1ed9-phuxe2n-loux1ea1i-thux1ec3-chux1ebf-ux111ux1eb7c-thuxf9-cho-khu-vux1ef1c-chuxe2u-uxe1-vuxe0-thuxe1i-buxecnh-dux1b0ux1a1ng}}

Đóng góp lý thuyết thứ ba là phát triển và kiểm định \textbf{khung ICRV
sáu chế độ} như một công cụ phân loại thể chế đặc thù cho khu vực châu Á
và Thái Bình Dương, thay thế cho các phân loại nhị phân (phát triển và
mới nổi) hoặc các thang điểm liên tục như chỉ số quản trị toàn cầu trong
tài liệu hiện tại.

Khung ICRV nắm bắt được sự đa dạng thể chế chưa từng được vận hành hóa
đầy đủ: sự khác biệt giữa Singapore (chế độ tiên tiến đổi mới, kết hợp
pháp quyền mạnh và hệ sinh thái đổi mới) với các nền vùng Vịnh (chế độ
tiên tiến tài nguyên, giàu tài nguyên nhưng hệ sinh thái đổi mới khác
biệt), và sự khác biệt giữa Trung Quốc (chế độ trung bình cao, chính phủ
can thiệp mạnh) với Việt Nam (chế độ chuyển đổi thu nhập trung bình
thấp, đang trong quá trình cải cách) và với Campuchia (chế độ đang nổi,
khoảng trống thể chế cao). Khung ICRV không chỉ là một phân loại mô tả
mà còn là một phân loại có sức dự đoán: vị trí và sự tồn tại của điểm
uốn trong quan hệ quốc tế hóa và hiệu quả thay đổi có hệ thống theo ba
vùng của gradient thể chế: chữ U ngược định hình rõ nhất ở vùng giữa
(Nhóm IV, điểm uốn 43,0\%), duỗi thẳng gần tuyến tính ở biên trên nơi
thể chế mạnh đẩy điểm uốn ra mức cường độ xuất khẩu rất cao (Singapore
khoảng 88,6\%), và tan rã ở biên dưới; nghĩa là thể chế càng mạnh, doanh
nghiệp càng duy trì lợi ích từ quốc tế hóa đến mức cường độ cao hơn
trước khi đạt đỉnh, xác nhận giá trị dự đoán của phân loại.

\begin{center}\rule{0.5\linewidth}{0.5pt}\end{center}

\hypertarget{ux111iux1ec1u-kiux1ec7n-ux1ed5n-ux111ux1ecbnh-thux1ec3-chux1ebf-nhux1b0-tiux1ec1n-ux111ux1ec1-ngux1ea7m-cux1ee7a-khung-cdcm}{%
\subsection{5.3.4 Điều kiện ổn định thể chế như tiền đề ngầm của khung
CDCM}\label{ux111iux1ec1u-kiux1ec7n-ux1ed5n-ux111ux1ecbnh-thux1ec3-chux1ebf-nhux1b0-tiux1ec1n-ux111ux1ec1-ngux1ea7m-cux1ee7a-khung-cdcm}}

Đóng góp lý thuyết thứ tư, nổi lên từ bằng chứng Ấn Độ (P9), là xác định
một tiền đề ngầm mà phiên bản ban đầu của khung CDCM chưa nêu rõ: bản
thân sự tồn tại của chữ U ngược giả định một môi trường thể chế ổn định
tương đối trong cửa sổ quan sát. Phát hiện ngưỡng tan biến tại Ấn Độ,
với điểm uốn dịch chuyển từ 61,8\% năm 2014 xuống 40,7\% năm 2022 rồi
mất ý nghĩa độ cong năm 2025 dưới áp lực tích lũy của một chuỗi cú sốc
thể chế cực đoan, cho thấy biến động thể chế đủ lớn có thể hòa tan chứ
không chỉ dịch chuyển điểm tối ưu. Sự đối lập trực tiếp với Trung Quốc,
nơi điểm uốn bền vững cấu trúc qua mười hai năm (Paternoster z = 0,82, p
= 0,412), làm rõ rằng tính bền vững của ngưỡng không phải đặc tính cố
hữu của một chế độ thể chế mà phụ thuộc vào tốc độ và biên độ thay đổi
thể chế tích lũy.

Phát hiện này tinh chỉnh khung CDCM theo hướng đệ quy: thể chế không chỉ
là biến điều tiết quyết định vị trí điểm uốn (tầng tĩnh) mà còn là điều
kiện ổn định nền cho phép một điểm tối ưu được hình thành (tầng động).
Khi thể chế biến động quá nhanh, doanh nghiệp không kịp thiết lập kỳ
vọng ổn định về phần thưởng của quốc tế hóa, và quan hệ quốc tế hóa và
hiệu quả mất cấu trúc bất kể mức thu nhập của nền kinh tế. Hệ quả là
khung CDCM không khẳng định chữ U ngược là quan hệ phổ quát mà là quan
hệ có điều kiện kép: điều kiện về mức chất lượng thể chế (định vị điểm
uốn) và điều kiện về độ ổn định thể chế (cho phép điểm uốn tồn tại).

Bằng chứng Ấn Độ cũng tinh chỉnh một mệnh đề phụ trợ của khung về vai
trò thay thế của hạ tầng số công cộng. Tương tác âm giữa chấp nhận số
Tầng 2 (hạ tầng thanh toán hợp nhất UPI) và cường độ xuất khẩu (hệ số
−4,02, p = 0,004) cho thấy hạ tầng số công cộng chỉ thay thế năng lực số
tư nhân khi hai loại hạ tầng có cùng định hướng mục đích. Hạ tầng thanh
toán nội địa vận hành trên đường ray nội tệ và yêu cầu định danh cư trú,
không thể thay thế hạ tầng tài chính xuyên biên giới mà doanh nghiệp
xuất khẩu yêu cầu. Mệnh đề tinh chỉnh này, rằng tính thay thế giữa hạ
tầng công và năng lực tư là có điều kiện về định hướng mục đích chứ
không phổ quát, là một đóng góp khái niệm bổ sung cho tài liệu về vai
trò của thể chế công trong chuyển đổi số doanh nghiệp.

\hypertarget{tuxedch-hux1ee3p-bux1ed1n-luxfd-thuyux1ebft-nux1ec1n-qua-trux1ee5c-thux1ec3-chux1ebf}{%
\subsection{5.3.5 Tích hợp bốn lý thuyết nền qua trục thể
chế}\label{tuxedch-hux1ee3p-bux1ed1n-luxfd-thuyux1ebft-nux1ec1n-qua-trux1ee5c-thux1ec3-chux1ebf}}

Đóng góp lý thuyết thứ năm mang tính tích hợp: luận án cho thấy bốn lý
thuyết nền của khung CDCM không phải là bốn cách giải thích cạnh tranh
nhau cho cùng một hiện tượng, mà là bốn lớp giải thích vận hành đồng
thời ở các tầng khác nhau và được liên kết bởi một trục chung là chế độ
thể chế. Mô hình quá trình quốc tế hóa của Johanson và Vahlne (1977)
giải thích vòng học hỏi qua đó doanh nghiệp tích lũy tri thức thị trường
nước ngoài và tăng cam kết, cơ chế làm nền cho nhánh đi lên của đường
cong ở các chế độ thể chế trung bình và cao. Lý thuyết dựa trên nguồn
lực (Barney, 1991) giải thích vì sao năng lực công nghệ chiều sâu nâng
mặt bằng năng suất một cách phổ quát: TCI là nguồn lực có giá trị, hiếm,
khó sao chép và khó thay thế, nên hiệu ứng nâng sàn của nó không lệ
thuộc vào chế độ thể chế. Lý thuyết thể chế (North, 1990; Khanna \&
Palepu, 2010) giải thích vì sao điểm uốn dịch chuyển và vì sao quan hệ
tan rã hoặc đảo dấu ở các chế độ thể chế yếu, qua cơ chế chi phí giao
dịch và khoảng trống thể chế. Và lý thuyết thượng tầng quản trị
(Hambrick \& Mason, 1984) giải thích kênh cá nhân hóa qua đặc điểm nhà
quản trị cấp cao, kênh mà luận án cho thấy chỉ vận hành có điều kiện
theo bối cảnh chứ không đồng nhất.

Điểm tích hợp then chốt là cả bốn cơ chế đều bị điều kiện hóa bởi cùng
một biến trục là chế độ thể chế ICRV. Vòng học hỏi Uppsala vận hành trơn
tru khi đệm thể chế đủ dày nhưng bị gián đoạn khi khoảng trống thể chế
quá lớn; lợi thế nguồn lực theo lý thuyết dựa trên nguồn lực được hiện
thực hóa thành hiệu quả ở mọi chế độ nhưng phần thưởng biên của nó thay
đổi theo độ dày của hệ sinh thái hỗ trợ; và kênh quản trị cấp cao phụ
thuộc cấu trúc ngành và thể chế nơi doanh nghiệp hoạt động. Cách đọc
tích hợp này lý giải vì sao một khung đơn lý thuyết không thể nắm bắt
đầy đủ cấu trúc ba vùng: mỗi lý thuyết nền chiếu sáng một đoạn của họ
đường cong, và chỉ khi xếp chồng cả bốn lớp dọc theo trục thể chế thì
bức tranh liên tục từ chữ U ngược đến tan rã và đảo dấu mới hiện ra trọn
vẹn. Đây là một đóng góp về mặt kiến trúc lý thuyết: thay vì lựa chọn
một lý thuyết thắng cuộc, luận án định vị rõ phạm vi áp dụng và điều
kiện biên của từng lý thuyết, và chỉ ra thể chế là biến liên kết chúng
thành một hệ giải thích nhất quán.

\hypertarget{huxe0m-uxfd-phux1b0ux1a1ng-phuxe1p-luux1eadn-cho-nghiuxean-cux1ee9u-phi-tuyux1ebfn-ux111a-bux1ed1i-cux1ea3nh}{%
\subsection{5.3.6 Hàm ý phương pháp luận cho nghiên cứu phi tuyến đa bối
cảnh}\label{huxe0m-uxfd-phux1b0ux1a1ng-phuxe1p-luux1eadn-cho-nghiuxean-cux1ee9u-phi-tuyux1ebfn-ux111a-bux1ed1i-cux1ea3nh}}

Đóng góp lý thuyết thứ sáu mang tính phương pháp luận và có hệ quả vượt
ra ngoài chủ đề quốc tế hóa và hiệu quả. Luận án minh họa cụ thể một vấn
đề mà tài liệu kiểm định quan hệ phi tuyến đã cảnh báo ở mức nguyên tắc
(Haans và cộng sự, 2016; Lind \& Mehlum, 2010): khi một dạng hàm phi
tuyến được ước lượng trên một mẫu gộp nhiều bối cảnh không đồng nhất về
cơ chế nền, hệ số bậc hai gộp và điểm uốn gộp có thể là một tạo tác của
phép trộn chứ không phản ánh bất kỳ bối cảnh đơn lẻ nào. Bằng chứng của
luận án cho thấy ba dạng quan hệ khác nhau về bản chất (chữ U ngược sắc
nét, gần tuyến tính dương, âm đơn điệu) cùng tồn tại trong một mẫu khu
vực, và việc gộp chúng tạo ra một tín hiệu nén và bất ổn đúng như các
nghiên cứu nhân rộng quy mô lớn ghi nhận (Pisani và cộng sự, 2020; Berry
\& Kaul, 2016).

Hệ quả phương pháp luận trực tiếp là: trong nghiên cứu phi tuyến đa bối
cảnh, việc phân tầng theo một biến điều tiết có cơ sở lý thuyết trước
khi ước lượng dạng hàm không phải là một bước tinh chỉnh tùy chọn mà là
điều kiện cần để diễn giải đúng. Khung phân loại thể chế ICRV mà luận án
phát triển có thể coi là một ví dụ về cách vận hành hóa nguyên tắc này:
thay vì hỏi ``dạng hàm trung bình là gì'', câu hỏi đúng là ``dạng hàm
thay đổi như thế nào dọc theo chiều điều tiết có cơ sở lý thuyết''. Cách
tiếp cận này cũng nhất quán với quan điểm về nhận dạng nhân quả tích lũy
qua một khối nghiên cứu (Shaver, 2020): bằng chứng vững không đến từ một
ước lượng gộp đơn lẻ mà từ sự hội tụ của nhiều nguồn dữ liệu và phương
pháp với các điểm yếu không trùng nhau, đúng như tam giác hóa giữa phân
tích tổng hợp P6 và ước lượng gộp toàn mẫu P7 đã trình bày ở Mục 5.2.5.
Đây là một đóng góp mà giá trị của nó vượt ra ngoài chủ đề luận án, áp
dụng được cho bất kỳ lĩnh vực nào mà quan hệ phi tuyến được kiểm định
trên dữ liệu gộp từ nhiều bối cảnh thể chế khác biệt.

\hypertarget{ux111ux1ecbnh-vux1ecb-trong-chux1b0ux1a1ng-truxecnh-nghux1ecb-sux1ef1-nux1ec1n-tux1ea3ng-vi-muxf4-cux1ee7a-kinh-doanh-quux1ed1c-tux1ebf}{%
\subsection{5.3.7 Định vị trong chương trình nghị sự nền tảng vi mô của
kinh doanh quốc
tế}\label{ux111ux1ecbnh-vux1ecb-trong-chux1b0ux1a1ng-truxecnh-nghux1ecb-sux1ef1-nux1ec1n-tux1ea3ng-vi-muxf4-cux1ee7a-kinh-doanh-quux1ed1c-tux1ebf}}

Đóng góp lý thuyết thứ bảy nằm ở cách luận án bắc cầu giữa cấp độ vĩ mô
và cấp độ vi mô của phân tích, một hướng đang được giới kinh doanh quốc
tế quan tâm mạnh dưới tên gọi nền tảng vi mô. Khung CDCM của luận án vốn
đã phân tầng tường minh thành tầng vĩ mô (chế độ thể chế ICRV), tầng tổ
chức (năng lực công nghệ và số), và tầng cá nhân (đặc điểm nhà quản trị
cấp cao). Chính tương tác giữa tầng vĩ mô và tầng cá nhân, tức câu hỏi
bối cảnh thể chế định hình hay ràng buộc nhận thức và quyết định của nhà
quản trị như thế nào, là nội dung cốt lõi của chương trình nghị sự nền
tảng vi mô.

Lý thuyết thượng tầng quản trị (Hambrick \& Mason, 1984) từ lâu đã đặt
nền cho cách đọc này khi lập luận rằng kết quả tổ chức là sự phản chiếu
nhận thức và giá trị của nhà quản trị; bước tiến gần đây của Grego và
cộng sự (2025) bổ sung chiều nhận thức cảm tính, cho thấy quyết định
quốc tế hóa không quy giản về tính toán duy lý lạnh lùng. Luận án đóng
góp vào mạch này một bằng chứng đặc thù: hiệu ứng của đặc điểm nhà quản
trị không phải một hằng số phổ quát mà thay đổi theo thước đo hiệu quả
và bối cảnh thể chế, biểu hiện rõ ở việc hệ số của biến quản trị nữ đổi
dấu giữa thước đo lợi nhuận đơn quốc gia (dương, trong chương sách Ấn
Độ) và thước đo năng suất gộp đa quốc gia (âm, trong P7 và P10). Phát
hiện này hàm ý rằng việc đưa đặc trưng cá nhân vào mô hình quan hệ quốc
tế hóa và hiệu quả mà không kiểm soát đồng thời bối cảnh thể chế và
thước đo dễ dẫn đến kết luận mâu thuẫn, qua đó củng cố lời kêu gọi của
tài liệu nền tảng vi mô về thiết kế đa cấp độ. Định vị này không chỉ làm
rõ giá trị của tầng cá nhân trong khung CDCM mà còn kết nối luận án với
một trong những hướng lý thuyết thời sự nhất của ngành. Cách định vị này
nhất quán với hướng dẫn của giới biên tập kinh doanh quốc tế về đóng góp
lý thuyết: Santangelo và Verbeke (2022) xác định việc làm rõ các điều
kiện biên của một lý thuyết và việc đánh giá thận trọng tính hợp lệ khi
áp dụng lý thuyết qua các cấp độ phân tích khác nhau là hai con đường
cốt lõi để tạo đóng góp lý thuyết cho kinh doanh quốc tế. Luận án thực
hiện đúng cả hai: vận hành hóa điều kiện biên thể chế bằng khung ICRV
thay vì giả định một quy luật phổ quát, và nối tầng vĩ mô thể chế với
tầng cá nhân nhà quản trị trong một khung đa cấp được lập luận tường
minh thay vì chuyển dịch lý thuyết giữa các cấp độ một cách máy móc.

\hypertarget{ux111ux1ec1-xuux1ea5t-chuxednh-suxe1ch-vuxe0-quux1ea3n-trux1ecb}{%
\section{5.4 Đề xuất chính sách và quản
trị}\label{ux111ux1ec1-xuux1ea5t-chuxednh-suxe1ch-vuxe0-quux1ea3n-trux1ecb}}

Cần phân định rõ tính chất của phần này so với các phần trước. Chương 4
và các mục 5.1 đến 5.3 mang tính thực chứng: chúng mô tả, giải thích và
dự đoán quan hệ giữa quốc tế hóa và hiệu quả dựa trên bằng chứng. Các đề
xuất dưới đây mang tính quy chuẩn: chúng là suy luận chính sách rút ra
từ các phát hiện thực chứng nhưng không đồng nhất với chúng, vì việc
chuyển từ một quy luật quan sát được sang một khuyến nghị hành động luôn
đòi hỏi thêm các giả định về mục tiêu, chi phí và năng lực thực thi. Vì
luận án dựa trên thiết kế lát cắt ngang và phần lớn suy luận là về liên
hệ chứ chưa phải nhân quả đã chứng minh trực tiếp (Mục 5.5.3.2), các đề
xuất nên được đọc như định hướng ưu tiên có cơ sở thực nghiệm, cần được
kiểm chứng và hiệu chỉnh trong từng bối cảnh triển khai, chứ không phải
như công thức nhân quả chắc chắn.

\hypertarget{ux111ux1ec1-xuux1ea5t-cho-nhuxe0-houx1ea1ch-ux111ux1ecbnh-chuxednh-suxe1ch-nuxe2ng-cao-tci-vuxe0-cux1ea3i-thiux1ec7n-thux1ec3-chux1ebf}{%
\subsection{5.4.1 Đề xuất cho nhà hoạch định chính sách: Nâng cao TCI và
cải thiện thể
chế}\label{ux111ux1ec1-xuux1ea5t-cho-nhuxe0-houx1ea1ch-ux111ux1ecbnh-chuxednh-suxe1ch-nuxe2ng-cao-tci-vuxe0-cux1ea3i-thiux1ec7n-thux1ec3-chux1ebf}}

Kết quả về TCI có hàm ý chính sách rõ ràng: \textbf{đầu tư vào năng lực
công nghệ của doanh nghiệp} là can thiệp chính sách có thể nâng cao lợi
ích từ quốc tế hóa, đặc biệt trong các bối cảnh ICRV trung bình như Việt
Nam và Trung Quốc. Chính sách hỗ trợ doanh nghiệp tiếp cận công nghệ
nước ngoài (thông qua chương trình kết nối với MNCs, giảm thuế nhập khẩu
công nghệ), khuyến khích R\&D, và hỗ trợ chứng nhận chất lượng quốc tế
sẽ nâng cao TCI của toàn bộ doanh nghiệp trong nền kinh tế, từ đó dịch
chuyển đường quốc tế hóa và hiệu quả lên trên và có thể mở rộng điểm uốn
về phía phải.

Gradient ICRV gợi ý rằng \textbf{cải thiện chất lượng thể chế} (pháp
quyền, bảo vệ quyền sở hữu trí tuệ, giảm chi phí tuân thủ thương mại) là
can thiệp dài hạn quan trọng nhất. Doanh nghiệp trong nền kinh tế thể
chế tốt hơn có thể quốc tế hóa đến mức độ cao hơn trước khi chi phí vượt
qua lợi ích, điều này có nghĩa là cải cách thể chế không chỉ có lợi về
mặt kinh doanh nội địa mà còn mở rộng không gian lợi ích từ xuất khẩu.

Đối với Việt Nam cụ thể: điểm uốn dịch chuyển từ 46,2\% (2009) xuống
39,3\% (2015) là tín hiệu đáng lo ngại. Nếu điểm uốn tiếp tục dịch trái,
điều đó hàm ý rằng lợi ích từ xuất khẩu đạt đỉnh sớm hơn và suy giảm
nhanh hơn, có thể do cấu trúc xuất khẩu tập trung vào gia công với biên
lợi nhuận mỏng. Chính sách nâng cao giá trị gia tăng trong xuất khẩu
(moving up the value chain) kết hợp với tăng cường TCI có thể giúp đảo
ngược xu hướng này.

\hypertarget{ux111ux1ec1-xuux1ea5t-cho-nhuxe0-quux1ea3n-trux1ecb-doanh-nghiux1ec7p-ux111iux1ec3m-uux1ed1n-theo-icrv}{%
\subsection{5.4.2 Đề xuất cho nhà quản trị doanh nghiệp: Điểm uốn theo
ICRV}\label{ux111ux1ec1-xuux1ea5t-cho-nhuxe0-quux1ea3n-trux1ecb-doanh-nghiux1ec7p-ux111iux1ec3m-uux1ed1n-theo-icrv}}

Kết quả thực nghiệm cung cấp \textbf{hướng dẫn thực tế} cho nhà quản trị
về mức độ quốc tế hóa tối ưu theo bối cảnh:

\begin{itemize}
\item
  \textbf{Doanh nghiệp Singapore (Nhóm I)}: không quan sát thấy ngưỡng
  bất lợi của quốc tế hóa trong miền dữ liệu (điểm uốn ước lượng ≈
  88,6\% nhưng không định vị tin cậy) cho thấy không gian quốc tế hóa
  rất rộng. Chiến lược mở rộng quốc tế mạnh mẽ (high-commitment
  internationalization) phù hợp với bối cảnh này, đặc biệt khi kết hợp
  với DAI cao.
\item
  \textbf{Doanh nghiệp Trung Quốc (Nhóm III)}: điểm uốn TP ≈ 48,78\% ổn
  định qua thập kỷ 2012--2024, gợi ý rằng chiến lược xuất khẩu tối ưu là
  duy trì tỷ lệ FSTS trong khoảng 40--50\%, kết hợp giữa thị trường nội
  địa khổng lồ và hoạt động xuất khẩu có chọn lọc.
\item
  \textbf{Doanh nghiệp Việt Nam (Nhóm IV)}: điểm uốn thấp hơn
  (39,3--46,2\%), hàm ý rằng nhà quản trị cần thận trọng hơn khi đẩy tỷ
  lệ xuất khẩu vượt quá 40\% tổng doanh thu. Thay vào đó, \textbf{nâng
  cao TCI} trước khi mở rộng quốc tế sẽ hiệu quả hơn là tăng FSTS thuần
  túy.
\item
  \textbf{Doanh nghiệp SIDS (Nhóm VI)}: kết quả FIP là cảnh báo nghiêm
  túc, đẩy mạnh xuất khẩu trong bối cảnh thiếu năng lực và hỗ trợ thể
  chế không cải thiện mà còn có thể làm xấu đi hiệu quả hoạt động. Doanh
  nghiệp SIDS nên ưu tiên xây dựng năng lực nội địa và nâng cao TCI
  trước khi mở rộng xuất khẩu đáng kể.
\end{itemize}

\hypertarget{ux111ux1ec1-xuux1ea5t-ux111ux1eb7c-biux1ec7t-cho-sids-truxe1nh-bux1eaby-quux1ed1c-tux1ebf-huxf3a-bux1eaft-buux1ed9c}{%
\subsection{5.4.3 Đề xuất đặc biệt cho SIDS: Tránh Bẫy Quốc Tế Hóa Bắt
Buộc}\label{ux111ux1ec1-xuux1ea5t-ux111ux1eb7c-biux1ec7t-cho-sids-truxe1nh-bux1eaby-quux1ed1c-tux1ebf-huxf3a-bux1eaft-buux1ed9c}}

Phát hiện hình phạt quốc tế hóa cưỡng bức tại bảy nền kinh tế đảo nhỏ
Thái Bình Dương có hàm ý chính sách cụ thể và cấp bách. Khác với năm bối
cảnh còn lại, nơi quan hệ có dạng chữ U ngược với một ngưỡng tối ưu, ở
nhóm đảo nhỏ hệ số nhánh cường độ xuất khẩu mang dấu âm và có ý nghĩa
cao (\(\beta = -1{,}339\) trên thang căn giữa), nghĩa là mọi mức tăng
cường độ xuất khẩu đều đi kèm năng suất thấp hơn, không tồn tại điểm uốn
để doanh nghiệp hướng tới. Hiện nay nhiều chương trình phát triển quốc
tế và viện trợ thương mại cho các nền đảo nhỏ tập trung vào khuyến khích
xuất khẩu như một chiến lược tăng trưởng tự thân. Kết quả này đặt câu
hỏi nghiêm túc về hiệu quả của cách tiếp cận đó khi ba điều kiện cấu
trúc nền tảng của một quan hệ chữ U ngược, gồm thị trường nội địa đủ
lớn, chi phí thương mại chấp nhận được, và hỗ trợ thể chế tối thiểu,
đồng thời bị vi phạm.

Thay vào đó, hàm ý chính sách là: \textbf{không phải xuất khẩu nhiều
hơn, mà phải xuất khẩu tốt hơn}. Cụ thể, ưu tiên nên là, thứ nhất, xây
dựng hạ tầng vận tải và kết nối để giảm chi phí cố định của xuất khẩu,
vốn là ràng buộc gay gắt nhất với các nền đảo phân tán địa lý; thứ hai,
phát triển năng lực doanh nghiệp thông qua năng lực công nghệ chiều sâu
chứ không chỉ hiện diện số bề mặt, vì bằng chứng cho thấy năng lực công
nghệ nâng mặt bằng năng suất ngay cả ở nhóm đảo nhỏ nhưng không tự nó
đảo được độ dốc âm; và thứ ba, thiết kế chương trình hỗ trợ xuất khẩu có
chọn lọc, tập trung vào doanh nghiệp đã đủ năng lực để hưởng lợi, thay
vì khuyến khích đại trà. Nói tóm lại, tránh bẫy quốc tế hóa cưỡng bức
bằng cách bảo đảm rằng các can thiệp chính sách xây dựng năng lực hấp
thụ trước khi thúc đẩy mở rộng thị trường. Một hệ quả thực tiễn quan
trọng là các chỉ số đánh giá hiệu quả viện trợ thương mại cho nhóm này
không nên lấy mức tăng kim ngạch xuất khẩu làm thước đo thành công
chính, vì trong điều kiện hình phạt quốc tế hóa cưỡng bức, mở rộng xuất
khẩu mà không kèm nâng cấp năng lực có thể đồng nghĩa với suy giảm năng
suất.

\hypertarget{huxe0m-uxfd-chuxednh-suxe1ch-trong-bux1ed1i-cux1ea3nh-phuxe1t-triux1ec3n-ux111ux1b0ux1a1ng-ux111ux1ea1i}{%
\subsection{5.4.4 Hàm ý chính sách trong bối cảnh phát triển đương
đại}\label{huxe0m-uxfd-chuxednh-suxe1ch-trong-bux1ed1i-cux1ea3nh-phuxe1t-triux1ec3n-ux111ux1b0ux1a1ng-ux111ux1ea1i}}

Đặt các đề xuất trên vào bối cảnh phát triển vĩ mô gần nhất làm rõ tính
cấp bách và phạm vi áp dụng của chúng. Thách thức việc làm là điểm khởi
đầu: với hơn 230 triệu người dự kiến bước vào tuổi lao động ở các thị
trường biên giai đoạn 2025--2035, năng lực tạo việc làm năng suất cao
của khu vực doanh nghiệp trở thành ưu tiên phát triển hàng đầu (World
Bank, 2026c). Phát hiện của luận án rằng quốc tế hóa chỉ tạo phần thưởng
năng suất khi đi kèm năng lực công nghệ và đệm thể chế đủ mạnh hàm ý
rằng chương trình nghị sự việc làm cần đặt nâng cấp năng lực doanh
nghiệp làm trụ cột, thay vì chỉ theo đuổi mở rộng thương mại như mục
tiêu tự thân.

Thứ hai, bằng chứng vĩ mô cho thấy không tồn tại một con đường phát
triển phổ quát: các nền kinh tế biên thành công đã đi những lối khác
nhau, từ đầu tư năng lượng và khoáng sản, đến chế tạo giá trị gia tăng,
đến dịch vụ và du lịch (World Bank, 2026c). Điều này củng cố nguyên tắc
chính sách phái sinh từ khung ICRV: gói can thiệp phải được hiệu chỉnh
theo chế độ thể chế và mô hình tăng trưởng đặc thù của từng nhóm, thay
vì áp dụng một công thức chung. Với Việt Nam (Nhóm IV), con đường chế
tạo giá trị gia tăng và dịch chuyển lên chuỗi giá trị là phù hợp với cả
bằng chứng cấp doanh nghiệp của luận án lẫn kinh nghiệm của các thị
trường biên thành công.

Thứ ba, đối với các nền kinh tế tài nguyên dẫn dắt (Nhóm II), cú sốc giá
hàng hóa năm 2026 vừa là cơ hội thu nhập ngắn hạn vừa là lời nhắc về rủi
ro phụ thuộc tài nguyên (World Bank, 2026d). Hàm ý chính sách là tận
dụng thặng dư tô lợi giai đoạn giá cao để đầu tư vào đa dạng hóa năng
lực công nghệ và phi tài nguyên, đúng như chiến lược Vision của các nền
GCC, thay vì củng cố cấu trúc nhà nước tô vốn san bằng động lực năng
suất doanh nghiệp.

Thứ tư, trong kỷ nguyên trí tuệ nhân tạo, hạ tầng số công và quản trị dữ
liệu trở thành điều kiện nền tảng để năng lực số cấp doanh nghiệp được
chuyển hóa thành giá trị (World Bank, 2026e). Phát hiện của luận án rằng
phần thưởng của áp dụng số phụ thuộc bối cảnh thể chế và tầng năng lực
hàm ý rằng chính sách số hóa cần phân tầng: ở các nền đang chuyển đổi,
ưu tiên là phổ cập hạ tầng số nền tảng và kỹ năng số cơ bản; ở các nền
tiên tiến đã bão hòa Tầng 1, trọng tâm dịch sang năng lực số chiều sâu
và quản trị dữ liệu, AI. Cách tiếp cận phân tầng này tránh được cả hai
sai lầm: đầu tư quá sớm vào công nghệ cao ở nơi chưa có nền tảng hấp
thụ, và dừng lại ở hiện diện số bề mặt ở nơi đã sẵn sàng cho tầng năng
lực cao hơn.

\hypertarget{khung-can-thiux1ec7p-phuxe2n-tux1ea7ng-theo-chux1ebf-ux111ux1ed9-thux1ec3-chux1ebf-icrv}{%
\subsection{5.4.5 Khung can thiệp phân tầng theo chế độ thể chế
ICRV}\label{khung-can-thiux1ec7p-phuxe2n-tux1ea7ng-theo-chux1ebf-ux111ux1ed9-thux1ec3-chux1ebf-icrv}}

Tổng hợp các đề xuất rời rạc theo từng bối cảnh thành một nguyên lý can
thiệp thống nhất làm rõ giá trị ứng dụng của khung ICRV: không tồn tại
một gói chính sách quốc tế hóa phổ quát, và bản thân việc áp dụng máy
móc bài học từ một chế độ thể chế sang chế độ khác là sai lầm chiến lược
căn bản. Bằng chứng ba vùng của luận án dịch trực tiếp thành ba định
hướng can thiệp khác biệt về bản chất.

Đối với chế độ thể chế mạnh nhất (Nhóm I, các nền như Singapore và Nhật
Bản), nơi quan hệ quốc tế hóa và hiệu quả gần tuyến tính dương và cơ chế
vận hành qua biên tham gia, ưu tiên chính sách không phải là khuyến
khích doanh nghiệp xuất khẩu nhiều hơn mà là hạ thấp ngưỡng tham gia
xuất khẩu cho nhóm doanh nghiệp đã có năng lực nhưng chưa vượt ngưỡng.
Vì phần thưởng năng suất tập trung ở việc có hay không xuất khẩu, can
thiệp hiệu quả nhất là giảm chi phí cố định gia nhập thị trường nước
ngoài và kết nối doanh nghiệp có năng lực với cơ hội quốc tế hóa. Đối
với các nền đã bão hòa hiện diện số bề mặt, trọng tâm số hóa dịch sang
năng lực số chiều sâu và quản trị dữ liệu, thay vì phổ cập hạ tầng nền
tảng vốn đã đạt trần.

Đối với chế độ chuyển đổi bậc thấp--trung (Nhóm IV, gồm Việt Nam và Ấn
Độ), nơi chữ U ngược sắc nét với điểm uốn quanh 40--46\%, nguyên lý can
thiệp là nâng cao năng lực công nghệ trước khi đẩy cường độ quốc tế hóa.
Vì năng lực công nghệ nâng mặt bằng năng suất và có thể dịch điểm uốn về
phía phải, đầu tư vào tiếp cận công nghệ nước ngoài, nghiên cứu và phát
triển, và chứng nhận chất lượng quốc tế sẽ mở rộng không gian quốc tế
hóa có lợi. Riêng với Việt Nam, việc điểm uốn dịch từ 46,2\% (2009)
xuống 39,3\% (2015) là tín hiệu cảnh báo về cấu trúc xuất khẩu gia công
biên mỏng; chính sách dịch chuyển lên chuỗi giá trị kết hợp tăng cường
năng lực công nghệ là cần thiết để đảo ngược xu hướng dịch trái của điểm
uốn. Đối với chế độ này, bài học Ấn Độ cũng cảnh báo rằng cải cách thể
chế dồn dập trong thời gian ngắn, dù có ý định tốt, có thể tạm thời phá
vỡ cơ chế chuyển hóa quốc tế hóa thành hiệu quả, nên trình tự và tốc độ
cải cách cần được cân nhắc.

Đối với chế độ thể chế yếu nhất (Nhóm V và VI, các nền đang nổi và đảo
nhỏ), nơi quan hệ tan rã hoặc đảo sang âm đơn điệu, nguyên lý can thiệp
đảo ngược hoàn toàn so với khuyến nghị thông thường: ưu tiên không phải
xuất khẩu nhiều hơn mà xuất khẩu tốt hơn. Phát hiện gánh nặng quốc tế
hóa bắt buộc tại SIDS đặt câu hỏi nghiêm túc về các chương trình viện
trợ thương mại tập trung thúc đẩy xuất khẩu đại trà ở những nền thiếu
nền tảng năng lực. Trình tự đúng là xây dựng hạ tầng vận tải và kết nối
để giảm chi phí cố định xuất khẩu, phát triển năng lực công nghệ chiều
sâu thay vì chỉ hiện diện số bề mặt, và thiết kế hỗ trợ xuất khẩu có
chọn lọc cho doanh nghiệp đã đủ năng lực hưởng lợi. Nguyên lý xuyên suốt
ba vùng là xây dựng năng lực hấp thụ trước khi thúc đẩy mở rộng thị
trường, một thứ tự ưu tiên mà khung ICRV giúp hiệu chỉnh cụ thể theo
từng chế độ.

\hypertarget{quy-tux1eafc-quyux1ebft-ux111ux1ecbnh-cho-nhuxe0-quux1ea3n-trux1ecb-doanh-nghiux1ec7p-theo-vux1ecb-truxed-thux1ec3-chux1ebf}{%
\subsection{5.4.6 Quy tắc quyết định cho nhà quản trị doanh nghiệp theo
vị trí thể
chế}\label{quy-tux1eafc-quyux1ebft-ux111ux1ecbnh-cho-nhuxe0-quux1ea3n-trux1ecb-doanh-nghiux1ec7p-theo-vux1ecb-truxed-thux1ec3-chux1ebf}}

Khung ba vùng cũng dịch thành một quy tắc quyết định gọn cho nhà quản
trị doanh nghiệp, thay vì một khuyến nghị chung chung về ``quốc tế hóa
nhiều hơn''. Bước thứ nhất là định vị nền kinh tế của doanh nghiệp trong
sáu chế độ ICRV, vì chính vị trí này quyết định dạng quan hệ giữa cường
độ xuất khẩu và năng suất mà doanh nghiệp đang vận hành bên trong. Bước
thứ hai là đọc tư thế chiến lược tương ứng với vùng đó. Ở chế độ thể chế
mạnh, nơi quan hệ gần tuyến tính dương, doanh nghiệp đã có năng lực nên
ưu tiên vượt qua ngưỡng tham gia xuất khẩu vì phần thưởng nằm ở việc có
mặt trên thị trường nước ngoài hơn là ở mức độ thâm dụng xuất khẩu. Ở
chế độ chuyển đổi, nơi chữ U ngược sắc nét với điểm uốn quanh 40--46\%,
doanh nghiệp cần so sánh cường độ xuất khẩu hiện tại của mình với ngưỡng
này: nếu còn dưới ngưỡng, dư địa mở rộng có lợi vẫn còn; nếu đã vượt
ngưỡng, ràng buộc không còn là tiếp cận thị trường mà là năng lực điều
phối, và đầu tư đúng là nâng năng lực công nghệ để dịch ngưỡng về phía
phải trước khi đẩy thêm cường độ. Ở chế độ thể chế yếu, nơi quan hệ tan
rã hoặc đảo dấu, tư thế đúng là xuất khẩu có chọn lọc dựa trên năng lực
thực có thay vì mở rộng đại trà.

Một cảnh báo diễn giải đi kèm quy tắc này. Bằng chứng về đặc trưng quản
lý cho thấy hệ số người quản lý cấp cao là nữ mang dấu âm trên toàn mẫu,
nhưng luận án đã lập luận rằng đây là quy luật mô tả phản ánh chọn lọc
cấu thành chứ không phải bằng chứng nhân quả về hiệu quả của lãnh đạo.
Do đó nhà quản trị và nhà đầu tư không nên đọc hệ số này như một chỉ báo
về năng lực lãnh đạo; suy luận đúng đòi hỏi phân tích nhận biết chọn lọc
trong nội bộ từng nền kinh tế. Quy tắc quyết định vì vậy đặt trọng tâm
vào vị trí thể chế và cường độ xuất khẩu so với điểm uốn, là những đại
lượng có cơ sở nhân quả cấu trúc rõ hơn, thay vì vào các đặc trưng nhân
khẩu của người quản lý.

\hypertarget{ux111ux1ec1-xuux1ea5t-cho-cuxe1c-tux1ed5-chux1ee9c-phuxe1t-triux1ec3n-vuxe0-tuxe0i-trux1ee3-quux1ed1c-tux1ebf}{%
\subsection{5.4.7 Đề xuất cho các tổ chức phát triển và tài trợ quốc
tế}\label{ux111ux1ec1-xuux1ea5t-cho-cuxe1c-tux1ed5-chux1ee9c-phuxe1t-triux1ec3n-vuxe0-tuxe0i-trux1ee3-quux1ed1c-tux1ebf}}

Bên cạnh nhà hoạch định chính sách quốc gia và nhà quản trị doanh
nghiệp, một nhóm đối tượng thứ ba có khả năng tác động trực tiếp đến quỹ
đạo quốc tế hóa của doanh nghiệp ở các nền kinh tế đang phát triển là
các tổ chức phát triển và tài trợ quốc tế. Phát hiện của luận án có hàm
ý cụ thể và đôi khi đi ngược lại thực hành phổ biến của nhóm này. Hàm ý
thứ nhất liên quan đến thước đo thành công của các chương trình hỗ trợ
thương mại. Trong điều kiện gánh nặng quốc tế hóa bắt buộc ở nhóm đảo
nhỏ, việc lấy mức tăng kim ngạch xuất khẩu làm chỉ tiêu thành công chính
là một sai lệch khuyến khích nghiêm trọng: nó có thể thưởng cho chính
các can thiệp đẩy doanh nghiệp vào vùng năng suất thấp hơn. Đề xuất là
chuyển trọng tâm đánh giá từ chỉ số quy mô thương mại sang chỉ số nâng
cấp năng lực, chẳng hạn tỷ lệ doanh nghiệp đạt chứng nhận chất lượng
quốc tế, cường độ chi cho nghiên cứu và phát triển, và mức cải thiện
năng suất lao động, là những đại lượng gắn với nhánh nâng mặt bằng bền
vững thay vì với cường độ xuất khẩu thuần túy.

Hàm ý thứ hai liên quan đến trình tự can thiệp. Bằng chứng rằng năng lực
công nghệ nâng mặt bằng năng suất ở mọi chế độ thể chế kể cả nhóm đảo
nhỏ, nhưng không tự nó đảo được độ dốc âm, gợi ý một nguyên tắc trình tự
cho thiết kế chương trình: xây dựng năng lực hấp thụ trước, thúc đẩy mở
rộng thị trường sau. Các tổ chức tài trợ nên ưu tiên cấu phần phát triển
năng lực công nghệ chiều sâu và hạ tầng kết nối giảm chi phí cố định
xuất khẩu, trước khi hoặc song song với các cấu phần xúc tiến xuất khẩu,
thay vì đặt xúc tiến xuất khẩu làm cấu phần dẫn dắt. Hàm ý thứ ba liên
quan đến phân bổ nguồn lực theo chế độ thể chế: vì cùng một đồng tài trợ
tạo ra phần thưởng năng suất rất khác nhau tùy chế độ thể chế của nền
kinh tế thụ hưởng, các tổ chức quốc tế có thể nâng hiệu quả phân bổ bằng
cách hiệu chỉnh gói can thiệp theo vị trí ICRV thay vì áp dụng một mô
hình chuẩn hóa xuyên quốc gia. Cách tiếp cận phân tầng này nhất quán với
khung năng lực thể chế cho thực thi chính sách mà tài liệu phát triển
gần đây đề xuất (Kim và cộng sự, 2026), theo đó năng lực thực thi của tổ
chức công là điều kiện trung gian quyết định liệu một can thiệp chính
sách có chuyển hóa thành kết quả phát triển hay không.

\hypertarget{phuxe2n-biux1ec7t-nux1ed9i-bux1ed9-nhuxf3m-tiuxean-tiux1ebfn-chux1ebf-ux111ux1ed9-ux111ux1ed5i-mux1edbi-so-vux1edbi-chux1ebf-ux111ux1ed9-tuxe0i-nguyuxean}{%
\subsection{5.4.8 Phân biệt nội bộ nhóm tiên tiến: chế độ đổi mới so với
chế độ tài
nguyên}\label{phuxe2n-biux1ec7t-nux1ed9i-bux1ed9-nhuxf3m-tiuxean-tiux1ebfn-chux1ebf-ux111ux1ed9-ux111ux1ed5i-mux1edbi-so-vux1edbi-chux1ebf-ux111ux1ed9-tuxe0i-nguyuxean}}

Một sắc thái chính sách mà phân loại nhị phân truyền thống bỏ sót là sự
phân hóa trong nội bộ nhóm thu nhập cao giữa chế độ tiên tiến đổi mới và
chế độ tiên tiến tài nguyên. Bằng chứng mô tả của luận án cho thấy hai
phân nhóm cùng mức thu nhập nhưng đối lập về cấu trúc năng lực, và độ
phân tán năng suất rất hẹp ở các nền tài nguyên lớn là biểu hiện của
kinh tế nhà nước tô nơi phân phối lại từ tô tài nguyên thu hẹp khoảng
cách năng suất giữa doanh nghiệp. Hàm ý chính sách cho chế độ tiên tiến
tài nguyên (Nhóm II) khác hẳn với chế độ tiên tiến đổi mới (Nhóm I): rủi
ro không phải là thiếu nguồn lực mà là cấu trúc nhà nước tô san bằng
động lực năng suất, làm yếu chính cơ chế chọn lọc theo năng lực mà quan
hệ quốc tế hóa và hiệu quả tích cực dựa vào. Đề xuất là tận dụng thặng
dư tô trong các giai đoạn giá hàng hóa cao để đầu tư vào đa dạng hóa
năng lực công nghệ phi tài nguyên, đúng theo định hướng các chiến lược
tầm nhìn của khối, thay vì củng cố cấu trúc phân phối lại làm xói mòn
khuyến khích nâng cấp năng lực doanh nghiệp. Phân biệt này là một ví dụ
cụ thể về giá trị ứng dụng của khung ICRV sáu chế độ so với phân loại
nhị phân: hai nền cùng được xếp là phát triển nhưng đòi hỏi hai gói
chính sách quốc tế hóa khác nhau về bản chất.

\hypertarget{hux1ea1n-chux1ebf-vuxe0-hux1b0ux1edbng-nghiuxean-cux1ee9u-tux1b0ux1a1ng-lai}{%
\section{5.5 Hạn chế và hướng nghiên cứu tương
lai}\label{hux1ea1n-chux1ebf-vuxe0-hux1b0ux1edbng-nghiuxean-cux1ee9u-tux1b0ux1a1ng-lai}}

\hypertarget{hux1ea1n-chux1ebf-vux1ec1-ux111o-lux1b0ux1eddng-dai}{%
\subsection{5.5.1 Hạn chế về đo lường
DAI}\label{hux1ea1n-chux1ebf-vux1ec1-ux111o-lux1b0ux1eddng-dai}}

Hạn chế quan trọng nhất của luận án liên quan đến đo lường DAI. Dữ liệu
WBES chỉ cung cấp số lượng item hạn chế để xây dựng chỉ số DAI, và các
item này thay đổi qua các thế hệ lược đồ WBES (PICS3, Standardized,
BREADY/BEE). Tại Việt Nam (P3), chỉ Tầng 1 DAI (website, email, online
sales) được sử dụng do hạn chế của các làn sóng 2009 và 2015, trong khi
Tầng 2 (e-payment, e-supplier) chỉ có từ 2023. Điều này tạo ra khả năng
đo lường không đầy đủ, đặc biệt là không nắm bắt được chiều sâu năng lực
số ở bối cảnh ICRV thấp.

Quan trọng hơn, tất cả các đo lường DAI đều là \textbf{tĩnh}, tức tại
một thời điểm khảo sát, không thể nắm bắt tính động của việc cập nhật và
tái cấu hình năng lực số theo thời gian. Điều này có nghĩa là các phát
hiện về DAI trong luận án phản ánh mức độ chấp nhận số tại một thời điểm
chứ không phải năng lực số năng động theo nghĩa của lý thuyết năng lực
động (Teece et al., 1997).

\hypertarget{hux1ea1n-chux1ebf-vux1ec1-thiuxean-lux1ec7ch-chux1ecdn-lux1ef1a-trong-mux1eabu-phux1ee5-doanh-nghiux1ec7p-xuux1ea5t-khux1ea9u}{%
\subsection{5.5.2 Hạn chế về thiên lệch chọn lựa trong mẫu phụ doanh
nghiệp xuất
khẩu}\label{hux1ea1n-chux1ebf-vux1ec1-thiuxean-lux1ec7ch-chux1ecdn-lux1ef1a-trong-mux1eabu-phux1ee5-doanh-nghiux1ec7p-xuux1ea5t-khux1ea9u}}

Các phân tích sử dụng mẫu chỉ gồm doanh nghiệp xuất khẩu (ví dụ: mẫu phụ
xuất khẩu Singapore N = 84, mẫu phụ chỉ doanh nghiệp xuất khẩu SIDS N =
26) đối mặt với rủi ro thiên lệch chọn lựa nghiêm trọng: doanh nghiệp
xuất khẩu không phải là mẫu ngẫu nhiên từ tổng thể doanh nghiệp mà có
khả năng là các doanh nghiệp đã có một số lợi thế ban đầu về năng suất,
nguồn lực và quan hệ kết nối. Kết quả từ chỉ doanh nghiệp xuất khẩu
không thể tổng quát hóa trực tiếp cho toàn bộ doanh nghiệp, và có thể
đánh giá thấp hoặc đánh giá cao các hiệu ứng tùy thuộc vào chiều của
chọn lọc.

Phân tích Heckman hai bước (hoặc các mô hình hiệu ứng xử lý) trong các
nghiên cứu tương lai có thể khắc phục phần nào hạn chế này bằng cách mô
hình hóa quá trình lựa chọn tham gia xuất khẩu như giai đoạn thứ nhất
trước khi ước lượng hiệu quả trong giai đoạn thứ hai.

\hypertarget{hux1ea1n-chux1ebf-vux1ec1-nhuxe2n-quux1ea3}{%
\subsection{5.5.3 Hạn chế về nhân
quả}\label{hux1ea1n-chux1ebf-vux1ec1-nhuxe2n-quux1ea3}}

Thiết kế lát cắt ngang, hay dữ liệu bảng gộp nhưng với WBES không phải
là dữ liệu theo dõi cùng doanh nghiệp qua thời gian, giới hạn khả năng
suy luận nhân quả. Quan hệ quốc tế hóa và hiệu quả có thể bị ảnh hưởng
bởi nhân quả ngược: doanh nghiệp hiệu quả hơn có thể có nhiều nguồn lực
để quốc tế hóa hơn, tạo ra chiều nhân quả ngược lại. Hiệu ứng cố định
nền kinh tế và hiệu ứng cố định năm giúp kiểm soát một phần, nhưng không
giải quyết hoàn toàn vấn đề nội sinh.

\hypertarget{hux1ea1n-chux1ebf-vux1ec1-bux1ea5t-biux1ebfn-ux111o-lux1b0ux1eddng-qua-cuxe1c-thux1ebf-hux1ec7-lux1b0ux1ee3c-ux111ux1ed3-vuxe0-buxe3o-huxf2a-thux1b0ux1edbc-ux111o}{%
\subsection{5.5.3.1 Hạn chế về bất biến đo lường qua các thế hệ lược đồ
và bão hòa thước
đo}\label{hux1ea1n-chux1ebf-vux1ec1-bux1ea5t-biux1ebfn-ux111o-lux1b0ux1eddng-qua-cuxe1c-thux1ebf-hux1ec7-lux1b0ux1ee3c-ux111ux1ed3-vuxe0-buxe3o-huxf2a-thux1b0ux1edbc-ux111o}}

Một nhóm hạn chế đo lường bổ sung cần được nêu rõ vì nó định khung mức
độ chắc chắn của các so sánh xuyên thời gian và xuyên bối cảnh. Thứ nhất
là vấn đề bất biến đo lường (measurement invariance) qua các thế hệ lược
đồ khảo sát. Bộ chỉ báo cấu thành DAI thay đổi giữa các thế hệ công cụ
WBES (PICS3, Standardized, BREADY/BEE), nghĩa là cùng một nhãn cấu trúc
DAI được vận hành hóa bằng các tập chỉ báo không hoàn toàn tương đương
qua các đợt. Khi diễn giải các dịch chuyển theo thời gian, chẳng hạn quỹ
đạo không đơn điệu của DAI tại Việt Nam hay sự xuất hiện của Tầng 2 tại
Ấn Độ 2025, cần thận trọng phân biệt giữa thay đổi thực của hiện tượng
và thay đổi do cấu trúc thước đo. Đây là lý do luận án diễn giải các so
sánh chéo đợt một cách dè dặt và dựa vào kiểm định Paternoster để phân
biệt các dịch chuyển có ý nghĩa thống kê khỏi nhiễu đo lường.

Thứ hai là vấn đề bão hòa thước đo Tầng 1. Khi tỷ lệ doanh nghiệp có
website tiệm cận bão hòa, như ở Nhật Bản 83,8\% và các nền tiên tiến nói
chung, phương sai của chỉ báo nhị phân này co lại và nhóm doanh nghiệp
không có website còn lại mang tính chọn lọc đặc thù. Hệ quả là hệ số đo
được phản ánh giới hạn của thước đo Tầng 1 nhiều hơn là giới hạn thực
của số hóa, một vấn đề đo lường cần phân biệt với quan hệ nhân quả. Bằng
chứng biến công cụ tại Việt Nam, nơi hệ số DAI được công cụ hóa bằng
không trong khi hệ số TCI được công cụ hóa dương mạnh, củng cố cách đọc
rằng kết quả không có ý nghĩa về uốn đường cong của DAI ở nhiều bối cảnh
phản ánh tính lỗi thời của biến đại diện Tầng 1 chứ không phải sự vắng
mặt của cơ chế. Hạn chế này có hệ quả trực tiếp: các kết luận về DAI ở
những bối cảnh chỉ đo được Tầng 1 nên được hiểu là chưa kết luận hơn là
bác bỏ, và việc kiểm định lại với thước đo Tầng 2 và Tầng 3 phong phú
hơn là điều kiện cần trước khi khẳng định vai trò của DAI ở các chế độ
thể chế thấp.

Thứ ba, cần nhắc lại rằng mọi đo lường DAI trong luận án đều là chỉ số
chấp nhận số tĩnh tại một thời điểm khảo sát, không nắm bắt được tính
động của việc cập nhật và tái cấu hình năng lực số theo nghĩa của lý
thuyết năng lực động (Teece và cộng sự, 1997; Teece, 2007). Do đó các
phát hiện về DAI phản ánh mức độ chấp nhận số tại thời điểm quan sát chứ
không phải năng lực số năng động, và mọi suy luận về cơ chế động cần
được xem là gợi mở chứ chưa được kiểm chứng trực tiếp trong khuôn khổ dữ
liệu hiện có.

\hypertarget{giux1edbi-hux1ea1n-cux1ee7a-suy-luux1eadn-nhuxe2n-quux1ea3-vuxe0-phux1ea1m-vi-khuxe1i-quuxe1t-huxf3a}{%
\subsection{5.5.3.2 Giới hạn của suy luận nhân quả và phạm vi khái quát
hóa}\label{giux1edbi-hux1ea1n-cux1ee7a-suy-luux1eadn-nhuxe2n-quux1ea3-vuxe0-phux1ea1m-vi-khuxe1i-quuxe1t-huxf3a}}

Cần định khung rõ bản chất suy luận mà thiết kế dữ liệu cho phép. Dữ
liệu doanh nghiệp WBES là lát cắt ngang trong từng đợt khảo sát; do đó
khái niệm ``hiệu quả'' xuyên suốt luận án chỉ một liên hệ đồng thời giữa
cường độ xuất khẩu và năng suất lao động trong một năm khảo sát nhất
định, không phải một hiệu ứng nhân quả được theo dõi theo thời gian trên
cùng doanh nghiệp. Chiến lược biến công cụ dùng để xử lý chọn lọc tham
gia xuất khẩu trong nghiên cứu đơn quốc gia Việt Nam (P3) không khả thi
ở quy mô 50 nền kinh tế, nên ở các phân tích gộp, chọn lọc tham gia xuất
khẩu dựa trên năng lực vừa là một phần của hiện tượng mà điểm uốn hiệu
chỉnh năng lực bộc lộ, vừa là nguồn nội sinh không thể loại bỏ hoàn toàn
bằng hiệu ứng cố định nền kinh tế và năm. Vì vậy cấu trúc ba vùng nên
được đọc như một quy luật liên hệ vững chắc về mặt thống kê, chứ chưa
phải một cơ chế nhân quả đã được chứng minh trực tiếp.

Phạm vi khái quát hóa cũng cần được giới hạn tương xứng. Mẫu bao gồm các
nền kinh tế châu Á và Thái Bình Dương đã được ánh xạ vào khung ICRV, với
cấu trúc thể chế và mô thức thương mại đặc thù của khu vực; việc mở rộng
kết luận ba vùng sang các khu vực khác, chẳng hạn châu Phi hạ Sahara hay
Mỹ Latinh, là một giả thuyết cần kiểm định độc lập chứ không phải một
suy rộng đương nhiên. Tương tự, cấu trúc ba vùng được nhận diện trên dải
cường độ xuất khẩu và phân bố quy mô doanh nghiệp của mẫu WBES, vốn
thiên về doanh nghiệp vừa và nhỏ ở các nền đang phát triển và mới nổi;
tính ổn định của điểm uốn ngoài dải này, đặc biệt ở các tập đoàn đa quốc
gia lớn không nằm trong khung lấy mẫu WBES, vẫn là câu hỏi mở. Những
giới hạn này không làm suy yếu các kết luận trong phạm vi đã định mà xác
định đúng đường biên mà bên trong đó các kết luận được hỗ trợ bởi bằng
chứng.

\hypertarget{hux1ea1n-chux1ebf-vux1ec1-thux1b0ux1edbc-ux111o-quux1ed1c-tux1ebf-huxf3a-vuxe0-thux1b0ux1edbc-ux111o-hiux1ec7u-quux1ea3}{%
\subsection{5.5.3.3 Hạn chế về thước đo quốc tế hóa và thước đo hiệu
quả}\label{hux1ea1n-chux1ebf-vux1ec1-thux1b0ux1edbc-ux111o-quux1ed1c-tux1ebf-huxf3a-vuxe0-thux1b0ux1edbc-ux111o-hiux1ec7u-quux1ea3}}

Một nhóm hạn chế nữa nằm ở chính việc vận hành hóa hai cấu trúc trung
tâm. Về phía quốc tế hóa, luận án dùng nhất quán cường độ xuất khẩu trên
doanh thu (FSTS) làm thước đo chính. Lựa chọn này có ưu điểm về tính so
sánh xuyên bối cảnh và tính sẵn có trong WBES, nhưng nó chỉ nắm bắt một
chiều của quốc tế hóa là cường độ thị trường nước ngoài, mà bỏ qua các
chiều khác như phạm vi địa lý (số thị trường) và chiều sâu cam kết (đầu
tư trực tiếp ra nước ngoài, lập cơ sở ở nước ngoài). Chính phân tích
tổng hợp P6 cho thấy cỡ ảnh hưởng quan sát thay đổi theo cách đo lường
quốc tế hóa, nên không loại trừ khả năng rằng dạng hàm ba vùng được nhận
diện cho cường độ xuất khẩu có thể khác đi nếu đo quốc tế hóa bằng phạm
vi địa lý hoặc chỉ số đa chiều. Cảnh báo này được củng cố bởi bằng chứng
gần đây của Verbeke, Eschmann và Gugler (2025): khi đo mức độ toàn cầu
hóa của doanh nghiệp đa quốc gia trên bốn chiều khác nhau (doanh số theo
không gian địa lý, đầu tư theo tài sản, vốn nhân lực theo phân bố lao
động, và vốn tri thức theo sáng chế), các tác giả phát hiện mức độ quốc
tế hóa đo được chênh lệch đáng kể giữa các chiều, đến mức một doanh
nghiệp có thể được xếp là toàn cầu theo chiều này nhưng định hướng nội
vùng theo chiều khác. Phát hiện này hàm ý trực tiếp rằng kết luận của
luận án, vốn dựa trên một chiều cường độ xuất khẩu, cần được nhân rộng
bằng các thước đo quốc tế hóa đa chiều trước khi khái quát hóa. Về phía
hiệu quả, luận án dùng năng suất lao động dạng logarit, một thước đo vận
hành mạnh nhưng hẹp hơn so với các thước đo hiệu quả đa chiều (lợi
nhuận, tăng trưởng, giá trị thị trường) mà tài liệu thường dùng; phân
tích tổng hợp cho thấy thước đo hiệu quả cũng tác động đáng kể đến cỡ
ảnh hưởng. Do đó các kết luận của luận án nên được đọc là gắn với cặp
thước đo cụ thể (cường độ xuất khẩu và năng suất lao động) và việc nhân
rộng với các cặp thước đo khác là một bước kiểm định cần thiết.

\hypertarget{hux1ea1n-chux1ebf-vux1ec1-phuxe2n-loux1ea1i-icrv-vuxe0-ux111o-lux1b0ux1eddng-thux1ec3-chux1ebf}{%
\subsection{5.5.3.4 Hạn chế về phân loại ICRV và đo lường thể
chế}\label{hux1ea1n-chux1ebf-vux1ec1-phuxe2n-loux1ea1i-icrv-vuxe0-ux111o-lux1b0ux1eddng-thux1ec3-chux1ebf}}

Khung phân loại ICRV sáu chế độ, dù là một đóng góp trung tâm, cũng mang
những hạn chế đo lường cần nêu rõ. Thứ nhất, việc gán một nền kinh tế
vào một chế độ là một thao tác rời rạc hóa một thực tại liên tục: các
nền kinh tế nằm gần ranh giới giữa hai chế độ có thể bị phân loại nhạy
cảm với ngưỡng, và một số nền có thể dịch chuyển chế độ trong giai đoạn
quan sát dài 2006--2026 mà phân loại tĩnh không nắm bắt được. Thứ hai,
bản thân các chỉ số thể chế nền tảng dùng để xây dựng ICRV là các đại
lượng đo lường gián tiếp với sai số riêng, và một bộ phận doanh nghiệp
trong mẫu gộp không được gán chế độ do thiếu chỉ số thể chế, tạo khả
năng thiên lệch nếu các quan sát thiếu chỉ số không ngẫu nhiên. Thứ ba,
hệ phân loại dùng ở cấp phân tích tổng hợp P6 (năm chế độ rút gọn) khác
với hệ sáu nhóm ở cấp doanh nghiệp, nên việc đối chiếu trực tiếp giữa
hai cấp cần thận trọng. Những hạn chế này không phủ định giá trị dự đoán
của ICRV mà chỉ ra rằng phân loại nên được hiểu là một xấp xỉ có cấu
trúc của gradient thể chế chứ không phải một thang đo chính xác tuyệt
đối.

\hypertarget{hux1ea1n-chux1ebf-vux1ec1-suy-luux1eadn-quanh-ux111ux1eb7c-trux1b0ng-quux1ea3n-trux1ecb}{%
\subsection{5.5.3.5 Hạn chế về suy luận quanh đặc trưng quản
trị}\label{hux1ea1n-chux1ebf-vux1ec1-suy-luux1eadn-quanh-ux111ux1eb7c-trux1b0ng-quux1ea3n-trux1ecb}}

Cuối cùng, các kết quả liên quan đặc trưng nhà quản trị cấp cao mang một
giới hạn suy luận đặc thù cần được nhấn mạnh để tránh diễn giải sai. Hệ
số mức âm nhẹ của nữ quản lý cấp cao (\(\hat{\beta} = -0{,}088\),
\(p = {,}026\); xem Mục 5.1.4) được luận án diễn giải là phản ánh sự
phân bố không ngẫu nhiên của nữ quản lý vào các ngành thâm dụng lao động
biên thấp, tức một dạng chọn lọc cấu thành, chứ không phải bằng chứng
nhân quả về hiệu quả của lãnh đạo nữ. Thiết kế dữ liệu hiện tại không
cho phép bóc tách kênh chọn lọc ngành khỏi kênh hiệu quả lãnh đạo, vì
WBES không có đủ biến để mô hình hóa quá trình phân bổ nhà quản trị vào
doanh nghiệp và ngành. Tương tự, kết quả không có ý nghĩa về kinh nghiệm
nhà quản trị phản ánh hạn chế đo lường là số năm kinh nghiệm không nắm
bắt được chất lượng và bề rộng kinh nghiệm quốc tế. Do đó mọi kết luận
về tầng quản trị của khung CDCM trong luận án nên được đọc là có điều
kiện và cần thiết kế nhận dạng riêng, một giới hạn mà chương sách bổ
sung về quản trị cấp cao tại Ấn Độ chỉ giải quyết được ở phạm vi đơn
quốc gia chứ chưa ở quy mô khu vực.

\hypertarget{hux1b0ux1edbng-nghiuxean-cux1ee9u-tux1b0ux1a1ng-lai}{%
\subsection{5.5.4 Hướng nghiên cứu tương
lai}\label{hux1b0ux1edbng-nghiuxean-cux1ee9u-tux1b0ux1a1ng-lai}}

Dựa trên các hạn chế trên, luận án đề xuất bảy hướng nghiên cứu ưu tiên:

\textbf{Thứ nhất}, xây dựng đo lường DAI phong phú hơn (Tầng 2 và Tầng
3) trên dữ liệu WBES mới nhất (2023--2026) và các nguồn dữ liệu bổ sung
(ITU Digital Development Index, OECD Digital Economy Outlook) để kiểm
định lại vai trò DAI trong các bối cảnh ICRV thấp hơn.

\textbf{Thứ hai}, sử dụng phương pháp nhận dạng nhân quả (causal
identification): instrumental variables (ví dụ: sử dụng các cải cách
thương mại ngoại sinh như FTA ký kết), difference-in-differences, hoặc
regression discontinuity design để ước lượng hiệu ứng nhân quả của quốc
tế hóa lên hiệu quả.

\textbf{Thứ ba}, mở rộng phân tích FIP tại SIDS với các công cụ đo lường
và dữ liệu panel thực sự (theo dõi cùng doanh nghiệp qua nhiều năm) để
kiểm định liệu FIP là trạng thái cấu trúc dài hạn hay chỉ là giai đoạn
tạm thời khi doanh nghiệp SIDS mới bắt đầu xuất khẩu.

\textbf{Thứ tư}, kiểm định vai trò của top manager characteristics (H4,
kinh nghiệm, giới tính) trong từng bối cảnh ICRV riêng biệt, đặc biệt
trong bối cảnh đang nổi và SIDS nơi tầng năng lực và thể chế yếu hơn,
liệu đặc điểm nhà quản trị có thể bù đắp một phần cho những thiếu hụt đó
không?

\textbf{Thứ năm}, khai thác sâu hơn Nhật Bản, nền kinh tế tiên tiến lớn
lần đầu được WBES khảo sát năm 2025. Bước đầu tiên đã được thực hiện
trong nghiên cứu thành phần P10: trên đợt khảo sát khai mở (N = 2.168),
quan hệ quốc tế hóa và hiệu quả tại Nhật gần tuyến tính dương, không có
điểm uốn trong miền quan sát, xác nhận dự đoán biên trên của khung ICRV;
các bước tiếp theo gồm tích hợp Nhật Bản vào các mô hình điều tiết ba
chiều và khai thác các module đặc thù (trường thọ doanh nghiệp, khoảng
cách giới trong lãnh đạo) khi có thêm đợt khảo sát thứ hai.

\textbf{Thứ sáu}, một hướng nghiên cứu ưu tiên xuất phát trực tiếp từ
phát hiện ngưỡng tan biến tại Ấn Độ là kiểm định nhân quả của biến động
thể chế lên hình dạng quan hệ quốc tế hóa và hiệu quả. Các cú sốc thể
chế rời rạc và có thể định thời điểm trong cửa sổ 2014--2025 (bãi bỏ
tiền mặt, áp dụng thuế hàng hóa và dịch vụ, ban hành luật phá sản, ra
mắt hạ tầng thanh toán hợp nhất) tạo điều kiện cho thiết kế khác biệt
trong khác biệt hoặc nghiên cứu sự kiện nhằm phân lập tác động nhân quả
của từng cú sốc lên độ cong và vị trí điểm uốn. Hướng này cũng cho phép
kiểm định trực tiếp mệnh đề về điều kiện ổn định thể chế của khung CDCM,
tức liệu sự tan biến ngưỡng là hệ quả của tổng biến động tích lũy hay
của một cú sốc cụ thể nào đó. Việc kết hợp dữ liệu chính sách thương mại
ngoại sinh, chẳng hạn các hiệp định thương mại tự do được ký kết, như
công cụ nhận dạng sẽ giúp tách bạch chiều nhân quả từ quốc tế hóa đến
hiệu quả khỏi chiều ngược lại, giải quyết một phần hạn chế nội sinh đã
nêu ở Mục 5.5.3.

\textbf{Thứ bảy}, sự dịch chuyển trọng tâm giữa biên tham gia và biên
cường độ dọc theo gradient thể chế, được phát hiện khi đối chiếu bằng
chứng Việt Nam (cơ chế biên tham gia ở chế độ chuyển đổi) với Nhật Bản
(cơ chế tự chọn lọc Melitz ở biên trên), mở ra một câu hỏi lý thuyết
chưa được giải quyết: chế độ thể chế nào đánh dấu sự chuyển giao giữa
hai cơ chế, và liệu có một vùng quá độ nơi cả hai cùng vận hành. Nghiên
cứu tương lai có thể vận dụng mô hình hai phần (two-part model) phân
tách rõ ràng quyết định tham gia xuất khẩu khỏi quyết định cường độ trên
toàn gradient ICRV, qua đó lập bản đồ chính xác sự dịch chuyển cơ chế và
làm giàu thêm tầng vi mô của khung CDCM, nhất quán với logic tự chọn lọc
theo năng suất của Melitz (2003) và với bằng chứng về quan hệ xuất khẩu
và năng suất ở cấp doanh nghiệp (Wagner, 2007).

\textbf{Thứ tám}, kiểm định tính khái quát hóa của cấu trúc ba vùng ra
ngoài khu vực châu Á và Thái Bình Dương. Vì WBES có độ phủ toàn cầu, một
hướng khả thi và trực tiếp là ánh xạ khung ICRV lên các nền kinh tế châu
Phi hạ Sahara và Mỹ Latinh rồi ước lượng lại gradient điểm uốn trên cùng
cặp thước đo, qua đó xác định liệu trật tự ba vùng có phải là đặc thù
khu vực hay là một quy luật phổ quát hơn về điều tiết thể chế. Hướng này
có giá trị vì nó kiểm định trực tiếp đường biên khái quát hóa mà Mục
5.5.3.2 đã nêu, và vì các khu vực này có những chế độ thể chế và mô thức
thương mại khác với châu Á, cung cấp một phép kiểm định khắt khe cho
tính bền của khung.

\textbf{Thứ chín}, mở rộng đo lường quốc tế hóa sang đa chiều. Vì cường
độ xuất khẩu chỉ là một chiều, nghiên cứu tương lai có thể xây dựng
thước đo phạm vi địa lý và chiều sâu cam kết từ các nguồn bổ sung (dữ
liệu hải quan, dữ liệu đầu tư trực tiếp ra nước ngoài) rồi kiểm định
liệu cấu trúc ba vùng có giữ nguyên hình dạng khi quốc tế hóa được đo đa
chiều. Bằng chứng phân tích tổng hợp rằng cỡ ảnh hưởng thay đổi theo
cách đo quốc tế hóa khiến đây là một hướng có khả năng làm thay đổi hoặc
tinh chỉnh kết luận, chứ không chỉ là một bài tập nhân rộng.

\textbf{Thứ mười}, một hướng đo lường ưu tiên là phát triển và kiểm định
thước đo năng lực số động thực sự. Vì mọi đo lường DAI hiện tại đều
tĩnh, một bước tiến quan trọng là kết hợp các đợt khảo sát liên tiếp của
cùng nền kinh tế để xây dựng đại lượng tái cấu hình năng lực số theo
thời gian, tiệm cận khái niệm năng lực động (Teece và cộng sự, 1997;
Teece, 2007), qua đó kiểm định trực tiếp giả thuyết rằng vai trò uốn
đường cong của năng lực số đến từ tính động chứ không chỉ từ mức chấp
nhận tĩnh. Kết hợp với các nguồn dữ liệu số chuyên biệt, hướng này cho
phép tách bạch hiệu ứng của chiều sâu năng lực số khỏi hiệu ứng của tính
động, một phân biệt mà dữ liệu WBES tĩnh hiện chưa cho phép.

\begin{center}\rule{0.5\linewidth}{0.5pt}\end{center}

\hypertarget{kux1ebft-luux1eadn}{%
\section{5.6 Kết luận}\label{kux1ebft-luux1eadn}}

\hypertarget{tux1ed5ng-kux1ebft-cuxe1c-nghiuxean-cux1ee9u-thuxe0nh-phux1ea7n}{%
\subsection{5.6.1 Tổng kết các nghiên cứu thành
phần}\label{tux1ed5ng-kux1ebft-cuxe1c-nghiuxean-cux1ee9u-thuxe0nh-phux1ea7n}}

Luận án này được xây dựng trên mười nghiên cứu thành phần và một chương
sách bổ sung cho nhau, kết hợp bằng chứng từ phân tích tổng hợp và phân
tích thực nghiệm sơ cấp:

\begin{itemize}
\item
  \textbf{P1}: cơ sở thực nghiệm đa quốc gia trên 17 nền kinh tế châu Á
  (\(N = 40.633\)), xác lập phân tầng năng suất vùng và hiệu ứng lá chắn
  số (tương tác rào cản và công nghệ \(+0{,}110\), \(p < {,}001\)), đặt
  nền cho khung CDCM.
\item
  \textbf{P2}: doanh nghiệp nhỏ và vừa Trung Quốc (\(N = 4.290\)), xác
  nhận chữ U ngược với điểm uốn 47,8\% và cơ chế bẫy vốn lưu động ở
  nhánh đi xuống.
\item
  \textbf{P3}: phân tích Việt Nam (chuyển đổi thu nhập trung bình thấp),
  xác nhận chữ U ngược với điểm uốn 39,3--46,2\%, TCI là biến điều tiết
  dương có ý nghĩa.
\item
  \textbf{P4}: phân tích Singapore (tiên tiến đổi mới), xác nhận chữ U
  ngược với điểm uốn khoảng 88,6\% (khoảng tin cậy rộng), DAI khuếch đại
  tại mức cường độ xuất khẩu cao.
\item
  \textbf{P5}: phân tích Trung Quốc 2012--2024 (trung bình cao), xác
  nhận chữ U ngược ổn định theo thời gian (điểm uốn khoảng 48,78\%, kiểm
  định Paternoster \(z = 0{,}82\), \(p = {,}412\)).
\item
  \textbf{P6}: phân tích tổng hợp ba tầng với \(k = 238\),
  \(\bar{r} = 0{,}074\), \(I^2 = 62{,}4\%\) (54,1\% tầng hai cộng 8,4\%
  tầng ba), ICRV điều tiết \(Q_M = 17{,}35\) (\(df = 4\),
  \(p = {,}002\)).
\item
  \textbf{P7}: kiểm định toàn mẫu châu Á và Thái Bình Dương trên khung
  50 nền kinh tế, \(N = 81.022\) (bậc hai) và \(N = 79.080\) (mô hình
  đầy đủ), điểm uốn 43,6--51,5\% (mô hình chính 43,6\%,
  \(p_\text{LM} < {,}001\)); chữ U ngược định hình rõ nhất ở Nhóm IV
  (43,0\%), gần tuyến tính ở Nhóm I, tan rã ở Nhóm V và VI; TCI nâng mặt
  bằng (\(+0{,}141\), \(p<{,}001\)), DAI nâng mặt bằng (\(+0{,}201\),
  \(p<{,}001\)) với tương tác đường cong cùng chiều nhưng không có ý
  nghĩa toàn mẫu.
\item
  \textbf{P8}: phân tích bảy nền kinh tế đảo nhỏ đang phát triển Thái
  Bình Dương, \(N = 959\), \(\beta(\text{FSTS}_c) = -1{,}339\), xác nhận
  hình phạt quốc tế hóa cưỡng bức, quan hệ âm đơn điệu không có điểm
  uốn.
\item
  \textbf{P9}: phân tích Ấn Độ qua ba đợt WBES 2014--2025 (khoảng 29.000
  doanh nghiệp), kiểm định độ bền ngưỡng qua thập kỷ chuyển đổi thể chế.
  Phát hiện sự \emph{tan biến} điểm uốn: 61,8\% (2014) rồi 40,7\% (2022)
  rồi mất ý nghĩa độ cong (2025: \(\beta_2 = -0{,}16\), \(p = {,}42\)),
  kèm tương tác Tầng 2 (UPI) âm, bằng chứng đầu tiên rằng chuyển đổi thể
  chế cực mạnh có thể hòa tan chứ không chỉ dịch chuyển ngưỡng.
\item
  \textbf{P10}: phân tích Nhật Bản, đợt khảo sát WBES đầu tiên năm 2025
  (tiên tiến đổi mới), xác nhận dự đoán biên trên: quan hệ duỗi thẳng về
  gần tuyến tính dương với điểm uốn nằm ngoài miền quan sát, năng lực và
  đặc trưng quản trị nâng mặt bằng nhưng không uốn đường cong.
\item
  \textbf{Chương sách} (Do \& Phan, 2025, IntechOpen): nghiên cứu tập
  trung vào quản trị cấp cao tại Ấn Độ (\(N = 380\)), tỷ suất lợi nhuận
  trên doanh thu là biến phụ thuộc; quốc tế hóa tác động âm tuyến tính
  (\(-36{,}66\)), kinh nghiệm (\(+1{,}55\)) và lãnh đạo nữ
  (\(+77{,}03\)) điều tiết dương, cung cấp bằng chứng đơn quốc gia cho
  H4.
\end{itemize}

\hypertarget{phuxe1t-hiux1ec7n-trung-tuxe2m}{%
\subsection{5.6.2 Phát hiện trung
tâm}\label{phuxe1t-hiux1ec7n-trung-tuxe2m}}

Phát hiện trung tâm của luận án có thể được tóm tắt trong một luận điểm
ngắn gọn:

\begin{quote}
\textbf{``Không có một mức độ quốc tế hóa tối ưu duy nhất, điểm uốn phụ
thuộc vào thể chế, và trong điều kiện thể chế cực yếu, quốc tế hóa có
thể gây ra bất lợi hiệu quả thay vì lợi ích.''}
\end{quote}

Bằng chứng từ 238 nghiên cứu (P6 meta-analysis) và 81.022 doanh nghiệp
tại 50 nền kinh tế châu Á và Thái Bình Dương (P7) đều chỉ đến kết luận
này. Quan hệ quốc tế hóa và hiệu quả không phải là một đường cong phổ
quát mà là một \textbf{họ đường cong phụ thuộc bối cảnh} (family of
context-dependent curves): chữ U ngược định hình rõ nhất ở vùng giữa
gradient thể chế (Nhóm IV, điểm uốn khoảng 43\%), duỗi thẳng gần tuyến
tính ở biên trên nơi thể chế mạnh đẩy điểm uốn ra mức cường độ xuất khẩu
rất cao (Singapore khoảng 88,6\%, gần hoặc vượt biên dữ liệu quan sát),
và tan rã ở biên dưới nơi thể chế yếu nhất; nghĩa là thể chế càng mạnh,
doanh nghiệp càng duy trì lợi ích từ quốc tế hóa đến mức cường độ cao
hơn trước khi đạt đỉnh.

\hypertarget{uxfd-nghux129a-ux111ux1ed1i-vux1edbi-luxfd-thuyux1ebft-vuxe0-thux1ef1c-tiux1ec5n}{%
\subsection{5.6.3 Ý nghĩa đối với lý thuyết và thực
tiễn}\label{uxfd-nghux129a-ux111ux1ed1i-vux1edbi-luxfd-thuyux1ebft-vuxe0-thux1ef1c-tiux1ec5n}}

Về lý thuyết, luận án đóng góp vào ba dòng nghiên cứu: (i) phân tích
tổng hợp tài liệu về quốc tế hóa và hiệu quả bằng cách phân tách cấu
trúc dị biệt và xác định ICRV như biến điều tiết có ý nghĩa; (ii) tài
liệu về năng lực số bằng cách làm rõ DAI là nguồn lực tình huống chứ
không phải năng lực nền tảng; và (iii) tài liệu về quốc tế hóa và SIDS
bằng cách đặt tên và cung cấp bằng chứng đầu tiên cho FIP.

Về thực tiễn, luận án cung cấp cơ sở để nhà hoạch định chính sách và nhà
quản trị doanh nghiệp tại từng bối cảnh ICRV cụ thể hiệu chỉnh chiến
lược quốc tế hóa phù hợp, không áp dụng máy móc bài học từ Singapore hay
Nhật Bản vào bối cảnh Việt Nam hay SIDS, mà nhận ra sự phụ thuộc bối
cảnh như là nguyên lý căn bản trong hoạch định chiến lược quốc tế hóa.

\hypertarget{lux1eddi-kux1ebft}{%
\subsection{5.6.4 Lời kết}\label{lux1eddi-kux1ebft}}

Luận án này bắt đầu từ câu hỏi đơn giản: quốc tế hóa có cải thiện hiệu
quả của doanh nghiệp không? Câu trả lời, sau mười nghiên cứu thành phần
và một chương sách, phân tích trên hàng chục ngàn doanh nghiệp ở gần năm
mươi nền kinh tế, là: \textbf{có, thường là như vậy, nhưng không phải
lúc nào, và không phải ở mọi nơi}. Giới hạn của lợi ích là có thực, phụ
thuộc vào bối cảnh thể chế, và trong trường hợp các nền kinh tế đảo nhỏ,
quốc tế hóa bắt buộc có thể gây ra thiệt hại. Hiểu biết này, rằng quan
hệ quốc tế hóa và hiệu quả là quan hệ có điều kiện theo thể chế chứ
không phải quan hệ phổ quát, là đóng góp cốt lõi mà luận án mang lại cho
khoa học quản trị kinh doanh quốc tế.

\begin{center}\rule{0.5\linewidth}{0.5pt}\end{center}

\newpage

\hypertarget{tuxe0i-liux1ec7u-tham-khux1ea3o}{%
\chapter{TÀI LIỆU THAM KHẢO}\label{tuxe0i-liux1ec7u-tham-khux1ea3o}}

\begin{quote}
Trình bày theo chuẩn APA 7th, một danh sách thống nhất xếp theo thứ tự
bảng chữ cái tên tác giả (hoặc tên tổ chức). Tài liệu cùng tác giả--cùng
năm phân biệt bằng hậu tố a, b, c.
\end{quote}

Abdi, M., \& Aulakh, P. S. (2018). Internationalization and performance:
Degree, duration, and scale of operations in foreign markets.
\emph{Journal of International Business Studies, 49}(7), 832--857.
\url{https://doi.org/10.1057/s41267-017-0117-9}

Abreha, K., \& Lopez Acevedo, G. (2024). \emph{Trade restructuring:
Assessing labor market and welfare effects} (Policy Research Working
Paper No.~10955). World Bank Group.

ADB. (2016). \emph{Asian Economic Integration Report 2016 Highlights}.
Asian Development Bank.
\url{https://www.adb.org/publications/asian-economic-integration-report-2016}

ADB. (2024). \emph{Asian Development Outlook 2024 --- Recovery and
resilience in Asia}. Asian Development Bank.

ADB. (2025). \emph{Annual Report 2025 --- Online Appendix 1: Resolutions
of the Board of Governors Adopted in 2025}. Asian Development Bank.**
\url{https://www.adb.org/documents/adb-annual-report-2025}

ADB. (2026). \emph{Asia-Pacific Development Report 2026: Highlights ---
Global Value Chain Participation in Developing Asia}. Asian Development
Bank.

ADB. (2026, April). \emph{Asian Development Outlook April 2026: Special
Topic --- Artificial Intelligence in Asia and the Pacific}. Asian
Development Bank.

ADB. (2026, April). \emph{Asian Development Outlook April 2026: The
Middle East Conflict Challenges Resilience in Asia and the Pacific}.
Asian Development Bank.

ADB. (2026, April). \emph{Vietnam Economic Outlook 2026: Navigating the
Crosscurrents}. Asian Development Bank.**
\url{https://www.adb.org/viet-nam}

ADB. (2026, March). \emph{Asia Bond Monitor March 2026: Emerging East
Asian Local Currency Bond Markets}. Asian Development Bank.
\url{https://doi.org/10.22617/SGP260088-2}

ADB. (2026, May). \emph{Market Assessment of Critical Minerals in
Mongolia}. Asian Development Bank.
\url{https://doi.org/10.22617/TCS260157-2}

Agarwal, A., Barattieri, A., \& Mattoo, A. (2026, April). \emph{Foreign
demand shocks and domestic value added: Firm-level evidence from Viet
Nam} (Policy Research Working Paper No.~11352). World Bank Group.

Agnolucci, P., Kenworthy, P., \& Kose, M. A. (2026, May 7). Five
questions on how the war in the Middle East is affecting commodity
markets. \emph{Let's Talk Development}. World Bank Blogs.
\url{https://blogs.worldbank.org/en/developmenttalk/five-questions-on-how-the-war-in-the-middle-east-is-affecting-co}

Aguinis, H., Dalton, D. R., Bosco, F. A., Pierce, C. A., \& Dalton, C.
M. (2011). Meta-analytic choices and judgment calls: Implications for
theory and research. \emph{Journal of Management, 37}(1), 5--38.
\url{https://doi.org/10.1177/0149206310370708}

Aguinis, H., Hill, N. S., \& Bailey, J. R. (2019). Best practices in
data collection and preparation: Recommendations for reviewers, editors,
and authors. \emph{Organizational Research Methods, 24}(4), 678--693.**
\url{https://doi.org/10.1177/1094428119836485}

Ahsan, F. M., \& Sinha, A. (2022). Internationalization motives,
location advantages and performance: The case of Indian firms from
knowledge-intensive industries. \emph{Cross Cultural \& Strategic
Management}. \url{https://doi.org/10.1108/ccsm-07-2021-0119}

Aiken, L. S., \& West, S. G. (1991). \emph{Multiple regression: Testing
and interpreting interactions}. Sage.

Al-Najjar, B., Salama, A., \& Abed, A. (2025). Gender and cultural
diversity on boards: Their impact on internationalisation. \emph{The
International Journal of Human Resource Management, 36}(13),
2305--2341.** \url{https://doi.org/10.1080/09585192.2025.2550765}

Almodóvar, P., \& Rugman, A. M. (2014). The M curve and the performance
of Spanish international new ventures. \emph{British Journal of
Management, 25}(S1), S6--S23.
\url{https://doi.org/10.1111/1467-8551.12048}

Altaf, M., \& Shah, S. A. A. (2016). Internationalization and
performance relationship in Pakistan: Moderating role of innovation.
\emph{Pakistan Journal of Commerce and Social Sciences, 10}(1), 76--90.

Angrist, J. D., \& Pischke, J.-S. (2009). \emph{Mostly harmless
econometrics: An empiricist's companion}. Princeton University Press.

Anil, K., \& Misra, A. (2022). Artificial intelligence in peer-to-peer
lending in India: A cross-case analysis. \emph{International Journal of
Emerging Markets, 17}(4), 1085--1106.

Arbelo, A., Arbelo-Pérez, M., \& Pérez-Gómez, P. (2024).
Internationalization and individual firm performance: A resource-based
view. \emph{Eurasian Business Review, 14}(4), 881--915.

Arte, L., \& Larimo, J. (2022). Moderating influence of psychic distance
on the international diversification--performance relationship: Insights
from a meta-analysis. \emph{Journal of Business Research, 139},
230--245. \url{https://doi.org/10.1016/j.jbusres.2021.09.060}

Assaf, A. G., Josiassen, A., Knezevic Cvelbar, L., \& Woo, L. (2012).
The effects of internationalization on hotel financial performance.
\emph{International Journal of Hospitality Management, 31}(1), 39--48.
\url{https://doi.org/10.1016/j.ijhm.2011.05.003}

Assaf, A. G., Josiassen, A., Knezevic Cvelbar, L., \& Woo, L. (2016).
The role of internationalization in the performance of hotel firms.
\emph{Annals of Tourism Research, 57}, 100--112.
\url{https://doi.org/10.1016/j.annals.2015.12.008}

Aterido, R., Hallward-Driemeier, M., \& Pagés, C. (2011). Big
constraints to small firms' growth? Business environment and employment
growth across firms. \emph{Economic Development and Cultural Change,
59}(3), 609--647. \url{https://doi.org/10.1086/658349}

Avenyo, E. K., Tregenna, F., \& Kraemer-Mbula, E. (2021). Do productive
capabilities affect export performance? Evidence from African firms.
\emph{European Journal of Development Research, 33}(2), 304--329.
\url{https://doi.org/10.1057/s41287-021-00364-6}

Aydemir, M., \& Alper, D. (2024). The internationalization-firm
performance relationship: A meta-analysis. \emph{International Journal
of Management Economics and Business}.
\url{https://doi.org/10.17130/ijmeb.1458989}

Ayyagari, M., Demirgüç-Kunt, A., \& Maksimovic, V. (2011). Firm
innovation in emerging markets. \emph{Journal of Financial and
Quantitative Analysis, 46}(6), 1545--1580.

Azman, M. A. H., Hon, C. K. H., Xia, B., Lee, B. L., \& Skitmore, M.
(2022). Internationalization and productivity of construction firms in a
developing country: The effect of institutional environment and
ownership on Malaysian construction firms. \emph{Journal of Management
in Engineering, 38}(3).
\url{https://doi.org/10.1061/(asce)me.1943-5479.0000985}

Babina, T., Fedyk, A., He, A., \& Hodson, J. (2024). Artificial
intelligence, firm growth, and product innovation. \emph{Journal of
Financial Economics, 151}, 103745.
\url{https://doi.org/10.1016/j.jfineco.2023.103745}

Bae, S.-C., Park, B. J., \& Wang, X. (2008). Multinationality, R\&D
intensity, and firm performance: Evidence from U.S. manufacturing firms.
\emph{Multinational Business Review, 16}(1), 53--77.
\url{https://doi.org/10.1108/1525383X200800003}

Banalieva, E. R., \& Dhanaraj, C. (2019). Internalization theory for the
digital economy. \emph{Journal of International Business Studies,
50}(8), 1372--1387.

Banerjee, A. V., \& Duflo, E. (2014). Do firms want to borrow more?
Testing credit constraints using a directed lending program.
\emph{Review of Economic Studies, 81}(2), 572--607.
\url{https://doi.org/10.1093/restud/rdt046}

Barattieri, A., Mattoo, A., \& Signoret, J. (2026, April). \emph{Foreign
tariffs as catalysts for reform? The potential development dividend}
(Policy Research Working Paper No.~11354). World Bank Group.

Barney, J. (1991). Firm resources and sustained competitive advantage.
\emph{Journal of Management, 17}(1), 99--120.
\url{https://doi.org/10.1177/014920639101700108}

Barłożewski, K., \& Trąpczyński, P. (2021a). Internationalisation
motives and the multinationality-performance relationship: The case of
Polish firms. \emph{Entrepreneurial Business and Economics Review,
9}(2), 85--104. \url{https://doi.org/10.15678/eber.2021.090206}

Barłożewski, K., \& Trąpczyński, P. (2021b). Is internationalisation
beneficial for novice internationalisers? The performance effects of
firm-specific advantages, internationalisation degree and firm size
revisited. \emph{Oeconomia Copernicana, 12}(1), 53--75.
\url{https://doi.org/10.24136/oc.2021.003}

Bausch, A., \& Krist, M. (2007). The effect of context-related
moderators on the internationalization--performance relationship:
Evidence from meta-analysis. \emph{Management International Review,
47}(3), 319--347. \url{https://doi.org/10.1007/s11575-007-0019-z}

Beblawi, H. (1987). The rentier state in the Arab world. In H. Beblawi
\& G. Luciani (Eds.), \emph{The rentier state} (pp.~49--62). Croom Helm.

Begg, C. B., \& Mazumdar, M. (1994). Operating characteristics of a rank
correlation test for publication bias. \emph{Biometrics, 50}(4),
1088--1101. \url{https://doi.org/10.2307/2533446}

Bello, D. C., \& Kostova, T. (2012). From the editors: Conducting high
impact international business research: The role of theory.
\emph{Journal of International Business Studies, 43}(6), 537--543.**
\url{https://doi.org/10.1057/jibs.2012.14}

Benito-Osorio, D., Guerras-Martín, L. Á., Zuñiga-Vicente, J. Á., \&
Rivero-Meneses, A. (2015). Diversification and performance: The
moderating effect of firm's international activity. \emph{International
Business Review, 24}(5), 893--907.
\url{https://doi.org/10.1016/j.ibusrev.2015.03.003}

Berry, H., \& Kaul, A. (2016). Replicating the
multinationality-performance relationship: Is there an S-curve?
\emph{Strategic Management Journal, 37}(11), 2275--2290.
\url{https://doi.org/10.1002/smj.2429}

Bertram, G. (2006). Introduction: The MIRAB model in the twenty-first
century. \emph{Asia Pacific Viewpoint, 47}(1), 1--13.
\url{https://doi.org/10.1111/j.1467-8373.2006.00297.x}

Bhandari, K. R., Zámborský, P., Ranta, M., \& Salo, J. (2023).
Digitalization, internationalization, and firm performance.
\emph{International Business Review, 32}(4), 102027.

Bharadwaj, A., El Sawy, O. A., Pavlou, P. A., \& Venkatraman, N. (2013).
Digital business strategy: Toward a next generation of insights.
\emph{MIS Quarterly, 37}(2), 471--482.
\url{https://doi.org/10.25300/MISQ/2013/37:2.3}

Bloom, N., Genakos, C., Sadun, R., \& Van Reenen, J. (2012). Management
practices across firms and countries. \emph{Academy of Management
Perspectives, 26}(1), 12--33.
\url{https://doi.org/10.5465/amp.2011.0077}

Boeddu, G., Feyen, E., Martinez Jaramillo, S., Mesquita, S., Palta, Y.,
Sarkar, A., Sinha, S., \& Gutiérrez Traverso, A. (2025).
\emph{Artificial Intelligence for Financial Sector Supervision: An
Emerging Market and Developing Economies Perspective} (Prosperity
Insight Series). World Bank Group.
\url{https://doi.org/10.1596/978-1-4648-2132-7}

Booltink, L. W. A., \& Saka-Helmhout, A. (2018). The effects of R\&D
intensity and internationalization on firm performance. \emph{Journal of
Small Business Management, 56}(S1), 170--185.
\url{https://doi.org/10.1111/jsbm.12393}

Borda, A., Geleilate, J.-M. G., Newburry, W., \& Kundu, S. K. (2017).
Firm internationalization, business group diversification and firm
performance: The case of Latin American firms. \emph{Journal of Business
Research, 72}, 104--113.
\url{https://doi.org/10.1016/j.jbusres.2016.11.021}

Borenstein, M., Hedges, L. V., Higgins, J. P. T., \& Rothstein, H. R.
(2021). \emph{Introduction to meta-analysis} (2nd ed.). Wiley.

Breusch, T. S., \& Pagan, A. R. (1979). A simple test for
heteroscedasticity and random coefficient variation. \emph{Econometrica,
47}(5), 1287--1294. \url{https://doi.org/10.2307/1911963}

Brida, J. G., Driha, O. M., \& Scuderi, R. (2016). Internationalization
of hotels: Effects on performance. \emph{International Journal of
Contemporary Hospitality Management, 28}(1), 179--202.
\url{https://doi.org/10.1108/IJCHM-10-2014-0517}

Briguglio, L. (1995). Small island developing states and their economic
vulnerabilities. \emph{World Development, 23}(9), 1615--1632.
\url{https://doi.org/10.1016/0305-750X(95)00065-K}

Brock, D. M., \& Yaffe, T. (2008). International diversification and
performance: The mediating role of implementation. \emph{International
Business Review, 17}(5), 600--615.
\url{https://doi.org/10.1016/j.ibusrev.2008.08.001}

Brynjolfsson, E., \& McAfee, A. (2014). \emph{The second machine age:
Work, progress, and prosperity in a time of brilliant technologies}. W.
W. Norton \& Company.

Brynjolfsson, E., Rock, D., \& Syverson, C. (2021). The productivity
J-curve: How intangibles complement general purpose technologies.
\emph{American Economic Journal: Macroeconomics, 13}(1), 333--372.
\url{https://doi.org/10.1257/mac.20180386}

Buchhave, H., Nguyen, C. V., Vu, C., Nguyen, G. T., \& Zumbyte, I.
(2026). \emph{Care for growth: Making industrial jobs work for women in
Viet Nam} (Policy Summary Note). World Bank Group \& Australian Aid.

Buckley, P. J., \& Casson, M. C. (1976). \emph{The future of the
multinational enterprise}. Macmillan.

Buckley, P. J., Clegg, L. J., Cross, A. R., Liu, X., Voss, H., \& Zheng,
P. (2007). The determinants of Chinese outward foreign direct
investment. \emph{Journal of International Business Studies, 38}(4),
499--518. \url{https://doi.org/10.1057/palgrave.jibs.8400277}

Calabrese, G., \& Manello, A. (2018). Firm internationalization and
performance: Evidence for designing policies. \emph{Journal of Policy
Modeling, 40}(5), 1022--1039.
\url{https://doi.org/10.1016/j.jpolmod.2018.01.008}

Cameron, A. C., \& Miller, D. L. (2015). A practitioner's guide to
cluster-robust inference. \emph{Journal of Human Resources, 50}(2),
317--372. \url{https://doi.org/10.3368/jhr.50.2.317}

Cameron, A. C., \& Trivedi, P. K. (2005). \emph{Microeconometrics:
Methods and applications}. Cambridge University Press.

Cameron, A. C., Gelbach, J. B., \& Miller, D. L. (2008). Bootstrap-based
improvements for inference with clustered errors. \emph{Review of
Economics and Statistics, 90}(3), 414--427.
\url{https://doi.org/10.1162/rest.90.3.414}

Cantele, S., Vernizzi, S., \& Zardini, A. (2016). Does export affect the
performance of Italian SMEs? An empirical analysis. \emph{Economia e
Diritto del Terziario, 28}(1), 5--22.

Cao, X., Li, P., Huang, X., \& Fan, L. (2022). The dual mechanism of
social networks on the relationship between internationalization and
firm performance: Empirical evidence from China. \emph{PLoS ONE, 17}(2),
Article e0263674. \url{https://doi.org/10.1371/journal.pone.0263674}

Capar, N., \& Kotabe, M. (2003). The relationship between international
diversification and performance in service firms. \emph{Journal of
International Business Studies, 34}(4), 345--355.
\url{https://doi.org/10.1057/palgrave.jibs.8400036}

Carpenter, M. A., \& Sanders, W. G. (2004). The effects of top
management team pay and firm internationalization on MNC performance.
\emph{Journal of Management, 30}(4), 509--528.
\url{https://doi.org/10.1016/j.jm.2004.02.001}

Carr, J. C., Haggard, K. S., Hmieleski, K. M., \& Zahra, S. A. (2010). A
study of the moderating effects of firm age at internationalization on
firm survival and short-term growth. \emph{Strategic Entrepreneurship
Journal, 4}(2), 183--192. \url{https://doi.org/10.1002/sej.90}

Caves, R. E. (1971). International corporations: The industrial
economics of foreign investment. \emph{Economica, 38}(149), 1--27.

Cazzaniga, M., Jaumotte, F., Li, L., Melina, G., Panton, A. J.,
Pizzinelli, C., Rockall, E., \& Tavares, M. M. (2024). Gen-AI:
Artificial intelligence and the future of work. \emph{IMF Staff
Discussion Note}, SDN/2024/001.** International Monetary Fund.

Chao, M. C.-H., \& Kumar, V. (2010). The impact of institutional
distance on the international diversity-performance relationship.
\emph{Journal of World Business, 45}(1), 93--103.
\url{https://doi.org/10.1016/j.jwb.2009.04.005}

Chari, M. D. R., Devaraj, S., \& David, P. (2007). International
diversification and firm performance: Role of information technology
investments. \emph{Journal of World Business, 42}(2), 184--197.
\url{https://doi.org/10.1016/j.jwb.2007.02.007}

Chen, H., \& Hsu, C. (2010). Internationalization, resource allocation
and firm performance. \emph{Industrial Marketing Management, 39}(7),
1103--1110. \url{https://doi.org/10.1016/j.indmarman.2009.10.001}

Chen, J., \& Meng, Q. (2022). Institutional constraints and exporting of
emerging-market firms. \emph{Managerial and Decision Economics, 43}(5),
1706--1720.

Chen, S., \& Tan, H. (2012). Region effects in the
internationalization--performance relationship in Chinese firms.
\emph{Journal of World Business, 47}(1), 73--80.

Chiao, Y.-C., \& Yang, K.-P. (2011). Internationalization and
performance of Taiwanese SMEs: Moderating effects of firm age and
managerial characteristics. \emph{International Journal of
Entrepreneurship and Small Business, 14}(4), 496--518.
\url{https://doi.org/10.1504/IJESB.2011.043273}

Chiao, Y.-C., Yang, K.-P., \& Yu, C.-M. J. (2006). Performance,
internationalization, and firm-specific advantages of SMEs in a
newly-industrialized economy. \emph{Small Business Economics, 26}(5),
475--492. \url{https://doi.org/10.1007/s11187-005-5604-6}

Cho, H.-J., \& Lee, J. (2018). Firm performance and
internationalization: The moderation of family ownership.
\emph{Management Research Review, 41}(9), 1053--1071.
\url{https://doi.org/10.1108/MRR-06-2017-0175}

Chodorow-Reich, G., Gopinath, G., Mishra, P., \& Narayanan, A. (2020).
Cash and the economy: Evidence from India's demonetization.
\emph{Quarterly Journal of Economics, 135}(1), 57--103.
\url{https://doi.org/10.1093/qje/qjz027}

CIEM. (2023). \emph{Báo cáo Năng lực cạnh tranh và hiệu quả doanh nghiệp
Việt Nam 2023} {[}Central Institute for Economic Management{]}. Nhà xuất
bản Khoa học và Kỹ thuật.

Ciszewska-Mlinarič, M., \& Mlinarič, F. (2010). Internationalization and
performance of small and medium-sized enterprises. \emph{Organizacija,
43}(5), 175--185. \url{https://doi.org/10.2478/v10051-010-0017-y}

Cohen, J. (1988). \emph{Statistical power analysis for the behavioral
sciences} (2nd ed.). Lawrence Erlbaum.

Cohen, W. M., \& Levinthal, D. A. (1990). Absorptive capacity: A new
perspective on learning and innovation. \emph{Administrative Science
Quarterly, 35}(1), 128--152. \url{https://doi.org/10.2307/2393553}

Coltman, T., Devinney, T. M., Midgley, D. F., \& Venaik, S. (2008).
Formative versus reflective measurement models: Two applications of
formative measurement. \emph{Journal of Business Research, 61}(12),
1250--1262. \url{https://doi.org/10.1016/j.jbusres.2008.01.013}

Combs, J. G., Crook, T. R., \& Shook, C. L. (2005). The dimensionality
of organizational performance and its implications for strategic
management research. In D. J. Ketchen \& D. D. Bergh (Eds.),
\emph{Research methodology in strategy and management} (Vol. 2,
pp.~259--286). Emerald.
\url{https://doi.org/10.1016/S1479-8387(05)02011-4}

Combs, J. G., Ketchen, D. J., Crook, T. R., \& Roth, P. L. (2011).
Assessing cumulative evidence within ``macro'' research: Why
meta-analysis should be preferred over vote counting. \emph{Journal of
Management Studies, 48}(1), 178--197.
\url{https://doi.org/10.1111/j.1467-6486.2009.00899.x}

Contractor, F. J. (2012). Why do multinational firms exist? A theory
note about the effect of multinational expansion on performance and
recent methodological critiques. \emph{Global Strategy Journal, 2}(4),
318--331. \url{https://doi.org/10.1111/j.2042-5805.2012.01043.x}

Cooper, H. (2010). \emph{Research synthesis and meta-analysis: A
step-by-step approach} (4th ed.). Sage.

Cuervo-Cazurra, A., Ciravegna, L., Melgarejo, M., \& Lopez, L. (2017).
Home country uncertainty and the internationalization-performance
relationship: Building an uncertainty management capability.
\emph{Journal of World Business, 53}(2), 209--221.
\url{https://doi.org/10.1016/j.jwb.2017.11.002}

Cuervo-Cazurra, A., Ciravegna, L., Melgarejo, M., \& Lopez, L. (2018).
Home country uncertainty and the internationalization--performance
relationship. \emph{Journal of World Business, 53}(2), 209--221.

Cusolito, A. P., \& Maloney, W. F. (2018). \emph{Productivity revisited:
Shifting paradigms in analysis and policy}. World Bank.
\url{https://doi.org/10.1596/978-1-4648-1219-6}

Da Cunha, H. M., Amal, M., Floriani, D. E., \& Fleury, M. T. L. (2023).
Firm internationalization approaches and performance: The moderating
role of the home country's formal institutions. \emph{International
Journal of Emerging Markets}.
\url{https://doi.org/10.1108/ijoem-08-2021-1299}

Dawson, J. F. (2014). Moderation in management research: What, why,
when, and how. \emph{Journal of Business and Psychology, 29}(1), 1--19.
\url{https://doi.org/10.1007/s10869-013-9308-7}

Deaton, A. (1985). Panel data from a time series of cross-sections.
\emph{Journal of Econometrics, 30}(1--2), 109--126.
\url{https://doi.org/10.1016/0304-4076(85)90134-4} \emph{(Cơ sở repeated
cross-sections --- Phụ lục A.1, A.6.)}

Debicki, B. J., Miao, C., \& Qian, S. (2020). Internationalization and
family firm performance. \emph{Cross Cultural \& Strategic Management,
27}(1), 1--25. \url{https://doi.org/10.1108/ccsm-04-2019-0075}

Demir, B., \& Javorcik, B. (2018). Don't throw in the towel, throw in
trade credit! \emph{Journal of International Economics, 111}, 177--189.
\url{https://doi.org/10.1016/j.jinteco.2018.01.008}

Dikova, D., \& Veselova, A. (2020). Performance effects of
internationalization: Contingency theory analysis of Russian
internationalized firms. \emph{Management and Organization Review,
17}(1), 173--197. \url{https://doi.org/10.1017/mor.2020.39}

Do, T. H., \& Phan, A. T. (2025). Internationalization and firm
performance of firms in India: The role of top management. In M.
Bartekova (Ed.), \emph{International business research: Traditional and
creative approaches}. IntechOpen.
\url{https://doi.org/10.5772/intechopen.1011012}

Drori, N., Alessandri, T., Bart, Y., \& Herstein, R. (2024). The impact
of digitalization on internationalization from an internalization theory
lens. \emph{Long Range Planning, 57}(1), 102395.
\url{https://doi.org/10.1016/j.lrp.2023.102395}

Dunning, J. H. (1988). The eclectic paradigm of international
production: A restatement and some possible extensions. \emph{Journal of
International Business Studies, 19}(1), 1--31.**
\url{https://doi.org/10.1057/palgrave.jibs.8490372}

Duval, S., \& Tweedie, R. (2000). Trim and fill: A simple
funnel-plot--based method of testing and adjusting for publication bias
in meta-analysis. \emph{Biometrics, 56}(2), 455--463.
\url{https://doi.org/10.1111/j.0006-341X.2000.00455.x}

Eden, L., \& Nielsen, B. B. (2020). Research methods in international
business: The challenge of complexity. \emph{Journal of International
Business Studies, 51}(9), 1609--1620.
\url{https://doi.org/10.1057/s41267-020-00374-2}

Egger, M., Davey Smith, G., Schneider, M., \& Minder, C. (1997). Bias in
meta-analysis detected by a simple, graphical test. \emph{British
Medical Journal, 315}(7109), 629--634.
\url{https://doi.org/10.1136/bmj.315.7109.629}

Eisenhardt, K. M., \& Martin, J. A. (2000). Dynamic capabilities: What
are they? \emph{Strategic Management Journal, 21}(10--11), 1105--1121.
\url{https://doi.org/10.1002/smj.118}

Elango, B. (2006). An empirical analysis of the relationship between
international activity and firm performance. \emph{Multinational
Business Review, 14}(3), 21--45.
\url{https://doi.org/10.1108/1525383X200600015}

Elango, B. (2012). How industry dynamics influence the
internationalization--performance relationship: Evidence from
technology-intensive firms. \emph{Thunderbird International Business
Review, 54}(5), 653--665. \url{https://doi.org/10.1002/tie.21491}

Endo, N., \& Ozaki, T. (2011). International expansion and performance:
Evidence from Japanese manufacturing firms. \emph{Asian Business \&
Management, 10}(4), 515--535. \url{https://doi.org/10.1057/abm.2011.27}

Enterprise Singapore. (2025, November). \emph{SWITCH 2025 concludes 10th
anniversary milestone with record 25,000 attendees} {[}Press release{]}.
Enterprise Singapore.
\url{https://www.enterprisesg.gov.sg/resources/media-centre/media-releases/2025/november/mr05225_switch-2025-concludes-10th-anniversary-milestone-with-record-25000-attendees-strengthens-position-as-premiere-platform-for-deep-tech-innovation}

Espinosa-Méndez, C., \& Jara, M. (2021). International diversification
and performance in family firm: Exploring nonlinear relationships with
the governance structure in an emerging economy. \emph{Spanish Journal
of Finance and Accounting, 50}(4), 441--468.
\url{https://doi.org/10.1080/02102412.2021.1886453}

Falahat, M., Ramayah, T., Soto-Acosta, P., \& Lee, Y.-Y. (2020). SMEs
internationalization: The role of product innovation, market
intelligence, pricing and marketing communication capabilities as
drivers of SMEs' international performance. \emph{Technological
Forecasting and Social Change, 152}, Article 119908.
\url{https://doi.org/10.1016/j.techfore.2020.119908}

Fariborzi, H., Verbeke, A., \& Steel, P. (2025). Early
internationalization: A meta-analysis of antecedents, dimensions, and
performance. \emph{Journal of Management Studies}.
\url{https://doi.org/10.1111/joms.70008}

Feenstra, R. C., Inklaar, R., \& Timmer, M. P. (2015). The next
generation of the Penn World Table. \emph{American Economic Review,
105}(10), 3150--3182. \url{https://doi.org/10.1257/aer.20130954}
\emph{(Điều chỉnh PPP cho so sánh mức xuyên quốc gia --- Phụ lục A.5.)}

Feng, D., Chen, Q., Song, M., \& Cui, L. (2019). Relationship between
the degree of internationalization and performance in manufacturing
enterprises of the Yangtze River Delta region. \emph{Emerging Markets
Finance and Trade, 55}(7), 1455--1471.

Fernhaber, S. A., \& Li, D. (2013). International exposure through
network relationships: Implications for new venture
internationalization. \emph{Journal of Business Venturing, 28}(2),
316--334. \url{https://doi.org/10.1016/j.jbusvent.2012.05.002}

Filatotchev, I., Wei, L. Q., Sarala, R. M., Dick, P., \& Prescott, J. E.
(2020). Connecting eastern and western perspectives on management:
Translation of practices across organizations, institution and
geographies. \emph{Journal of Management Studies, 57}(1), 1--24.
\url{https://doi.org/10.1111/joms.12526}

Fortanier, F., Muller, A., \& Van Tulder, R. (2007).
Internationalization and performance: Accounting for (host) country
effects. \emph{Academy of Management Proceedings, 2007}(1), 1--6.
\url{https://doi.org/10.5465/ambpp.2007.26508474}

Freixanet, J., \& Rialp, J. (2021). Disentangling the relationship
between internationalization, incremental and radical innovation, and
firm performance. \emph{Global Strategy Journal}.
\url{https://doi.org/10.1002/gsj.1412}

Fu, Y. (2024). Enterprises' internationalization, R\&D investment and
enterprise performance. \emph{Finance Research Letters, 68}.
\url{https://doi.org/10.1016/j.frl.2024.105721}

Gallegos Mardones, J., \& Ibáñez, M. J. (2025). Women's participation in
top management business positions and the internationalization process:
Evidence for Chilean companies. \emph{SAGE Open, 15}(2), 1--17.**
\url{https://doi.org/10.1177/21582440251335079}

Ganvir, V., \& Dwivedi, M. (2017). International diversification and
performance of Indian firms: An empirical analysis. \emph{South Asian
Journal of Business Studies, 6}(3), 259--282.
\url{https://doi.org/10.1108/SAJBS-03-2016-0020}

Garbe, J. N., \& Richter, N. F. (2009). Causal analysis of the
internationalization and firm performance relationship based on neural
networks --- advocating the transnational structure. \emph{Journal of
International Management, 15}(4), 413--431.
\url{https://doi.org/10.1016/j.intman.2009.01.003}

Gaur, A. S., \& Kumar, V. (2009). International diversification,
business group affiliation and firm performance: Empirical evidence from
India. \emph{British Journal of Management, 20}(2), 172--186.
\url{https://doi.org/10.1111/j.1467-8551.2008.00558.x}

Geginat, C., \& Ramalho, R. (2015). \emph{Electricity connections and
firm performance in 183 countries} (Policy Research Working Paper
No.~7460). World Bank Group.

Glaum, M., \& Oesterle, M. (2007). 40 years of research on
internationalization and firm performance: More questions than answers?
\emph{Management International Review, 47}(3), 307--317.
\url{https://doi.org/10.1007/s11575-007-0018-0}

Gomes, L., \& Ramaswamy, K. (1999). An empirical examination of the form
of the relationship between multinationality and performance.
\emph{Journal of International Business Studies, 30}(1), 173--187.
\url{https://doi.org/10.1057/palgrave.jibs.8490065}

GovTech Singapore. (2024). \emph{Our digital government rankings}.
Government Technology Agency of Singapore.
\url{https://www.tech.gov.sg/about-us/our-achievements/our-digital-government-rankings/}

Grant, R. M. (1987). Multinationality and performance among British
manufacturing companies. \emph{Journal of International Business
Studies, 18}(3), 79--89.
\url{https://doi.org/10.1057/palgrave.jibs.8490413}

Graves, C., \& Shan, Y. G. (2014). An empirical analysis of the effect
of internationalization on the performance of unlisted family and
nonfamily firms in Australia. \emph{Family Business Review, 27}(2),
142--160. \url{https://doi.org/10.1177/0894486513510534}

Grego, M., Buckley, P. J., Munjal, S., Voss, H., \& Wang, E. Y. (2025).
Managers and internationalization decisions: An affect-enacted model.
\emph{Global Strategy Journal, 15}(3), 317--345.
\url{https://doi.org/10.1002/gsj.1516}

Haans, R. F. J., Pieters, C., \& He, Z.-L. (2016). Thinking about U:
Theorizing and testing U- and inverted U-shaped relationships in
strategy research. \emph{Strategic Management Journal, 37}(7),
1177--1195. \url{https://doi.org/10.1002/smj.2399}

Hair, J. F., Black, W. C., Babin, B. J., \& Anderson, R. E. (2019).
\emph{Multivariate data analysis} (8th ed.). Cengage.

Hall, P. A., \& Soskice, D. (2001). \emph{Varieties of capitalism: The
institutional foundations of comparative advantage}. Oxford University
Press. \url{https://doi.org/10.1093/0199247757.001.0001}

Hambrick, D. C. (2007). Upper echelons theory: An update. \emph{Academy
of Management Review, 32}(2), 334--343.
\url{https://doi.org/10.5465/amr.2007.24345254}

Hambrick, D. C., \& Mason, P. A. (1984). Upper echelons: The
organization as a reflection of its top managers. \emph{Academy of
Management Review, 9}(2), 193--206.
\url{https://doi.org/10.5465/amr.1984.4277628}

Hedges, L. V., \& Olkin, I. (1985). \emph{Statistical methods for
meta-analysis}. Academic Press.

Helpman, E., Melitz, M. J., \& Yeaple, S. R. (2004). Export versus FDI
with heterogeneous firms. \emph{American Economic Review, 94}(1),
300--316. \url{https://doi.org/10.1257/000282804322970814}

Hemmert, M., \& Jackson, K. (2016). Is there an East Asian model of MNC
internationalization? A comparative analysis of Japanese and Korean
firms. \emph{Asia Pacific Business Review, 22}(4), 567--594.
\url{https://doi.org/10.1080/13602381.2016.1168617}

Hennart, J.-F. (2007). The theoretical rationale for a
multinationality--performance relationship. \emph{Management
International Review, 47}(3), 423--452.
\url{https://doi.org/10.1007/s11575-007-0022-4}

Hessels, J., \& Terjesen, S. (2010). Resource dependency and
institutional theory perspectives on direct and indirect export choices.
\emph{Small Business Economics, 34}(2), 203--220.
\url{https://doi.org/10.1007/s11187-008-9156-4}

Higgins, J. P. T., Thompson, S. G., Deeks, J. J., \& Altman, D. G.
(2003). Measuring inconsistency in meta-analyses. \emph{BMJ, 327}(7414),
557--560. \url{https://doi.org/10.1136/bmj.327.7414.557}

Hilmersson, M. (2014). Small and medium-sized enterprise
internationalisation strategy and performance in times of market
turbulence. \emph{International Small Business Journal, 32}(4),
386--400. \url{https://doi.org/10.1177/0266242613485778}

Hitt, M. A., Hoskisson, R. E., \& Kim, H. (1997). International
diversification: Effects on innovation and firm performance in
product-diversified firms. \emph{Academy of Management Journal, 40}(4),
767--798. \url{https://doi.org/10.2307/256948}

Hojnik, J., Ruzzier, M., \& Konecnik Ruzzier, M. (2018).
Internationalisation and economic performance: The mediating role of
innovation. \emph{Economics Research-Ekonomska Istraživanja, 31}(1),
1770--1789. \url{https://doi.org/10.1080/1331677X.2018.1504673}

Hosseini, M., Brege, S., \& Nord, T. (2018). A combined focused industry
and company size investigation of the internationalization-performance
relationship: The case of small and medium-sized enterprises (SMEs)
within the Swedish wood manufacturing industry. \emph{Forest Policy and
Economics}. \url{https://doi.org/10.1016/j.forpol.2018.09.013}

Hsieh, C.-T., \& Klenow, P. J. (2009). Misallocation and manufacturing
TFP in China and India. \emph{Quarterly Journal of Economics, 124}(4),
1403--1448. \url{https://doi.org/10.1162/qjec.2009.124.4.1403}

Hsu, M.-L., \& Boggs, D. J. (2003). Internationalization and
performance: Traditional measures and their decomposition.
\emph{Multinational Business Review, 11}(3), 23--49.
\url{https://doi.org/10.1108/1525383X200300015}

Hsu, W. T., Chen, H. L., \& Cheng, C. Y. (2013). Internationalization
and firm performance of SMEs: The moderating effects of CEO attributes.
\emph{Journal of World Business, 48}(1), 1--12.
\url{https://doi.org/10.1016/j.jwb.2012.06.001}

Huang, L., \& Marciano, D. (2020). Interdependence relationship of
internationalization-performance in manufacturing firms listed in
Indonesia Stock Exchange and Chinese Stock Exchanges. In
\emph{Proceedings of the 17th International Symposium on Management
(INSYMA 2020)}. \url{https://doi.org/10.2991/aebmr.k.200127.085}

Hymer, S. H. (1976). \emph{The international operations of national
firms: A study of direct foreign investment}. MIT Press.

IMDA. (2025). \emph{Asia Tech x Singapore turns five} {[}Press
release{]}. Infocomm Media Development Authority.
\url{https://www.imda.gov.sg/resources/press-releases-factsheets-and-speeches/press-releases/2025/asia-tech-x-singapore-turns-five}

IMF. (2026, April). \emph{World Economic Outlook April 2026: Global
Economy in the Shadow of War}. International Monetary Fund.

International Finance Corporation. (2022, December). \emph{Mind the
gaps: Women in leadership in Viet Nam's banking sector}. IFC, World Bank
Group.

International Finance Corporation. (n.d.). \emph{Private Sector
Development Blueprint: Decoding the World Bank Enterprise Surveys}.
World Bank Group.

Jacquemin, A. P., \& Berry, C. H. (1979). Entropy measure of
diversification and corporate growth. \emph{Journal of Industrial
Economics, 27}(4), 359--369. \url{https://doi.org/10.2307/2097958}

Johanson, J., \& Vahlne, J.-E. (1977). The internationalization process
of the firm --- A model of knowledge development and increasing foreign
market commitments. \emph{Journal of International Business Studies,
8}(1), 23--32. \url{https://doi.org/10.1057/palgrave.jibs.8490676}

Juniati, S., \& Wahyuni, S. (2019). Internationalization and
performance: Evidence from Indonesian firms. \emph{International Journal
of Business and Society, 20}(3), 1134--1152.

Kafouros, M., Aliyev, M., Piperopoulos, P., Au, A. K. M., Ho, J. W. Y.,
\& Wong, S. Y. N. (2023). The role of institutional quality and industry
dynamism in explaining firm performance in emerging economies.
\emph{Global Strategy Journal, 14}(1), 56--83.**
\url{https://doi.org/10.1002/gsj.1479}

Kalubandi, S. C., Sharma, D., Edacherian, S., \& Karna, A. (2025). Role
of foreign market knowledge and home country institutions in
international entrepreneurship-performance relationship: A
meta-analysis. \emph{Asia Pacific Journal of Management}.
\url{https://doi.org/10.1007/s10490-025-10076-8}

Kano, L., Tsang, E. W. K., \& Yeung, H. W.-C. (2020). Global value
chains: A review of the multi-disciplinary literature. \emph{Journal of
International Business Studies, 51}(4), 577--622.
\url{https://doi.org/10.1057/s41267-020-00304-2}

Karna, A., Richter, A., \& Riesenkampff, E. (2016). Revisiting the role
of the environment in the capabilities--financial performance
relationship: A meta-analysis. \emph{Strategic Management Journal,
37}(6), 1154--1173. \url{https://doi.org/10.1002/smj.2379}

Katz, R., \& Callorda, F. (2018). \emph{The economic contribution of
broadband, digitization and ICT regulation}. ITU Telecommunication
Development Sector.

Kaufmann, D., Kraay, A., \& Mastruzzi, M. (2011). The Worldwide
Governance Indicators: Methodology and analytical issues. \emph{Hague
Journal on the Rule of Law, 3}(2), 220--246.
\url{https://doi.org/10.1017/S1876404511200046}

Kayaci, A. (2022). The relationship between internationalization level
and financial performance: A study with BİST companies. \emph{Cumhuriyet
Üniversitesi İktisadi ve İdari Bilimler Dergisi}.
\url{https://doi.org/10.37880/cumuiibf.1077104}

Khanna, T., \& Palepu, K. G. (2010). \emph{Winning in emerging markets:
A road map for strategy and execution}. Harvard Business Press.

Khatua, A., Mondal, A., \& Dhanda, S. (2024). The internationalization
and performance of INVs: Liability or learning advantages of newness?
\emph{Journal of International Entrepreneurship}. Advance online
publication. \url{https://doi.org/10.1007/s10843-023-00345-2}

Kim, G., Kumar, T., Ramalho, R., \& Russell, S. (2026, January).
\emph{Institutional capacity for policy implementation: An analytical
framework} (Policy Research Working Paper No.~11279). World Bank Group.

Kim, W. C., Hwang, P., \& Burgers, W. P. (1989). Global diversification
strategy and corporate profit performance. \emph{Strategic Management
Journal, 10}(1), 45--57. \url{https://doi.org/10.1002/smj.4250100105}

Kirca, A. H., Roth, K., Hult, G. T. M., \& Cavusgil, S. T. (2012). The
role of context in the multinationality--performance relationship: A
meta-analytic review. \emph{Global Strategy Journal, 2}(2), 108--121.
\url{https://doi.org/10.1111/j.2042-5805.2012.01028.x}

Knight, G. A., \& Cavusgil, S. T. (2004). Innovation, organizational
capabilities, and the born-global firm. \emph{Journal of International
Business Studies, 35}(2), 124--141.
\url{https://doi.org/10.1057/palgrave.jibs.8400071}

Kongkaew, W., Kraiwanit, T., \& Limna, P. (2021). Internationalization
and business performance of Thai SMEs. \emph{Journal of International
Business and Economy, 22}(2), 1--22.

Kraus, S., Breier, M., Lim, W. M., Dabić, M., Kumar, S., Kanbach, D.,
Mukherjee, D., Corvello, V., Piñeiro-Chousa, J., Liguori, E.,
Palacios-Marqués, D., Schiavone, F., Ferraris, A., Fernandes, C., \&
Ferreira, J. J. (2022). Literature reviews as independent studies:
Guidelines for academic practice. \emph{Review of Managerial Science,
16}(8), 2577--2595. \url{https://doi.org/10.1007/s11846-022-00588-8}

Kumar, N., \& Singh, N. (2008). Internationalization and performance of
Indian pharmaceutical firms. \emph{Vikalpa: The Journal for Decision
Makers, 33}(4), 13--24.

Kumar, V., \& Gaur, A. S. (2007). Internationalization and performance
of Indian firms. \emph{International Business: Research, Teaching, and
Practice, 1}(1), 1--20.

La Porta, R., \& Shleifer, A. (2014). Informality and development.
\emph{Journal of Economic Perspectives, 28}(3), 109--126.
\url{https://doi.org/10.1257/jep.28.3.109}

Lakens, D., Scheel, A. M., \& Isager, P. M. (2018). Equivalence testing
for psychological research: A tutorial. \emph{Advances in Methods and
Practices in Psychological Science, 1}(2), 259--269.
\url{https://doi.org/10.1177/2515245918770963}

Lall, S. (1992). Technological capabilities and industrialization.
\emph{World Development, 20}(2), 165--186.
\url{https://doi.org/10.1016/0305-750X(92)90097-F}

Landis, J. R., \& Koch, G. G. (1977). The measurement of observer
agreement for categorical data. \emph{Biometrics, 33}(1), 159--174.
\url{https://doi.org/10.2307/2529310}

Lee, C. Y. (2013). The effect of internationalization on business
performance: The case of Korean firms. \emph{Asian Review of Accounting,
21}(2), 115--133. \url{https://doi.org/10.1108/ARA-09-2012-0061}

Lee, J.-Y., Woo, C.-W., \& Shim, W.-S. (2019). The effect of
internationalization on performance of Korean firms: Moderating effect
of technological capability. \emph{International Journal of Trade,
Economics and Finance, 10}(3), 84--89.
\url{https://doi.org/10.18178/ijtef.2019.10.3.645}

Lee, S. M., Lim, S.-B., \& Pathak, R. D. (2014). Culture and
entrepreneurial orientation: A multi-country study. \emph{International
Entrepreneurship and Management Journal, 10}(1), 1--15.
\url{https://doi.org/10.1007/s11365-011-0218-8}

Leung, T., \& Sharma, P. (2021). Differences in the impact of R\&D
intensity and R\&D internationalization on firm performance: Mediating
role of innovation performance. \emph{Journal of Business Research,
131}, 81--91. \url{https://doi.org/10.1016/j.jbusres.2021.03.060}

Li, W., Li, C., \& Wei, G. (2022). The dual mechanism of social networks
on the relationship between internationalization and firm performance.
\emph{PLOS ONE, 17}(11), e0277421.

Li, Y., Liu, Y., \& Qian, C. (2022). Digitalization,
internationalization, and firm performance: Evidence from Chinese
manufacturing firms. \emph{International Business Review, 31}(3),
101953. \url{https://doi.org/10.1016/j.ibusrev.2021.101953}

Likitwongkajon, N., \& Vithessonthi, C. (2020). Do foreign investments
increase firm value and firm performance? Evidence from Japan.
\emph{Research in International Business and Finance, 51}, 101099.

Lin, W. T., Liu, Y., \& Cheng, K. (2011). The internationalization and
performance of a firm: Moderating effect of a firm's behavior.
\emph{Journal of International Management, 17}(1), 83--95.
\url{https://doi.org/10.1016/j.intman.2010.12.004}

Lind, J. T., \& Mehlum, H. (2010). With or without U? The appropriate
test for a U-shaped relationship. \emph{Oxford Bulletin of Economics and
Statistics, 72}(1), 109--118.
\url{https://doi.org/10.1111/j.1468-0084.2009.00569.x}

Little, R. J. A., \& Rubin, D. B. (2019). \emph{Statistical analysis
with missing data} (3rd ed.). John Wiley \& Sons.
\url{https://doi.org/10.1002/9781119482260} \emph{(Xử lý dữ liệu thiếu,
listwise/MAR --- Phụ lục A.9.)}

Liu, X., \& Zhang, H. (2024). How geopolitical risk affects firms'
internationalization performance: Evidence from China. \emph{Heliyon,
10}(6), e27512.

Long, J. S., \& Ervin, L. H. (2000). Using heteroscedasticity consistent
standard errors in the linear regression model. \emph{The American
Statistician, 54}(3), 217--224.
\url{https://doi.org/10.1080/00031305.2000.10474549}

Lu, J. W., \& Beamish, P. W. (2001). The internationalization and
performance of SMEs. \emph{Strategic Management Journal, 22}(6--7),
565--586.

Lu, J. W., \& Beamish, P. W. (2004). International diversification and
firm performance: The S-curve hypothesis. \emph{Academy of Management
Journal, 47}(4), 598--609. \url{https://doi.org/10.2307/20159604}

Luo, Y., \& Tung, R. L. (2007). International expansion of emerging
market enterprises: A springboard perspective. \emph{Journal of
International Business Studies, 38}(4), 481--498.
\url{https://doi.org/10.1057/palgrave.jibs.8400275}

MacKinnon, J. G., \& Webb, M. D. (2017). Wild bootstrap inference for
wildly different cluster sizes. \emph{Journal of Applied Econometrics,
32}(2), 233--254. \url{https://doi.org/10.1002/jae.2508}

MacKinnon, J. G., \& White, H. (1985). Some
heteroskedasticity-consistent covariance matrix estimators with improved
finite sample properties. \emph{Journal of Econometrics, 29}(3),
305--325. \url{https://doi.org/10.1016/0304-4076(85)90158-7}

Mahler, D. G., Serajuddin, U., Wadhwa, D., \& Yonzan, N. (2026, April).
\emph{The world is developing at its slowest pace in 75 years} (Policy
Research Working Paper No.~11350). World Bank Group.

Mahler, D., Wang, H., \& Weber, M. (2026, May). Inequalities in use of
and exposure to artificial intelligence. In \emph{Atlas of Global
Development 2026}. World Bank Data360. World Bank Group.
\url{https://data360.worldbank.org/en/int/atlas/artificial-intelligence}

Majocchi, A., \& Zucchella, A. (2003). Internationalization and
performance: Findings from a set of Italian SMEs. \emph{International
Small Business Journal, 21}(3), 249--268.
\url{https://doi.org/10.1177/0266242603021003001}

Malerba, F., \& Orsenigo, L. (1995). Schumpeterian patterns of
innovation. \emph{Cambridge Journal of Economics, 19}(1), 47--65.**
\url{https://doi.org/10.1093/oxfordjournals.cje.a035301}

Manotas, E. C., \& Gonzalez-Perez, M. A. (2020). Internationalization
and performance of small and medium-sized enterprises from emerging
economies. \emph{Competitiveness Review: An International Business
Journal, 31}(4), 642--662. \url{https://doi.org/10.1108/cr-03-2019-0028}

Manova, K. (2013). Credit constraints, heterogeneous firms, and
international trade. \emph{Review of Economic Studies, 80}(2), 711--744.
\url{https://doi.org/10.1093/restud/rds036}

Mar, T. T., Do, T. H., \& Phan, A. T. (2026). \emph{Digital adoption and
the internationalization--performance curve: Evidence from Singapore
firms} {[}Manuscript under review{]}. \emph{Management International
Review}.

Marano, V., Arregle, J.-L., Hitt, M. A., Spadafora, E., \& van Essen, M.
(2016). Home country institutions and the
internationalization--performance relationship: A meta-analytic review.
\emph{Journal of Management, 42}(5), 1075--1110.
\url{https://doi.org/10.1177/0149206315624963}

McDougall, P. P., \& Oviatt, B. M. (1996). New venture
internationalization, strategic change, and performance: A follow-up
study. \emph{Journal of Business Venturing, 11}(1), 23--40.
\url{https://doi.org/10.1016/0883-9026(95)00081-X}

Melitz, M. J. (2003). The impact of trade on intra-industry
reallocations and aggregate industry productivity. \emph{Econometrica,
71}(6), 1695--1725. \url{https://doi.org/10.1111/1468-0262.00467}

Meyer, K. E., van Witteloostuijn, A., \& Beugelsdijk, S. (2017). What's
in a p? Reassessing best practices for conducting and reporting
hypothesis-testing research. \emph{Journal of International Business
Studies, 48}(5), 535--551.**
\url{https://doi.org/10.1057/s41267-017-0078-8}

Miller, S. R., Lavie, D., \& Delios, A. (2016). International intensity,
diversity, and distance: Unpacking the internationalization-performance
relationship. \emph{International Business Review, 25}(4), 907--920.
\url{https://doi.org/10.1016/j.ibusrev.2015.12.001}

Mizik, N., \& Jacobson, R. (2003). Trading off between value creation
and value appropriation: The financial implications of shifts in
strategic emphasis. \emph{Journal of Marketing, 67}(1), 63--76.
\url{https://doi.org/10.1509/jmkg.67.1.63.18595}

Mohr, A., \& Batsakis, G. (2016). Internationalization speed and firm
performance: A study of the market-seeking expansion of retail MNEs.
\emph{Management International Review, 57}(2), 153--177.
\url{https://doi.org/10.1007/s11575-016-0284-7}

Mondal, A., Ray, S., \& Lahiri, S. (2022). Family ownership, family
management, and multinationality: Evidence from India. \emph{Journal of
Business Research, 138}, 347--359.
\url{https://doi.org/10.1016/j.jbusres.2021.09.017}

Murhadi, W. R., Setiawan, S. S., \& Ernawati, E. (2025). Governance and
internationalisation on firm performance: A case study in Indonesia and
Malaysia. \emph{Revista de Métodos Cuantitativos para la Economía y la
Empresa, 40}, 1--15.
\url{https://doi.org/10.46661/rev.metodoscuant.econ.empresa.10865}

Nambisan, S., Wright, M., \& Feldman, M. (2019). The digital
transformation of innovation and entrepreneurship: Progress, challenges
and key themes. \emph{Research Policy, 48}(8), 103773.
\url{https://doi.org/10.1016/j.respol.2019.03.018}

Nath, P., Kirca, A. H., \& Kim, S. (2021). A study of the
internationalization--performance relationship in global retailing.
\emph{Journal of International Marketing, 29}(1), 57--76.

Nguyen, Q., \& Kim, S. (2020). The multinationality and performance
relationship: Revisiting the literature and exploring the implications.
\emph{International Business Review, 29}(6), Article 101670.
\url{https://doi.org/10.1016/j.ibusrev.2020.101670}

Nielsen, S. (2010). Top management team internationalization and firm
performance: The mediating role of foreign market entry.
\emph{Management International Review, 50}(2), 185--206.
\url{https://doi.org/10.1007/s11575-010-0029-0}

North, D. C. (1990). \emph{Institutions, institutional change and
economic performance}. Cambridge University Press.

Nosi, C., Pucci, T., Silvestri, L., \& Zanni, L. (2017). Firm
performance and legitimacy in internationalization processes of Italian
SMEs. \emph{Sinergie Italian Journal of Management, 35}(102), 49--70.
\url{https://doi.org/10.7433/s102.2017.04}

OECD. (2026, April). \emph{FDI in Figures, April 2026}. Organisation for
Economic Co-operation and Development.

OECD. (2026, March). \emph{OECD Economic Outlook, Volume 2026 Issue 1}.
Organisation for Economic Co-operation and Development.

Olmos, M. F., \& Díez-Vial, I. (2015). Internationalization pathways and
the performance of SMEs. \emph{European Journal of Marketing, 49}(3/4),
420--443. \url{https://doi.org/10.1108/ejm-06-2012-0365}

Oster, E. (2019). Unobservable selection and coefficient stability:
Theory and evidence. \emph{Journal of Business \& Economic Statistics,
37}(2), 187--204. \url{https://doi.org/10.1080/07350015.2016.1227711}

Page, M. J., McKenzie, J. E., Bossuyt, P. M., Boutron, I., Hoffmann, T.
C., Mulrow, C. D., Shamseer, L., Tetzlaff, J. M., Akl, E. A., Brennan,
S. E., Chou, R., Glanville, J., Grimshaw, J. M., Hróbjartsson, A., Lalu,
M. M., Li, T., Loder, E. W., Mayo-Wilson, E., McDonald, S., \& Moher, D.
(2021). The PRISMA 2020 statement: An updated guideline for reporting
systematic reviews. \emph{BMJ, 372}, n71.
\url{https://doi.org/10.1136/bmj.n71}

Pan, Y., Verbeke, A., \& Yuan, W. (2021). CEO transformational
leadership and corporate entrepreneurship in China. \emph{Management and
Organization Review, 17}(1), 45--76.
\url{https://doi.org/10.1017/mor.2020.59}

Pangarkar, N. (2008). Internationalization and performance of small- and
medium-sized enterprises. \emph{Journal of World Business, 43}(4),
475--485.

Paternoster, R., Brame, R., Mazerolle, P., \& Piquero, A. (1998). Using
the correct statistical test for the equality of regression
coefficients. \emph{Criminology, 36}(4), 859--866.
\url{https://doi.org/10.1111/j.1745-9125.1998.tb01268.x}

Pattnaik, C., \& Elango, B. (2009). The impact of firm resources on the
internationalization and performance relationship: A study of Indian
manufacturing firms. \emph{Multinational Business Review, 17}(2),
69--90. \url{https://doi.org/10.1108/1525383X200900012}

Peng, M. W. (2003). Institutional transitions and strategic choices.
\emph{Academy of Management Review, 28}(2), 275--296.
\url{https://doi.org/10.5465/amr.2003.9416341}

Peng, M. W., \& Ilinitch, A. Y. (1998). Export intermediary firms: A
note on export development research. \emph{Journal of International
Business Studies, 29}(3), 609--620.
\url{https://doi.org/10.1057/palgrave.jibs.8490011}

Peng, M. W., Wang, D. Y. L., \& Jiang, Y. (2008). An institution-based
view of international business strategy: A focus on emerging economies.
\emph{Journal of International Business Studies, 39}(5), 920--936.
\url{https://doi.org/10.1057/palgrave.jibs.8400377}

Peterson, R. A., \& Brown, S. P. (2005). On the use of beta coefficients
in meta-analysis. \emph{Journal of Applied Psychology, 90}(1), 175--181.
\url{https://doi.org/10.1037/0021-9010.90.1.175}

Phan, A. T., Do, T. H., \& Phan, M. T. (2020b). The moderating effects
of managers' experience and gender on internationalization and firm
performance of manufacturing enterprises in Turkey. \emph{Accounting,
6}(7), 1209--1216. \url{https://doi.org/10.5267/j.ac.2020.9.006}

Phan, A. T., Nguyen, V. S., Dinh, V. T., Thai, V. H., Nguyen, C. T.,
Dao, Q. T., Nguyen, V. L., \& Do, T. H. (2021). The moderating effects
of the manager's characteristics and financial obstacles on
internationalization and firm performance in Poland. \emph{PalArch's
Journal of Archaeology of Egypt/Egyptology, 18}(4), 4063--4081.

Phan, A. T., Tran, T. M., \& Lu, T. M. (2020a). The existence of
W-shaped relationship between internationalization and firm performance:
An empirical study of enterprises in Cameroon. \emph{Decision Science
Letters, 9}(4), 573--580. \url{https://doi.org/10.5267/j.dsl.2020.7.001}

Pierce, J. R., \& Aguinis, H. (2013). The too-much-of-a-good-thing
effect in management. \emph{Journal of Management, 39}(2), 313--338.
\url{https://doi.org/10.1177/0149206311410060}

Pirlea, A. F., \& Corso, L. (2026, May). The unfinished digital
revolution: Expanding internet access. In \emph{Atlas of Global
Development 2026}. World Bank Data360. World Bank Group.
\url{https://data360.worldbank.org/en/int/atlas/internet-access}

Pirlea, A. F., Wadhwa, D., Mahler, D., Serajuddin, U., Welch, M., Thudt,
A., \& Lambrechts, M. (2026, May). Global progress. In \emph{Atlas of
Global Development 2026}. World Bank Data360. World Bank Group.
\url{https://data360.worldbank.org/en/int/atlas/global-progress}

Pisani, N., Garcia-Bernardo, J., \& Heemskerk, E. (2020). Does it pay to
be a multinational? A large-sample, cross-national replication assessing
the multinationality--performance relationship. \emph{Strategic
Management Journal, 41}(1), 152--172.
\url{https://doi.org/10.1002/smj.3087}

Polat, G., \& Mutlu, H. M. (2012). Internationalization performance of
construction companies: Evidence from Turkey. \emph{Engineering,
Construction and Architectural Management, 19}(2), 202--220.
\url{https://doi.org/10.1108/09699981211205849}

Post, C., \& Byron, K. (2015). Women on boards and firm financial
performance: A meta-analysis. \emph{Academy of Management Journal,
58}(5), 1546--1571. \url{https://doi.org/10.5465/amj.2013.0319}

Pouresmaeili, M., Salamzadeh, A., Salamzadeh, Y., \& Braga, V. (2018).
Internationalization of Iranian SMEs and its impact on firm performance.
\emph{International Journal of Business and Globalisation, 21}(3),
388--404. \url{https://doi.org/10.1504/IJBG.2018.095267}

Purkayastha, A., Karna, A., Sharma, S., \& Bhadra, D. (2024). Role of
resource investment management and strategic resource deployment
capabilities in internationalization-performance relationship.
\emph{Journal of International Management, 30}(3), Article 101122.
\url{https://doi.org/10.1016/j.intman.2024.101122}

Purkayastha, A., Kumar, V., \& Gupta, V. (2021). Emerging market
internationalizing firms: Learning through internationalization to
achieve entrepreneurial orientation. \emph{Journal of World Business,
56}(4), Article 101207. \url{https://doi.org/10.1016/j.jwb.2021.101207}

Purkayastha, S., \& Gupta, V. K. (2022). Business group affiliation and
internationalization-performance relationship: The moderating role of
product diversification. \emph{International Business Review, 31}(3),
Article 101968. \url{https://doi.org/10.1016/j.ibusrev.2021.101968}

Purkayastha, S., Manolova, T. S., \& Edelman, L. F. (2018). Business
group effects on the R\&D intensity-internationalization relationship:
Empirical evidence from India. \emph{Journal of World Business, 53}(5),
645--657. \url{https://doi.org/10.1016/j.jwb.2018.03.006}

Putri, W. H., \& Pan, L. (2022). The relationship between
internationalization and performance: Evidence from Indonesian
manufacturing firms. \emph{Cogent Business \& Management, 9}(1), Article
2079802. \url{https://doi.org/10.1080/23311975.2022.2079802}

Qian, G. (2002). Multinationality, product diversification, and
profitability of emerging US small- and medium-sized enterprises.
\emph{Journal of Business Venturing, 17}(6), 611--633.
\url{https://doi.org/10.1016/S0883-9026(01)00079-5}

Raudenbush, S. W., \& Bryk, A. S. (2002). \emph{Hierarchical linear
models: Applications and data analysis methods} (2nd ed.). Sage.

Riahi-Belkaoui, A. (1998). The effects of the degree of
internationalization on firm performance. \emph{International Business
Review, 7}(3), 315--321.
\url{https://doi.org/10.1016/S0969-5931(98)00013-4}

Richard, P. J., Devinney, T. M., Yip, G. S., \& Johnson, G. (2009).
Measuring organizational performance: Towards methodological best
practice. \emph{Journal of Management, 35}(3), 718--804.
\url{https://doi.org/10.1177/0149206308330560}

Rienda, L., \& Andreu, R. (2021). The role of family firms heterogeneity
on the internationalisation and performance relationship. \emph{European
Journal of Family Business, 11}(2).
\url{https://doi.org/10.24310/ejfbejfb.v11i2.10591}

Rosenthal, R. (1991). \emph{Meta-analytic procedures for social
research} (rev. ed.). Sage.

Rosenthal, R. (1994). Parametric measures of effect size. In H. Cooper
\& L. V. Hedges (Eds.), \emph{The handbook of research synthesis}
(pp.~231--244). Sage.

Ruigrok, W., \& Wagner, H. (2003). Internationalization and performance:
An organizational learning perspective. \emph{Management International
Review, 43}(1), 63--83.

Ruigrok, W., Amann, W., \& Wagner, H. (2007). The
internationalization-performance relationship at Swiss firms: A test of
the S-shape and extreme degrees of internationalization.
\emph{Management International Review, 47}(3), 349--368.
\url{https://doi.org/10.1007/s11575-007-0015-4}

Ruzzier, M., \& Ruzzier, M. K. (2014). On the relationship between firm
size, resources, age at entry and internationalization: The case of
Slovenian SMEs. \emph{Journal of Business Economics and Management,
16}(1), 52--73. \url{https://doi.org/10.3846/16111699.2012.745812}

Sachs, J. D., \& Warner, A. M. (2001). The curse of natural resources.
\emph{European Economic Review, 45}(4--6), 827--838.
\url{https://doi.org/10.1016/S0014-2921(01)00125-8}

Sahay, R., von Allmen, U. E., Lahreche, A., Khera, P., Ogawa, S.,
Bazarbash, M., \& Beaton, K. (2020). \emph{The promise of fintech:
Financial inclusion in the post COVID-19 era} (IMF Departmental Paper
No.~DP/2020/09). International Monetary Fund.

Sambharya, R. B. (1996). Foreign experience of top management teams and
international diversification strategies of U.S. multinational
corporations. \emph{Strategic Management Journal, 17}(9), 739--746.
\url{https://doi.org/10.1002/(SICI)1097-0266(199611)17:9}\textless739::AID-SMJ846\textgreater3.0.CO;2-K

Santangelo, G. D., \& Verbeke, A. (2022). Actionable guidelines to
improve `theory-related' contributions to international business
research. \emph{Journal of International Business Studies, 53}(9),
1843--1855. \url{https://doi.org/10.1057/s41267-022-00567-x}

Sasabuchi, S. (1980). A test of a multivariate normal mean with
composite hypotheses determined by linear inequalities.
\emph{Biometrika, 67}(2), 429--439.
\url{https://doi.org/10.1093/biomet/67.2.429}

Schmuck, A., Lagerström, K., \& Sallis, J. F. (2022). Turning the
tables: The relationship between performance and multinationality.
\emph{Management International Review, 62}(1), 3--26.
\url{https://doi.org/10.1007/s11575-022-00464-3}

Scott, W. R. (1995). \emph{Institutions and organizations}. Sage.

Shaver, J. M. (2020). Causal identification through a cumulative body of
research in the study of strategy and organizations. \emph{Journal of
Management, 46}(7), 1244--1256.
\url{https://doi.org/10.1177/0149206319846272}

Siddharthan, N. S., \& Lall, S. (1982). The recent growth of the largest
US multinationals. \emph{Oxford Bulletin of Economics and Statistics,
44}(1), 1--13.
\url{https://doi.org/10.1111/j.1468-0084.1982.mp44001001.x}

Sikdar, S., \& Mukhopadhyay, A. (2026). Productivity gains in Indian
manufacturing: Insights from formal and informal sectors. \emph{Asian
Development Review, 43}(1), Article 12.** Asian Development Bank.

Singla, C., \& George, R. (2013). Internationalization and performance:
A contextual analysis of Indian firms. \emph{Journal of Business
Research, 66}(12), 2500--2506.
\url{https://doi.org/10.1016/j.jbusres.2013.05.041}

Solon, G., Haider, S. J., \& Wooldridge, J. M. (2015). What are we
weighting for? \emph{Journal of Human Resources, 50}(2), 301--316.
\url{https://doi.org/10.3368/jhr.50.2.301}

Stallkamp, M., \& Schotter, A. P. J. (2021). Platforms without borders?
\emph{Global Strategy Journal, 11}(1), 58--80.

Stanley, T. D., \& Doucouliagos, H. (2014). Meta-regression
approximations to reduce publication selection bias. \emph{Research
Synthesis Methods, 5}(1), 60--78.
\url{https://doi.org/10.1002/jrsm.1095}

StartupBlink. (2025). \emph{Global Startup Ecosystem Index 2025:
Singapore profile}. StartupBlink Research.
\url{https://www.startupblink.com/startup-ecosystem/singapore}

StartupBlink. (2026). \emph{Innovators Business Environment Index 2026}.
StartupBlink Research.
\url{https://www.startupblink.com/blog/innovators-business-environment-index/}

Sullivan, D. (1994). Measuring the degree of internationalization of a
firm. \emph{Journal of International Business Studies, 25}(2), 325--342.
\url{https://doi.org/10.1057/palgrave.jibs.8490203}

Sun, W., Price, J. M., \& Ding, Y. (2019). The longitudinal effects of
internationalization on firm performance: The moderating role of
marketing capability. \emph{Journal of Business Research, 95}, 326--337.
\url{https://doi.org/10.1016/j.jbusres.2018.08.014}

Sutton, A. J., \& Higgins, J. P. T. (2008). Recent developments in
meta-analysis. \emph{Statistics in Medicine, 27}(5), 625--650.
\url{https://doi.org/10.1002/sim.2982}

Suurmond, R., van Rhee, H., \& Hak, T. (2017). Introduction, comparison,
and validation of Meta-Essentials: A free and simple tool for
meta-analysis. \emph{Research Synthesis Methods, 8}(4), 537--553.
\url{https://doi.org/10.1002/jrsm.1260}

Tallman, S., \& Li, J. (1996). Effects of international diversity and
product diversity on the performance of multinational firms.
\emph{Academy of Management Journal, 39}(1), 179--196.
\url{https://doi.org/10.2307/256635}

Tashman, P., Marano, V., \& Kostova, T. (2019). Walking the walk or
talking the talk? Corporate social responsibility and
internationalization intensity in the Fortune Global 500. \emph{Journal
of International Business Studies, 50}(2), 253--261.
\url{https://doi.org/10.1057/s41267-018-0197-4}

Teece, D. J. (1986). Profiting from technological innovation:
Implications for integration, collaboration, licensing and public
policy. \emph{Research Policy, 15}(6), 285--305.
\url{https://doi.org/10.1016/0048-7333(86)90027-2}

Teece, D. J. (2007). Explicating dynamic capabilities: The nature and
microfoundations of (sustainable) enterprise performance.
\emph{Strategic Management Journal, 28}(13), 1319--1350.
\url{https://doi.org/10.1002/smj.640}

Teece, D. J., Pisano, G., \& Shuen, A. (1997). Dynamic capabilities and
strategic management. \emph{Strategic Management Journal, 18}(7),
509--533. \url{https://doi.org/10.1002/smj.4250180702}

Thomas, D. E. (2006). International diversification and firm performance
in Mexican firms: A curvilinear relationship. \emph{Journal of Business
Research, 59}(4), 501--507.
\url{https://doi.org/10.1016/j.jbusres.2005.09.008}

Thomas, D. E., \& Eden, L. (2004). What is the shape of the
multinationality-performance relationship? \emph{Multinational Business
Review, 12}(1), 89--110. \url{https://doi.org/10.1108/1525383X200400005}

Tongurai, J., \& Vithessonthi, C. (2022). Learning, foreign operations
and operating performance. \emph{Global Finance Journal, 54}, Article
100721. \url{https://doi.org/10.1016/j.gfj.2022.100721}

Torraco, R. J. (2005). Writing integrative literature reviews:
Guidelines and examples. \emph{Human Resource Development Review, 4}(3),
356--367. \url{https://doi.org/10.1177/1534484305278283}

Tran, T. B., \& Pham, T. H. (2024). Tác động của FDI đến năng suất doanh
nghiệp Việt Nam: Phân tích đa thế hệ điều tra. \emph{Tạp chí Kinh tế và
Phát triển, 311}, 24--38.

Tsao, S.-M., \& Chen, G.-Z. (2012). The impact of internationalization
on performance and innovation: The moderating effects of ownership
concentration. \emph{Asia Pacific Journal of Management, 29}(3),
617--642. \url{https://doi.org/10.1007/s10490-010-9240-5}

Tsao, S.-M., \& Lien, W.-H. (2013). Family management and
internationalization: The impact on firm performance and innovation.
\emph{Management International Review, 53}(2), 175--193.
\url{https://doi.org/10.1007/s11575-012-0135-1}

Tsionas, M. G., \& Tzeremes, N. G. (2021). The degree of
internationalization and firm productivity: Empirical evidence from
large multinationals. \emph{British Journal of Management, 32}(4),
1032--1052. \url{https://doi.org/10.1111/1467-8551.12560}

UNCTAD. (2023). \emph{World Investment Report 2023: Investing in
sustainable energy for all}. United Nations Conference on Trade and
Development.

Upadhyayula, R. S., Bhatt, D., \& Gaur, A. S. (2016). International
expansion and performance of Indian firms: Moderating role of board
structure. \emph{Journal of Indian Business Research, 8}(2), 130--149.
\url{https://doi.org/10.1108/JIBR-10-2015-0087}

Vahlne, J.-E. (2020). Development of the Uppsala model of
internationalization process: From internationalization to evolution.
\emph{Global Strategy Journal, 10}(2), 239--250.
\url{https://doi.org/10.1002/gsj.1383}

Vélez-Calle, A., Sánchez-Henriquez, F., \& MacLennan, A. F. (2018). Peer
influence and internationalization intentions of Latin American firms.
\emph{International Journal of Emerging Markets, 13}(1), 95--118.
\url{https://doi.org/10.1108/IJoEM-11-2016-0298}

Verbeek, M. (2008). Pseudo-panels and repeated cross-sections. In L.
Mátyás \& P. Sevestre (Eds.), \emph{The econometrics of panel data} (3rd
ed., pp.~369--383). Springer.
\url{https://doi.org/10.1007/978-3-540-75892-1_11} \emph{(Khung lát cắt
ngang lặp lại --- Phụ lục A.1, A.6.)}

Verbeke, A., Eschmann, D., \& Gugler, P. (2025). Asymmetries in
firm-level globalization: The case of Swiss multinational enterprises.
\emph{Management International Review, 65}(2), 305--328.
\url{https://doi.org/10.1007/s11575-025-00568-6}

Verbeke, A., Oh, C. H., \& Jain, R. (2025). What is the future of
regional multinational enterprises? \emph{International Business
Review}. \url{https://doi.org/10.1016/j.ibusrev.2025.102442}

Verhoef, P. C., Broekhuizen, T., Bart, Y., Bhattacharya, A., Dong, J.
Q., Fabian, N., \& Haenlein, M. (2021). Digital transformation: A
multidisciplinary reflection and research agenda. \emph{Journal of
Business Research, 122}, 889--901.
\url{https://doi.org/10.1016/j.jbusres.2019.09.022}

Vernon, R. (1966). International investment and international trade in
the product cycle. \emph{The Quarterly Journal of Economics, 80}(2),
190--207.

Viechtbauer, W. (2010). Conducting meta-analyses in R with the metafor
package. \emph{Journal of Statistical Software, 36}(3), 1--48.
\url{https://doi.org/10.18637/jss.v036.i03}

Vilas-Boas da Silva, J., \& González Pérez, M. A. (2007).
Internationalization and performance: Evidence from Spanish firms.
\emph{Cuadernos de Economía y Dirección de la Empresa, 10}(33), 31--62.

Vithessonthi, C., \& Racela, O. C. (2016). Short- and long-run effects
of internationalization and R\&D intensity on firm performance.
\emph{Journal of Multinational Financial Management, 34}, 28--45.
\url{https://doi.org/10.1016/j.mulfin.2015.12.001}

Wach, K. (2021). The evolution of the Uppsala model: Towards
non-linearity of internationalization of firms. \emph{International
Entrepreneurship Review, 7}(2), 7--19.
\url{https://doi.org/10.15678/IER.2021.0702.01}

Wagner, J. (2007). Exports and productivity: A survey of the evidence
from firm-level data. \emph{The World Economy, 30}(1), 60--82.
\url{https://doi.org/10.1111/j.1467-9701.2007.00872.x}

Wang, G., Zhang, H., Xia, B., Wu, G., \& Han, Y. (2020). Relationship
between internationalization and financial performance: Evidence from
ENR-listed Chinese firms. \emph{Journal of Management in Engineering,
36}(2), Article 04019044.
\url{https://doi.org/10.1061/(asce)me.1943-5479.0000736}

Wang, L., Huang, Y., \& Hong, Z. (2024). Digitalization as a
double-edged sword: A deep learning analysis of risk management in
Chinese banks. \emph{International Review of Financial Analysis, 94},
103249. \url{https://doi.org/10.1016/j.irfa.2024.103249}

Wei, S., \& Lin, L. (2021). The relationship between extent of
internationalization and firm performance (Taiwan 1992--2017).
\emph{Mathematical Problems in Engineering}.
\url{https://doi.org/10.1155/2021/5528972}

Williamson, O. E. (1985). \emph{The economic institutions of
capitalism}. Free Press.

Williamson, O. E. (2000). The new institutional economics: Taking stock,
looking ahead. \emph{Journal of Economic Literature, 38}(3), 595--613.**
\url{https://doi.org/10.1257/jel.38.3.595}

WIPO. (2024). \emph{Global Innovation Index 2024: Unlocking the promise
of social entrepreneurship}. World Intellectual Property Organization.

Wooldridge, J. M. (2010). \emph{Econometric analysis of cross section
and panel data} (2nd ed.). MIT Press.

World Bank Enterprise Surveys. (2025). \emph{Enterprise Surveys: What
businesses experience}. The World Bank.
\url{https://www.enterprisesurveys.org/en/enterprisesurveys}

World Bank. (2013). \emph{China Enterprise Survey 2012} {[}Data file{]}.
World Bank Enterprise Surveys. \url{https://www.enterprisesurveys.org}

World Bank. (2019). \emph{Understanding the questionnaire: Enterprise
Survey questionnaire manual}. World Bank Enterprise Surveys.

World Bank. (2023). \emph{Vietnam 2023 Enterprise Surveys country
profile}. World Bank Enterprise Surveys.

World Bank. (2024). \emph{Mongolia Economic Update: Towards diversified
growth}. World Bank.

World Bank. (2024, December). \emph{Tonga 2024 Enterprise Surveys
country profile}. World Bank Enterprise Surveys.

World Bank. (2024, December). \emph{Viet Nam economy snapshot: Justice}
(Prosperity Data360). World Bank Group.
\url{https://www.worldbank.org/ext/en/country/vietnam}

World Bank. (2025, December). \emph{Kiribati 2025 Enterprise Surveys
country profile}. World Bank Enterprise Surveys.

World Bank. (2025, July 1). \emph{Hướng dẫn học tập: Hệ thống phân loại
quốc gia và khảo sát doanh nghiệp của Ngân hàng Thế giới (FY26)}
{[}Learning guide: World Bank country classification system FY26{]}.
World Bank.

World Bank. (2025a, March). \emph{Digital public infrastructure and
development: A World Bank Group approach}. World Bank Group.

World Bank. (2025b, September). \emph{Viet Nam economic update:
Nurturing high-tech talents} (Taking Stock, September 2025). World Bank
Group.

World Bank. (2025c, October). \emph{Viet Nam economy snapshot: Finance,
competitiveness \& innovation} (Prosperity Data360). World Bank Group.
\url{https://www.worldbank.org/ext/en/country/vietnam}

World Bank. (2025d, December). \emph{Korea global digital knowledge
center}. World Bank Group.

World Bank. (2026). \emph{Indicator descriptions}. World Bank Enterprise
Surveys.

World Bank. (2026a). \emph{ID4D \& G2Px annual report 2023: Putting
people at the center of digital public infrastructure}. World Bank
Group.

World Bank. (2026b, April). \emph{Viet Nam macro monitoring: April
2026}. World Bank Group, Fiscal Policy and Growth Viet Nam Team.

World Bank. (2026c). Frontier market economies: Promise, performance,
and prospects. In \emph{Global Economic Prospects, January 2026}
(Chapter 4). World Bank Group.
\url{https://doi.org/10.1596/978-1-4648-2148-6}

World Bank. (2026d, April). \emph{Commodity markets outlook, April
2026}. World Bank Group.

World Bank. (2026e). \emph{Public institutions in the age of AI:
Emerging practices} (Prosperity Insight Series). World Bank Group,
GovTech and Public Sector Innovation.

World Bank. (n.d.). \emph{Tài liệu tóm tắt: Các chỉ số khảo sát doanh
nghiệp Enterprise Surveys} {[}Summary document: Enterprise Survey
indicators{]}. World Bank.** \url{https://www.enterprisesurveys.org}

Wu, J., Wood, G., \& Khan, Z. (2022). Internationalization and firm
performance: Evidence from a meta-analysis. \emph{International Business
Review, 31}(2), 101920.

Xiao, S. S., Jeong, I., Moon, J. J., Chung, C. C., \& Chung, J. (2013).
Internationalization and performance of firms in China. \emph{Journal of
International Management, 19}(2), 118--137.

Xiong, Z., Sun, C., \& Zhao, M. (2024). Internationalization speed and
corporate long-term performance: Exploring the inverted U-shaped
relationship. \emph{Finance Research Letters}.
\url{https://doi.org/10.1016/j.frl.2024.106247}

Xu, J. (2024). Rethinking institutional arbitrage: De jure exposure and
de facto enforcement. \emph{Global Strategy Journal, 14}(4), 754--799.**
\url{https://doi.org/10.1002/gsj.1510}

Yang et al.~(2025) JBR: hiện chưa verify được volume/pages --- cần kiểm
tra Crossref hoặc loại bỏ.

Yang, B., Bai, W., Chen, Y., \& Rong, K. (2025). Internationalization of
digital firms: A systematic review and research agenda. \emph{Journal of
Business Research}.

Yip, G. S., Biscarri, J. G., \& Monti, J. A. (2000). The role of the
internationalization process in the performance of newly
internationalizing firms. \emph{Journal of International Marketing,
8}(3), 10--35. \url{https://doi.org/10.1509/jimk.8.3.10.19635}

Zaheer, S. (1995). Overcoming the liability of foreignness.
\emph{Academy of Management Journal, 38}(2), 341--363.
\url{https://doi.org/10.2307/256411}

Zeng, S. X., Xie, X. M., Tam, C. M., \& Wan, T. W. (2009). Relationships
between business factors and performance in internationalization: An
empirical study in China. \emph{Management Decision, 47}(2), 308--329.
\url{https://doi.org/10.1108/00251740910938939}

Zhang, Y., \& Wei, J. (2022). The impact of the internationalization of
China's new retail industry on corporate performance. \emph{PLOS ONE,
17}(2), e0262891.

Zhang, Y., Qiu, Z., Park, D., \& Tian, S. (2026). The role of artificial
intelligence in finance: A selective review and implications for Asia's
financial stability. In \emph{Asia Bond Monitor} (March 2026, Special
Section, pp.~38--43). Asian Development Bank.

Zhou, L., \& Wu, A. (2014). Earliness of internationalization and
performance outcomes: Exploring the moderating effects of venture age
and international commitment. \emph{Journal of World Business, 49}(1),
132--142. \url{https://doi.org/10.1016/j.jwb.2013.04.003}

Zhou, X. (2018). Effects of foreign direct investment on the efficiency
of Chinese firms. \emph{Applied Economics, 50}(2), 135--148.
\url{https://doi.org/10.1080/00036846.2017.1321839}

\newpage

\hypertarget{phux1ee5-lux1ee5c-a-quy-truxecnh-hux1ee3p-nhux1ea5t-vuxe0-huxe0i-huxf2a-huxf3a-dux1eef-liux1ec7u-doanh-nghiux1ec7p-ux111a-quux1ed1c-gia-wbes}{%
\chapter{PHỤ LỤC A: QUY TRÌNH HỢP NHẤT VÀ HÀI HÒA HÓA DỮ LIỆU DOANH
NGHIỆP ĐA QUỐC GIA
(WBES)}\label{phux1ee5-lux1ee5c-a-quy-truxecnh-hux1ee3p-nhux1ea5t-vuxe0-huxe0i-huxf2a-huxf3a-dux1eef-liux1ec7u-doanh-nghiux1ec7p-ux111a-quux1ed1c-gia-wbes}}

\begin{quote}
Phụ lục này trình bày chi tiết, có luận chứng khoa học, quy trình hợp
nhất (pooling) và hài hòa hóa (harmonisation) các bộ khảo sát doanh
nghiệp đơn lẻ theo quốc gia--năm thành một bộ dữ liệu cấp doanh nghiệp
thống nhất phục vụ phân tích thực nghiệm đa bối cảnh của luận án (đặc
biệt Nghiên cứu thành phần P7). Đây là \textbf{một cấu phần phương pháp
luận chính thức} của luận án: nó bảo đảm tính minh bạch, khả năng tái
lập (replicability) và khả năng phản biện trước hội đồng. Toàn bộ quy
trình được mã hóa thành tập lệnh tái lập
(\texttt{scripts/build\_pooled\_dataset.py},
\texttt{scripts/wbes\_canon.py}) và báo cáo dòng dữ liệu kèm theo
(\texttt{data\_wbes/analysis/pooled\_dataflow.csv}).
\end{quote}

\hypertarget{a.1-mux1ee5c-ux111uxedch-vuxe0-vux1ecb-truxed-trong-thiux1ebft-kux1ebf-nghiuxean-cux1ee9u}{%
\section{A.1 Mục đích và vị trí trong thiết kế nghiên
cứu}\label{a.1-mux1ee5c-ux111uxedch-vuxe0-vux1ecb-truxed-trong-thiux1ebft-kux1ebf-nghiuxean-cux1ee9u}}

Luận án kiểm định mối quan hệ quốc tế hóa--hiệu quả (I--P) trên một mẫu
doanh nghiệp trải rộng nhiều nền kinh tế và nhiều đợt khảo sát. Bản chất
dữ liệu là \textbf{lát cắt ngang lặp lại} (repeated cross-sections) chứ
không phải bảng cân bằng (balanced panel): mỗi đợt khảo sát WBES rút một
mẫu doanh nghiệp \textbf{độc lập} từ khung mẫu của nền kinh tế tại thời
điểm đó, không theo dõi cùng một doanh nghiệp qua thời gian (Deaton,
1985; Verbeek, 2008). Đặc tính này quy định toàn bộ chiến lược hợp nhất
và nhận dạng (identification) trình bày dưới đây.

Phụ lục A bổ sung cho Chương 3 (Mục 3.3) ở cấp độ \textbf{thao tác chi
tiết} (operational detail): trong khi Mục 3.3 nêu nguồn dữ liệu và định
nghĩa biến, Phụ lục A tài liệu hóa từng bước chuyển hóa từ 119 tệp khảo
sát thô đến mẫu phân tích cuối cùng, cùng luận chứng cho mỗi quyết định
lọc và hài hòa.

\hypertarget{a.2-nguux1ed3n-dux1eef-liux1ec7u-vuxe0-ux111ux1a1n-vux1ecb-phuxe2n-tuxedch}{%
\section{A.2 Nguồn dữ liệu và đơn vị phân
tích}\label{a.2-nguux1ed3n-dux1eef-liux1ec7u-vuxe0-ux111ux1a1n-vux1ecb-phuxe2n-tuxedch}}

Dữ liệu sơ cấp là \textbf{World Bank Enterprise Surveys (WBES)}, bộ khảo
sát doanh nghiệp cấp cơ sở (establishment-level) do Nhóm Ngân hàng Thế
giới thực hiện theo một bộ công cụ toàn cầu chuẩn hóa, áp dụng
\textbf{chọn mẫu ngẫu nhiên phân tầng} (stratified random sampling) theo
ba chiều: ngành (sản xuất/dịch vụ theo ISIC), quy mô lao động
(nhỏ/vừa/lớn) và vùng địa lý trong mỗi nền kinh tế (World Bank
Enterprise Surveys, 2025; World Bank, 2019). Thiết kế phân tầng này bảo
đảm tính đại diện cho khu vực tư nhân chính thức phi nông nghiệp của mỗi
nền kinh tế và cho phép so sánh giữa các nước nhờ bộ câu hỏi cốt lõi
đồng nhất (Aterido et al., 2011). Đơn vị phân tích là \textbf{cơ sở kinh
doanh} (establishment); biến phụ thuộc, biến độc lập và biến điều tiết
đều được đo ở cấp này.

Một ràng buộc then chốt về tính so sánh: WBES chỉ áp dụng \textbf{phương
pháp luận toàn cầu thống nhất từ năm 2006} trở đi; các đợt trước 2006
dùng bộ công cụ và khung chọn mẫu không tương thích (World Bank, 2019).
Do đó luận án giới hạn ở các đợt khảo sát \textbf{năm ≥ 2006} để bảo
toàn tính so sánh giữa các đợt và các nước (xem A.3, Bước S1).

Tiêu chí ``năm ≥ 2006'' là \textbf{điều kiện cần nhưng chưa đủ}: một số
nền kinh tế vẫn còn các đợt khảo sát sau 2006 chạy trên \textbf{bộ công
cụ khu vực/quốc gia cũ} (ví dụ bộ PICS khu vực MENA hay các bảng hỏi
Đông Dương đời đầu) với hệ thống mã biến riêng. Tiêu chí vận hành để một
đợt khảo sát được nhận vào khung so sánh là \textbf{sự hiện diện của các
biến lõi đã hài hòa} của bộ công cụ Standardized toàn cầu, cụ thể là
\texttt{d2} (doanh thu), \texttt{l1} (lao động),
\texttt{d3b}/\texttt{d3c} (xuất khẩu trực tiếp/gián tiếp). Ba đợt khảo
sát đáp ứng năm ≥ 2006 nhưng \textbf{không} chứa các biến lõi này,
Lebanon 2009 (bộ PICS MENA, mã \texttt{q\#\#\#\#}), Cambodia 2007 và
Cambodia 2013 (bảng hỏi đời đầu), được loại theo đúng nguyên tắc này,
song song với việc loại các bộ công cụ phi chuẩn (Informal, ISBS, ISES,
Micro, panel). Cần phân biệt: việc loại trừ này là do \textbf{bất tương
thích bộ công cụ}, không phải do thiếu hay hỏng dữ liệu; các đợt
Standardized của chính những nền này vẫn được giữ (Lebanon 2013 và 2019;
Cambodia 2016 và 2023), và các đợt 2006--2012 của nền khác dùng bộ công
cụ Standardized (ví dụ Jordan 2006) đều được nhận.

\hypertarget{a.3-quy-truxecnh-hux1ee3p-nhux1ea5t-suxe1u-bux1b0ux1edbc-s1s6-vuxe0-duxf2ng-dux1eef-liux1ec7u}{%
\section{A.3 Quy trình hợp nhất sáu bước (S1--S6) và dòng dữ
liệu}\label{a.3-quy-truxecnh-hux1ee3p-nhux1ea5t-suxe1u-bux1b0ux1edbc-s1s6-vuxe0-duxf2ng-dux1eef-liux1ec7u}}

Quy trình được thiết kế theo nguyên tắc minh bạch kiểu PRISMA cho dữ
liệu thứ cấp: mỗi bước ghi rõ tiêu chí và số quan sát còn lại, cho phép
tái dựng và kiểm tra (Page et al., 2021).

\begin{longtable}[]{@{}
  >{\raggedright\arraybackslash}p{(\columnwidth - 6\tabcolsep) * \real{0.2500}}
  >{\raggedright\arraybackslash}p{(\columnwidth - 6\tabcolsep) * \real{0.2500}}
  >{\raggedright\arraybackslash}p{(\columnwidth - 6\tabcolsep) * \real{0.2500}}
  >{\raggedright\arraybackslash}p{(\columnwidth - 6\tabcolsep) * \real{0.2500}}@{}}
\toprule\noalign{}
\begin{minipage}[b]{\linewidth}\raggedright
Bước
\end{minipage} & \begin{minipage}[b]{\linewidth}\raggedright
Thao tác
\end{minipage} & \begin{minipage}[b]{\linewidth}\raggedright
Tiêu chí
\end{minipage} & \begin{minipage}[b]{\linewidth}\raggedright
Kết quả (mẫu gộp chính thức)
\end{minipage} \\
\midrule\noalign{}
\endhead
\bottomrule\noalign{}
\endlastfoot
\textbf{S1} & Thu thập \& lọc loại khảo sát & Giữ lát cắt ngang doanh
nghiệp tư nhân chuẩn; \textbf{loại} bảng panel, Informal Sector, ISBS,
ISES, Micro, và khảo sát theo dõi TGS; giữ \textbf{năm ≥ 2006} & --- \\
\textbf{S2} & Khử trùng lặp & Một khảo sát cho mỗi cặp (nền kinh tế ×
năm) & --- \\
\textbf{S3} & Hài hòa hóa biến & Ánh xạ trường câu hỏi dị biệt về một
lược đồ chung (A.4) & --- \\
\textbf{S4} & Xử lý mã thiếu & Mã hóa mã không phản hồi WBES (−9
\emph{Don't Know}, −8 \emph{Refusal}, −7 \emph{N/A}) thành khuyết; kiểm
tra miền giá trị & --- \\
\textbf{S5} & Phân tầng thể chế ICRV & Gán chế độ thể chế; giới hạn
khung phân tích \textbf{50 nền} Á--Thái Bình Dương (gồm Nhật Bản 2025) &
Khung mô tả 50 nền = \textbf{88.869} DN / 103 cặp \\
\textbf{S6} & Mẫu phân tích listwise & Loại quan sát thiếu biến phụ
thuộc hoặc biến giải thích trọng tâm & Mẫu phân tích mô tả =
\textbf{88.869} (50 nền); N hồi quy M2 = \textbf{81.022}; mô hình đầy đủ
M5 = \textbf{79.080} \\
\end{longtable}

\emph{Nguồn số chính thức: tệp mẫu gộp của luận án
(\texttt{data\_wbes/p7/p7\_pooled\_clean.csv}), 106.765 bản ghi doanh
nghiệp--năm từ các đợt khảo sát giai đoạn 2006--2026.}

\textbf{Hình A.1.} \emph{Sơ đồ dòng dữ liệu hợp nhất WBES (kiểu
PRISMA).}

\begin{Shaded}
\begin{Highlighting}[]
\NormalTok{flowchart TD}
\NormalTok{    A["119 tệp khảo sát WBES\textless{}br/\textgreater{}quốc gia–năm (thô)"] {-}{-}\textgreater{}|"S1 · lọc cross{-}section chuẩn;\textless{}br/\textgreater{}loại panel/Informal/ISBS/ISES/Micro/TGS; năm ≥ 2006"| B["Lát cắt ngang chuẩn ≥ 2006"]}
\NormalTok{    B {-}{-}\textgreater{}|"S2 · khử trùng lặp"| C["Một khảo sát / cặp (nền × năm)"]}
\NormalTok{    C {-}{-}\textgreater{}|"S3 · hài hòa hóa biến\textless{}br/\textgreater{}(LP, FSTS, TCI/DAI, controls)"| D["Lược đồ biến chung"]}
\NormalTok{    D {-}{-}\textgreater{}|"S4 · mã thiếu −9/−8/−7 thành khuyết;\textless{}br/\textgreater{}kiểm tra miền giá trị"| E["106.765 bản ghi doanh nghiệp–năm"]}
\NormalTok{    E {-}{-}\textgreater{}|"S5 · gán chế độ thể chế ICRV\textless{}br/\textgreater{}(đủ chỉ số WGI)"| F["96.415 DN phân loại · 52 nền"]}
\NormalTok{    F {-}{-}\textgreater{}|"loại Comoros / Cyprus / Türkiye\textless{}br/\textgreater{}(ngoài Á–Thái Bình Dương); thêm Nhật Bản 2025"| G["88.869 DN, 50 nền\textless{}br/\textgreater{}(KHUNG PHÂN TÍCH)"]}
\NormalTok{    G {-}{-}\textgreater{}|"S6 · listwise tập biến trọng tâm"| H["81.022, N hồi quy M2"]}
\NormalTok{    H {-}{-}\textgreater{}|"+ đầy đủ biến kiểm soát"| I["79.080, M5 (hiệu ứng cố định hai chiều)"]}
\end{Highlighting}
\end{Shaded}

\textbf{Lưu ý tái lập (replication note).} Kho dữ liệu thô tái dựng được
công khai trong repo (\texttt{data\_wbes/raw\_dta/}, 106 sóng chuẩn từ
2006 của 50 nền, gồm Nhật Bản 2025) là \textbf{tập con} của bộ dữ liệu
đầy đủ mà nhóm tác giả đã lắp ráp (một số sóng lịch sử chưa được tải
lại). Tập lệnh \texttt{scripts/build\_pooled\_dataset.py} tái lập
\textbf{quy trình} S1--S6 trên kho hiện có và cho ra dòng dữ liệu kiểm
chứng; các cỡ mẫu \textbf{chính thức} của luận án trên khung 50 nền
(88.869 mẫu mô tả / 81.022 N hồi quy M2 / 79.080 M5) được ước lượng trực
tiếp từ kho dữ liệu thô và là con số được trích dẫn trong Chương 3--4.
Sai khác giữa hai lớp được công bố minh bạch để hội đồng kiểm chứng.

\hypertarget{a.4-huxe0i-huxf2a-huxf3a-biux1ebfn-sux1ed1-bux1b0ux1edbc-s3}{%
\section{A.4 Hài hòa hóa biến số (Bước
S3)}\label{a.4-huxe0i-huxf2a-huxf3a-biux1ebfn-sux1ed1-bux1b0ux1edbc-s3}}

Thách thức cốt lõi của hợp nhất đa quốc gia là \textbf{dị biệt tên
trường và mã hóa} giữa các đợt khảo sát (ví dụ, một số đợt tách xuất
khẩu trực tiếp/gián tiếp, một số gộp). Luận án xây dựng một \textbf{lược
đồ biến chung} ánh xạ về các module cốt lõi bất biến của bộ công cụ WBES
toàn cầu (World Bank, 2019). Các biến trọng tâm và quy tắc xây dựng:

\begin{itemize}
\tightlist
\item
  \textbf{Năng suất lao động (biến phụ thuộc chính):}
  \(\ln(\text{LP}) = \ln(\text{doanh thu hàng năm} / \text{số lao động toàn thời gian})\),
  từ trường \texttt{d2} (tổng doanh thu năm tài chính gần nhất) và
  \texttt{l1} (lao động thường xuyên cuối kỳ). Lựa chọn năng suất lao
  động làm thước đo hiệu quả là phù hợp với dữ liệu WBES vốn không có
  giá thị trường (Combs et al., 2005; Richard et al., 2009).
\item
  \textbf{Cường độ quốc tế hóa (FSTS):} tổng tỷ trọng xuất khẩu trực
  tiếp và gián tiếp trên doanh thu, từ \texttt{d3b} + \texttt{d3c}, giới
  hạn miền {[}0, 100{]}. Doanh nghiệp không xuất khẩu nhận giá trị FSTS
  = 0 (giá trị hợp lệ, không phải khuyết).
\item
  \textbf{Năng lực công nghệ (TCI) và chấp nhận số (DAI):} xây dựng từ
  các chỉ báo nhị phân (chứng nhận chất lượng quốc tế \texttt{b8}; công
  nghệ nước ngoài cấp phép; website \texttt{c22b}; email; thanh toán
  điện tử), tách biệt theo lập luận thuần khiết cấu trúc (construct
  purity) ở Mục 2.1.3 và Coltman et al.~(2008).
\item
  \textbf{Biến kiểm soát:} tuổi doanh nghiệp (từ \texttt{b5}), quy mô
  (log lao động), sở hữu nước ngoài, đặc điểm nhà quản trị, ngành ISIC.
\end{itemize}

\hypertarget{a.5-xux1eed-luxfd-tuxednh-so-suxe1nh-tiux1ec1n-tux1ec7-vux1ea5n-ux111ux1ec1-vuxe0-giux1ea3i-phuxe1p}{%
\section{A.5 Xử lý tính so sánh tiền tệ: vấn đề và giải
pháp}\label{a.5-xux1eed-luxfd-tuxednh-so-suxe1nh-tiux1ec1n-tux1ec7-vux1ea5n-ux111ux1ec1-vuxe0-giux1ea3i-phuxe1p}}

Một mối đe dọa hiệu lực (validity threat) đặc thù của mẫu gộp đa quốc
gia là \textbf{năng suất lao động thô được biểu thị bằng nội tệ khác
nhau} (ví dụ doanh thu/lao động của Việt Nam tính bằng VND, của
Singapore bằng SGD), khiến \textbf{mức} năng suất \textbf{không so sánh
trực tiếp} giữa các nước, một dạng ``ảo ảnh tiền tệ'' (currency
artefact). Luận án xử lý vấn đề này theo hai lớp bổ trợ:

\begin{enumerate}
\def\labelenumi{\arabic{enumi}.}
\tightlist
\item
  \textbf{Chuẩn hóa within country--year:} năng suất lao động được chuẩn
  hóa z (trừ trung bình, chia độ lệch chuẩn) \textbf{trong từng cặp nền
  kinh tế × năm} trước khi gộp, đưa mọi quan sát về cùng thang vị thế
  tương đối trong phân phối nội bộ nước--năm. Phép biến đổi này loại bỏ
  khác biệt đơn vị tiền tệ và mặt bằng giá, đồng thời bảo toàn cấu trúc
  biến thiên bên trong mỗi nước (nơi nhận dạng quan hệ I--P diễn ra).
\item
  \textbf{Hiệu ứng cố định hai chiều (two-way fixed effects):} mọi mô
  hình hồi quy P7 đưa vào \textbf{hiệu ứng cố định nền kinh tế} và
  \textbf{hiệu ứng cố định năm}. Hiệu ứng cố định nền hấp thụ trọn vẹn
  mọi khác biệt \textbf{bất biến theo nước} (gồm đơn vị tiền tệ, mặt
  bằng giá, thể chế nền, văn hóa), do đó hệ số quan hệ I--P được nhận
  dạng \textbf{từ biến thiên bên trong nước} (within-country variation),
  miễn nhiễm với ảo ảnh tiền tệ ở mức (Wooldridge, 2010).
\end{enumerate}

Hệ quả phương pháp luận quan trọng: vì lý do trên, \textbf{các so sánh
mức năng suất thô giữa các nước bị loại khỏi suy diễn}; các so sánh giữa
nhóm thể chế ICRV ở Chương 4 chỉ dùng \textbf{độ phân tán} (đã chuẩn
hóa/PPP--winsorize theo country--year) và \textbf{gradient điểm uốn},
những đại lượng bất biến với đơn vị tiền tệ (xem Mục 4.1.1). Khi cần so
sánh \textbf{mức} xuyên quốc gia, luận án dùng tỷ suất không thứ nguyên
(ví dụ ROS) hoặc điều chỉnh sức mua tương đương (PPP) theo Bảng Penn
World (Feenstra et al., 2015).

\hypertarget{a.6-chiux1ebfn-lux1b0ux1ee3c-hux1ee3p-nhux1ea5t-kinh-tux1ebf-lux1b0ux1ee3ng-dux1eef-liux1ec7u-gux1ed9p-luxe1t-cux1eaft-ngang-lux1eb7p-lux1ea1i}{%
\section{A.6 Chiến lược hợp nhất kinh tế lượng (dữ liệu gộp lát cắt
ngang lặp
lại)}\label{a.6-chiux1ebfn-lux1b0ux1ee3c-hux1ee3p-nhux1ea5t-kinh-tux1ebf-lux1b0ux1ee3ng-dux1eef-liux1ec7u-gux1ed9p-luxe1t-cux1eaft-ngang-lux1eb7p-lux1ea1i}}

Vì dữ liệu là lát cắt ngang lặp lại chứ không phải panel, ước lượng tuân
theo khung \textbf{pooled cross-section với hiệu ứng cố định} thay vì mô
hình panel động (Deaton, 1985; Wooldridge, 2010). Cụ thể, mô hình lõi
của P7 có dạng bậc hai để nắm bắt quan hệ chữ U ngược (H1):

\[\ln(\text{LP})_{ijt} = \beta_1 \text{FSTS}_{ijt} + \beta_2 \text{FSTS}^2_{ijt} + \mathbf{X}_{ijt}\boldsymbol{\gamma} + \alpha_j + \tau_t + \varepsilon_{ijt}\]

trong đó \(i\) là doanh nghiệp, \(j\) là nền kinh tế, \(t\) là năm khảo
sát; \(\alpha_j\) và \(\tau_t\) lần lượt là hiệu ứng cố định nền và năm;
\(\mathbf{X}\) là vector biến kiểm soát. Điểm uốn (turning point) được
tính \(\text{TP} = -\beta_1/(2\beta_2)\) và kiểm định bằng thủ tục
Lind--Mehlum (2010) cùng Sasabuchi để xác nhận tồn tại cực trị trong
miền dữ liệu. Sai số chuẩn được hiệu chỉnh phương sai thay đổi
(heteroskedasticity-robust) và phân cụm theo nền kinh tế nhằm phản ánh
tương quan nội cụm của thiết kế khảo sát phân tầng (Cameron \& Trivedi,
2005; White, 1980). Bậc hai (quadratic) được ưu tiên hơn bậc ba (cubic)
khi khối lượng phân phối ở vùng FSTS cao không đủ để nhận dạng điểm uốn
bậc ba một cách chính xác.

\hypertarget{a.7-phuxe2n-tux1ea7ng-thux1ec3-chux1ebf-icrv-bux1b0ux1edbc-s5}{%
\section{A.7 Phân tầng thể chế ICRV (Bước
S5)}\label{a.7-phuxe2n-tux1ea7ng-thux1ec3-chux1ebf-icrv-bux1b0ux1edbc-s5}}

Mỗi nền kinh tế được gán vào một trong sáu chế độ thể chế (Institutional
Context Regime Variation, ICRV) dựa trên tổ hợp các Chỉ số Quản trị Toàn
cầu (Worldwide Governance Indicators, WGI) của Kaufmann et al.~(2011),
mức phát triển kinh tế và đặc trưng cấu trúc thể chế, mở rộng từ khung
khoảng trống thể chế của Khanna và Palepu (2010). Việc gán nhóm đòi hỏi
đủ chỉ số WGI; các doanh nghiệp ở nước--năm thiếu chỉ số thể chế không
được phân loại (chênh lệch giữa 106.765 bản ghi thô và 96.415 bản ghi
phân loại được ở Bước S5). Khung phân tích cuối cùng giới hạn ở
\textbf{50 nền kinh tế châu Á và Thái Bình Dương} (gồm Nhật Bản 2025),
loại ba nền ngoài khu vực (Comoros, Cyprus, Türkiye) khỏi mẫu phân tích
chính trong khi vẫn giữ chúng ở mẫu gộp phân loại 52 nền cho mục đích độ
vững (riêng Comoros được dùng trong kiểm định độ vững mở rộng của P8).

\hypertarget{a.8-trux1ecdng-sux1ed1-mux1eabu-vuxe0-tuxednh-ux111ux1ea1i-diux1ec7n}{%
\section{A.8 Trọng số mẫu và tính đại
diện}\label{a.8-trux1ecdng-sux1ed1-mux1eabu-vuxe0-tuxednh-ux111ux1ea1i-diux1ec7n}}

Vì thiết kế WBES là phân tầng với xác suất chọn không đều, các thống kê
\textbf{mô tả cấp quần thể} (population-level) cần dùng trọng số chọn
mẫu để bảo đảm tính đại diện. Tuy nhiên, đối với \textbf{suy diễn hồi
quy có hiệu ứng cố định}, luận án theo khuyến nghị của Solon, Haider và
Wooldridge (2015): khi mô hình đã kiểm soát các chiều phân tầng (quy mô,
ngành qua biến kiểm soát; nền qua hiệu ứng cố định), ước lượng OLS không
trọng số cho hệ số nhất quán và thường hiệu quả hơn; trọng số được dùng
trong phân tích độ vững để kiểm tra tính ổn định. Cách tiếp cận này cân
bằng giữa tính đại diện và hiệu quả ước lượng một cách có luận chứng.

\hypertarget{a.9-xux1eed-luxfd-dux1eef-liux1ec7u-thiux1ebfu-bux1b0ux1edbc-s4-s6}{%
\section{A.9 Xử lý dữ liệu thiếu (Bước S4 \&
S6)}\label{a.9-xux1eed-luxfd-dux1eef-liux1ec7u-thiux1ebfu-bux1b0ux1edbc-s4-s6}}

Mã không phản hồi của WBES (−9 \emph{Don't Know}, −8 \emph{Refusal}, −7
\emph{Not Applicable}) được mã hóa nhất quán thành giá trị khuyết trước
mọi phép tính. Mẫu phân tích áp dụng \textbf{xóa theo danh sách}
(listwise deletion) trên tập biến trọng tâm, khiến N giảm dần khi thêm
biến kiểm soát (rõ nhất ở biến sở hữu nước ngoài, khuyết
\textasciitilde50\% ở các nền nhỏ và sóng sớm). Giả định nền tảng là dữ
liệu khuyết ngẫu nhiên có điều kiện (missing at random) trên các hiệp
biến quan sát được; luận án báo cáo minh bạch sự suy giảm N qua các đặc
tả và kiểm tra tính ổn định của điểm uốn để bảo đảm kết quả không bị chi
phối bởi cơ chế khuyết (Little \& Rubin, 2019).

\hypertarget{a.10-khux1ea3-nux103ng-tuxe1i-lux1eadp-vuxe0-hux1ea1n-chux1ebf}{%
\section{A.10 Khả năng tái lập và hạn
chế}\label{a.10-khux1ea3-nux103ng-tuxe1i-lux1eadp-vuxe0-hux1ea1n-chux1ebf}}

\textbf{Tái lập.} Toàn bộ quy trình được mã hóa xác định (deterministic)
và công khai: \texttt{scripts/wbes\_canon.py} (chuẩn hóa định danh
nước--năm, quy tắc lọc), \texttt{scripts/build\_pooled\_dataset.py} (hợp
nhất + hài hòa + dòng dữ liệu), với báo cáo
\texttt{data\_wbes/analysis/pooled\_dataflow.csv}. Việc xử lý do thiếu
Stata được thực hiện bằng Python (pandas/numpy), một giải pháp thay thế
hợp lệ cho phân tích kinh tế lượng lát cắt ngang gộp khi không có
license Stata.

\textbf{Hạn chế.} (i) Dữ liệu lát cắt ngang lặp lại không cho phép theo
dõi quỹ đạo doanh nghiệp, hạn chế suy diễn nhân quả động; (ii) năng suất
lao động là proxy cho hiệu quả, không phải TFPR trực tiếp; (iii) đo
lường DAI bị giới hạn ở proxy website cấp nền tảng do ràng buộc dữ liệu;
(iv) kho thô tái lập hiện là tập con của bộ lắp ráp đầy đủ. Các hạn chế
này được nêu nhất quán với phần hạn chế ở Chương 5.

\begin{center}\rule{0.5\linewidth}{0.5pt}\end{center}

\emph{Liên kết tập lệnh và dữ liệu (phục vụ tái lập):}

\begin{itemize}
\tightlist
\item
  \texttt{scripts/build\_pooled\_dataset.py}
\item
  \texttt{scripts/wbes\_canon.py}
\item
  \texttt{scripts/cd1\_descriptives\_pipeline.py}
\item
  \texttt{data\_wbes/analysis/pooled\_dataflow.csv}
\item
  \texttt{DATA\_UPDATE\_MANIFEST.md}
\item
  \texttt{CANONICAL\_RECONCILIATION.md}
\end{itemize}

\newpage

\end{document}
